% Options for packages loaded elsewhere
% Options for packages loaded elsewhere
\PassOptionsToPackage{unicode}{hyperref}
\PassOptionsToPackage{hyphens}{url}
\PassOptionsToPackage{dvipsnames,svgnames,x11names}{xcolor}
%
\documentclass[
  japanese,
  a4paper,
  DIV=11,
  numbers=noendperiod]{scrreprt}
\usepackage{xcolor}
\usepackage[margin=25mm]{geometry}
\usepackage{amsmath,amssymb}
\setcounter{secnumdepth}{-\maxdimen} % remove section numbering
\usepackage{iftex}
\ifPDFTeX
  \usepackage[T1]{fontenc}
  \usepackage[utf8]{inputenc}
  \usepackage{textcomp} % provide euro and other symbols
\else % if luatex or xetex
  \usepackage{unicode-math} % this also loads fontspec
  \defaultfontfeatures{Scale=MatchLowercase}
  \defaultfontfeatures[\rmfamily]{Ligatures=TeX,Scale=1}
\fi
\usepackage{lmodern}
\ifPDFTeX\else
  % xetex/luatex font selection
\fi
% Use upquote if available, for straight quotes in verbatim environments
\IfFileExists{upquote.sty}{\usepackage{upquote}}{}
\IfFileExists{microtype.sty}{% use microtype if available
  \usepackage[]{microtype}
  \UseMicrotypeSet[protrusion]{basicmath} % disable protrusion for tt fonts
}{}
\makeatletter
\@ifundefined{KOMAClassName}{% if non-KOMA class
  \IfFileExists{parskip.sty}{%
    \usepackage{parskip}
  }{% else
    \setlength{\parindent}{0pt}
    \setlength{\parskip}{6pt plus 2pt minus 1pt}}
}{% if KOMA class
  \KOMAoptions{parskip=half}}
\makeatother
% Make \paragraph and \subparagraph free-standing
\makeatletter
\ifx\paragraph\undefined\else
  \let\oldparagraph\paragraph
  \renewcommand{\paragraph}{
    \@ifstar
      \xxxParagraphStar
      \xxxParagraphNoStar
  }
  \newcommand{\xxxParagraphStar}[1]{\oldparagraph*{#1}\mbox{}}
  \newcommand{\xxxParagraphNoStar}[1]{\oldparagraph{#1}\mbox{}}
\fi
\ifx\subparagraph\undefined\else
  \let\oldsubparagraph\subparagraph
  \renewcommand{\subparagraph}{
    \@ifstar
      \xxxSubParagraphStar
      \xxxSubParagraphNoStar
  }
  \newcommand{\xxxSubParagraphStar}[1]{\oldsubparagraph*{#1}\mbox{}}
  \newcommand{\xxxSubParagraphNoStar}[1]{\oldsubparagraph{#1}\mbox{}}
\fi
\makeatother

\usepackage{color}
\usepackage{fancyvrb}
\newcommand{\VerbBar}{|}
\newcommand{\VERB}{\Verb[commandchars=\\\{\}]}
\DefineVerbatimEnvironment{Highlighting}{Verbatim}{commandchars=\\\{\}}
% Add ',fontsize=\small' for more characters per line
\usepackage{framed}
\definecolor{shadecolor}{RGB}{241,243,245}
\newenvironment{Shaded}{\begin{snugshade}}{\end{snugshade}}
\newcommand{\AlertTok}[1]{\textcolor[rgb]{0.68,0.00,0.00}{#1}}
\newcommand{\AnnotationTok}[1]{\textcolor[rgb]{0.37,0.37,0.37}{#1}}
\newcommand{\AttributeTok}[1]{\textcolor[rgb]{0.40,0.45,0.13}{#1}}
\newcommand{\BaseNTok}[1]{\textcolor[rgb]{0.68,0.00,0.00}{#1}}
\newcommand{\BuiltInTok}[1]{\textcolor[rgb]{0.00,0.23,0.31}{#1}}
\newcommand{\CharTok}[1]{\textcolor[rgb]{0.13,0.47,0.30}{#1}}
\newcommand{\CommentTok}[1]{\textcolor[rgb]{0.37,0.37,0.37}{#1}}
\newcommand{\CommentVarTok}[1]{\textcolor[rgb]{0.37,0.37,0.37}{\textit{#1}}}
\newcommand{\ConstantTok}[1]{\textcolor[rgb]{0.56,0.35,0.01}{#1}}
\newcommand{\ControlFlowTok}[1]{\textcolor[rgb]{0.00,0.23,0.31}{\textbf{#1}}}
\newcommand{\DataTypeTok}[1]{\textcolor[rgb]{0.68,0.00,0.00}{#1}}
\newcommand{\DecValTok}[1]{\textcolor[rgb]{0.68,0.00,0.00}{#1}}
\newcommand{\DocumentationTok}[1]{\textcolor[rgb]{0.37,0.37,0.37}{\textit{#1}}}
\newcommand{\ErrorTok}[1]{\textcolor[rgb]{0.68,0.00,0.00}{#1}}
\newcommand{\ExtensionTok}[1]{\textcolor[rgb]{0.00,0.23,0.31}{#1}}
\newcommand{\FloatTok}[1]{\textcolor[rgb]{0.68,0.00,0.00}{#1}}
\newcommand{\FunctionTok}[1]{\textcolor[rgb]{0.28,0.35,0.67}{#1}}
\newcommand{\ImportTok}[1]{\textcolor[rgb]{0.00,0.46,0.62}{#1}}
\newcommand{\InformationTok}[1]{\textcolor[rgb]{0.37,0.37,0.37}{#1}}
\newcommand{\KeywordTok}[1]{\textcolor[rgb]{0.00,0.23,0.31}{\textbf{#1}}}
\newcommand{\NormalTok}[1]{\textcolor[rgb]{0.00,0.23,0.31}{#1}}
\newcommand{\OperatorTok}[1]{\textcolor[rgb]{0.37,0.37,0.37}{#1}}
\newcommand{\OtherTok}[1]{\textcolor[rgb]{0.00,0.23,0.31}{#1}}
\newcommand{\PreprocessorTok}[1]{\textcolor[rgb]{0.68,0.00,0.00}{#1}}
\newcommand{\RegionMarkerTok}[1]{\textcolor[rgb]{0.00,0.23,0.31}{#1}}
\newcommand{\SpecialCharTok}[1]{\textcolor[rgb]{0.37,0.37,0.37}{#1}}
\newcommand{\SpecialStringTok}[1]{\textcolor[rgb]{0.13,0.47,0.30}{#1}}
\newcommand{\StringTok}[1]{\textcolor[rgb]{0.13,0.47,0.30}{#1}}
\newcommand{\VariableTok}[1]{\textcolor[rgb]{0.07,0.07,0.07}{#1}}
\newcommand{\VerbatimStringTok}[1]{\textcolor[rgb]{0.13,0.47,0.30}{#1}}
\newcommand{\WarningTok}[1]{\textcolor[rgb]{0.37,0.37,0.37}{\textit{#1}}}

\usepackage{longtable,booktabs,array}
\newcounter{none} % for unnumbered tables
\usepackage{calc} % for calculating minipage widths
% Correct order of tables after \paragraph or \subparagraph
\usepackage{etoolbox}
\makeatletter
\patchcmd\longtable{\par}{\if@noskipsec\mbox{}\fi\par}{}{}
\makeatother
% Allow footnotes in longtable head/foot
\IfFileExists{footnotehyper.sty}{\usepackage{footnotehyper}}{\usepackage{footnote}}
\makesavenoteenv{longtable}
\usepackage{graphicx}
\makeatletter
\newsavebox\pandoc@box
\newcommand*\pandocbounded[1]{% scales image to fit in text height/width
  \sbox\pandoc@box{#1}%
  \Gscale@div\@tempa{\textheight}{\dimexpr\ht\pandoc@box+\dp\pandoc@box\relax}%
  \Gscale@div\@tempb{\linewidth}{\wd\pandoc@box}%
  \ifdim\@tempb\p@<\@tempa\p@\let\@tempa\@tempb\fi% select the smaller of both
  \ifdim\@tempa\p@<\p@\scalebox{\@tempa}{\usebox\pandoc@box}%
  \else\usebox{\pandoc@box}%
  \fi%
}
% Set default figure placement to htbp
\def\fps@figure{htbp}
\makeatother



\ifLuaTeX
\usepackage[bidi=basic,shorthands=off,provide=*]{babel}
\else
\usepackage[bidi=default,shorthands=off,provide=*]{babel}
\fi
\ifLuaTeX
  \usepackage{selnolig} % disable illegal ligatures
\fi


\setlength{\emergencystretch}{3em} % prevent overfull lines

\providecommand{\tightlist}{%
  \setlength{\itemsep}{0pt}\setlength{\parskip}{0pt}}



 


\usepackage{luatexja-fontspec}
% PDF 本文の和文：Noto Sans CJK JP（サンセリフ、CI/Mac 共に確実）。
% HTML は IBM Plex Sans JP（_assets/ja-theme.scss 経由）で組まれるので、
% Web 読者にはモダンサンセリフ、PDF 読者にはサンセリフの Noto という
% 二系統だが、いずれも「明朝で読みにくい」問題は解消する。
% IBM Plex Sans JP の PDF 直接利用は luatexja-fontspec の
% ウェイトサフィックス命名（Family-Bold 等）と
% IBM Plex のファイル命名（IBMPlexSansJP-SemiBold）が
% 噛み合わず、現状はフォールバック運用とする（v0.3 で再挑戦予定）。
\setmainjfont{Noto Sans CJK JP}
\setsansjfont{Noto Sans CJK JP}
% 欧文側は IBM Plex Sans が CI にあれば使う、無ければ既定。
\usepackage{fontspec}
\IfFontExistsTF{IBM Plex Sans}{\setmainfont{IBM Plex Sans}}{}
\IfFontExistsTF{IBM Plex Mono}{\setmonofont{IBM Plex Mono}}{}
\usepackage{xcolor}
\usepackage{tcolorbox}
% Quarto の callout が要求するライブラリ群
% （enhanced jigsaw / breakable / theorems / raster / skins）
\tcbuselibrary{skins,breakable,theorems,raster}
\usepackage{amsmath,amssymb,physics}
% 旧 TeX フォント切替コマンド（\rm \bf \it \sf \tt）の互換性確保。
% scrreprt は標準でこれらをエラー化するため、無条件で再定義する。
% 数式モード内に混入していても \mathrm 等として扱われる。
\AtBeginDocument{%
  \let\rm\mathrm
  \let\bf\mathbf
  \let\it\mathit
  \let\sf\mathsf
  \let\tt\mathtt
}
\KOMAoption{captions}{tableheading}
\makeatletter
\@ifpackageloaded{tcolorbox}{}{\usepackage[skins,breakable]{tcolorbox}}
\@ifpackageloaded{fontawesome5}{}{\usepackage{fontawesome5}}
\definecolor{quarto-callout-color}{HTML}{909090}
\definecolor{quarto-callout-note-color}{HTML}{0758E5}
\definecolor{quarto-callout-important-color}{HTML}{CC1914}
\definecolor{quarto-callout-warning-color}{HTML}{EB9113}
\definecolor{quarto-callout-tip-color}{HTML}{00A047}
\definecolor{quarto-callout-caution-color}{HTML}{FC5300}
\definecolor{quarto-callout-color-frame}{HTML}{acacac}
\definecolor{quarto-callout-note-color-frame}{HTML}{4582ec}
\definecolor{quarto-callout-important-color-frame}{HTML}{d9534f}
\definecolor{quarto-callout-warning-color-frame}{HTML}{f0ad4e}
\definecolor{quarto-callout-tip-color-frame}{HTML}{02b875}
\definecolor{quarto-callout-caution-color-frame}{HTML}{fd7e14}
\makeatother
\makeatletter
\@ifpackageloaded{bookmark}{}{\usepackage{bookmark}}
\makeatother
\makeatletter
\@ifpackageloaded{caption}{}{\usepackage{caption}}
\AtBeginDocument{%
\ifdefined\contentsname
  \renewcommand*\contentsname{目次}
\else
  \newcommand\contentsname{目次}
\fi
\ifdefined\listfigurename
  \renewcommand*\listfigurename{図一覧}
\else
  \newcommand\listfigurename{図一覧}
\fi
\ifdefined\listtablename
  \renewcommand*\listtablename{表一覧}
\else
  \newcommand\listtablename{表一覧}
\fi
\ifdefined\figurename
  \renewcommand*\figurename{図}
\else
  \newcommand\figurename{図}
\fi
\ifdefined\tablename
  \renewcommand*\tablename{表}
\else
  \newcommand\tablename{表}
\fi
}
\@ifpackageloaded{float}{}{\usepackage{float}}
\floatstyle{ruled}
\@ifundefined{c@chapter}{\newfloat{codelisting}{h}{lop}}{\newfloat{codelisting}{h}{lop}[chapter]}
\floatname{codelisting}{コード}
\newcommand*\listoflistings{\listof{codelisting}{コード一覧}}
\makeatother
\makeatletter
\makeatother
\makeatletter
\@ifpackageloaded{caption}{}{\usepackage{caption}}
\@ifpackageloaded{subcaption}{}{\usepackage{subcaption}}
\makeatother
\usepackage{bookmark}
\IfFileExists{xurl.sty}{\usepackage{xurl}}{} % add URL line breaks if available
\urlstyle{same}
\hypersetup{
  pdftitle={宇宙の熱的放射},
  pdfauthor={和田 桂一},
  pdflang={ja},
  colorlinks=true,
  linkcolor={blue},
  filecolor={Maroon},
  citecolor={Blue},
  urlcolor={Blue},
  pdfcreator={LaTeX via pandoc}}


\title{宇宙の熱的放射}
\usepackage{etoolbox}
\makeatletter
\providecommand{\subtitle}[1]{% add subtitle to \maketitle
  \apptocmd{\@title}{\par {\large #1 \par}}{}{}
}
\makeatother
\subtitle{連続放射と線スペクトルの背景物理}
\author{和田 桂一}
\date{2026-06-26}
\begin{document}
\maketitle

\renewcommand*\contentsname{目次}
{
\setcounter{tocdepth}{2}
\tableofcontents
}

\bookmarksetup{startatroot}

\chapter{宇宙の熱的放射}\label{ux5b87ux5b99ux306eux71b1ux7684ux653eux5c04}

連続放射と線スペクトルの背景物理

\hfill\break

\bookmarksetup{startatroot}

\chapter*{版履歴}\label{sec-version-history}
\addcontentsline{toc}{chapter}{版履歴}

\markboth{版履歴}{版履歴}

\begin{tcolorbox}[enhanced jigsaw, arc=.35mm, bottomrule=.15mm, breakable, colback=white, colframe=quarto-callout-note-color-frame, left=2mm, leftrule=.75mm, opacityback=0, rightrule=.15mm, toprule=.15mm]
\begin{minipage}[t]{5.5mm}
\textcolor{quarto-callout-note-color}{\faInfo}
\end{minipage}%
\begin{minipage}[t]{\textwidth - 5.5mm}

本書はドラフト段階にあり、レビュアーからのフィードバックを反映しながら継続的に改訂中です。

\end{minipage}%
\end{tcolorbox}

{\def\LTcaptype{none} % do not increment counter
\begin{longtable}[]{@{}
  >{\raggedright\arraybackslash}p{(\linewidth - 6\tabcolsep) * \real{0.2500}}
  >{\raggedright\arraybackslash}p{(\linewidth - 6\tabcolsep) * \real{0.2500}}
  >{\raggedright\arraybackslash}p{(\linewidth - 6\tabcolsep) * \real{0.2500}}
  >{\raggedright\arraybackslash}p{(\linewidth - 6\tabcolsep) * \real{0.2500}}@{}}
\toprule\noalign{}
\begin{minipage}[b]{\linewidth}\raggedright
版
\end{minipage} & \begin{minipage}[b]{\linewidth}\raggedright
公開日／生成日
\end{minipage} & \begin{minipage}[b]{\linewidth}\raggedright
主な変更
\end{minipage} & \begin{minipage}[b]{\linewidth}\raggedright
引用
\end{minipage} \\
\midrule\noalign{}
\endhead
\bottomrule\noalign{}
\endlastfoot
\textbf{次版（作業中）} & 2026 秋目標 & 図版充実（v0.2 ターゲット 7/8
枚反映済）、用語統一、序文刷新、フォント刷新 等 & --- \\
\textbf{v0.1.2} & 2026-06-23（生成日） & 第19章に図2点を追加（§19.2
自由-自由スペクトル、§19.3--19.4 連続吸収端と再結合連続光） & --- \\
\textbf{v0.1.1} & 2026-06-23（生成日） & 第19章「自由電子が作る連続光 ―
自由-自由・束縛-自由・再結合」を新設、第VII・VIII部を1章ずつ繰り下げ（全
8 部 26 章） & --- \\
\textbf{v0.1} & 2026-06-11 & \textbf{初版 preprint} ― 全 8 部 25 章 +
付録、CC BY 4.0 & \href{https://doi.org/10.5281/zenodo.20635563}{DOI
10.5281/zenodo.20635563} \\
\end{longtable}
}

詳細な改訂内容は
\href{https://github.com/wadakeiichi/inverse_text_radiation/blob/main/CHANGELOG.md}{CHANGELOG.md}
と
\href{https://github.com/wadakeiichi/inverse_text_radiation/releases}{GitHub
Releases}。

オンライン版（GitHub Pages）は \texttt{main} ブランチへの push
のたびに自動更新されるため、本ページの内容は
\textbf{常に最新の作業中ドラフト}
です。「\textbf{固定版を引用したい}」場合は上記表の
DOI（特定の版）を使用してください。

\begin{center}\rule{0.5\linewidth}{0.5pt}\end{center}

\bookmarksetup{startatroot}

\chapter*{はじめに ― この本の三つの原則}\label{sec-three-principles}
\addcontentsline{toc}{chapter}{はじめに ― この本の三つの原則}

\markboth{はじめに ― この本の三つの原則}{はじめに ― この本の三つの原則}

\begin{tcolorbox}[enhanced jigsaw, arc=.35mm, bottomrule=.15mm, breakable, colback=white, colframe=quarto-callout-note-color-frame, left=2mm, leftrule=.75mm, opacityback=0, rightrule=.15mm, toprule=.15mm]
\begin{minipage}[t]{5.5mm}
\textcolor{quarto-callout-note-color}{\faInfo}
\end{minipage}%
\begin{minipage}[t]{\textwidth - 5.5mm}

\textbf{英題（参考）}：\emph{Thermal Radiation in the Universe: An
Inverse Approach to the Physics}

\end{minipage}%
\end{tcolorbox}

本書は、宇宙物理の中心的観測対象である
\textbf{連続スペクトル（黒体放射）} と \textbf{線スペクトル}
を入口に、基礎物理を逆向きに学び直すための教科書である。次の三つの原則を貫く。

\begin{tcolorbox}[enhanced jigsaw, arc=.35mm, bottomrule=.15mm, bottomtitle=1mm, breakable, colback=white, colbacktitle=quarto-callout-note-color!10!white, colframe=quarto-callout-note-color-frame, coltitle=black, left=2mm, leftrule=.75mm, opacityback=0, opacitybacktitle=0.6, rightrule=.15mm, title=\textcolor{quarto-callout-note-color}{\faInfo}\hspace{0.5em}{ノート}, titlerule=0mm, toprule=.15mm, toptitle=1mm]

\textbf{本書の射程について}：書名の「熱的スペクトル」は、本書が
\textbf{熱的連続放射（プランク分布とそこからのずれ）と熱的線スペクトル（LTE
近似が成り立つ条件下の線形成）を主軸とする}
ことを示す。これは「逆引きの出発点」を限定することで、現代物理学の主要分野を一筆書きでつなぐ本書の方法論を可能にしている。

整理すると、本書の \textbf{主軸} は次の二つである。

\begin{itemize}
\tightlist
\item
  \textbf{熱的連続放射}（第I〜VI部）：プランク分布と、そこからのずれ
\item
  \textbf{熱的線スペクトル}（第VII部）：LTE
  近似が成り立つ条件下の線形成・原子物理・線輪郭（line
  profile）・観測応用
\end{itemize}

これに対し、読者が主軸の射程と限界を見定められるよう、次のものを
\textbf{並走対象}（主軸を相対化する比較対象）として扱う。

\begin{itemize}
\tightlist
\item
  \textbf{非熱的放射}（第15章
  §15.6、第18章）：シンクロトロン・逆コンプトン散乱の存在と、熱的放射との区別法
\item
  \textbf{non-LTE 線形成}（第21章 §21.5）：物質側統計が Boltzmann/Saha
  分布から逸脱する場合
\item
  \textbf{21 cm 線・AGN
  線}（第20・24章）：低密度・非平衡環境という、熱的線形成の極限
\item
  \textbf{QED
  の見通し}（第VIII部）：連続と線を統一する電磁場--物質相互作用
  Hamiltonian
\end{itemize}

並走対象は主軸そのものではないが、「どこまでが LTE
の熱的描像で語れ、どこから別の枠組みが要るか」を明確にするために置いている。本格的な非熱的天体物理（高エネルギー宇宙線、GRB、AGN
ジェット）は Rybicki \& Lightman、Longair 等の専門書を併読してほしい。

\end{tcolorbox}

\textbf{原則 1 ― 逆引き}

観測されるスペクトルという完成形から、その背後にある物理（放射輸送・統計力学・量子力学・電磁気学・原子物理・QED）を逆向きにたどる。

\textbf{原則 2 ― 背景物理を曖昧にしない}

各観測量の背後にどの物理が、どんな前提のもとで効いているかを、章を横断して明示する。ブラックボックス化や「ここでは扱わない」を避ける。

\textbf{原則 3 ― 天下り的に式を与えない}

プランク分布、リュードベリ公式、振動子強度、Voigt
線輪郭などを結果として与えるのではなく、観測から逆向きに辿って
\textbf{必然性} として理解できるようにする。

\section*{本書の二つの中心地図}\label{sec-two-maps}
\addcontentsline{toc}{section}{本書の二つの中心地図}

\markright{本書の二つの中心地図}

本書は二つの「中心地図」を持つ。各部の冒頭でこの地図を再掲し、強調する因子を切り替えながら、章群がどの物理を担っているかを明示する。

\subsection{連続：プランク関数}\label{ux9023ux7d9aux30d7ux30e9ux30f3ux30afux95a2ux6570}

\textbf{プランク関数（連続スペクトルの地図）}

\begin{equation}\protect\phantomsection\label{eq-planck-function}{
B_\nu(T) = \frac{2 h \nu^3}{c^2} \cdot \frac{1}{e^{h\nu/kT} - 1}
}\end{equation}

{\def\LTcaptype{none} % do not increment counter
\begin{longtable}[]{@{}lll@{}}
\toprule\noalign{}
因子 & 物理的意味 & 結びつく部 \\
\midrule\noalign{}
\endhead
\bottomrule\noalign{}
\endlastfoot
\(B_\nu\) & 観測量（比強度） & 第II部 \\
\(1/(e^{h\nu/kT}-1)\) & 熱平衡・光子統計 & 第III部 \\
\(h\nu\) & エネルギー量子 & 第IV部 \\
\(\nu^3/c^2\) & 電磁場のモード密度 & 第V部 \\
黒体からのずれ & 現実の天体 & 第VI部 \\
\end{longtable}
}

\begin{tcolorbox}[enhanced jigsaw, arc=.35mm, bottomrule=.15mm, breakable, colback=white, colframe=quarto-callout-tip-color-frame, left=2mm, leftrule=.75mm, opacityback=0, rightrule=.15mm, toprule=.15mm]

\textbf{地図の使い方}：各部の冒頭でこの式が再掲され、その部で扱う因子が強調表示される。読者は自分が今どこを学んでいるかをこの式の上で確認できる。

\end{tcolorbox}

\subsection{線：水素の線吸収係数}\label{ux7ddaux6c34ux7d20ux306eux7ddaux5438ux53ceux4fc2ux6570}

\textbf{水素線スペクトルの地図 ― 二層構造}

\textbf{第一層：イコニック・アンカー（リュードベリ公式）}

\begin{equation}\protect\phantomsection\label{eq-rydberg}{
\frac{1}{\lambda} = R_\mathrm{H} \left( \frac{1}{n_1^2} - \frac{1}{n_2^2} \right)
}\end{equation}

線スペクトルの「骨格」を一目で示す。\(R_\mathrm{H}\)
は基本定数の組合せ、\(1/n^2\)
は束縛状態のエネルギー、離散性が「線が存在する理由」を語る。

\textbf{第二層：構造的マップ（Hα を標準例とした線吸収係数）}

\begin{equation}\protect\phantomsection\label{eq-line-absorption}{
\alpha_\nu = \frac{\pi e^2}{m_e c} \, f_{lu} \, n_l \, \phi(\nu - \nu_0)
}\end{equation}

{\def\LTcaptype{none} % do not increment counter
\begin{longtable}[]{@{}lll@{}}
\toprule\noalign{}
因子 & 物理的意味 & 対応する章 \\
\midrule\noalign{}
\endhead
\bottomrule\noalign{}
\endlastfoot
\(\alpha_\nu\) & 観測量（線吸収係数） & 第21章 \\
\(\pi e^2 / m_e c\) & 古典極限の線強度 & 第22章＋第VIII部 \\
\(f_{lu}\) & 振動子強度 & 第22章 \\
\(n_l\) & 下準位の占有数 & 第21章＋Part III \\
\(\phi(\nu - \nu_0)\) & 線輪郭（自然・Doppler・Voigt） & 第23章 \\
\(\nu_0\) & 線位置（リュードベリ） & 第22章 \\
\end{longtable}
}

\textbf{適用範囲（重要）}：この式の前因子 \(\pi e^2/m_e c\) と振動子強度
\(f_{lu}\) は、古典 Lorentz 振動子に対応する
\textbf{電気双極子（E1）許容線} に対するものである。21 cm
線（基底状態超微細構造の \textbf{磁気双極子 M1 遷移}）や
\textbf{禁制線（E2・M1 など）} では、この前因子・\(f_{lu}\)
をそのまま使えない。これらの線では、測定された自発放出率 \(A_{ul}\) から
Einstein
関係を介して有効遷移強度（\(B_{lu}\)）に読み替える（式~\ref{eq-line-absorption}
の構造は保ったまま、強度因子だけを差し替える）。

\begin{tcolorbox}[enhanced jigsaw, arc=.35mm, bottomrule=.15mm, breakable, colback=white, colframe=quarto-callout-tip-color-frame, left=2mm, leftrule=.75mm, opacityback=0, rightrule=.15mm, toprule=.15mm]

\textbf{地図の使い方}：第VII部の各章はこの式の異なる因子を解き明かす。プランク関数地図と並べて読めば、本書を貫く二つの逆引きが同時に把握できる。第VIII部の
QED で、両者は前因子 \(2h\nu^3/c^2\)
を介して統一される（式~\ref{eq-einstein-A} 参照）。

\end{tcolorbox}

両者は、第VIII部の量子電磁力学において、同じ電磁場--物質相互作用の二側面として統合される。

\section*{読者タイプ別おすすめルート}\label{sec-routes}
\addcontentsline{toc}{section}{読者タイプ別おすすめルート}

\markright{読者タイプ別おすすめルート}

通読も可能だが、目的に応じて読む順序を変えられる。詳細は
\hyperref[sec-prologue-reading-routes]{序章 §0.4} 参照。

{\def\LTcaptype{none} % do not increment counter
\begin{longtable}[]{@{}
  >{\raggedright\arraybackslash}p{(\linewidth - 4\tabcolsep) * \real{0.3333}}
  >{\raggedright\arraybackslash}p{(\linewidth - 4\tabcolsep) * \real{0.3333}}
  >{\raggedright\arraybackslash}p{(\linewidth - 4\tabcolsep) * \real{0.3333}}@{}}
\toprule\noalign{}
\begin{minipage}[b]{\linewidth}\raggedright
ルート
\end{minipage} & \begin{minipage}[b]{\linewidth}\raggedright
想定読者
\end{minipage} & \begin{minipage}[b]{\linewidth}\raggedright
おおまかな順序
\end{minipage} \\
\midrule\noalign{}
\endhead
\bottomrule\noalign{}
\endlastfoot
\textbf{A} & まず宇宙の現象をつかみたい & 第I → II → VI →
VII（20・24章） \\
\textbf{B} & 標準的に一通り学びたい & 第I 〜 VI（線も読むなら VII） \\
\textbf{C} & 場の量子論まで見通したい & 第I 〜 V →
VIII（VII経由が望ましい） \\
\textbf{D} & 連続スペクトルを実践的に & 第I（1〜3章）→ II（4〜5章）→
VI \\
\textbf{E} & 分光学・線スペクトルから & 第I（1章）→ II → VII（→
IV・VIII） \\
\end{longtable}
}

\section*{難易度表示}\label{sec-difficulty}
\addcontentsline{toc}{section}{難易度表示}

\markright{難易度表示}

各節見出しに「\textbf{本文目安}」（節の本文を読むのに必要な学年）と、必要に応じて「\textbf{導出}」（背後の計算・証明を追えるレベル）を付す。

想定する標準カリキュラム：

{\def\LTcaptype{none} % do not increment counter
\begin{longtable}[]{@{}lll@{}}
\toprule\noalign{}
表示 & 想定 & 内容 \\
\midrule\noalign{}
\endhead
\bottomrule\noalign{}
\endlastfoot
\textbf{高校} & 高校生 & 高校物理・数学で概念的に読める \\
\textbf{B1} & 大学1年 & 微積分・線形代数・力学入門 \\
\textbf{B2} & 大学2年 & 電磁気・熱力学・解析力学の入口 \\
\textbf{B3} & 大学3年 & 量子力学・統計力学・電磁気の標準内容 \\
\textbf{B4} & 卒研前後 & 天体物理・放射過程・量子統計を使う \\
\textbf{M1} & 大学院初年 & 放射輸送・統計力学・場の量子論の入口 \\
\textbf{M2+} & 修士後半〜博士前期 & 専門的応用・論文読解 \\
\textbf{D} & 博士後期 & 研究レベルの補足・発展 \\
\end{longtable}
}

\textbf{本書全体の方針}：本文は最大でも \textbf{M1}
までに抑える。\textbf{M2+} や \textbf{D}
の内容は、発展コラム・導出ノートに分離し、本筋を読むのに必須でないことを明示する。大学によってカリキュラムは違うので、これはあくまで一つの目安であり、理解の優劣を示すものではない。

\textbf{表記例}：

\begin{itemize}
\tightlist
\item
  \texttt{{[}本文目安：B2{]}}：本文を読むのに大学2年程度
\item
  \texttt{{[}本文目安：B3｜導出：M1{]}}：本文は B3、背後の導出は M1
  レベル
\end{itemize}

\begin{tcolorbox}[enhanced jigsaw, arc=.35mm, bottomrule=.15mm, bottomtitle=1mm, breakable, colback=white, colbacktitle=quarto-callout-tip-color!10!white, colframe=quarto-callout-tip-color-frame, coltitle=black, left=2mm, leftrule=.75mm, opacityback=0, opacitybacktitle=0.6, rightrule=.15mm, title=\textcolor{quarto-callout-tip-color}{\faLightbulb}\hspace{0.5em}{オンライン版の使い方}, titlerule=0mm, toprule=.15mm, toptitle=1mm]

用語・概念・式にリンクが埋め込まれている。「プランク分布」から「光子統計」「モード密度」「CMB」へ、「Hα」から「リュードベリ公式」「振動子強度」「Voigt
線輪郭」「HII 領域」へ、連続側と線側の両方で網目状に行き来できる。

\end{tcolorbox}

\bookmarksetup{startatroot}

\chapter{目次}\label{ux76eeux6b21}

\begin{tcolorbox}[enhanced jigsaw, arc=.35mm, bottomrule=.15mm, breakable, colback=white, colframe=quarto-callout-note-color-frame, left=2mm, leftrule=.75mm, opacityback=0, rightrule=.15mm, toprule=.15mm]
\begin{minipage}[t]{5.5mm}
\textcolor{quarto-callout-note-color}{\faInfo}
\end{minipage}%
\begin{minipage}[t]{\textwidth - 5.5mm}

本書の全体像を一望できるよう、全部・全章を一覧する。各タイトルからその章へ直接ジャンプできる。

\end{minipage}%
\end{tcolorbox}

\href{00-prologue.qmd}{序章　この本の目的と歩き方}

\section*{第I部　宇宙に満ちる黒体放射}\label{ux7b2ciux90e8-ux5b87ux5b99ux306bux6e80ux3061ux308bux9ed2ux4f53ux653eux5c04}
\addcontentsline{toc}{section}{第I部　宇宙に満ちる黒体放射}

\markright{第I部　宇宙に満ちる黒体放射}

\begin{itemize}
\tightlist
\item
  \href{part1/01-where-blackbody.qmd}{第1章　宇宙で黒体放射はどこに現れるか}
\item
  \href{part1/02-reading-spectra.qmd}{第2章　黒体スペクトルをどう読むか}
\item
  \href{part1/03-universality.qmd}{第3章　黒体放射はなぜ普遍的なのか}
\end{itemize}

\section*{第II部　放射をどう記述するか}\label{ux7b2ciiux90e8-ux653eux5c04ux3092ux3069ux3046ux8a18ux8ff0ux3059ux308bux304b}
\addcontentsline{toc}{section}{第II部　放射をどう記述するか}

\markright{第II部　放射をどう記述するか}

\begin{itemize}
\tightlist
\item
  \href{part2/04-radiation-field.qmd}{第4章　放射場の基本量}
\item
  \href{part2/05-radiative-transfer.qmd}{第5章　放射輸送の基礎}
\end{itemize}

\section*{第III部　逆引き1：熱力学・統計力学へ}\label{ux7b2ciiiux90e8-ux9006ux5f15ux304d1ux71b1ux529bux5b66ux7d71ux8a08ux529bux5b66ux3078}
\addcontentsline{toc}{section}{第III部　逆引き1：熱力学・統計力学へ}

\markright{第III部　逆引き1：熱力学・統計力学へ}

\begin{itemize}
\tightlist
\item
  \href{part3/06-thermal-equilibrium.qmd}{第6章　熱平衡とは何か}
\item
  \href{part3/07-statistical-mechanics.qmd}{第7章　統計力学の基本}
\item
  \href{part3/08-photon-statistics.qmd}{第8章　光子の統計}
\end{itemize}

\section*{第IV部　逆引き2：量子論へ}\label{ux7b2civux90e8-ux9006ux5f15ux304d2ux91cfux5b50ux8ad6ux3078}
\addcontentsline{toc}{section}{第IV部　逆引き2：量子論へ}

\markright{第IV部　逆引き2：量子論へ}

\begin{itemize}
\tightlist
\item
  \href{part4/09-classical-failure.qmd}{第9章　古典論ではなぜ失敗するのか}
\item
  \href{part4/10-planck-quantum.qmd}{第10章　プランクの量子仮説}
\item
  \href{part4/11-einstein-method.qmd}{第11章　アインシュタインの方法}
\item
  \href{part4/12-quantum-foundations.qmd}{第12章　量子力学の基礎へ戻る}
\end{itemize}

\section*{第V部　逆引き3：電磁気学へ}\label{ux7b2cvux90e8-ux9006ux5f15ux304d3ux96fbux78c1ux6c17ux5b66ux3078}
\addcontentsline{toc}{section}{第V部　逆引き3：電磁気学へ}

\markright{第V部　逆引き3：電磁気学へ}

\begin{itemize}
\tightlist
\item
  \href{part5/13-electromagnetic-waves.qmd}{第13章　電磁波としての光}
\item
  \href{part5/14-cavity-modes.qmd}{第14章　空洞中の電磁モード}
\end{itemize}

\section*{第VI部　現実の宇宙へ戻る}\label{ux7b2cviux90e8-ux73feux5b9fux306eux5b87ux5b99ux3078ux623bux308b}
\addcontentsline{toc}{section}{第VI部　現実の宇宙へ戻る}

\markright{第VI部　現実の宇宙へ戻る}

\begin{itemize}
\tightlist
\item
  \href{part6/15-non-blackbody.qmd}{第15章　現実の天体は完全な黒体ではない}
\item
  \href{part6/16-cmb.qmd}{第16章　宇宙マイクロ波背景放射}
\item
  \href{part6/17-stars-planets-dust.qmd}{第17章　恒星・惑星・ダスト・電離ガスの熱放射}
\item
  \href{part6/18-accretion-high-energy.qmd}{第18章　降着円盤と高エネルギー天体}
\item
  \href{part6/19-continuum-microprocesses.qmd}{第19章　自由電子が作る連続光
  ― 自由-自由・束縛-自由・再結合}
\end{itemize}

\section*{第VII部　線スペクトル ―
もう一つの逆引き}\label{ux7b2cviiux90e8-ux7ddaux30b9ux30daux30afux30c8ux30eb-ux3082ux3046ux4e00ux3064ux306eux9006ux5f15ux304d}
\addcontentsline{toc}{section}{第VII部　線スペクトル ― もう一つの逆引き}

\markright{第VII部　線スペクトル ― もう一つの逆引き}

\begin{itemize}
\tightlist
\item
  \href{part7/20-where-lines.qmd}{第20章　線スペクトルはどこに現れるか}
\item
  \href{part7/21-line-formation.qmd}{第21章　線形成 ―
  吸収線と輝線の物理}
\item
  \href{part7/22-atomic-physics.qmd}{第22章　線スペクトルから原子物理へ}
\item
  \href{part7/23-line-shape.qmd}{第23章　線輪郭（line profile）の物理}
\item
  \href{part7/24-line-applications.qmd}{第24章　線スペクトルから天体物理を読む}
\end{itemize}

\section*{第VIII部　QEDからの俯瞰 ―
連続と線を結ぶ見通し}\label{ux7b2cviiiux90e8-qedux304bux3089ux306eux4fefux77b0-ux9023ux7d9aux3068ux7ddaux3092ux7d50ux3076ux898bux901aux3057}
\addcontentsline{toc}{section}{第VIII部　QEDからの俯瞰 ―
連続と線を結ぶ見通し}

\markright{第VIII部　QEDからの俯瞰 ― 連続と線を結ぶ見通し}

\begin{itemize}
\tightlist
\item
  \href{part8/25-em-field-quantization.qmd}{第25章　電磁場を量子化するとは何か
  ― 直観的描像}
\item
  \href{part8/26-qed-minimum.qmd}{第26章　Einstein係数・自然幅・選択則を
  QED で見通す}
\end{itemize}

\section*{参考文献・付録}\label{ux53c2ux8003ux6587ux732eux4ed8ux9332}
\addcontentsline{toc}{section}{参考文献・付録}

\markright{参考文献・付録}

\begin{itemize}
\tightlist
\item
  \href{references.qmd}{参考文献}
\item
  \href{appendix/A-math-notes.qmd}{付録 A　数学補章}
\item
  \href{appendix/B-units-constants.qmd}{付録 B　単位と基本定数}
\item
  \href{appendix/C-glossary.qmd}{用語索引（和英・英和）}
\end{itemize}

\bookmarksetup{startatroot}

\chapter{序章　この本の目的と歩き方}\label{ux5e8fux7ae0-ux3053ux306eux672cux306eux76eeux7684ux3068ux6b69ux304dux65b9}

\section{0.1 宇宙の熱的放射から何を学ぶのか}\label{sec-prologue-purpose}

夜空の天体をスペクトルに分解すると、そこにはほぼ例外なく二種類の成分が現れる。振動数に対して滑らかに変化する
\textbf{連続スペクトル}
と、原子や分子の離散的なエネルギー準位間の遷移に由来し、特定の振動数の周りに集中して現れる
\textbf{線スペクトル}
である。後者は、水素のバルマー線のように鋭い線として見えることもあれば、多数の線が密集して分子バンドを形作ることも、PAH（多環芳香族炭化水素）の放射のように幅の広い山や谷として現れることもある。太陽や恒星の連続光、惑星や星間塵の赤外放射、宇宙マイクロ波背景放射（CMB）は前者の、恒星大気の吸収線、星雲の輝線、中性水素の
21 cm
線、分子線、AGN・クェーサーの輝線・吸収線は後者の代表例であり、宇宙の物理状態について私たちが知っていることの大半は、この二つの成分を読み解くことから得られてきた。

本書が扱うのは、このうち \textbf{熱的放射} ―
おおまかに言えば、熱平衡（またはそれに近い状態）にある物質が出す放射 ―
である。連続側の典型が黒体放射であり、線側も多くは熱的なガスの中で形成される。なぜ「熱的」に絞るのか。熱的放射はスペクトルの形が温度というただ一つの量で（線ならば温度・密度・組成という少数の量で）統制されており、その背後には熱平衡・統計力学・量子論・電磁気学という基礎物理の構造が、最も純粋な形で現れているからである。シンクロトロン放射に代表される非熱的放射は、観測スペクトルの解釈に必要な範囲で触れるにとどめる（この区別の正確な定義は第1章
§1.9 で与える）。

本書の目的は、これらの放射を支配する式を、暗記すべき公式として与えることではない。観測されるスペクトルの形・強度・線位置・線形から出発し、「\textbf{なぜこの形なのか}」を逆向きにたどる。連続側ではプランク関数を、線側では水素の線吸収係数を道しるべとして、放射の記述、熱平衡、光子の統計、エネルギー量子化、電磁波のモード構造、原子の量子状態、線形成をへて、最終的には両者が
QED
の中で一つに合流するところまでを、一つの見通しの中で理解することを目指す。

\section{0.2 対象読者と前提知識}\label{sec-prologue-audience}

本書は、物理学科、天文学科、宇宙物理学科などの学部学生、および大学院修士課程の学生を主な対象とする。観測スペクトル（連続・線の両方）と基礎物理を結びつけて理解したい読者、既存教科書で個々の式を学んだあと「結局この式はなぜこうなのか」を観測現象から再確認したい読者にも向いている。

前提知識としては、微積分、指数関数と対数、基本的な力学、初歩的な電磁気学。熱力学・量子力学・原子物理の素養があると望ましいが、第I・II部から第III部以降への遡り方は読者の進度に合わせて調整できる。

\section{0.3 本書の構成}\label{sec-prologue-structure}

本書は次の構造をとる。

\begin{enumerate}
\def\labelenumi{\arabic{enumi}.}
\tightlist
\item
  \textbf{第I・II部}：宇宙のスペクトル（連続＋線）の現象論と、放射場の記述
\item
  \textbf{第III〜V部}：連続スペクトル（黒体放射）の逆引き ―
  統計力学・量子論・電磁気学へ
\item
  \textbf{第VI部}：現実の天体で連続放射がどう現れ、どうずれるか
\item
  \textbf{第VII部}：線スペクトルの逆引き ―
  線形成・原子物理・線形・観測応用
\item
  \textbf{第VIII部}：QED で連続と線を統一する
\end{enumerate}

\section{0.4
読者タイプ別おすすめルート}\label{sec-prologue-reading-routes}

本書は通読も可能だが、目的に応じて読み方を変えられる。各部の
\textbf{どこから入るか} ・ \textbf{どこで深掘りするか}
を、読者の関心に合わせて選んでほしい。

\subsection{ルート
A：まず宇宙の現象をつかみたい}\label{ux30ebux30fcux30c8-aux307eux305aux5b87ux5b99ux306eux73feux8c61ux3092ux3064ux304bux307fux305fux3044}

\textbf{第I → II → VI → VII（20・24章）→ 必要に応じて III〜V・VII
の理論章へ}

観測される宇宙のスペクトル（連続も線も）の \textbf{現象論}
を先に押さえ、必要になったら基礎物理に戻る。「逆引き」の精神そのままを、より大胆に実践するルート。

\subsection{ルート
B：標準的に一通り学びたい}\label{ux30ebux30fcux30c8-bux6a19ux6e96ux7684ux306bux4e00ux901aux308aux5b66ux3073ux305fux3044}

\textbf{第I 〜 VI（線も読むなら VII）}

本書の章番号順に通読。連続スペクトル側を完結させてから、必要に応じて線スペクトル側（第VII部）へ進む。

\subsection{ルート
C：場の量子論まで見通したい}\label{ux30ebux30fcux30c8-cux5834ux306eux91cfux5b50ux8ad6ux307eux3067ux898bux901aux3057ux305fux3044}

\textbf{第I 〜 V → VIII（連続と線の統一視点で読むなら VII 経由）}

QED で連続と線が統一される第VIII部までを \textbf{見通す}
ルート。第VI部の応用章を後回しにし、骨格となる理論部分（I-V）と最終章（VIII）の二点で本書の構造を把握する。

\subsection{ルート
D：連続スペクトルを実践的に}\label{ux30ebux30fcux30c8-dux9023ux7d9aux30b9ux30daux30afux30c8ux30ebux3092ux5b9fux8df5ux7684ux306b}

\textbf{第I（1〜3章）→ II（4〜5章）→ VI}

CMB・恒星・降着円盤の \textbf{実践的な解析}
にすぐ取り組みたい読者向け。理論章（III〜V）を最小限の参照に留め、観測解析の道具立てだけを先に固める。

\subsection{ルート
E：分光学・線スペクトルから}\label{ux30ebux30fcux30c8-eux5206ux5149ux5b66ux7ddaux30b9ux30daux30afux30c8ux30ebux304bux3089}

\textbf{第I（1章 特に §1.8）→ II（4〜5章）→ VII（→ IV・VIII で深掘り）}

線スペクトル（恒星吸収線・星雲輝線・分子線等）を中心に学びたい読者向け。第VII部を主軸とし、必要に応じて量子論（第IV部）と
QED（第VIII部）に遡る。

\subsection{どのルートでも}\label{ux3069ux306eux30ebux30fcux30c8ux3067ux3082}

中心地図（プランク関数と水素線吸収係数）は、ほぼすべての章で再登場する。\textbf{どの章から読み始めても、まずは中心地図に戻る}
ことを習慣にしてほしい。これが本書の三原則（① 逆引き ②
背景物理を曖昧にしない ③
天下り的に与えない）を読者の側から実装する最も簡単な方法である。

\section{0.5
難易度表示と本文・導出・発展の区別}\label{sec-prologue-difficulty}

各節見出しに、その節の \textbf{読み手のレベル目安}
を物理学科の標準的な学年スケールで示す：

{\def\LTcaptype{none} % do not increment counter
\begin{longtable}[]{@{}ll@{}}
\toprule\noalign{}
表記 & レベル \\
\midrule\noalign{}
\endhead
\bottomrule\noalign{}
\endlastfoot
\textbf{高校} & 高校物理 \\
\textbf{B1}〜\textbf{B2} & 学部 1〜2
年（力学・電磁気学・熱力学の基礎） \\
\textbf{B3}〜\textbf{B4} & 学部 3〜4
年（量子力学・統計力学の基礎を経た段階） \\
\textbf{M1}〜\textbf{M2+} & 大学院修士課程 \\
\textbf{D} & 博士後期課程 \\
\end{longtable}
}

そして節ごとに、本書の三段の役割を区別する：

\begin{itemize}
\tightlist
\item
  \textbf{本文目安：X} ― その節の本文を理解できる学年レベル
\item
  \textbf{導出：Y} ―
  厳密な導出を追うのに必要な学年レベル（本文目安より高いことが多い）
\item
  \textbf{発展} ― 飛ばしても本筋を追える、より深い議論を含む節
\end{itemize}

例：「{[}本文目安：B3／導出：M1{]}」は、本文は学部 3
年で読めるが、導出を逐一追うには修士 1
年の道具が必要、の意。「{[}本文目安：B4／発展：M2+{]}」のように
\textbf{本文目安と発展レベル} を併記している場合もある。

「本文目安：X」のみの節は、本文と導出のレベルがほぼ等しい節である。

\section{0.6 オンライン版の使い方}\label{sec-prologue-online}

オンライン版では、用語・記号・主要な式にリンクを付し、関連節へ移動できるようにする。本文のほかに、導出を丁寧にたどるノート、観測データの実例、演習問題の補助解説、必要な数学への短い解説を階層的に配置する。本書のねらいの一つである「\textbf{背景物理を曖昧にしない}」を支える仕組みでもある。

\section{0.7 図の一覧}\label{sec-prologue-figures}

本書に登場する主要な図を一覧する。各図番号はクリックすると該当箇所へ移動できる（オンライン版・PDF版とも）。「再掲」は、他章で初出した図を文脈を変えて再登場させたものである。中心地図（図~\ref{fig-planck-map}・図~\ref{fig-hydrogen-map}）は、ほぼすべての章で姿を変えて再登場する。

\textbf{第I部　宇宙に満ちる黒体放射}

\begin{itemize}
\tightlist
\item
  図~\ref{fig-planck-map}：プランク関数の各因子と、それを逆引きする各部の対応（中心地図）
\item
  図~\ref{fig-planck-curves}：多様な温度のプランク曲線
\item
  図~\ref{fig-planck-three-temps}：三つの温度のプランク曲線（ピーク移動と
  \(T^4\) 則）
\end{itemize}

\textbf{第II部　放射をどう記述するか}

\begin{itemize}
\tightlist
\item
  図~\ref{fig-specific-intensity}：比強度 \(I_\nu\) の定義の幾何
\end{itemize}

\textbf{第III部　逆引き1：熱力学・統計力学へ}

\begin{itemize}
\tightlist
\item
  図~\ref{fig-quantum-freezeout}：量子化振動子の平均エネルギーと等分配の比較（高振動数モードの凍結）
\end{itemize}

\textbf{第IV部　逆引き2：量子論へ}

\begin{itemize}
\tightlist
\item
  図~\ref{fig-uv-catastrophe}：紫外破綻 ― Rayleigh--Jeans
  則の発散とプランク分布
\item
  図~\ref{fig-quantum-freezeout-ch10}：高振動数モードの凍結（第7章の再掲）
\item
  図~\ref{fig-einstein-three-processes}：二準位系の三つの素過程と
  Einstein 係数
\end{itemize}

\textbf{第V部　逆引き3：電磁気学へ}

\begin{itemize}
\tightlist
\item
  図~\ref{fig-k-space}：\(k\) 空間の格子点の数え上げ（前因子
  \(8\pi\nu^2/c^3\) の幾何）
\end{itemize}

\textbf{第VII部　線スペクトル ― もう一つの逆引き}

\begin{itemize}
\tightlist
\item
  図~\ref{fig-hydrogen-map}：水素線吸収係数の各因子と、それを逆引きする章群の対応（中心地図）
\item
  図~\ref{fig-line-geometry}：同じ遷移が吸収線にも輝線にも見える幾何（Kirchhoff
  の三法則）
\item
  図~\ref{fig-equivalent-width}：等価幅 \(W_\lambda\) の幾何的定義
\item
  図~\ref{fig-einstein-three-processes-line}：三つの素過程の準位図（第IV部の再掲）
\item
  図~\ref{fig-line-geometry-lte}：吸収線・輝線の幾何（第20章の再掲、ソース関数の文脈）
\item
  図~\ref{fig-voigt-profile}：Voigt プロファイルと Gaussian／Lorentzian
  の比較
\end{itemize}

\part{第I部　宇宙に満ちる黒体放射}

\chapter{第1章　宇宙で黒体放射はどこに現れるか}\label{ux7b2c1ux7ae0-ux5b87ux5b99ux3067ux9ed2ux4f53ux653eux5c04ux306fux3069ux3053ux306bux73feux308cux308bux304b}

\textbf{この章の主題}：宇宙のさまざまな場所に現れる黒体放射の実例を概観する。出発点は理論ではなく、観測されるスペクトルそのものである。本書全体を貫く三原則のうち、本章は「①
逆引き」の入口に当たる。

\section*{この章の中心地図}\label{sec-where-blackbody-map}
\addcontentsline{toc}{section}{この章の中心地図}

\markright{この章の中心地図}

\begin{figure}[H]

\centering{

\includegraphics[width=0.9\linewidth,height=\textheight,keepaspectratio]{index_files/mediabag/part1/../images/planck-map.pdf}

}

\caption{\label{fig-planck-map}プランク関数の各因子と、それを逆引きする各部の対応。\textbf{因子の色＝担当する部の色}で結ばれており、矢印に頼らず読み取れる設計になっている。観測（第I-II部）から出発し、各因子の物理を逆引きしたうえで、最終的に第VIII部
QED
で同じ式が「電磁場の熱状態」として書き直される。各部の冒頭でこの地図が再掲され、その部が担当する因子のハイライトが切り替わる。}

\end{figure}%

\textbf{プランク関数（連続スペクトルの地図）}

\begin{equation}\protect\phantomsection\label{eq-planck-function}{
B_\nu(T) = \frac{2 h \nu^3}{c^2} \cdot \frac{1}{e^{h\nu/kT} - 1}
}\end{equation}

{\def\LTcaptype{none} % do not increment counter
\begin{longtable}[]{@{}lll@{}}
\toprule\noalign{}
因子 & 物理的意味 & 結びつく部 \\
\midrule\noalign{}
\endhead
\bottomrule\noalign{}
\endlastfoot
\(B_\nu\) & 観測量（比強度） & 第II部 \\
\(1/(e^{h\nu/kT}-1)\) & 熱平衡・光子統計 & 第III部 \\
\(h\nu\) & エネルギー量子 & 第IV部 \\
\(\nu^3/c^2\) & 電磁場のモード密度 & 第V部 \\
黒体からのずれ & 現実の天体 & 第VI部 \\
\end{longtable}
}

\begin{tcolorbox}[enhanced jigsaw, arc=.35mm, bottomrule=.15mm, breakable, colback=white, colframe=quarto-callout-tip-color-frame, left=2mm, leftrule=.75mm, opacityback=0, rightrule=.15mm, toprule=.15mm]

\textbf{地図の使い方}：各部の冒頭でこの式が再掲され、その部で扱う因子が強調表示される。読者は自分が今どこを学んでいるかをこの式の上で確認できる。

\end{tcolorbox}

\section*{この章で答える問い}\label{sec-where-blackbody-questions}
\addcontentsline{toc}{section}{この章で答える問い}

\markright{この章で答える問い}

\begin{itemize}
\tightlist
\item
  なぜ太陽光のスペクトルは、ある「一つの形」に近いのか（→
  第7-10章で逆引き、§17.1 で応用）
\item
  同じ「黒体に近い形」が、6000 K の恒星と 2.7 K
  の宇宙背景放射でなぜ同時に成り立つのか（→ 第3章、第16章）
\item
  「黒体」というのは観測される実物のことなのか、それとも理想化された概念なのか（→
  §1.7、第6章）
\item
  観測されるスペクトルが「黒体に近い」とは、具体的にどう判定するのか（→
  第15章）
\item
  なぜ降着円盤のような激しい場所でも、局所的には黒体的なスペクトルが見えるのか（→
  第6章 LTE、§18.1）
\end{itemize}

\section*{到達目標}\label{sec-where-blackbody-goals}
\addcontentsline{toc}{section}{到達目標}

\markright{到達目標}

この章を読み終えると、読者は次のことができるようになる：

\begin{itemize}
\tightlist
\item
  宇宙のさまざまな天体（恒星・惑星・ダスト・CMB・降着円盤）に、どのような形で黒体放射が現れるかを列挙できる
\item
  「黒体」が観測対象ではなく理想化された概念であることを説明できる
\item
  観測スペクトルが黒体に「近い」とき、どの量を比較すれば判定できるかを述べられる
\end{itemize}

\begin{center}\rule{0.5\linewidth}{0.5pt}\end{center}

\section{1.1 黒体放射とは何か ―
観測の言葉で}\label{sec-where-blackbody-definition}

{[}本文目安：B1{]}

「黒体放射」（\textbf{blackbody
radiation}）という言葉は、文脈によって二つの意味で使われる。一つは\textbf{観測される現象}、もう一つは\textbf{理想化された概念}である。両者は密接に関係しているが、混同すると後の議論で迷いが生じやすい。本書ではまず観測の言葉で導入する。

観測の立場から見れば、黒体放射は次の特徴を持つ\textbf{連続スペクトル}（\textbf{continuous
spectrum}）として現れる：

\begin{enumerate}
\def\labelenumi{\arabic{enumi}.}
\tightlist
\item
  \textbf{連続性}：振動数（\textbf{frequency}）\(\nu\)
  の関数として滑らかに分布し、特定の振動数に鋭いピークを持たない
\item
  \textbf{温度依存性}：たった一つの量 ―
  温度（\textbf{temperature}）\(T\) ― がスペクトル全体の形を決める
\item
  \textbf{普遍性}：放射する物体の材質・形・履歴に依存しない
\end{enumerate}

これら三つの特徴を持つスペクトルは、宇宙のさまざまな場所で観測される。以下、その実例を一つずつ眺めていく。各事例で重要なのは、「なぜそうなっているか」をすぐ問うのではなく、\textbf{まず観測される事実そのものを正確に把握すること}である。なぜそうなるか、つまり「物理」は、第II部以降で逆向きに辿っていく。

\begin{tcolorbox}[enhanced jigsaw, arc=.35mm, bottomrule=.15mm, bottomtitle=1mm, breakable, colback=white, colbacktitle=quarto-callout-note-color!10!white, colframe=quarto-callout-note-color-frame, coltitle=black, left=2mm, leftrule=.75mm, opacityback=0, opacitybacktitle=0.6, rightrule=.15mm, title=\textcolor{quarto-callout-note-color}{\faInfo}\hspace{0.5em}{ノート}, titlerule=0mm, toprule=.15mm, toptitle=1mm]

\textbf{用語に注意}：「黒体」(black body)
という言葉の「黒」は、入射光をすべて吸収する性質に由来する。完全吸収体は同時に最も効率の良い放出体になる（第3章の
Kirchhoff
の法則）。色が黒いという見た目の話ではない。実際、最も「黒体に近い」太陽は、目には眩しい白〜黄に見える。

\end{tcolorbox}

\section{1.2 太陽の連続スペクトル}\label{sec-where-blackbody-sun}

{[}本文目安：B1{]}

太陽は、地上から最も詳細に観測できる恒星である。可視光から近赤外にかけてのスペクトルを描くと、約
500 nm 付近にピークを持つ滑らかな連続スペクトルが現れる。

観測される事実：

\begin{itemize}
\tightlist
\item
  \textbf{ピーク波長}：約 500 nm（可視光の緑〜青緑領域。これは
  \(B_\lambda\) で見たときのピーク。\(B_\nu\)
  で見たときのピーク振動数は別の値に対応する ―
  振動数表示と波長表示でピーク位置が違うのが Wien 則のひと癖で、第2章
  §2.2-§2.3 で扱う）
\item
  \textbf{対応温度}：表面温度
  \(T_\odot \simeq 5778\ \mathrm{K}\)（有効温度。詳細は第17章）
\item
  \textbf{細部}：連続スペクトルの上に多数の細い暗線（Fraunhofer
  線）が刻まれている
\end{itemize}

太陽スペクトルは、ピークの位置とスペクトル全体の形状を見るだけなら、温度
5800 K
の黒体に非常によく似ている。ところが拡大して見ると、無数の吸収線が走っている。

この「\textbf{連続成分と線成分の二重構造}」は、後の章で詳しく扱う。本書では連続成分（黒体に近い側）を第I〜VI部で、線成分を第VII部で並行に追跡する。両者は別物に見えるが、第5章の放射輸送方程式の上では同じ枠組みで記述される（節~\ref{sec-radiative-transfer-bridge}
参照）。

\section{1.3 恒星の色と温度}\label{sec-where-blackbody-stars}

{[}本文目安：B1{]}

太陽以外の恒星も、それぞれ特有の「色」を持つ。スペクトル型（OBAFGKM）の順に温度が下がり、色がシフトする：

{\def\LTcaptype{none} % do not increment counter
\begin{longtable}[]{@{}
  >{\raggedright\arraybackslash}p{(\linewidth - 6\tabcolsep) * \real{0.2500}}
  >{\raggedright\arraybackslash}p{(\linewidth - 6\tabcolsep) * \real{0.2500}}
  >{\raggedright\arraybackslash}p{(\linewidth - 6\tabcolsep) * \real{0.2500}}
  >{\raggedright\arraybackslash}p{(\linewidth - 6\tabcolsep) * \real{0.2500}}@{}}
\toprule\noalign{}
\begin{minipage}[b]{\linewidth}\raggedright
スペクトル型
\end{minipage} & \begin{minipage}[b]{\linewidth}\raggedright
表面温度 {[}K{]}
\end{minipage} & \begin{minipage}[b]{\linewidth}\raggedright
見た目の色
\end{minipage} & \begin{minipage}[b]{\linewidth}\raggedright
代表的な恒星
\end{minipage} \\
\midrule\noalign{}
\endhead
\bottomrule\noalign{}
\endlastfoot
O & 30,000〜50,000 & 青 & \(\theta^1\) Ori
C（オリオン大星雲の励起源、O7） \\
B & 10,000〜30,000 & 青白 & Rigel \\
A & 7,500〜10,000 & 白 & Sirius \\
F & 6,000〜7,500 & 黄白 & Procyon \\
G & 5,200〜6,000 & 黄 & \textbf{太陽} \\
K & 3,700〜5,200 & 橙 & Aldebaran \\
M & 2,400〜3,700 & 赤 & Betelgeuse \\
\end{longtable}
}

「色」と「温度」が一対一に対応するのは、\textbf{すべての恒星が「黒体に近い」連続スペクトルを示すからこそ}成り立つ。温度が高いとピーク波長が短く（青側に）、低いと長く（赤側に）なる。この対応は、Wien
の法則として第2章で定量化する。

\begin{tcolorbox}[enhanced jigsaw, arc=.35mm, bottomrule=.15mm, bottomtitle=1mm, breakable, colback=white, colbacktitle=quarto-callout-note-color!10!white, colframe=quarto-callout-note-color-frame, coltitle=black, left=2mm, leftrule=.75mm, opacityback=0, opacitybacktitle=0.6, rightrule=.15mm, title=\textcolor{quarto-callout-note-color}{\faInfo}\hspace{0.5em}{ノート}, titlerule=0mm, toprule=.15mm, toptitle=1mm]

\textbf{注意（観測）}：実際の恒星は表層大気の組成・電離状態によって、黒体スペクトルからのずれが大きい。それでも、概略的な分類や全体光度の推定には黒体近似が極めて有効である。「\textbf{理想からのずれを定量化するために、まず理想がよく定まっている}」という構図に注意したい。

\end{tcolorbox}

\section{1.4 惑星・地球・ダストの熱放射}\label{sec-where-blackbody-cool}

{[}本文目安：B1{]}

宇宙の「冷たい側」も黒体放射を出す。低温では Wien
の法則（第2章）からピークが長波長側に来る：

{\def\LTcaptype{none} % do not increment counter
\begin{longtable}[]{@{}lll@{}}
\toprule\noalign{}
天体 & 温度 {[}K{]} & 主な放射帯 \\
\midrule\noalign{}
\endhead
\bottomrule\noalign{}
\endlastfoot
地球の昼面 & \(\sim 290\) & 中間赤外（10〜15 μm） \\
木星 & \(\sim 130\) & 遠赤外 \\
太陽系外縁天体 & 30〜50 & 遠赤外 \\
星間ダスト & 10〜30 & 遠赤外〜サブミリ \\
原始惑星系円盤 & 100〜数千 & 近赤外〜遠赤外 \\
\end{longtable}
}

赤外線天文衛星（Herschel・JWST
など）は、まさにこの低温熱放射を観測するために設計されている。

ただし、ダストは厳密には黒体ではない。波長によって\textbf{放射効率}（\textbf{emissivity}）が変化するため、\textbf{修正黒体}（\textbf{modified
blackbody}）と呼ばれる形に従う。修正の仕組み自体は第15章で扱うが、ここでは「現実の天体は、しばしば黒体スペクトルに\textbf{スペクトル指数}を掛けた形で見える」ことを覚えておけば十分である。

\section{1.5
宇宙マイクロ波背景放射（CMB）}\label{sec-where-blackbody-cmb}

{[}本文目安：B1{]}

宇宙のスペクトルの中で、最も完璧な黒体スペクトルを示すのは ―
\textbf{宇宙そのもの}である。

宇宙マイクロ波背景放射（CMB）の観測事実：

\begin{itemize}
\tightlist
\item
  \textbf{温度}：\(T_{\mathrm{CMB}} = 2.725 \pm 0.001\ \mathrm{K}\)
\item
  \textbf{ピーク波長}：約 1.06 mm（マイクロ波領域。\(B_\lambda\)
  ピーク表示）
\item
  \textbf{黒体スペクトルからのずれ（スペクトル歪み）}：全天平均スペクトルが完全黒体形からずれる量で、COBE/FIRAS（1990年代）は
  rms \(\Delta I/I \sim 5 \times 10^{-5}\)、ピーク値
  \(|\Delta I/I| < 10^{-4}\) という上限を与えた
\end{itemize}

\begin{tcolorbox}[enhanced jigsaw, arc=.35mm, bottomrule=.15mm, bottomtitle=1mm, breakable, colback=white, colbacktitle=quarto-callout-important-color!10!white, colframe=quarto-callout-important-color-frame, coltitle=black, left=2mm, leftrule=.75mm, opacityback=0, opacitybacktitle=0.6, rightrule=.15mm, title=\textcolor{quarto-callout-important-color}{\faExclamation}\hspace{0.5em}{重要}, titlerule=0mm, toprule=.15mm, toptitle=1mm]

\textbf{二種類の「ずれ」を混同しない}：CMB
には性質の異なる二つの「ずれ」がある。

\begin{itemize}
\tightlist
\item
  \textbf{スペクトル歪み}：全天平均した\textbf{スペクトルの形}（強度 vs
  振動数）が、完全黒体からどれだけずれるか。FIRAS
  の残差や、将来観測が狙う \(\mu\) 歪み・\(y\)
  歪みがこれにあたる。早期宇宙の熱化過程の不完全さを反映する。
\item
  \textbf{温度ゆらぎ（異方性）}：観測する\textbf{方向ごと}に黒体温度がわずかに違う量で、\(\Delta T/T \sim 10^{-5}\)。各方向のスペクトルは依然として黒体形であり、形そのもののずれではない。初期密度ゆらぎを反映する（第16章）。
\end{itemize}

本節で扱うのは前者（スペクトル歪み）である。後者の異方性は §1.9
と第16章で扱う。

\end{tcolorbox}

CMB の黒体性は驚異的である。観測精度の限界まで、全天平均の CMB
スペクトルは理論的な完全黒体スペクトルと一致する。これは「宇宙史上で最も完璧な黒体放射」と呼んでよい。

なぜこれほど完璧な黒体スペクトルが宇宙全体に存在するのか ―
この問いは本書全体の重要な動機の一つである。

熱化（thermalization）と自由伝搬の開始を分けて理解する必要がある：

\begin{itemize}
\tightlist
\item
  \textbf{熱化}：宇宙誕生から再結合期に至るまでの早い時代に、\textbf{二重
  Compton 散乱・制動放射・Compton
  散乱}などの光子と物質の強い相互作用によって、放射場のスペクトルが
  Planck 形に整えられた。この過程はもっと早期（典型的には宇宙年齢
  \(\lesssim 10^7\) 秒）に終わる。
\item
  \textbf{再結合期}（宇宙年齢 約 38
  万年）：電子と陽子が結合して中性水素になり、光が自由伝搬を始めた時期。これ以降、すでに完成していた黒体スペクトルが膨張で温度を下げながら（形を保ったまま）今日まで届く。
\end{itemize}

つまり「再結合期に黒体になった」というよりは「\textbf{再結合期までに完成した黒体スペクトルが、再結合期に解放された}」と理解するのが正確である。詳しくは第16章で扱う。

\begin{tcolorbox}[enhanced jigsaw, arc=.35mm, bottomrule=.15mm, bottomtitle=1mm, breakable, colback=white, colbacktitle=quarto-callout-tip-color!10!white, colframe=quarto-callout-tip-color-frame, coltitle=black, left=2mm, leftrule=.75mm, opacityback=0, opacitybacktitle=0.6, rightrule=.15mm, title=\textcolor{quarto-callout-tip-color}{\faLightbulb}\hspace{0.5em}{ヒント}, titlerule=0mm, toprule=.15mm, toptitle=1mm]

\textbf{観測との対応}：CMB
の黒体性は、宇宙が「熱平衡を経た過去」を持つことの直接的証拠である。これは
ビッグバン宇宙論の最強の観測的支柱の一つになっている。

\end{tcolorbox}

\begin{tcolorbox}[enhanced jigsaw, arc=.35mm, bottomrule=.15mm, breakable, colback=white, colframe=quarto-callout-tip-color-frame, left=2mm, leftrule=.75mm, opacityback=0, rightrule=.15mm, toprule=.15mm]
\begin{minipage}[t]{5.5mm}
\textcolor{quarto-callout-tip-color}{\faLightbulb}
\end{minipage}%
\begin{minipage}[t]{\textwidth - 5.5mm}

\textbf{問い}：なぜ CMB だけがこれほど完全な黒体スペクトルになるのか？

\textbf{短答}：宇宙が再結合期（誕生後約 38
万年）まで、放射場と物質が密に相互作用し、長時間・高密度の条件下で完璧な熱平衡を達成したから。光学的にきわめて厚い状況が、宇宙という最大スケールで実現していた。

\textbf{もう一歩}：再結合以降、放射と物質が分離（decoupling）したあとも、宇宙膨張は黒体スペクトルの「形」を保存する（断熱不変量としての性質）。だから現在の
2.725 K の CMB も完璧な黒体形を保つ。詳しくは第16章。

\end{minipage}%
\end{tcolorbox}

\section{1.6
降着円盤と高温天体の熱的放射}\label{sec-where-blackbody-hot}

{[}本文目安：B2{]}

激しい現象も、局所的にはしばしば黒体に近い。例：

{\def\LTcaptype{none} % do not increment counter
\begin{longtable}[]{@{}
  >{\raggedright\arraybackslash}p{(\linewidth - 4\tabcolsep) * \real{0.3333}}
  >{\raggedright\arraybackslash}p{(\linewidth - 4\tabcolsep) * \real{0.3333}}
  >{\raggedright\arraybackslash}p{(\linewidth - 4\tabcolsep) * \real{0.3333}}@{}}
\toprule\noalign{}
\begin{minipage}[b]{\linewidth}\raggedright
天体
\end{minipage} & \begin{minipage}[b]{\linewidth}\raggedright
温度 {[}K{]}
\end{minipage} & \begin{minipage}[b]{\linewidth}\raggedright
主な放射帯
\end{minipage} \\
\midrule\noalign{}
\endhead
\bottomrule\noalign{}
\endlastfoot
白色矮星表面 & \(\sim 10^4\) & 紫外〜可視 \\
中性子星表面 & \(10^6\)〜\(10^7\) & 軟 X 線 \\
降着円盤（連星系・AGN） & \(10^4\)〜\(10^8\) & 紫外〜硬 X 線 \\
ガンマ線バースト（プロンプト準熱的成分） & \(10^9\) 以上 &
ガンマ線（解釈は議論中） \\
\end{longtable}
}

ガンマ線バースト（GRB）の \textbf{残光（afterglow）}
はシンクロトロン優勢の \textbf{非熱的放射} であり、本表の対象外（§1.9
と整合）。本表に含めた GRB プロンプト相は、Band
関数で記述される連続光のうち準熱的成分の解釈に対応する。

降着円盤の場合、半径ごとに異なる温度の黒体スペクトルが合成される。これを
\textbf{多温度黒体} または \textbf{多温度円盤スペクトル}
と呼ぶ。各半径では「局所的に黒体近似」が成り立つ、というのが暗黙の仮定であり、第18章で詳しく扱う。

激しい場でも局所的に黒体近似が成り立つのは、\textbf{熱化時間}が\textbf{流体力学的時間スケール}より短い限り、放射場と物質が十分に相互作用して熱平衡に近づくからである。この「局所熱平衡」（LTE）という概念は、第5章の放射輸送と、第6章の熱平衡で本格的に扱う。

\section{1.7
「黒体」とは何が理想化されているのか}\label{sec-where-blackbody-idealization}

{[}本文目安：B2{]}

これまでの例を振り返ると、「黒体に近い」と言われる天体には、いずれも近似が入っていることがわかる。観測対象としての「黒体」と、概念としての「黒体」を、ここで明確に区別しておきたい。

\textbf{理想化された黒体に課されている三つの条件}：

\begin{enumerate}
\def\labelenumi{\arabic{enumi}.}
\tightlist
\item
  \textbf{完全吸収}：入射するすべての電磁波を、すべての振動数で 100\%
  吸収する
\item
  \textbf{完全熱化}：内部で十分に物質と光が相互作用し、放射場が単一の温度
  \(T\) で熱平衡に達している
\item
  \textbf{平衡温度の一意性}：物体全体が同じ温度 \(T\) で記述できる
\end{enumerate}

現実の天体は、これらの条件をどこまで満たすかで「黒体に近い」かが決まる：

{\def\LTcaptype{none} % do not increment counter
\begin{longtable}[]{@{}llll@{}}
\toprule\noalign{}
天体 & 条件 1（完全吸収） & 条件 2（熱化） & 条件 3（単一 T） \\
\midrule\noalign{}
\endhead
\bottomrule\noalign{}
\endlastfoot
太陽 & △（外層は不完全） & ◯（光球は熱化） & △（線吸収あり） \\
恒星一般 & △ & ◯ & △ \\
ダスト & ✗（修正黒体） & ◯ & △ \\
CMB & ◯ & ◯ & ◯ \\
降着円盤 & ◯（局所） & ◯（局所） & ✗（多温度） \\
\end{longtable}
}

つまり、\textbf{黒体は実在する物体ではなく、観測スペクトルを比較するための基準}である。本書では各章で、観測スペクトルがこの基準からどれだけ・なぜずれるかを順に解き明かしていく。とりわけ第VI部（第15--18章）は、まさに「黒体からのずれ」を主題にしている。

\section{1.8 連続と線 ―
本書が扱う二つの主柱}\label{sec-where-blackbody-two-pillars}

{[}本文目安：B1{]}

ここまで「連続スペクトル」を中心に話してきたが、上の例（特に 1.2
の太陽）で見たように、観測スペクトルにはしばしば \textbf{線スペクトル}
が重なっている。両者は本書の二つの主柱である：

\begin{itemize}
\tightlist
\item
  \textbf{連続スペクトル}（\textbf{continuous spectrum},
  本書の前半・第I〜VI部の主題）
  黒体放射、修正黒体、自由--自由放射（\textbf{free-free emission},
  bremsstrahlung）、シンクロトロン放射（\textbf{synchrotron
  radiation}）など、振動数 \(\nu\) の滑らかな関数として現れる成分
\item
  \textbf{線スペクトル}（\textbf{line spectrum},
  本書の後半・第VII部の主題） 特定の振動数 \(\nu_0\)
  にピークを持つ細い線として現れる成分。原子・分子の量子準位間の遷移に由来
\end{itemize}

両者は別物に見えるが、第II部の放射輸送方程式（第5章）の上では、\textbf{同じ枠組み}で記述される。違いは吸収係数
\(\alpha_\nu\) の構造だけである：

\begin{itemize}
\tightlist
\item
  連続吸収（\textbf{continuous absorption}）：\(\alpha_\nu\) が \(\nu\)
  に対して滑らかに変化
\item
  線吸収（\textbf{line absorption}）：\(\alpha_\nu\) が \(\nu_0\)
  の周りに鋭いピーク \(\phi(\nu - \nu_0)\) を持つ
\end{itemize}

詳しくは 節~\ref{sec-radiative-transfer-bridge}
で扱う。本書の最終章（第VIII部）では、QED
の言葉で連続と線が同じ電磁場--物質相互作用の二側面であることが明らかになる。

\section{1.9 熱的放射と非熱的放射 ―
本書の焦点}\label{sec-where-blackbody-thermal}

{[}本文目安：B1{]}

連続スペクトルにはもう一つ重要な区別がある。\textbf{熱的放射}（thermal
radiation）と\textbf{非熱的放射}（non-thermal
radiation）である。前節で「連続と線」を本書の二本柱として整理したが、その\textbf{連続側にもう一つの軸}が走っている。

\begin{itemize}
\tightlist
\item
  \textbf{熱的放射}：熱平衡（または
  LTE）にある物質が出す放射。スペクトル形は物質の温度 \(T\)
  だけで決まる。黒体放射が代表例。他に、Maxwell
  分布のプラズマからの\textbf{自由--自由放射}（熱的制動放射）、再結合放射、LTE
  下の線放射などが含まれる。
\item
  \textbf{非熱的放射}：粒子分布が熱平衡から大きく外れた状態 ―
  典型的には\textbf{べき則分布} \(f(E) \propto E^{-p}\) ―
  の高エネルギー粒子から出る放射。スペクトル形は温度ではなく粒子分布関数で決まる。\textbf{シンクロトロン放射}（磁場中の相対論的電子）、逆
  Compton 散乱、曲率放射、宇宙線起源の放射などが含まれる。
\end{itemize}

{\def\LTcaptype{none} % do not increment counter
\begin{longtable}[]{@{}lll@{}}
\toprule\noalign{}
区別 & スペクトル形を決める量 & 代表例 \\
\midrule\noalign{}
\endhead
\bottomrule\noalign{}
\endlastfoot
\textbf{熱的} & 温度 \(T\) のみ & 黒体放射、熱的自由--自由、LTE 線 \\
\textbf{非熱的} & 粒子分布 \(f(E)\) & シンクロトロン、逆
Compton、曲率放射 \\
\end{longtable}
}

\textbf{本書の焦点}：本書は\textbf{主として熱的放射}を扱う。プランク分布の三層構造（熱平衡・量子化・モード密度）を逆引きするのが第I〜V部、現実の天体での熱的放射の「ずれ」が第VI部、線スペクトル（多くは
LTE / non-LTE の枠で扱う）が第VII部、QED 統合が第VIII部。

非熱的放射は本書の中心テーマではないが、観測スペクトルの解釈に必要な範囲で触れる：

\begin{itemize}
\tightlist
\item
  第15章 §15.6 で「非熱的放射との区別」として
\item
  第18章でAGN・降着円盤・コンプトン化・シンクロトロンを軽く
\end{itemize}

非熱的放射の本格的扱いには、本書を超えた専門書（Rybicki \& Lightman
の第6章以降、Longair \emph{High Energy Astrophysics}
など）に進むのがよい。

\begin{tcolorbox}[enhanced jigsaw, arc=.35mm, bottomrule=.15mm, bottomtitle=1mm, breakable, colback=white, colbacktitle=quarto-callout-note-color!10!white, colframe=quarto-callout-note-color-frame, coltitle=black, left=2mm, leftrule=.75mm, opacityback=0, opacitybacktitle=0.6, rightrule=.15mm, title=\textcolor{quarto-callout-note-color}{\faInfo}\hspace{0.5em}{ノート}, titlerule=0mm, toprule=.15mm, toptitle=1mm]

\textbf{用語}：「熱的」「非熱的」は astro-dic.jp の標準訳。\textbf{LTE /
non-LTE} とは別の概念なので注意：

\begin{itemize}
\tightlist
\item
  \textbf{熱的 / 非熱的}：粒子の\textbf{エネルギー分布}が Maxwell
  分布か否か
\item
  \textbf{LTE /
  non-LTE}：原子の\textbf{準位占有数}が局所温度だけで決まるか、放射場
  \(J_\nu\) にも依存するか
\end{itemize}

非熱的電子（べき則）の集団が出すシンクロトロンの\textbf{周りにある原子ガス}は、独立に
LTE でも non-LTE でもありうる。両軸は直交する。

\end{tcolorbox}

\begin{tcolorbox}[enhanced jigsaw, arc=.35mm, bottomrule=.15mm, breakable, colback=white, colframe=quarto-callout-tip-color-frame, left=2mm, leftrule=.75mm, opacityback=0, rightrule=.15mm, toprule=.15mm]
\begin{minipage}[t]{5.5mm}
\textcolor{quarto-callout-tip-color}{\faLightbulb}
\end{minipage}%
\begin{minipage}[t]{\textwidth - 5.5mm}

\textbf{問い}：黒体放射は「熱的」だが、それなら CMB はどう分類すべきか？

\textbf{短答}：CMB
は完全に熱的。再結合期に光子と物質が完全な熱平衡に達し、その黒体スペクトルが膨張で温度を下げながら（形を保ったまま）今日まで残っている。

\textbf{もう一歩}：「2.725 K
黒体スペクトル」が宇宙最大スケールの熱的放射の典型例である。一方、CMB
の\textbf{温度ゆらぎ}（\(\Delta T/T \sim 10^{-5}\)）は初期宇宙の量子ゆらぎが膨張で増幅されたもので、\textbf{熱平衡の統計ゆらぎ起源ではない}が、各方向の\textbf{スペクトル形そのもの}は依然として熱的（プランク分布）のまま。「非熱的（粒子のエネルギー分布がべき則）」とは別の意味で「熱起源でない」点に注意
― \textbf{スペクトル（縦軸：強度 vs 振動数）} と
\textbf{空間ゆらぎ（方向ごとの温度の違い）} を混同しないことが重要。

\end{minipage}%
\end{tcolorbox}

\section{1.10 この本の読み方 ―
観測から理論へ逆にたどる}\label{sec-where-blackbody-howto}

{[}本文目安：B1{]}

通常の物理教科書は、まず統計力学・量子論・電磁気学を学び、その応用例として黒体放射を最後に扱う。本書はその順序を\textbf{反転}させる：

\begin{verbatim}
通常の教科書：
  基礎物理（統計・量子・電磁気）→ 応用：黒体放射

本書：
  観測：黒体スペクトル → 逆引き → 基礎物理
\end{verbatim}

具体的には次の流れになる：

\begin{enumerate}
\def\labelenumi{\arabic{enumi}.}
\tightlist
\item
  \textbf{第I部}（本章＋第2--3章）：観測の現象論
\item
  \textbf{第II部}（第4--5章）：放射場の記述
\item
  \textbf{第III部}（第6--8章）：熱平衡と統計力学への逆引き
\item
  \textbf{第IV部}（第9--12章）：量子論への逆引き
\item
  \textbf{第V部}（第13--14章）：電磁気学への逆引き
\item
  \textbf{第VI部}（第15--18章）：現実の宇宙へ戻る
\item
  \textbf{第VII部}（第20--24章）：線スペクトルの逆引き
\item
  \textbf{第VIII部}（第25--26章）：QED で連続と線を統一
\end{enumerate}

各部の冒頭には「\textbf{この部で答える問い}」と「\textbf{中心地図}」が再掲され、読者は常に自分が中心地図のどの因子を解き明かしている途中かを確認できる。

\begin{tcolorbox}[enhanced jigsaw, arc=.35mm, bottomrule=.15mm, bottomtitle=1mm, breakable, colback=white, colbacktitle=quarto-callout-tip-color!10!white, colframe=quarto-callout-tip-color-frame, coltitle=black, left=2mm, leftrule=.75mm, opacityback=0, opacitybacktitle=0.6, rightrule=.15mm, title=\textcolor{quarto-callout-tip-color}{\faLightbulb}\hspace{0.5em}{ヒント}, titlerule=0mm, toprule=.15mm, toptitle=1mm]

\textbf{読み方のコツ}：本書は通読できるが、目的によって読む順序を変えてもよい。表紙ページ（\hyperref[ux5b87ux5b99ux306eux71b1ux7684ux653eux5c04]{\texttt{index.qmd}}）の「読者タイプ別ルート」と、\hyperref[sec-prologue-reading-routes]{序章
§0.4} を参照。

\end{tcolorbox}

\begin{center}\rule{0.5\linewidth}{0.5pt}\end{center}

\section{計算例：プランク曲線を描いてみる}\label{sec-where-blackbody-numerical}

{[}本文目安：B1{]}

「同じ関数形が温度だけで形を変える」ことを実際に見てみよう。次のコードは、太陽（5800
K）、低温の恒星（3000 K）、人体（310 K）、CMB（2.7
K）のプランク曲線を比較する。

\begin{tcolorbox}[enhanced jigsaw, arc=.35mm, bottomrule=.15mm, bottomtitle=1mm, breakable, colback=white, colbacktitle=quarto-callout-note-color!10!white, colframe=quarto-callout-note-color-frame, coltitle=black, left=2mm, leftrule=.75mm, opacityback=0, opacitybacktitle=0.6, rightrule=.15mm, title=\textcolor{quarto-callout-note-color}{\faInfo}\hspace{0.5em}{ノート}, titlerule=0mm, toprule=.15mm, toptitle=1mm]

\textbf{注意}：プランク関数 \(B_\nu(T)\)
の具体的な式は本章では「天下り的には」与えない。式の構造と各因子の意味は第2章で扱う。ここでは、観測される多種多様な「黒体スペクトル」が一つの曲線族にすぎないことを、視覚的に確認するのが目的である。

\end{tcolorbox}

\begin{Shaded}
\begin{Highlighting}[]
\ImportTok{import}\NormalTok{ sys}\OperatorTok{;}\NormalTok{ sys.path.insert(}\DecValTok{0}\NormalTok{, }\StringTok{"../code"}\NormalTok{)}
\ImportTok{from}\NormalTok{ planck\_curve }\ImportTok{import}\NormalTok{ plot\_planck\_with\_limits}

\NormalTok{fig, ax }\OperatorTok{=}\NormalTok{ plot\_planck\_with\_limits(temperatures}\OperatorTok{=}\NormalTok{(}\FloatTok{2.7}\NormalTok{, }\DecValTok{310}\NormalTok{, }\DecValTok{3000}\NormalTok{, }\DecValTok{5800}\NormalTok{))}
\NormalTok{fig}
\end{Highlighting}
\end{Shaded}

\begin{figure}[H]

\centering{

\pandocbounded{\includegraphics[keepaspectratio]{part1/01-where-blackbody_files/figure-pdf/fig-planck-curves-output-1.png}}

}

\caption{\label{fig-planck-curves}多様な温度のプランク曲線。温度だけが形を決め、すべて同じ関数族に属する。}

\end{figure}%

ここで重要なのは、4つの曲線が異なる温度（4桁の範囲）にもかかわらず、すべて同じ関数族
\(B_\nu(T)\)
に属する、という一点である。本書はこの関数族を「中心地図」として、その背後の物理を逆引きで解き明かしていく。

\begin{center}\rule{0.5\linewidth}{0.5pt}\end{center}

\section*{この章で何がわかったか}\label{sec-where-blackbody-summary}
\addcontentsline{toc}{section}{この章で何がわかったか}

\markright{この章で何がわかったか}

\textbf{中心地図に戻る}

本章では、本書の中心地図であるプランク関数 \(B_\nu(T)\)
について、その\textbf{観測される現れ方}を整理した：

\begin{itemize}
\tightlist
\item
  太陽・恒星 ― 数千 K
\item
  惑星・ダスト ― 数十〜数百 K
\item
  CMB ― 2.7 K
\item
  降着円盤・中性子星 ― 数百万〜数億 K
\end{itemize}

すべてが同じ「黒体に近い」スペクトル形を示し、温度 \(T\)
だけが形を決めている。ただし「黒体」は実在する物体ではなく、観測スペクトルを比較するための理想化された基準である。

\textbf{次章へ}：第2章では、観測スペクトルから温度・半径・光度をどう読み取るかを、Wien
則・Stefan--Boltzmann
則・色温度などの道具を使って具体的に学ぶ。中心地図プランク関数のどの因子がどの観測量に対応するかが、第2章の中心テーマである。

\section*{演習問題}\label{sec-where-blackbody-exercises}
\addcontentsline{toc}{section}{演習問題}

\markright{演習問題}

\textbf{模範解答}
は折りたたみを展開して確認できる（オンライン版）。まず自力で解いてから確認すると効果的である。

\subsection*{問題 1-1　散らばったデータから Wien
変位定数を逆算する}\label{ex-1-1}
\addcontentsline{toc}{subsection}{問題 1-1　散らばったデータから Wien
変位定数を逆算する}

{[}★ 難易度：☆ {]}

本章は §1.2〜§1.5 で温度とピーク波長の組を別々に与えている（太陽
\(T_\odot\simeq5778\) K で \(\lambda_{\max}\simeq500\) nm、CMB
\(T=2.725\) K で \(\lambda_{\max}\simeq1.06\) mm、地球 \(T\simeq290\) K
で中間赤外
\(\sim10\ \mu\)m）。式を使わずに、これらが一つの法則に従うことを確かめる。

\begin{enumerate}
\def\labelenumi{\arabic{enumi}.}
\tightlist
\item
  3つの組について積 \(\lambda_{\max}\cdot T\)
  を計算し、ほぼ一定値になることを確かめよ。
\item
  その共通値（経験的な Wien 変位定数）\(b\) を概算せよ。
\item
  「温度が4桁違う天体が同じ \(b\)
  を共有する」という事実は、第1章の主張「すべて同じ関数族
  \(B_\nu(T)\)」をどう支持するか述べよ。
\end{enumerate}

\textbf{関連}：\hyperref[sec-where-blackbody-cmb]{§1.5
CMB}／変位則の正式な扱いは \hyperref[sec-reading-spectra-wien]{§2.3 Wien
の変位則}（導出は\hyperref[ex-2-2]{演習 2-2}）。

\section{模範解答（問題
1-1）}\label{ux6a21ux7bc4ux89e3ux7b54ux554fux984c-1-1}

\textbf{(1)} 単位を m·K に揃える：

{\def\LTcaptype{none} % do not increment counter
\begin{longtable}[]{@{}llll@{}}
\toprule\noalign{}
天体 & \(\lambda_{\max}\) & \(T\) {[}K{]} & \(\lambda_{\max}T\)
{[}m·K{]} \\
\midrule\noalign{}
\endhead
\bottomrule\noalign{}
\endlastfoot
太陽 & \(5.0\times10^{-7}\) m & 5778 & \(2.89\times10^{-3}\) \\
CMB & \(1.06\times10^{-3}\) m & 2.725 & \(2.89\times10^{-3}\) \\
地球 & \(1.0\times10^{-5}\) m & 290 & \(2.9\times10^{-3}\) \\
\end{longtable}
}

3桁の精度で同じ値。

\textbf{(2)} \(b\approx2.9\times10^{-3}\) m·K。これは §2.3 の Wien
変位定数 \(b=2.898\times10^{-3}\) m·K と一致する。

\textbf{(3)}
もし天体ごとにスペクトルの関数形が違えば、「ピーク波長×温度」が天体ごとにバラバラになるはず。4桁の温度差を超えて単一の
\(b\) を共有することは、ピーク位置の温度依存が全天体で共通＝同一の関数族
\(B_\nu(T)\)
が温度だけを変えて現れている、という第1章の主張を定量的に裏づける。

\textbf{答え}：\(\lambda_{\max}T\approx2.9\times10^{-3}\) m·K
で一定。単一の関数族が温度パラメータだけで全天体を記述することの証拠。

\subsection*{問題
1-2　理想黒体の3条件で現実天体を採点する}\label{ex-1-2}
\addcontentsline{toc}{subsection}{問題
1-2　理想黒体の3条件で現実天体を採点する}

{[}★ 難易度：☆☆ {]}

§1.7 は理想黒体に課す3条件（① 完全吸収、② 完全熱化、③
単一温度）を挙げ、天体ごとの達成度を表にした。各天体について「どの条件が破れ、その結果スペクトルにどんな『ずれ』が現れるか」を結びつける。

\begin{enumerate}
\def\labelenumi{\arabic{enumi}.}
\tightlist
\item
  \textbf{CMB}
  がほぼ完璧な黒体である理由を、3条件すべての観点から述べよ。
\item
  \textbf{降着円盤}
  が破る条件を1つ挙げ、その帰結である「多温度黒体」とは観測上どういうスペクトルになるか述べよ。
\item
  \textbf{星間ダスト}
  が破る条件を1つ挙げ、その帰結（修正黒体／放射効率の振動数依存）を述べよ。
\end{enumerate}

\textbf{関連}：\hyperref[sec-where-blackbody-idealization]{§1.7
理想化の3条件}／「ずれ」の整理は
\hyperref[sec-reading-spectra-deviation]{§2.7}、修正黒体は
\href{../part6/15-non-blackbody.qmd}{第15章}、多温度円盤は
\href{../part6/18-accretion-high-energy.qmd}{第18章}。

\section{模範解答（問題
1-2）}\label{ux6a21ux7bc4ux89e3ux7b54ux554fux984c-1-2}

\textbf{(1)} CMB は ①
再結合期以前に光と物質が密に相互作用して完全吸収・完全熱化に相当する状況を達成、②
長時間・高密度で完璧な熱平衡に到達、③ 宇宙全体が単一温度 2.725 K
で記述できる ― 3条件すべてを最もよく満たす。だから
\(\Delta I/I<10^{-4}\) という観測精度の限界まで黒体形からずれない。

\textbf{(2)}
降着円盤は条件③（単一温度）を破る。半径ごとに温度が異なり、各半径の
\(B_\nu(T(r))\) が重なる。観測スペクトルは単一の \(T\)
では再現できず、ピークが鈍り高振動数側に裾を引く「多温度黒体（多温度円盤スペクトル）」になる。各半径では局所的に黒体近似が成り立つことが前提。

\textbf{(3)} ダストは条件①（完全吸収）を破る。波長によって放射効率
\(\epsilon_\nu<1\)
が変化するため、黒体に振動数依存のスペクトル指数を掛けた「修正黒体」として観測される。

\textbf{答え}：CMB＝3条件すべて満たす。円盤＝③単一温度を破る→多温度合成。ダスト＝①完全吸収を破る→修正黒体。

\subsection*{問題 1-3　熱的/非熱的 と LTE/非LTE
は直交する}\label{ex-1-3}
\addcontentsline{toc}{subsection}{問題 1-3　熱的/非熱的 と LTE/非LTE
は直交する}

{[}★ 難易度：☆☆ {]}

§1.9 は「熱的/非熱的」（粒子のエネルギー分布が Maxwell
分布か）と「LTE/非LTE」（準位占有が局所温度だけで決まるか）が別概念で直交すると強調した。

\begin{enumerate}
\def\labelenumi{\arabic{enumi}.}
\tightlist
\item
  「熱的 vs 非熱的」を分ける物理量と、「LTE vs
  非LTE」を分ける物理量をそれぞれ一言で述べよ。
\item
  次の各放射を「熱的/非熱的」で分類せよ：熱的制動放射、シンクロトロン放射、CMB、LTE
  線放射。
\item
  CMB の温度ゆらぎ \(\Delta T/T\sim10^{-5}\)
  は「非熱的」放射ではない。なぜか。スペクトル（強度 vs
  振動数）と空間ゆらぎ（方向ごとの温度差）の違いに触れて説明せよ。
\end{enumerate}

\textbf{関連}：\hyperref[sec-where-blackbody-thermal]{§1.9
熱的放射と非熱的放射}／非熱的放射の応用は
\href{../part6/18-accretion-high-energy.qmd}{第18章}、LTE/非LTE は
\hyperref[sec-universality-thermal-equilibrium]{§3.2}。

\section{模範解答（問題
1-3）}\label{ux6a21ux7bc4ux89e3ux7b54ux554fux984c-1-3}

\textbf{(1)} 熱的/非熱的を分けるのは
\textbf{粒子のエネルギー分布}（Maxwell
分布なら熱的、べき則など非Maxwellなら非熱的）。LTE/非LTE を分けるのは
\textbf{原子の準位占有数}（局所温度 \(T\) だけで決まれば LTE、放射場
\(J_\nu\) にも依存すれば非LTE）。

\textbf{(2)} 熱的制動放射＝熱的（Maxwell
分布プラズマ）。シンクロトロン＝非熱的（べき則の相対論的電子）。CMB＝熱的（完全な熱平衡）。LTE
線放射＝熱的。

\textbf{(3)} CMB
は各方向で見ると依然として完全なプランク分布（熱的スペクトル）。温度ゆらぎとは「方向ごとに黒体温度がわずかに違う」という
\textbf{空間的} な差であって、各点の \textbf{スペクトル形}
が非熱的（べき則）になっているわけではない。「非熱的」は強度 vs
振動数のスペクトル形の話、ゆらぎは方向ごとの温度の話で、軸が異なる。ゆらぎの起源（初期宇宙の量子ゆらぎ）も熱平衡の統計ゆらぎとは別物。

\textbf{答え}：分布関数 vs
準位占有という別軸。制動・CMB・LTE線＝熱的、シンクロトロン＝非熱的。CMB
ゆらぎはスペクトルでなく空間温度差なので非熱的とは無関係。

\subsection*{問題 1-4　CMB
の黒体性：「熱化」と「自由伝搬の開始」を分ける}\label{ex-1-4}
\addcontentsline{toc}{subsection}{問題 1-4　CMB
の黒体性：「熱化」と「自由伝搬の開始」を分ける}

{[}★ 難易度：☆☆ {]}

§1.5 は「再結合期に黒体になった」という言い方は不正確で、「再結合期
\textbf{までに}
完成した黒体スペクトルが再結合期に解放された」と理解せよ、と注意している。

\begin{enumerate}
\def\labelenumi{\arabic{enumi}.}
\tightlist
\item
  「熱化（thermalization）」と「自由伝搬の開始（再結合・脱結合）」は、それぞれ宇宙史のいつ起き、何を担う過程か区別して述べよ。
\item
  熱化を担った物理過程を §1.5 から挙げよ。
\item
  再結合以降、宇宙膨張で温度が下がってもスペクトルが黒体形を保つのはなぜか（断熱不変としての性質に触れて）述べよ。
\end{enumerate}

\textbf{関連}：\hyperref[sec-where-blackbody-cmb]{§1.5
CMB}／本格的扱いは \href{../part6/16-cmb.qmd}{第16章
CMB}、熱平衡の条件は \hyperref[sec-universality-real-objects]{§3.6}。

\section{模範解答（問題
1-4）}\label{ux6a21ux7bc4ux89e3ux7b54ux554fux984c-1-4}

\textbf{(1)} \textbf{熱化}：宇宙ごく初期（典型的に宇宙年齢
\(\lesssim10^7\) 秒）に、放射場のスペクトルを Planck
形へ整える過程。\textbf{自由伝搬の開始（再結合期、約38万年）}：電子と陽子が結合して中性水素になり、光が散乱されずに直進できるようになる過程。前者がスペクトルの形を作り、後者がそれを「解放」する。両者は時代も役割も別。

\textbf{(2)} 二重 Compton 散乱、制動放射、Compton
散乱など、光子と物質の強い相互作用。これらが早期の高密度環境でスペクトルを
Planck 形に整えた。

\textbf{(3)}
宇宙膨張は黒体スペクトルの「形」を保存する（断熱不変量としての性質）。膨張で各光子の振動数が一様に赤方偏移し、温度パラメータ
\(T\) が下がるだけで、プランク分布の関数形 \(B_\nu(T)\)
はそのまま保たれる。だから現在 2.725 K でも完璧な黒体形を示す。

\textbf{答え}：熱化＝初期にスペクトルを作る、再結合＝それを解放する、で別過程。膨張は形を保存（断熱不変）するので低温でも黒体のまま。

\section*{さらに学ぶための参考文献}\label{sec-where-blackbody-further}
\addcontentsline{toc}{section}{さらに学ぶための参考文献}

\markright{さらに学ぶための参考文献}

\begin{itemize}
\tightlist
\item
  太陽スペクトルの詳細な観測：Stix, \emph{The Sun} (Springer, 2nd ed.,
  2004)
\item
  恒星のスペクトル分類の歴史：Kaler, \emph{Stars and their Spectra}
  (Cambridge, 2nd ed., 2011)
\item
  CMB の黒体性の発見：Partridge, \emph{3K: The Cosmic Microwave
  Background Radiation} (Cambridge, 1995)
\item
  観測天文学全般の入門：縣 秀彦・嶺重 慎 編『シリーズ現代の天文学』各巻
\end{itemize}

\chapter{第2章　黒体スペクトルをどう読むか}\label{ux7b2c2ux7ae0-ux9ed2ux4f53ux30b9ux30daux30afux30c8ux30ebux3092ux3069ux3046ux8aadux3080ux304b}

\textbf{この章の主題}：プランク分布の基本的な形と振る舞いを観測の言葉で読み解く。観測スペクトルから温度・半径・光度を引き出すための道具（Wien
の法則、Stefan--Boltzmann
則、色温度など）を、観測の必要から立ち上げていく。

本章では、第1章で「観測される事実」として並べた黒体スペクトルを、\textbf{どう測り、どう読むか}という観測の側に踏み込む。式は「なぜそうなるか」（第III〜V部）に踏み込む前に、まず「観測量との対応関係」だけで扱う。

\section*{この章の中心地図}\label{sec-reading-spectra-map}
\addcontentsline{toc}{section}{この章の中心地図}

\markright{この章の中心地図}

\textbf{プランク関数（連続スペクトルの地図）}

\begin{equation}\protect\phantomsection\label{eq-planck-function}{
B_\nu(T) = \frac{2 h \nu^3}{c^2} \cdot \frac{1}{e^{h\nu/kT} - 1}
}\end{equation}

{\def\LTcaptype{none} % do not increment counter
\begin{longtable}[]{@{}lll@{}}
\toprule\noalign{}
因子 & 物理的意味 & 結びつく部 \\
\midrule\noalign{}
\endhead
\bottomrule\noalign{}
\endlastfoot
\(B_\nu\) & 観測量（比強度） & 第II部 \\
\(1/(e^{h\nu/kT}-1)\) & 熱平衡・光子統計 & 第III部 \\
\(h\nu\) & エネルギー量子 & 第IV部 \\
\(\nu^3/c^2\) & 電磁場のモード密度 & 第V部 \\
黒体からのずれ & 現実の天体 & 第VI部 \\
\end{longtable}
}

\begin{tcolorbox}[enhanced jigsaw, arc=.35mm, bottomrule=.15mm, breakable, colback=white, colframe=quarto-callout-tip-color-frame, left=2mm, leftrule=.75mm, opacityback=0, rightrule=.15mm, toprule=.15mm]

\textbf{地図の使い方}：各部の冒頭でこの式が再掲され、その部で扱う因子が強調表示される。読者は自分が今どこを学んでいるかをこの式の上で確認できる。

\end{tcolorbox}

\begin{tcolorbox}[enhanced jigsaw, arc=.35mm, bottomrule=.15mm, bottomtitle=1mm, breakable, colback=white, colbacktitle=quarto-callout-tip-color!10!white, colframe=quarto-callout-tip-color-frame, coltitle=black, left=2mm, leftrule=.75mm, opacityback=0, opacitybacktitle=0.6, rightrule=.15mm, title=\textcolor{quarto-callout-tip-color}{\faLightbulb}\hspace{0.5em}{ヒント}, titlerule=0mm, toprule=.15mm, toptitle=1mm]

本章では、地図のプランク関数 \(B_\nu(T)\)
そのものを「観測される曲線」として扱う。各因子の物理的起源は第III〜V部で逆引きしていく。本章はその準備として、\textbf{曲線の形と読み方}を共通言語にする。

\end{tcolorbox}

\section*{この章で答える問い}\label{sec-reading-spectra-questions}
\addcontentsline{toc}{section}{この章で答える問い}

\markright{この章で答える問い}

\begin{itemize}
\tightlist
\item
  観測スペクトルのピーク位置から、その天体の温度がわかるのはなぜか
\item
  スペクトルの面積（全エネルギー）と温度の関係は、観測で何の役に立つのか
\item
  振動数表示 \(B_\nu\) と波長表示 \(B_\lambda\)
  でピーク位置がずれるのはなぜか
\item
  「色温度」「有効温度」「輝度温度」は何が違い、どれを使うべきか
\item
  観測されるスペクトルからどうやって温度・半径・光度を引き出すのか
\end{itemize}

\section*{到達目標}\label{sec-reading-spectra-goals}
\addcontentsline{toc}{section}{到達目標}

\markright{到達目標}

この章を読み終えると、読者は次のことができるようになる：

\begin{itemize}
\tightlist
\item
  プランク分布の形を見て、温度を一桁の精度で推定できる
\item
  ウィーンの法則と Stefan--Boltzmann
  則を、観測量から温度・光度を逆算する道具として使える
\item
  振動数表示と波長表示の違いを理解し、観測データを正しい単位系で扱える
\end{itemize}

\begin{center}\rule{0.5\linewidth}{0.5pt}\end{center}

\section{2.1 プランク分布の形}\label{sec-reading-spectra-shape}

{[}本文目安：B2{]}

第1章では、宇宙のさまざまな天体が「同じ関数族」に属する連続スペクトルを示すことを観察した。この関数族は、振動数
\(\nu\) の関数として書くと次の形をもつ：

\begin{equation}\protect\phantomsection\label{eq-planck-nu}{
B_\nu(T) = \frac{2 h \nu^3}{c^2} \cdot \frac{1}{e^{h\nu / kT} - 1}
}\end{equation}

ここで \(B_\nu\) は\textbf{比強度}（\textbf{specific
intensity}）、つまり「単位面積・単位時間・単位立体角・単位振動数あたりのエネルギー」である。単位は
erg s\(^{-1}\) cm\(^{-2}\) Hz\(^{-1}\) sr\(^{-1}\)（SI なら W m\(^{-2}\)
Hz\(^{-1}\)
sr\(^{-1}\)）。本章では「面を通過するエネルギーの目安」として簡略に用いるが、\textbf{正式な定義}
― 面要素を法線から \(\theta\) 傾けて通過する成分に対する \(\cos\theta\)
を含む形 ― は第4章 §4.1 で扱う。

\begin{tcolorbox}[enhanced jigsaw, arc=.35mm, bottomrule=.15mm, bottomtitle=1mm, breakable, colback=white, colbacktitle=quarto-callout-important-color!10!white, colframe=quarto-callout-important-color-frame, coltitle=black, left=2mm, leftrule=.75mm, opacityback=0, opacitybacktitle=0.6, rightrule=.15mm, title=\textcolor{quarto-callout-important-color}{\faExclamation}\hspace{0.5em}{重要}, titlerule=0mm, toprule=.15mm, toptitle=1mm]

\textbf{本章での扱い方}：率直に言えば、本章では
\(B_\nu\)（式~\ref{eq-planck-nu}）を\textbf{一度ひとつの経験式・観測曲線として受け入れる}。第1章で「観測される現象」として温度ごとに同じ形をとることを確認したので、いまはその観測される曲線を表す式が上のように書ける、という事実だけ使う。式の各因子（\(2h\nu^3/c^2\)、\(1/(e^{h\nu/kT}-1)\)）がなぜこの形になるかは、第8章（光子統計）、第10章（プランクの量子仮説）、第14章（モード密度）で順に逆引きする。

\end{tcolorbox}

この曲線の特徴を、まず観測の言葉で整理しよう：

\begin{enumerate}
\def\labelenumi{\arabic{enumi}.}
\tightlist
\item
  \textbf{低振動数側}：\(B_\nu \propto \nu^2\) で立ち上がる
\item
  \textbf{ピーク}：\(h\nu \simeq 2.82\, kT\) 付近に最大値
\item
  \textbf{高振動数側}：\(e^{-h\nu/kT}\) で指数関数的に減衰
\end{enumerate}

温度を変えると、曲線全体がピークごと平行移動し、面積が急激に変わる。

これら観測で測る三つの特徴は、プランク関数 式~\ref{eq-planck-nu}
の三つの因子に一対一で対応する。本章は「観測から式へ」の対応を読む章なので、その対応をあらかじめ表にしておく：

{\def\LTcaptype{none} % do not increment counter
\begin{longtable}[]{@{}
  >{\raggedright\arraybackslash}p{(\linewidth - 4\tabcolsep) * \real{0.3333}}
  >{\raggedright\arraybackslash}p{(\linewidth - 4\tabcolsep) * \real{0.3333}}
  >{\raggedright\arraybackslash}p{(\linewidth - 4\tabcolsep) * \real{0.3333}}@{}}
\toprule\noalign{}
\begin{minipage}[b]{\linewidth}\raggedright
観測で測る特徴
\end{minipage} & \begin{minipage}[b]{\linewidth}\raggedright
プランク関数の因子
\end{minipage} & \begin{minipage}[b]{\linewidth}\raggedright
逆引きする部
\end{minipage} \\
\midrule\noalign{}
\endhead
\bottomrule\noalign{}
\endlastfoot
低振動数の傾き（\(\propto \nu^2\)） & 前因子 \(2h\nu^3/c^2\)
のうちモード密度由来の \(\nu^2\) & 第V部（モード密度） \\
ピーク位置（\(h\nu \simeq 2.82\,kT\)） & 指数因子 \(1/(e^{h\nu/kT}-1)\)
と前因子の競合 & 第III・IV部（光子統計・量子化） \\
Wien 側の指数減衰（\(e^{-h\nu/kT}\)） & 指数因子の高振動数極限 &
第IV部（エネルギー量子化） \\
\end{longtable}
}

各因子の物理的起源を逆引きするのが第III〜V部であり、本章はその「読み方」を先に固める。

\begin{Shaded}
\begin{Highlighting}[]
\ImportTok{import}\NormalTok{ sys}\OperatorTok{;}\NormalTok{ sys.path.insert(}\DecValTok{0}\NormalTok{, }\StringTok{"../code"}\NormalTok{)}
\ImportTok{from}\NormalTok{ planck\_curve }\ImportTok{import}\NormalTok{ plot\_planck\_with\_limits}

\NormalTok{fig, ax }\OperatorTok{=}\NormalTok{ plot\_planck\_with\_limits(temperatures}\OperatorTok{=}\NormalTok{(}\DecValTok{3000}\NormalTok{, }\DecValTok{5800}\NormalTok{, }\DecValTok{10000}\NormalTok{))}
\NormalTok{fig}
\end{Highlighting}
\end{Shaded}

\begin{figure}[H]

\centering{

\pandocbounded{\includegraphics[keepaspectratio]{part1/02-reading-spectra_files/figure-pdf/fig-planck-three-temps-output-1.png}}

}

\caption{\label{fig-planck-three-temps}三つの温度のプランク曲線。温度上昇でピークは高振動数側に移り、面積（全エネルギー）は
\(T^4\) に比例して大きくなる。}

\end{figure}%

\section{2.2 振動数表示と波長表示 ―
同じ物理、異なる関数形}\label{sec-reading-spectra-nu-lambda}

{[}本文目安：B3｜導出：B3{]}

スペクトルは振動数 \(\nu\) で表すこともあれば、波長 \(\lambda\)
で表すこともある。同じ物理現象なのだから、同じ放射エネルギーを表しているはずだが、\(B_\nu\)
と \(B_\lambda\) の関数形は異なる。

両者は\textbf{同じエネルギー}を\textbf{異なる変数の刻み}で表現している：

\begin{equation}\protect\phantomsection\label{eq-nu-lambda-conversion}{
B_\nu \, d\nu = -B_\lambda \, d\lambda
}\end{equation}

ここで \(\nu = c/\lambda\) から \(d\nu = -c/\lambda^2 \, d\lambda\)
なので、

\begin{equation}\protect\phantomsection\label{eq-planck-lambda}{
B_\lambda(T) = B_\nu(T) \cdot \frac{c}{\lambda^2} = \frac{2 h c^2}{\lambda^5} \cdot \frac{1}{e^{hc/\lambda kT} - 1}
}\end{equation}

ジャコビアン \(c/\lambda^2\) がかかることで、\(B_\lambda\) の関数形は
\(B_\nu\) と異なる。

\textbf{ピーク位置のずれ}：

{\def\LTcaptype{none} % do not increment counter
\begin{longtable}[]{@{}
  >{\raggedright\arraybackslash}p{(\linewidth - 4\tabcolsep) * \real{0.3333}}
  >{\raggedright\arraybackslash}p{(\linewidth - 4\tabcolsep) * \real{0.3333}}
  >{\raggedright\arraybackslash}p{(\linewidth - 4\tabcolsep) * \real{0.3333}}@{}}
\toprule\noalign{}
\begin{minipage}[b]{\linewidth}\raggedright
表示
\end{minipage} & \begin{minipage}[b]{\linewidth}\raggedright
ピーク条件
\end{minipage} & \begin{minipage}[b]{\linewidth}\raggedright
太陽（\(T = 5778\) K）での値
\end{minipage} \\
\midrule\noalign{}
\endhead
\bottomrule\noalign{}
\endlastfoot
\(B_\nu\) & \(h\nu / kT \simeq 2.82\) &
\(\nu_{\max} \simeq 3.4 \times 10^{14}\) Hz（\(\lambda \simeq 880\)
nm） \\
\(B_\lambda\) & \(hc/\lambda kT \simeq 4.97\) &
\(\lambda_{\max} \simeq 502\) nm \\
\end{longtable}
}

この\textbf{「ピーク位置がずれる」}現象は、観測のミスを生みやすい。同じ温度の太陽について、振動数で見ると「赤外にピーク」、波長で見ると「可視（緑）にピーク」となる。両者に矛盾はない。\textbf{スペクトルのピーク位置を語るときは、必ずどちらの変数で語っているかを明示する}のが鉄則である。

\begin{tcolorbox}[enhanced jigsaw, arc=.35mm, bottomrule=.15mm, bottomtitle=1mm, breakable, colback=white, colbacktitle=quarto-callout-warning-color!10!white, colframe=quarto-callout-warning-color-frame, coltitle=black, left=2mm, leftrule=.75mm, opacityback=0, opacitybacktitle=0.6, rightrule=.15mm, title=\textcolor{quarto-callout-warning-color}{\faExclamationTriangle}\hspace{0.5em}{警告}, titlerule=0mm, toprule=.15mm, toptitle=1mm]

\textbf{本書の流儀}：本書では特に断らない限り、振動数表示 \(B_\nu(T)\)
を主とする。理由は、放射輸送方程式やモード密度の議論（第II・V部）が振動数表示で素直に書けるからである。波長表示が便利な場面（観測装置の特性、可視光天文学）では適宜変換する。

\end{tcolorbox}

\section{2.3 ピークと温度 ― Wien
の変位則}\label{sec-reading-spectra-wien}

{[}本文目安：B2{]}

式~\ref{eq-planck-lambda} のピーク条件から、波長表示のピーク
\(\lambda_{\max}\) と温度 \(T\) の間に簡単な関係が成り立つ：

\begin{equation}\protect\phantomsection\label{eq-wien-displacement}{
\lambda_{\max} \cdot T = b \simeq 2.898 \times 10^{-3}\ \mathrm{m \cdot K}
}\end{equation}

これが \textbf{Wien の変位則}（\textbf{Wien's displacement
law}）である。\(b\) は Wien の変位定数。

実例：

{\def\LTcaptype{none} % do not increment counter
\begin{longtable}[]{@{}lll@{}}
\toprule\noalign{}
天体 & 温度 \(T\) {[}K{]} & \(\lambda_{\max}\) \\
\midrule\noalign{}
\endhead
\bottomrule\noalign{}
\endlastfoot
CMB & 2.7 & \(\sim 1.07\) mm（マイクロ波） \\
星間ダスト & 30 & \(\sim 100\) μm（遠赤外） \\
木星 & 130 & \(\sim 22\) μm（中間赤外） \\
地球 & 290 & \(\sim 10\) μm（中間赤外） \\
太陽 & 5800 & \(\sim 500\) nm（可視） \\
O 型星 & 30000 & \(\sim 100\) nm（紫外） \\
\end{longtable}
}

\textbf{観測との対応}：可視光から赤外までで観測した SED（spectral energy
distribution）のピーク波長を読み取れば、Wien
の式で逆算するだけで、温度を一桁の精度で推定できる。これが天文学で最もよく使う「色から温度」のステップである。

\begin{tcolorbox}[enhanced jigsaw, arc=.35mm, bottomrule=.15mm, bottomtitle=1mm, breakable, colback=white, colbacktitle=quarto-callout-tip-color!10!white, colframe=quarto-callout-tip-color-frame, coltitle=black, left=2mm, leftrule=.75mm, opacityback=0, opacitybacktitle=0.6, rightrule=.15mm, title=\textcolor{quarto-callout-tip-color}{\faLightbulb}\hspace{0.5em}{ヒント}, titlerule=0mm, toprule=.15mm, toptitle=1mm]

\textbf{振動数表示の Wien}：振動数表示のピーク位置
\(\nu_{\max}/T \simeq 5.88 \times 10^{10}\) Hz/K も Wien
の式の別表現として有用。ただし上述のとおり、波長表示と振動数表示でピーク位置の温度依存は同じ「逆比例」だが、定数の値が違うので両者を混ぜないように。

\end{tcolorbox}

\section{2.4 面積と全放射 ― Stefan--Boltzmann
則}\label{sec-reading-spectra-stefan-boltzmann}

{[}本文目安：B2{]}

プランク曲線 \(B_\nu(T)\)
から、黒体表面の単位面積・単位時間あたりに放出される全エネルギー（全放射フラックス）が得られる。手順は二段階：

\textbf{(i) 振動数別フラックス
\(F_\nu\)}：黒体表面は等方放出をするが、観測者がいるのは\textbf{外向きの半球}である。表面から半球の全方向に立体角積分すると、射影因子
\(\cos\theta\) がかかって（式~\ref{eq-flux-def} 参照）

\begin{equation}\protect\phantomsection\label{eq-fnu-pi-bnu}{
F_\nu = \int_{\text{半球}} B_\nu \cos\theta\, d\Omega = \pi B_\nu(T)
}\end{equation}

ここで「全立体角 \(4\pi\) ではなく
\(\pi\)」になるのは、\textbf{半球積分}かつ\textbf{\(\cos\theta\)
射影}の二つの効果（\(\int_0^{2\pi} d\varphi \int_0^{\pi/2} \cos\theta \sin\theta\, d\theta = \pi\)）から。第II部・第4章で改めて整理する。

\textbf{(ii)
全振動数積分}：\(F = \int_0^\infty F_\nu d\nu = \pi \int_0^\infty B_\nu d\nu\)
を実行すると

\begin{equation}\protect\phantomsection\label{eq-stefan-boltzmann}{
F = \sigma T^4
}\end{equation}

ここで \(\sigma = 5.670 \times 10^{-5}\) erg s\(^{-1}\) cm\(^{-2}\)
K\(^{-4}\) は \textbf{Stefan--Boltzmann 定数}（\textbf{Stefan--Boltzmann
constant}）である。

「\(T^4\)」という強い温度依存性は天文学のスケール感を支配する：

\begin{itemize}
\tightlist
\item
  太陽（\(T_\odot = 5778\)
  K）の単位面積あたりフラックス：\(\sigma T_\odot^4 \simeq 6.3 \times 10^{10}\)
  erg s\(^{-1}\) cm\(^{-2}\)
\item
  2 倍温度の星（\(T = 11560\) K）：単位面積あたり \(2^4 = 16\) 倍
\item
  半分温度の星（\(T = 2900\) K）：単位面積あたり \(1/16\)
\end{itemize}

光度（luminosity, 単位時間に天体全体から放出されるエネルギー）は、半径
\(R\) の球を仮定すれば

\begin{equation}\protect\phantomsection\label{eq-luminosity-sb}{
L = 4\pi R^2 \cdot \sigma T^4
}\end{equation}

これと観測される光度から、半径と温度の組み合わせ \(R^2 T^4\)
を読み出せる。距離が分かっていれば、観測されるフラックス \(f\)
と光度の関係 \(L = 4\pi d^2 f\) から、独立に \(L\) を決められる。

\begin{tcolorbox}[enhanced jigsaw, arc=.35mm, bottomrule=.15mm, bottomtitle=1mm, breakable, colback=white, colbacktitle=quarto-callout-note-color!10!white, colframe=quarto-callout-note-color-frame, coltitle=black, left=2mm, leftrule=.75mm, opacityback=0, opacitybacktitle=0.6, rightrule=.15mm, title=\textcolor{quarto-callout-note-color}{\faInfo}\hspace{0.5em}{ノート}, titlerule=0mm, toprule=.15mm, toptitle=1mm]

\textbf{\(T^4\) の物理}：Stefan--Boltzmann 則の温度依存 \(T^4\)
は、本章では「観測事実」として受け入れる。理由は第6章（熱平衡における放射場のエネルギー密度
\(u = aT^4\)）と第8章（光子の統計力学）で逆引きする。

\end{tcolorbox}

\section{2.5 低振動数極限と高振動数極限 ―
古典側と量子側}\label{sec-reading-spectra-limits}

{[}本文目安：B3｜導出：B3{]}

プランク分布の両端の極限は、それぞれ別個の式として古典物理の中で扱われていた。本書ではこれを「同じプランク分布の両極限」として統一的に見る。

\subsection{低振動数極限：Rayleigh--Jeans
則}\label{ux4f4eux632fux52d5ux6570ux6975ux9650rayleighjeans-ux5247}

\(h\nu \ll kT\) では、\(e^{h\nu/kT} - 1 \simeq h\nu/kT\) なので

\begin{equation}\protect\phantomsection\label{eq-rayleigh-jeans}{
B_\nu(T) \underset{h\nu \ll kT}{\longrightarrow} \frac{2 \nu^2}{c^2} kT
}\end{equation}

ここに \textbf{\(h\)
が現れていない}ことに注意。これは古典物理から（量子論なしで）導ける式である。電波天文学ではほぼ常にこの極限で観測している。

\subsection{高振動数極限：Wien
の式}\label{ux9ad8ux632fux52d5ux6570ux6975ux9650wien-ux306eux5f0f}

\(h\nu \gg kT\) では、\(e^{h\nu/kT} - 1 \simeq e^{h\nu/kT}\) なので

\begin{equation}\protect\phantomsection\label{eq-wien-limit}{
B_\nu(T) \underset{h\nu \gg kT}{\longrightarrow} \frac{2 h \nu^3}{c^2} e^{-h\nu/kT}
}\end{equation}

こちらは指数関数で急激に落ちる側。\textbf{\(h\)
が現れる}ことが本質的で、古典物理だけでは出てこない。

\subsection{この二つの極限から見える「物理の対立」}\label{ux3053ux306eux4e8cux3064ux306eux6975ux9650ux304bux3089ux898bux3048ux308bux7269ux7406ux306eux5bfeux7acb}

低振動数では \(h\) なしで OK、高振動数では \(h\) が本質的に効く ―
この事実は、\textbf{古典統計と量子の境目}を観測スペクトルが直接示していることを意味する。第IV部・第9章で「古典論ではなぜ失敗するのか（紫外線破綻）」を本格的に扱うが、そこへの伏線がここに置かれている。

\begin{tcolorbox}[enhanced jigsaw, arc=.35mm, bottomrule=.15mm, bottomtitle=1mm, breakable, colback=white, colbacktitle=quarto-callout-tip-color!10!white, colframe=quarto-callout-tip-color-frame, coltitle=black, left=2mm, leftrule=.75mm, opacityback=0, opacitybacktitle=0.6, rightrule=.15mm, title=\textcolor{quarto-callout-tip-color}{\faLightbulb}\hspace{0.5em}{ヒント}, titlerule=0mm, toprule=.15mm, toptitle=1mm]

\textbf{目安}：\(h\nu/kT \sim 1\)、つまり \(\nu \sim kT/h\)
を境に、Rayleigh--Jeans 側か Wien
側かが切り替わる。太陽光（\(T \sim 5800\) K）なら境界は
\(\nu \sim 10^{14}\) Hz（近赤外）。電波域は完全に RJ
側、可視光は遷移域、紫外〜X 線は Wien 側。

\end{tcolorbox}

\begin{tcolorbox}[enhanced jigsaw, arc=.35mm, bottomrule=.15mm, breakable, colback=white, colframe=quarto-callout-tip-color-frame, left=2mm, leftrule=.75mm, opacityback=0, rightrule=.15mm, toprule=.15mm]
\begin{minipage}[t]{5.5mm}
\textcolor{quarto-callout-tip-color}{\faLightbulb}
\end{minipage}%
\begin{minipage}[t]{\textwidth - 5.5mm}

\textbf{問い}：なぜ古典論で導けるのは Rayleigh--Jeans 側だけなのか？

\textbf{短答}：古典論ではエネルギー等分配定理が成り立ち、放射場の各モードに
\(kT\) のエネルギーが分配される。これがそのまま
\(B_\nu \propto \nu^2\, kT\) という RJ 形を生む。一方、Wien
側で必要な指数関数的減衰は、エネルギーが \(h\nu\)
という離散単位でやりとりされるという量子的性質がないと出てこない。

\textbf{もう一歩}：高振動数側で古典論を信じ続けると、全エネルギーが発散する。これが「\textbf{紫外線破綻}」で、観測される黒体スペクトルの形と決定的に矛盾する。この矛盾こそが量子論を要請した歴史的契機である。詳しくは第9章。

\end{minipage}%
\end{tcolorbox}

\section{2.6 色温度・有効温度・輝度温度 ―
どの温度か}\label{sec-reading-spectra-temperatures}

{[}本文目安：B3｜導出：B3{]}

「天体の温度」と言っても、現実の天体は完全な黒体ではない。観測値から「温度」を引き出す手段はいくつかあり、それぞれ別物として扱う必要がある。

{\def\LTcaptype{none} % do not increment counter
\begin{longtable}[]{@{}
  >{\raggedright\arraybackslash}p{(\linewidth - 4\tabcolsep) * \real{0.3333}}
  >{\raggedright\arraybackslash}p{(\linewidth - 4\tabcolsep) * \real{0.3333}}
  >{\raggedright\arraybackslash}p{(\linewidth - 4\tabcolsep) * \real{0.3333}}@{}}
\toprule\noalign{}
\begin{minipage}[b]{\linewidth}\raggedright
温度
\end{minipage} & \begin{minipage}[b]{\linewidth}\raggedright
定義
\end{minipage} & \begin{minipage}[b]{\linewidth}\raggedright
用途
\end{minipage} \\
\midrule\noalign{}
\endhead
\bottomrule\noalign{}
\endlastfoot
\textbf{色温度} \(T_c\)（\textbf{color temperature}） &
観測スペクトルの\textbf{形}（特に二点の比）が一致する黒体の温度 &
連続スペクトル形を黒体と比較するとき \\
\textbf{有効温度} \(T_\mathrm{eff}\)（\textbf{effective temperature}） &
全光度が等しくなる黒体の温度。\(L = 4\pi R^2 \sigma T_\mathrm{eff}^4\)
で定義 & 恒星の全放射エネルギーを基準にするとき \\
\textbf{輝度温度} \(T_b\)（\textbf{brightness temperature}） &
ある振動数 \(\nu\) で観測される比強度 \(I_\nu\)
が、その振動数のプランク関数と等しくなる温度。\(I_\nu = B_\nu(T_b)\) &
電波観測でほぼ常用 \\
\end{longtable}
}

完全な黒体ならこれら三つはすべて一致する。\textbf{ずれの大きさが、その天体がどれだけ黒体からはずれているかの指標}にもなる。

\begin{tcolorbox}[enhanced jigsaw, arc=.35mm, bottomrule=.15mm, bottomtitle=1mm, breakable, colback=white, colbacktitle=quarto-callout-note-color!10!white, colframe=quarto-callout-note-color-frame, coltitle=black, left=2mm, leftrule=.75mm, opacityback=0, opacitybacktitle=0.6, rightrule=.15mm, title=\textcolor{quarto-callout-note-color}{\faInfo}\hspace{0.5em}{ノート}, titlerule=0mm, toprule=.15mm, toptitle=1mm]

\textbf{注意（観測）}：太陽の場合、色温度・有効温度・輝度温度はおおよそ似た値（5500〜5800
K）になるが、波長帯によって 100 K
程度のずれが現れる。これは「太陽が完全な黒体ではない」ことの定量的な証拠であり、第15章（恒星大気と吸収線）でこのずれの起源を扱う。

\end{tcolorbox}

\section{2.7 黒体からのずれ ―
観測スペクトルが教えてくれるもの}\label{sec-reading-spectra-deviation}

{[}本文目安：B3｜導出：B3{]}

第1章 1.7
でみたように、現実の天体スペクトルは黒体から多少なりともずれる。観測で出会う主な「ずれ」を整理しておく：

\begin{enumerate}
\def\labelenumi{\arabic{enumi}.}
\tightlist
\item
  \textbf{吸収線・輝線}：連続スペクトルに細い線が重なる（第VII部）
\item
  \textbf{連続不透明度の振動数依存性}：滑らかな形が振動数によって変形（第15章、修正黒体）
\item
  \textbf{多温度合成}：複数の温度成分が重なって、単一の \(T\)
  では再現できない（第17・18章）
\item
  \textbf{光学的に薄い領域からの放射}：黒体まで熱化していない（第5章、第8章）
\item
  \textbf{非熱的放射の混入}：シンクロトロン（磁場中の相対論的電子）、逆
  Compton
  散乱、曲率放射など。\textbf{べき則分布}の高エネルギー電子から出るため、温度
  \(T\) ではなく粒子分布で決まる。第1章 §1.9
  で見たように、これは「黒体からの単なる量的ずれ」を超えた\textbf{質的に異なる放射}である（第15章
  §15.6、第18章）
\end{enumerate}

\textbf{重要なポイント}は、1〜4
のずれを測ることが\textbf{観測天文学の主要な仕事}である、ということである。黒体スペクトルは「比較の基準」として、ずれの定量化を可能にしてくれる。5
の非熱的放射は本書の中心テーマではないが、観測スペクトルに重なる場面では明示的に区別する必要がある。

\section{2.8
観測データから温度・半径・光度を推定する}\label{sec-reading-spectra-estimation}

{[}本文目安：B3｜導出：B3{]}

ここまでの道具を使って、実際に観測データから物理量を引き出す手順を整理する。

\textbf{典型的なワークフロー}（恒星を例に）：

\begin{enumerate}
\def\labelenumi{\arabic{enumi}.}
\tightlist
\item
  \textbf{観測}：複数の波長帯でフラックス \(f_\nu\) または \(f_\lambda\)
  を測る（SED の構築）
\item
  \textbf{温度の推定}：

  \begin{itemize}
  \tightlist
  \item
    Wien の式（式~\ref{eq-wien-displacement}）でピーク波長から \(T\)
    を概算、または
  \item
    色（バンド間のフラックス比）を黒体と比較して色温度 \(T_c\) を決定
  \end{itemize}
\item
  \textbf{距離の特定}：年周視差、その他の距離指標を使う
\item
  \textbf{光度の計算}：観測フラックスと距離から \(L = 4\pi d^2 f\)
\item
  \textbf{半径の推定}：Stefan--Boltzmann (式~\ref{eq-luminosity-sb})
  を逆算して \(R = \sqrt{L / (4\pi \sigma T^4)}\)
\end{enumerate}

この一連の手順は、\textbf{プランク分布が中心地図として}機能しているからこそ意味を持つ。観測量と物理量を結ぶ橋渡しが、すべて
\(B_\nu(T)\) の形（とそのモーメント）に集約される。

\begin{tcolorbox}[enhanced jigsaw, arc=.35mm, bottomrule=.15mm, bottomtitle=1mm, breakable, colback=white, colbacktitle=quarto-callout-tip-color!10!white, colframe=quarto-callout-tip-color-frame, coltitle=black, left=2mm, leftrule=.75mm, opacityback=0, opacitybacktitle=0.6, rightrule=.15mm, title=\textcolor{quarto-callout-tip-color}{\faLightbulb}\hspace{0.5em}{ヒント}, titlerule=0mm, toprule=.15mm, toptitle=1mm]

\textbf{例（数値）：太陽}

\begin{itemize}
\tightlist
\item
  観測されるピーク波長：約 500 nm → Wien の式から \(T \simeq 5800\) K
\item
  観測される地球での太陽定数：\(f_\odot \simeq 1.4 \times 10^6\) erg
  s\(^{-1}\) cm\(^{-2}\)
\item
  距離：\(d = 1\) AU \(= 1.5 \times 10^{13}\) cm
\item
  光度：\(L_\odot = 4\pi d^2 f_\odot \simeq 3.8 \times 10^{33}\) erg/s
\item
  半径：\(R_\odot = \sqrt{L_\odot / (4\pi \sigma T^4)} \simeq 7 \times 10^{10}\)
  cm
\end{itemize}

実測値ともよく一致する。これが「太陽が黒体に近い」ことの定量的な証拠になっている。

\end{tcolorbox}

\begin{center}\rule{0.5\linewidth}{0.5pt}\end{center}

\section*{この章で何がわかったか}\label{sec-reading-spectra-summary}
\addcontentsline{toc}{section}{この章で何がわかったか}

\markright{この章で何がわかったか}

\textbf{中心地図に戻る}

本章で、中心地図のプランク関数 \(B_\nu(T)\)
について、\textbf{観測の側からの「読み方」が一通り定まった}：

\begin{itemize}
\tightlist
\item
  \textbf{ピーク}（Wien）→ 温度を引き出す
\item
  \textbf{面積}（Stefan--Boltzmann）→ 光度と半径を引き出す
\item
  \textbf{両極限}（RJ・Wien）→ 古典と量子の境目を見せる
\item
  \textbf{三種の温度}（\(T_c, T_\mathrm{eff}, T_b\)）→
  黒体からのずれを定量化する
\end{itemize}

ただし「なぜそうなるか」は、まだ何も答えていない。なぜプランク分布の前因子が
\(2h\nu^3/c^2\) なのか、なぜ指数関数因子が \(1/(e^{h\nu/kT}-1)\)
なのか、なぜ Stefan--Boltzmann の指数が \(T^4\)
なのか。これらは第III部（統計力学）以降の逆引きで明らかになる。

\textbf{次章へ}：第3章では、もう一つの観測事実
―「\textbf{なぜ材質によらず同じスペクトル形になるのか}」―
を扱う。この「普遍性」が、第III部の熱平衡・統計力学への入口になる。

\section*{演習問題}\label{sec-reading-spectra-exercises}
\addcontentsline{toc}{section}{演習問題}

\markright{演習問題}

\textbf{模範解答}
は折りたたみを展開して確認できる（オンライン版）。まず自力で解いてから確認すると効果的である。

\subsection*{問題 2-1　Bν から Bλ への変換とヤコビアン}\label{ex-2-1}
\addcontentsline{toc}{subsection}{問題 2-1　Bν から Bλ
への変換とヤコビアン}

{[}★ 難易度：☆☆ {]}

§2.2 は \(B_\lambda=B_\nu\cdot c/\lambda^2\)
を与えた。これを導く。出発点は「同じエネルギーを異なる変数の刻みで表す」という関係
\(B_\nu\,d\nu=-B_\lambda\,d\lambda\)。

\begin{enumerate}
\def\labelenumi{\arabic{enumi}.}
\tightlist
\item
  \(\nu=c/\lambda\) から \(d\nu=-(c/\lambda^2)d\lambda\)
  を示し、\(B_\lambda=B_\nu\,c/\lambda^2\)
  を導け。負号がなぜ現れ、なぜ最終式で消えるか説明せよ。
\item
  \(B_\nu=\dfrac{2h\nu^3}{c^2}\dfrac{1}{e^{h\nu/kT}-1}\) に代入して
  \(B_\lambda=\dfrac{2hc^2}{\lambda^5}\dfrac{1}{e^{hc/\lambda kT}-1}\)
  を導け。
\item
  なぜ \(B_\nu\) は \(\nu^3\)、\(B_\lambda\) は \(\lambda^{-5}\)
  なのか（冪が3と5でずれる理由）をヤコビアンの観点から一言で述べよ。
\end{enumerate}

\textbf{関連}：\hyperref[sec-reading-spectra-nu-lambda]{§2.2
振動数表示と波長表示}／ピーク位置のずれは\hyperref[ex-2-2]{演習
2-2}、比強度の定義は \href{../part2/04-radiation-field.qmd}{第4章}。

\section{模範解答（問題
2-1）}\label{ux6a21ux7bc4ux89e3ux7b54ux554fux984c-2-1}

\textbf{(1)} \(\nu=c/\lambda\) を微分して
\(d\nu=-(c/\lambda^2)d\lambda\)。\(\nu\) が増えると \(\lambda\)
が減るので負号がつく。物理量（エネルギー）は正なので、刻みの向きを合わせるために
\(B_\nu\,d\nu=-B_\lambda\,d\lambda\) と書く。これに
\(d\nu=-(c/\lambda^2)d\lambda\) を入れると

\[
B_\nu\left(-\frac{c}{\lambda^2}\right)d\lambda=-B_\lambda\,d\lambda\ \Rightarrow\ B_\lambda=B_\nu\frac{c}{\lambda^2}.
\]

両辺の負号が打ち消し、最終式は正になる。

\textbf{(2)} \(\nu=c/\lambda\)
を代入。前因子：\(\dfrac{2h\nu^3}{c^2}\cdot\dfrac{c}{\lambda^2}=\dfrac{2h}{c^2}\dfrac{c^3}{\lambda^3}\dfrac{c}{\lambda^2}=\dfrac{2hc^2}{\lambda^5}\)。指数：\(h\nu/kT=hc/\lambda kT\)。よって

\[
B_\lambda=\frac{2hc^2}{\lambda^5}\frac{1}{e^{hc/\lambda kT}-1}.
\]

\textbf{(3)} \(B_\nu\propto\nu^3\) にヤコビアン
\(c/\lambda^2\propto\lambda^{-2}\)
が掛かる。\(\nu^3\propto\lambda^{-3}\) なので
\(\lambda^{-3}\cdot\lambda^{-2}=\lambda^{-5}\)。冪の差2はヤコビアン由来。

\textbf{答え}：\(B_\lambda=B_\nu c/\lambda^2=\dfrac{2hc^2}{\lambda^5}\dfrac{1}{e^{hc/\lambda kT}-1}\)。\(\lambda^{-5}\)
の「5」はヤコビアン \(\lambda^{-2}\) が \(\nu^3=\lambda^{-3}\)
に掛かるため。

\subsection*{問題 2-2　Wien 変位則の導出 ― 超越方程式と b
の値}\label{ex-2-2}
\addcontentsline{toc}{subsection}{問題 2-2　Wien 変位則の導出 ―
超越方程式と b の値}

{[}★ 難易度：☆☆☆ {]}

§2.3 は \(\lambda_{\max}T=b\simeq2.898\times10^{-3}\) m·K
を結果として与えた。\(B_\lambda\)
を最大化して導く。\(x\equiv hc/\lambda kT\) とおく。

\begin{enumerate}
\def\labelenumi{\arabic{enumi}.}
\tightlist
\item
  \(B_\lambda\propto\lambda^{-5}(e^{x}-1)^{-1}\) を \(\lambda\)
  で微分してゼロとおき、\((5-x)e^{x}=5\)、すなわち \(x=5(1-e^{-x})\)
  を導け。
\item
  数値解 \(x_{\max}\approx4.965\) を用いて \(b=\dfrac{hc}{4.965\,k}\)
  を示し、数値を計算して \(2.90\times10^{-3}\) m·K を確認せよ。
\item
  振動数表示のピークは
  \(h\nu/kT\approx2.821\)（§2.2、第10章の問題で導出）。なぜ
  \(\lambda_{\max}\nu_{\max}\ne c\)
  となるのか、ヤコビアンの観点から説明せよ。
\end{enumerate}

\textbf{関連}：\hyperref[sec-reading-spectra-wien]{§2.3 Wien
の変位則}／振動数表示ピーク 2.821 の導出は\hyperref[ex-10-4]{演習
10-4}、データ照合は\hyperref[ex-1-1]{演習 1-1}。

\section{模範解答（問題
2-2）}\label{ux6a21ux7bc4ux89e3ux7b54ux554fux984c-2-2}

\textbf{(1)}
\(f(\lambda)=\lambda^{-5}(e^{x}-1)^{-1}\)、\(x=hc/\lambda kT\) なので
\(dx/d\lambda=-x/\lambda\)。積の微分：

\[
f'=-5\lambda^{-6}(e^x-1)^{-1}-\lambda^{-5}\frac{e^x}{(e^x-1)^2}\frac{dx}{d\lambda}=0.
\]

\(dx/d\lambda=-x/\lambda\) を代入し、\(\lambda^{-6}(e^x-1)^{-1}\)
で割ると

\[
-5+\frac{x e^x}{e^x-1}=0\ \Rightarrow\ x e^x=5(e^x-1)\ \Rightarrow\ (5-x)e^x=5,\ \text{即ち}\ x=5(1-e^{-x}).
\]

\textbf{(2)} 非自明解は
\(x_{\max}\approx4.965\)。\(x=hc/\lambda_{\max}kT\) より

\[
\lambda_{\max}T=\frac{hc}{x_{\max}k}=\frac{hc}{4.965\,k}.
\]

数値：\(hc=6.626\times10^{-34}\times3.00\times10^8=1.988\times10^{-25}\)
J·m。\(4.965k=4.965\times1.381\times10^{-23}=6.857\times10^{-23}\)。よって
\(b=1.988\times10^{-25}/6.857\times10^{-23}=2.90\times10^{-3}\) m·K。✓

\textbf{(3)} ピークを与える条件は最大化する関数（\(B_\lambda\) か
\(B_\nu\)）で異なる。両者はヤコビアン \(c/\lambda^2\)
だけずれた別関数なので、最大点も別位置になる（\(x=4.965\) と
\(2.821\)）。\(\lambda_{\max}=hc/4.965kT\)、\(\nu_{\max}=2.821kT/h\)
ゆえ
\(\lambda_{\max}\nu_{\max}=(2.821/4.965)c\approx0.568\,c\ne c\)。同じスペクトルでも「どの変数でピークを測るか」で答えが変わる。

\textbf{答え}：\((5-x)e^x=5\Rightarrow x\approx4.965\)、\(b=hc/4.965k\approx2.90\times10^{-3}\)
m·K。\(\lambda_{\max}\nu_{\max}=0.568c\ne c\)（ヤコビアンでピークがずれるため）。

\subsection*{問題 2-3　Fν=πBν ― 半球積分はなぜ 4π でなく π
か}\label{ex-2-3}
\addcontentsline{toc}{subsection}{問題 2-3　Fν=πBν ― 半球積分はなぜ 4π
でなく π か}

{[}★ 難易度：☆☆ {]}

§2.4 は黒体表面の射出フラックスが
\(F_\nu=\displaystyle\int_{\text{半球}}B_\nu\cos\theta\,d\Omega=\pi B_\nu\)
になると述べ、立体角積分の値だけ示した。これを自分で実行する。

\begin{enumerate}
\def\labelenumi{\arabic{enumi}.}
\tightlist
\item
  立体角要素 \(d\Omega=\sin\theta\,d\theta\,d\varphi\)
  を使い、外向き半球（\(0\le\theta\le\pi/2\)）で
  \(\displaystyle\int\cos\theta\,d\Omega\) を計算して \(\pi\) を得よ。
\item
  「全立体角 \(4\pi\)」でも「半球 \(2\pi\)」でもなく \(\pi\)
  になる理由を、（i）半球に限ること と（ii）\(\cos\theta\)
  射影、の二つに分けて説明せよ。
\item
  もし \(\cos\theta\) を落として \(\int_{\text{半球}}B_\nu d\Omega\)
  としたら値はいくつか。なぜそれは物理的に正しくないか。
\end{enumerate}

\textbf{関連}：\hyperref[sec-reading-spectra-stefan-boltzmann]{§2.4
Stefan--Boltzmann 則}／フラックスと比強度の正式な定義（\(\cos\theta\)
の由来）は \href{../part2/04-radiation-field.qmd}{第4章 放射場}。

\section{模範解答（問題
2-3）}\label{ux6a21ux7bc4ux89e3ux7b54ux554fux984c-2-3}

\textbf{(1)} 黒体は等方なので \(B_\nu\) は方向に依らず積分の外へ：

\[
\int_{\text{半球}}\cos\theta\,d\Omega=\int_0^{2\pi}d\varphi\int_0^{\pi/2}\cos\theta\sin\theta\,d\theta=2\pi\left[\frac{\sin^2\theta}{2}\right]_0^{\pi/2}=2\pi\cdot\frac12=\pi.
\]

よって \(F_\nu=\pi B_\nu\)。

\textbf{(2)} （i）表面から外へ出る放射だけを数えるので積分は上半球
\(2\pi\) sr
に限る（下向きは表面内部へ向かい射出に寄与しない）。（ii）フラックスは面を
\textbf{垂直に} 通る成分なので、法線から \(\theta\)
傾いた方向の寄与には射影因子 \(\cos\theta\) が掛かる。\(\cos\theta\)
の角度平均が \(1/2\) なので、\(2\pi\times\tfrac12=\pi\) となる。

\textbf{(3)} \(\int_{\text{半球}}d\Omega=2\pi\) なので値は
\(2\pi B_\nu\)。だがこれは「すべての方向の放射が面を垂直に通る」と誤って数えたことになり、斜め方向の放射が面を通る実効断面が
\(\cos\theta\)
で減る効果を無視している。フラックス（面を通過する正味エネルギー流）の定義に反するので物理的に正しくない。

\textbf{答え}：\(\int_{\text{半球}}\cos\theta\,d\Omega=\pi\)。半球に限る効果（\(2\pi\)）×\(\cos\theta\)
射影の平均（\(1/2\)）＝\(\pi\)。\(\cos\theta\) を落とすと \(2\pi B_\nu\)
となり射影を無視した誤り。

\subsection*{問題 2-4　Stefan--Boltzmann 定数 σ
を基本定数で表す}\label{ex-2-4}
\addcontentsline{toc}{subsection}{問題 2-4　Stefan--Boltzmann 定数 σ
を基本定数で表す}

{[}★ 難易度：☆☆☆ {]}

§2.4 は \(\sigma\)
の数値だけ与え、導出を第6・8章に委ねた。\(F=\pi\displaystyle\int_0^\infty B_\nu\,d\nu\)
から直接 \(\sigma\) を基本定数で表す。

\begin{enumerate}
\def\labelenumi{\arabic{enumi}.}
\tightlist
\item
  \(F=\pi\displaystyle\int_0^\infty\frac{2h\nu^3}{c^2}\frac{d\nu}{e^{h\nu/kT}-1}\)
  で \(x=h\nu/kT\)
  と変数変換し、\(F=\dfrac{2\pi k^4}{c^2h^3}T^4\displaystyle\int_0^\infty\frac{x^3}{e^x-1}dx\)
  を導け。
\item
  \(\displaystyle\int_0^\infty\frac{x^3}{e^x-1}dx=\frac{\pi^4}{15}\)
  を用いて \(\sigma=\dfrac{2\pi^5 k^4}{15c^2h^3}\) を示し、SI
  で数値（\(\approx5.67\times10^{-8}\) W
  m\(^{-2}\)K\(^{-4}\)）を確認せよ。
\item
  第IV部の問題で導いた放射定数 \(a=\dfrac{8\pi^5k^4}{15c^3h^3}\)
  と、\(\sigma=ac/4\) の関係を確かめよ。
\end{enumerate}

\textbf{関連}：\hyperref[sec-reading-spectra-stefan-boltzmann]{§2.4
Stefan--Boltzmann 則}／\(a=4\sigma/c\) と \(a\)
の導出は\hyperref[ex-10-3]{演習 10-3}、本文
\hyperref[sec-planck-quantum-derivation]{§10.3}。積分公式は
\href{../appendix/A-math-notes.qmd}{付録A 数学ノート}。

\section{模範解答（問題
2-4）}\label{ux6a21ux7bc4ux89e3ux7b54ux554fux984c-2-4}

\textbf{(1)} \(x=h\nu/kT\)、\(\nu=kTx/h\)、\(d\nu=(kT/h)dx\)。前因子
\(\dfrac{2h}{c^2}\nu^3=\dfrac{2h}{c^2}\left(\dfrac{kT}{h}\right)^3x^3\)。\(d\nu\)
の \((kT/h)\) と合わせて

\[
F=\pi\cdot\frac{2h}{c^2}\left(\frac{kT}{h}\right)^4\int_0^\infty\frac{x^3}{e^x-1}dx=\frac{2\pi k^4}{c^2h^3}T^4\int_0^\infty\frac{x^3}{e^x-1}dx.
\]

\textbf{(2)} 積分 \(=\pi^4/15\) を代入：

\[
F=\frac{2\pi k^4}{c^2h^3}\frac{\pi^4}{15}T^4=\underbrace{\frac{2\pi^5k^4}{15c^2h^3}}_{\sigma}T^4.
\]

数値（SI）：分子
\(2\pi^5k^4=2\times306.0\times(1.381\times10^{-23})^4=612.0\times3.64\times10^{-92}=2.23\times10^{-89}\)。分母
\(15c^2h^3=15\times(3.00\times10^8)^2\times(6.626\times10^{-34})^3=15\times9.00\times10^{16}\times2.91\times10^{-100}=3.93\times10^{-82}\)。よって
\(\sigma\approx5.67\times10^{-8}\) W m\(^{-2}\)K\(^{-4}\)。✓

\textbf{(3)}
\(\dfrac{ac}{4}=\dfrac{1}{4}\cdot\dfrac{8\pi^5k^4}{15c^3h^3}\cdot c=\dfrac{2\pi^5k^4}{15c^2h^3}=\sigma\)。確かに
\(\sigma=ac/4\)。エネルギー密度 \(u=aT^4\) と射出フラックス
\(F=\sigma T^4\) が \(\sigma=ac/4\) で結ばれる（等方場
\(u=4\pi B/c\)、\(F=\pi B\) から \(F=(c/4)u\)）。

\textbf{答え}：\(\sigma=\dfrac{2\pi^5k^4}{15c^2h^3}\approx5.67\times10^{-8}\)
W m\(^{-2}\)K\(^{-4}\)。放射定数とは \(\sigma=ac/4\)。

\subsection*{問題
2-5　観測データから太陽の温度・光度・半径を引き出す}\label{ex-2-5}
\addcontentsline{toc}{subsection}{問題
2-5　観測データから太陽の温度・光度・半径を引き出す}

{[}★ 難易度：☆☆☆ {]}

§2.8 のワークフローを数値で実行する。観測量：ピーク波長
\(\lambda_{\max}\simeq500\) nm、地球での太陽定数
\(f\simeq1.36\times10^3\) W m\(^{-2}\)、距離 \(d=1\) AU
\(=1.496\times10^{11}\) m。

\begin{enumerate}
\def\labelenumi{\arabic{enumi}.}
\tightlist
\item
  Wien の式（\(b=2.898\times10^{-3}\) m·K）から有効温度 \(T\) を求めよ。
\item
  \(L=4\pi d^2 f\) から太陽光度 \(L_\odot\) を求めよ。
\item
  Stefan--Boltzmann \(L=4\pi R^2\sigma T^4\) を逆算して半径
  \(R_\odot=\sqrt{L/(4\pi\sigma T^4)}\)
  を求め、既知値（\(R_\odot\simeq6.96\times10^8\) m）と比べよ。
\end{enumerate}

\textbf{関連}：\hyperref[sec-reading-spectra-estimation]{§2.8
観測データから推定}／各道具は
\hyperref[sec-reading-spectra-wien]{§2.3}・\hyperref[sec-reading-spectra-stefan-boltzmann]{§2.4}、\(\sigma\)
の導出は\hyperref[ex-2-4]{演習 2-4}。

\section{模範解答（問題
2-5）}\label{ux6a21ux7bc4ux89e3ux7b54ux554fux984c-2-5}

\textbf{(1)}
\(T=b/\lambda_{\max}=2.898\times10^{-3}/5.00\times10^{-7}=5.80\times10^3\)
K。（有効温度 \(\simeq5800\) K）

\textbf{(2)}
\(L=4\pi d^2 f=4\pi(1.496\times10^{11})^2\times1.36\times10^3\)。\(d^2=2.238\times10^{22}\)、\(4\pi d^2=2.812\times10^{23}\)、\(\times f=3.82\times10^{26}\)
W。すなわち \(L_\odot\simeq3.8\times10^{26}\) W（既知
\(3.83\times10^{26}\) W と一致）。

\textbf{(3)}
\(\sigma T^4=5.67\times10^{-8}\times(5800)^4=5.67\times10^{-8}\times1.13\times10^{15}=6.41\times10^7\)
W m\(^{-2}\)。

\[
R=\sqrt{\frac{L}{4\pi\sigma T^4}}=\sqrt{\frac{3.82\times10^{26}}{4\pi\times6.41\times10^7}}=\sqrt{\frac{3.82\times10^{26}}{8.05\times10^8}}=\sqrt{4.74\times10^{17}}\approx6.9\times10^8\ \text{m}.
\]

既知値 \(6.96\times10^8\) m
とよく一致。観測量（ピーク波長・フラックス・距離）だけから太陽の物理量がプランク分布を介して再現できる。

\textbf{答え}：\(T\simeq5800\) K、\(L_\odot\simeq3.8\times10^{26}\)
W、\(R_\odot\simeq6.9\times10^8\) m。いずれも実測とよく一致。

\section*{さらに学ぶための参考文献}\label{sec-reading-spectra-further}
\addcontentsline{toc}{section}{さらに学ぶための参考文献}

\markright{さらに学ぶための参考文献}

\begin{itemize}
\tightlist
\item
  Rybicki \& Lightman, \emph{Radiative Processes in Astrophysics}
  (Wiley, 1979) --- 1.5 節「Blackbody Radiation」
\item
  嶺重 慎 編『宇宙流体力学の基礎』（日本評論社, 2014）---
  放射場の記述に関する章
\item
  Karttunen et al., \emph{Fundamental Astronomy} (Springer, 6th ed.,
  2017) --- 黒体放射と光度・色温度の基礎
\end{itemize}

\chapter{第3章　黒体放射はなぜ普遍的なのか}\label{ux7b2c3ux7ae0-ux9ed2ux4f53ux653eux5c04ux306fux306aux305cux666eux904dux7684ux306aux306eux304b}

\textbf{この章の主題}：第1章で挙げた多彩な天体 ― 太陽から
CMB、ダスト、降着円盤まで ― が同じ関数族 \(B_\nu(T)\)
のスペクトルを示すのはなぜか、を考える。答えの核心は「熱平衡」と「空洞放射」の概念にある。この章で本書の三原則の「①
逆引き」が一段階深くなり、観測の \textbf{普遍性} が \textbf{熱平衡}
という条件に結びつく。

\section*{この章の中心地図}\label{sec-universality-map}
\addcontentsline{toc}{section}{この章の中心地図}

\markright{この章の中心地図}

\textbf{プランク関数（連続スペクトルの地図）}

\begin{equation}\protect\phantomsection\label{eq-planck-function}{
B_\nu(T) = \frac{2 h \nu^3}{c^2} \cdot \frac{1}{e^{h\nu/kT} - 1}
}\end{equation}

{\def\LTcaptype{none} % do not increment counter
\begin{longtable}[]{@{}lll@{}}
\toprule\noalign{}
因子 & 物理的意味 & 結びつく部 \\
\midrule\noalign{}
\endhead
\bottomrule\noalign{}
\endlastfoot
\(B_\nu\) & 観測量（比強度） & 第II部 \\
\(1/(e^{h\nu/kT}-1)\) & 熱平衡・光子統計 & 第III部 \\
\(h\nu\) & エネルギー量子 & 第IV部 \\
\(\nu^3/c^2\) & 電磁場のモード密度 & 第V部 \\
黒体からのずれ & 現実の天体 & 第VI部 \\
\end{longtable}
}

\begin{tcolorbox}[enhanced jigsaw, arc=.35mm, bottomrule=.15mm, breakable, colback=white, colframe=quarto-callout-tip-color-frame, left=2mm, leftrule=.75mm, opacityback=0, rightrule=.15mm, toprule=.15mm]

\textbf{地図の使い方}：各部の冒頭でこの式が再掲され、その部で扱う因子が強調表示される。読者は自分が今どこを学んでいるかをこの式の上で確認できる。

\end{tcolorbox}

\begin{tcolorbox}[enhanced jigsaw, arc=.35mm, bottomrule=.15mm, bottomtitle=1mm, breakable, colback=white, colbacktitle=quarto-callout-tip-color!10!white, colframe=quarto-callout-tip-color-frame, coltitle=black, left=2mm, leftrule=.75mm, opacityback=0, opacitybacktitle=0.6, rightrule=.15mm, title=\textcolor{quarto-callout-tip-color}{\faLightbulb}\hspace{0.5em}{ヒント}, titlerule=0mm, toprule=.15mm, toptitle=1mm]

本章は中心地図のうち、特に「\textbf{この式が熱平衡で成り立つこと}」 ―
つまりプランク関数が普遍的なスペクトルになる理由 ―
を扱う。これは第III部全体への入口でもある。

\end{tcolorbox}

\section*{この章で答える問い}\label{sec-universality-questions}
\addcontentsline{toc}{section}{この章で答える問い}

\markright{この章で答える問い}

\begin{itemize}
\tightlist
\item
  なぜ全く異なる物質（恒星表面・星間塵・初期宇宙）が、同じ関数形のスペクトルを示すのか
\item
  「黒体」と「完全吸収体」はなぜ同じものとして扱えるのか
\item
  キルヒホッフの法則は「ある物質が良い吸収体なら良い放出体でもある」と言うが、その必然性はどこから来るのか
\item
  現実の天体が「黒体に近づく」のは、どんな条件が満たされるからなのか
\end{itemize}

\section*{到達目標}\label{sec-universality-goals}
\addcontentsline{toc}{section}{到達目標}

\markright{到達目標}

この章を読み終えると、読者は次のことができるようになる：

\begin{itemize}
\tightlist
\item
  黒体放射の普遍性が、物質の詳細ではなく熱平衡という条件から来ていることを説明できる
\item
  空洞放射の考え方を用いて、なぜ黒体スペクトルが物質によらないかを論じられる
\item
  実際の天体が黒体に近づくための条件（光学的厚さ・熱化時間など）を列挙できる
\end{itemize}

\begin{center}\rule{0.5\linewidth}{0.5pt}\end{center}

\section{3.1 材質によらないスペクトル ―
観測の不思議}\label{sec-universality-material-independence}

{[}本文目安：B2-B3{]}

第1章で並べた天体は、化学組成・密度・サイズが極端に異なる：

{\def\LTcaptype{none} % do not increment counter
\begin{longtable}[]{@{}
  >{\raggedright\arraybackslash}p{(\linewidth - 4\tabcolsep) * \real{0.3333}}
  >{\raggedright\arraybackslash}p{(\linewidth - 4\tabcolsep) * \real{0.3333}}
  >{\raggedright\arraybackslash}p{(\linewidth - 4\tabcolsep) * \real{0.3333}}@{}}
\toprule\noalign{}
\begin{minipage}[b]{\linewidth}\raggedright
天体
\end{minipage} & \begin{minipage}[b]{\linewidth}\raggedright
主な物質
\end{minipage} & \begin{minipage}[b]{\linewidth}\raggedright
密度のオーダー
\end{minipage} \\
\midrule\noalign{}
\endhead
\bottomrule\noalign{}
\endlastfoot
太陽 & 水素プラズマ & \(10^{-7}\) g/cm³（光球） \\
星間ダスト & 炭素・ケイ酸塩固体 & 個別粒子は \(\sim 1\) g/cm³ \\
CMB & 光子気体（電子・陽子はほぼ全て中性化済） & 光子数密度 \(\sim 400\)
cm\(^{-3}\) \\
降着円盤 & 高温電離プラズマ & \(10^{-7}\)〜\(10^{-3}\) g/cm³ \\
\end{longtable}
}

それでも、これらが示すスペクトルは（細かな線吸収・放出効率の振動数依存性を別にすれば）\textbf{同じ関数族
\(B_\nu(T)\)} に従う。

これは、初学者にとっては驚くべき事実である。\textbf{化学・物理状態が全く違うのに、なぜ放射スペクトルだけが普遍的になるのか}？
その答えこそが本章のテーマである。先に結論を言えば：

\begin{quote}
\textbf{黒体放射の普遍性は、放射場が「熱平衡」にあるという条件だけから決まる。物質の詳細は、平衡を保証するために必要なだけで、平衡そのもののスペクトル形には影響しない。}
\end{quote}

これは深い命題で、本書全体を貫く中心的なメッセージのひとつである。以下、その理屈を空洞放射と
Kirchhoff の法則を通じて辿る。

\section{3.2
熱平衡にある放射場}\label{sec-universality-thermal-equilibrium}

{[}本文目安：B2-B3{]}

「熱平衡」を観測の言葉で定義すると、次のようになる：

\begin{quote}
ある系（物質＋放射場）が、十分に長い時間にわたって相互作用した結果、エネルギーの流れが正味でゼロになった状態。各部分が同じ温度
\(T\) で記述できる。
\end{quote}

ここで本書では二段階の概念を区別する：

\begin{itemize}
\tightlist
\item
  \textbf{完全熱平衡}（complete thermodynamic
  equilibrium）：物質も放射場も\textbf{両方}が同じ温度 \(T\)
  で平衡。次の三つが同時に成り立つ。

  \begin{enumerate}
  \def\labelenumi{\arabic{enumi}.}
  \tightlist
  \item
    \textbf{詳細釣り合い}：放射の吸収と放出が各振動数で正味ゼロ
  \item
    \textbf{温度の一意性}：物質と放射場が同じ単一の温度 \(T\)
    で記述できる
  \item
    \textbf{放射場の等方性}：\(I_\nu\) が方向に依存しない
  \end{enumerate}
\item
  \textbf{LTE（局所熱力学平衡, local thermodynamic equilibrium;
  LTE）}：物質の状態（速度分布・電離・励起）が局所温度 \(T\)
  で決まる近似。放射場 \(J_\nu\)
  は\textbf{必ずしも等方や黒体である必要はない}。たとえば恒星大気の表層近傍では、物質は
  LTE のままで \(J_\nu\) は異方化していく。
\end{itemize}

本章では、第3節以降の空洞放射の議論において\textbf{完全熱平衡}を想定する。LTE
と完全熱平衡の区別は、第5章 §5.5 で本格的に扱う。

完全熱平衡の場での放射スペクトルは、\textbf{物質の種類によらず一意に決まる}
―
これが第3章の核心である。なぜそうなるかは、次節の空洞放射の議論で明らかになる。

\begin{tcolorbox}[enhanced jigsaw, arc=.35mm, bottomrule=.15mm, bottomtitle=1mm, breakable, colback=white, colbacktitle=quarto-callout-note-color!10!white, colframe=quarto-callout-note-color-frame, coltitle=black, left=2mm, leftrule=.75mm, opacityback=0, opacitybacktitle=0.6, rightrule=.15mm, title=\textcolor{quarto-callout-note-color}{\faInfo}\hspace{0.5em}{ノート}, titlerule=0mm, toprule=.15mm, toptitle=1mm]

\textbf{熱平衡の定義は本章では仮定}：「熱平衡」を出発点として受け入れ、その帰結を見るのが本章の立場。熱平衡自体の起源（最大エントロピー原理）は第6章で逆引きする。

\end{tcolorbox}

\section{3.3 空洞放射の考え方 ― 思考実験}\label{sec-universality-cavity}

{[}本文目安：B2-B3{]}

黒体放射の普遍性を理解する最良の道具が、\textbf{空洞放射}（\textbf{cavity
radiation}）の思考実験である。

\textbf{設定}：

\begin{enumerate}
\def\labelenumi{\arabic{enumi}.}
\tightlist
\item
  任意の物質でできた壁を持つ、内部が真空（または希薄物質）の空洞を用意する
\item
  壁を温度 \(T\) に保つ
\item
  十分に長い時間放置し、内部の電磁波と壁が熱平衡に達するまで待つ
\end{enumerate}

\textbf{観察}：内部の放射場は、その温度 \(T\)
だけで決まる特定のスペクトル \(u_\nu(T)\)
を持つ。\textbf{このスペクトルは、壁の材質によらない}。

\textbf{証明（背理法）}：もし壁の材質によってスペクトルが違うと仮定する。すると：

\begin{itemize}
\tightlist
\item
  異なる材質の二つの空洞 A・B を同じ温度 \(T\)
  で熱平衡にしたあと、両者を細い導管でつなぐ
\item
  もしスペクトルが違えば、ある振動数帯でエネルギーが A → B（あるいは B →
  A）に流れる
\item
  これは「熱平衡で温度が同じなのに正味のエネルギー輸送がある」ことを意味し、熱力学第二法則に反する
\end{itemize}

したがって、\textbf{空洞放射のスペクトルは壁の材質によらず一意である}。これが「黒体放射の普遍性」の正体である。

\begin{tcolorbox}[enhanced jigsaw, arc=.35mm, bottomrule=.15mm, bottomtitle=1mm, breakable, colback=white, colbacktitle=quarto-callout-important-color!10!white, colframe=quarto-callout-important-color-frame, coltitle=black, left=2mm, leftrule=.75mm, opacityback=0, opacitybacktitle=0.6, rightrule=.15mm, title=\textcolor{quarto-callout-important-color}{\faExclamation}\hspace{0.5em}{重要}, titlerule=0mm, toprule=.15mm, toptitle=1mm]

\textbf{重要}：壁に小さな穴を開けると、その穴は完璧な吸収体である（入射光は穴を通って空洞内に入り、ほぼ確実に吸収される）。同時に穴からは空洞内部の熱平衡放射
― つまり黒体スペクトル ―
が出てくる。\textbf{完全吸収体は完全放射体である}、これが Kirchhoff
の法則の本質である。

\end{tcolorbox}

\begin{tcolorbox}[enhanced jigsaw, arc=.35mm, bottomrule=.15mm, breakable, colback=white, colframe=quarto-callout-tip-color-frame, left=2mm, leftrule=.75mm, opacityback=0, rightrule=.15mm, toprule=.15mm]
\begin{minipage}[t]{5.5mm}
\textcolor{quarto-callout-tip-color}{\faLightbulb}
\end{minipage}%
\begin{minipage}[t]{\textwidth - 5.5mm}

\textbf{問い}：壁が完全鏡（吸収率 0）なら、空洞放射はどうなるか？

\textbf{短答}：壁が光を吸収しないので、放射場と物質が相互作用しない。結果として放射場は熱平衡に達せず、内部の放射スペクトルは「最初に入っていた光」のまま決まらない。\textbf{空洞放射の議論には、少なくとも一部の吸収が必須}である。

\textbf{もう一歩}：この状況は実在する。たとえばレーザー共振器（高反射率ミラー）はこれに近く、定常的な黒体放射は出さない。逆に小穴を開けた空洞は実効的な完全吸収体になり、必ず黒体スペクトルを出す（§3.4）。「黒体は熱平衡に達した放射場」という命題は、平衡に達するための相互作用の存在を前提にしている。

\end{minipage}%
\end{tcolorbox}

\section{3.4
完全吸収体と完全放射体}\label{sec-universality-perfect-emitter}

{[}本文目安：B2-B3{]}

黒体（black
body）の名前の由来は、\textbf{入射光をすべて吸収する物体}にある。色が黒く見えるのは反射光がないからだ。前節の空洞放射の議論から、次のことがわかる：

\begin{quote}
\textbf{完全吸収体は、温度 \(T\)
で平衡を保てば、必ず完全な黒体スペクトルを放出する。}
\end{quote}

この対応関係は、しばしば次のように整理される：

\begin{itemize}
\tightlist
\item
  \textbf{吸収率} \(a_\nu\)（absorptivity,
  dimensionless）：入射光のうち、その物体が吸収する割合（\(0 \le a_\nu \le 1\)）
\item
  \textbf{放出率} \(\epsilon_\nu\)（emissivity,
  dimensionless）：その物体が同じ温度 \(T\)
  で出す放射の、完全黒体に対する比
\end{itemize}

両者の関係は \textbf{Kirchhoff
の法則}（次節で本格的に）で結ばれ、熱平衡では

\begin{equation}\protect\phantomsection\label{eq-kirchhoff}{
a_\nu = \epsilon_\nu
}\end{equation}

が成り立つ。つまり吸収が強い振動数では放出も強く、吸収が弱い振動数では放出も弱い。

具体例：

\begin{itemize}
\tightlist
\item
  黒体：\(a_\nu = \epsilon_\nu = 1\) ですべての \(\nu\)
\item
  灰色物体：\(a_\nu = \epsilon_\nu = \mathrm{const} < 1\) ですべての
  \(\nu\)
\item
  ダスト：\(a_\nu = \epsilon_\nu = a(\nu) < 1\) で振動数依存（→
  修正黒体、第15章）
\item
  完全鏡：\(a_\nu = \epsilon_\nu = 0\) で何も吸収せず何も放出しない
\end{itemize}

\begin{tcolorbox}[enhanced jigsaw, arc=.35mm, bottomrule=.15mm, bottomtitle=1mm, breakable, colback=white, colbacktitle=quarto-callout-note-color!10!white, colframe=quarto-callout-note-color-frame, coltitle=black, left=2mm, leftrule=.75mm, opacityback=0, opacitybacktitle=0.6, rightrule=.15mm, title=\textcolor{quarto-callout-note-color}{\faInfo}\hspace{0.5em}{ノート}, titlerule=0mm, toprule=.15mm, toptitle=1mm]

\textbf{記号に注意}：本章の \textbf{吸収率} \(a_\nu\)
は無次元の比であり、放射輸送（第5章）に出てくる \textbf{吸収係数}
\(\alpha_\nu\)（cm\(^{-1}\)）とは別の量である。両者を混同しないよう、本書では：

\begin{itemize}
\tightlist
\item
  \(a_\nu\)（小文字 a）：吸収率、無次元
\item
  \(\alpha_\nu\)：吸収係数、単位 cm\(^{-1}\)
\end{itemize}

として区別する。

\end{tcolorbox}

\section{3.5 Kirchhoff の法則の意味 ―
観測スペクトルへの効き方}\label{sec-universality-kirchhoff}

{[}本文目安：B2-B3{]}

Kirchhoff の法則は、放射輸送（第5章）の
\textbf{ソース関数}（\textbf{source function}）の表式に直結する。

放射輸送方程式は

\begin{equation}\protect\phantomsection\label{eq-rt-kirchhoff-context}{
\frac{dI_\nu}{ds} = -\alpha_\nu I_\nu + j_\nu
}\end{equation}

の形で、ソース関数を \(S_\nu \equiv j_\nu / \alpha_\nu\) と定義する。

\textbf{局所熱平衡}（\textbf{LTE}: \textbf{L}ocal \textbf{T}hermodynamic
\textbf{E}quilibrium）では、Kirchhoff の法則 式~\ref{eq-kirchhoff} から

\begin{equation}\protect\phantomsection\label{eq-source-equals-planck}{
S_\nu = B_\nu(T)
}\end{equation}

が成り立つ。つまり「\textbf{熱平衡ではソース関数がプランク関数になる}」。これは第5章の放射輸送の議論の中心的な結果で、観測スペクトルが黒体に近づく条件と直結する。

\begin{tcolorbox}[enhanced jigsaw, arc=.35mm, bottomrule=.15mm, bottomtitle=1mm, breakable, colback=white, colbacktitle=quarto-callout-tip-color!10!white, colframe=quarto-callout-tip-color-frame, coltitle=black, left=2mm, leftrule=.75mm, opacityback=0, opacitybacktitle=0.6, rightrule=.15mm, title=\textcolor{quarto-callout-tip-color}{\faLightbulb}\hspace{0.5em}{ヒント}, titlerule=0mm, toprule=.15mm, toptitle=1mm]

\textbf{観測との対応}：観測される連続スペクトルが黒体に近いかどうかを判定するとき、私たちは暗黙のうちに
\(S_\nu = B_\nu(T)\)
が成り立っているか（=その天体が局所熱平衡か）を判定している。星の光球は概ね
LTE、星間ダストもダスト粒子内部では LTE。一方、希薄な星間ガスや AGN
の輝線領域は非LTE で、\(S_\nu \ne B_\nu\) になる ―
これが線スペクトルの主役となる。詳しくは第21章・第7部で扱う。

\end{tcolorbox}

\section{3.6
現実の天体はなぜ黒体に近づくのか}\label{sec-universality-real-objects}

{[}本文目安：B2-B3{]}

黒体放射の「普遍性」は理想化された熱平衡の話だった。では、現実の天体はどんな条件で黒体に近づくのか？
三つの条件を整理する：

\subsection{条件
1：光学的に厚い}\label{ux6761ux4ef6-1ux5149ux5b66ux7684ux306bux539aux3044}

光が物体内部を通り抜けるとき、十分に何度も吸収・再放出が繰り返されれば、放射場と物質が熱化する。これを「\textbf{光学的に厚い}」（\textbf{optically
thick}, \(\tau \gg 1\)）と表現する。

\begin{itemize}
\tightlist
\item
  太陽の光球：可視光で \(\tau \gg 1\) → ほぼ黒体
\item
  大気圏外の希薄ガス：\(\tau \ll 1\) →
  黒体にならない（線スペクトルになる）
\end{itemize}

\subsection{条件 2：熱化時間 \textless{}
流体力学的時間}\label{ux6761ux4ef6-2ux71b1ux5316ux6642ux9593-ux6d41ux4f53ux529bux5b66ux7684ux6642ux9593}

光と物質が熱平衡に達するのに必要な時間（熱化時間
\(t_\mathrm{therm}\)）が、温度・密度が変わるスケール（流体力学的時間
\(t_\mathrm{dyn}\)）より短ければ、局所的に熱平衡が成り立つ。

\begin{itemize}
\tightlist
\item
  恒星内部：\(t_\mathrm{therm} \ll t_\mathrm{dyn}\) → ほぼ LTE
\item
  早期宇宙（再結合期以前）：同様 → 完璧な熱平衡 → CMB の完全な黒体性
\item
  衝撃波直後：\(t_\mathrm{therm} \sim t_\mathrm{dyn}\) → 部分的に非平衡
\end{itemize}

\subsection{条件
3：単一温度で記述できる}\label{ux6761ux4ef6-3ux5358ux4e00ux6e29ux5ea6ux3067ux8a18ux8ff0ux3067ux304dux308b}

物体全体が単一の温度 \(T\) で記述できれば、スペクトルは単一の
\(B_\nu(T)\)。複数の温度成分が混ざれば、合成スペクトル（多温度黒体）になる。

\begin{itemize}
\tightlist
\item
  CMB：単一温度 2.725 K → 単一黒体
\item
  降着円盤：半径ごとに温度が異なる → 多温度合成（第18章）
\end{itemize}

\textbf{観測上の判定}：これら三条件のうち、どれが破れているかで「黒体からのずれ」の正体が変わる。これは第II部の放射輸送・第VI部の応用で具体的に扱う。

\section{3.7
黒体放射が「特別な基準」になる理由}\label{sec-universality-benchmark}

{[}本文目安：B2-B3{]}

ここまでの議論を振り返ろう：

\begin{enumerate}
\def\labelenumi{\arabic{enumi}.}
\tightlist
\item
  観測される連続スペクトルは、しばしば \(B_\nu(T)\) の形に近い（第1章）
\item
  その形は Wien の式・Stefan--Boltzmann
  則などで「読む」道具を持つ（第2章）
\item
  \textbf{形が普遍的になるのは、熱平衡という条件のおかげ}（本章）
\end{enumerate}

これらを合わせると、\textbf{黒体放射は宇宙物理における「基準スペクトル」}として、二重の意味で特別である：

\textbf{第一に、観測の「ものさし」として}。
黒体スペクトルは、観測されるスペクトルがどれだけ「熱化された」状態にあるかを測る基準を提供する。観測スペクトルを黒体に当てはめてみて、どれだけずれているかを定量化することで、放射の起源を判定できる。

\textbf{第二に、基礎物理の「結節点」として}。 プランク関数 \(B_\nu(T)\)
の形は、統計力学（第III部）、量子論（第IV部）、電磁気学（第V部）が交わる一点で決まる。観測される黒体スペクトル
―「あの曲線」―
を完全に説明するには、これら三つの基礎物理すべてが必要となる。逆に言えば、\(B_\nu(T)\)
の各因子をひとつずつ「逆引き」していけば、現代物理の中核を順に再構成できる。これが本書のロードマップである。

\begin{tcolorbox}[enhanced jigsaw, arc=.35mm, bottomrule=.15mm, bottomtitle=1mm, breakable, colback=white, colbacktitle=quarto-callout-tip-color!10!white, colframe=quarto-callout-tip-color-frame, coltitle=black, left=2mm, leftrule=.75mm, opacityback=0, opacitybacktitle=0.6, rightrule=.15mm, title=\textcolor{quarto-callout-tip-color}{\faLightbulb}\hspace{0.5em}{ヒント}, titlerule=0mm, toprule=.15mm, toptitle=1mm]

\textbf{本章の到達点}：第I部の三つの章で、黒体放射という現象を「観測 →
読み方 →
普遍性の起源」と段階的に深掘りしてきた。普遍性の起源が「熱平衡」にあると示せたところで、第I部は終わる。第II部では、放射場そのものを記述する言語を整え、第III部で熱平衡の中身を逆引きする
― これが本書のリズムである。

\end{tcolorbox}

\begin{center}\rule{0.5\linewidth}{0.5pt}\end{center}

\section*{この章で何がわかったか}\label{sec-universality-summary}
\addcontentsline{toc}{section}{この章で何がわかったか}

\markright{この章で何がわかったか}

\textbf{中心地図に戻る}

本章で、中心地図 \(B_\nu(T)\) の \textbf{「式が成り立つことそのもの」} ―
つまり熱平衡で物質の詳細によらず一意のスペクトル形になる、という事実 ―
の理由が明らかになった：

\begin{itemize}
\tightlist
\item
  空洞放射の思考実験（熱力学第二法則）→ スペクトルは壁の材質によらず一意
\item
  Kirchhoff の法則 → 吸収率 = 放出率（熱平衡）
\item
  ソース関数 \(S_\nu = B_\nu(T)\)（LTE）→
  観測スペクトルが黒体に近づく条件
\end{itemize}

ただし「なぜプランク関数の \textbf{具体的な関数形}
がこれなのか」は、まだ示されていない。これは第III部（光子統計）・第IV部（量子化）・第V部（モード密度）の逆引きで段階的に明らかになる。

\textbf{次章へ}：第I部はここで終わり、第II部に移る。第4章では、放射場を記述するための基本量（エネルギー密度・フラックス・比強度・立体角）を整理し、観測量と理論量を結ぶ言語を整える。第5章では放射輸送方程式を扱い、本章で触れたソース関数
\(S_\nu\) と LTE の意味をより精密にする。

\section*{演習問題}\label{sec-universality-exercises}
\addcontentsline{toc}{section}{演習問題}

\markright{演習問題}

\textbf{模範解答}
は折りたたみを展開して確認できる（オンライン版）。まず自力で解いてから確認すると効果的である。

\subsection*{問題
3-1　空洞放射の普遍性を第二法則で証明する}\label{ex-3-1}
\addcontentsline{toc}{subsection}{問題
3-1　空洞放射の普遍性を第二法則で証明する}

{[}★ 難易度：☆☆ {]}

§3.3 は背理法の骨子だけ示した。これを各振動数について厳密化する。

\begin{enumerate}
\def\labelenumi{\arabic{enumi}.}
\tightlist
\item
  異なる材質の空洞 A・B を同じ温度 \(T\) に保ち、振動数
  \([\nu,\nu+d\nu]\) だけ通すフィルターを挟んで導管でつなぐ。もし
  \(u_\nu^A\ne u_\nu^B\) なら何が起こるか述べよ。
\item
  それが熱力学第二法則のどの主張に反するか説明せよ（同一温度の二系の間に正味のエネルギー流＝仕事を取り出せる永久機関）。
\item
  フィルターを使う議論が「全エネルギーだけでなく各振動数で \(u_\nu\)
  が普遍」という強い結論を与える理由を述べよ。
\end{enumerate}

\textbf{関連}：\hyperref[sec-universality-cavity]{§3.3
空洞放射}／熱平衡そのものの起源（最大エントロピー）は
\href{../part3/06-thermal-equilibrium.qmd}{第6章}。

\section{模範解答（問題
3-1）}\label{ux6a21ux7bc4ux89e3ux7b54ux554fux984c-3-1}

\textbf{(1)} フィルターは \([\nu,\nu+d\nu]\) の放射だけを通す。もし
\(u_\nu^A>u_\nu^B\) なら、この帯域で A から B
へ正味のエネルギーが流れる（逆なら逆向き）。両空洞は同じ温度 \(T\)
なのに、放射エネルギーが一方向に流れ続ける。

\textbf{(2)}
同一温度の二つの系の間に外部からの仕事なしに正味のエネルギー（熱）が流れることになる。この流れを使えば永久に仕事を取り出せ、「外部に変化を残さず、低温部へ熱を移す以外に熱を仕事に変換し続ける機関は作れない」という熱力学第二法則（Clausius/Kelvin
の主張）に反する。よって仮定 \(u_\nu^A\ne u_\nu^B\) は誤り。

\textbf{(3)} フィルターで振動数を1つずつ切り出せるので、「もし
\textbf{ある} 振動数で \(u_\nu\)
が違えば」その帯域だけで第二法則違反が起こる。したがって違反を避けるには
\textbf{すべての} 振動数で \(u_\nu^A=u_\nu^B\)
でなければならない。全エネルギー（積分値）の一致よりも強い、スペクトルの各点での普遍性が示せる。

\textbf{答え}：\(u_\nu^A\ne u_\nu^B\)
だと同一温度の二系間に正味エネルギー流＝第二法則違反。フィルター議論で各振動数ごとに
\(u_\nu(T)\) が材質によらず一意と結論。

\subsection*{問題 3-2　Kirchhoff の法則 aν=εν
を平衡から導く}\label{ex-3-2}
\addcontentsline{toc}{subsection}{問題 3-2　Kirchhoff の法則 aν=εν
を平衡から導く}

{[}★ 難易度：☆☆ {]}

§3.4--3.5 は
\(a_\nu=\epsilon_\nu\)（吸収率＝放出率）を結果として与えた。空洞内の平衡から導く。温度
\(T\) の空洞（内部は黒体放射 \(B_\nu(T)\)）に、吸収率 \(a_\nu\)・放出率
\(\epsilon_\nu\) の小物体を置く。

\begin{enumerate}
\def\labelenumi{\arabic{enumi}.}
\tightlist
\item
  単位時間・単位面積・単位振動数あたり、この物体が吸収するエネルギーと放出するエネルギーをそれぞれ
  \(a_\nu, \epsilon_\nu, B_\nu\) で書け。
\item
  物体が空洞と同じ温度 \(T\)
  を保つ（熱平衡）には吸収＝放出が必要。これから \(a_\nu=\epsilon_\nu\)
  を導け。
\item
  なぜこの等式が「各振動数ごと」に成り立たねばならないのか（詳細釣り合い）を述べよ。
\end{enumerate}

\textbf{関連}：\hyperref[sec-universality-perfect-emitter]{§3.4
完全吸収体と完全放射体}、\hyperref[sec-universality-kirchhoff]{§3.5
Kirchhoff の法則}／詳細釣り合いの一般論は
\href{../part3/06-thermal-equilibrium.qmd}{第6章 §6.6}。

\section{模範解答（問題
3-2）}\label{ux6a21ux7bc4ux89e3ux7b54ux554fux984c-3-2}

\textbf{(1)} 空洞内の入射放射は黒体
\(B_\nu(T)\)。吸収エネルギー（率）は入射のうち吸収される割合：\(\propto a_\nu B_\nu(T)\)。放出エネルギーは、完全黒体の放出
\(B_\nu(T)\) に放出率を掛けた
\(\propto \epsilon_\nu B_\nu(T)\)。（比例係数＝立体角・面積因子は両者共通。）

\textbf{(2)} 物体が温度 \(T\) を保つには、各振動数で吸収＝放出：

\[
a_\nu B_\nu(T)=\epsilon_\nu B_\nu(T)\ \Rightarrow\ a_\nu=\epsilon_\nu.
\]

\(B_\nu(T)\ne0\) なので約せる。吸収が強い振動数では放出も強い。

\textbf{(3)}
もしある振動数で吸収＞放出、別の振動数で吸収＜放出だと、全体では釣り合っていても、物体は前者の振動数で正味エネルギーを受け取り後者で失う。これは空洞放射のスペクトルを各振動数で変えてしまい（あるいは物体内に振動数間のエネルギー輸送を生み）、平衡と矛盾する。第二法則と整合する平衡（詳細釣り合い）は各振動数ごとの釣り合いを要求するので、\(a_\nu=\epsilon_\nu\)
は振動数ごとに成り立つ。

\textbf{答え}：平衡で
\(a_\nu B_\nu=\epsilon_\nu B_\nu\Rightarrow a_\nu=\epsilon_\nu\)。詳細釣り合いゆえ各振動数で成立。

\subsection*{問題 3-3　LTE でソース関数が Sν=Bν になる}\label{ex-3-3}
\addcontentsline{toc}{subsection}{問題 3-3　LTE でソース関数が Sν=Bν
になる}

{[}★ 難易度：☆☆ {]}

§3.5 は放射輸送方程式 \(\dfrac{dI_\nu}{ds}=-\alpha_\nu I_\nu+j_\nu\)
とソース関数 \(S_\nu\equiv j_\nu/\alpha_\nu\) を導入し、LTE で
\(S_\nu=B_\nu(T)\) になると述べた。これを空洞平衡から導く。

\begin{enumerate}
\def\labelenumi{\arabic{enumi}.}
\tightlist
\item
  熱平衡の空洞内では放射場が等方かつ一様で
  \(I_\nu=B_\nu(T)\)、したがって
  \(dI_\nu/ds=0\)。これを輸送方程式に入れて
  \(j_\nu/\alpha_\nu=B_\nu(T)\)、すなわち \(S_\nu=B_\nu(T)\) を示せ。
\item
  \(S_\nu=B_\nu\) が成り立つのは「物質が局所温度 \(T\)
  で熱化している（LTE）」ときである。なぜ放射場 \(J_\nu\)
  が黒体でなくても \(S_\nu=B_\nu\) が言えるのか、Kirchhoff
  の法則（\hyperref[ex-3-2]{演習 3-2}）が \(j_\nu/\alpha_\nu\)
  を物質の温度だけの関数にすることに触れて述べよ。
\item
  逆に希薄ガスや AGN 輝線領域で \(S_\nu\ne B_\nu\)
  になるのはなぜか、一言で。
\end{enumerate}

\textbf{関連}：\hyperref[sec-universality-kirchhoff]{§3.5 Kirchhoff
とソース関数}／放射輸送の本格的扱いは
\href{../part2/05-radiative-transfer.qmd}{第5章}、線形成での再導出は
\href{../part7/21-line-formation.qmd}{第21章 §21.4}。

\section{模範解答（問題
3-3）}\label{ux6a21ux7bc4ux89e3ux7b54ux554fux984c-3-3}

\textbf{(1)} 平衡の空洞では \(I_\nu=B_\nu(T)\)
が一様等方なので位置によらず \(dI_\nu/ds=0\)。輸送方程式に代入：

\[
0=-\alpha_\nu B_\nu(T)+j_\nu\ \Rightarrow\ j_\nu=\alpha_\nu B_\nu(T)\ \Rightarrow\ S_\nu\equiv\frac{j_\nu}{\alpha_\nu}=B_\nu(T).
\]

\textbf{(2)} Kirchhoff の法則は放出と吸収の比 \(j_\nu/\alpha_\nu\) を
\textbf{物質の温度 \(T\) だけ} の関数（\(=B_\nu(T)\)）にする。\(j_\nu\)
も \(\alpha_\nu\) も物質の局所状態（温度 \(T\)
で熱化した占有数）で決まり、その比は放射場 \(I_\nu\) や \(J_\nu\)
に依らない。だから周囲の放射場が黒体でなく異方的でも、物質さえ局所温度
\(T\) で熱化（LTE）していれば \(S_\nu=j_\nu/\alpha_\nu=B_\nu(T)\)
が成り立つ。

\textbf{(3)} 希薄ガスや AGN
輝線領域では準位占有数が局所温度だけでは決まらず放射場 \(J_\nu\)
に依存する（非LTE）。すると \(j_\nu/\alpha_\nu\)
が物質温度だけの関数でなくなり \(S_\nu\ne B_\nu\) となる。

\textbf{答え}：平衡 \(dI_\nu/ds=0\) から
\(S_\nu=j_\nu/\alpha_\nu=B_\nu(T)\)。Kirchhoff
により比は物質温度だけの関数なので、LTE
なら放射場が黒体でなくても成立。非LTEだと崩れる。

\subsection*{問題
3-4　完全鏡の空洞と「平衡に必要な相互作用」}\label{ex-3-4}
\addcontentsline{toc}{subsection}{問題
3-4　完全鏡の空洞と「平衡に必要な相互作用」}

{[}★ 難易度：☆☆ {]}

§3.3
の問いは「壁が完全鏡（\(a_\nu=0\)）なら空洞放射はどうなるか」を扱った。これを現実天体の黒体化条件（§3.6）とつなぐ。

\begin{enumerate}
\def\labelenumi{\arabic{enumi}.}
\tightlist
\item
  完全鏡の空洞が黒体スペクトルを出さない理由を、平衡に達するメカニズムの有無から説明せよ。
\item
  レーザー共振器（高反射率ミラー）がこれに近いこと、逆に小穴を開けた空洞が実効的な完全吸収体になることを述べよ。
\item
  §3.6
  の条件1「光学的に厚い（\(\tau\gg1\)）」が、なぜ黒体への接近を保証するのか。完全鏡の例（吸収ゼロ）と対比して述べよ。
\end{enumerate}

\textbf{関連}：\hyperref[sec-universality-cavity]{§3.3
空洞放射（完全鏡の問い）}、\hyperref[sec-universality-real-objects]{§3.6
現実の天体が黒体に近づく条件}／光学的厚さと放射輸送は
\href{../part2/05-radiative-transfer.qmd}{第5章}。

\section{模範解答（問題
3-4）}\label{ux6a21ux7bc4ux89e3ux7b54ux554fux984c-3-4}

\textbf{(1)}
黒体スペクトルは「放射場と物質が相互作用して熱平衡に達した」結果である。完全鏡は光を吸収も放出もしないので、放射場と壁（物質）がエネルギーをやりとりできない。よって放射場は熱平衡に向かわず、内部のスペクトルは「最初に入っていた光」のまま決まらない。平衡に達するには少なくとも一部の吸収（相互作用）が不可欠。

\textbf{(2)}
レーザー共振器は高反射率ミラーで吸収が極小なので、定常的な黒体放射を出さず、特定モードの非平衡な光を蓄える。一方、空洞に小穴を開けると、入射光は穴から入って内部で何度も反射・吸収されほぼ確実に吸収される＝実効的に完全吸収体（\(a_\nu\to1\)）。同時に穴からは内部の熱平衡放射＝黒体スペクトルが出る。

\textbf{(3)}
光学的に厚い（\(\tau\gg1\)）とは、光が抜ける間に多数回の吸収・再放出を受けること。これにより放射場と物質が十分相互作用して局所的に熱化し、ソース関数が
\(B_\nu(T)\) に近づく（\hyperref[ex-3-3]{演習
3-3}）。完全鏡（吸収ゼロ、\(\tau=0\)）は相互作用が皆無で熱化しないのと正反対で、\(\tau\gg1\)
こそが黒体への接近を保証する。

\textbf{答え}：吸収ゼロ＝相互作用なし＝熱化せず黒体にならない。小穴空洞は実効完全吸収体で黒体を出す。\(\tau\gg1\)
は多数回の吸収再放出で熱化を保証し \(B_\nu\) に近づける。

\section*{さらに学ぶための参考文献}\label{sec-universality-further}
\addcontentsline{toc}{section}{さらに学ぶための参考文献}

\markright{さらに学ぶための参考文献}

\begin{itemize}
\tightlist
\item
  Rybicki \& Lightman, \emph{Radiative Processes in Astrophysics}
  (Wiley, 1979) --- 1.5--1.6 節（Kirchhoff の法則と LTE）
\item
  Mihalas, \emph{Stellar Atmospheres} (Freeman, 2nd ed., 1978) --- LTE
  と非 LTE の体系的扱い
\item
  Reif, \emph{Fundamentals of Statistical and Thermal Physics}
  (McGraw-Hill, 1965) --- 熱力学第二法則と空洞放射
\end{itemize}

\part{第II部　放射をどう記述するか}

\chapter{第4章　放射場の基本量}\label{ux7b2c4ux7ae0-ux653eux5c04ux5834ux306eux57faux672cux91cf}

\textbf{この章の主題}：放射場を記述するための基本量を、観測装置で何を測っているのかという視点から整理する。比強度・立体角・エネルギー密度・フラックス・それらのモーメント。これらは連続スペクトルにも線スペクトルにも共通の言語を与え、第II部全体の土台となる。

本章は第I部で扱ってきた「観測されるスペクトル」を、これから定量的に扱うための\textbf{記述の枠組み}を整える章であり、本書の哲学に立ち返れば「①
逆引き」の橋渡し位置にある。

\section*{この章の中心地図}\label{sec-radiation-field-map}
\addcontentsline{toc}{section}{この章の中心地図}

\markright{この章の中心地図}

（この章は地図の再掲なし。第II部全体は、観測されるスペクトルを記述するための言語を整える章）

\begin{tcolorbox}[enhanced jigsaw, arc=.35mm, bottomrule=.15mm, bottomtitle=1mm, breakable, colback=white, colbacktitle=quarto-callout-note-color!10!white, colframe=quarto-callout-note-color-frame, coltitle=black, left=2mm, leftrule=.75mm, opacityback=0, opacitybacktitle=0.6, rightrule=.15mm, title=\textcolor{quarto-callout-note-color}{\faInfo}\hspace{0.5em}{ノート}, titlerule=0mm, toprule=.15mm, toptitle=1mm]

\textbf{方針}：本章で導入する記号と定義は、第III部以降のすべての逆引きで使い続ける。本章は概念定義の章なので、難易度は基本的に
☆☆ で統一する。

\end{tcolorbox}

\section*{この章で答える問い}\label{sec-radiation-field-questions}
\addcontentsline{toc}{section}{この章で答える問い}

\markright{この章で答える問い}

\begin{itemize}
\tightlist
\item
  「明るさ」と「光度」と「フラックス」と「比強度」は何がどう違うのか
\item
  観測装置（望遠鏡＋検出器）が直接測っているのは、放射場のどの量なのか
\item
  なぜ立体角という概念が、放射の記述に必須なのか
\item
  同じ放射場を表す量が、振動数表示と波長表示で形が変わるのはなぜか
\item
  連続スペクトルも線スペクトルも、同じ基本量（\(I_\nu\), \(J_\nu\)
  など）で記述できるのか
\end{itemize}

\section*{到達目標}\label{sec-radiation-field-goals}
\addcontentsline{toc}{section}{到達目標}

\markright{到達目標}

この章を読み終えると、読者は次のことができるようになる：

\begin{itemize}
\tightlist
\item
  エネルギー密度・フラックス・比強度・モーメントの定義と相互関係を述べられる
\item
  観測量がこれらのどれに対応するかを判別できる
\item
  同じ枠組みが連続スペクトルにも線スペクトルにも適用できることを説明できる
\end{itemize}

\begin{center}\rule{0.5\linewidth}{0.5pt}\end{center}

\section{4.1 比強度 ―
放射場の最も基本的な量}\label{sec-radiation-field-intensity}

{[}本文目安：B3{]}

放射場を記述する基本量はいくつもあるが、本書ではそのうちの一つ ―
\textbf{比強度} \(I_\nu\)（\textbf{specific intensity}）―
を出発点とする。比強度は次のように定義される：

\begin{quote}
ある位置 \(\mathbf{r}\) で、ある方向 \(\mathbf{n}\)
に向かって、ある立体角内に、ある振動数帯にわたって、単位面積・単位時間あたりに通過するエネルギー
\end{quote}

形式的には、空間に固定された面要素 \(dA\)（法線方向
\(\hat{\mathbf{k}}\)）、時間素 \(dt\)、\(\hat{\mathbf{k}}\) から角度
\(\theta\) の方向に開いた立体角素 \(d\Omega\)、振動数素 \(d\nu\)
あたりのエネルギー \(dE\) として：

\begin{equation}\protect\phantomsection\label{eq-specific-intensity-def}{
dE = I_\nu(\mathbf{r}, \mathbf{n}, t)\, \cos\theta\, dA\, dt\, d\Omega\, d\nu
}\end{equation}

単位は erg s\(^{-1}\) cm\(^{-2}\) Hz\(^{-1}\) sr\(^{-1}\)（CGS）。

\begin{figure}[H]

\centering{

\includegraphics[width=0.85\linewidth,height=\textheight,keepaspectratio]{index_files/mediabag/part2/../images/specific-intensity.pdf}

}

\caption{\label{fig-specific-intensity}比強度 \(I_\nu\)
の定義の幾何。面要素 \(dA\) の法線方向 \(\hat{\mathbf{k}}\) から角度
\(\theta\) の方向 \(\mathbf{n}\) に向かう光線束を、立体角素
\(d\Omega\)・振動数素 \(d\nu\) に区切って測る。\textbf{\(\cos\theta\)
は「面に垂直に通過する成分」だけを取り出す射影因子}
で、第2章で簡略に扱っていた『面を通過するエネルギー』の正式な定義に対応する。}

\end{figure}%

\begin{tcolorbox}[enhanced jigsaw, arc=.35mm, bottomrule=.15mm, bottomtitle=1mm, breakable, colback=white, colbacktitle=quarto-callout-note-color!10!white, colframe=quarto-callout-note-color-frame, coltitle=black, left=2mm, leftrule=.75mm, opacityback=0, opacitybacktitle=0.6, rightrule=.15mm, title=\textcolor{quarto-callout-note-color}{\faInfo}\hspace{0.5em}{ノート}, titlerule=0mm, toprule=.15mm, toptitle=1mm]

\textbf{用語}：\(\cos\theta\) という\textbf{射影因子}は、「面要素 \(dA\)
を通過する光線の有効面積は
\(dA \cos\theta\)」という幾何の帰結である。\(dA\)
の取り方を\textbf{光線方向に垂直}にする流儀（\(dA_\perp = dA \cos\theta\)）もあり、その場合は
\(\cos\theta\)
が定義式から消える。両流儀は等価だが、本書では「空間に固定された面」を使う前者を採る。これによりフラックスの定義式
式~\ref{eq-flux-def} に \(\cos\theta\) が自然に引き継がれる。

\end{tcolorbox}

\begin{tcolorbox}[enhanced jigsaw, arc=.35mm, bottomrule=.15mm, bottomtitle=1mm, breakable, colback=white, colbacktitle=quarto-callout-note-color!10!white, colframe=quarto-callout-note-color-frame, coltitle=black, left=2mm, leftrule=.75mm, opacityback=0, opacitybacktitle=0.6, rightrule=.15mm, title=\textcolor{quarto-callout-note-color}{\faInfo}\hspace{0.5em}{ノート}, titlerule=0mm, toprule=.15mm, toptitle=1mm]

\textbf{方針}：\(I_\nu\) は
\textbf{位置・方向・時刻・振動数の関数}である。第I部で見たプランク関数
\(B_\nu(T)\) は、この \(I_\nu\) の一つの特別な場合 ― 「\textbf{温度
\(T\) の熱平衡放射場の比強度}」 ―
にほかならない。本書を通じて、観測量と物理を結ぶ橋として \(I_\nu\)
を最も基本的な量として扱う。

\end{tcolorbox}

なぜ \(I_\nu\) を出発点にするか。理由は次の三つ：

\begin{enumerate}
\def\labelenumi{\arabic{enumi}.}
\tightlist
\item
  \textbf{観測装置との対応が明快}：望遠鏡が空のある方向を見て、ある振動数帯のエネルギーを単位時間に拾うとき、その値は本質的に
  \(I_\nu\) である。
\item
  \textbf{保存則として性質がよい}：何もない真空中を伝播する間、\(I_\nu\)
  は経路に沿って一定に保たれる（\textbf{比強度の保存則}）。これは観測の基本前提になる。
\item
  \textbf{他のすべての量が \(I_\nu\)
  から導ける}：フラックス・エネルギー密度・放射圧などはすべて \(I_\nu\)
  の角度モーメントとして表される（節~\ref{sec-radiation-field-moments}
  参照）。
\end{enumerate}

\begin{tcolorbox}[enhanced jigsaw, arc=.35mm, bottomrule=.15mm, breakable, colback=white, colframe=quarto-callout-tip-color-frame, left=2mm, leftrule=.75mm, opacityback=0, rightrule=.15mm, toprule=.15mm]
\begin{minipage}[t]{5.5mm}
\textcolor{quarto-callout-tip-color}{\faLightbulb}
\end{minipage}%
\begin{minipage}[t]{\textwidth - 5.5mm}

\textbf{問い}：比強度は真空中の伝播でなぜ保存されるのか？

\textbf{短答}：「光線一本あたりのエネルギーが、伝播方向に何もぶつからない限り変わらない」というシンプルな保存則。位相空間体積の保存（Liouville
の定理の光学版）の帰結でもある。

\textbf{もう一歩}：これにより、距離 \(d\)
の遠方の天体から地球で観測する比強度 \(I_\nu\)
は、天体表面での比強度と等しい。「遠くの星ほど暗く見える」と感じるのは、距離で減衰する\textbf{フラックス}
\(F_\nu = I_\nu \cdot \Omega\) の方であって（立体角
\(\Omega \propto 1/d^2\)）、\(I_\nu\)
そのものではない。表面輝度は距離によらないという表現と等価である。

\end{minipage}%
\end{tcolorbox}

\section{4.2 立体角 ―
観測装置との橋渡し}\label{sec-radiation-field-solid-angle}

{[}本文目安：B2{]}

比強度の定義に \textbf{立体角}（\textbf{solid angle}）\(d\Omega\)
が現れる以上、これを正確に扱えないと放射場の記述はできない。立体角は二次元球面上の「面積」に相当する量で、極座標で

\begin{equation}\protect\phantomsection\label{eq-solid-angle-element}{
d\Omega = \sin\theta\, d\theta\, d\varphi
}\end{equation}

と書ける。全球面の立体角は \(4\pi\) sr である。

天文学で出会う代表的な立体角の見積もり：

{\def\LTcaptype{none} % do not increment counter
\begin{longtable}[]{@{}
  >{\raggedright\arraybackslash}p{(\linewidth - 4\tabcolsep) * \real{0.3333}}
  >{\raggedright\arraybackslash}p{(\linewidth - 4\tabcolsep) * \real{0.3333}}
  >{\raggedright\arraybackslash}p{(\linewidth - 4\tabcolsep) * \real{0.3333}}@{}}
\toprule\noalign{}
\begin{minipage}[b]{\linewidth}\raggedright
天体
\end{minipage} & \begin{minipage}[b]{\linewidth}\raggedright
視直径 \(\theta\)
\end{minipage} & \begin{minipage}[b]{\linewidth}\raggedright
立体角 \(\Omega \simeq \pi(\theta/2)^2\)
\end{minipage} \\
\midrule\noalign{}
\endhead
\bottomrule\noalign{}
\endlastfoot
太陽（地球から） & 約 32 arcmin & \(\sim 6.8 \times 10^{-5}\) sr \\
月 & 約 31 arcmin & \(\sim 6.4 \times 10^{-5}\) sr \\
木星 & 約 50 arcsec & \(\sim 4.6 \times 10^{-8}\) sr \\
アルファ・ケンタウリ A（最近傍恒星） & 約 9 mas &
\(\sim 1.5 \times 10^{-15}\) sr \\
全天 & \(4\pi\) sr & \(\simeq 12.57\) sr \\
\end{longtable}
}

（各立体角は円形視野 \(\Omega = \pi(\theta/2)^2\)
で計算した。視直径から立体角を求める統一公式である。）

立体角がないと「強度」と「フラックス」が混乱する。\textbf{点光源}（立体角が限りなく小さい）と\textbf{広がった光源}（立体角が有限）では、観測装置との関係が本質的に違う。

\begin{tcolorbox}[enhanced jigsaw, arc=.35mm, bottomrule=.15mm, bottomtitle=1mm, breakable, colback=white, colbacktitle=quarto-callout-note-color!10!white, colframe=quarto-callout-note-color-frame, coltitle=black, left=2mm, leftrule=.75mm, opacityback=0, opacitybacktitle=0.6, rightrule=.15mm, title=\textcolor{quarto-callout-note-color}{\faInfo}\hspace{0.5em}{ノート}, titlerule=0mm, toprule=.15mm, toptitle=1mm]

\textbf{用語}：天文学では「立体角」「ビーム角」「視野角」など類義語が並ぶ。詳しい用語は
\href{https://astro-dic.jp}{astro-dic.jp}
の標準訳を参照。本書では「立体角」（solid angle）を統一して用いる。

\end{tcolorbox}

\section{4.3 エネルギー密度 ―
放射場のエネルギー含有量}\label{sec-radiation-field-energy-density}

{[}本文目安：B3{]}

ある場所での放射場の\textbf{単位体積あたりのエネルギー}を、振動数別に表したのが\textbf{エネルギー密度}（\textbf{energy
density}）\(u_\nu\) である。比強度 \(I_\nu\) と次の関係で結ばれる：

\begin{equation}\protect\phantomsection\label{eq-energy-density-def}{
u_\nu(\mathbf{r}, t) = \frac{1}{c} \oint I_\nu(\mathbf{r}, \mathbf{n}, t)\, d\Omega
}\end{equation}

ここで積分は全立体角 \(4\pi\) にわたる。意味は単純で、\(I_\nu/c\)
が「ある方向から単位立体角あたりに供給されるエネルギー密度」になっており、それを全方向に足し上げている。

\textbf{等方放射場}（\textbf{isotropic radiation field}, \(I_\nu\)
が方向に依らない）の場合、式~\ref{eq-energy-density-def} は

\begin{equation}\protect\phantomsection\label{eq-energy-density-isotropic}{
u_\nu = \frac{4\pi}{c} I_\nu
}\end{equation}

となる。たとえば熱平衡の放射場 ― 黒体放射 ― では \(I_\nu = B_\nu(T)\)
なので、

\begin{equation}\protect\phantomsection\label{eq-energy-density-bb}{
u_\nu^{\text{BB}}(T) = \frac{4\pi}{c} B_\nu(T)
}\end{equation}

である。これを全振動数で積分すると、放射場の全エネルギー密度
\(u = \int u_\nu d\nu\) が \(a T^4\) となる（\(a = 4\sigma/c\)）。これは
Stefan-Boltzmann 則の別表現で、第6章で詳しく扱う。

\begin{tcolorbox}[enhanced jigsaw, arc=.35mm, bottomrule=.15mm, bottomtitle=1mm, breakable, colback=white, colbacktitle=quarto-callout-tip-color!10!white, colframe=quarto-callout-tip-color-frame, coltitle=black, left=2mm, leftrule=.75mm, opacityback=0, opacitybacktitle=0.6, rightrule=.15mm, title=\textcolor{quarto-callout-tip-color}{\faLightbulb}\hspace{0.5em}{ヒント}, titlerule=0mm, toprule=.15mm, toptitle=1mm]

\textbf{目安}：CMB（\(T = 2.725\) K）の全エネルギー密度は
\(u = aT^4 \simeq 4.2 \times 10^{-13}\) erg
cm\(^{-3}\)。これは銀河系内の宇宙線や星間磁場のエネルギー密度（いずれも
\(\sim 10^{-12}\) erg
cm\(^{-3}\)）と同程度のスケールで、宇宙物理の重要な「ベースライン」になる。

\end{tcolorbox}

\section{4.4 フラックス ―
観測される「明るさ」}\label{sec-radiation-field-flux}

{[}本文目安：B3{]}

ある面を\textbf{正味で}通過する単位面積・単位時間あたりのエネルギーを表すのが
\textbf{フラックス}（\textbf{flux}）\(F_\nu\) である。\(I_\nu\)
から次のように構成される：

\begin{equation}\protect\phantomsection\label{eq-flux-def}{
F_\nu(\mathbf{r}, t) = \oint I_\nu(\mathbf{r}, \mathbf{n}, t) \cos\theta\, d\Omega
}\end{equation}

ここで \(\theta\) は面の法線と方向ベクトル \(\mathbf{n}\)
の間の角。重要なのは \(\cos\theta\)
の符号：法線方向に出ていく光は正、戻ってくる光は負を寄与する。

\textbf{完全に等方的な放射場}では、\(\oint \cos\theta\, d\Omega = 0\)
なので
\(F_\nu = 0\)。つまり熱平衡で閉じている空洞の内部からは、正味のフラックスは出ない。

\textbf{観測者から見た天体のフラックス}：地球から距離 \(d\)
の場所にある光度 \(L_\nu\) の天体からは、

\begin{equation}\protect\phantomsection\label{eq-flux-observed}{
F_\nu^{\text{obs}} = \frac{L_\nu}{4\pi d^2}
}\end{equation}

が観測される。これが「明るさ」と俗に呼ばれる量である。

\begin{tcolorbox}[enhanced jigsaw, arc=.35mm, bottomrule=.15mm, bottomtitle=1mm, breakable, colback=white, colbacktitle=quarto-callout-note-color!10!white, colframe=quarto-callout-note-color-frame, coltitle=black, left=2mm, leftrule=.75mm, opacityback=0, opacitybacktitle=0.6, rightrule=.15mm, title=\textcolor{quarto-callout-note-color}{\faInfo}\hspace{0.5em}{ノート}, titlerule=0mm, toprule=.15mm, toptitle=1mm]

\textbf{観測との対応}：望遠鏡が口径 \(A\) で立体角
\(\Omega_{\text{beam}}\) のビームを持つとき：

\begin{itemize}
\tightlist
\item
  \textbf{点光源}（角度サイズ
  \(< \Omega_{\text{beam}}\)）の場合：装置が拾うエネルギーは
  \(F_\nu \cdot A \cdot d\nu\)
\item
  \textbf{広がった光源}（角度サイズ
  \(> \Omega_{\text{beam}}\)）の場合：拾うエネルギーは
  \(I_\nu \cdot A \cdot \Omega_{\text{beam}} \cdot d\nu\)
\end{itemize}

つまり\textbf{点光源ではフラックスが、広がった光源では比強度が}、それぞれの場で意味を持つ。観測の現場で「何を測っているか」を意識すると、この区別が常に問題になる。

\end{tcolorbox}

\begin{tcolorbox}[enhanced jigsaw, arc=.35mm, bottomrule=.15mm, breakable, colback=white, colframe=quarto-callout-tip-color-frame, left=2mm, leftrule=.75mm, opacityback=0, rightrule=.15mm, toprule=.15mm]
\begin{minipage}[t]{5.5mm}
\textcolor{quarto-callout-tip-color}{\faLightbulb}
\end{minipage}%
\begin{minipage}[t]{\textwidth - 5.5mm}

\textbf{問い}：点光源と広がった光源の境界はどこで切り替わるのか？

\textbf{短答}：観測装置の角度分解能（ビーム角
\(\Omega_\text{beam}\)）と天体の角度サイズ \(\Omega_\text{source}\)
の比較で決まる。\(\Omega_\text{source} \ll \Omega_\text{beam}\)
なら点光源として扱える。

\textbf{もう一歩}：同じ天体でも装置によって扱いが変わる。たとえばアルファ・ケンタウリ
A は通常の地上望遠鏡では点光源だが、VLTI 干渉計（角度分解能 \(\sim 1\)
mas）では円盤として分解される。「点光源か広がった光源か」は天体固有の性質ではなく、\textbf{観測装置との関係}である。

\end{minipage}%
\end{tcolorbox}

\section{\texorpdfstring{4.5 比強度のモーメント ― \(J\), \(F\),
\(P\)}{4.5 比強度のモーメント ― J, F, P}}\label{sec-radiation-field-moments}

{[}本文目安：B3{]}

比強度 \(I_\nu\)
から導かれる重要な量は、次の三つの\textbf{角度モーメント}として統一的に扱える。

\begin{equation}\protect\phantomsection\label{eq-mean-intensity}{
J_\nu = \frac{1}{4\pi} \oint I_\nu\, d\Omega \quad \text{（平均強度、0 次モーメント）}
}\end{equation}

\begin{equation}\protect\phantomsection\label{eq-flux-moment}{
F_\nu = \oint I_\nu \cos\theta\, d\Omega \quad \text{（フラックス、1 次モーメント）}
}\end{equation}

\begin{equation}\protect\phantomsection\label{eq-radiation-pressure}{
P_\nu = \frac{1}{c} \oint I_\nu \cos^2\theta\, d\Omega \quad \text{（放射圧、2 次モーメント）}
}\end{equation}

これらの幾何学的意味は明快である：

\begin{itemize}
\tightlist
\item
  \(J_\nu\)（\textbf{平均強度}, \textbf{mean intensity}）は、\(I_\nu\)
  の全方向平均。エネルギー密度との関係は \(u_\nu = (4\pi/c) J_\nu\)。
\item
  \(F_\nu\)（\textbf{フラックス},
  \textbf{flux}）は、方向性をもったエネルギーの正味流。観測量。
\item
  \(P_\nu\)（\textbf{放射圧}, \textbf{radiation
  pressure}）は、面の単位面積あたりの運動量フラックス。恒星内部や降着円盤で重要。
\end{itemize}

\begin{tcolorbox}[enhanced jigsaw, arc=.35mm, bottomrule=.15mm, bottomtitle=1mm, breakable, colback=white, colbacktitle=quarto-callout-note-color!10!white, colframe=quarto-callout-note-color-frame, coltitle=black, left=2mm, leftrule=.75mm, opacityback=0, opacitybacktitle=0.6, rightrule=.15mm, title=\textcolor{quarto-callout-note-color}{\faInfo}\hspace{0.5em}{ノート}, titlerule=0mm, toprule=.15mm, toptitle=1mm]

\textbf{方針}：本書では特に断らない限り「\textbf{フラックス}」を
\(F_\nu\) で書く。文献によっては \(H_\nu = F_\nu/(4\pi)\) や
\(\pi F_\nu\)
という記法も使われるので、別文献を読むときは定義の確認が必要。

\end{tcolorbox}

熱平衡（等方）の場合、これらは簡単に：

{\def\LTcaptype{none} % do not increment counter
\begin{longtable}[]{@{}ll@{}}
\toprule\noalign{}
量 & 等方放射場での値 \\
\midrule\noalign{}
\endhead
\bottomrule\noalign{}
\endlastfoot
\(J_\nu\) & \(B_\nu(T)\) \\
\(F_\nu\) & \(0\) \\
\(P_\nu\) & \(\frac{1}{3} u_\nu = \frac{4\pi}{3c} B_\nu(T)\) \\
\end{longtable}
}

放射圧が \(u/3\) になるのは光子気体の状態方程式
\(P = u/3\)（光子は質量ゼロで等方的）を反映している。第6章・第8章で詳しく扱う。

\section{4.6
振動数表示と波長表示の変換}\label{sec-radiation-field-nu-lambda}

{[}本文目安：B3{]}

第2章 §2.2
でプランク関数について見たのと同じ理屈で、\textbf{任意のスペクトル量}は振動数表示と波長表示で関数形が異なる。鍵は\textbf{同じエネルギーを異なる変数で表す}という事実：

\begin{equation}\protect\phantomsection\label{eq-nu-lambda-equivalence}{
I_\nu\, |d\nu| = I_\lambda\, |d\lambda|
}\end{equation}

これと \(\nu = c/\lambda\) から
\(|d\nu/d\lambda| = c/\lambda^2\)、よって

\begin{equation}\protect\phantomsection\label{eq-nu-lambda-conversion-i}{
I_\lambda = I_\nu \cdot \frac{c}{\lambda^2}, \qquad I_\nu = I_\lambda \cdot \frac{\lambda^2}{c}
}\end{equation}

同様にフラックスでも

\begin{equation}\protect\phantomsection\label{eq-nu-lambda-conversion-f}{
F_\lambda = F_\nu \cdot \frac{c}{\lambda^2}
}\end{equation}

である。エネルギー密度・モーメント全般について同じ関係が成り立つ。

\begin{tcolorbox}[enhanced jigsaw, arc=.35mm, bottomrule=.15mm, bottomtitle=1mm, breakable, colback=white, colbacktitle=quarto-callout-note-color!10!white, colframe=quarto-callout-note-color-frame, coltitle=black, left=2mm, leftrule=.75mm, opacityback=0, opacitybacktitle=0.6, rightrule=.15mm, title=\textcolor{quarto-callout-note-color}{\faInfo}\hspace{0.5em}{ノート}, titlerule=0mm, toprule=.15mm, toptitle=1mm]

\textbf{注意（観測）}：観測装置の波長分解能・振動数分解能の区別は、本質的にここから来る。電波天文学（GHz
帯）では振動数表示が、光学天文学（nm
帯）では波長表示が、それぞれ自然に好まれる。

\end{tcolorbox}

\section{4.7 異方性と等方性 ―
放射場の方向依存}\label{sec-radiation-field-anisotropy}

{[}本文目安：B3{]}

放射場が\textbf{等方的}であるとは、\(I_\nu\) が方向ベクトル
\(\mathbf{n}\)
に依らないということ。これは強い理想化条件で、現実の天体放射ではほぼ常に破られる。

実例で比較：

{\def\LTcaptype{none} % do not increment counter
\begin{longtable}[]{@{}ll@{}}
\toprule\noalign{}
場 & 異方性の程度 \\
\midrule\noalign{}
\endhead
\bottomrule\noalign{}
\endlastfoot
不透明体内部（深部） & 完全等方（LTE） \\
CMB & \(\Delta I / I \sim 10^{-5}\)（ほぼ等方） \\
太陽の表面光球から見る放射場 & 中心と縁で異なる（縁減光） \\
地球から見た太陽光 & 一方向に集中（点光源近似） \\
\end{longtable}
}

\textbf{等方放射場の特性}：等方なら
\(F_\nu = 0\)、\(J_\nu = I_\nu\)、\(P_\nu = (1/3)(4\pi/c) I_\nu\)
となり、量の間の関係が極端に単純になる。逆に言えば、\textbf{異方性こそが「方向性のあるエネルギー流」を生み出す}。観測される放射は必ず異方的である。

\begin{tcolorbox}[enhanced jigsaw, arc=.35mm, bottomrule=.15mm, bottomtitle=1mm, breakable, colback=white, colbacktitle=quarto-callout-tip-color!10!white, colframe=quarto-callout-tip-color-frame, coltitle=black, left=2mm, leftrule=.75mm, opacityback=0, opacitybacktitle=0.6, rightrule=.15mm, title=\textcolor{quarto-callout-tip-color}{\faLightbulb}\hspace{0.5em}{ヒント}, titlerule=0mm, toprule=.15mm, toptitle=1mm]

\textbf{観測との対応}：CMB の \(10^{-5}\)
という小さな異方性は、宇宙論的に巨大な情報量を持つ。本書では第16章で扱う。一方で、太陽光の高度な異方性は、地球で日射量を一方向のフラックスとして計算できることを意味する。

\end{tcolorbox}

\section{4.8 観測で何が測られるか ―
装置と放射場の対応}\label{sec-radiation-field-observation}

{[}本文目安：B3{]}

最後に、観測装置で実際に何が測られているかを整理する。これは本章全体を、観測の言葉に翻訳し直す節である。

\textbf{点光源}（角度サイズが装置のビームより小さい）の場合：

\begin{itemize}
\tightlist
\item
  観測される量は \textbf{フラックス} \(F_\nu\)（あるいはその波長表示
  \(F_\lambda\)）
\item
  単位エネルギー量は
  \(F_\nu \cdot d\nu \cdot A_\text{eff}\)（\(A_\text{eff}\)
  は望遠鏡の有効集光面積）
\item
  例：遠方銀河、QSO、恒星
\end{itemize}

\textbf{広がった光源}（角度サイズが装置のビームより大きい）の場合：

\begin{itemize}
\tightlist
\item
  観測される量は \textbf{比強度} \(I_\nu\)（あるいは輝度温度
  \(T_b\)、第2章 §2.6 で導入）
\item
  装置の角度分解能で空間的に分解できる
\item
  例：CMB、HII 領域、星間ダスト、銀河系
\end{itemize}

\textbf{スペクトロスコピー}：

\begin{itemize}
\tightlist
\item
  入射光を振動数（あるいは波長）に分解する装置を使う
\item
  各振動数チャンネルで \(F_\nu\) または \(I_\nu\) を測る
\item
  線スペクトルと連続スペクトルの双方が同じ装置で観測できる
\end{itemize}

本書を通じて、観測スペクトルの量と物理量を結ぶ橋を見ていくが、その橋はすべて
\(I_\nu\) から派生した量を介して張られる。

\begin{center}\rule{0.5\linewidth}{0.5pt}\end{center}

\section*{この章で何がわかったか}\label{sec-radiation-field-summary}
\addcontentsline{toc}{section}{この章で何がわかったか}

\markright{この章で何がわかったか}

\textbf{中心地図に戻る}

本章で、放射場を記述する\textbf{共通の言語}が整った：

\begin{itemize}
\tightlist
\item
  \textbf{比強度} \(I_\nu\)：放射場の最も基本的な量、方向と振動数に依存
\item
  \textbf{モーメント}（\(J_\nu\), \(F_\nu\),
  \(P_\nu\)）：物理的状況ごとに使い分け
\item
  \textbf{エネルギー密度} \(u_\nu\)：体積あたりのエネルギー含量
\item
  \textbf{観測量}：点光源ではフラックス、広がった光源では比強度
\end{itemize}

中心地図のプランク関数 \(B_\nu(T)\)
は、これらの記述枠組みの中で「\textbf{熱平衡の比強度}」として位置づけられた。今後第III部以降では、なぜ熱平衡で
\(I_\nu = B_\nu(T)\)
になるのか（光子統計・量子化・モード密度）を逆引きしていく。

\textbf{次章へ}：第5章では、これらの量がどう時間・空間的に変化するか ―
つまり光が物質と相互作用しながら伝播するときの\textbf{放射輸送方程式} ―
を扱う。この方程式は、観測スペクトルを「天体内部の歴史」と結びつける鍵であり、連続スペクトル（第VI部）と線スペクトル（第VII部）の両方の土台となる。

\section*{演習問題}\label{sec-radiation-field-exercises}
\addcontentsline{toc}{section}{演習問題}

\markright{演習問題}

\textbf{模範解答}
は折りたたみを展開して確認できる（オンライン版）。まず自力で解いてから確認すると効果的である。

\subsection*{問題 4-1　等方場で
uν=(4π/c)Iν、Fν=0、Pν=uν/3}\label{ex-4-1}
\addcontentsline{toc}{subsection}{問題 4-1　等方場で
uν=(4π/c)Iν、Fν=0、Pν=uν/3}

{[}★ 難易度：☆☆ {]}

§4.3〜4.5 は等方放射場での値を表で与えた。定義から自分で導く。立体角要素
\(d\Omega=\sin\theta\,d\theta\,d\varphi\)。等方なら \(I_\nu\)
は方向に依らず積分の外へ出せる。

\begin{enumerate}
\def\labelenumi{\arabic{enumi}.}
\tightlist
\item
  \(u_\nu=\dfrac{1}{c}\displaystyle\oint I_\nu\,d\Omega\) から、等方場で
  \(u_\nu=\dfrac{4\pi}{c}I_\nu\) を示せ。
\item
  \(F_\nu=\displaystyle\oint I_\nu\cos\theta\,d\Omega\)
  が等方場でゼロになることを示せ。
\item
  \(P_\nu=\dfrac{1}{c}\displaystyle\oint I_\nu\cos^2\theta\,d\Omega\)
  を計算し、\(\displaystyle\oint\cos^2\theta\,d\Omega=\dfrac{4\pi}{3}\)
  を使って \(P_\nu=\dfrac{1}{3}u_\nu\)
  を導け。これは光子気体の状態方程式 \(P=u/3\) に対応する。
\end{enumerate}

\textbf{関連}：\hyperref[sec-radiation-field-energy-density]{§4.3
エネルギー密度}、\hyperref[sec-radiation-field-moments]{§4.5
モーメント}／\(P=u/3\)
の光子気体としての扱いは\href{../part3/08-photon-statistics.qmd}{第8章}。

\section{模範解答（問題
4-1）}\label{ux6a21ux7bc4ux89e3ux7b54ux554fux984c-4-1}

\textbf{(1)} \(I_\nu\) を外に出すと
\(u_\nu=\dfrac{I_\nu}{c}\displaystyle\oint d\Omega=\dfrac{I_\nu}{c}\cdot4\pi=\dfrac{4\pi}{c}I_\nu\)。

\textbf{(2)}
\(\displaystyle\oint\cos\theta\,d\Omega=\int_0^{2\pi}d\varphi\int_0^{\pi}\cos\theta\sin\theta\,d\theta=2\pi\left[\dfrac{\sin^2\theta}{2}\right]_0^{\pi}=2\pi\cdot0=0\)。（上半球の正と下半球の負が打ち消す。）よって
\(F_\nu=I_\nu\cdot0=0\)。

\textbf{(3)}
\(\displaystyle\oint\cos^2\theta\,d\Omega=\int_0^{2\pi}d\varphi\int_0^{\pi}\cos^2\theta\sin\theta\,d\theta=2\pi\left[-\dfrac{\cos^3\theta}{3}\right]_0^{\pi}=2\pi\cdot\dfrac{2}{3}=\dfrac{4\pi}{3}\)。よって

\[
P_\nu=\frac{I_\nu}{c}\cdot\frac{4\pi}{3}=\frac{1}{3}\cdot\frac{4\pi}{c}I_\nu=\frac{1}{3}u_\nu.
\]

\textbf{答え}：等方場で
\(u_\nu=\frac{4\pi}{c}I_\nu\)、\(F_\nu=0\)、\(P_\nu=\frac13 u_\nu\)。最後は光子気体の状態方程式
\(P=u/3\)。

\subsection*{問題 4-2　表面輝度は距離によらない ― Iν 不変、Fν は 1/d²
に従う}\label{ex-4-2}
\addcontentsline{toc}{subsection}{問題 4-2　表面輝度は距離によらない ―
Iν 不変、Fν は 1/d² に従う}

{[}★ 難易度：☆☆ {]}

§4.1
は「遠くの星が暗く見えるのは比強度ではなくフラックスが減るから」と述べた。これを定量的に確かめる。

\begin{enumerate}
\def\labelenumi{\arabic{enumi}.}
\tightlist
\item
  真空中で比強度 \(I_\nu\) が経路に沿って保存することを前提に、距離
  \(d\) にある半径 \(R\) の一様な円盤天体について、観測される比強度
  \(I_\nu\) が \(d\) に依らないことを述べよ。
\item
  その天体が張る立体角が \(\Omega\simeq\pi R^2/d^2\propto1/d^2\)
  であることを示し、観測フラックス
  \(F_\nu\simeq I_\nu\,\Omega\propto1/d^2\) を導け。
\item
  点光源の式 \(F_\nu^\mathrm{obs}=L_\nu/4\pi d^2\) と (2)
  が整合することを、\(L_\nu\) を \(I_\nu\) と \(R\) で表して確かめよ。
\end{enumerate}

\textbf{関連}：\hyperref[sec-radiation-field-intensity]{§4.1
比強度（保存則）}、\hyperref[sec-radiation-field-flux]{§4.4
フラックス}／立体角は\hyperref[sec-radiation-field-solid-angle]{§4.2}、光度の定義は\hyperref[sec-reading-spectra-stefan-boltzmann]{§2.4}。

\section{模範解答（問題
4-2）}\label{ux6a21ux7bc4ux89e3ux7b54ux554fux984c-4-2}

\textbf{(1)}
比強度は真空伝播で保存する（光線一本あたりのエネルギーが不変、位相空間体積保存の帰結）。よって天体表面の比強度と、遠方で観測する比強度は等しく、距離
\(d\) に依らない（表面輝度の距離不変）。

\textbf{(2)} 半径 \(R\) の円盤を距離 \(d\) から見ると、見込む角の半径は
\(\theta\simeq R/d\)（\(R\ll d\)）。立体角は
\(\Omega\simeq\pi\theta^2=\pi R^2/d^2\propto1/d^2\)。観測フラックスは「比強度×天体の立体角」なので
\(F_\nu\simeq I_\nu\Omega=\pi I_\nu R^2/d^2\propto1/d^2\)。

\textbf{(3)} 表面フラックス
\(F_\nu^\mathrm{surf}=\pi I_\nu\)（半球射出、§2.4）、光度
\(L_\nu=4\pi R^2\cdot F_\nu^\mathrm{surf}=4\pi R^2\cdot\pi I_\nu=4\pi^2 R^2 I_\nu\)。点光源式に入れると

\[
F_\nu^\mathrm{obs}=\frac{L_\nu}{4\pi d^2}=\frac{4\pi^2R^2 I_\nu}{4\pi d^2}=\frac{\pi I_\nu R^2}{d^2},
\]

これは (2) と完全に一致する。\(I_\nu\) は不変、\(F_\nu\) だけが
\(1/d^2\) で落ちる。

\textbf{答え}：\(I_\nu\) は距離不変、\(\Omega\propto1/d^2\) ゆえ
\(F_\nu=\pi I_\nu R^2/d^2\propto1/d^2\)。点光源式とも一致。

\subsection*{問題 4-3　CMB のエネルギー密度を見積もる}\label{ex-4-3}
\addcontentsline{toc}{subsection}{問題 4-3　CMB
のエネルギー密度を見積もる}

{[}★ 難易度：☆ {]}

§4.3 は熱平衡で
\(u=\displaystyle\int u_\nu\,d\nu=aT^4\)（\(a=4\sigma/c\)）になると述べ、CMB
で \(u\simeq4.2\times10^{-13}\) erg cm\(^{-3}\)
と見積もった。これを確かめる。

\begin{enumerate}
\def\labelenumi{\arabic{enumi}.}
\tightlist
\item
  放射定数を CGS で \(a=7.566\times10^{-15}\) erg cm\(^{-3}\) K\(^{-4}\)
  とし、\(T_\mathrm{CMB}=2.725\) K で \(u=aT^4\) を計算せよ。
\item
  等方黒体場では \(u_\nu=\dfrac{4\pi}{c}B_\nu(T)\)、\(J_\nu=B_\nu(T)\)
  となる。なぜ \(J_\nu=B_\nu\) なのに \(u_\nu\) には \(4\pi/c\)
  が付くのか、定義の違いから一言で述べよ。
\item
  得られた \(u\)
  が、銀河系内の宇宙線・星間磁場のエネルギー密度（\(\sim10^{-12}\) erg
  cm\(^{-3}\)）と同程度であることを確認し、その含意を一言。
\end{enumerate}

\textbf{関連}：\hyperref[sec-radiation-field-energy-density]{§4.3
エネルギー密度}、\hyperref[sec-radiation-field-moments]{§4.5
平均強度}／\(u=aT^4\) の導出は\hyperref[ex-10-3]{演習
10-3}、\href{../part3/06-thermal-equilibrium.qmd}{第6章}。

\section{模範解答（問題
4-3）}\label{ux6a21ux7bc4ux89e3ux7b54ux554fux984c-4-3}

\textbf{(1)} \(T^4=(2.725)^4=55.1\) K\(^4\)。よって

\[
u=aT^4=7.566\times10^{-15}\times55.1\simeq4.2\times10^{-13}\ \mathrm{erg\,cm^{-3}}.
\]

本文の値と一致。

\textbf{(2)} \(J_\nu\)
は比強度の\textbf{方向平均}（\(\frac{1}{4\pi}\oint I_\nu d\Omega\)）なので等方場で
\(J_\nu=I_\nu=B_\nu\)。一方 \(u_\nu\)
は単位体積あたりの\textbf{エネルギー}（\(\frac1c\oint I_\nu d\Omega\)）で、全立体角の積分
\(4\pi\) と \(1/c\) が付く。平均か総和か、という定義の違い。実際
\(u_\nu=\frac{4\pi}{c}J_\nu\)。

\textbf{(3)} \(4.2\times10^{-13}\) と \(10^{-12}\) erg cm\(^{-3}\)
はオーダーで同程度。CMB・宇宙線・星間磁場のエネルギー密度が拮抗していることは、銀河系の星間媒質がこれらの間でエネルギーをやりとりする「準平衡」的な系であることを示唆する（重要なベースライン）。

\textbf{答え}：\(u\simeq4.2\times10^{-13}\) erg
cm\(^{-3}\)。\(J_\nu\)＝方向平均、\(u_\nu\)＝体積エネルギー（\(4\pi/c\)
倍）。宇宙線・磁場と同オーダー。

\subsection*{問題 4-4　点光源か広がった光源か ―
装置との関係で決まる}\label{ex-4-4}
\addcontentsline{toc}{subsection}{問題 4-4　点光源か広がった光源か ―
装置との関係で決まる}

{[}★ 難易度：☆☆ {]}

§4.4・4.8
は「点光源か広がった光源か」は天体固有でなく観測装置との関係で決まると述べた。\(\alpha\)
Cen A（視直径 \(\sim9\) mas）を例に確かめる。

\begin{enumerate}
\def\labelenumi{\arabic{enumi}.}
\tightlist
\item
  \(\alpha\) Cen A を円盤とみて立体角
  \(\Omega_\mathrm{src}\simeq\pi(\theta/2)^2\) を sr
  で概算せよ（\(\theta=9\) mas）。
\item
  通常の地上望遠鏡（シーイング \(\sim1\) arcsec）のビーム立体角
  \(\Omega_\mathrm{beam}\simeq\pi(0.5'')^2\)
  と比べ、点光源として扱える理由を述べよ。
\item
  VLTI 干渉計（分解能 \(\sim1\)
  mas）ではどうなるか。同じ天体で扱いが変わることが何を意味するか、観測される量（\(F_\nu\)
  か \(I_\nu\) か）の観点から述べよ。
\end{enumerate}

\textbf{関連}：\hyperref[sec-radiation-field-flux]{§4.4
フラックス（点光源/広がった光源）}、\hyperref[sec-radiation-field-observation]{§4.8
装置と放射場}／立体角の見積りは\hyperref[sec-radiation-field-solid-angle]{§4.2}。

\section{模範解答（問題
4-4）}\label{ux6a21ux7bc4ux89e3ux7b54ux554fux984c-4-4}

\textbf{(1)} \(\theta=9\) mas \(=9\times10^{-3}\) arcsec。1 arcsec
\(=4.85\times10^{-6}\) rad なので
\(\theta=9\times10^{-3}\times4.85\times10^{-6}=4.4\times10^{-8}\)
rad。半径 \(\theta/2=2.2\times10^{-8}\) rad。立体角

\[
\Omega_\mathrm{src}\simeq\pi(2.2\times10^{-8})^2\simeq1.5\times10^{-15}\ \mathrm{sr}.
\]

（本文の表の値と一致。）

\textbf{(2)} シーイング 1 arcsec のビーム半径 \(0.5''=2.4\times10^{-6}\)
rad、\(\Omega_\mathrm{beam}\simeq\pi(2.4\times10^{-6})^2\simeq1.8\times10^{-11}\)
sr。\(\Omega_\mathrm{src}/\Omega_\mathrm{beam}\sim10^{-4}\ll1\)。天体はビームよりはるかに小さく、空間的に分解できない＝点光源。測れるのはフラックス
\(F_\nu\) のみ。

\textbf{(3)} VLTI の分解能 \(\sim1\) mas は天体視直径 9 mas
より小さいので、円盤を空間的に分解できる＝広がった光源として扱え、面上の比強度
\(I_\nu\)（表面輝度分布）が測れる。同じ天体が装置の分解能次第で点光源にも広がった光源にもなる。すなわち「点か広がりか」は天体の性質でなく装置との関係で決まる。

\textbf{答え}：\(\Omega_\mathrm{src}\simeq1.5\times10^{-15}\)
sr。シーイング1″では \(\ll\Omega_\mathrm{beam}\) ゆえ点光源（\(F_\nu\)
を測る）。VLTI 1 masでは分解でき広がった光源（\(I_\nu\)）。装置依存。

\section*{さらに学ぶための参考文献}\label{sec-radiation-field-further}
\addcontentsline{toc}{section}{さらに学ぶための参考文献}

\markright{さらに学ぶための参考文献}

\begin{itemize}
\tightlist
\item
  Rybicki \& Lightman, \emph{Radiative Processes in Astrophysics}
  (Wiley, 1979) --- 第1章「Fundamentals of Radiative Transfer」
\item
  Mihalas, \emph{Stellar Atmospheres} (Freeman, 2nd ed., 1978) ---
  放射場の量の体系的整理
\item
  Chandrasekhar, \emph{Radiative Transfer} (Dover, 1960) ---
  古典的だが現在も標準
\end{itemize}

\chapter{第5章　放射輸送の基礎}\label{ux7b2c5ux7ae0-ux653eux5c04ux8f38ux9001ux306eux57faux790e}

\textbf{この章の主題}：観測されるスペクトルは「天体内部のある瞬間の状態」を直接見ているのではなく、「光が経た歴史」を映している。この歴史を記述する方程式が\textbf{放射輸送方程式}である。本章では、吸収・放出・散乱の三つの基本過程を整理し、光学的厚さとソース関数の概念を導入したうえで、第I部で見た「熱平衡で黒体になる」「光学的に厚いと黒体に近づく」という観測事実が、放射輸送方程式の極限としてどう現れるかを示す。

\section*{この章の中心地図}\label{sec-radiative-transfer-map}
\addcontentsline{toc}{section}{この章の中心地図}

\markright{この章の中心地図}

（この章は地図の再掲なし。第II部全体は、観測されるスペクトルを記述するための言語を整える章）

\begin{tcolorbox}[enhanced jigsaw, arc=.35mm, bottomrule=.15mm, bottomtitle=1mm, breakable, colback=white, colbacktitle=quarto-callout-note-color!10!white, colframe=quarto-callout-note-color-frame, coltitle=black, left=2mm, leftrule=.75mm, opacityback=0, opacitybacktitle=0.6, rightrule=.15mm, title=\textcolor{quarto-callout-note-color}{\faInfo}\hspace{0.5em}{ノート}, titlerule=0mm, toprule=.15mm, toptitle=1mm]

\textbf{方針}：本章で導入する放射輸送方程式は、第VI部の連続スペクトル応用、および第VII部の線スペクトル全章で繰り返し使う\textbf{最も汎用的な道具}である。本章の最後では、この方程式が連続と線の両方を同じ枠組みで扱えることを明示する。

\end{tcolorbox}

\section*{この章で答える問い}\label{sec-radiative-transfer-questions}
\addcontentsline{toc}{section}{この章で答える問い}

\markright{この章で答える問い}

\begin{itemize}
\tightlist
\item
  観測されるスペクトルが「天体内部の何か」を直接見ているのではなく、「光が経た歴史」を見ているのはなぜか
\item
  光学的厚さが大きい／小さいで、観測されるスペクトル形がどう変わるのか
\item
  なぜ「光学的に厚い」と黒体に近づくのか
\item
  連続吸収と線吸収は、同じ輸送方程式で扱えるのか
\item
  線スペクトルの形は、放射輸送のどの段階で決まるのか
\end{itemize}

\section*{到達目標}\label{sec-radiative-transfer-goals}
\addcontentsline{toc}{section}{到達目標}

\markright{到達目標}

この章を読み終えると、読者は次のことができるようになる：

\begin{itemize}
\tightlist
\item
  放射輸送方程式を、観測スペクトルの形と結びつけて読める
\item
  光学的に厚い極限で比強度がソース関数に等しくなる帰結を導ける
\item
  連続と線がどちらも同じ輸送方程式に乗ることを理解し、両者の違いを吸収係数の構造で説明できる
\end{itemize}

\begin{center}\rule{0.5\linewidth}{0.5pt}\end{center}

\section{5.1 放射と物質の相互作用 ―
吸収・放出・散乱}\label{sec-radiative-transfer-interactions}

{[}本文目安：B3-B4{]}

光が物質の中を伝播するとき、本質的に三つの相互作用が起きる：

\begin{enumerate}
\def\labelenumi{\arabic{enumi}.}
\tightlist
\item
  \textbf{吸収}（\textbf{absorption}）：光子が消え、エネルギーが物質に渡る（熱化、または励起）
\item
  \textbf{放出}（\textbf{emission}）：物質が光子を作り出す（脱励起、熱的放射、自由電子の制動放射など）
\item
  \textbf{散乱}（\textbf{scattering}）：光子の方向（時に振動数）が変わる
\end{enumerate}

これらをそれぞれ係数として表す：

{\def\LTcaptype{none} % do not increment counter
\begin{longtable}[]{@{}
  >{\raggedright\arraybackslash}p{(\linewidth - 6\tabcolsep) * \real{0.2500}}
  >{\raggedright\arraybackslash}p{(\linewidth - 6\tabcolsep) * \real{0.2500}}
  >{\raggedright\arraybackslash}p{(\linewidth - 6\tabcolsep) * \real{0.2500}}
  >{\raggedright\arraybackslash}p{(\linewidth - 6\tabcolsep) * \real{0.2500}}@{}}
\toprule\noalign{}
\begin{minipage}[b]{\linewidth}\raggedright
過程
\end{minipage} & \begin{minipage}[b]{\linewidth}\raggedright
係数
\end{minipage} & \begin{minipage}[b]{\linewidth}\raggedright
単位
\end{minipage} & \begin{minipage}[b]{\linewidth}\raggedright
意味
\end{minipage} \\
\midrule\noalign{}
\endhead
\bottomrule\noalign{}
\endlastfoot
吸収 & \(\alpha_\nu\) & cm\(^{-1}\) &
単位経路長あたりの吸収率（absorption coefficient） \\
放出 & \(j_\nu\) & erg s\(^{-1}\) cm\(^{-3}\) Hz\(^{-1}\) sr\(^{-1}\) &
単位体積・単位立体角あたりの放出係数（emission coefficient） \\
散乱 & \(\sigma_\nu\) & cm\(^{-1}\) &
単位経路長あたりの散乱率（scattering coefficient） \\
\end{longtable}
}

吸収係数 \(\alpha_\nu\)（\textbf{absorption
coefficient}）は、しばしば\textbf{質量吸収係数}（\textbf{mass absorption
coefficient}, opacity）\(\kappa_\nu\)（cm\(^2\) g\(^{-1}\)）と密度
\(\rho\) の積として \(\alpha_\nu = \kappa_\nu \rho\) と書かれる。

\begin{tcolorbox}[enhanced jigsaw, arc=.35mm, bottomrule=.15mm, bottomtitle=1mm, breakable, colback=white, colbacktitle=quarto-callout-note-color!10!white, colframe=quarto-callout-note-color-frame, coltitle=black, left=2mm, leftrule=.75mm, opacityback=0, opacitybacktitle=0.6, rightrule=.15mm, title=\textcolor{quarto-callout-note-color}{\faInfo}\hspace{0.5em}{ノート}, titlerule=0mm, toprule=.15mm, toptitle=1mm]

\textbf{方針}：本章では純粋な吸収・放出を主に扱い、散乱は形式的にしか入れない。散乱が本質的に効く場面（共鳴散乱、コンプトン化、Thomson
散乱）は第15章・第18章・第VII部で個別に扱う。

参考までに、散乱込みの放射輸送方程式は、\textbf{平面平行大気}を仮定して鉛直座標
\(z\) で書くと

\[
\mu\frac{dI_\nu}{dz} = -(\alpha_\nu + \sigma_\nu) I_\nu + j_\nu + \sigma_\nu \langle J_\nu \rangle
\]

の形になる。ここで \(\mu \equiv \cos\theta\) は光線方向と鉛直方向 \(z\)
のなす角の余弦で、光線経路 \(s\) と鉛直座標 \(z\) を \(dz = \mu\, ds\)
で結ぶことで現れる（光線そのものに沿って書く
式~\ref{eq-radiative-transfer} では \(\mu\)
は現れない）。右辺の最後の項が「他方向から来た光が散乱で本方向へ振り向けられる」効果（\(\langle J_\nu \rangle\)
は平均強度、散乱角分布が等方の場合）。これに対応するソース関数
\(S_\nu = (\alpha_\nu B_\nu + \sigma_\nu \langle J_\nu \rangle)/(\alpha_\nu + \sigma_\nu)\)
が第15章 §15.5 で再登場する。

\end{tcolorbox}

\begin{tcolorbox}[enhanced jigsaw, arc=.35mm, bottomrule=.15mm, bottomtitle=1mm, breakable, colback=white, colbacktitle=quarto-callout-note-color!10!white, colframe=quarto-callout-note-color-frame, coltitle=black, left=2mm, leftrule=.75mm, opacityback=0, opacitybacktitle=0.6, rightrule=.15mm, title=\textcolor{quarto-callout-note-color}{\faInfo}\hspace{0.5em}{ノート}, titlerule=0mm, toprule=.15mm, toptitle=1mm]

\textbf{方針（本書の焦点）}：本章および本書全体で扱う放射輸送は、放出源が\textbf{熱的}（Maxwell
分布の電子・LTE の原子）である場合を主に想定する。第1章 §1.9
で述べたように、本書は熱的放射に焦点を絞る。

非熱的放出源（シンクロトロン、逆 Compton など）でも放射輸送方程式
式~\ref{eq-radiative-transfer} の\textbf{形は同じ}である ― 違いは
\(j_\nu, \alpha_\nu\) の中身が、温度ではなく粒子分布関数 \(f(E)\)
で決まる点。詳しくは第18章。

\end{tcolorbox}

\section{5.2 光学的厚さ ―
経路の累積効果}\label{sec-radiative-transfer-optical-depth}

{[}本文目安：B3-B4{]}

光が経路 \(s\) を進むときの累積吸収を表すのが
\textbf{光学的厚さ}（\textbf{optical depth}）\(\tau_\nu\) である：

\begin{equation}\protect\phantomsection\label{eq-optical-depth}{
d\tau_\nu = \alpha_\nu\, ds, \qquad \tau_\nu = \int_0^s \alpha_\nu(s')\, ds'
}\end{equation}

光が物質を通過したあとに残る強度は、吸収のみが効く場合 \(e^{-\tau_\nu}\)
の減衰係数で抑えられる。直観的な解釈：

\begin{itemize}
\tightlist
\item
  \textbf{光学的に厚い}（\textbf{optically thick},
  \(\tau_\nu \gg 1\)）：光は完全に吸収・再放出を繰り返し、起源を忘れる
\item
  \textbf{光学的に薄い}（\textbf{optically thin},
  \(\tau_\nu \ll 1\)）：光はほぼ通過し、背景の光がほとんど見える
\item
  \textbf{平均自由行程}（\textbf{mean free
  path}）：\(\ell_{\mathrm{mfp}} = 1/\alpha_\nu\)
\end{itemize}

\begin{tcolorbox}[enhanced jigsaw, arc=.35mm, bottomrule=.15mm, bottomtitle=1mm, breakable, colback=white, colbacktitle=quarto-callout-tip-color!10!white, colframe=quarto-callout-tip-color-frame, coltitle=black, left=2mm, leftrule=.75mm, opacityback=0, opacitybacktitle=0.6, rightrule=.15mm, title=\textcolor{quarto-callout-tip-color}{\faLightbulb}\hspace{0.5em}{ヒント}, titlerule=0mm, toprule=.15mm, toptitle=1mm]

\textbf{目安}：太陽の光球で可視光の \(\tau \sim 1\)
となる深さが「光球表面」の定義。これが「太陽の表面」と呼ぶ場所である。地表の大気では晴天時に
\(\tau_\text{可視} \sim 0.1\)、曇天で \(\tau \sim 1\) 以上。

\end{tcolorbox}

\section{5.3 放射輸送方程式 ―
光の伝播の支配方程式}\label{sec-radiative-transfer-equation}

{[}本文目安：B3-B4{]}

吸収と放出を組み合わせると、比強度 \(I_\nu\) が経路 \(s\)
に沿ってどう変化するかが書ける：

\begin{equation}\protect\phantomsection\label{eq-radiative-transfer}{
\frac{dI_\nu}{ds} = -\alpha_\nu I_\nu + j_\nu
}\end{equation}

これが \textbf{放射輸送方程式}（\textbf{radiative transfer equation},
RTE）である。右辺第一項は吸収による損失、第二項は放出による獲得を表す。

光学的厚さを使って同じ式を書き直すと、

\begin{equation}\protect\phantomsection\label{eq-radiative-transfer-tau}{
\frac{dI_\nu}{d\tau_\nu} = -I_\nu + S_\nu
}\end{equation}

ただし \(S_\nu \equiv j_\nu / \alpha_\nu\) をここで定義した（次節）。

式~\ref{eq-radiative-transfer-tau}
は形式的に解ける一階線形常微分方程式で、

\begin{equation}\protect\phantomsection\label{eq-rt-formal-solution}{
I_\nu(\tau_\nu) = I_\nu(0)\, e^{-\tau_\nu} + \int_0^{\tau_\nu} S_\nu(\tau')\, e^{-(\tau_\nu - \tau')}\, d\tau'
}\end{equation}

第一項は「初期光がそのまま \(e^{-\tau}\)
倍に減衰したもの」、第二項は「経路上で放出された光が、観測点までの吸収で減衰しながら積算されたもの」と読める。\textbf{観測されるスペクトルが「光が経た歴史」を映す}とは、この第二項の意味である。

\section{5.4 ソース関数 ―
輸送方程式の心臓部}\label{sec-radiative-transfer-source}

{[}本文目安：B3-B4{]}

放射輸送方程式に現れた \textbf{ソース関数}（\textbf{source function}）

\begin{equation}\protect\phantomsection\label{eq-source-function}{
S_\nu \equiv \frac{j_\nu}{\alpha_\nu}
}\end{equation}

は、放射輸送において中心的な役割を果たす量である。形式解
式~\ref{eq-rt-formal-solution} を見ると、\(I_\nu\) は「ソース関数
\(S_\nu\) への緩和の途中にある」量と読める：

\begin{tcolorbox}[enhanced jigsaw, arc=.35mm, bottomrule=.15mm, bottomtitle=1mm, breakable, colback=white, colbacktitle=quarto-callout-note-color!10!white, colframe=quarto-callout-note-color-frame, coltitle=black, left=2mm, leftrule=.75mm, opacityback=0, opacitybacktitle=0.6, rightrule=.15mm, title=\textcolor{quarto-callout-note-color}{\faInfo}\hspace{0.5em}{ノート}, titlerule=0mm, toprule=.15mm, toptitle=1mm]

\textbf{用語}：英語 \emph{source function}
の日本語訳として「ソース関数」と「\textbf{源泉関数}」の両方が使われる（Mihalas
訳系では「源泉関数」、Rybicki-Lightman
翻訳系では「ソース関数」が多い）。指す量は同一。本書では「ソース関数」を主に用いる。

\end{tcolorbox}

\begin{itemize}
\tightlist
\item
  光学的に薄い極限（\(\tau \ll 1\)）：\(I_\nu \approx I_\nu(0)(1 - \tau_\nu) + S_\nu\, \tau_\nu\)
  ＝
  \(I_\nu(0) + \tau_\nu (S_\nu - I_\nu(0))\)（背景光の少しの吸収＋少しの放出。\(I_\nu(0)\)
  と \(S_\nu\) の大小関係で \textbf{吸収線か輝線か} が決まる）
\item
  光学的に厚い極限（\(\tau \gg 1\)）：\(I_\nu \to S_\nu\)（\textbf{ソース関数の値そのもの}になる）
\end{itemize}

つまり \(S_\nu\) とは「比強度が漸近する目標値」である。

\begin{tcolorbox}[enhanced jigsaw, arc=.35mm, bottomrule=.15mm, bottomtitle=1mm, breakable, colback=white, colbacktitle=quarto-callout-note-color!10!white, colframe=quarto-callout-note-color-frame, coltitle=black, left=2mm, leftrule=.75mm, opacityback=0, opacitybacktitle=0.6, rightrule=.15mm, title=\textcolor{quarto-callout-note-color}{\faInfo}\hspace{0.5em}{ノート}, titlerule=0mm, toprule=.15mm, toptitle=1mm]

\textbf{観測との対応}：観測される連続スペクトルが「黒体に近い」かどうかを判定するとき、私たちは暗黙のうちに「光学的に厚い」かつ「\(S_\nu = B_\nu(T)\)」が成り立つかを判定している。後者の条件こそが、次節で扱う
\textbf{LTE}（Local Thermodynamic Equilibrium）である。

\end{tcolorbox}

\begin{tcolorbox}[enhanced jigsaw, arc=.35mm, bottomrule=.15mm, breakable, colback=white, colframe=quarto-callout-tip-color-frame, left=2mm, leftrule=.75mm, opacityback=0, rightrule=.15mm, toprule=.15mm]
\begin{minipage}[t]{5.5mm}
\textcolor{quarto-callout-tip-color}{\faLightbulb}
\end{minipage}%
\begin{minipage}[t]{\textwidth - 5.5mm}

\textbf{問い}：ソース関数は常にプランク関数 \(B_\nu(T)\) なのか？

\textbf{短答}：いいえ。\(S_\nu = B_\nu(T)\) となるのは熱平衡 ―
より厳密には \textbf{LTE}（次節）― が成り立つときに限る。一般には
\(S_\nu \neq B_\nu\) で、その内容は物質の準位占有数 \(n_l\)
から決める必要がある。

\textbf{もう一歩}：第3章の Kirchhoff の法則 式~\ref{eq-kirchhoff}
は「熱平衡における吸収率 =
放出率」を主張するが、これは詳細釣り合いという\textbf{強い}条件である。これが破れる場面で何が起きるかを、次節
§5.5 で扱う。

\end{minipage}%
\end{tcolorbox}

\section{5.5 LTE と non-LTE ―
局所温度で状態が決まるか}\label{sec-radiative-transfer-lte}

{[}本文目安：B3-B4{]}

ソース関数 \(S_\nu = j_\nu / \alpha_\nu\) の値は、物質の状態 ―
具体的には準位占有数 \(n_l\) ―
に依存する。準位占有数の決まり方によって、放射輸送には決定的に異なる二つのレジームがある。それが
\textbf{LTE} と \textbf{non-LTE} である。

{\def\LTcaptype{none} % do not increment counter
\begin{longtable}[]{@{}
  >{\raggedright\arraybackslash}p{(\linewidth - 4\tabcolsep) * \real{0.3333}}
  >{\raggedright\arraybackslash}p{(\linewidth - 4\tabcolsep) * \real{0.3333}}
  >{\raggedright\arraybackslash}p{(\linewidth - 4\tabcolsep) * \real{0.3333}}@{}}
\toprule\noalign{}
\begin{minipage}[b]{\linewidth}\raggedright
略号
\end{minipage} & \begin{minipage}[b]{\linewidth}\raggedright
フル名
\end{minipage} & \begin{minipage}[b]{\linewidth}\raggedright
意味
\end{minipage} \\
\midrule\noalign{}
\endhead
\bottomrule\noalign{}
\endlastfoot
LTE & \textbf{L}ocal \textbf{T}hermodynamic \textbf{E}quilibrium &
\textbf{局所}熱平衡：物質の速度分布・電離・励起が、その点の局所温度
\(T\) だけで決まる \\
non-LTE & \textbf{not in LTE}（LTE でない） & LTE
が破れた状態。典型的には放射場 \(J_\nu\)
が準位を決めることで\textbf{非局所性}が入る \\
\end{longtable}
}

区別の鍵は「準位占有数がその点の \textbf{局所量}
だけで決まるか、それとも
\textbf{遠方の情報}（放射場経由）に依存するか」である。「non-」は単に「LTE
でない」の意味だが、その\textbf{典型的な起源}が放射場経由の非局所性にあると理解するとよい。

\subsection{LTE ―
局所温度だけで物質の状態が決まる}\label{lte-ux5c40ux6240ux6e29ux5ea6ux3060ux3051ux3067ux7269ux8ceaux306eux72b6ux614bux304cux6c7aux307eux308b}

ある点の物質の状態 ― 速度分布、電離度、励起準位の占有数 ―
が、\textbf{その点の局所的な温度 \(T\) だけで一意に決まる}近似が LTE
である。具体的には次が成り立つ：

\begin{itemize}
\tightlist
\item
  速度分布が Maxwell 分布（局所 \(T\)）
\item
  電離平衡が Saha 式（局所 \(T\)）
\item
  励起平衡が Boltzmann 分布（局所 \(T\)）
\end{itemize}

これらが満たされれば、第3章で見た Kirchhoff の法則 式~\ref{eq-kirchhoff}
によりソース関数は

\begin{equation}\protect\phantomsection\label{eq-source-lte}{
S_\nu = B_\nu(T) \quad \text{（LTE）}
}\end{equation}

となる。重要なのは、\textbf{LTE
は放射場が黒体である必要を含まない}点である。放射場 \(J_\nu\)
が異方的でも、\(J_\nu \neq B_\nu\)
でも、物質側の状態が局所温度で決まっていれば LTE
であり、ソース関数（放出と吸収の比）はちょうど \(B_\nu(T)\)
になる。放射場と物質が\textbf{両方}プランク的になっている強い状況は「\textbf{完全熱平衡}」（第6章）と呼んで区別する。

LTE
が成り立つのは、\textbf{衝突遷移率が放射遷移率より十分速い}とき。衝突は局所量（局所密度
\(n\)、局所温度
\(T\)）だけで決まるので、それが支配的なら占有数は局所量の関数になる。典型場面：

\begin{itemize}
\tightlist
\item
  高密度・高温度のプラズマ（恒星内部、光球深部）
\item
  早期宇宙（再結合期以前の電離プラズマ）
\end{itemize}

\subsection{non-LTE ―
放射場が占有数を変えることで非局所性が入る}\label{non-lte-ux653eux5c04ux5834ux304cux5360ux6709ux6570ux3092ux5909ux3048ux308bux3053ux3068ux3067ux975eux5c40ux6240ux6027ux304cux5165ux308b}

LTE の前提が崩れる ― すなわち準位占有数を局所温度だけで Boltzmann
分布に決められない ― 状態を \textbf{non-LTE}
と呼ぶ。多くの場面で、占有数は放射場 \(J_\nu\)
にも依存し、その結果\textbf{遠方の情報が入り込む}ことになる。

なぜ「非局所」になるのか。形式解 式~\ref{eq-rt-formal-solution}
が示すように、ある点での放射場 \(J_\nu(\mathbf{r})\)
は、\textbf{媒質全体の構造}（経路上の温度・密度・組成）を積分した量である。光は遠方からやってきて、その点の物質と相互作用する。したがって、放射遷移率（\(J_\nu\)
に比例する誘導吸収・誘導放出）が無視できないとき、占有数を決める\textbf{統計平衡方程式}は

\begin{equation}\protect\phantomsection\label{eq-statistical-equilibrium}{
\sum_{u} \bigl( C_{lu} + B_{lu} J_{\nu_{lu}} \bigr) n_l \;=\; \sum_{u} \bigl( C_{ul} + A_{ul} + B_{ul} J_{\nu_{lu}} \bigr) n_u
}\end{equation}

の形になる。これは \textbf{多準位原子}
に対する統計平衡方程式で、ある一つの下準位 \(l\) に注目し、その占有数
\(n_l\) の流入・流出を釣り合わせたものである。和 \(\sum_u\) は、\(l\)
と遷移で結ばれるすべての上準位 \(u\) にわたる総和を表す（\(\nu_{lu}\)
は遷移 \(l \leftrightarrow u\) の振動数）。左辺は \(l\) から各 \(u\)
への遷移による \(n_l\) の流出、右辺は各 \(u\) から \(l\)
への遷移による流入である（\(C_{ij}\) は衝突遷移率、\(A_{ul}, B_{ij}\) は
Einstein 係数。詳細は第11章で扱う）。最も簡単な
\textbf{二準位原子}（\(u\)
が一つだけ）では和が一項に縮約する。\textbf{右辺・左辺に
\(J_{\nu_{lu}}\) が入っていることが、典型的な non-LTE
で非局所性が現れる正体}である。占有数は媒質全体の構造の関数となり、その結果ソース関数も
\(B_\nu(T_\text{local})\) から外れる。

衝突率が無視できる極限が「強い
non-LTE」、衝突率が放射率より十分速い極限が LTE である。

\begin{tcolorbox}[enhanced jigsaw, arc=.35mm, bottomrule=.15mm, bottomtitle=1mm, breakable, colback=white, colbacktitle=quarto-callout-note-color!10!white, colframe=quarto-callout-note-color-frame, coltitle=black, left=2mm, leftrule=.75mm, opacityback=0, opacitybacktitle=0.6, rightrule=.15mm, title=\textcolor{quarto-callout-note-color}{\faInfo}\hspace{0.5em}{ノート}, titlerule=0mm, toprule=.15mm, toptitle=1mm]

\textbf{用語}：non-LTE という名前は「LTE
でない状態」という消極的な意味で用いられることが多い。Mihalas
以降の標準的解釈は、上記のように「準位占有数が放射場 \(J_\nu\)
を通じて非局所的になる」状態を指す。本書はこの解釈を採る。

\end{tcolorbox}

\subsection{どこで LTE
が破れるか}\label{ux3069ux3053ux3067-lte-ux304cux7834ux308cux308bux304b}

衝突率は密度に比例し、放射率は放射場に比例する。したがって\textbf{密度が低いほど
non-LTE が効きやすい}。

{\def\LTcaptype{none} % do not increment counter
\begin{longtable}[]{@{}ll@{}}
\toprule\noalign{}
領域 & LTE の妥当性 \\
\midrule\noalign{}
\endhead
\bottomrule\noalign{}
\endlastfoot
恒星内部・光球深部 & LTE がよい近似 \\
恒星上層大気 & しばしばnon-LTE が重要（特に強い線） \\
希薄な星間ガス（HII 領域、惑星状星雲） & 強いnon-LTE、禁制線が顕著 \\
AGN 輝線領域・QSO 散光雲 & 強いnon-LTE \\
早期宇宙（再結合期以前） & 完全 LTE → CMB の完璧な黒体性 \\
\end{longtable}
}

\begin{tcolorbox}[enhanced jigsaw, arc=.35mm, bottomrule=.15mm, bottomtitle=1mm, breakable, colback=white, colbacktitle=quarto-callout-note-color!10!white, colframe=quarto-callout-note-color-frame, coltitle=black, left=2mm, leftrule=.75mm, opacityback=0, opacitybacktitle=0.6, rightrule=.15mm, title=\textcolor{quarto-callout-note-color}{\faInfo}\hspace{0.5em}{ノート}, titlerule=0mm, toprule=.15mm, toptitle=1mm]

\textbf{注意（観測）}：「LTE か non-LTE
か」は観測スペクトルの解釈で常に最初に問わなければならない。第III部（統計力学）でこの区別の意味が深まり、第VII部の線スペクトルでは
non-LTE の扱いが本格的な主題になる（特に第21章 §21.7）。

\end{tcolorbox}

\begin{tcolorbox}[enhanced jigsaw, arc=.35mm, bottomrule=.15mm, breakable, colback=white, colframe=quarto-callout-tip-color-frame, left=2mm, leftrule=.75mm, opacityback=0, rightrule=.15mm, toprule=.15mm]
\begin{minipage}[t]{5.5mm}
\textcolor{quarto-callout-tip-color}{\faLightbulb}
\end{minipage}%
\begin{minipage}[t]{\textwidth - 5.5mm}

\textbf{問い}：non-LTE のスペクトルはプランク関数にはならないのか？

\textbf{短答}：そのまま読めば「ならない」。non-LTE ではソース関数が
\(S_\nu \neq B_\nu(T_\text{local})\)
なので、光学的に厚い領域から出てくる比強度も
\(I_\nu \to S_\nu \neq B_\nu(T_\text{local})\)
となり、観測スペクトルは局所温度のプランク関数から外れる。

\textbf{もう一歩}：ただし「\textbf{熱化深さ}（thermalization
depth）」と呼ばれる十分な深さまで沈み込んだ領域では、放射場 \(J_\nu\)
がその場所で熱化して \(J_\nu \to B_\nu(T_\text{local})\)
となり、結果として統計平衡方程式 式~\ref{eq-statistical-equilibrium}
の占有数も LTE 値に漸近、\(S_\nu \to B_\nu(T_\text{local})\)
となる。つまり\textbf{十分に深い non-LTE 媒質は実質 LTE}になる。non-LTE
効果が顕著に観測スペクトルに現れるのは、光学的に薄い領域や中程度の
\(\tau\) の上層大気である。詳しくは第VII部 §21.7 で扱う。

\end{minipage}%
\end{tcolorbox}

\section{5.6 光学的に厚い極限 ―
黒体に近づく仕組み}\label{sec-radiative-transfer-thick}

{[}本文目安：B3-B4{]}

第3章で「黒体放射の普遍性は熱平衡から来る」と論じたが、これは放射輸送方程式の極限として明示的に示せる。

LTE のもとで \(\tau_\nu \gg 1\)
の領域から出てくる光を考える。式~\ref{eq-rt-formal-solution}
の積分は、\(\tau' \sim \tau_\nu - 1\)
付近に強く重み付けされる（実効的な放出領域）。\(S_\nu = B_\nu(T)\)
がこの領域でほぼ一定なら、

\begin{equation}\protect\phantomsection\label{eq-rt-thick-limit}{
I_\nu \to B_\nu(T)
}\end{equation}

\textbf{これが黒体放射である}。恒星の光球が黒体に近いのは、まさにこの極限が成り立つ場所だからだ。

\begin{tcolorbox}[enhanced jigsaw, arc=.35mm, bottomrule=.15mm, bottomtitle=1mm, breakable, colback=white, colbacktitle=quarto-callout-tip-color!10!white, colframe=quarto-callout-tip-color-frame, coltitle=black, left=2mm, leftrule=.75mm, opacityback=0, opacitybacktitle=0.6, rightrule=.15mm, title=\textcolor{quarto-callout-tip-color}{\faLightbulb}\hspace{0.5em}{ヒント}, titlerule=0mm, toprule=.15mm, toptitle=1mm]

\textbf{観測との対応}：太陽の光球は可視光で \(\tau \sim 1\)
の付近を見ており、温度が \(\sim 5800\) K で LTE
近似がよく成り立つ。結果としてスペクトルは \(B_\nu(5800\text{ K})\)
に近い。第1章で見た「太陽スペクトルがほぼ黒体」の理由を、放射輸送の言葉で説明し直したことになる。

\end{tcolorbox}

\section{5.7 光学的に薄い極限 ―
放出の単純な積算}\label{sec-radiative-transfer-thin}

{[}本文目安：B3-B4{]}

逆の極限 \(\tau_\nu \ll 1\) では、式~\ref{eq-rt-formal-solution} は

\begin{equation}\protect\phantomsection\label{eq-rt-thin-limit}{
I_\nu \approx I_\nu(0) + \int j_\nu\, ds = I_\nu(0) + \tau_\nu S_\nu
}\end{equation}

となり、観測される強度は「背景光＋経路上の放出の総和」になる。背景光がなければ
\(I_\nu \approx \tau_\nu S_\nu\)
で、ガスの局所的な状態（密度・温度・組成）を直接反映する。

光学的に薄い極限が成り立つ場面：

\begin{itemize}
\tightlist
\item
  \textbf{希薄な星間ガス}：HII 領域、惑星状星雲、超新星残骸の輝線
\item
  \textbf{コロナや高温プラズマ}：X 線輝線
\item
  \textbf{AGN 輝線領域}：広輝線・狭輝線の放出
\item
  \textbf{CMB 観測}：宇宙再結合以降の光は基本的に光学的に薄い
\end{itemize}

このとき観測されるスペクトルは黒体ではなく、線スペクトルや特定の連続放射機構（自由--自由放射、シンクロトロン）の形になる。\textbf{観測スペクトルが黒体からどう「ずれるか」は、しばしばこの光学的に薄い極限の特徴で決まる}。

\begin{tcolorbox}[enhanced jigsaw, arc=.35mm, bottomrule=.15mm, breakable, colback=white, colframe=quarto-callout-tip-color-frame, left=2mm, leftrule=.75mm, opacityback=0, rightrule=.15mm, toprule=.15mm]
\begin{minipage}[t]{5.5mm}
\textcolor{quarto-callout-tip-color}{\faLightbulb}
\end{minipage}%
\begin{minipage}[t]{\textwidth - 5.5mm}

\textbf{問い}：背景光がなく \(\tau \ll 1\) なら、何も見えないのか？

\textbf{短答}：いや、\textbf{自己発光}が見える。式~\ref{eq-rt-thin-limit}
から \(I_\nu \approx \tau_\nu S_\nu = \int j_\nu\, ds\)
となり、経路上で物質が放出した光が薄く ― しかし確かに ― 見える。これが
HII 領域や星雲の輝線スペクトルの正体である。

\textbf{もう一歩}：見える明るさは \(\tau_\nu\)
に比例するので連続部分は暗いが、特定の振動数（線の中心 \(\nu_0\)）では
\(\alpha_\nu^\text{line}\) が大きく、\(\tau_\nu\)
が局所的に厚くなって\textbf{強い線}として見える。これが連続スペクトルがほとんどない場での「\textbf{輝線天体}」の見え方であり、第VII部で本格的に扱う線スペクトルの基本構造である。

\end{minipage}%
\end{tcolorbox}

\section{5.8 連続吸収と線吸収 ―
第VII部への橋渡し}\label{sec-radiative-transfer-bridge}

{[}本文目安：B3-B4{]}

ここまで放射輸送方程式 式~\ref{eq-radiative-transfer}
を、おもに連続スペクトル（黒体放射と「黒体からのずれ」）を念頭に書いてきた。しかしこの方程式は、観測されるスペクトルが連続か線かを区別していない。違いはすべて、吸収係数
\(\alpha_\nu\) と放出係数 \(j_\nu\) が振動数 \(\nu\)
に対してどんな構造を持つかに集約される。

\begin{itemize}
\tightlist
\item
  \textbf{連続吸収}：\(\alpha_\nu\) が \(\nu\)
  に対して滑らかに変化する成分（自由--自由・自由--束縛・ダストの吸収など）。第VI部で本格的に扱う。
\item
  \textbf{線吸収}：\(\alpha_\nu\) が特定の振動数 \(\nu_0\)
  の周りに鋭いピーク \(\phi(\nu - \nu_0)\)
  を持つ成分（束縛--束縛遷移）。これが線スペクトルの正体である。
\end{itemize}

つまり\textbf{線スペクトルは「連続スペクトルとは別の物理」ではなく、同じ放射場の同じ輸送方程式の上に乗る、\(\alpha_\nu\)
の構造が違うだけの現象}である。第VII部はこの観点から、線吸収係数を

\begin{equation}\protect\phantomsection\label{eq-line-absorption-preview}{
\alpha_\nu^{\text{line}} = \frac{\pi e^2}{m_e c}\, f_{lu}\, n_l\, \phi(\nu - \nu_0)
}\end{equation}

の形に分解し、各因子（\(\nu_0\), \(f_{lu}\), \(n_l\),
\(\phi\)）の背後にある原子物理・統計力学・線形機構を、それぞれ逆引きしていく。

\begin{tcolorbox}[enhanced jigsaw, arc=.35mm, bottomrule=.15mm, bottomtitle=1mm, breakable, colback=white, colbacktitle=quarto-callout-note-color!10!white, colframe=quarto-callout-note-color-frame, coltitle=black, left=2mm, leftrule=.75mm, opacityback=0, opacitybacktitle=0.6, rightrule=.15mm, title=\textcolor{quarto-callout-note-color}{\faInfo}\hspace{0.5em}{ノート}, titlerule=0mm, toprule=.15mm, toptitle=1mm]

\textbf{読みのコツ}：本書を通読するなら、第II部の輸送方程式を「連続と線の共通の土台」と意識しておくと、第VI部（連続のずれ）と第VII部（線スペクトルの逆引き）の対称性が見えやすい。

\end{tcolorbox}

\section{5.9
現実の天体スペクトルへの第一歩}\label{sec-radiative-transfer-real-spectra}

{[}本文目安：B3-B4{]}

放射輸送方程式は形式的にはシンプルだが、現実の天体スペクトルを再現するには、次の三つを具体的に与える必要がある：

\begin{enumerate}
\def\labelenumi{\arabic{enumi}.}
\tightlist
\item
  \textbf{吸収係数の振動数依存性} \(\alpha_\nu(\nu)\) ―
  連続成分と線成分の合成
\item
  \textbf{ソース関数の物理} \(S_\nu(\nu, \mathbf{r})\) ― LTE かnon-LTE
  か、占有数はどう決まるか
\item
  \textbf{境界条件} ― 「背景光」「天体の内部温度構造」など
\end{enumerate}

これら三つを与えるためには、本書の後半が必要になる：

\begin{itemize}
\tightlist
\item
  \textbf{吸収係数の中身}（式~\ref{eq-line-absorption-preview}
  の各因子）→ 第VII部の原子物理（第22章）と量子論（第IV部）
\item
  \textbf{ソース関数の構造}（占有数 \(n_l\) の決定）→ 第III部の統計力学
\item
  \textbf{モード密度の前因子}（\(2h\nu^3/c^2\)）→ 第V部の電磁気学
\item
  \textbf{熱史と境界条件}（早期宇宙の CMB、恒星の進化）→ 第VI部の応用
\end{itemize}

\begin{tcolorbox}[enhanced jigsaw, arc=.35mm, bottomrule=.15mm, bottomtitle=1mm, breakable, colback=white, colbacktitle=quarto-callout-tip-color!10!white, colframe=quarto-callout-tip-color-frame, coltitle=black, left=2mm, leftrule=.75mm, opacityback=0, opacitybacktitle=0.6, rightrule=.15mm, title=\textcolor{quarto-callout-tip-color}{\faLightbulb}\hspace{0.5em}{ヒント}, titlerule=0mm, toprule=.15mm, toptitle=1mm]

\textbf{観測との対応}：観測天文学者が実際にスペクトルを解析するとき、彼らは（明示的または暗黙のうちに）次の手順を踏んでいる：

\begin{enumerate}
\def\labelenumi{(\roman{enumi})}
\tightlist
\item
  観測スペクトルから「黒体に近い連続成分＋線特徴＋背景」に分解する
\item
  連続成分から温度・光度を引き出す（第2章の道具）
\item
  線特徴から組成・速度・密度を引き出す（第VII部の道具）
\end{enumerate}

本書はこの逆向きに ― 観測から各物理量を引き出す道具立てを ―
順に整える、というロードマップを取っている。

\end{tcolorbox}

\begin{center}\rule{0.5\linewidth}{0.5pt}\end{center}

\section*{この章で何がわかったか}\label{sec-radiative-transfer-summary}
\addcontentsline{toc}{section}{この章で何がわかったか}

\markright{この章で何がわかったか}

\textbf{中心地図に戻る}

本章で、本書全体を貫く\textbf{放射輸送方程式}

\[
\frac{dI_\nu}{d\tau_\nu} = -I_\nu + S_\nu
\]

が立ち上がった。この方程式は次の役割を果たす：

\begin{itemize}
\tightlist
\item
  観測される比強度 \(I_\nu\) を、光学的厚さ \(\tau_\nu\) と ソース関数
  \(S_\nu\) の関数として与える
\item
  \textbf{光学的に厚い + LTE} の極限で \(I_\nu \to B_\nu(T)\)
  となり、\textbf{第3章の「黒体は熱平衡で普遍」}が放射輸送の極限として再現される
\item
  \textbf{光学的に薄い}極限では、観測スペクトルは経路上の放出の単純な積算となる
\item
  連続と線の区別は、吸収係数 \(\alpha_\nu(\nu)\)
  の構造の違いに集約される（\textbf{第VII部への橋渡し}）
\end{itemize}

ここで第I部（観測の現象論）と第II部（放射場の記述と輸送）が一通り揃った。第III部からは、この上に立って、プランク関数の各因子を順に「逆引き」していく。

\textbf{次章へ}：第6章では、放射輸送のソース関数 \(S_\nu = B_\nu(T)\)
が成り立つための条件 ― \textbf{熱平衡} ―
を、改めて統計力学の言葉で取り出す。なぜ温度 \(T\)
だけが平衡放射場の形を決めるのか。なぜエントロピー最大の状態でプランク分布になるのか。これらは第III部全体を貫く問いになる。

\section*{演習問題}\label{sec-radiative-transfer-exercises}
\addcontentsline{toc}{section}{演習問題}

\markright{演習問題}

\textbf{模範解答}
は折りたたみを展開して確認できる（オンライン版）。まず自力で解いてから確認すると効果的である。

\subsection*{問題 5-1　放射輸送方程式の形式解を導く}\label{ex-5-1}
\addcontentsline{toc}{subsection}{問題
5-1　放射輸送方程式の形式解を導く}

{[}★ 難易度：☆☆ {]}

§5.3 は形式解
\(I_\nu(\tau_\nu)=I_\nu(0)e^{-\tau_\nu}+\displaystyle\int_0^{\tau_\nu}S_\nu(\tau')e^{-(\tau_\nu-\tau')}d\tau'\)
を結果として与えた。これを導く。出発点は
\(\dfrac{dI_\nu}{d\tau_\nu}=-I_\nu+S_\nu\)。

\begin{enumerate}
\def\labelenumi{\arabic{enumi}.}
\tightlist
\item
  積分因子 \(e^{\tau_\nu}\) を掛けて左辺を
  \(\dfrac{d}{d\tau_\nu}(I_\nu e^{\tau_\nu})\) にまとめよ。
\item
  \(0\) から \(\tau_\nu\) まで積分して形式解を導け。
\item
  第一項と第二項をそれぞれ物理的に解釈せよ（「初期光の減衰」「経路上の放出の積算」）。
\end{enumerate}

\textbf{関連}：\hyperref[sec-radiative-transfer-equation]{§5.3
放射輸送方程式}、\hyperref[sec-radiative-transfer-optical-depth]{§5.2
光学的厚さ}。

\section{模範解答（問題
5-1）}\label{ux6a21ux7bc4ux89e3ux7b54ux554fux984c-5-1}

\textbf{(1)} 両辺に \(e^{\tau_\nu}\)
を掛ける：\(e^{\tau_\nu}\dfrac{dI_\nu}{d\tau_\nu}+e^{\tau_\nu}I_\nu=e^{\tau_\nu}S_\nu\)。左辺は積の微分でまとまる：

\[
\frac{d}{d\tau_\nu}\left(I_\nu e^{\tau_\nu}\right)=S_\nu e^{\tau_\nu}.
\]

\textbf{(2)} \(0\) から \(\tau_\nu\) まで積分：

\[
I_\nu(\tau_\nu)e^{\tau_\nu}-I_\nu(0)=\int_0^{\tau_\nu}S_\nu(\tau')e^{\tau'}d\tau'.
\]

両辺を \(e^{\tau_\nu}\) で割る：

\[
I_\nu(\tau_\nu)=I_\nu(0)e^{-\tau_\nu}+\int_0^{\tau_\nu}S_\nu(\tau')e^{-(\tau_\nu-\tau')}d\tau'.
\]

\textbf{(3)} 第一項：入射光 \(I_\nu(0)\) が経路全体の吸収で
\(e^{-\tau_\nu}\) 倍に減衰したもの。第二項：各深さ \(\tau'\)
で放出された光 \(S_\nu d\tau'\) が、そこから観測点までの吸収
\(e^{-(\tau_\nu-\tau')}\)
で減衰しながら積算されたもの。観測スペクトルが「光が経た歴史」を映すとは、この第二項のこと。

\textbf{答え}：積分因子 \(e^{\tau}\) で
\(\frac{d}{d\tau}(Ie^\tau)=Se^\tau\)、積分して形式解。第一項＝初期光の減衰、第二項＝放出の減衰付き積算。

\subsection*{問題 5-2　厚い極限と薄い極限を形式解から}\label{ex-5-2}
\addcontentsline{toc}{subsection}{問題
5-2　厚い極限と薄い極限を形式解から}

{[}★ 難易度：☆☆ {]}

§5.6・5.7 は両極限を結果として示した。ソース関数 \(S_\nu\)
が経路上ほぼ一定の場合について、形式解（問題 5-1）から導け。

\begin{enumerate}
\def\labelenumi{\arabic{enumi}.}
\tightlist
\item
  一定 \(S_\nu\)
  で形式解の積分を実行し、\(I_\nu(\tau)=I_\nu(0)e^{-\tau}+S_\nu(1-e^{-\tau})\)
  を示せ。
\item
  光学的に厚い極限 \(\tau\gg1\) で \(I_\nu\to S_\nu\) を導け。
\item
  光学的に薄い極限 \(\tau\ll1\) で
  \(I_\nu\approx I_\nu(0)+\tau(S_\nu-I_\nu(0))\)
  を導け（\(e^{-\tau}\approx1-\tau\)）。
\end{enumerate}

\textbf{関連}：\hyperref[sec-radiative-transfer-thick]{§5.6
光学的に厚い極限}、\hyperref[sec-radiative-transfer-thin]{§5.7
光学的に薄い極限}。

\section{模範解答（問題
5-2）}\label{ux6a21ux7bc4ux89e3ux7b54ux554fux984c-5-2}

\textbf{(1)} \(S_\nu\) 一定なら積分は

\[
\int_0^{\tau}S_\nu e^{-(\tau-\tau')}d\tau'=S_\nu e^{-\tau}\int_0^{\tau}e^{\tau'}d\tau'=S_\nu e^{-\tau}(e^{\tau}-1)=S_\nu(1-e^{-\tau}).
\]

よって \(I_\nu(\tau)=I_\nu(0)e^{-\tau}+S_\nu(1-e^{-\tau})\)。

\textbf{(2)} \(\tau\gg1\) で \(e^{-\tau}\to0\) ゆえ
\(I_\nu\to S_\nu\)。初期光は完全に消え、比強度はソース関数の値そのものになる。

\textbf{(3)} \(\tau\ll1\) で \(e^{-\tau}\approx1-\tau\)：

\[
I_\nu\approx I_\nu(0)(1-\tau)+S_\nu\tau=I_\nu(0)+\tau(S_\nu-I_\nu(0)).
\]

\textbf{答え}：\(I_\nu=I_\nu(0)e^{-\tau}+S_\nu(1-e^{-\tau})\)。厚い→\(S_\nu\)、薄い→\(I_\nu(0)+\tau(S_\nu-I_\nu(0))\)。

\subsection*{問題 5-3　吸収線か輝線か ―
ソース関数と背景光の大小}\label{ex-5-3}
\addcontentsline{toc}{subsection}{問題 5-3　吸収線か輝線か ―
ソース関数と背景光の大小}

{[}★ 難易度：☆☆ {]}

問題 5-2(3) の薄い極限
\(I_\nu\approx I_\nu(0)+\tau_\nu(S_\nu-I_\nu(0))\)
を使って、線が吸収で見えるか輝線で見えるかを論じる。線の中心 \(\nu_0\)
では \(\alpha_\nu\) が大きく \(\tau_\nu\) が局所的に増す。

\begin{enumerate}
\def\labelenumi{\arabic{enumi}.}
\tightlist
\item
  \(S_\nu<I_\nu(0)\) のとき、線中心で \(I_\nu\)
  は連続光より大きいか小さいか。これは吸収線・輝線のどちらか。
\item
  \(S_\nu>I_\nu(0)\)（特に背景光 \(I_\nu(0)\approx0\)）のときはどうか。
\item
  太陽の Fraunhofer 吸収線（光球の外側ほど温度が下がる）と、HII
  領域の輝線スペクトル（背景光なし）が、それぞれどちらの条件に対応するか述べよ。
\end{enumerate}

\textbf{関連}：\hyperref[sec-radiative-transfer-source]{§5.4
ソース関数}、\hyperref[sec-radiative-transfer-thin]{§5.7
薄い極限（輝線天体）}／太陽の連続+線の二重構造は\hyperref[sec-where-blackbody-sun]{§1.2}、線スペクトルは\href{../part7/21-line-formation.qmd}{第21章}。

\section{模範解答（問題
5-3）}\label{ux6a21ux7bc4ux89e3ux7b54ux554fux984c-5-3}

\textbf{(1)} \(S_\nu<I_\nu(0)\) なら \(S_\nu-I_\nu(0)<0\)。線中心で
\(\tau_\nu\) が大きいほど \(I_\nu=I_\nu(0)+\tau_\nu(S_\nu-I_\nu(0))\)
は連続光 \(I_\nu(0)\) より \textbf{小さく}
なる。すなわち連続光を背景にした \textbf{吸収線} として見える。

\textbf{(2)} \(S_\nu>I_\nu(0)\) なら \(S_\nu-I_\nu(0)>0\)。線中心で
\(I_\nu\) は背景より \textbf{大きく}、\textbf{輝線}
になる。特に背景光がない（\(I_\nu(0)\approx0\)）と
\(I_\nu\approx\tau_\nu S_\nu\)
で、線中心だけ明るい純粋な輝線スペクトルになる。

\textbf{(3)} 太陽の Fraunhofer 線：明るい光球（連続光 \(I_\nu(0)\)
大）を背景に、外側の低温層（\(S_\nu=B_\nu(T_\mathrm{out})<I_\nu(0)\)）が吸収する
→ \(S_\nu<I_\nu(0)\)、吸収線。HII
領域：背景となる連続光がほぼなく、ガス自身の放出だけが見える →
\(S_\nu>I_\nu(0)\approx0\)、輝線。線の出方は線中心のソース関数と背景光の大小で決まる。

\textbf{答え}：\(S_\nu<I_\nu(0)\)→吸収線、\(S_\nu>I_\nu(0)\)→輝線。太陽
Fraunhofer＝前者、HII 領域＝後者。

\subsection*{問題 5-4　τ がほぼ 1
の層はなぜ「光球（見える表面）」か}\label{ex-5-4}
\addcontentsline{toc}{subsection}{問題 5-4　τ がほぼ 1
の層はなぜ「光球（見える表面）」か}

{[}★ 難易度：☆☆ {]}

§5.2・5.6 は「太陽の表面とは可視光で \(\tau\sim1\)
の深さ」と述べた。形式解の重み付けから理由を考える。

\begin{enumerate}
\def\labelenumi{\arabic{enumi}.}
\tightlist
\item
  外向きに観測する場合、形式解の第二項の重み \(e^{-(\tau_\nu-\tau')}\)
  は \(\tau'\)
  がどの辺りで最大か。なぜ「観測点から光学的厚さ約1の深さ」が実効的な放出層になるか述べよ。
\item
  平均自由行程 \(\ell_\mathrm{mfp}=1/\alpha_\nu\)
  を使い、「\(\tau\sim1\)
  は光が散乱・吸収されずに脱出できる最後の深さ」という描像で (1)
  を言い換えよ。
\item
  温度が深さとともに上がる恒星大気で、\(\tau\sim1\) の層が LTE
  なら観測スペクトルが \(B_\nu(T(\tau\!\sim\!1))\) に近づく理由を、問題
  5-2(2) と結びつけて述べよ。
\end{enumerate}

\textbf{関連}：\hyperref[sec-radiative-transfer-optical-depth]{§5.2
光学的厚さ}、\hyperref[sec-radiative-transfer-thick]{§5.6
厚い極限}／太陽光球が黒体に近い理由は\hyperref[sec-universality-real-objects]{§3.6}（\hyperref[ex-3-4]{演習
3-4}）。

\section{模範解答（問題
5-4）}\label{ux6a21ux7bc4ux89e3ux7b54ux554fux984c-5-4}

\textbf{(1)} 観測点（\(\tau_\nu=0\)
を外側にとる座標で考えると）に届く放出の重みは \(e^{-(\tau_\nu-\tau')}\)
で、観測点から光学的厚さが約1以内の層（\(\tau\lesssim1\)）からの寄与が支配的。それより深い層の光は手前の吸収
\(e^{-\tau}\) で強く減衰して届かない。だから「観測点から \(\tau\sim1\)
の深さ」が実効的な放出層＝「見える表面」になる。

\textbf{(2)} \(\ell_\mathrm{mfp}=1/\alpha_\nu\)
は光が一回吸収・散乱されずに進める距離。\(\tau\sim1\)
は観測点から測って平均自由行程ほぼ1回分の深さで、そこから出た光子はもう吸収されずに脱出できる。それより深い光子は途中で吸収・再放出され「直接は見えない」。よって
\(\tau\sim1\) が脱出できる最後の深さ＝光球。

\textbf{(3)} \(\tau\sim1\) 層が LTE なら
\(S_\nu=B_\nu(T(\tau\!\sim\!1))\)。その層は内側に向かって光学的に厚い（\(\tau\gg1\)）ので、問題
5-2(2) の厚い極限により出てくる比強度は
\(I_\nu\to S_\nu=B_\nu(T(\tau\!\sim\!1))\)。すなわち観測スペクトルは
\(\tau\sim1\) 層の温度のプランク関数に近づく。太陽が
\(B_\nu(5800\text{K})\) に近いのはこのため。

\textbf{答え}：重み \(e^{-(\tau-\tau')}\) ゆえ \(\tau\sim1\)
層が実効放出層＝脱出できる最後の深さ＝光球。そこが LTE＋厚いので
\(I_\nu\to B_\nu(T(\tau\!\sim\!1))\)。

\subsection*{問題 5-5　LTE で Sν=Bν ―
放射場は黒体でなくてよい}\label{ex-5-5}
\addcontentsline{toc}{subsection}{問題 5-5　LTE で Sν=Bν ―
放射場は黒体でなくてよい}

{[}★ 難易度：☆☆☆ {]}

§5.5 の微妙な点を整理する。LTE は「物質の状態が局所温度 \(T\)
だけで決まる」近似であって、放射場 \(J_\nu\)
が黒体である必要は含まない。

\begin{enumerate}
\def\labelenumi{\arabic{enumi}.}
\tightlist
\item
  Kirchhoff の法則が「比 \(j_\nu/\alpha_\nu\) を物質の局所温度 \(T\)
  だけの関数（\(=B_\nu(T)\)）にする」ことを使って、\(J_\nu\ne B_\nu\)
  でも LTE なら \(S_\nu=B_\nu(T)\) となる理由を述べよ。
\item
  「LTE」と「完全熱平衡」の違いを、（物質の状態）と（放射場の状態）の二つの観点で述べよ。
\item
  統計平衡方程式に放射場 \(J_\nu\) が入る（誘導吸収・放出
  \(B_{lu}J_\nu, B_{ul}J_\nu\)）と、なぜ占有数に「非局所性」が入るのか。逆に「熱化深さ」より深いとなぜ実質
  LTE に戻るのか述べよ。
\end{enumerate}

\textbf{関連}：\hyperref[sec-radiative-transfer-lte]{§5.5 LTE と
non-LTE}／Kirchhoff
とソース関数は\hyperref[sec-universality-kirchhoff]{§3.5}（\hyperref[ex-3-3]{演習
3-3}）、完全熱平衡は\href{../part3/06-thermal-equilibrium.qmd}{第6章}、統計平衡と
Einstein 係数は\hyperref[ex-11-1]{演習 11-1}。

\section{模範解答（問題
5-5）}\label{ux6a21ux7bc4ux89e3ux7b54ux554fux984c-5-5}

\textbf{(1)} ソース関数は \(S_\nu=j_\nu/\alpha_\nu\)。Kirchhoff
の法則により、物質が局所温度 \(T\) で熱化していれば、放出と吸収の比
\(j_\nu/\alpha_\nu\) は \textbf{その物質の温度 \(T\) だけの関数}
\(B_\nu(T)\) になる。\(j_\nu\) も \(\alpha_\nu\) も局所占有数（Boltzmann
分布、局所 \(T\)）で決まり、その比は周囲の放射場 \(J_\nu\)
に依らない。だから \(J_\nu\) が異方的でも \(J_\nu\ne B_\nu\)
でも、物質さえ LTE なら \(S_\nu=B_\nu(T)\)。

\textbf{(2)} LTE：\textbf{物質} の速度分布・電離・励起が局所 \(T\)
で決まる（Maxwell・Saha・Boltzmann）。放射場 \(J_\nu\)
は黒体でも等方でもなくてよい。完全熱平衡：\textbf{物質も放射場も両方}
が同一温度で平衡（\(J_\nu=B_\nu\) かつ等方、\(F_\nu=0\)）。LTE
は物質だけの条件、完全熱平衡は放射場まで含む強い条件。

\textbf{(3)} 誘導吸収・放出の率は \(J_\nu\)
に比例する。形式解が示すように \(J_\nu(\mathbf r)\)
は媒質全体の構造（経路上の温度・密度）を積分した非局所量。だから占有数を決める統計平衡方程式に
\(J_\nu\)
が入ると、占有数＝媒質全体の関数となり非局所性が入る。一方「熱化深さ」より深い領域では放射場がその場で熱化し
\(J_\nu\to B_\nu(T_\mathrm{local})\) となるため、統計平衡の解も
Boltzmann（局所値）に漸近し \(S_\nu\to B_\nu(T_\mathrm{local})\)、実質
LTE に戻る。non-LTE が観測に効くのは薄い領域や中程度 \(\tau\)
の上層大気。

\textbf{答え}：Kirchhoff が \(j/\alpha\) を局所 \(T\)
だけの関数にするので、放射場が黒体でなくても物質 LTE なら
\(S_\nu=B_\nu\)。完全熱平衡は放射場まで黒体な強い条件。\(J_\nu\)
経由で非局所性が入るが、熱化深さより深いと \(J_\nu\to B_\nu\) で実質
LTE。

\section*{さらに学ぶための参考文献}\label{sec-radiative-transfer-further}
\addcontentsline{toc}{section}{さらに学ぶための参考文献}

\markright{さらに学ぶための参考文献}

\begin{itemize}
\tightlist
\item
  Rybicki \& Lightman, \emph{Radiative Processes in Astrophysics}
  (Wiley, 1979) --- 第1章「Fundamentals of Radiative Transfer」
\item
  Mihalas, \emph{Stellar Atmospheres} (Freeman, 2nd ed., 1978) --- LTE /
  non-LTE の本格的扱い
\item
  Hubeny \& Mihalas, \emph{Theory of Stellar Atmospheres} (Princeton,
  2014) --- 現代的な放射輸送の決定版
\end{itemize}

\part{第III部　逆引き1：熱力学・統計力学へ}

\chapter{第6章　熱平衡とは何か}\label{ux7b2c6ux7ae0-ux71b1ux5e73ux8861ux3068ux306fux4f55ux304b}

\textbf{この章の主題}：第I部・第II部で「観測される連続スペクトルが黒体に近いのは熱平衡だから」「LTE
ではソース関数がプランク関数になる」と繰り返し述べてきた。その\textbf{熱平衡という条件}を、本章で改めて統計力学側から取り出す。温度・エントロピー・最大エントロピー原理を整理し、なぜ熱平衡という状況がスペクトルの形を一意に決めるのか、その理由を明らかにする。第III部はここから始まり、第7・8章で光子の統計に進んで、最終的にプランク分布の\textbf{占有因子}
\(1/(e^{h\nu/kT}-1)\) を導出する。

\section*{この章の中心地図}\label{sec-thermal-equilibrium-map}
\addcontentsline{toc}{section}{この章の中心地図}

\markright{この章の中心地図}

\textbf{中心地図：プランク関数 ― 本部で扱う因子}

\[
B_\nu(T) = \frac{2 h \nu^3}{c^2} \cdot \boxed{\frac{1}{e^{h\nu/kT} - 1}}
\]

強調されている \textbf{占有因子} \(1/(e^{h\nu/kT}-1)\)
が、第III部で逆引きする対象である。これは熱平衡にある光子（Bose
粒子）の統計から立ち上がる量で、低振動数極限では \(kT/(h\nu)\)
に、高振動数極限では \(e^{-h\nu/kT}\)
に漸近する。第8章で本格的に導出する。

\begin{tcolorbox}[enhanced jigsaw, arc=.35mm, bottomrule=.15mm, bottomtitle=1mm, breakable, colback=white, colbacktitle=quarto-callout-note-color!10!white, colframe=quarto-callout-note-color-frame, coltitle=black, left=2mm, leftrule=.75mm, opacityback=0, opacitybacktitle=0.6, rightrule=.15mm, title=\textcolor{quarto-callout-note-color}{\faInfo}\hspace{0.5em}{ノート}, titlerule=0mm, toprule=.15mm, toptitle=1mm]

\textbf{方針}：本章は中心地図のうち、特に\textbf{「この式が熱平衡で成り立つこと」そのもの}に焦点を絞る。なぜ熱平衡という条件が、放射場の分布を一意に決めてしまうのか。形そのもの（占有因子・前因子）は第7・8章で扱う。

\end{tcolorbox}

\section*{この章で答える問い}\label{sec-thermal-equilibrium-questions}
\addcontentsline{toc}{section}{この章で答える問い}

\markright{この章で答える問い}

\begin{itemize}
\tightlist
\item
  「熱平衡」という言葉が観測スペクトルにどう関係するのか
\item
  なぜ熱平衡にある放射場だけが、プランク分布の特別な形をとるのか
\item
  温度とは、観測の言葉では何を測っているのか
\item
  放射場が熱平衡に「近い」と判定するには、どの量を比較すればよいのか
\end{itemize}

\section*{到達目標}\label{sec-thermal-equilibrium-goals}
\addcontentsline{toc}{section}{到達目標}

\markright{到達目標}

この章を読み終えると、読者は次のことができるようになる：

\begin{itemize}
\tightlist
\item
  熱平衡という条件が、なぜ放射場の分布を一意に決めるのかを論じられる
\item
  温度・エントロピー・最大エントロピー原理の関係を整理できる
\item
  観測される CMB
  がほぼ完全な熱平衡である事実が何を意味するかを説明できる
\end{itemize}

\begin{center}\rule{0.5\linewidth}{0.5pt}\end{center}

\section{6.1 温度と熱平衡 ―
観測の言葉での定義}\label{sec-thermal-equilibrium-temperature}

{[}本文目安：B2{]}

「温度」という言葉は日常的に使われる。しかし統計力学の文脈で正確に扱うには、観測の言葉での定義が必要である。

\textbf{熱力学的な定義}：二つの系を熱接触させたとき、エネルギーの正味の流れが止まれば、両者は同じ\textbf{温度}
\(T\) にあると言う。

\textbf{統計力学的な定義}：熱平衡で実現される確率分布が
\textbf{Boltzmann 形}（\(\propto e^{-E/kT}\)）になるとき、その指数の係数
\(1/kT\) で温度を定める。第7章で本格的に扱う。

両者の同等性は、熱力学第二法則と統計力学の基礎仮定から導かれる。本章ではまず観測の側
― 「温度とは何が等しいか」 ― から入る。

\textbf{熱平衡の三条件}（第3章 §3.2 の再掲）：

\begin{enumerate}
\def\labelenumi{\arabic{enumi}.}
\tightlist
\item
  \textbf{詳細釣り合い}：振動数ごとに、吸収率と放出率が局所的に釣り合う
\item
  \textbf{温度の一意性}：系全体が単一の \(T\) で記述できる
\item
  \textbf{空間的等方性}：少なくとも局所的に放射場が等方的
\end{enumerate}

これらが満たされる状況が「熱平衡」であり、そこでの放射場のスペクトルは（第3章の空洞放射の議論から）\textbf{物質の詳細によらず一意}に決まる。本章の目的は、この「一意性」を統計力学の言葉で根拠づけることである。

\begin{tcolorbox}[enhanced jigsaw, arc=.35mm, bottomrule=.15mm, bottomtitle=1mm, breakable, colback=white, colbacktitle=quarto-callout-note-color!10!white, colframe=quarto-callout-note-color-frame, coltitle=black, left=2mm, leftrule=.75mm, opacityback=0, opacitybacktitle=0.6, rightrule=.15mm, title=\textcolor{quarto-callout-note-color}{\faInfo}\hspace{0.5em}{ノート}, titlerule=0mm, toprule=.15mm, toptitle=1mm]

\textbf{用語}：熱平衡（thermal
equilibrium）と熱力学的平衡（thermodynamic
equilibrium）はしばしば同じ意味で用いられるが、後者は化学平衡・力学的平衡も含むより広い概念。本書では放射と物質が同じ温度を共有する状態を「熱平衡」と呼ぶ。

\end{tcolorbox}

\section{6.2 マクロ状態とミクロ状態 ―
「同じ温度」の中身}\label{sec-thermal-equilibrium-macro-micro}

{[}本文目安：B3{]}

統計力学の出発点は、\textbf{マクロ状態}と\textbf{ミクロ状態}の区別である。

\begin{itemize}
\tightlist
\item
  \textbf{マクロ状態}（\textbf{macrostate}）：観測者が見る巨視的な記述（温度
  \(T\)、体積 \(V\)、粒子数 \(N\)、エネルギー \(E\)、\ldots）
\item
  \textbf{ミクロ状態}（\textbf{microstate}）：系を構成するすべての粒子・モードの個別の状態（古典なら座標と運動量の組、量子なら波動関数）
\end{itemize}

ひとつのマクロ状態に対して、多数のミクロ状態が対応する。この対応個数を\textbf{多重度}（\textbf{multiplicity},
状態数）\(\Omega\) と呼ぶ。

\textbf{例：気体分子の分布}

理想気体 \(N\) 個の分子を二つの体積に分けるとき、分子が左 \(N_L\)、右
\(N_R = N - N_L\) で分かれるマクロ状態の多重度は

\begin{equation}\protect\phantomsection\label{eq-multiplicity-gas}{
\Omega(N_L) = \binom{N}{N_L} = \frac{N!}{N_L!\, N_R!}
}\end{equation}

これは \(N_L = N/2\)
で最大になる。\textbf{等しく分かれる状態が圧倒的に多いミクロ状態を持つ}ことが、気体が均一に広がる理由である。

\textbf{例：放射場のモード}

第V部で本格的に扱うが、空洞中の電磁場は離散的なモードの集まりとみなせる。各モードに何個の光子が入っているかで、放射場のミクロ状態が決まる。同じ全エネルギー
\(E\)
を、これらのモードに\textbf{どう分配するか}でミクロ状態が変わる。マクロ状態（全エネルギー、温度）が決まっても、ミクロ状態には多くの選択肢がある。

\begin{tcolorbox}[enhanced jigsaw, arc=.35mm, bottomrule=.15mm, breakable, colback=white, colframe=quarto-callout-tip-color-frame, left=2mm, leftrule=.75mm, opacityback=0, rightrule=.15mm, toprule=.15mm]
\begin{minipage}[t]{5.5mm}
\textcolor{quarto-callout-tip-color}{\faLightbulb}
\end{minipage}%
\begin{minipage}[t]{\textwidth - 5.5mm}

\textbf{問い}：なぜマクロ状態だけを見て温度が決まるのか？
ミクロ状態を全部知らないとダメではないのか？

\textbf{短答}：熱平衡では、\textbf{ほぼすべてのミクロ状態が事実上同じマクロ状態}（最大多重度を持つマクロ状態）に対応する。だから観測者はマクロ状態だけで物理を語れる。

\textbf{もう一歩}：これは「\textbf{ほとんど確実に}」起きることを「\textbf{確実に}」起きると同一視する近似である。アボガドロ数
\(N \sim 10^{23}\) のスケールでは、揺らぎは \(1/\sqrt{N} \sim 10^{-11}\)
で、観測誤差より遥かに小さい。

\end{minipage}%
\end{tcolorbox}

\section{6.3 エントロピー ―
多重度の対数}\label{sec-thermal-equilibrium-entropy}

{[}本文目安：B3{]}

多重度 \(\Omega\)
を直接扱うと、桁が極端に大きくなって扱いにくい（\(10^{10^{23}}\)
など）。対数をとった量

\begin{equation}\protect\phantomsection\label{eq-entropy-boltzmann}{
S = k_B \ln \Omega
}\end{equation}

を \textbf{エントロピー}（\textbf{entropy}）と呼ぶ。これが Boltzmann
の関係である（彼の墓碑にこの式が刻まれている）。

エントロピーには重要な性質がある：

\begin{enumerate}
\def\labelenumi{\arabic{enumi}.}
\tightlist
\item
  \textbf{加法性}：独立な二つの系のエントロピーは和になる（\(\Omega = \Omega_1 \Omega_2\)
  から \(S = S_1 + S_2\)）
\item
  \textbf{拡張性}：粒子数 \(N\) に比例する（示量的）
\item
  \textbf{熱力学第二法則}：孤立系のエントロピーは時間とともに増加し、平衡で最大となる
\end{enumerate}

第三の性質が\textbf{最大エントロピー原理}（\textbf{principle of maximum
entropy}）であり、平衡状態のミクロ的特徴づけを与える。

\begin{tcolorbox}[enhanced jigsaw, arc=.35mm, bottomrule=.15mm, bottomtitle=1mm, breakable, colback=white, colbacktitle=quarto-callout-note-color!10!white, colframe=quarto-callout-note-color-frame, coltitle=black, left=2mm, leftrule=.75mm, opacityback=0, opacitybacktitle=0.6, rightrule=.15mm, title=\textcolor{quarto-callout-note-color}{\faInfo}\hspace{0.5em}{ノート}, titlerule=0mm, toprule=.15mm, toptitle=1mm]

\textbf{用語}：エントロピーは「乱雑さ」や「情報の欠如」と訳されることがあるが、定義の本質は「\textbf{マクロ状態に対応するミクロ状態の数の対数}」である。乱雑さや情報という比喩は、この定義から導かれる二次的な解釈である。

\end{tcolorbox}

\textbf{温度との関係}：エントロピー \(S\) をエネルギー \(E\)
の関数として書くと、温度はその偏微分で定まる：

\begin{equation}\protect\phantomsection\label{eq-temperature-entropy}{
\frac{1}{T} = \left(\frac{\partial S}{\partial E}\right)_{V, N}
}\end{equation}

つまり「温度」とは、エネルギーを少し加えたときにエントロピーがどれだけ増えるかの\textbf{逆数}である。これが熱力学的な温度の最も本質的な定義になる。

\section{6.4 最大エントロピー原理 ―
平衡分布の決定}\label{sec-thermal-equilibrium-max-entropy}

{[}本文目安：B3-B4｜導出：M1{]}

熱平衡状態は、孤立系のエントロピー \(S\)
が最大となる状態である。この単純な要請から、Boltzmann 分布や Planck
分布の関数形が一意に決まる。

\textbf{設定}：エネルギー準位 \(\{E_i\}\)
を持つ系を考え、各準位を占めるミクロ状態の数（あるいは粒子の個数）\(\{n_i\}\)
を未知数とする。これらは次の制約を満たす：

\[
\sum_i n_i = N \quad \text{（全粒子数）}
\]

\[
\sum_i n_i E_i = E \quad \text{（全エネルギー）}
\]

これらの制約のもとで多重度 \(\Omega(\{n_i\})\)
を最大化するのが最大エントロピー原理。具体的にはラグランジュの未定乗数法（乗数
\(\alpha, \beta\)）を使い、

\begin{equation}\protect\phantomsection\label{eq-lagrange}{
\frac{\partial}{\partial n_i}\left[ \ln \Omega - \alpha \sum_j n_j - \beta \sum_j n_j E_j \right] = 0
}\end{equation}

を解く。\(\Omega\)
の具体形は粒子の統計（古典・Bose・Fermi）で変わるが、いずれも

\[
n_i \propto f(\beta E_i, \alpha)
\]

の形になる。

\begin{itemize}
\tightlist
\item
  \textbf{古典統計}（Maxwell--Boltzmann）：\(n_i \propto e^{-(\alpha + \beta E_i)}\)（第7章）
\item
  \textbf{Bose 統計}：\(n_i = 1/(e^{\alpha + \beta E_i} - 1)\)（第8章）
\item
  \textbf{Fermi 統計}：\(n_i = 1/(e^{\alpha + \beta E_i} + 1)\)
\end{itemize}

ここで \(\beta = 1/k_B T\)
が温度の逆数になることが、式~\ref{eq-temperature-entropy}
と整合する。\(\alpha\) は化学ポテンシャル \(\mu\) と
\(\alpha = -\mu/k_B T\) で対応する。

\begin{tcolorbox}[enhanced jigsaw, arc=.35mm, bottomrule=.15mm, breakable, colback=white, colframe=quarto-callout-tip-color-frame, left=2mm, leftrule=.75mm, opacityback=0, rightrule=.15mm, toprule=.15mm]
\begin{minipage}[t]{5.5mm}
\textcolor{quarto-callout-tip-color}{\faLightbulb}
\end{minipage}%
\begin{minipage}[t]{\textwidth - 5.5mm}

\textbf{問い}：最大エントロピー原理は「単なる数学的処方箋」なのか？

\textbf{短答}：いいえ、物理的根拠がある。\textbf{ほとんどすべてのミクロ状態が最大多重度のマクロ状態に対応する}という事実が、最大エントロピーの状態が圧倒的に実現される理由である。

\textbf{もう一歩}：この原理は情報理論（Shannon,
Jaynes）の立場から再導出することもできる：「制約を満たすあらゆる確率分布のうち、最も先入観のない（情報量が最小の）分布が、熱平衡の分布である」。物理の最大エントロピーと情報理論の最大エントロピーは、本質的に同じことを言っている。

\end{minipage}%
\end{tcolorbox}

\section{6.5 放射場を含む熱平衡 ―
光子と物質の共通の温度}\label{sec-thermal-equilibrium-radiation}

{[}本文目安：B3{]}

物質の熱平衡は分子・原子の運動エネルギー分配で記述される。では\textbf{放射場}の熱平衡とは何か？

放射場を空洞中の電磁モードの集まりと見なすと（第V部・第14章）、各モードはエネルギー量子
\(h\nu\)
で何個の光子を保有するかという「占有数」で特徴づけられる。物質と放射場の熱平衡は、\textbf{両者が共通の温度
\(T\) を持つ}ことを意味する。

具体的には：

\begin{itemize}
\tightlist
\item
  物質の準位占有数：Boltzmann 分布 \(n_i \propto g_i e^{-E_i/k_B T}\)
\item
  放射場の各モードの平均光子数：Bose--Einstein 分布
  \(\bar{n}_\text{mode} = 1/(e^{h\nu/k_B T} - 1)\)
\item
  ソース関数 \(S_\nu = B_\nu(T)\)（Kirchhoff 第3章）
\end{itemize}

これらが\textbf{同じ \(T\)
で結ばれている}ことが「物質と放射場の熱平衡」の中身である。各々が違う温度なら、相互作用によりエネルギーが流れ、やがて共通の
\(T\) に達する。

\begin{tcolorbox}[enhanced jigsaw, arc=.35mm, bottomrule=.15mm, bottomtitle=1mm, breakable, colback=white, colbacktitle=quarto-callout-note-color!10!white, colframe=quarto-callout-note-color-frame, coltitle=black, left=2mm, leftrule=.75mm, opacityback=0, opacitybacktitle=0.6, rightrule=.15mm, title=\textcolor{quarto-callout-note-color}{\faInfo}\hspace{0.5em}{ノート}, titlerule=0mm, toprule=.15mm, toptitle=1mm]

\textbf{用語の整理}：

\begin{itemize}
\tightlist
\item
  \textbf{完全熱平衡}（complete thermodynamic
  equilibrium）：物質も放射場も\textbf{両方}が同じ温度 \(T\)
  で熱平衡。\(J_\nu = B_\nu(T)\) かつ物質の準位は Boltzmann。CMB
  の初期宇宙のように、両者が極めて長時間相互作用した状況で実現する。
\item
  \textbf{LTE}（局所熱平衡, 第5章
  §5.5）：物質の状態（速度分布・電離・励起）が局所温度 \(T\)
  で決まる近似。放射場 \(J_\nu\)
  は\textbf{必ずしも黒体である必要はない}。たとえば恒星大気では、深部では
  \(J_\nu \to B_\nu(T)\) だが、表層へ向かうと \(J_\nu\)
  は異方化しつつ物質は LTE のまま、という遷移領域がある。
\item
  \textbf{non-LTE}：上記 LTE の前提が崩れた状態。多くは放射場 \(J_\nu\)
  が占有数を決めることで非局所性が入る。
\end{itemize}

本書で「熱平衡」と言うときは通常 \textbf{完全熱平衡}を意味する。LTE
は弱い意味の「物質側だけの熱平衡近似」であり、両者を混同しないよう注意。

\end{tcolorbox}

\begin{tcolorbox}[enhanced jigsaw, arc=.35mm, bottomrule=.15mm, bottomtitle=1mm, breakable, colback=white, colbacktitle=quarto-callout-note-color!10!white, colframe=quarto-callout-note-color-frame, coltitle=black, left=2mm, leftrule=.75mm, opacityback=0, opacitybacktitle=0.6, rightrule=.15mm, title=\textcolor{quarto-callout-note-color}{\faInfo}\hspace{0.5em}{ノート}, titlerule=0mm, toprule=.15mm, toptitle=1mm]

\textbf{方針（本書の焦点）}：本章で導入する「熱平衡」は、粒子の\textbf{エネルギー分布が
Maxwell 分布}である状況、つまり\textbf{熱的}な粒子分布を前提とする。

ジェット・AGN
コロナ・宇宙線のような\textbf{非熱的}な粒子（べき則分布の高エネルギー電子）の集団では、温度という概念そのものが意味を失い、本章の枠組みはそのままでは適用できない。本書は熱的放射に焦点を絞るため、非熱的系の統計は第18章で軽く触れる程度に留める（第1章
§1.9 参照）。

なお「熱的 / 非熱的」（粒子の分布）と「LTE /
non-LTE」（原子準位の占有）は別の軸である ―
非熱的電子が出すシンクロトロンの周囲の原子ガスは独立に LTE / non-LTE
を取りうる。

\end{tcolorbox}

\section{6.6 詳細釣り合い ―
平衡の局所判定}\label{sec-thermal-equilibrium-detailed-balance}

{[}本文目安：B3-B4｜導出：M1{]}

熱平衡を「ミクロ的」に特徴づける条件が
\textbf{詳細釣り合い}（\textbf{detailed balance}）である：

\begin{quote}
平衡状態では、任意の二つの状態 \(i \to j\) への遷移確率と、その逆
\(j \to i\) への遷移確率が等しい。
\end{quote}

すなわち

\begin{equation}\protect\phantomsection\label{eq-detailed-balance}{
n_i\, W_{i \to j} = n_j\, W_{j \to i}
}\end{equation}

ここで \(W_{i \to j}\) は単位時間あたりの遷移確率、\(n_i\) は状態 \(i\)
の占有数。この関係は、各準位ペア \((i, j)\) について個別に成立する。

\textbf{なぜ詳細釣り合いが効くのか}：

熱平衡を「全粒子の正味の流れがゼロ」と定義しても矛盾はないが、それだけでは「ある複雑な閉じた経路を回る流れ」を排除できない。詳細釣り合いはこれより強い条件で、\textbf{任意のペアで個別に釣り合う}ことを要請する。

物理的には「時間反転対称性」と結びつく。マクロな熱平衡は、ミクロな時間反転対称性の帰結として理解できる。

\begin{tcolorbox}[enhanced jigsaw, arc=.35mm, bottomrule=.15mm, bottomtitle=1mm, breakable, colback=white, colbacktitle=quarto-callout-note-color!10!white, colframe=quarto-callout-note-color-frame, coltitle=black, left=2mm, leftrule=.75mm, opacityback=0, opacitybacktitle=0.6, rightrule=.15mm, title=\textcolor{quarto-callout-note-color}{\faInfo}\hspace{0.5em}{ノート}, titlerule=0mm, toprule=.15mm, toptitle=1mm]

\textbf{方針}：本章では詳細釣り合いを定義しただけだが、第11章で Einstein
がこの原理を放射と物質の相互作用に適用し、\textbf{プランク分布を再導出}する場面が出てくる。Einstein
の議論は、本書の中心地図のプランク関数を「もう一つの逆引きルート」で立ち上げる重要な道筋である。

\end{tcolorbox}

\begin{center}\rule{0.5\linewidth}{0.5pt}\end{center}

\section*{この章で何がわかったか}\label{sec-thermal-equilibrium-summary}
\addcontentsline{toc}{section}{この章で何がわかったか}

\markright{この章で何がわかったか}

\textbf{中心地図に戻る}

本章で、中心地図 \(B_\nu(T)\) の
\textbf{「式が熱平衡で成り立つ」という条件の中身}が言語化された：

\begin{itemize}
\tightlist
\item
  \textbf{温度} \(T\) は、エントロピーをエネルギーで微分したものの逆数
\item
  \textbf{熱平衡}
  は、最大エントロピー状態であり、詳細釣り合いが各遷移ペアで成り立つ
\item
  \textbf{物質と放射場の完全な熱平衡（CTE）} は、両者が共通の \(T\)
  で結ばれる状態。これは LTE（物質側だけが局所の \(T\) で Boltzmann
  分布、放射場は自由に流れる）より \textbf{強い条件} で、空洞放射や CMB
  のような閉じた系でのみ厳密に実現される（区別の整理は §6.5）
\end{itemize}

ただし「プランク分布の具体形」 ― 占有因子 \(1/(e^{h\nu/kT}-1)\)
がなぜこの形なのか ― はまだ示されていない。これには光子が「区別できない
Bose 粒子」であることを使う必要がある。

\textbf{次章へ}：第7章では、まず古典統計（Boltzmann
分布）を立ち上げ、分配関数と熱力学量の関係を整理する。古典統計の枠組みでどこまでわかり、どこで破綻するかを見たうえで、第8章で「光子の統計」へと進む。

\section*{演習問題}\label{sec-thermal-equilibrium-exercises}
\addcontentsline{toc}{section}{演習問題}

\markright{演習問題}

\textbf{模範解答}
は折りたたみを展開して確認できる（オンライン版）。まず自力で解いてから開くこと。

\subsection*{問題
6-1　多重度の最大と「マクロ状態で十分」な理由}\label{ex-6-1}
\addcontentsline{toc}{subsection}{問題
6-1　多重度の最大と「マクロ状態で十分」な理由}

{[}★ 難易度：☆☆ {]}

§6.2 は気体分子を左右に分ける多重度 \(\Omega(N_L)=\binom{N}{N_L}\) が
\(N_L=N/2\) で最大になると述べた。これを示し、揺らぎの幅を見積もる。

\begin{enumerate}
\def\labelenumi{\arabic{enumi}.}
\tightlist
\item
  Stirling 近似 \(\ln N!\simeq N\ln N-N\) を使い、\(\ln\Omega(N_L)\)
  を最大化して \(N_L=N/2\) を導け。
\item
  \(N_L=N/2+x\) のまわりで \(\ln\Omega\) を \(x\)
  の2次まで展開し、\(\ln\Omega\simeq\ln\Omega_{\max}-\dfrac{2x^2}{N}\)
  を示せ（分布が幅 \(\sim\sqrt{N}\) のガウス型）。
\item
  相対的な揺らぎ \(\sqrt{\langle x^2\rangle}/N\sim1/\sqrt{N}\)
  が、\(N\sim10^{23}\) で \(10^{-11}\)
  程度になることを確かめ、「マクロ状態だけで温度が語れる」理由を述べよ。
\end{enumerate}

\textbf{関連}：\hyperref[sec-thermal-equilibrium-macro-micro]{§6.2
マクロ状態とミクロ状態}／揺らぎの一般論は\hyperref[sec-statistical-mechanics-fluctuations]{§7.4}（\hyperref[ex-7-3]{演習
7-3}）。

\section{模範解答（問題
6-1）}\label{ux6a21ux7bc4ux89e3ux7b54ux554fux984c-6-1}

\textbf{(1)} \(\ln\Omega=\ln N!-\ln N_L!-\ln(N-N_L)!\)。Stirling
で微分（\(N\) 固定）：

\[
\frac{d\ln\Omega}{dN_L}=-\ln N_L+\ln(N-N_L)=\ln\frac{N-N_L}{N_L}=0\ \Rightarrow\ N_L=N-N_L\ \Rightarrow\ N_L=\frac N2.
\]

\textbf{(2)}
\(N_L=N/2+x\)、\(N_R=N/2-x\)。\(\ln\Omega=N\ln N-N_L\ln N_L-N_R\ln N_R\)（定数項略）。2次まで展開すると、各
\(\ln(N/2\pm x)\simeq\ln(N/2)\pm\frac{2x}{N}-\frac{2x^2}{N^2}\)
を使って整理し

\[
\ln\Omega\simeq\ln\Omega_{\max}-\frac{2x^2}{N}.
\]

すなわち \(\Omega\propto e^{-2x^2/N}\)。これはガウス型で標準偏差
\(\sigma_x=\sqrt{N}/2\)。

\textbf{(3)} \(\sqrt{\langle x^2\rangle}=\sigma_x=\sqrt N/2\)
ゆえ相対揺らぎ \(\sigma_x/N=1/(2\sqrt N)\sim1/\sqrt N\)。\(N=10^{23}\)
で
\(\sim10^{-11.5}\)。観測誤差よりはるかに小さく、ほぼすべてのミクロ状態が
\(N_L\simeq N/2\)
のマクロ状態に集中する。だからマクロ状態（最大多重度の状態）だけで物理を語れる。

\textbf{答え}：\(N_L=N/2\) で最大、\(\Omega\propto e^{-2x^2/N}\)（幅
\(\sqrt N/2\)）。相対揺らぎ \(\sim1/\sqrt N\sim10^{-11}\)
ゆえマクロ状態で十分。

\subsection*{問題 6-2　エントロピーの加法性と温度
1/T=dS/dE}\label{ex-6-2}
\addcontentsline{toc}{subsection}{問題 6-2　エントロピーの加法性と温度
1/T=dS/dE}

{[}★ 難易度：☆☆ {]}

§6.3 はエントロピー \(S=k_B\ln\Omega\) の加法性と温度の定義
\(1/T=(\partial S/\partial E)_{V,N}\) を与えた。両方を導く。

\begin{enumerate}
\def\labelenumi{\arabic{enumi}.}
\tightlist
\item
  独立な二系 \(\Omega=\Omega_1\Omega_2\) から
  \(S=S_1+S_2\)（加法性）を示せ。
\item
  エネルギー \(E=E_1+E_2\) を分け合う二系を熱接触させ、全多重度
  \(\Omega_1(E_1)\Omega_2(E_2)\)
  を最大化する。\(\partial S_1/\partial E_1=\partial S_2/\partial E_2\)
  が平衡条件であることを示せ。
\item
  この共通の値を \(1/T\)
  と定義する根拠（「温度が等しいとき平衡」という熱力学的描像との整合）を述べよ。
\end{enumerate}

\textbf{関連}：\hyperref[sec-thermal-equilibrium-entropy]{§6.3
エントロピー}、\hyperref[sec-thermal-equilibrium-temperature]{§6.1
温度と熱平衡}。

\section{模範解答（問題
6-2）}\label{ux6a21ux7bc4ux89e3ux7b54ux554fux984c-6-2}

\textbf{(1)} 独立な二系の全多重度は積
\(\Omega=\Omega_1\Omega_2\)。対数をとると
\(S=k_B\ln(\Omega_1\Omega_2)=k_B\ln\Omega_1+k_B\ln\Omega_2=S_1+S_2\)。エントロピーは加法的（示量的）。

\textbf{(2)} \(E_1\) を変数に全エントロピー \(S=S_1(E_1)+S_2(E-E_1)\)
を最大化：

\[
\frac{dS}{dE_1}=\frac{\partial S_1}{\partial E_1}-\frac{\partial S_2}{\partial E_2}=0\ \Rightarrow\ \frac{\partial S_1}{\partial E_1}=\frac{\partial S_2}{\partial E_2}.
\]

（\(dE_2/dE_1=-1\) を使った。）平衡では両系で \(\partial S/\partial E\)
が一致する。

\textbf{(3)}
熱力学では「二系を接触させて正味のエネルギー流が止まる＝温度が等しい」が平衡。(2)
の共通量 \(\partial S/\partial E\)
がまさに平衡で一致する量なので、これを温度と結びつけるのが自然。次元と符号（エネルギーを加えるとエントロピーが増える）から
\(1/T\equiv\partial S/\partial E\)
と定義すると、熱力学の温度と一致する。

\textbf{答え}：\(\Omega=\Omega_1\Omega_2\Rightarrow S=S_1+S_2\)。平衡で
\(\partial S_1/\partial E_1=\partial S_2/\partial E_2\)、これを \(1/T\)
と定義。

\subsection*{問題 6-3　最大エントロピーから Boltzmann
分布を導く}\label{ex-6-3}
\addcontentsline{toc}{subsection}{問題 6-3　最大エントロピーから
Boltzmann 分布を導く}

{[}★ 難易度：☆☆☆ {]}

§6.4 は Lagrange 未定乗数法で \(n_i\propto f(\beta E_i,\alpha)\)
になると述べた。識別可能粒子（古典）について具体的に実行する。準位
\(\{E_i\}\) に \(\{n_i\}\) 個を配る多重度
\(\Omega=\dfrac{N!}{\prod_i n_i!}\)。

\begin{enumerate}
\def\labelenumi{\arabic{enumi}.}
\tightlist
\item
  \(\ln\Omega\) に Stirling を使い、制約
  \(\sum n_i=N\)、\(\sum n_i E_i=E\) のもとで
  \(\dfrac{\partial}{\partial n_i}[\ln\Omega-\alpha\sum n_j-\beta\sum n_j E_j]=0\)
  を解いて \(n_i\propto e^{-\beta E_i}\) を導け。
\item
  規格化して \(n_i=\dfrac{N e^{-\beta E_i}}{\sum_j e^{-\beta E_j}}\)
  を示し、分配関数 \(Z=\sum_j e^{-\beta E_j}\)
  が分母に現れることを確認せよ。
\item
  \(\beta=1/k_BT\) が問題 6-2 の温度定義と整合する理由を一言で。
\end{enumerate}

\textbf{関連}：\hyperref[sec-thermal-equilibrium-max-entropy]{§6.4
最大エントロピー原理}／分配関数の活用は\hyperref[sec-statistical-mechanics-partition]{§7.2}、Bose
版は\hyperref[sec-photon-statistics-bose-einstein]{§8.3}（\hyperref[ex-8-2]{演習
8-2}）。

\section{模範解答（問題
6-3）}\label{ux6a21ux7bc4ux89e3ux7b54ux554fux984c-6-3}

\textbf{(1)}
\(\ln\Omega=\ln N!-\sum_i\ln n_i!\simeq N\ln N-N-\sum_i(n_i\ln n_i-n_i)\)。\(n_i\)
で微分すると \(\partial\ln\Omega/\partial n_i=-\ln n_i\)。変分条件：

\[
-\ln n_i-\alpha-\beta E_i=0\ \Rightarrow\ \ln n_i=-\alpha-\beta E_i\ \Rightarrow\ n_i=e^{-\alpha}e^{-\beta E_i}\propto e^{-\beta E_i}.
\]

\textbf{(2)} \(\sum_i n_i=N\) より
\(e^{-\alpha}\sum_j e^{-\beta E_j}=N\)、よって
\(e^{-\alpha}=N/Z\)、\(Z=\sum_j e^{-\beta E_j}\)。したがって

\[
n_i=\frac{N e^{-\beta E_i}}{\sum_j e^{-\beta E_j}}=\frac{N}{Z}e^{-\beta E_i}.
\]

分母にちょうど分配関数 \(Z\) が現れる。

\textbf{(3)} エネルギー制約の Lagrange 乗数 \(\beta\)
は、エントロピー最大化で「エネルギーを1単位動かしたときのエントロピー変化」に対応する。これは問題
6-2 の \(1/T=\partial S/\partial E\) そのもので、\(\beta=1/k_BT\)
と同定される。

\textbf{答え}：変分で \(n_i\propto e^{-\beta E_i}\)、規格化で
\(n_i=(N/Z)e^{-\beta E_i}\)、\(Z=\sum e^{-\beta E_j}\)。\(\beta=1/k_BT\)
は温度定義と整合。

\subsection*{問題 6-4　詳細釣り合いと遷移率の比}\label{ex-6-4}
\addcontentsline{toc}{subsection}{問題 6-4　詳細釣り合いと遷移率の比}

{[}★ 難易度：☆☆ {]}

§6.6 は詳細釣り合い \(n_i W_{i\to j}=n_j W_{j\to i}\) を導入した。これと
Boltzmann 分布を組み合わせる。

\begin{enumerate}
\def\labelenumi{\arabic{enumi}.}
\tightlist
\item
  熱平衡で \(n_i/n_j=(g_i/g_j)e^{-(E_i-E_j)/k_BT}\) を使い、遷移率の比
  \(W_{i\to j}/W_{j\to i}\) を求めよ。
\item
  「全粒子の正味流れがゼロ」という弱い条件と、「各ペアで個別に釣り合う」詳細釣り合いの違いを、3状態を巡回する流れ（\(1\to2\to3\to1\)）を例に説明せよ。
\item
  この原理が第11章で Einstein
  係数の関係式を導くのにどう使われるか、一言で予告せよ。
\end{enumerate}

\textbf{関連}：\hyperref[sec-thermal-equilibrium-detailed-balance]{§6.6
詳細釣り合い}／Einstein
への応用は\hyperref[sec-einstein-method-detailed-balance]{§11.3}（\hyperref[ex-11-1]{演習
11-1}）、Boltzmann
比は\hyperref[sec-statistical-mechanics-boltzmann]{§7.3}。

\section{模範解答（問題
6-4）}\label{ux6a21ux7bc4ux89e3ux7b54ux554fux984c-6-4}

\textbf{(1)} 詳細釣り合いから
\(W_{i\to j}/W_{j\to i}=n_j/n_i\)。Boltzmann 比を代入：

\[
\frac{W_{i\to j}}{W_{j\to i}}=\frac{n_j}{n_i}=\frac{g_j}{g_i}e^{-(E_j-E_i)/k_BT}.
\]

（上準位ほど占有が少ないので、上向き遷移率と下向き遷移率の比がこの形になる。）

\textbf{(2)} 「正味流れゼロ」は、たとえば \(1\to2\to3\to1\)
と一方向に循環する定常流があっても各状態の占有数が一定なら満たされてしまう。詳細釣り合いはより強く、各ペア（\(1\!\leftrightarrow\!2\),
\(2\!\leftrightarrow\!3\),
\(3\!\leftrightarrow\!1\)）で個別に往復が釣り合うことを要求し、そうした循環流を禁じる。これは微視的時間反転対称性の帰結で、真の熱平衡を特徴づける。

\textbf{(3)} 第11章で Einstein
は、放射と二準位原子の吸収・自発放出・誘導放出に詳細釣り合いを課し、\(J_\nu=B_\nu\)
を要求することで Einstein 関係式 \(A_{ul}/B_{ul}=2h\nu^3/c^2\) と
\(g_lB_{lu}=g_uB_{ul}\)
を導く。本問の「遷移率の比」が原子と放射の係数の比に具体化される。

\textbf{答え}：\(W_{i\to j}/W_{j\to i}=(g_j/g_i)e^{-(E_j-E_i)/k_BT}\)。詳細釣り合いは循環流を禁じる強い条件。第11章の
Einstein 関係式の基礎。

\section*{さらに学ぶための参考文献}\label{sec-thermal-equilibrium-further}
\addcontentsline{toc}{section}{さらに学ぶための参考文献}

\markright{さらに学ぶための参考文献}

\begin{itemize}
\tightlist
\item
  Reif, \emph{Fundamentals of Statistical and Thermal Physics}
  (McGraw-Hill, 1965) --- 熱平衡とエントロピーの古典的扱い
\item
  Landau \& Lifshitz, \emph{Statistical Physics, Part 1} (Pergamon, 3rd
  ed., 1980) --- エントロピー最大原理の体系的扱い
\item
  Jaynes, ``Information Theory and Statistical Mechanics''
  (\emph{Physical Review} 106, 620, 1957) --- 情報理論側からの再構成
\end{itemize}

\chapter{第7章　統計力学の基本}\label{ux7b2c7ux7ae0-ux7d71ux8a08ux529bux5b66ux306eux57faux672c}

\textbf{この章の主題}：第6章で立ち上げた「最大エントロピー状態が熱平衡」という原理を、具体的な計算道具として整える。状態数・分配関数・Boltzmann
分布・揺らぎ ―
これらは古典統計力学の中心的な道具立てであり、第8章で光子の統計（Bose--Einstein）へ拡張するための土台となる。最後に、古典統計を放射場にそのまま適用すると何が起きるかを予告する。これは第IV部・第9章「古典論ではなぜ失敗するのか」への入口でもある。

\section*{この章の中心地図}\label{sec-statistical-mechanics-map}
\addcontentsline{toc}{section}{この章の中心地図}

\markright{この章の中心地図}

\textbf{中心地図：プランク関数 ― 本部で扱う因子}

\[
B_\nu(T) = \frac{2 h \nu^3}{c^2} \cdot \boxed{\frac{1}{e^{h\nu/kT} - 1}}
\]

強調されている \textbf{占有因子} \(1/(e^{h\nu/kT}-1)\)
が、第III部で逆引きする対象である。これは熱平衡にある光子（Bose
粒子）の統計から立ち上がる量で、低振動数極限では \(kT/(h\nu)\)
に、高振動数極限では \(e^{-h\nu/kT}\)
に漸近する。第8章で本格的に導出する。

\begin{tcolorbox}[enhanced jigsaw, arc=.35mm, bottomrule=.15mm, bottomtitle=1mm, breakable, colback=white, colbacktitle=quarto-callout-note-color!10!white, colframe=quarto-callout-note-color-frame, coltitle=black, left=2mm, leftrule=.75mm, opacityback=0, opacitybacktitle=0.6, rightrule=.15mm, title=\textcolor{quarto-callout-note-color}{\faInfo}\hspace{0.5em}{ノート}, titlerule=0mm, toprule=.15mm, toptitle=1mm]

\textbf{方針}：本章で扱う \textbf{Boltzmann 分布}
は、プランク分布の親戚であり、量子効果が無視できる極限では Boltzmann
分布に近づく。中心地図の指数関数因子 \(1/(e^{h\nu/kT}-1)\)
の「指数の中身」がここで準備される。

\end{tcolorbox}

\section*{この章で答える問い}\label{sec-statistical-mechanics-questions}
\addcontentsline{toc}{section}{この章で答える問い}

\markright{この章で答える問い}

\begin{itemize}
\tightlist
\item
  なぜ多数の粒子の集団は、熱力学量（温度・圧力）のような少数の変数で記述できてしまうのか
\item
  分配関数を計算すると、なぜ観測されるスペクトル形が出てくるのか
\item
  古典統計を放射場にそのまま適用すると、何が起こるのか（→
  第9章への伏線）
\item
  揺らぎは観測スペクトルにどんな形で現れるのか
\end{itemize}

\section*{到達目標}\label{sec-statistical-mechanics-goals}
\addcontentsline{toc}{section}{到達目標}

\markright{到達目標}

この章を読み終えると、読者は次のことができるようになる：

\begin{itemize}
\tightlist
\item
  状態数・分配関数・Boltzmann 分布の論理を整理して述べられる
\item
  平均値と揺らぎが熱力学量とどう結びつくかを説明できる
\item
  古典統計を放射場に適用したときに何が起こるかを予告的に把握できる
\end{itemize}

\begin{center}\rule{0.5\linewidth}{0.5pt}\end{center}

\section{7.1 状態数と確率 ―
等確率の原理}\label{sec-statistical-mechanics-states}

{[}本文目安：B3{]}

統計力学の出発点は、ミクロ状態に対する\textbf{等確率の仮定}である：

\begin{quote}
孤立系のミクロ状態のうち、与えられたエネルギー \(E\)
を満たすものは、すべて等しい確率で実現する。
\end{quote}

これだけで多くの結果が出る。粒子数 \(N\)、体積 \(V\)、エネルギー \(E\)
が指定されたマクロ状態に対応するミクロ状態の数 \(\Omega(E, V, N)\)
を計算すれば、エントロピー
\(S = k_B \ln \Omega\)（式~\ref{eq-entropy-boltzmann}）からすべての熱力学量が導かれる。

これを \textbf{ミクロカノニカル分布} と呼ぶ。

しかし実用上はもう一段使いやすい設定が便利である。系を\textbf{熱浴}（温度
\(T\) で固定）と接触させた状況を考える。エネルギー \(E_i\)
のミクロ状態が実現する確率は

\begin{equation}\protect\phantomsection\label{eq-boltzmann-prob}{
P_i = \frac{1}{Z}\, e^{-E_i / k_B T}
}\end{equation}

となる。これが \textbf{Boltzmann 分布}（\textbf{Boltzmann distribution},
あるいは canonical distribution）で、\(Z\) は規格化定数：

\begin{equation}\protect\phantomsection\label{eq-partition-function}{
Z = \sum_i e^{-E_i / k_B T}
}\end{equation}

これが \textbf{分配関数}（\textbf{partition function}, あるいは独語
\emph{Zustandssumme}「状態和」）。\(Z\)
は系のすべての情報を含む量で、ここから物理量がすべて引き出せる。

\begin{tcolorbox}[enhanced jigsaw, arc=.35mm, bottomrule=.15mm, bottomtitle=1mm, breakable, colback=white, colbacktitle=quarto-callout-note-color!10!white, colframe=quarto-callout-note-color-frame, coltitle=black, left=2mm, leftrule=.75mm, opacityback=0, opacitybacktitle=0.6, rightrule=.15mm, title=\textcolor{quarto-callout-note-color}{\faInfo}\hspace{0.5em}{ノート}, titlerule=0mm, toprule=.15mm, toptitle=1mm]

\textbf{用語}：「分配関数」は英語 partition function
の和訳だが、原語ともに物理的意味が直接には汲み取りにくい。ドイツ語の
Zustandssumme（\textbf{状態和}）の方が明快で、エネルギー重み
\(e^{-E_i/k_B T}\)
を全状態にわたって足し上げた量、という意味がそのまま読み取れる。

\end{tcolorbox}

\section{7.2 分配関数 ―
すべての熱力学量を生成する母関数}\label{sec-statistical-mechanics-partition}

{[}本文目安：B3{]}

分配関数 \(Z(T, V, N)\) から、平均エネルギーは

\begin{equation}\protect\phantomsection\label{eq-mean-energy}{
\langle E \rangle = -\frac{\partial \ln Z}{\partial \beta}
}\end{equation}

ただし \(\beta = 1/k_B T\)。同様に他の熱力学量も：

{\def\LTcaptype{none} % do not increment counter
\begin{longtable}[]{@{}
  >{\raggedright\arraybackslash}p{(\linewidth - 2\tabcolsep) * \real{0.5000}}
  >{\raggedright\arraybackslash}p{(\linewidth - 2\tabcolsep) * \real{0.5000}}@{}}
\toprule\noalign{}
\begin{minipage}[b]{\linewidth}\raggedright
物理量
\end{minipage} & \begin{minipage}[b]{\linewidth}\raggedright
分配関数からの計算式
\end{minipage} \\
\midrule\noalign{}
\endhead
\bottomrule\noalign{}
\endlastfoot
平均エネルギー \(\langle E \rangle\) &
\(-\partial \ln Z / \partial \beta\) \\
Helmholtz 自由エネルギー \(F\) & \(-k_B T \ln Z\) \\
エントロピー \(S\) & \(-\partial F / \partial T\) \\
圧力 \(P\) & \(-\partial F / \partial V\) \\
比熱 \(C_V\) & \(\partial \langle E \rangle / \partial T\) \\
\end{longtable}
}

つまり、\textbf{分配関数を計算できれば熱力学はほぼ完結する}。逆に言えば、実際の計算の核心はミクロ状態の和
式~\ref{eq-partition-function} を実行する点にある。

\textbf{例：N 個の独立な調和振動子}

エネルギー \(\epsilon_n = (n + 1/2)\hbar\omega\) を持つ独立な振動子
\(N\) 個の系では、

\begin{equation}\protect\phantomsection\label{eq-partition-oscillator}{
Z = \prod_{i=1}^N Z_1, \quad Z_1 = \sum_{n=0}^\infty e^{-\beta(n + 1/2)\hbar\omega} = \frac{e^{-\beta\hbar\omega/2}}{1 - e^{-\beta\hbar\omega}}
}\end{equation}

ここから平均エネルギー：

\begin{equation}\protect\phantomsection\label{eq-mean-energy-oscillator}{
\langle E \rangle = N \left( \frac{\hbar\omega}{2} + \frac{\hbar\omega}{e^{\beta\hbar\omega} - 1} \right)
}\end{equation}

\begin{figure}[H]

\centering{

\includegraphics[width=0.88\linewidth,height=\textheight,keepaspectratio]{index_files/mediabag/part3/../images/quantum-freeze-out.pdf}

}

\caption{\label{fig-quantum-freezeout}量子化振動子 1 個の平均エネルギー
\(\langle\epsilon\rangle\) と古典の等分配 \(k_B T\)
の比較。低振動数（\(h\nu \ll k_BT\)）では古典と一致するが、\textbf{\(h\nu \approx k_BT\)
付近から急速に等分配を下回り}、\(h\nu \gg k_BT\)
では指数的に「凍結」する。これが Rayleigh--Jeans
紫外破綻からプランク分布への移行を量子論側から見たもので、第10章 §10.2
で本格的に扱う。}

\end{figure}%

第二項に \textbf{\(1/(e^{\hbar\omega/k_B T} - 1)\)}
という形が現れていることに注意。これがプランク分布の占有因子に他ならない。プランクが量子仮説でこの式に辿り着いたのは、第10章で本格的に扱う。

\begin{tcolorbox}[enhanced jigsaw, arc=.35mm, bottomrule=.15mm, bottomtitle=1mm, breakable, colback=white, colbacktitle=quarto-callout-note-color!10!white, colframe=quarto-callout-note-color-frame, coltitle=black, left=2mm, leftrule=.75mm, opacityback=0, opacitybacktitle=0.6, rightrule=.15mm, title=\textcolor{quarto-callout-note-color}{\faInfo}\hspace{0.5em}{ノート}, titlerule=0mm, toprule=.15mm, toptitle=1mm]

\textbf{零点エネルギー \(\hbar\omega/2\) の扱いについて}：本章では
\(\epsilon_n = (n+1/2)\hbar\omega\) と書いて零点項を残したが、Planck
の原論文（1900）は \(\epsilon_n = n h\nu\)
と書いた（第10章でこの形を再現する）。零点項の有無は
\textbf{観測量である占有因子 \(1/(e^{\hbar\omega/k_BT}-1)\)
には影響しない}（零点エネルギーは温度によらない定数で、温度微分・分配関数の対数微分で消える）。両者の物理的意味（古典的振動子の量子化
vs
エネルギー量子の数え上げ）と、零点項が観測可能になる場面（真空揺らぎ、Casimir
効果、自発放出、§25.3）は第IV部と第VIII部で扱う

\end{tcolorbox}

\begin{tcolorbox}[enhanced jigsaw, arc=.35mm, bottomrule=.15mm, bottomtitle=1mm, breakable, colback=white, colbacktitle=quarto-callout-warning-color!10!white, colframe=quarto-callout-warning-color-frame, coltitle=black, left=2mm, leftrule=.75mm, opacityback=0, opacitybacktitle=0.6, rightrule=.15mm, title=\textcolor{quarto-callout-warning-color}{\faExclamationTriangle}\hspace{0.5em}{警告}, titlerule=0mm, toprule=.15mm, toptitle=1mm]

\textbf{先取りについて}：本節では「エネルギーが \(\hbar\omega\)
の整数倍に
\textbf{量子化されている}」という前提を、第IV部（第9〜10章）の逆引きより
\textbf{先に}
使っている。これは「観測されるプランク分布の指数因子を、最も簡潔に再現できる仮定として一度受け入れる」立場であり、\textbf{第8章のモード密度の借用}（§8.5）と同じ精神の先取りである。なぜ量子化されなくてはならないか（古典統計の破綻
→ Rayleigh-Jeans の紫外発散 → Planck
仮説）は第9-10章で逆引きされる。第7章で先取りした分は、第10章で必ず回収される

\end{tcolorbox}

\begin{tcolorbox}[enhanced jigsaw, arc=.35mm, bottomrule=.15mm, breakable, colback=white, colframe=quarto-callout-tip-color-frame, left=2mm, leftrule=.75mm, opacityback=0, rightrule=.15mm, toprule=.15mm]
\begin{minipage}[t]{5.5mm}
\textcolor{quarto-callout-tip-color}{\faLightbulb}
\end{minipage}%
\begin{minipage}[t]{\textwidth - 5.5mm}

\textbf{問い}：いま導いた振動子の式 式~\ref{eq-mean-energy-oscillator}
は、プランク分布そのものではないのか？

\textbf{短答}：構造はほぼ同じだが、まだプランク分布ではない。式~\ref{eq-mean-energy-oscillator}
は「ある \textbf{一つの}
モードの平均エネルギー」を与える式である。プランク分布 \(B_\nu(T)\)
は、これに「振動数 \(\nu\)
近傍にどれだけのモードがあるか」（モード密度）を掛けたもの。

\textbf{もう一歩}：モード密度の話は第V部・第14章で扱う。中心地図のプランク関数の
\textbf{指数関数因子}（占有因子）が本章で説明され、\textbf{前因子}
\(2h\nu^3/c^2\) は第V部で説明される、という構造になっている。

\end{minipage}%
\end{tcolorbox}

\begin{tcolorbox}[enhanced jigsaw, arc=.35mm, bottomrule=.15mm, bottomtitle=1mm, breakable, colback=white, colbacktitle=quarto-callout-note-color!10!white, colframe=quarto-callout-note-color-frame, coltitle=black, left=2mm, leftrule=.75mm, opacityback=0, opacitybacktitle=0.6, rightrule=.15mm, title=\textcolor{quarto-callout-note-color}{\faInfo}\hspace{0.5em}{ノート}, titlerule=0mm, toprule=.15mm, toptitle=1mm]

\textbf{使えるようになった道具}（§7.1〜7.2）：

\begin{itemize}
\tightlist
\item
  分配関数 \(Z\)
  から、平均エネルギー・自由エネルギー・エントロピー・圧力・比熱までを微分一発で取り出せる
\item
  量子化された調和振動子の平均エネルギー
  式~\ref{eq-mean-energy-oscillator}
  に、観測されるプランク分布の指数関数因子
  \(1/(e^{\hbar\omega/k_BT}-1)\) がすでに現れる ―
  観測される高振動数の指数減衰の起源がここに芽生えている
\end{itemize}

\end{tcolorbox}

\section{7.3 Boltzmann 分布 ―
古典統計の中心定理}\label{sec-statistical-mechanics-boltzmann}

{[}本文目安：B3{]}

式~\ref{eq-boltzmann-prob} を観測量で表すと、\(N\)
個の同種粒子のうち、エネルギー \(E_i\) の状態 \(i\) にある粒子数の平均は

\begin{equation}\protect\phantomsection\label{eq-boltzmann-distribution}{
\langle n_i \rangle = N \cdot \frac{g_i e^{-E_i/k_B T}}{\sum_j g_j e^{-E_j/k_B T}}
}\end{equation}

ここで \(g_i\) は状態 \(i\)
の\textbf{縮退度}（\textbf{degeneracy}）。これが（古典）\textbf{Boltzmann
分布} である。

\textbf{例：水素原子の準位占有数}

水素原子の主量子数 \(n\) の準位は、エネルギー
\(E_n = -13.6 \text{ eV} / n^2\) をもち、縮退度は
\(g_n = 2 n^2\)（電子スピン込み）。温度 \(T\) の熱浴と平衡にあるとき、

\begin{equation}\protect\phantomsection\label{eq-boltzmann-hydrogen}{
\frac{n_2}{n_1} = \frac{g_2}{g_1} e^{-(E_2 - E_1)/k_B T}
}\end{equation}

この比は、第VII部で線スペクトルの強度比から温度を求める基本式になる。

\begin{tcolorbox}[enhanced jigsaw, arc=.35mm, bottomrule=.15mm, bottomtitle=1mm, breakable, colback=white, colbacktitle=quarto-callout-note-color!10!white, colframe=quarto-callout-note-color-frame, coltitle=black, left=2mm, leftrule=.75mm, opacityback=0, opacitybacktitle=0.6, rightrule=.15mm, title=\textcolor{quarto-callout-note-color}{\faInfo}\hspace{0.5em}{ノート}, titlerule=0mm, toprule=.15mm, toptitle=1mm]

\textbf{観測との対応}：観測される線スペクトル強度比から温度を読むのは、まさにこの
Boltzmann 分布が成り立っているという仮定（LTE）に基づいている。LTE
が破れる場面（non-LTE、第5章 §5.5）では 式~\ref{eq-boltzmann-hydrogen}
は使えない。

\end{tcolorbox}

\begin{tcolorbox}[enhanced jigsaw, arc=.35mm, bottomrule=.15mm, bottomtitle=1mm, breakable, colback=white, colbacktitle=quarto-callout-note-color!10!white, colframe=quarto-callout-note-color-frame, coltitle=black, left=2mm, leftrule=.75mm, opacityback=0, opacitybacktitle=0.6, rightrule=.15mm, title=\textcolor{quarto-callout-note-color}{\faInfo}\hspace{0.5em}{ノート}, titlerule=0mm, toprule=.15mm, toptitle=1mm]

\textbf{使えるようになった道具}（§7.3）：Boltzmann 分布
式~\ref{eq-boltzmann-distribution}
は、観測される複数の線の強度比から温度を読み出す「\textbf{励起温度}」推定の基礎式である。第VII部・第24章で本格的に応用する。

\end{tcolorbox}

\section{7.4
平均値と揺らぎ}\label{sec-statistical-mechanics-fluctuations}

{[}本文目安：B3{]}

統計力学は平均値だけでなく、\textbf{揺らぎ}（\textbf{fluctuations},
統計的なばらつき）も予言する。揺らぎは温度・密度・エネルギー・光子数などのあらゆる量に現れる。

エネルギー揺らぎは分配関数からこう求まる：

\begin{equation}\protect\phantomsection\label{eq-energy-fluctuation}{
\langle (\Delta E)^2 \rangle = \langle E^2 \rangle - \langle E \rangle^2 = k_B T^2 \frac{\partial \langle E \rangle}{\partial T} = k_B T^2 C_V
}\end{equation}

\(C_V\) は定積熱容量。\(N \sim 10^{23}\) では、相対的な揺らぎは
\(\Delta E / \langle E \rangle \sim 1/\sqrt{N} \sim 10^{-11}\)
で、観測に出てこない（マクロ熱力学が成り立つ理由）。

しかし\textbf{少数粒子系}や\textbf{特定の振動数帯}では揺らぎが重要になる。

\textbf{CMB のスペクトル歪み}：CMB
の黒体スペクトルは完璧だが、初期宇宙の様々な過程（粒子崩壊、再結合、Compton
散乱の不完全さ）がスペクトルに微小な歪みを残しうる。これらの歪みは「光子数の揺らぎ」「エネルギーの揺らぎ」を反映する量であり、第16章で扱う。

\begin{tcolorbox}[enhanced jigsaw, arc=.35mm, bottomrule=.15mm, breakable, colback=white, colframe=quarto-callout-tip-color-frame, left=2mm, leftrule=.75mm, opacityback=0, rightrule=.15mm, toprule=.15mm]
\begin{minipage}[t]{5.5mm}
\textcolor{quarto-callout-tip-color}{\faLightbulb}
\end{minipage}%
\begin{minipage}[t]{\textwidth - 5.5mm}

\textbf{問い}：揺らぎが小さいのは、なぜ「自然法則」のように見えるのか？

\textbf{短答}：\(N\) が大きいから。揺らぎは \(1/\sqrt{N}\)
で減衰するので、アボガドロ数規模の系では事実上ゼロに見える。

\textbf{もう一歩}：宇宙論的スケールの観測では話が違う。空間的に異なる場所の
CMB 温度ゆらぎは \(\Delta T / T \sim 10^{-5}\)
で観測されているが、これは「熱平衡の揺らぎ」ではなく、初期宇宙の量子ゆらぎが膨張で増幅された結果である（第16章）。両者は混同しないこと。

\end{minipage}%
\end{tcolorbox}

\section{7.5
熱力学関数とのつながり}\label{sec-statistical-mechanics-thermodynamics}

{[}本文目安：B3-B4｜導出：M1{]}

ここまで導入した分配関数 \(Z\)
は、熱力学の標準的関数とすべて対応している。鍵となる関係：

\begin{equation}\protect\phantomsection\label{eq-free-energy}{
F = -k_B T \ln Z
}\end{equation}

ここから

\[
S = -\frac{\partial F}{\partial T}, \quad P = -\frac{\partial F}{\partial V}, \quad \mu = \frac{\partial F}{\partial N}
\]

\textbf{化学ポテンシャル}（\textbf{chemical potential}）\(\mu\)
は粒子数に対する Lagrange 乗数で、式~\ref{eq-lagrange} の
\(\alpha = -\mu/k_B T\) に対応する。

\textbf{化学ポテンシャルの物理的意味}：粒子を一つ加えたときの自由エネルギーの変化。「粒子数を保存しないといけない系」では化学ポテンシャルが現れ、保存しない系（光子のように）では
\(\mu = 0\) になる。

なぜ光子は粒子数を保存しないか。電磁場のモードと物質の相互作用（吸収・放出）で光子は生成・消滅し、その数は決まっていない。エネルギーは保存するが個数は保存しない。これが\textbf{光子の化学ポテンシャルがゼロ}になる物理的理由で、第8章で詳しく扱う。

\section{7.6 古典統計の限界 ―
放射場での破綻}\label{sec-statistical-mechanics-classical-limit}

{[}本文目安：B3{]}

古典統計力学を放射場にそのまま適用すると、本書が逆引きで辿る\textbf{最初の決定的な失敗}に遭遇する。

\textbf{古典の結論（Rayleigh--Jeans）}：振動数 \(\nu\)
あたりのモード数を \(g(\nu)\)
とすると、エネルギー等分配定理から各モードには平均エネルギー \(k_B T\)
が配分される。よって振動数 \(\nu\) あたりの放射エネルギー密度は

\begin{equation}\protect\phantomsection\label{eq-rj-classical}{
u_\nu^\text{RJ} = g(\nu) \cdot k_B T = \frac{8\pi \nu^2}{c^3} k_B T
}\end{equation}

これは振動数の二乗に比例して増加し、全エネルギー \(\int u_\nu d\nu\)
は発散する。

\textbf{観測との矛盾}：実際の黒体スペクトルは、高振動数で指数関数的に減衰し、全エネルギー
\(u = aT^4\) は有限である。古典統計は決定的に間違っている。

\textbf{何が問題か}：

\begin{itemize}
\tightlist
\item
  古典統計の\textbf{エネルギー等分配定理}（\textbf{equipartition
  theorem}, 各自由度に
  \(k_B T\)）は、エネルギーが\textbf{連続的に}やりとりされる前提に立つ
\item
  高振動数のモード（\(h\nu \gg k_B T\)）では、エネルギー \(h\nu\)
  という大きな単位でしかエネルギーが分配されえない →
  等分配が成り立たない
\end{itemize}

これを正確に取り扱うには、\textbf{量子的な性質}（エネルギー量子化＋粒子の区別不能性）が必要である。第8章で「光子の統計」として導入し、これがプランク分布の占有因子
\(1/(e^{h\nu/kT}-1)\) を生む。

\begin{tcolorbox}[enhanced jigsaw, arc=.35mm, bottomrule=.15mm, bottomtitle=1mm, breakable, colback=white, colbacktitle=quarto-callout-note-color!10!white, colframe=quarto-callout-note-color-frame, coltitle=black, left=2mm, leftrule=.75mm, opacityback=0, opacitybacktitle=0.6, rightrule=.15mm, title=\textcolor{quarto-callout-note-color}{\faInfo}\hspace{0.5em}{ノート}, titlerule=0mm, toprule=.15mm, toptitle=1mm]

\textbf{注意（理論）}：古典統計の破綻は「\textbf{紫外線破綻}」（\textbf{ultraviolet
catastrophe}）と呼ばれる。これは決定的な歴史的契機で、量子論誕生（Planck
1900, Einstein 1905）に直結した。第9章で本格的に扱う。

\end{tcolorbox}

\begin{tcolorbox}[enhanced jigsaw, arc=.35mm, bottomrule=.15mm, breakable, colback=white, colframe=quarto-callout-tip-color-frame, left=2mm, leftrule=.75mm, opacityback=0, rightrule=.15mm, toprule=.15mm]
\begin{minipage}[t]{5.5mm}
\textcolor{quarto-callout-tip-color}{\faLightbulb}
\end{minipage}%
\begin{minipage}[t]{\textwidth - 5.5mm}

\textbf{問い}：Rayleigh--Jeans 則は完全に間違いなのか？

\textbf{短答}：いいえ、低振動数極限（\(h\nu \ll k_B T\)）では正しい。電波天文学ではほぼ常にこの極限で観測しており、Rayleigh--Jeans
は実用的に最も使われる近似である（第2章 §2.5）。

\textbf{もう一歩}：プランク分布の物理を学ぶときは、「古典で完全に間違っている」ではなく「\textbf{低振動数では古典が正しく、高振動数で量子効果が効く}」というニュアンスで理解するのが正確である。これが第IV部全体の論点になる。

\end{minipage}%
\end{tcolorbox}

\begin{tcolorbox}[enhanced jigsaw, arc=.35mm, bottomrule=.15mm, bottomtitle=1mm, breakable, colback=white, colbacktitle=quarto-callout-note-color!10!white, colframe=quarto-callout-note-color-frame, coltitle=black, left=2mm, leftrule=.75mm, opacityback=0, opacitybacktitle=0.6, rightrule=.15mm, title=\textcolor{quarto-callout-note-color}{\faInfo}\hspace{0.5em}{ノート}, titlerule=0mm, toprule=.15mm, toptitle=1mm]

\textbf{使えるようになった道具}（§7.4〜7.6）：

\begin{itemize}
\tightlist
\item
  揺らぎ \(\Delta E / \langle E \rangle \sim 1/\sqrt{N}\)
  から、観測上いつ熱力学的決定論が信用できるかを判断できる
\item
  古典統計の枠組みでは観測される Rayleigh--Jeans
  形（低振動数）は説明できる。観測される高振動数の指数減衰は古典では出ない
  ― これが量子論を要請する観測上の根拠
\end{itemize}

\end{tcolorbox}

\begin{center}\rule{0.5\linewidth}{0.5pt}\end{center}

\section*{この章で何がわかったか}\label{sec-statistical-mechanics-summary}
\addcontentsline{toc}{section}{この章で何がわかったか}

\markright{この章で何がわかったか}

\textbf{中心地図に戻る}

本章で、古典統計力学の主要道具 ― 分配関数 \(Z\)、Boltzmann 分布
\(\langle n_i \rangle \propto g_i e^{-E_i/k_B T}\)、揺らぎ、熱力学関数 ―
が整った。

中心地図 \(B_\nu(T)\) の
\textbf{「式が成り立つ条件」}（第6章）に続き、本章では
\textbf{「指数の中身
\(E/k_B T\)」が温度の正体である}ことが言語化された。

しかし古典統計をそのまま放射場に適用しようとすると、Rayleigh--Jeans
則の発散という決定的な破綻に出会う。\textbf{指数関数因子
\(1/(e^{h\nu/kT}-1)\) が出てこない}のである。

\textbf{次章へ}：第8章で「光子の統計」 ― 区別できない Bose
粒子の統計力学 ―
を扱う。ここで初めて、プランク分布の指数関数因子（占有因子）が必然として現れる。光子の化学ポテンシャルがなぜゼロかも、ここで明らかになる。これが第III部全体のクライマックスである。

\section*{演習問題}\label{sec-statistical-mechanics-exercises}
\addcontentsline{toc}{section}{演習問題}

\markright{演習問題}

\textbf{模範解答}
は折りたたみを展開して確認できる（オンライン版）。まず自力で解いてから開くこと。

\subsection*{問題 7-1　熱浴との接触から Boltzmann
確率を導く}\label{ex-7-1}
\addcontentsline{toc}{subsection}{問題 7-1　熱浴との接触から Boltzmann
確率を導く}

{[}★ 難易度：☆☆ {]}

§7.1 は \(P_i=e^{-E_i/k_BT}/Z\) を結果として与えた。系を温度 \(T\)
の大きな熱浴と接触させた設定から導く。

\begin{enumerate}
\def\labelenumi{\arabic{enumi}.}
\tightlist
\item
  系がエネルギー \(E_i\) の状態にある確率は、熱浴が残りのエネルギー
  \(E_\mathrm{tot}-E_i\)
  を取る多重度に比例する：\(P_i\propto\Omega_\mathrm{bath}(E_\mathrm{tot}-E_i)\)。これを
  \(S_\mathrm{bath}=k_B\ln\Omega_\mathrm{bath}\) で書き換えよ。
\item
  \(E_i\ll E_\mathrm{tot}\) として
  \(S_\mathrm{bath}(E_\mathrm{tot}-E_i)\) を \(E_i\)
  について1次まで展開し、\(1/T=\partial S_\mathrm{bath}/\partial E\)
  を使って \(P_i\propto e^{-E_i/k_BT}\) を導け。
\item
  規格化して \(Z=\sum_i e^{-E_i/k_BT}\) が現れることを示せ。
\end{enumerate}

\textbf{関連}：\hyperref[sec-statistical-mechanics-states]{§7.1
状態数と確率}／温度の定義は\hyperref[sec-thermal-equilibrium-entropy]{§6.3}（\hyperref[ex-6-2]{演習
6-2}）。

\section{模範解答（問題
7-1）}\label{ux6a21ux7bc4ux89e3ux7b54ux554fux984c-7-1}

\textbf{(1)} 系が状態 \(i\)（エネルギー
\(E_i\)）にあるとき、全体のミクロ状態数は熱浴の多重度
\(\Omega_\mathrm{bath}(E_\mathrm{tot}-E_i)\) に等しい。等確率の原理から
\(P_i\propto\Omega_\mathrm{bath}(E_\mathrm{tot}-E_i)=e^{S_\mathrm{bath}(E_\mathrm{tot}-E_i)/k_B}\)。

\textbf{(2)} \(E_i\ll E_\mathrm{tot}\) で Taylor 展開：

\[
S_\mathrm{bath}(E_\mathrm{tot}-E_i)\simeq S_\mathrm{bath}(E_\mathrm{tot})-E_i\frac{\partial S_\mathrm{bath}}{\partial E}=S_\mathrm{bath}(E_\mathrm{tot})-\frac{E_i}{T}.
\]

よって
\(P_i\propto e^{S_\mathrm{bath}(E_\mathrm{tot})/k_B}e^{-E_i/k_BT}\propto e^{-E_i/k_BT}\)。

\textbf{(3)} \(\sum_i P_i=1\) より
\(P_i=\dfrac{e^{-E_i/k_BT}}{\sum_j e^{-E_j/k_BT}}=\dfrac{e^{-E_i/k_BT}}{Z}\)、\(Z=\sum_j e^{-E_j/k_BT}\)。

\textbf{答え}：\(P_i\propto\Omega_\mathrm{bath}(E_\mathrm{tot}-E_i)\)
を1次展開し \(1/T=\partial S/\partial E\) を使うと
\(P_i=e^{-E_i/k_BT}/Z\)。

\subsection*{問題
7-2　振動子の分配関数と、零点項が占有因子に効かないこと}\label{ex-7-2}
\addcontentsline{toc}{subsection}{問題
7-2　振動子の分配関数と、零点項が占有因子に効かないこと}

{[}★ 難易度：☆☆ {]}

§7.2 は調和振動子の \(Z_1\) と \(\langle E\rangle\)
を与え、零点項が観測量に影響しないと注記した。これを確かめる。準位
\(\epsilon_n=(n+\tfrac12)\hbar\omega\)。

\begin{enumerate}
\def\labelenumi{\arabic{enumi}.}
\tightlist
\item
  \(Z_1=\displaystyle\sum_{n=0}^\infty e^{-\beta(n+1/2)\hbar\omega}=\dfrac{e^{-\beta\hbar\omega/2}}{1-e^{-\beta\hbar\omega}}\)
  を示せ。
\item
  \(\langle\epsilon\rangle=-\partial\ln Z_1/\partial\beta=\dfrac{\hbar\omega}{2}+\dfrac{\hbar\omega}{e^{\beta\hbar\omega}-1}\)
  を導け。
\item
  零点項 \(\hbar\omega/2\) は (i) 占有因子
  \(1/(e^{\beta\hbar\omega}-1)\) と (ii) 比熱
  \(C=\partial\langle\epsilon\rangle/\partial T\)
  のどちらにも影響しないことを示せ。なぜ Planck の \(\epsilon_n=nh\nu\)
  と Schrödinger の \(\epsilon_n=(n+\tfrac12)\hbar\omega\)
  が同じスペクトルを与えるのか。
\end{enumerate}

\textbf{関連}：\hyperref[sec-statistical-mechanics-partition]{§7.2
分配関数}／Planck
の量子仮説からの導出は\hyperref[sec-planck-quantum-oscillator]{§10.2}（\hyperref[ex-10-1]{演習
10-1}）、零点エネルギーは\hyperref[sec-quantum-foundations-oscillator]{§12.2}。

\section{模範解答（問題
7-2）}\label{ux6a21ux7bc4ux89e3ux7b54ux554fux984c-7-2}

\textbf{(1)} 共通因子 \(e^{-\beta\hbar\omega/2}\) を括り出すと等比級数：

\[
Z_1=e^{-\beta\hbar\omega/2}\sum_{n=0}^\infty e^{-\beta n\hbar\omega}=e^{-\beta\hbar\omega/2}\cdot\frac{1}{1-e^{-\beta\hbar\omega}}.
\]

\textbf{(2)}
\(\ln Z_1=-\dfrac{\beta\hbar\omega}{2}-\ln(1-e^{-\beta\hbar\omega})\)。\(\beta\)
で微分：

\[
\langle\epsilon\rangle=-\frac{\partial\ln Z_1}{\partial\beta}=\frac{\hbar\omega}{2}+\frac{\hbar\omega e^{-\beta\hbar\omega}}{1-e^{-\beta\hbar\omega}}=\frac{\hbar\omega}{2}+\frac{\hbar\omega}{e^{\beta\hbar\omega}-1}.
\]

\textbf{(3)} (i) 第二項 \(1/(e^{\beta\hbar\omega}-1)\)
が占有因子で、零点項 \(\hbar\omega/2\) とは独立に現れる。(ii)
\(\hbar\omega/2\) は温度に依らない定数なので
\(C=\partial\langle\epsilon\rangle/\partial T\)
で微分すると消える。よって観測量（占有因子・比熱・放射スペクトル）は零点項に依らない。Planck（\(nh\nu\)）と
Schrödinger（\((n+\tfrac12)\hbar\omega\)）は零点という定数オフセットだけが違い、温度微分・対数微分で消えるので、同じ占有因子＝同じスペクトルを与える。

\textbf{答え}：\(Z_1=e^{-\beta\hbar\omega/2}/(1-e^{-\beta\hbar\omega})\)、\(\langle\epsilon\rangle=\tfrac12\hbar\omega+\hbar\omega/(e^{\beta\hbar\omega}-1)\)。零点項は定数ゆえ占有因子・比熱に効かず、\(nh\nu\)
と \((n+\tfrac12)\hbar\omega\) は同じスペクトル。

\subsection*{問題 7-3　エネルギー揺らぎ (ΔE)²=kBT²Cv}\label{ex-7-3}
\addcontentsline{toc}{subsection}{問題 7-3　エネルギー揺らぎ
(ΔE)²=kBT²Cv}

{[}★ 難易度：☆☆ {]}

§7.4
は揺らぎの式を結果として与えた。分配関数から導く。\(\beta=1/k_BT\)。

\begin{enumerate}
\def\labelenumi{\arabic{enumi}.}
\tightlist
\item
  \(\langle E\rangle=-\dfrac{\partial\ln Z}{\partial\beta}\)、\(\langle E^2\rangle=\dfrac{1}{Z}\dfrac{\partial^2 Z}{\partial\beta^2}\)
  を使い、\(\langle(\Delta E)^2\rangle=\dfrac{\partial^2\ln Z}{\partial\beta^2}=-\dfrac{\partial\langle E\rangle}{\partial\beta}\)
  を示せ。
\item
  \(\dfrac{\partial}{\partial\beta}=-k_BT^2\dfrac{\partial}{\partial T}\)
  を使い、\(\langle(\Delta E)^2\rangle=k_BT^2 C_V\)
  を導け（\(C_V=\partial\langle E\rangle/\partial T\)）。
\item
  \(\langle E\rangle\propto N\)、\(C_V\propto N\) から相対揺らぎ
  \(\sqrt{\langle(\Delta E)^2\rangle}/\langle E\rangle\propto1/\sqrt N\)
  を示せ。
\end{enumerate}

\textbf{関連}：\hyperref[sec-statistical-mechanics-fluctuations]{§7.4
平均値と揺らぎ}／多重度の揺らぎ幅は\hyperref[ex-6-1]{演習 6-1}、CMB
スペクトル歪みは\href{../part6/16-cmb.qmd}{第16章}。

\section{模範解答（問題
7-3）}\label{ux6a21ux7bc4ux89e3ux7b54ux554fux984c-7-3}

\textbf{(1)} \(\langle E^2\rangle-\langle E\rangle^2\)
を計算する。\(\langle E\rangle=-\partial_\beta\ln Z=-\dfrac{1}{Z}\partial_\beta Z\)。もう一度微分：

\[
\frac{\partial^2\ln Z}{\partial\beta^2}=\frac{1}{Z}\frac{\partial^2 Z}{\partial\beta^2}-\left(\frac{1}{Z}\frac{\partial Z}{\partial\beta}\right)^2=\langle E^2\rangle-\langle E\rangle^2=\langle(\Delta E)^2\rangle.
\]

また
\(\partial^2\ln Z/\partial\beta^2=-\partial\langle E\rangle/\partial\beta\)。

\textbf{(2)} \(\partial_\beta=-k_BT^2\partial_T\)（\(\beta=1/k_BT\) より
\(d\beta=-dT/k_BT^2\)）。よって

\[
\langle(\Delta E)^2\rangle=-\frac{\partial\langle E\rangle}{\partial\beta}=k_BT^2\frac{\partial\langle E\rangle}{\partial T}=k_BT^2 C_V.
\]

\textbf{(3)} \(\langle E\rangle\propto N\)、\(C_V\propto N\) なので
\(\langle(\Delta E)^2\rangle\propto N\)、\(\sqrt{\langle(\Delta E)^2\rangle}\propto\sqrt N\)。相対揺らぎは
\(\sqrt N/N=1/\sqrt N\)。マクロ系で揺らぎが消える理由。

\textbf{答え}：\(\langle(\Delta E)^2\rangle=\partial^2\ln Z/\partial\beta^2=k_BT^2C_V\)。相対揺らぎ
\(\propto1/\sqrt N\)。

\subsection*{問題 7-4　水素準位の Boltzmann 占有と Balmer
線の温度依存}\label{ex-7-4}
\addcontentsline{toc}{subsection}{問題 7-4　水素準位の Boltzmann 占有と
Balmer 線の温度依存}

{[}★ 難易度：☆☆☆ {]}

§7.3 の Boltzmann 分布 \(n_2/n_1=(g_2/g_1)e^{-(E_2-E_1)/k_BT}\)
を水素に応用する。\(E_n=-13.6/n^2\)
eV、\(g_n=2n^2\)、\(k_B=8.617\times10^{-5}\) eV/K。

\begin{enumerate}
\def\labelenumi{\arabic{enumi}.}
\tightlist
\item
  \(n=1\to2\) の励起エネルギー \(E_2-E_1\) と縮退度比 \(g_2/g_1\)
  を求めよ。
\item
  \(T=5800\) K（太陽光球）と \(T=10000\) K（A 型星）で占有数比
  \(n_2/n_1\) を計算し、何桁違うか述べよ。
\item
  Balmer 線（\(n=2\) から始まる可視吸収線）が低温の太陽で弱く、A
  型星（\(\sim10000\)
  K）で強くなる理由を述べよ。さらに高温では弱まる（電離が効く）ことに一言触れよ。
\end{enumerate}

\textbf{関連}：\hyperref[sec-statistical-mechanics-boltzmann]{§7.3
Boltzmann
分布}／恒星の色と温度は\hyperref[sec-where-blackbody-stars]{§1.3}、線強度からの温度推定は\href{../part7/24-line-applications.qmd}{第24章}、原子準位は\href{../part7/22-atomic-physics.qmd}{第22章}。

\section{模範解答（問題
7-4）}\label{ux6a21ux7bc4ux89e3ux7b54ux554fux984c-7-4}

\textbf{(1)} \(E_2-E_1=-13.6/4-(-13.6/1)=-3.4+13.6=10.2\)
eV。\(g_2/g_1=2\cdot2^2/(2\cdot1^2)=8/2=4\)。

\textbf{(2)} \(T=5800\) K：\(k_BT=8.617\times10^{-5}\times5800=0.500\)
eV。指数
\(-10.2/0.500=-20.4\)、\(e^{-20.4}=1.37\times10^{-9}\)。\(n_2/n_1=4\times1.37\times10^{-9}\approx5.5\times10^{-9}\)。
\(T=10000\) K：\(k_BT=0.862\) eV。指数
\(-10.2/0.862=-11.8\)、\(e^{-11.8}=7.2\times10^{-6}\)。\(n_2/n_1=4\times7.2\times10^{-6}\approx2.9\times10^{-5}\)。
比は \(2.9\times10^{-5}/5.5\times10^{-9}\approx5300\) 倍。約3〜4桁違う。

\textbf{(3)} Balmer 線は \(n=2\) にいる原子が吸収して生じる。太陽（5800
K）では \(n=2\) の占有が \(\sim10^{-9}\) と極めて少なく、Balmer
線は弱い。A 型星（\(\sim10000\) K）では \(n=2\)
占有が数千倍に増え、Balmer 線が最強になる。さらに高温（O・B
型）では水素が電離して中性水素自体が減り、Balmer 線は再び弱まる（電離は
Saha 式、本問の Boltzmann だけでは表せない別効果）。

\textbf{答え}：\(E_2-E_1=10.2\)
eV、\(g_2/g_1=4\)。\(n_2/n_1\approx5.5\times10^{-9}\)（5800
K）→\(2.9\times10^{-5}\)（10000 K）で約5300倍。Balmer 線は A
型で最強、さらに高温では電離で弱まる。

\section*{さらに学ぶための参考文献}\label{sec-statistical-mechanics-further}
\addcontentsline{toc}{section}{さらに学ぶための参考文献}

\markright{さらに学ぶための参考文献}

\begin{itemize}
\tightlist
\item
  Reif, \emph{Fundamentals of Statistical and Thermal Physics}
  (McGraw-Hill, 1965) --- Boltzmann 分布と分配関数の標準教科書
\item
  Pathria \& Beale, \emph{Statistical Mechanics} (Academic Press, 3rd
  ed., 2011) --- より現代的な体系化
\item
  久保亮五『大学演習　熱学・統計力学』（裳華房, 1961）---
  古典的な日本の教科書、演習豊富
\end{itemize}

\chapter{第8章　光子の統計}\label{ux7b2c8ux7ae0-ux5149ux5b50ux306eux7d71ux8a08}

\textbf{この章の主題}：第7章で立ち上げた古典統計（Boltzmann
分布）は、放射場に適用すると Rayleigh--Jeans
則の発散という決定的な破綻を起こした。本章では、光子を「区別できない
Bose 粒子」として扱い、Bose--Einstein
統計を立ち上げる。これによりプランク分布の占有因子 \(1/(e^{h\nu/kT}-1)\)
が必然として現れる。本章は第III部のクライマックスであり、中心地図のもっとも特徴的な構造を解き明かす。

\section*{この章の中心地図}\label{sec-photon-statistics-map}
\addcontentsline{toc}{section}{この章の中心地図}

\markright{この章の中心地図}

\textbf{中心地図：プランク関数 ― 本部で扱う因子}

\[
B_\nu(T) = \frac{2 h \nu^3}{c^2} \cdot \boxed{\frac{1}{e^{h\nu/kT} - 1}}
\]

強調されている \textbf{占有因子} \(1/(e^{h\nu/kT}-1)\)
が、第III部で逆引きする対象である。これは熱平衡にある光子（Bose
粒子）の統計から立ち上がる量で、低振動数極限では \(kT/(h\nu)\)
に、高振動数極限では \(e^{-h\nu/kT}\)
に漸近する。第8章で本格的に導出する。

\begin{tcolorbox}[enhanced jigsaw, arc=.35mm, bottomrule=.15mm, bottomtitle=1mm, breakable, colback=white, colbacktitle=quarto-callout-note-color!10!white, colframe=quarto-callout-note-color-frame, coltitle=black, left=2mm, leftrule=.75mm, opacityback=0, opacitybacktitle=0.6, rightrule=.15mm, title=\textcolor{quarto-callout-note-color}{\faInfo}\hspace{0.5em}{ノート}, titlerule=0mm, toprule=.15mm, toptitle=1mm]

\textbf{方針}：本章でターゲットとするのは中心地図プランク関数
\(B_\nu(T) = (2h\nu^3/c^2) \cdot \boxed{1/(e^{h\nu/kT}-1)}\)
の\textbf{第二因子}（占有因子）である。前因子 \(2h\nu^3/c^2\)
は第V部のモード密度で扱う。

\end{tcolorbox}

\section*{この章で答える問い}\label{sec-photon-statistics-questions}
\addcontentsline{toc}{section}{この章で答える問い}

\markright{この章で答える問い}

\begin{itemize}
\tightlist
\item
  なぜ光子の集団は古典粒子の統計（Maxwell--Boltzmann）ではなく、Bose--Einstein
  統計に従うのか
\item
  なぜ光子の化学ポテンシャルはゼロでなければならないのか
\item
  占有因子 \(1/(e^{h\nu/kT}-1)\)
  は、観測スペクトルのどの特徴に対応しているのか
\item
  光子気体の統計を仮定すると、なぜプランク分布が出てくるのか
\end{itemize}

\section*{到達目標}\label{sec-photon-statistics-goals}
\addcontentsline{toc}{section}{到達目標}

\markright{到達目標}

この章を読み終えると、読者は次のことができるようになる：

\begin{itemize}
\tightlist
\item
  Bose--Einstein 分布を、識別不能粒子の統計力学から導ける
\item
  光子の化学ポテンシャルがゼロである理由を、粒子数非保存の観点から説明できる
\item
  プランク分布の指数関数因子 \(1/(e^{h\nu/kT}-1)\)
  が、光子統計の帰結であることを論じられる
\end{itemize}

\begin{center}\rule{0.5\linewidth}{0.5pt}\end{center}

\section{8.1 区別できない粒子 ―
古典統計はなぜ破綻するか}\label{sec-photon-statistics-indistinguishable}

{[}本文目安：B3-B4{]}

第7章で扱った古典統計は、\textbf{粒子を識別できる}前提に立っていた。\(N\)
個の粒子に「1番、2番、\ldots、\(N\) 番」とラベルをつけ、それぞれが状態
\(i_1, i_2, \ldots\) にあるとき、ミクロ状態の数を数えた。

しかし、ミクロ世界の同種粒子は\textbf{区別できない}。電子1個と電子2個を入れ替えても、物理的に区別できる状況にはならない。これは古典物理の枠組みを超える性質で、量子力学（第IV部）から要請される事実である。

区別できない粒子の場合、ミクロ状態の数え方が変わる。状態 \(i\)
にある粒子数 \(n_i\) の組 \(\{n_i\}\)
だけがミクロ状態を特徴づける。\textbf{どの個別粒子がどこにいるかは問題にならない}。

具体例で比較する。状態が二つ（A, B）、粒子が二つの場合：

{\def\LTcaptype{none} % do not increment counter
\begin{longtable}[]{@{}ll@{}}
\toprule\noalign{}
識別可能（古典） & 識別不能（量子） \\
\midrule\noalign{}
\endhead
\bottomrule\noalign{}
\endlastfoot
1 が A、2 が A: 1 通り & (2, 0): 1 通り \\
1 が A、2 が B: 1 通り & (1, 1): 1 通り \\
1 が B、2 が A: 1 通り & (0, 2): 1 通り \\
1 が B、2 が B: 1 通り & \\
\textbf{合計: 4 状態} & \textbf{合計: 3 状態} \\
\end{longtable}
}

ここから「1 が A、2 が B」と「1 が B、2 が
A」が同じ状態に潰れることに注目したい。この潰れが Boltzmann 分布から
Bose 統計への移行を生む。

\begin{tcolorbox}[enhanced jigsaw, arc=.35mm, bottomrule=.15mm, bottomtitle=1mm, breakable, colback=white, colbacktitle=quarto-callout-note-color!10!white, colframe=quarto-callout-note-color-frame, coltitle=black, left=2mm, leftrule=.75mm, opacityback=0, opacitybacktitle=0.6, rightrule=.15mm, title=\textcolor{quarto-callout-note-color}{\faInfo}\hspace{0.5em}{ノート}, titlerule=0mm, toprule=.15mm, toptitle=1mm]

\textbf{用語}：区別できない粒子の統計は、整数スピン粒子なら
\textbf{Bose--Einstein 統計}、半整数スピン粒子なら \textbf{Fermi--Dirac
統計}になる。光子はスピン1の Bose 粒子なので、Bose
統計に従う。詳しくは第IV部で扱う。

\end{tcolorbox}

\section{8.2 Bose 粒子と Fermi 粒子 ―
二種類の量子統計}\label{sec-photon-statistics-bose-fermi}

{[}本文目安：B3-B4{]}

区別できない粒子は、状態を共有できるかどうかでさらに二つに分かれる：

\begin{itemize}
\tightlist
\item
  \textbf{Bose 粒子}（\textbf{boson},
  ボソン）：同じ状態に\textbf{何個でも}入れる。光子・グルーオン・\(W^\pm\)・\(Z^0\)・Higgs・\(^4\)He
  など。スピンが整数。
\item
  \textbf{Fermi 粒子}（\textbf{fermion},
  フェルミオン）：同じ状態には\textbf{1個しか}入れない（\textbf{Pauli
  の排他原理}, \textbf{Pauli exclusion
  principle}）。電子・陽子・中性子・クォーク・\(^3\)He
  など。スピンが半整数。
\end{itemize}

この区別は粒子の本質的な性質（スピン統計の定理）に基づくが、本書では結果だけを使う。証明は第IV部・第VIII部の場の量子論で扱われる。

光子は Bose 粒子なので、本章では Bose 統計（Bose--Einstein
分布）を扱う。

\section{8.3 Bose--Einstein 分布 ―
占有数の最確値}\label{sec-photon-statistics-bose-einstein}

{[}本文目安：B4｜導出：M1{]}

Bose 粒子の集団に最大エントロピー原理（第6章 §6.4）を適用する。

各「単一粒子状態」\(i\) に何個の粒子が入るかを \(n_i\) とする。\(n_i\)
は 0, 1, 2, \ldots{} と自由。状態 \(i\) の縮退度（同じエネルギー \(E_i\)
の状態数）を \(g_i\)、その状態を占有する粒子の総数を \(N_i\)
とすると、これらの占有を実現する場合の数（ミクロ状態数）は、

\begin{equation}\protect\phantomsection\label{eq-bose-multiplicity}{
W_\text{BE} = \prod_i \frac{(N_i + g_i - 1)!}{N_i!\,(g_i - 1)!}
}\end{equation}

これは「\(g_i\) 個の箱に \(N_i\)
個の区別できない玉を入れる組合せ」である。

エントロピー \(S = k_B \ln W\) を、制約 \(\sum N_i = N\),
\(\sum N_i E_i = E\) のもとで最大化する。Lagrange 未定乗数法（第6章 §6.4
の手順）と Stirling 近似で、結果は

\begin{equation}\protect\phantomsection\label{eq-bose-einstein}{
\bar{n}_i = \frac{1}{e^{(E_i - \mu)/k_B T} - 1}
}\end{equation}

これが \textbf{Bose--Einstein 分布}（\textbf{Bose--Einstein
distribution}）である。\(\bar{n}_i\) は状態 \(i\)
あたりの平均\textbf{占有数}（\textbf{occupation number}）。

\section{導出ノート（オンライン版限定）}\label{ux5c0eux51faux30ceux30fcux30c8ux30aaux30f3ux30e9ux30a4ux30f3ux7248ux9650ux5b9a}

式~\ref{eq-bose-multiplicity} の対数をとり、Stirling 近似
\(\ln N! \simeq N \ln N - N\) を使うと

\[
\ln W = \sum_i \left[ (N_i + g_i) \ln(N_i + g_i) - N_i \ln N_i - g_i \ln g_i \right]
\]

これに Lagrange 乗数
\(\alpha\)（粒子数保存）、\(\beta\)（エネルギー保存）を加えて

\[
\frac{\partial}{\partial N_i}\left[ \ln W - \alpha (\sum_j N_j - N) - \beta (\sum_j N_j E_j - E) \right] = 0
\]

から \(\ln(N_i + g_i) - \ln N_i = \alpha + \beta E_i\)、すなわち
\(\bar{n}_i = N_i/g_i = 1/(e^{\alpha + \beta E_i} - 1)\)
を得る。\(\beta = 1/k_B T\)、\(\alpha = -\mu/k_B T\) で
式~\ref{eq-bose-einstein} に至る。

式~\ref{eq-bose-einstein} はプランク分布の占有因子の形と一致する ―
ただし化学ポテンシャル \(\mu\) がついている。光子の場合は次節で
\(\mu = 0\) となる。

\begin{tcolorbox}[enhanced jigsaw, arc=.35mm, bottomrule=.15mm, breakable, colback=white, colframe=quarto-callout-tip-color-frame, left=2mm, leftrule=.75mm, opacityback=0, rightrule=.15mm, toprule=.15mm]
\begin{minipage}[t]{5.5mm}
\textcolor{quarto-callout-tip-color}{\faLightbulb}
\end{minipage}%
\begin{minipage}[t]{\textwidth - 5.5mm}

\textbf{問い}：Bose 分布の分母にある「\(-1\)」は、Fermi
分布の「\(+1\)」と何が違うのか？

\textbf{短答}：粒子の数え方の違いから来る。Bose
では「同じ状態に複数入れる」ため Σ\_n = 1/(1−x)（無限等比和）の因子から
\(-1\)、Fermi では「0 か 1 だけ」だから 1+x から \(+1\) になる。

\textbf{もう一歩}：分母の符号は、低温（\(E - \mu \ll k_B T\)）で巨大な違いを生む。Bose
では \(\bar{n} \to \infty\)（Bose--Einstein 凝縮の可能性）、Fermi では
\(\bar{n} \to 1\)（縮退、白色矮星・中性子星）。本書では光子のみ扱うが、第18章の中性子星表面で
Fermi 縮退に再会する。

\end{minipage}%
\end{tcolorbox}

\begin{tcolorbox}[enhanced jigsaw, arc=.35mm, bottomrule=.15mm, bottomtitle=1mm, breakable, colback=white, colbacktitle=quarto-callout-note-color!10!white, colframe=quarto-callout-note-color-frame, coltitle=black, left=2mm, leftrule=.75mm, opacityback=0, opacitybacktitle=0.6, rightrule=.15mm, title=\textcolor{quarto-callout-note-color}{\faInfo}\hspace{0.5em}{ノート}, titlerule=0mm, toprule=.15mm, toptitle=1mm]

\textbf{使えるようになった道具}（§8.3）：Bose--Einstein 占有数
式~\ref{eq-bose-einstein}
を、温度と化学ポテンシャルから具体的に書き下せる。これが観測されるプランク分布の指数関数因子の正体を与える。

\end{tcolorbox}

\section{8.4 光子の化学ポテンシャル ―
なぜゼロか}\label{sec-photon-statistics-mu-zero}

{[}本文目安：B4｜導出：M1{]}

光子は 式~\ref{eq-bose-einstein} で \textbf{\(\mu = 0\)}
が成り立つ。これがプランク分布のあの単純な形

\begin{equation}\protect\phantomsection\label{eq-photon-occupation}{
\bar{n}_\text{photon}(\nu, T) = \frac{1}{e^{h\nu/k_B T} - 1}
}\end{equation}

を生む。なぜ \(\mu = 0\) になるのか。

化学ポテンシャル \(\mu\) は、粒子を一個加えたときの自由エネルギーの変化
\(\partial F / \partial N\) である（第7章 §7.5）。これは\textbf{粒子数
\(N\)
が保存量である場合}にのみ意味を持つ。粒子数が保存しない系では、\(N\)
は熱平衡で「他の量と整合するように決まる」量となり、\(\mu\)
は粒子追加の代償なし ― すなわち \(\mu = 0\) ― になる。

光子の場合：

\begin{itemize}
\tightlist
\item
  電磁場と物質の相互作用（吸収・放出）で光子は\textbf{生成・消滅する}
\item
  光子の総数 \(N_\gamma\) は熱平衡で自動的に決まる量
\item
  よって \(\mu_\gamma = 0\)
\end{itemize}

これに対して、原子・分子・電子のような物質粒子は通常の温度・密度では数が保存し（化学反応がないかぎり）、\(\mu \ne 0\)
となる。

\begin{tcolorbox}[enhanced jigsaw, arc=.35mm, bottomrule=.15mm, bottomtitle=1mm, breakable, colback=white, colbacktitle=quarto-callout-note-color!10!white, colframe=quarto-callout-note-color-frame, coltitle=black, left=2mm, leftrule=.75mm, opacityback=0, opacitybacktitle=0.6, rightrule=.15mm, title=\textcolor{quarto-callout-note-color}{\faInfo}\hspace{0.5em}{ノート}, titlerule=0mm, toprule=.15mm, toptitle=1mm]

\textbf{観測との対応}：CMB
のスペクトルが完璧な黒体（\(\mu = 0\)）であることは、再結合以前に光子と物質が十分に長く相互作用して光子数を「自由に」調整できたことの証拠。もし
\(\mu \ne 0\) なら、CMB スペクトルは Bose--Einstein
分布の形（プランクから外れる）を残すはずで、これを \textbf{「\(\mu\)
歪み」} と呼ぶ。第16章で扱う。

\end{tcolorbox}

\begin{tcolorbox}[enhanced jigsaw, arc=.35mm, bottomrule=.15mm, breakable, colback=white, colframe=quarto-callout-tip-color-frame, left=2mm, leftrule=.75mm, opacityback=0, rightrule=.15mm, toprule=.15mm]
\begin{minipage}[t]{5.5mm}
\textcolor{quarto-callout-tip-color}{\faLightbulb}
\end{minipage}%
\begin{minipage}[t]{\textwidth - 5.5mm}

\textbf{問い}：光子の数が保存しないなら、宇宙の光子数はどこから決まっているのか？

\textbf{短答}：「光子数 \(N_\gamma\)
は熱平衡で自動的に決まる」とは、温度 \(T\) と空洞の体積 \(V\)
が決まれば、\(N_\gamma(T, V)\) が一意になることを意味する。具体的には
\(N_\gamma \propto V T^3\)。

\textbf{もう一歩}：観測される CMB の光子数密度 \(n_\gamma \simeq 411\)
cm\(^{-3}\) は、CMB 温度 \(T = 2.725\) K と \(n_\gamma \propto T^3\)
から一意に決まる。早期宇宙の物質との結合は、この光子数を「自由に」作り出してくれた。

\end{minipage}%
\end{tcolorbox}

\begin{tcolorbox}[enhanced jigsaw, arc=.35mm, bottomrule=.15mm, bottomtitle=1mm, breakable, colback=white, colbacktitle=quarto-callout-note-color!10!white, colframe=quarto-callout-note-color-frame, coltitle=black, left=2mm, leftrule=.75mm, opacityback=0, opacitybacktitle=0.6, rightrule=.15mm, title=\textcolor{quarto-callout-note-color}{\faInfo}\hspace{0.5em}{ノート}, titlerule=0mm, toprule=.15mm, toptitle=1mm]

\textbf{使えるようになった道具}（§8.4）：化学ポテンシャルがゼロな系（光子）の占有数
式~\ref{eq-photon-occupation}
は、温度と振動数だけで書き下せる。これが観測される CMB
の黒体スペクトルの指数関数因子の正体である。CMB が完全な黒体に近い
(\(\mu \lesssim 10^{-5}\)) ことは、\(\mu = 0\)
の前提（光子数の非保存）が宇宙史で確かに実現したことの観測的証拠になる（第16章
\(\mu\) 歪み）。

\end{tcolorbox}

\section{8.5 光子気体の統計 ―
エネルギー密度から圧力まで}\label{sec-photon-statistics-photon-gas}

{[}本文目安：B4｜導出：M1{]}

式~\ref{eq-photon-occupation}
を使って、放射場の各種熱力学量を計算できる。

\textbf{エネルギー密度}（振動数あたり）：振動数 \(\nu\)
近傍のモード密度を \(g(\nu) d\nu\)、各モードのエネルギーを
\(h\nu\)、平均光子数を \(\bar{n}\) とすると

\begin{equation}\protect\phantomsection\label{eq-photon-energy-density}{
u_\nu d\nu = g(\nu)\, h\nu\, \bar{n}(\nu, T)\, d\nu = g(\nu) \cdot \frac{h\nu}{e^{h\nu/k_B T} - 1}\, d\nu
}\end{equation}

\begin{tcolorbox}[enhanced jigsaw, arc=.35mm, bottomrule=.15mm, bottomtitle=1mm, breakable, colback=white, colbacktitle=quarto-callout-note-color!10!white, colframe=quarto-callout-note-color-frame, coltitle=black, left=2mm, leftrule=.75mm, opacityback=0, opacitybacktitle=0.6, rightrule=.15mm, title=\textcolor{quarto-callout-note-color}{\faInfo}\hspace{0.5em}{ノート}, titlerule=0mm, toprule=.15mm, toptitle=1mm]

\textbf{方針（先取りの注意）}：以下の計算でモード密度
\(g(\nu) = 8\pi\nu^2/c^3\)
を\textbf{先取りして借用する}。これは電磁場が三次元空間で離散的な定在波モードを持つことから出る量で、その本格的な導出は\textbf{第V部・第14章で行う}。ここでは結果だけを\textbf{仮に借りる}位置づけであり、第14章でこの借用を返す（節~\ref{sec-photon-statistics-photon-gas}
↔ \href{../part5/14-cavity-modes.qmd}{第14章}の往復で確認できる）。

本章の中心課題は \textbf{占有因子 \(1/(e^{h\nu/kT}-1)\) の導出}（前節
式~\ref{eq-photon-occupation}）である。ただしこれを放射場のエネルギー密度として観測できる量にするには、モード密度が必要である。その導出は第V部・第14章で行う。占有因子だけでは観測量にならず、「占有因子＋モード密度」の組合せで初めて
Planck
分布全体が見える。\textbf{全体形が見えてから、第V部でモード密度の起源に逆引きする}、というのが本書の戦略である。

\end{tcolorbox}

第V部から \(g(\nu) = 8\pi\nu^2/c^3\) を借りて代入すれば

\begin{equation}\protect\phantomsection\label{eq-planck-energy-density}{
u_\nu(T) = \frac{8\pi \nu^2}{c^3} \cdot \frac{h\nu}{e^{h\nu/k_B T} - 1} = \frac{8\pi h \nu^3}{c^3} \cdot \frac{1}{e^{h\nu/k_B T} - 1}
}\end{equation}

そして比強度 \(B_\nu = c u_\nu / (4\pi)\) から

\begin{equation}\protect\phantomsection\label{eq-planck-rederived}{
B_\nu(T) = \frac{2 h \nu^3}{c^2} \cdot \frac{1}{e^{h\nu/k_B T} - 1}
}\end{equation}

\textbf{プランク分布の全体形が、光子気体の Bose--Einstein
統計と（第V部から借りた）モード密度から再構成された}。本章で逆引きしたのは右側の
\textbf{占有因子}であり、前因子 \(2h\nu^3/c^2\)
の起源は第V部で引き続き扱う。

\textbf{全エネルギー密度}：式~\ref{eq-planck-energy-density}
を全振動数で積分（\(x = h\nu/k_B T\) の置換と
\(\int_0^\infty x^3/(e^x - 1) dx = \pi^4/15\)）すると

\begin{equation}\protect\phantomsection\label{eq-energy-density-T4}{
u(T) = \int_0^\infty u_\nu d\nu = \frac{4\sigma}{c} T^4 \equiv a T^4
}\end{equation}

ここで \(\sigma\) は Stefan--Boltzmann 定数。\textbf{\(T^4\)
の起源は、占有因子の積分結果として現れる}。

\textbf{圧力}：光子気体の圧力は、第7章 §7.5 の熱力学関数の方法で

\begin{equation}\protect\phantomsection\label{eq-photon-pressure}{
P = \frac{1}{3} u = \frac{a T^4}{3}
}\end{equation}

「\(P = u/3\)」の係数は、光子が質量ゼロ・等方分布の Bose
気体である帰結。気体の状態方程式（理想気体
\(PV = N k_B T\)）とは異なる。

\begin{tcolorbox}[enhanced jigsaw, arc=.35mm, bottomrule=.15mm, bottomtitle=1mm, breakable, colback=white, colbacktitle=quarto-callout-note-color!10!white, colframe=quarto-callout-note-color-frame, coltitle=black, left=2mm, leftrule=.75mm, opacityback=0, opacitybacktitle=0.6, rightrule=.15mm, title=\textcolor{quarto-callout-note-color}{\faInfo}\hspace{0.5em}{ノート}, titlerule=0mm, toprule=.15mm, toptitle=1mm]

\textbf{観測との対応}：恒星内部や降着円盤の中心部では、\textbf{放射圧}
\(P_\text{rad} = aT^4/3\) がガス圧 \(P_\text{gas} = n k_B T\)
より大きくなることがある。Eddington
限界はこの放射圧と重力のバランスから決まる量で、星の質量上限を与える。第18章で扱う。

\end{tcolorbox}

\section{8.6 統計力学から見たプランク分布 ―
第III部の到達点}\label{sec-photon-statistics-summary-section}

{[}本文目安：B4｜導出：M1{]}

第I部で「観測される連続スペクトルは黒体に近い」と確認し、第II部で「LTE
ではソース関数がプランク関数」と整理し、第III部の本章で
\textbf{「Bose--Einstein 統計から プランク分布の占有因子
\(1/(e^{h\nu/kT}-1)\) が必然として出る」} ことを示した。

これは中心地図プランク関数

\[
B_\nu(T) = \boxed{\frac{2 h \nu^3}{c^2}} \cdot \boxed{\frac{1}{e^{h\nu/k_B T} - 1}}
\]

の \textbf{指数関数因子}（右側、本章で解き明かした）が、統計力学（光子の
Bose 統計＋熱平衡）の帰結であることを意味する。

残りは：

\begin{itemize}
\tightlist
\item
  \textbf{前因子 \(2h\nu^3/c^2\)} ：モード密度
  \(\propto \nu^2\)、エネルギー量子 \(h\nu\)、偏光自由度 2
  の積。\textbf{第V部（電磁気学）と第IV部（量子論）の合作}で解き明かす
\item
  \textbf{\(h\nu\) 自体の意味}：エネルギーが連続でなく \(h\nu\)
  の単位で離散化されること。\textbf{第IV部（量子論）の中心テーマ}
\end{itemize}

第IV部・第V部でこれらの「なぜ」が解明されれば、中心地図プランク関数のすべての因子が逆引きで説明されたことになる。これが本書の第I〜V部全体の流れである。

\begin{center}\rule{0.5\linewidth}{0.5pt}\end{center}

\section*{この章で何がわかったか}\label{sec-photon-statistics-summary}
\addcontentsline{toc}{section}{この章で何がわかったか}

\markright{この章で何がわかったか}

\textbf{中心地図に戻る}

本章で、中心地図 \(B_\nu(T)\) の \textbf{指数関数因子
\(1/(e^{h\nu/kT}-1)\)} が解き明かされた：

\begin{itemize}
\tightlist
\item
  \textbf{Bose 統計}（光子が区別できないボソンであること）から
  \(1/(e^{(E-\mu)/k_B T} - 1)\) という形が出る
\item
  \textbf{光子の化学ポテンシャル
  \(\mu = 0\)}（光子数が保存しないこと）が分母を簡単にする
\item
  結果として \textbf{占有因子 \(1/(e^{h\nu/k_B T} - 1)\)}
  が現れ、これが指数関数的減衰（高振動数）と \(\nu^2\)
  立ち上がり（低振動数）の両方を統一的に与える
\end{itemize}

第I〜III部のロードマップが完結した： \textbf{観測される連続スペクトル →
熱平衡という条件 → 光子の Bose 統計 → 占有因子}

\textbf{次章へ}：第9〜12章では「指数の中の
\(h\nu\)」を逆引きする。なぜエネルギーが連続でなく \(h\nu\)
の量子化された単位で扱われるのか。これが量子論の起源と直結する。

\section*{演習問題}\label{sec-photon-statistics-exercises}
\addcontentsline{toc}{section}{演習問題}

\markright{演習問題}

\textbf{模範解答}
は折りたたみを展開して確認できる（オンライン版）。まず自力で解いてから開くこと。

\subsection*{問題 8-1　識別可能 vs 不能 ― 数え上げの違い}\label{ex-8-1}
\addcontentsline{toc}{subsection}{問題 8-1　識別可能 vs 不能 ―
数え上げの違い}

{[}★ 難易度：☆☆ {]}

§8.1 は2状態2粒子で古典4通り・Bose 3通りになることを示した。一般化する。

\begin{enumerate}
\def\labelenumi{\arabic{enumi}.}
\tightlist
\item
  \(g\) 個の状態に \(N\) 個の区別できない Bose 粒子を入れる場合の数が
  \(\binom{N+g-1}{N}\)
  になることを「仕切り棒（玉と仕切り）」の論法で示せ。
\item
  \(g=2,N=2\) で値が 3 になることを確かめ、古典の \(g^N=4\) と比べよ。
\item
  「1 が A・2 が B」と「1 が B・2 が A」が同一状態に潰れることが、なぜ
  Boltzmann 統計から Bose 統計へのずれを生むのか述べよ。
\end{enumerate}

\textbf{関連}：\hyperref[sec-photon-statistics-indistinguishable]{§8.1
区別できない粒子}、\hyperref[sec-photon-statistics-bose-fermi]{§8.2 Bose
と
Fermi}／識別不能性の量子論的根拠は\href{../part4/12-quantum-foundations.qmd}{第12章}。

\section{模範解答（問題
8-1）}\label{ux6a21ux7bc4ux89e3ux7b54ux554fux984c-8-1}

\textbf{(1)} \(N\) 個の区別できない玉を \(g\)
個の状態に分けるのは、\(N\) 個の玉と \(g-1\)
本の仕切り棒を一列に並べる場合の数に等しい。全 \(N+g-1\)
個の位置から玉の位置 \(N\) 個を選ぶので \(\binom{N+g-1}{N}\)。

\textbf{(2)} \(g=2,N=2\)：\(\binom{2+2-1}{2}=\binom{3}{2}=3\)。Bose
は3通り（(2,0),(1,1),(0,2)）。古典は各粒子が独立に2状態を選べて
\(2^2=4\)。識別不能だと1つ少ない。

\textbf{(3)} 古典では「1 が A・2 が B」と「1 が B・2 が
A」を別状態と数えるが、識別不能だと粒子のラベルに意味がなく同一状態に潰れる。この「数え過ぎの打ち消し」が状態数を変え、最大エントロピー分布が
Boltzmann（\(\propto e^{-\beta E}\)）から Bose（\(1/(e^{\beta E}-1)\)
型）へずれる起源になる。特に同じ状態に複数入れることが許される効果が分母の
\(-1\) を生む。

\textbf{答え}：場合の数 \(\binom{N+g-1}{N}\)（仕切り棒法）、\(g=N=2\) で
3（古典は4）。ラベルの消失が Bose 統計へのずれを生む。

\subsection*{問題 8-2　Bose--Einstein 分布の導出}\label{ex-8-2}
\addcontentsline{toc}{subsection}{問題 8-2　Bose--Einstein 分布の導出}

{[}★ 難易度：☆☆☆ {]}

§8.3
は導出をオンライン版の畳んだノートに収めた。これを自力で展開する。状態群
\(i\)（縮退度 \(g_i\)、占有 \(N_i\)）のミクロ状態数
\(W=\displaystyle\prod_i\frac{(N_i+g_i-1)!}{N_i!\,(g_i-1)!}\)。

\begin{enumerate}
\def\labelenumi{\arabic{enumi}.}
\tightlist
\item
  \(\ln W\) に Stirling 近似を適用し（\(g_i\gg1\) とみて
  \(g_i-1\approx g_i\)）、\(\ln W\simeq\sum_i[(N_i+g_i)\ln(N_i+g_i)-N_i\ln N_i-g_i\ln g_i]\)
  を示せ。
\item
  制約 \(\sum N_i=N\)、\(\sum N_iE_i=E\) のもとで Lagrange 乗数
  \(\alpha,\beta\) を用い、\(\partial/\partial N_i=0\) から
  \(\ln\dfrac{N_i+g_i}{N_i}=\alpha+\beta E_i\) を導け。
\item
  占有数 \(\bar n_i=N_i/g_i\)
  について解き、\(\bar n_i=\dfrac{1}{e^{\alpha+\beta E_i}-1}=\dfrac{1}{e^{(E_i-\mu)/k_BT}-1}\)
  を示せ（\(\beta=1/k_BT\)、\(\alpha=-\mu/k_BT\)）。
\end{enumerate}

\textbf{関連}：\hyperref[sec-photon-statistics-bose-einstein]{§8.3
Bose--Einstein 分布}／Lagrange
法の枠組みは\hyperref[sec-thermal-equilibrium-max-entropy]{§6.4}（\hyperref[ex-6-3]{演習
6-3}）、\(\mu=0\) は\hyperref[ex-8-3]{演習 8-3}。

\section{模範解答（問題
8-2）}\label{ux6a21ux7bc4ux89e3ux7b54ux554fux984c-8-2}

\textbf{(1)}
\(\ln W=\sum_i[\ln(N_i+g_i-1)!-\ln N_i!-\ln(g_i-1)!]\)。Stirling
\(\ln x!\simeq x\ln x-x\) と \(g_i-1\approx g_i\) を使うと、\(-x\)
の項は制約で打ち消し合い

\[
\ln W\simeq\sum_i\left[(N_i+g_i)\ln(N_i+g_i)-N_i\ln N_i-g_i\ln g_i\right].
\]

\textbf{(2)} 変分関数
\(\Phi=\ln W-\alpha\sum_j N_j-\beta\sum_j N_jE_j\)。\(N_i\) で微分：

\[
\frac{\partial\ln W}{\partial N_i}=\ln(N_i+g_i)+1-\ln N_i-1=\ln\frac{N_i+g_i}{N_i}.
\]

\[
\frac{\partial\Phi}{\partial N_i}=\ln\frac{N_i+g_i}{N_i}-\alpha-\beta E_i=0\ \Rightarrow\ \ln\frac{N_i+g_i}{N_i}=\alpha+\beta E_i.
\]

\textbf{(3)}
\(\dfrac{N_i+g_i}{N_i}=1+\dfrac{g_i}{N_i}=1+\dfrac{1}{\bar n_i}=e^{\alpha+\beta E_i}\)
ゆえ \(\dfrac{1}{\bar n_i}=e^{\alpha+\beta E_i}-1\)、すなわち

\[
\bar n_i=\frac{1}{e^{\alpha+\beta E_i}-1}=\frac{1}{e^{(E_i-\mu)/k_BT}-1}.
\]

\textbf{答え}：Stirling＋Lagrange で
\(\ln\frac{N_i+g_i}{N_i}=\alpha+\beta E_i\)、解いて
\(\bar n_i=1/(e^{(E_i-\mu)/k_BT}-1)\)。分母の \(-1\) が Bose の特徴。

\subsection*{問題 8-3　光子の μ=0 と CMB 光子数密度（T³
則）}\label{ex-8-3}
\addcontentsline{toc}{subsection}{問題 8-3　光子の μ=0 と CMB
光子数密度（T³ 則）}

{[}★ 難易度：☆☆☆ {]}

§8.4 は光子の化学ポテンシャルが \(\mu=0\) で
\(\bar n=1/(e^{h\nu/kT}-1)\) になり、\(n_\gamma\propto T^3\)
と述べた。確かめる。

\begin{enumerate}
\def\labelenumi{\arabic{enumi}.}
\tightlist
\item
  なぜ光子は \(\mu=0\) なのか、粒子数の保存／非保存の観点から述べよ。
\item
  光子数密度
  \(n_\gamma=\displaystyle\int_0^\infty g(\nu)\bar n(\nu)\,d\nu\)（\(g(\nu)=8\pi\nu^2/c^3\)）で
  \(x=h\nu/k_BT\)
  と置換し、\(n_\gamma=\dfrac{8\pi}{c^3}\left(\dfrac{k_BT}{h}\right)^3\displaystyle\int_0^\infty\frac{x^2}{e^x-1}dx\)
  を導け。
\item
  \(\displaystyle\int_0^\infty\frac{x^2}{e^x-1}dx=2\zeta(3)\approx2.404\)
  を用い、\(T_\mathrm{CMB}=2.725\) K で \(n_\gamma\) を計算して
  \(\approx411\) cm\(^{-3}\) を確認せよ。
\end{enumerate}

\textbf{関連}：\hyperref[sec-photon-statistics-mu-zero]{§8.4
光子の化学ポテンシャル}／\(\mu\) 歪みと CMB
は\href{../part6/16-cmb.qmd}{第16章}、モード密度 \(g(\nu)\)
の導出は\href{../part5/14-cavity-modes.qmd}{第14章}、\(\int x^3\)
版（エネルギー密度）は\hyperref[ex-10-3]{演習 10-3}。

\section{模範解答（問題
8-3）}\label{ux6a21ux7bc4ux89e3ux7b54ux554fux984c-8-3}

\textbf{(1)} 化学ポテンシャルは粒子数 \(N\)
が保存量のときに意味を持つ「粒子1個追加の自由エネルギー代償」。光子は物質との吸収・放出で生成・消滅し、総数が保存しない。熱平衡では
\(N_\gamma\)
が温度・体積から自動的に決まる従属量となり、追加の代償がない＝\(\mu=0\)。

\textbf{(2)} \(x=h\nu/k_BT\)、\(\nu=k_BTx/h\)、\(d\nu=(k_BT/h)dx\)：

\[
n_\gamma=\int_0^\infty\frac{8\pi\nu^2}{c^3}\frac{d\nu}{e^{h\nu/k_BT}-1}=\frac{8\pi}{c^3}\left(\frac{k_BT}{h}\right)^3\int_0^\infty\frac{x^2}{e^x-1}dx.
\]

\textbf{(3)} 積分 \(=2.404\)。\(k_BT/h\)
を計算：\(k_BT=1.381\times10^{-23}\times2.725=3.76\times10^{-23}\)
J、\(/h=3.76\times10^{-23}/6.626\times10^{-34}=5.68\times10^{10}\)
s\(^{-1}\)。3乗
\(=1.83\times10^{32}\)。\(8\pi/c^3=25.13/(3.00\times10^8)^3=9.31\times10^{-25}\)。掛け合わせ：

\[
n_\gamma=9.31\times10^{-25}\times1.83\times10^{32}\times2.404\approx4.1\times10^8\ \mathrm{m^{-3}}=411\ \mathrm{cm^{-3}}.
\]

観測値とよく一致。\(n_\gamma\propto T^3\) は前因子 \((k_BT/h)^3\)
から明らか。

\textbf{答え}：光子数非保存ゆえ
\(\mu=0\)。\(n_\gamma=\frac{8\pi}{c^3}(k_BT/h)^3\cdot2.404\approx411\)
cm\(^{-3}\)（\(T=2.725\) K）、\(\propto T^3\)。

\subsection*{問題 8-4　光子気体の P=u/3 と放射圧 vs
ガス圧}\label{ex-8-4}
\addcontentsline{toc}{subsection}{問題 8-4　光子気体の P=u/3 と放射圧 vs
ガス圧}

{[}★ 難易度：☆☆ {]}

§8.5 は光子気体の圧力 \(P=u/3=aT^4/3\)
を与えた。理想気体と比べ、放射圧が効く条件を探る。

\begin{enumerate}
\def\labelenumi{\arabic{enumi}.}
\tightlist
\item
  等方放射場では \(P_\nu=u_\nu/3\)（演習 4-1）。これが理想気体
  \(P_\mathrm{gas}=nk_BT\) と「温度・密度への依存」でどう違うか述べよ。
\item
  放射圧 \(P_\mathrm{rad}=aT^4/3\) とガス圧 \(P_\mathrm{gas}=nk_BT\)
  が等しくなる条件 \(T^3/n=3k_B/a\)
  を導け。これは高温・低密度で放射圧が優勢になることを示す。
\item
  太陽中心（\(T\simeq1.5\times10^7\) K、\(n\simeq10^{32}\)
  m\(^{-3}\)）で \(P_\mathrm{rad}/P_\mathrm{gas}\)
  を概算し、放射圧が支配的でない（がより重い星では効く）ことを確認せよ。\(a=7.566\times10^{-16}\)
  J m\(^{-3}\)K\(^{-4}\)。
\end{enumerate}

\textbf{関連}：\hyperref[sec-photon-statistics-photon-gas]{§8.5
光子気体の統計}／\(P_\nu=u_\nu/3\) の導出は\hyperref[ex-4-1]{演習
4-1}、Eddington
限界・降着は\href{../part6/18-accretion-high-energy.qmd}{第18章}。

\section{模範解答（問題
8-4）}\label{ux6a21ux7bc4ux89e3ux7b54ux554fux984c-8-4}

\textbf{(1)} 放射圧 \(P_\mathrm{rad}=u/3=aT^4/3\)
は\textbf{温度の4乗}に効き密度には依らない（光子数自体が \(T^3\)
で決まる）。理想気体 \(P_\mathrm{gas}=nk_BT\) は密度 \(n\)
に比例し温度には1乗。よって温度を上げ密度を下げるほど放射圧が相対的に強くなる。

\textbf{(2)} \(P_\mathrm{rad}=P_\mathrm{gas}\)
とおく：\(\dfrac{aT^4}{3}=nk_BT\Rightarrow\dfrac{T^3}{n}=\dfrac{3k_B}{a}\)。左辺が大きい（高温・低密度）ほど放射圧優勢。

\textbf{(3)}
\(P_\mathrm{rad}=aT^4/3=7.566\times10^{-16}\times(1.5\times10^7)^4/3\)。\((1.5\times10^7)^4=5.06\times10^{28}\)。\(\times7.566\times10^{-16}=3.83\times10^{13}\)、\(/3=1.28\times10^{13}\)
Pa。
\(P_\mathrm{gas}=nk_BT=10^{32}\times1.381\times10^{-23}\times1.5\times10^7=2.07\times10^{16}\)
Pa。 比
\(P_\mathrm{rad}/P_\mathrm{gas}\approx1.28\times10^{13}/2.07\times10^{16}\approx6\times10^{-4}\)。太陽中心では放射圧は約0.06\%で支配的でない。質量の大きい星ほど中心温度が高く密度が相対的に低いため放射圧が増し、超大質量星では放射圧が支配的になって
Eddington 限界（星の質量上限）を与える。

\textbf{答え}：\(P_\mathrm{rad}=aT^4/3\)（密度非依存）、\(P_\mathrm{gas}=nk_BT\)。等しい条件
\(T^3/n=3k_B/a\)。太陽中心で
\(P_\mathrm{rad}/P_\mathrm{gas}\approx6\times10^{-4}\)（重い星では支配的）。

\section*{さらに学ぶための参考文献}\label{sec-photon-statistics-further}
\addcontentsline{toc}{section}{さらに学ぶための参考文献}

\markright{さらに学ぶための参考文献}

\begin{itemize}
\tightlist
\item
  Pathria \& Beale, \emph{Statistical Mechanics} (Academic Press, 3rd
  ed., 2011) --- 第6章「Boson Statistics」が標準的
\item
  Huang, \emph{Statistical Mechanics} (Wiley, 2nd ed., 1987) ---
  Bose--Einstein 凝縮の体系的扱い
\item
  Mandl, \emph{Statistical Physics} (Wiley, 2nd ed., 1988) ---
  学部向けに丁寧な扱い
\item
  田崎晴明『統計力学 II』（培風館, 2008）--- Bose 統計と Planck
  分布の和文教科書
\end{itemize}

\part{第IV部　逆引き2：量子論へ}

\chapter{第9章　古典論ではなぜ失敗するのか}\label{ux7b2c9ux7ae0-ux53e4ux5178ux8ad6ux3067ux306fux306aux305cux5931ux6557ux3059ux308bux306eux304b}

\textbf{この章の主題}：第III部までで、観測される黒体スペクトルが熱平衡にある光子気体の
\textbf{Bose 統計} から導かれることを見た。しかし、その手前で
\textbf{古典論（古典電磁気＋古典統計）だけで黒体スペクトルを説明しようとすると、なぜ失敗するのか}
を本章で明示的に追体験する。「紫外線破綻」（\textbf{ultraviolet
catastrophe}）と呼ばれる発散現象が、観測スペクトルとどう決定的に矛盾するのかを示し、量子論を要請せざるを得なくなった歴史的・論理的な瞬間を再現する。第10章のプランクの量子仮説、第11章の
Einstein の方法は、いずれもこの破綻から出発する。

\section*{この章の中心地図}\label{sec-classical-failure-map}
\addcontentsline{toc}{section}{この章の中心地図}

\markright{この章の中心地図}

\textbf{中心地図：プランク関数 ― 本部で扱う因子}

\[
B_\nu(T) = \frac{2\, \boxed{h}\, \nu^3}{c^2} \cdot \frac{1}{e^{\boxed{h}\nu/kT} - 1}
\]

強調されている \textbf{Planck 定数
\(h\)}（前因子の中と、指数の中の二箇所）が、第IV部で逆引きする対象である。エネルギーが連続でなく
\(h\nu\) という離散単位でやりとりされる ―
これが量子論の核心で、第10章（Planck の量子仮説）と第11章（Einstein
の方法）でその必然性を二つの異なるルートから示す。

\begin{tcolorbox}[enhanced jigsaw, arc=.35mm, bottomrule=.15mm, bottomtitle=1mm, breakable, colback=white, colbacktitle=quarto-callout-note-color!10!white, colframe=quarto-callout-note-color-frame, coltitle=black, left=2mm, leftrule=.75mm, opacityback=0, opacitybacktitle=0.6, rightrule=.15mm, title=\textcolor{quarto-callout-note-color}{\faInfo}\hspace{0.5em}{ノート}, titlerule=0mm, toprule=.15mm, toptitle=1mm]

\textbf{方針}：本章では、強調された \textbf{Planck 定数 \(h\)}
を\textbf{まだ持っていない}という設定で議論を進める。つまり「古典電磁気＋古典統計」だけで放射場のスペクトルを計算してみる。結果として、観測との決定的な矛盾が出る
― それが量子論を要請する観測上の根拠である。

\end{tcolorbox}

\section*{この章で答える問い}\label{sec-classical-failure-questions}
\addcontentsline{toc}{section}{この章で答える問い}

\markright{この章で答える問い}

\begin{itemize}
\tightlist
\item
  空洞中の電磁波を古典統計で扱うと、なぜ全エネルギーが発散してしまうのか
\item
  Rayleigh--Jeans
  則が低振動数側ではよく合うのに、高振動数側で失敗するのはなぜか
\item
  「紫外線破綻」と呼ばれる現象は、観測スペクトルのどの形と矛盾するのか
\item
  この破綻を回避するために、何を変えなければならなかったのか
\end{itemize}

\section*{到達目標}\label{sec-classical-failure-goals}
\addcontentsline{toc}{section}{到達目標}

\markright{到達目標}

この章を読み終えると、読者は次のことができるようになる：

\begin{itemize}
\tightlist
\item
  古典電磁気＋古典統計で Rayleigh--Jeans 則が導かれる筋道をたどれる
\item
  紫外線破綻の物理的起源を、エネルギー等分配定理から説明できる
\item
  量子仮説の必然性を、観測との矛盾から論じられる
\end{itemize}

\begin{center}\rule{0.5\linewidth}{0.5pt}\end{center}

\section{9.1 空洞モードの数え上げ ―
古典電磁気の準備}\label{sec-classical-failure-modes}

{[}本文目安：B3{]}

古典の議論を始めるには、まず\textbf{空洞中にどんな電磁波（モード,
mode）が存在できるか}を数える必要がある。

辺長 \(L\)
の立方体空洞を考え、壁が完全導体だと仮定する。電場が壁面で消えるという境界条件から、許される定在波の波数は

\begin{equation}\protect\phantomsection\label{eq-cavity-wavevector}{
\mathbf{k} = \frac{\pi}{L} (n_x, n_y, n_z), \qquad n_x, n_y, n_z = 1, 2, 3, \ldots
}\end{equation}

の形に制限される。波数 \(k = |\mathbf{k}|\) が \(k\)
未満であるモード数を数えると、立体角を考慮した結果として、\textbf{波数空間の体積}を使って数えられる。

詳細は\textbf{第14章で導出}する（ここでは仮に借りる）。\textbf{本章では結果だけを借用し}、第14章でこの借用を返す（節~\ref{sec-classical-failure-modes}
↔ \href{../part5/14-cavity-modes.qmd}{第14章
空洞モード}の往復で確認できる）。\textbf{振動数 \(\nu\)
あたりの単位体積あたりのモード数}は

\begin{equation}\protect\phantomsection\label{eq-mode-density-borrowed}{
g(\nu) = \frac{8\pi \nu^2}{c^3}
}\end{equation}

である。係数 \(8\pi\) には電磁波の\textbf{偏光自由度 \(\times 2\)}
が入っている。この \(\nu^2\)
という振動数の冪が、本章の議論の核心になる。

\begin{tcolorbox}[enhanced jigsaw, arc=.35mm, bottomrule=.15mm, bottomtitle=1mm, breakable, colback=white, colbacktitle=quarto-callout-note-color!10!white, colframe=quarto-callout-note-color-frame, coltitle=black, left=2mm, leftrule=.75mm, opacityback=0, opacitybacktitle=0.6, rightrule=.15mm, title=\textcolor{quarto-callout-note-color}{\faInfo}\hspace{0.5em}{ノート}, titlerule=0mm, toprule=.15mm, toptitle=1mm]

\textbf{方針（仮に借りる）}：式~\ref{eq-mode-density-borrowed}
の導出には\textbf{量子論は一切要らない}。古典電磁気と境界条件だけで出る幾何的な数え上げである。本格的な導出は\textbf{第14章で行う}ので、本章では「古典電磁気が許す結果」として一度受け入れ（仮に借り）、第14章で借用を返す。中心地図の各因子について「初出／仮に借りる場所／本当に導出する場所」は第14章で一覧化される。

\end{tcolorbox}

\begin{tcolorbox}[enhanced jigsaw, arc=.35mm, bottomrule=.15mm, breakable, colback=white, colframe=quarto-callout-tip-color-frame, left=2mm, leftrule=.75mm, opacityback=0, rightrule=.15mm, toprule=.15mm]
\begin{minipage}[t]{5.5mm}
\textcolor{quarto-callout-tip-color}{\faLightbulb}
\end{minipage}%
\begin{minipage}[t]{\textwidth - 5.5mm}

\textbf{問い}：なぜ三次元空間でモード密度が \(\nu^2\) に比例するのか？

\textbf{短答}：波数空間の球殻 \(4\pi k^2 dk\)
に許されるモードがあるから。\(k = 2\pi\nu/c\) なので
\(k^2 dk \propto \nu^2 d\nu\)。

\textbf{もう一歩}：一般に \(d\) 次元では
\(g(\nu) \propto \nu^{d-1}\)。三次元の宇宙では \(\nu^2\)、二次元なら
\(\nu\)、一次元なら定数。この幾何的事実が、観測される黒体スペクトルの低振動数側
\(B_\nu \propto \nu^2\) の振る舞いを直接決定している。

\end{minipage}%
\end{tcolorbox}

\section{9.2 エネルギー等分配定理と Rayleigh--Jeans
則}\label{sec-classical-failure-rayleigh-jeans}

{[}本文目安：B3{]}

次に、各モードに\textbf{平均でどれだけのエネルギー}が分配されるかを古典統計で求める。

古典統計力学の中心定理 ―
\textbf{エネルギー等分配定理}（\textbf{equipartition theorem}, 第7章
§7.6）― によれば、二次形式で書ける自由度ひとつあたりに
\(\frac{1}{2} k_B T\)
が分配される。電磁波の各モードは調和振動子と等価で、\(E = \frac{1}{2}(p^2/m) + \frac{1}{2}(m\omega^2 q^2)\)
という \textbf{二つの二次形式の和}（運動量項と位置項）―
あるいは電磁場の言葉では電場と磁場のエネルギー密度
\(\frac{\epsilon_0}{2}|\mathbf{E}|^2 + \frac{1}{2\mu_0}|\mathbf{B}|^2\)
の二つ ― として書ける。それぞれに \(\frac{1}{2}k_B T\)
が分配されるので、\textbf{各モードに合計 \(k_B T\) が分配される}：

\begin{equation}\protect\phantomsection\label{eq-equipartition-mode}{
\langle E_\text{mode} \rangle_\text{classical} = k_B T
}\end{equation}

これと 式~\ref{eq-mode-density-borrowed} を組み合わせると、振動数
\(\nu\) あたりのエネルギー密度は

\begin{equation}\protect\phantomsection\label{eq-rayleigh-jeans-energy-density}{
u_\nu^\text{RJ}(T) = g(\nu) \cdot \langle E_\text{mode} \rangle = \frac{8\pi \nu^2}{c^3} \cdot k_B T
}\end{equation}

そして比強度に直すと

\begin{equation}\protect\phantomsection\label{eq-rayleigh-jeans-bnu}{
B_\nu^\text{RJ}(T) = \frac{c}{4\pi} u_\nu^\text{RJ} = \frac{2\nu^2}{c^2} k_B T
}\end{equation}

これが \textbf{Rayleigh--Jeans 則}（\textbf{Rayleigh--Jeans
law}）である。古典電磁気＋古典統計だけから出た、放射場のスペクトルの予言である。

\begin{tcolorbox}[enhanced jigsaw, arc=.35mm, bottomrule=.15mm, bottomtitle=1mm, breakable, colback=white, colbacktitle=quarto-callout-note-color!10!white, colframe=quarto-callout-note-color-frame, coltitle=black, left=2mm, leftrule=.75mm, opacityback=0, opacitybacktitle=0.6, rightrule=.15mm, title=\textcolor{quarto-callout-note-color}{\faInfo}\hspace{0.5em}{ノート}, titlerule=0mm, toprule=.15mm, toptitle=1mm]

\textbf{観測との対応}：第2章 §2.5
で見たように、式~\ref{eq-rayleigh-jeans-bnu}
は\textbf{観測されるプランク分布の低振動数極限 \(h\nu \ll k_B T\)
と一致する}。電波天文学はほぼ常にこの極限で観測しており、Rayleigh--Jeans
は実用的に正しい。問題は次節で見るように、\textbf{高振動数側で破綻する}こと。

\end{tcolorbox}

\section{9.3 紫外線破綻 ―
全エネルギーの発散}\label{sec-classical-failure-uv-catastrophe}

{[}本文目安：B3{]}

\begin{figure}[H]

\centering{

\includegraphics[width=0.9\linewidth,height=\textheight,keepaspectratio]{index_files/mediabag/part4/../images/uv-catastrophe.pdf}

}

\caption{\label{fig-uv-catastrophe}紫外線破綻 ― 古典の Rayleigh--Jeans
則（\textbf{\(u_\nu \propto \nu^2\)、青破線}）が高振動数で発散する一方、観測されるプランク分布（\textbf{黒実線}）は
Wien 則的に指数的減衰する。両者は \textbf{\(h\nu \approx k_B T\)}
の境界帯（網かけ）で分かれる。本図は §2.5（Planck/RJ/Wien 比較）と §10
の量子仮説導入とも共用する基準図。}

\end{figure}%

式~\ref{eq-rayleigh-jeans-energy-density}
を\textbf{全振動数で積分}して、放射場の全エネルギー密度 \(u\)
を計算しようとすると

\begin{equation}\protect\phantomsection\label{eq-uv-catastrophe}{
u_\text{classical} = \int_0^\infty u_\nu^\text{RJ} \, d\nu = \frac{8\pi k_B T}{c^3} \int_0^\infty \nu^2 \, d\nu = \infty
}\end{equation}

積分は高振動数（紫外〜X線〜ガンマ線）側で\textbf{発散}する。これが
\textbf{紫外線破綻}（\textbf{ultraviolet catastrophe}）である。

物理的にこれが意味するのは、「温度 \(T\)
の空洞は無限大のエネルギー密度を持つ」「電球を点けたら無限の紫外線・X
線が出る」 ― 観測される現実とまったく合わない結論である。

観測との矛盾を具体的に見ると：

{\def\LTcaptype{none} % do not increment counter
\begin{longtable}[]{@{}
  >{\raggedright\arraybackslash}p{(\linewidth - 4\tabcolsep) * \real{0.3333}}
  >{\raggedright\arraybackslash}p{(\linewidth - 4\tabcolsep) * \real{0.3333}}
  >{\raggedright\arraybackslash}p{(\linewidth - 4\tabcolsep) * \real{0.3333}}@{}}
\toprule\noalign{}
\begin{minipage}[b]{\linewidth}\raggedright
量
\end{minipage} & \begin{minipage}[b]{\linewidth}\raggedright
観測
\end{minipage} & \begin{minipage}[b]{\linewidth}\raggedright
古典予言
\end{minipage} \\
\midrule\noalign{}
\endhead
\bottomrule\noalign{}
\endlastfoot
全エネルギー密度 \(u(T)\) & \(aT^4\)（有限） & 無限大 \\
ピーク振動数 & \(\nu_\text{peak} \propto T\)（Wien 則） &
ピーク自体が存在しない \\
高振動数側の振る舞い & \(e^{-h\nu/k_B T}\) で指数減衰 & \(\nu^2\)
で単調増加 \\
\end{longtable}
}

特に「ピークが存在しない」「指数減衰がない」という二点は、第1章で見た\textbf{観測される黒体スペクトル}と決定的に矛盾する。

\begin{tcolorbox}[enhanced jigsaw, arc=.35mm, bottomrule=.15mm, breakable, colback=white, colframe=quarto-callout-tip-color-frame, left=2mm, leftrule=.75mm, opacityback=0, rightrule=.15mm, toprule=.15mm]
\begin{minipage}[t]{5.5mm}
\textcolor{quarto-callout-tip-color}{\faLightbulb}
\end{minipage}%
\begin{minipage}[t]{\textwidth - 5.5mm}

\textbf{問い}：紫外線破綻はどの観測事実から致命的に証明されるのか？

\textbf{短答}：恒星のスペクトルがピークを持つこと、CMB
が完全な黒体形を取ること、地上で測られる空洞放射のスペクトル形 ―
いずれも有限のピークと指数減衰を示す。古典論ではこれらが一切説明できない。

\textbf{もう一歩}：歴史的には Rubens \& Kurlbaum (1900)
の長波長黒体スペクトルの精密測定が決め手となった。Planck
はこのデータを再現するために量子仮説を立てた（第10章）。観測が理論を駆動した古典的な例である。

\end{minipage}%
\end{tcolorbox}

\section{9.4 何が足りなかったのか ―
等分配の物理的前提}\label{sec-classical-failure-missing}

{[}本文目安：B3{]}

紫外線破綻はどこから来たか。式~\ref{eq-equipartition-mode}
の「各モードに
\(k_B T\)」という等分配定理に\textbf{穴がある}。等分配定理が成立する前提は：

\begin{enumerate}
\def\labelenumi{\arabic{enumi}.}
\tightlist
\item
  \textbf{エネルギーが連続的にやりとりできる}：粒子間の衝突や相互作用で、任意の小さなエネルギー単位でやりとりが起きる
\item
  \textbf{エネルギーが Maxwell-Boltzmann
  分布に従う}：状態数密度が均一で、状態間の遷移が自由
\end{enumerate}

ここで重要なのは前提 1
である。古典物理では「エネルギーは連続量」と暗黙に仮定している。しかしもしエネルギーが\textbf{離散単位}
\(\epsilon\) でしかやりとりできず、しかも \(\epsilon \gg k_B T\)
となる場面では：

\begin{itemize}
\tightlist
\item
  そのモードを励起するために必要なエネルギー \(\epsilon\)
  がほとんど常に手に入らない
\item
  平均してそのモードに行くエネルギーは \(k_B T\) よりも遥かに小さくなる
\item
  高振動数側のモードが「凍結」（freeze out）する
\end{itemize}

これが\textbf{紫外線破綻を救う鍵}である。すなわち、放射のエネルギーがある離散単位
\(\epsilon(\nu)\) でやりとりされ、しかも \(\epsilon(\nu)\) が振動数
\(\nu\) とともに増える ― 具体的には \(\epsilon \propto \nu\) ―
なら、高振動数側のモードが熱力学的に届かなくなる。これが Planck (1900)
のたどり着いた仮説で、第10章で本格的に扱う。

\begin{tcolorbox}[enhanced jigsaw, arc=.35mm, bottomrule=.15mm, bottomtitle=1mm, breakable, colback=white, colbacktitle=quarto-callout-note-color!10!white, colframe=quarto-callout-note-color-frame, coltitle=black, left=2mm, leftrule=.75mm, opacityback=0, opacitybacktitle=0.6, rightrule=.15mm, title=\textcolor{quarto-callout-note-color}{\faInfo}\hspace{0.5em}{ノート}, titlerule=0mm, toprule=.15mm, toptitle=1mm]

\textbf{注意（理論）}：「等分配が破綻する」という現象は、黒体放射以外にも現れる。低温での\textbf{固体の比熱}（Debye/Einstein
モデル）、二原子分子の\textbf{回転・振動自由度の凍結}、Bose--Einstein
凝縮など、量子効果が顔を出すすべての場面で同じ構造が現れる。黒体放射は最も歴史的に重要な例である。

\end{tcolorbox}

\section{9.5 量子仮説の必然性 ―
第10・11章への入口}\label{sec-classical-failure-necessity}

{[}本文目安：B3{]}

ここまでの議論をまとめると、観測される黒体スペクトルを古典論で説明できない理由は次のように要約される：

\begin{enumerate}
\def\labelenumi{\arabic{enumi}.}
\tightlist
\item
  空洞の電磁モード密度 \(g(\nu) \propto \nu^2\)
  そのものは古典電磁気で正しく出る（第V部）
\item
  各モードの平均エネルギーを古典統計で \(k_B T\)
  と置く（等分配定理）と、\(u_\nu \propto \nu^2\) が出てしまう
\item
  これを全振動数で積分すると \(u = \infty\)（紫外線破綻）
\item
  観測される高振動数の指数減衰を再現するには、\textbf{等分配を破る何か}が必要
\end{enumerate}

ここから出る論理的選択肢は限られる。一つは\textbf{エネルギー単位の離散化}（Planck
の量子仮説、第10章）。もう一つは\textbf{光が粒子としてふるまう}こと（Einstein
の光量子、第11章）。両者は等価で、結局は同じプランク分布に至るが、思考の出発点が違う。

両方のルートは、本書の三原則のうち「②
背景物理を曖昧にしない」を実装する重要な複線である。第10章でプランクの方法を、第11章で
Einstein の方法をたどることで、プランク分布の \textbf{占有因子
\(1/(e^{h\nu/k_B T}-1)\)}
がどちらのルートからも必然として出ることを確認する。第8章では Bose
統計から導いた占有因子が、二つの別の角度からも再導出される。

\begin{tcolorbox}[enhanced jigsaw, arc=.35mm, bottomrule=.15mm, bottomtitle=1mm, breakable, colback=white, colbacktitle=quarto-callout-note-color!10!white, colframe=quarto-callout-note-color-frame, coltitle=black, left=2mm, leftrule=.75mm, opacityback=0, opacitybacktitle=0.6, rightrule=.15mm, title=\textcolor{quarto-callout-note-color}{\faInfo}\hspace{0.5em}{ノート}, titlerule=0mm, toprule=.15mm, toptitle=1mm]

\textbf{使えるようになった道具}（§9.1〜9.5）：

\begin{itemize}
\tightlist
\item
  空洞の電磁モード密度 \(g(\nu) = 8\pi\nu^2/c^3\)
  は古典電磁気の結果として理解できる（第V部で本格導出）
\item
  Rayleigh--Jeans 則 式~\ref{eq-rayleigh-jeans-bnu}
  は等分配定理から出る古典の結果で、低振動数極限では観測と合う
\item
  紫外線破綻は等分配定理が高振動数モードに無条件で適用されることに由来する。これを救う鍵は
  \textbf{エネルギーの離散化}
\end{itemize}

\end{tcolorbox}

\begin{center}\rule{0.5\linewidth}{0.5pt}\end{center}

\section*{この章で何がわかったか}\label{sec-classical-failure-summary}
\addcontentsline{toc}{section}{この章で何がわかったか}

\markright{この章で何がわかったか}

\textbf{中心地図に戻る}

本章で、中心地図プランク関数の前因子に現れる \(\nu^2\)
が\textbf{古典電磁気のモード密度に由来する}ことを確認した。一方で、各モードの平均エネルギーを古典統計で
\(k_B T\) と置くと\textbf{観測スペクトルが再現できない}（紫外線破綻）。

中心地図の指数関数因子
\(1/(e^{h\nu/k_B T}-1)\)（占有因子）には、古典では出てこない\textbf{離散的なエネルギー量子
\(h\nu\)} が含まれる。この \(h\) こそが、本部で逆引きする対象である。

\textbf{次章へ}：第10章で Planck の量子仮説 ― 振動子のエネルギーが
\(h\nu\) の整数倍に限られるという思い切った仮定 ―
を導入し、これがプランク分布を再現することを示す。歴史的には Planck
(1900) の論文に当たる。

\section*{演習問題}\label{sec-classical-failure-exercises}
\addcontentsline{toc}{section}{演習問題}

\markright{演習問題}

\textbf{模範解答}
は折りたたみを展開して確認できる（オンライン版）。まず自力で解いてから開くこと。

\subsection*{問題 9-1　空洞モード密度 g(ν)=8πν²/c³
の数え上げ}\label{ex-9-1}
\addcontentsline{toc}{subsection}{問題 9-1　空洞モード密度 g(ν)=8πν²/c³
の数え上げ}

{[}★ 難易度：☆☆ {]}

本文 §9.1 はモード密度 \(g(\nu)=8\pi\nu^2/c^3\)
を「第14章で本格導出する結果」として借用している。この幾何的な数え上げを自分で再現せよ。辺長
\(L\) の立方体（完全導体壁）で許される波数は

\[
\mathbf{k}=\frac{\pi}{L}(n_x,n_y,n_z),\qquad n_x,n_y,n_z=1,2,3,\dots
\]

\begin{enumerate}
\def\labelenumi{\arabic{enumi}.}
\tightlist
\item
  波数の大きさが \(k\) 未満のモード総数 \(N(k)\) を、\(\mathbf{k}\)
  空間の「格子セル体積」と「正の八分球の体積」から求めよ。偏光自由度
  \(\times 2\) を忘れぬこと。
\item
  \(k=2\pi\nu/c\) を使って単位体積あたりのモード数を \(\nu\)
  の関数で表し、\(g(\nu)=dN/(L^3\,d\nu)=8\pi\nu^2/c^3\) を示せ。
\item
  係数 \(8\pi\)
  のうち、どの因子が「3次元」に、どの因子が「偏光」に由来するか説明せよ。
\end{enumerate}

\textbf{関連}：\hyperref[sec-classical-failure-modes]{§9.1
空洞モードの数え上げ}／本格導出は\href{../part5/14-cavity-modes.qmd}{第14章
空洞モード}。

\section{模範解答（問題
9-1）}\label{ux6a21ux7bc4ux89e3ux7b54ux554fux984c-9-1}

\textbf{(1)} 各 \((n_x,n_y,n_z)\)
は正の整数なので、許される波数ベクトルは \(\mathbf{k}\) 空間で間隔
\(\pi/L\)
の格子点を、第1八分象限（すべて正）に限って占める。1モードあたりの格子セル体積は
\((\pi/L)^3\)。大きさ \(k\) 未満の領域は半径 \(k\)
の球の正の八分の一、体積
\(\frac{1}{8}\cdot\frac{4}{3}\pi k^3\)。偏光2状態を掛けて

\[
N(k)=2\times\frac{\frac{1}{8}\cdot\frac{4}{3}\pi k^3}{(\pi/L)^3}=2\times\frac{\pi k^3 L^3}{6\pi^3}=\frac{L^3 k^3}{3\pi^2}.
\]

\textbf{(2)} 単位体積あたり \(n(k)=N/L^3=k^3/(3\pi^2)\)。\(k=2\pi\nu/c\)
を代入すると

\[
n(\nu)=\frac{1}{3\pi^2}\left(\frac{2\pi\nu}{c}\right)^3=\frac{8\pi^3\nu^3}{3\pi^2 c^3}=\frac{8\pi\nu^3}{3c^3}.
\]

\[
g(\nu)=\frac{dn}{d\nu}=\frac{8\pi\nu^2}{c^3}.
\]

\textbf{(3)} \(\nu^2\) は3次元波数空間の球殻 \(4\pi k^2\,dk\)（次元
\(d\) なら \(k^{d-1}\)）に由来する幾何因子。係数の
\(2\)（\(8\pi=2\times4\pi\) の \(2\)）は電磁波の横波2偏光に由来する。

\textbf{答え}：\(N(k)=L^3k^3/3\pi^2\)、よって
\(g(\nu)=8\pi\nu^2/c^3\)。\(\nu^2\)＝3次元の幾何、係数の
\(2\)＝偏光自由度。

\subsection*{問題 9-2　古典分配関数から 1モードの平均エネルギー
kBT}\label{ex-9-2}
\addcontentsline{toc}{subsection}{問題 9-2　古典分配関数から
1モードの平均エネルギー kBT}

{[}★ 難易度：☆ {]}

§9.2 は等分配定理で各モードに \(k_BT\)
が分配されると述べた。これを古典分配関数から直接導け。モードは調和振動子
\(E=\dfrac{p^2}{2m}+\dfrac{1}{2}m\omega^2 q^2\) と等価である。

\begin{enumerate}
\def\labelenumi{\arabic{enumi}.}
\tightlist
\item
  古典分配関数
  \(Z=\dfrac{1}{h_0}\displaystyle\int e^{-\beta E}\,dq\,dp\)（\(\beta=1/k_BT\)）を計算し、\(Z\propto\beta^{-1}\)
  を示せ。
\item
  \(\langle E\rangle=-\partial\ln Z/\partial\beta\) から
  \(\langle E\rangle=k_BT\) を導け。
\item
  この「2つの二次形式に各
  \(\frac12 k_BT\)」という結果を、3次元の自由粒子（並進運動エネルギーのみ）の
  \(\frac32 k_BT\) と対比して、なぜ振動子は \(k_BT\)
  なのかを一言で説明せよ。
\end{enumerate}

\textbf{関連}：\hyperref[sec-classical-failure-rayleigh-jeans]{§9.2
等分配と
Rayleigh--Jeans}／等分配定理は\hyperref[ux7b2c7ux7ae0-ux7d71ux8a08ux529bux5b66ux306eux57faux672c]{第7章
§7.6}。この古典結果を量子化したものが\hyperref[sec-planck-quantum-oscillator]{§10.2}の\hyperref[ex-10-1]{演習
10-1}。

\section{模範解答（問題
9-2）}\label{ux6a21ux7bc4ux89e3ux7b54ux554fux984c-9-2}

\textbf{(1)} 位置と運動量の積分が分離する：

\[
Z=\frac{1}{h_0}\int_{-\infty}^{\infty}e^{-\beta p^2/2m}dp\int_{-\infty}^{\infty}e^{-\beta m\omega^2 q^2/2}dq.
\]

ガウス積分 \(\int e^{-a x^2}dx=\sqrt{\pi/a}\) を用いると、各因子は
\(\sqrt{2\pi m/\beta}\) と \(\sqrt{2\pi/(\beta m\omega^2)}\)。積は

\[
Z=\frac{1}{h_0}\sqrt{\frac{2\pi m}{\beta}}\sqrt{\frac{2\pi}{\beta m\omega^2}}=\frac{2\pi}{h_0\omega}\cdot\frac{1}{\beta}\;\propto\;\beta^{-1}.
\]

\textbf{(2)} \(\ln Z=-\ln\beta+\text{const}\) なので

\[
\langle E\rangle=-\frac{\partial\ln Z}{\partial\beta}=\frac{1}{\beta}=k_BT.
\]

2つの二次形式（\(p^2\) と \(q^2\)）が各 \(\beta^{-1/2}\)
を生み、合わせて \(\beta^{-1}\) ＝ \(k_BT\)。

\textbf{(3)} 等分配は「ハミルトニアン中の二次形式1つあたり
\(\frac12 k_BT\)」を与える。自由粒子は運動量の二次形式が3つ（\(p_x^2,p_y^2,p_z^2\)）で
\(\frac32 k_BT\)。振動子は運動量＋位置で二次形式が2つなので
\(k_BT\)。位置に対する復元力（ポテンシャル項）があることが差を生む。

\textbf{答え}：\(Z\propto\beta^{-1}\Rightarrow\langle E\rangle=k_BT\)。二次形式2個（運動＋ポテンシャル）ゆえ自由粒子の
\(\frac32k_BT\) と異なる。

\subsection*{問題
9-3　紫外発散の「次数」とカットオフ依存性}\label{ex-9-3}
\addcontentsline{toc}{subsection}{問題
9-3　紫外発散の「次数」とカットオフ依存性}

{[}★ 難易度：☆☆ {]}

§9.3 は
\(u_\mathrm{classical}=\int_0^\infty u_\nu^\mathrm{RJ}d\nu=\infty\)
を示した。なぜ「紫外線」破綻と呼ぶのかを定量的に確かめる。

\begin{enumerate}
\def\labelenumi{\arabic{enumi}.}
\tightlist
\item
  人工的な上限 \(\nu_{\max}\) を導入し、全エネルギー密度が
  \(u(\nu_{\max})=\dfrac{8\pi k_BT}{3c^3}\nu_{\max}^3\)
  となることを示せ。発散は \(\nu_{\max}\to\infty\)
  でどんな冪で増えるか。
\item
  エネルギーがどの振動数帯に集中するかを「\(d u/d\nu\propto\nu^2\)
  が単調増加」という事実から論じ、なぜ高振動数（紫外）側が破綻の原因かを述べよ。
\item
  プランク分布 \(u_\nu\propto\nu^3/(e^{h\nu/k_BT}-1)\)
  ではこの積分が収束する理由を、被積分関数の \(\nu\to\infty\)
  での振る舞いから説明せよ。
\end{enumerate}

\textbf{関連}：\hyperref[sec-classical-failure-uv-catastrophe]{§9.3
紫外線破綻}／収束する全積分（Stefan--Boltzmann）は\hyperref[ex-10-3]{演習
10-3}、本文\hyperref[sec-planck-quantum-derivation]{§10.3}。

\section{模範解答（問題
9-3）}\label{ux6a21ux7bc4ux89e3ux7b54ux554fux984c-9-3}

\textbf{(1)} \(u_\nu^\mathrm{RJ}=\dfrac{8\pi\nu^2}{c^3}k_BT\) を \(0\)
から \(\nu_{\max}\) まで積分：

\[
u(\nu_{\max})=\frac{8\pi k_BT}{c^3}\int_0^{\nu_{\max}}\nu^2 d\nu=\frac{8\pi k_BT}{3c^3}\nu_{\max}^3.
\]

発散は \(\nu_{\max}^3\)
の3次冪。カットオフを2倍にすると寄与が8倍になる。

\textbf{(2)} エネルギー密度の振動数あたり寄与 \(du/d\nu\propto\nu^2\) は
\(\nu\)
とともに単調増加し、上限近傍が積分を支配する。すなわちエネルギーは常に「いちばん高い振動数帯」に集中する。古典論には上限がないので、紫外〜X線〜ガンマ線側がいくらでも寄与を積み増し、発散する。これが「紫外線」破綻の名の由来。

\textbf{(3)} プランクでは被積分関数が \(\nu\to\infty\) で

\[
\frac{\nu^3}{e^{h\nu/k_BT}-1}\;\sim\;\nu^3 e^{-h\nu/k_BT}\to 0
\]

と指数的に抑えられる。指数減衰は任意の冪 \(\nu^3\)
に勝つので積分は収束し、有限の \(u=aT^4\)
を与える（\hyperref[ex-10-3]{演習 10-3}）。量子化によるエネルギー単位
\(h\nu\)
が高振動数モードを「凍結」させることが、収束の物理的理由である。

\textbf{答え}：\(u\propto\nu_{\max}^3\)
で発散、寄与は高振動数に集中（＝紫外破綻）。プランクは
\(\nu^3e^{-h\nu/k_BT}\to0\) で収束。

\subsection*{問題 9-4　破綻の「責任の所在」―
モード密度か等分配か}\label{ex-9-4}
\addcontentsline{toc}{subsection}{問題 9-4　破綻の「責任の所在」―
モード密度か等分配か}

{[}★ 難易度：☆☆ {]}

§9.4--9.5 の論理を整理する問いである。

\begin{enumerate}
\def\labelenumi{\arabic{enumi}.}
\tightlist
\item
  Rayleigh--Jeans 則 \(u_\nu^\mathrm{RJ}=g(\nu)\langle E\rangle\)
  を構成する2因子のうち、量子論で「変えなくてよい」のはどちらで、「変えなければならない」のはどちらか。理由とともに述べよ。
\item
  観測される黒体スペクトル（第1章）が古典論を致命的に否定する2つの特徴を挙げよ。
\item
  もしエネルギー単位が \(\epsilon=h\nu\)（\(\propto\nu\)）でなく \(\nu\)
  に依らない定数 \(\epsilon_0\)
  だったら、高振動数モードの凍結は起きるか。簡潔に論ぜよ。
\end{enumerate}

\textbf{関連}：\hyperref[sec-classical-failure-missing]{§9.4
等分配の前提}、\hyperref[sec-classical-failure-necessity]{§9.5
量子仮説の必然性}／観測の特徴は\hyperref[ux7b2c1ux7ae0-ux5b87ux5b99ux3067ux9ed2ux4f53ux653eux5c04ux306fux3069ux3053ux306bux73feux308cux308bux304b]{第1章}。

\section{模範解答（問題
9-4）}\label{ux6a21ux7bc4ux89e3ux7b54ux554fux984c-9-4}

\textbf{(1)} \textbf{モード密度 \(g(\nu)=8\pi\nu^2/c^3\)
は変えなくてよい}。これは古典電磁気＋境界条件の幾何的結果で、量子論を要しない（実際プランク分布の前因子
\(\nu^2\)/\(\nu^3\)
はこれが起源）。\textbf{変えるべきは各モードの平均エネルギー
\(\langle E\rangle\)}。等分配の
\(k_BT\)（エネルギー連続の仮定に依存）を、量子化された
\(h\nu/(e^{h\nu/k_BT}-1)\) に置き換える必要がある。

\textbf{(2)} ①
スペクトルに有限のピークが存在する（古典は単調増加でピークなし）。②
高振動数側が指数的に減衰する（古典は \(\nu^2\) で増大）。CMB
の完全な黒体形や恒星スペクトルがこれを示す。

\textbf{(3)} 起きない。凍結は「エネルギー単位 \(\epsilon(\nu)\)
が高振動数で \(k_BT\) を超える」ことで生じる。\(\epsilon_0\)
が定数なら、\(\epsilon_0\ll k_BT\) を満たす限りすべての振動数で等分配
\(k_BT\) が回復し、\(\nu\to\infty\)
でもモードが凍結せず発散は残る。\(\epsilon\propto\nu\)
という\textbf{振動数に比例して増える単位}であることが、高振動数を選択的に止める鍵である。

\textbf{答え}：\(g(\nu)\) は不変、\(\langle E\rangle\)
を量子化。否定する観測＝ピークの存在と指数減衰。\(\epsilon\propto\nu\)
でないと凍結は起きない。

\section*{さらに学ぶための参考文献}\label{sec-classical-failure-further}
\addcontentsline{toc}{section}{さらに学ぶための参考文献}

\markright{さらに学ぶための参考文献}

\begin{itemize}
\tightlist
\item
  Pais, \emph{Subtle is the Lord} (Oxford, 1982) --- Planck・Einstein
  の量子論誕生の歴史的精査
\item
  Rybicki \& Lightman, \emph{Radiative Processes in Astrophysics}
  (Wiley, 1979) --- §1.5「Blackbody Radiation」
\item
  Kuhn, \emph{Black-Body Theory and the Quantum Discontinuity,
  1894-1912} (Oxford, 1978) --- 黒体放射と量子論誕生の科学史
\end{itemize}

\chapter{第10章　プランクの量子仮説}\label{ux7b2c10ux7ae0-ux30d7ux30e9ux30f3ux30afux306eux91cfux5b50ux4eeeux8aac}

\textbf{この章の主題}：第9章で見たように、古典論で黒体スペクトルを説明しようとすると紫外線破綻を起こす。Planck
(1900) はこの破綻を救うために、\textbf{振動子のエネルギーが \(h\nu\)
の整数倍にしか取れない}という思い切った仮定を立てた。本章では、この仮定
― \textbf{エネルギー量子化}（\textbf{energy quantization}）―
がプランク分布を必然的に与えることを示す。第8章で Bose
統計から導いた占有因子と、ここで導く Planck
のオリジナルな計算は、別の出発点から同じゴールに到達する。Planck
の歴史的歩みと、現代的解釈の両方を示す。

\section*{この章の中心地図}\label{sec-planck-quantum-map}
\addcontentsline{toc}{section}{この章の中心地図}

\markright{この章の中心地図}

\textbf{中心地図：プランク関数 ― 本部で扱う因子}

\[
B_\nu(T) = \frac{2\, \boxed{h}\, \nu^3}{c^2} \cdot \frac{1}{e^{\boxed{h}\nu/kT} - 1}
\]

強調されている \textbf{Planck 定数
\(h\)}（前因子の中と、指数の中の二箇所）が、第IV部で逆引きする対象である。エネルギーが連続でなく
\(h\nu\) という離散単位でやりとりされる ―
これが量子論の核心で、第10章（Planck の量子仮説）と第11章（Einstein
の方法）でその必然性を二つの異なるルートから示す。

\begin{tcolorbox}[enhanced jigsaw, arc=.35mm, bottomrule=.15mm, bottomtitle=1mm, breakable, colback=white, colbacktitle=quarto-callout-note-color!10!white, colframe=quarto-callout-note-color-frame, coltitle=black, left=2mm, leftrule=.75mm, opacityback=0, opacitybacktitle=0.6, rightrule=.15mm, title=\textcolor{quarto-callout-note-color}{\faInfo}\hspace{0.5em}{ノート}, titlerule=0mm, toprule=.15mm, toptitle=1mm]

\textbf{方針}：本章で扱う \textbf{Planck 定数
\(h\)}（中心地図の強調因子）は、本書全体で最も重要な量子論的定数である。本章は「\(h\)
を導入する」と「\(h\) がスペクトルの形に責任を持つ」ことを示す章である。

\end{tcolorbox}

\section*{この章で答える問い}\label{sec-planck-quantum-questions}
\addcontentsline{toc}{section}{この章で答える問い}

\markright{この章で答える問い}

\begin{itemize}
\tightlist
\item
  プランクは何を「やむを得ず」仮定したのか
\item
  なぜ \(h\nu\) という離散単位が必要だったのか
\item
  プランクの量子仮説は、観測スペクトルのどの部分を救ったのか
\item
  振動子のエネルギー量子化と、光子の存在は同じことを言っているのか
\end{itemize}

\section*{到達目標}\label{sec-planck-quantum-goals}
\addcontentsline{toc}{section}{到達目標}

\markright{到達目標}

この章を読み終えると、読者は次のことができるようになる：

\begin{itemize}
\tightlist
\item
  プランクの量子仮説を出発点として、プランク分布を導ける
\item
  \(h\nu\) が高振動数側の指数関数的減衰を生む仕組みを説明できる
\item
  プランクの歴史的試行錯誤と、現代的解釈（光子の存在）の関係を整理できる
\end{itemize}

\begin{center}\rule{0.5\linewidth}{0.5pt}\end{center}

\section{10.1 エネルギー量子 ― Planck
の大胆な仮定}\label{sec-planck-quantum-energy-quantum}

{[}本文目安：B3{]}

第9章で見た古典の破綻を救うには、エネルギーの離散単位 \(\epsilon(\nu)\)
を導入する必要があった。Planck (1900) はこの単位として

\begin{equation}\protect\phantomsection\label{eq-planck-quantum-epsilon}{
\epsilon = h\nu
}\end{equation}

を仮定した。\(h\) は新しい普遍定数で、後に \textbf{Planck
定数}（\textbf{Planck constant}）と呼ばれる。現在の値は

\[
h = 6.626 \times 10^{-27}\ \mathrm{erg \cdot s} = 6.626 \times 10^{-34}\ \mathrm{J \cdot s}
\]

である。Planck
はこの値を、彼の式が黒体スペクトルの観測データを再現するように決定した。

\textbf{仮定の内容}を正確に述べると次のようになる：

\begin{quote}
振動数 \(\nu\)
で振動する電磁モード（あるいは空洞壁の振動子）のエネルギーは、\(0, h\nu, 2h\nu, 3h\nu, \ldots\)
という離散的な値しかとらない。
\end{quote}

これは古典物理から見ると不自然な仮定である。古典では振動子のエネルギーは任意の連続値をとれるはずだから。しかしこの仮定なしには、観測される黒体スペクトルの高振動数側の減衰は出てこない。

\begin{tcolorbox}[enhanced jigsaw, arc=.35mm, bottomrule=.15mm, bottomtitle=1mm, breakable, colback=white, colbacktitle=quarto-callout-note-color!10!white, colframe=quarto-callout-note-color-frame, coltitle=black, left=2mm, leftrule=.75mm, opacityback=0, opacitybacktitle=0.6, rightrule=.15mm, title=\textcolor{quarto-callout-note-color}{\faInfo}\hspace{0.5em}{ノート}, titlerule=0mm, toprule=.15mm, toptitle=1mm]

\textbf{歴史}：Planck
自身は当初、この量子仮説を\textbf{計算上の便法}として導入した。「振動子のエネルギーが本質的に離散である」という物理的解釈には、彼自身ためらいがあった。エネルギー量子の物理的実在性を主張したのは、後の
Einstein (1905) ― 光電効果の論文 ― である。第11章で扱う。

\end{tcolorbox}

\section{10.2 振動子の量子化 ―
統計力学の組み直し}\label{sec-planck-quantum-oscillator}

{[}本文目安：B3{]}

エネルギー量子 式~\ref{eq-planck-quantum-epsilon} を仮定すると、第7章
§7.2 の手順で振動子の分配関数を組み直せる。エネルギー準位が
\(E_n = n h\nu\)（\(n = 0, 1, 2, \ldots\)）で離散化された場合、分配関数は

\begin{equation}\protect\phantomsection\label{eq-partition-quantized-oscillator}{
Z_1(\beta) = \sum_{n=0}^{\infty} e^{-\beta n h\nu} = \frac{1}{1 - e^{-\beta h\nu}}
}\end{equation}

ここで
\(\beta = 1/k_B T\)。これは無限等比級数の和である（古典の連続積分なら
\(\int_0^\infty e^{-\beta E} dE = 1/\beta = k_B T\) が出る）。

分配関数の対数から平均エネルギーが求まる（第7章 §7.2 の手法）：

\begin{equation}\protect\phantomsection\label{eq-mean-energy-quantized}{
\langle E \rangle = -\frac{\partial \ln Z_1}{\partial \beta} = \frac{h\nu}{e^{h\nu/k_B T} - 1}
}\end{equation}

これが\textbf{量子化された振動子のひとつあたりの平均エネルギー}である。第7章
§7.2 の 式~\ref{eq-mean-energy-oscillator} で既に見た形だが、ここでは
Planck の量子仮説から直接出ている。

式~\ref{eq-mean-energy-quantized} を古典の等分配
\(\langle E \rangle_\text{classical} = k_B T\) と比較すると：

\begin{itemize}
\tightlist
\item
  低振動数極限（\(h\nu \ll k_B T\)）：\(\langle E \rangle \to k_B T\) ―
  \textbf{古典と一致}
\item
  高振動数極限（\(h\nu \gg k_B T\)）：\(\langle E \rangle \to h\nu \, e^{-h\nu/k_B T} \to 0\)
  ― \textbf{古典より遥かに小さい}
\end{itemize}

\begin{figure}[H]

\centering{

\includegraphics[width=0.88\linewidth,height=\textheight,keepaspectratio]{index_files/mediabag/part4/../images/quantum-freeze-out.pdf}

}

\caption{\label{fig-quantum-freezeout-ch10}量子化振動子 1
個の平均エネルギーと古典の等分配の比較（第7章 §7.2
で初出した図を、Planck
量子仮説の文脈で再掲）。\textbf{高振動数モードが「凍結」する仕組み}
こそ、紫外破綻（第9章 §9.3）を回避してプランク分布を生む鍵である。}

\end{figure}%

これが\textbf{高振動数モードの「凍結」}である。エネルギー単位 \(h\nu\)
が熱エネルギー \(k_B T\)
より遥かに大きい場合、そのモードを励起するための熱が手に入らず、平均エネルギーが指数関数的に小さくなる。

\begin{tcolorbox}[enhanced jigsaw, arc=.35mm, bottomrule=.15mm, breakable, colback=white, colframe=quarto-callout-tip-color-frame, left=2mm, leftrule=.75mm, opacityback=0, rightrule=.15mm, toprule=.15mm]
\begin{minipage}[t]{5.5mm}
\textcolor{quarto-callout-tip-color}{\faLightbulb}
\end{minipage}%
\begin{minipage}[t]{\textwidth - 5.5mm}

\textbf{問い}：なぜ「凍結」と呼ぶのか？

\textbf{短答}：高振動数モードに熱的にアクセスできない ―
つまり熱エネルギーがそのモードに「行けない」 ―
状態を、低温物理の言葉で「凍結」（freeze-out）と呼ぶ。

\textbf{もう一歩}：同じ概念は、二原子分子の回転・振動自由度の低温での凍結、固体の比熱が低温で
\(T^3\) に従う Debye
則、宇宙論的にニュートリノが脱結合する状況など、さまざまな場面で再登場する。Planck
の議論は、その最初の例である。

\end{minipage}%
\end{tcolorbox}

\section{10.3 プランク分布の導出 ― 量子振動子 ×
モード密度}\label{sec-planck-quantum-derivation}

{[}本文目安：B3-B4｜導出：M1{]}

式~\ref{eq-mean-energy-quantized} と、第9章
式~\ref{eq-mode-density-borrowed} のモード密度
\(g(\nu) = 8\pi\nu^2/c^3\)
を組み合わせると、放射場の振動数あたりエネルギー密度は

\begin{equation}\protect\phantomsection\label{eq-planck-energy-density-quantum}{
u_\nu(T) = g(\nu) \cdot \langle E \rangle = \frac{8\pi h \nu^3}{c^3} \cdot \frac{1}{e^{h\nu/k_B T} - 1}
}\end{equation}

比強度に直すと

\begin{equation}\protect\phantomsection\label{eq-planck-quantum-rederived}{
B_\nu(T) = \frac{c}{4\pi} u_\nu = \frac{2 h \nu^3}{c^2} \cdot \frac{1}{e^{h\nu/k_B T} - 1}
}\end{equation}

\textbf{プランク分布が、振動子の量子化から導かれた}。これが Planck
(1900) の歴史的論文の核心である。

第8章で Bose
統計から導いた占有因子と、本章でエネルギー量子化から導いた占有因子は、\textbf{全く同じ形}になる：

{\def\LTcaptype{none} % do not increment counter
\begin{longtable}[]{@{}
  >{\raggedright\arraybackslash}p{(\linewidth - 2\tabcolsep) * \real{0.5000}}
  >{\raggedright\arraybackslash}p{(\linewidth - 2\tabcolsep) * \real{0.5000}}@{}}
\toprule\noalign{}
\begin{minipage}[b]{\linewidth}\raggedright
出発点
\end{minipage} & \begin{minipage}[b]{\linewidth}\raggedright
占有因子
\end{minipage} \\
\midrule\noalign{}
\endhead
\bottomrule\noalign{}
\endlastfoot
第8章：光子の Bose 統計（\(\mu = 0\)） &
\(\bar{n} = 1/(e^{h\nu/k_B T} - 1)\) \\
第10章：振動子のエネルギー量子化 &
\(\langle E \rangle / h\nu = 1/(e^{h\nu/k_B T} - 1)\) \\
\end{longtable}
}

これは偶然ではない。両者は同じ物理（電磁場の各モードを量子化された調和振動子と見る）の二つの言い方である。第25章で「電磁場の量子化」として整理し直す。

\begin{tcolorbox}[enhanced jigsaw, arc=.35mm, bottomrule=.15mm, bottomtitle=1mm, breakable, colback=white, colbacktitle=quarto-callout-note-color!10!white, colframe=quarto-callout-note-color-frame, coltitle=black, left=2mm, leftrule=.75mm, opacityback=0, opacitybacktitle=0.6, rightrule=.15mm, title=\textcolor{quarto-callout-note-color}{\faInfo}\hspace{0.5em}{ノート}, titlerule=0mm, toprule=.15mm, toptitle=1mm]

\textbf{観測との対応}：高振動数で
\(1/(e^{h\nu/k_B T}-1) \sim e^{-h\nu/k_B T}\)
が出ることが、観測される黒体スペクトルの\textbf{指数減衰}（第1章で見た）の起源である。これが
Planck の量子仮説が「観測の必然」として要請された理由である。

\end{tcolorbox}

\section{10.4 Wien・Rayleigh--Jeans
両極限の再導出}\label{sec-planck-quantum-limits}

{[}本文目安：B3-B4｜導出：M1{]}

式~\ref{eq-planck-quantum-rederived} から、第2章 §2.5
で見た両極限がきれいに出る。

\textbf{低振動数極限}（\(h\nu \ll k_B T\)）：分母
\(e^{h\nu/k_B T} - 1 \simeq h\nu/k_B T\) なので

\begin{equation}\protect\phantomsection\label{eq-rj-from-planck}{
B_\nu \to \frac{2\nu^2}{c^2} k_B T
}\end{equation}

これは第9章 式~\ref{eq-rayleigh-jeans-bnu} の Rayleigh--Jeans
則と一致する。\textbf{\(h\)
が消える}ことに注意。この極限では古典統計が正しい。

\textbf{高振動数極限}（\(h\nu \gg k_B T\)）：分母
\(e^{h\nu/k_B T} - 1 \simeq e^{h\nu/k_B T}\) なので

\begin{equation}\protect\phantomsection\label{eq-wien-from-planck}{
B_\nu \to \frac{2h\nu^3}{c^2} e^{-h\nu/k_B T}
}\end{equation}

これが \textbf{Wien の式}である（\textbf{Wien's distribution
law}）。\textbf{\(h\) が指数の中と前因子の両方に現れる}ことに注意。

両極限の境目は \(h\nu \sim k_B T\)、すなわち
\(\nu \sim k_B T/h \simeq 2 \times 10^{10} (T/\text{K})\) Hz
である。太陽光（\(T \sim 5800\) K）なら境界は近赤外、CMB（2.7
K）なら境界はマイクロ波領域。

\begin{tcolorbox}[enhanced jigsaw, arc=.35mm, bottomrule=.15mm, breakable, colback=white, colframe=quarto-callout-tip-color-frame, left=2mm, leftrule=.75mm, opacityback=0, rightrule=.15mm, toprule=.15mm]
\begin{minipage}[t]{5.5mm}
\textcolor{quarto-callout-tip-color}{\faLightbulb}
\end{minipage}%
\begin{minipage}[t]{\textwidth - 5.5mm}

\textbf{問い}：なぜ Wien は Planck より先に Wien
の式（経験式として）を発見できたのか？

\textbf{短答}：Wien (1896) は熱力学と次元解析だけで
\(B_\nu = A\nu^3 f(\nu/T)\) という形（\textbf{Wien displacement
function}）まで絞り、観測データに当てはめて \(f(x) = e^{-\beta x}\)
という近似を得た。これが高振動数側に正しい形だが、低振動数では
Rayleigh--Jeans の方が観測と合うことが後に判明した。

\textbf{もう一歩}：Planck の真の貢献は、Wien 式（高振動数で正しい）と
Rayleigh--Jeans
式（低振動数で正しい）の\textbf{両方を補間する単一の式}を、量子仮説という物理的根拠つきで提示したこと。経験式の補間ではなく、物理的に正しい全振動数スペクトルを与えた。

\end{minipage}%
\end{tcolorbox}

\section{10.5 歴史的背景と現代的解釈}\label{sec-planck-quantum-history}

{[}本文目安：B3{]}

Planck の量子仮説の歴史的・物理的位置づけを整理する。

\textbf{Planck (1900)} ―
黒体スペクトルを再現するために、空洞壁の物質振動子のエネルギーが
\(h\nu\) の整数倍に限られると仮定。Planck
自身は当初これを\textbf{計算上の便法}と見なし、エネルギーが本質的に離散である物理的解釈は留保した。

\textbf{Einstein (1905)} ―
光電効果の解析から、電磁場そのものが量子化された光（\textbf{光量子},
light quanta）から成ることを主張。Planck
の振動子の量子化と物理的に等価だが、解釈の強度がより強い。第11章で扱う。

\textbf{現代的解釈（場の量子論）}：電磁場の各モードを\textbf{量子化された調和振動子}として扱う（第VIII部・第25章）。各モードのエネルギー準位が
\(E_n = (n + 1/2) h\nu\)（零点エネルギー込み）となり、\(n\)
が「そのモードに含まれる光子の数」を表す。Planck の振動子と Einstein
の光量子は、この描像の中で統一される。

\textbf{観測の意義}：Planck (1900)
の論文は、20世紀物理学の出発点と広く認められている。一つの\textbf{観測}（黒体スペクトル）から、自然界に普遍的な定数
\(h\)
が導入され、量子力学・場の量子論が立ち上がる流れの起点となった。本書の三原則「①
逆引き」の象徴的な瞬間とも言える ―
観測が理論を駆動した、その典型例である。

\begin{tcolorbox}[enhanced jigsaw, arc=.35mm, bottomrule=.15mm, bottomtitle=1mm, breakable, colback=white, colbacktitle=quarto-callout-note-color!10!white, colframe=quarto-callout-note-color-frame, coltitle=black, left=2mm, leftrule=.75mm, opacityback=0, opacitybacktitle=0.6, rightrule=.15mm, title=\textcolor{quarto-callout-note-color}{\faInfo}\hspace{0.5em}{ノート}, titlerule=0mm, toprule=.15mm, toptitle=1mm]

\textbf{使えるようになった道具}（§10.1〜10.5）：

\begin{itemize}
\tightlist
\item
  エネルギー量子 \(\epsilon = h\nu\)
  という仮定から、量子化振動子の平均エネルギー
  式~\ref{eq-mean-energy-quantized} が出る
\item
  これにモード密度を組み合わせると、観測されるプランク分布全体形が再構成される
\item
  Wien・Rayleigh--Jeans の両極限が一つの式から自然に出る ―
  観測される異なる振動数帯のスペクトルが、単一の量子仮説で説明される
\end{itemize}

\end{tcolorbox}

\begin{center}\rule{0.5\linewidth}{0.5pt}\end{center}

\section*{この章で何がわかったか}\label{sec-planck-quantum-summary}
\addcontentsline{toc}{section}{この章で何がわかったか}

\markright{この章で何がわかったか}

\textbf{中心地図に戻る}

本章で、中心地図プランク関数の \textbf{指数の中の \(h\)}（\(h\nu/k_B T\)
の項）と、\textbf{前因子の \(h\)}（\(2h\nu^3/c^2\)）の両方が、Planck
の量子仮説 \(\epsilon = h\nu\)
から必然として出ることを示した。これは中心地図の指数関数因子の量子論的起源を明らかにする。

ただし、Planck
の議論は「空洞壁の物質振動子のエネルギーが量子化される」という形であり、\textbf{光そのものが量子化されているかどうか}は留保されたままだった。これを明確にするのが、第11章で扱う
Einstein の方法である。

\textbf{次章へ}：第11章では Einstein (1905, 1916)
の議論をたどる。光量子の導入と、自発放出・誘導放出・誘導吸収による詳細釣り合いから、プランク分布を再度導出する。Einstein
係数は線スペクトル（第VII部）の中心概念にもなり、本書の二つの地図（プランク関数地図と水素線地図）をつなぐ要となる。

\section*{演習問題}\label{sec-planck-quantum-exercises}
\addcontentsline{toc}{section}{演習問題}

\markright{演習問題}

\textbf{模範解答}
は折りたたみを展開して確認できる（オンライン版）。まず自力で解いてから開くこと。

\subsection*{問題
10-1　量子化振動子の分配関数と平均エネルギー}\label{ex-10-1}
\addcontentsline{toc}{subsection}{問題
10-1　量子化振動子の分配関数と平均エネルギー}

{[}★ 難易度：☆ {]}

§10.2 は分配関数 \(Z_1=1/(1-e^{-\beta h\nu})\) と平均エネルギー
\(\langle E\rangle=h\nu/(e^{h\nu/k_BT}-1)\) を「第7章 §7.2
の手法で」与えた。両式を自力で導け。準位は
\(E_n=nh\nu\ (n=0,1,2,\ldots)\)。

\begin{enumerate}
\def\labelenumi{\arabic{enumi}.}
\tightlist
\item
  等比級数の和として
  \(Z_1=\sum_{n=0}^\infty e^{-\beta nh\nu}=\dfrac{1}{1-e^{-\beta h\nu}}\)
  を示せ（収束条件も述べよ）。
\item
  \(\langle E\rangle=-\partial\ln Z_1/\partial\beta\) を実行して
  \(\langle E\rangle=\dfrac{h\nu}{e^{\beta h\nu}-1}\) を導け。
\item
  古典結果 \(\langle E\rangle=k_BT\)
  と比べ、低振動数極限で一致することを確認せよ。
\end{enumerate}

\textbf{関連}：\hyperref[sec-planck-quantum-oscillator]{§10.2
振動子の量子化}／手法は\hyperref[ux7b2c7ux7ae0-ux7d71ux8a08ux529bux5b66ux306eux57faux672c]{第7章
統計力学}、古典との対比は\hyperref[ex-9-2]{演習 9-2}。

\section{模範解答（問題
10-1）}\label{ux6a21ux7bc4ux89e3ux7b54ux554fux984c-10-1}

\textbf{(1)} 公比 \(r=e^{-\beta h\nu}\) の等比級数。\(\beta,h,\nu>0\)
より \(0<r<1\) で収束し、

\[
Z_1=\sum_{n=0}^\infty r^n=\frac{1}{1-r}=\frac{1}{1-e^{-\beta h\nu}}.
\]

\textbf{(2)} \(\ln Z_1=-\ln(1-e^{-\beta h\nu})\)。\(\beta\) で微分：

\[
\langle E\rangle=-\frac{\partial\ln Z_1}{\partial\beta}=\frac{h\nu\,e^{-\beta h\nu}}{1-e^{-\beta h\nu}}=\frac{h\nu}{e^{\beta h\nu}-1}=\frac{h\nu}{e^{h\nu/k_BT}-1}.
\]

\textbf{(3)} 低振動数 \(x\equiv h\nu/k_BT\ll1\) で \(e^x-1\simeq x\)
ゆえ \(\langle E\rangle\to h\nu/x=k_BT\)。古典の等分配と一致する。

\textbf{答え}：\(Z_1=1/(1-e^{-\beta h\nu})\)、\(\langle E\rangle=h\nu/(e^{h\nu/k_BT}-1)\overset{h\nu\ll k_BT}{\longrightarrow}k_BT\)。

\subsection*{問題 10-2　両極限と「最初の量子補正」}\label{ex-10-2}
\addcontentsline{toc}{subsection}{問題 10-2　両極限と「最初の量子補正」}

{[}★ 難易度：☆☆ {]}

§10.4
はプランク分布の低・高振動数極限を示した。導出を埋め、さらに一歩進める。\(x\equiv h\nu/k_BT\)
とおく。

\begin{enumerate}
\def\labelenumi{\arabic{enumi}.}
\tightlist
\item
  低振動数（\(x\ll1\)）で
  \(B_\nu\to\dfrac{2\nu^2}{c^2}k_BT\)（Rayleigh--Jeans）を導け。
\item
  平均エネルギー \(\langle E\rangle=\dfrac{h\nu}{e^x-1}\) を \(x\)
  の小ささで展開し、古典値 \(k_BT\) からの \textbf{最初の補正項}
  を求めよ。補正の符号は何を意味するか。
\item
  高振動数（\(x\gg1\)）で Wien の式
  \(B_\nu\to\dfrac{2h\nu^3}{c^2}e^{-h\nu/k_BT}\) を導け。
\end{enumerate}

\textbf{関連}：\hyperref[sec-planck-quantum-limits]{§10.4
Wien・Rayleigh--Jeans 両極限}／RJ
の古典導出は\hyperref[sec-classical-failure-rayleigh-jeans]{§9.2}、観測比較は\hyperref[sec-reading-spectra-limits]{§2.5}。

\section{模範解答（問題
10-2）}\label{ux6a21ux7bc4ux89e3ux7b54ux554fux984c-10-2}

\textbf{(1)} \(B_\nu=\dfrac{2h\nu^3}{c^2}\dfrac{1}{e^x-1}\)。\(x\ll1\)
で \(e^x-1\simeq x=h\nu/k_BT\) ゆえ
\(B_\nu\to\dfrac{2h\nu^3}{c^2}\cdot\dfrac{k_BT}{h\nu}=\dfrac{2\nu^2}{c^2}k_BT\)。\(h\)
が消える＝古典極限。

\textbf{(2)}
\(e^x-1=x\left(1+\dfrac{x}{2}+\dfrac{x^2}{6}+\cdots\right)\) ゆえ

\[
\langle E\rangle=\frac{h\nu}{e^x-1}=\frac{k_BT}{1+\frac{x}{2}+\cdots}=k_BT\left(1-\frac{x}{2}+\cdots\right)=k_BT-\frac{h\nu}{2}+\cdots.
\]

最初の補正は
\(-h\nu/2\)。\textbf{負号}は、量子化により各モードが古典の等分配値より小さいエネルギーしか受け取れない（高振動数ほど効く凍結の始まり）ことを意味する。

\textbf{(3)} \(x\gg1\) で \(e^x-1\simeq e^x\) ゆえ
\(B_\nu\to\dfrac{2h\nu^3}{c^2}e^{-x}=\dfrac{2h\nu^3}{c^2}e^{-h\nu/k_BT}\)。これが
Wien の式。

\textbf{答え}：低：\(2\nu^2k_BT/c^2\)；補正
\(\langle E\rangle\simeq k_BT-\tfrac12 h\nu\)（負＝凍結の兆し）；高：\(2h\nu^3e^{-h\nu/k_BT}/c^2\)。

\subsection*{問題 10-3　プランク分布から Stefan--Boltzmann
則}\label{ex-10-3}
\addcontentsline{toc}{subsection}{問題 10-3　プランク分布から
Stefan--Boltzmann 則}

{[}★ 難易度：☆☆☆ {]}

第9章の表で「全エネルギー密度＝\(aT^4\)（有限）」とあったが、本文は積分を実行していない。プランクのエネルギー密度
\(u_\nu=\dfrac{8\pi h\nu^3}{c^3}\dfrac{1}{e^{h\nu/k_BT}-1}\)
を全振動数で積分し、\(u=aT^4\) を導いて放射定数 \(a\) を基本定数で表せ。

\begin{enumerate}
\def\labelenumi{\arabic{enumi}.}
\tightlist
\item
  \(x=h\nu/k_BT\)
  で変数変換し、\(u=\dfrac{8\pi h}{c^3}\left(\dfrac{k_BT}{h}\right)^4\displaystyle\int_0^\infty\frac{x^3}{e^x-1}dx\)
  を導け。
\item
  \(\displaystyle\int_0^\infty\frac{x^3}{e^x-1}dx=\frac{\pi^4}{15}\)
  を用いて \(a=\dfrac{8\pi^5 k_B^4}{15c^3h^3}\)
  を示し、数値（SI）を概算せよ。
\item
  \(u=aT^4\) と射出強度 \(j=\sigma T^4\) の関係 \(a=4\sigma/c\)
  から、\(\sigma\) の値を確認せよ。
\end{enumerate}

\textbf{関連}：\hyperref[sec-planck-quantum-derivation]{§10.3
プランク分布の導出}／発散する古典版との対比は\hyperref[ex-9-3]{演習
9-3}、積分公式は\href{../appendix/A-math-notes.qmd}{付録A 数学ノート}。

\section{模範解答（問題
10-3）}\label{ux6a21ux7bc4ux89e3ux7b54ux554fux984c-10-3}

\textbf{(1)} \(x=h\nu/k_BT\) より
\(\nu=k_BTx/h,\ d\nu=(k_BT/h)dx\)。代入して

\[
u=\frac{8\pi h}{c^3}\left(\frac{k_BT}{h}\right)^4\int_0^\infty\frac{x^3}{e^x-1}dx.
\]

\textbf{(2)} 積分 \(=\pi^4/15\) を入れて

\[
u=\frac{8\pi^5 k_B^4}{15c^3h^3}T^4\equiv aT^4,\qquad a=\frac{8\pi^5k_B^4}{15c^3h^3}\approx7.57\times10^{-16}\ \mathrm{J\,m^{-3}\,K^{-4}}.
\]

\textbf{(3)} \(a=4\sigma/c\) より
\(\sigma=ac/4=7.57\times10^{-16}\times3.00\times10^8/4\approx5.67\times10^{-8}\ \mathrm{W\,m^{-2}\,K^{-4}}\)。既知の
Stefan--Boltzmann 定数と一致する。

\textbf{答え}：\(a=\dfrac{8\pi^5k_B^4}{15c^3h^3}\approx7.57\times10^{-16}\,\mathrm{J\,m^{-3}K^{-4}}\)、\(\sigma=ac/4\approx5.67\times10^{-8}\,\mathrm{W\,m^{-2}K^{-4}}\)。

\subsection*{問題 10-4　Wien 変位則：ピーク振動数は T
に比例}\label{ex-10-4}
\addcontentsline{toc}{subsection}{問題 10-4　Wien 変位則：ピーク振動数は
T に比例}

{[}★ 難易度：☆☆ {]}

§10.4 は両極限の境目を \(\nu\sim k_BT/h\)
と見積もった。プランク分布のピークを正確に求めて「変位則」を導く。

\begin{enumerate}
\def\labelenumi{\arabic{enumi}.}
\tightlist
\item
  \(B_\nu\propto\dfrac{\nu^3}{e^{x}-1}\ (x=h\nu/k_BT)\) を \(\nu\)
  で微分してゼロにおき、\(3(1-e^{-x})=x\) という超越方程式を導け。
\item
  数値解 \(x_\mathrm{peak}\approx2.821\) を用いて
  \(\nu_\mathrm{peak}=2.821\,k_BT/h\)
  を示し、\(\nu_\mathrm{peak}\propto T\) を確認せよ。
\item
  太陽 \(T\simeq5800\) K について \(\nu_\mathrm{peak}\)
  と対応波長を概算せよ（\(B_\lambda\)
  のピーク波長とは係数が異なることに触れよ）。
\end{enumerate}

\textbf{関連}：\hyperref[sec-planck-quantum-limits]{§10.4
両極限と境界}／波長表示の Wien 変位則の導出は\hyperref[ex-2-2]{演習
2-2}、観測スペクトルのピークは\hyperref[ux7b2c1ux7ae0-ux5b87ux5b99ux3067ux9ed2ux4f53ux653eux5c04ux306fux3069ux3053ux306bux73feux308cux308bux304b]{第1章}。

\section{模範解答（問題
10-4）}\label{ux6a21ux7bc4ux89e3ux7b54ux554fux984c-10-4}

\textbf{(1)} \(f(\nu)=\nu^3(e^x-1)^{-1}\)、\(x=h\nu/k_BT\) なので
\(dx/d\nu=h/k_BT\)。積の微分をゼロとおき、\(\nu^2/(e^x-1)\) で割って
\(h\nu/k_BT=x\) を使うと

\[
3-\frac{x e^x}{e^x-1}=0\ \Longrightarrow\ 3(1-e^{-x})=x.
\]

\textbf{(2)} この超越方程式の非自明解は
\(x_\mathrm{peak}\approx2.821\)。よって
\(\nu_\mathrm{peak}=x_\mathrm{peak}k_BT/h=2.821\,k_BT/h\propto T\)（Wien
変位則）。

\textbf{(3)} \(T=5800\) K で
\(\nu_\mathrm{peak}\approx3.4\times10^{14}\)
Hz、\(\lambda=c/\nu_\mathrm{peak}\approx8.8\times10^{2}\)
nm（近赤外）。これは \(B_\lambda\) を最大化して得る
\(\lambda_\mathrm{peak}\approx500\) nm とは別物。\(B_\nu\) と
\(B_\lambda\) ではヤコビアン \(d\nu=-(c/\lambda^2)d\lambda\)
のため最大化条件が異なり（\(\lambda\) 側は
\(x\approx4.965\)）、\(\nu_\mathrm{peak}\lambda_\mathrm{peak}\ne c\)
となる。

\textbf{答え}：\(3(1-e^{-x})=x\Rightarrow x\approx2.821\)、\(\nu_\mathrm{peak}=2.821k_BT/h\propto T\)。太陽で
\(\nu_\mathrm{peak}\approx3.4\times10^{14}\) Hz（\(\sim880\)
nm）。\(B_\nu\) と \(B_\lambda\) のピークは一致しない。

\section*{さらに学ぶための参考文献}\label{sec-planck-quantum-further}
\addcontentsline{toc}{section}{さらに学ぶための参考文献}

\markright{さらに学ぶための参考文献}

\begin{itemize}
\tightlist
\item
  Planck, ``Ueber das Gesetz der Energieverteilung im Normalspectrum'',
  \emph{Annalen der Physik} 309 (1901) --- 量子仮説の原論文（独語）
\item
  Pais, \emph{Subtle is the Lord} (Oxford, 1982) --- Planck・Einstein
  の物理学的位置づけ
\item
  Reif, \emph{Fundamentals of Statistical and Thermal Physics}
  (McGraw-Hill, 1965) --- 量子化振動子の標準的扱い
\end{itemize}

\chapter{第11章　アインシュタインの方法}\label{ux7b2c11ux7ae0-ux30a2ux30a4ux30f3ux30b7ux30e5ux30bfux30a4ux30f3ux306eux65b9ux6cd5}

\textbf{この章の主題}：第10章では Planck (1900) の量子仮説
―「振動子のエネルギーが \(h\nu\) の整数倍に限られる」―
からプランク分布を導いた。本章では Einstein (1916/1917) の別ルート
―「光量子と原子の遷移確率の詳細釣り合い」―
から同じプランク分布を再導出する。この道筋で導入される \textbf{Einstein
係数}（\textbf{Einstein coefficients}）\(A_{ul}\), \(B_{lu}\),
\(B_{ul}\)
は、線スペクトル（第VII部）の中心的概念にもなり、本書の二つの中心地図（プランク関数地図と水素線地図）を結ぶ要となる。Einstein
の議論は「観測 → 物理」の逆引きの第二のルートを示すと同時に、第VIII部
QED の入口にもつながる。

\section*{この章の中心地図}\label{sec-einstein-method-map}
\addcontentsline{toc}{section}{この章の中心地図}

\markright{この章の中心地図}

\textbf{中心地図：プランク関数 ― 本部で扱う因子}

\[
B_\nu(T) = \frac{2\, \boxed{h}\, \nu^3}{c^2} \cdot \frac{1}{e^{\boxed{h}\nu/kT} - 1}
\]

強調されている \textbf{Planck 定数
\(h\)}（前因子の中と、指数の中の二箇所）が、第IV部で逆引きする対象である。エネルギーが連続でなく
\(h\nu\) という離散単位でやりとりされる ―
これが量子論の核心で、第10章（Planck の量子仮説）と第11章（Einstein
の方法）でその必然性を二つの異なるルートから示す。

\begin{tcolorbox}[enhanced jigsaw, arc=.35mm, bottomrule=.15mm, bottomtitle=1mm, breakable, colback=white, colbacktitle=quarto-callout-note-color!10!white, colframe=quarto-callout-note-color-frame, coltitle=black, left=2mm, leftrule=.75mm, opacityback=0, opacitybacktitle=0.6, rightrule=.15mm, title=\textcolor{quarto-callout-note-color}{\faInfo}\hspace{0.5em}{ノート}, titlerule=0mm, toprule=.15mm, toptitle=1mm]

\textbf{方針}：本章は中心地図プランク関数の
\textbf{量子論的逆引きの第二ルート}である。第10章（Planck
の量子仮説）と並行に、第8章（Bose
統計）とも合流する三つ目の経路。三つの異なる出発点から同じ占有因子
\(1/(e^{h\nu/kT}-1)\) が出ることが、その必然性を強く示す。

\end{tcolorbox}

\section*{この章で答える問い}\label{sec-einstein-method-questions}
\addcontentsline{toc}{section}{この章で答える問い}

\markright{この章で答える問い}

\begin{itemize}
\tightlist
\item
  Einstein はなぜ「光量子」を導入する必要があったのか
\item
  自発放出・誘導放出・誘導吸収はどう違い、どう関係しているのか
\item
  なぜ詳細釣り合いだけからプランク分布を再導出できるのか
\item
  Einstein の A 係数と B 係数は、観測のどの量にどう対応するのか
\end{itemize}

\section*{到達目標}\label{sec-einstein-method-goals}
\addcontentsline{toc}{section}{到達目標}

\markright{到達目標}

この章を読み終えると、読者は次のことができるようになる：

\begin{itemize}
\tightlist
\item
  自発放出と誘導放出の区別を、観測スペクトルとの関係で説明できる
\item
  詳細釣り合いの議論からプランク分布を再導出できる
\item
  Einstein
  係数が、連続スペクトル（第III〜V部）と線スペクトル（第VII部）の両方をつなぐ鍵であることを理解できる
\end{itemize}

\begin{center}\rule{0.5\linewidth}{0.5pt}\end{center}

\section{11.1 光量子 ― Einstein
の革命的提案}\label{sec-einstein-method-photon}

{[}本文目安：B3-B4{]}

Einstein (1905) は、Planck の量子仮説をより強い形で押し出した。Planck
は\textbf{空洞壁の物質振動子のエネルギー}が量子化されるとしたが、Einstein
は

\begin{quote}
\textbf{電磁場そのものが、エネルギー \(h\nu\) の粒子（光量子, light
quantum）の集まりとして振る舞う}
\end{quote}

と主張した。これが現代でいう
\textbf{光子}（\textbf{photon}）の概念である。光量子の証拠として
Einstein が引いたのは、\textbf{光電効果}（\textbf{photoelectric
effect}）― 光を金属に当てると電子が飛び出す現象 ― である。

光電効果の観測事実：

{\def\LTcaptype{none} % do not increment counter
\begin{longtable}[]{@{}
  >{\raggedright\arraybackslash}p{(\linewidth - 4\tabcolsep) * \real{0.3333}}
  >{\raggedright\arraybackslash}p{(\linewidth - 4\tabcolsep) * \real{0.3333}}
  >{\raggedright\arraybackslash}p{(\linewidth - 4\tabcolsep) * \real{0.3333}}@{}}
\toprule\noalign{}
\begin{minipage}[b]{\linewidth}\raggedright
観測される事実
\end{minipage} & \begin{minipage}[b]{\linewidth}\raggedright
古典波動論での予想
\end{minipage} & \begin{minipage}[b]{\linewidth}\raggedright
観測される実際
\end{minipage} \\
\midrule\noalign{}
\endhead
\bottomrule\noalign{}
\endlastfoot
入射光強度を増やすと飛び出す電子のエネルギー & 増える &
\textbf{変わらない}（電子数は増える） \\
入射光振動数を上げると & 関係なし & 電子のエネルギーが線形に増える \\
閾値振動数 \(\nu_0\) 以下では & 連続的に減衰 &
\textbf{電子が一切出ない} \\
\end{longtable}
}

これらは「電子が \(h\nu\)
のエネルギーを持つ光量子を一つずつ吸収する」と考えると自然に説明される。光量子の運動エネルギーが金属の仕事関数
\(\Phi\) を上回るとき電子が飛び出す：\(E_e = h\nu - \Phi\)。

Einstein はこの解析で 1921 年の Nobel
賞を受賞した（相対論ではなく！）。これが「光は量子化されている」という主張の最初の決定的証拠である。

\begin{tcolorbox}[enhanced jigsaw, arc=.35mm, bottomrule=.15mm, bottomtitle=1mm, breakable, colback=white, colbacktitle=quarto-callout-note-color!10!white, colframe=quarto-callout-note-color-frame, coltitle=black, left=2mm, leftrule=.75mm, opacityback=0, opacitybacktitle=0.6, rightrule=.15mm, title=\textcolor{quarto-callout-note-color}{\faInfo}\hspace{0.5em}{ノート}, titlerule=0mm, toprule=.15mm, toptitle=1mm]

\textbf{注意（理論）}：Einstein (1905)
の段階では、「光量子」は数学的な計算便法ではなく、\textbf{物理的に実在する粒子}として提案された。当時
Planck
自身もこの解釈には抵抗があったほど、革命的な主張だった。光量子の実在は
Compton 散乱（1923）で完全に確立される。

\end{tcolorbox}

\section{11.2 自発放出・誘導放出・誘導吸収 ―
三つの素過程}\label{sec-einstein-method-three-processes}

{[}本文目安：B3-B4{]}

Einstein (1916)
は、原子・分子の準位間遷移と光（光量子の集団）の相互作用を、\textbf{三つの基本過程}に分解した。下準位
\(l\) とエネルギーの高い上準位 \(u\)
の二準位系を考える（\(E_u - E_l = h\nu\)）：

\begin{figure}[H]

\centering{

\includegraphics[width=0.92\linewidth,height=\textheight,keepaspectratio]{index_files/mediabag/part4/../images/einstein-processes.pdf}

}

\caption{\label{fig-einstein-three-processes}二準位系における三つの素過程と
Einstein 係数の対応。自発放出は外部放射場に依存せず確率
\(A_{ul}\)、誘導吸収・誘導放出は放射場 \(J_\nu\) に比例する確率
\(B_{lu}J_\nu\), \(B_{ul}J_\nu\) で起こる。第VII部（線スペクトル）の
§21.2 で本図がそのまま再登場する。}

\end{figure}%

\begin{enumerate}
\def\labelenumi{\arabic{enumi}.}
\tightlist
\item
  \textbf{自発放出}（\textbf{spontaneous emission}）：上準位 \(u\)
  にある原子が、外部からの光なしに\textbf{自発的}に光量子を放出して下準位
  \(l\) へ落ちる。単位時間当たりの確率は \textbf{Einstein A 係数}
  \(A_{ul}\) で表される。
\item
  \textbf{誘導吸収}（\textbf{stimulated absorption}, あるいは単に
  \textbf{absorption}）：下準位 \(l\) にある原子が、外部の放射場
  \(J_\nu\) を吸収して上準位 \(u\) へ遷移する。単位時間当たりの確率は
  \(B_{lu} J_\nu\)。\textbf{Einstein B 係数} \(B_{lu}\)
  は遷移強度を、\(J_\nu\) は外部放射場の強さを表す。
\item
  \textbf{誘導放出}（\textbf{stimulated emission}）：上準位 \(u\)
  にある原子が、外部の放射場 \(J_\nu\)
  に\textbf{刺激されて}光量子を放出して下準位 \(l\) へ落ちる。確率は
  \(B_{ul} J_\nu\)。\(B_{ul}\) も Einstein B 係数の一つ。
\end{enumerate}

{\def\LTcaptype{none} % do not increment counter
\begin{longtable}[]{@{}llll@{}}
\toprule\noalign{}
過程 & 上下準位 & 外部光の役割 & 単位時間遷移率 \\
\midrule\noalign{}
\endhead
\bottomrule\noalign{}
\endlastfoot
自発放出 & \(u \to l\) & 不要 & \(A_{ul}\) \\
誘導吸収 & \(l \to u\) & 必要 & \(B_{lu} J_\nu\) \\
誘導放出 & \(u \to l\) & 必要 & \(B_{ul} J_\nu\) \\
\end{longtable}
}

\textbf{重要な点}：Einstein
が「\textbf{誘導放出}」を予言したのは、その後のレーザー（\textbf{LASER}
= Light Amplification by Stimulated Emission of
Radiation）の理論的基盤となる重要な貢献である。

\begin{tcolorbox}[enhanced jigsaw, arc=.35mm, bottomrule=.15mm, bottomtitle=1mm, breakable, colback=white, colbacktitle=quarto-callout-note-color!10!white, colframe=quarto-callout-note-color-frame, coltitle=black, left=2mm, leftrule=.75mm, opacityback=0, opacitybacktitle=0.6, rightrule=.15mm, title=\textcolor{quarto-callout-note-color}{\faInfo}\hspace{0.5em}{ノート}, titlerule=0mm, toprule=.15mm, toptitle=1mm]

\textbf{用語}：Einstein 係数 \(A_{ul}\), \(B_{lu}\), \(B_{ul}\)
の添字は遷移の向きを示す。\(u\) → \(l\) なら \(A_{ul}\) や
\(B_{ul}\)、\(l\) → \(u\) なら
\(B_{lu}\)。混同しないよう、本書では遷移を表す矢印（\(u \to l\)）と添字を並行に書くようにする。

\end{tcolorbox}

\begin{tcolorbox}[enhanced jigsaw, arc=.35mm, bottomrule=.15mm, bottomtitle=1mm, breakable, colback=white, colbacktitle=quarto-callout-important-color!10!white, colframe=quarto-callout-important-color-frame, coltitle=black, left=2mm, leftrule=.75mm, opacityback=0, opacitybacktitle=0.6, rightrule=.15mm, title=\textcolor{quarto-callout-important-color}{\faExclamation}\hspace{0.5em}{重要}, titlerule=0mm, toprule=.15mm, toptitle=1mm]

\textbf{本章では「三つの素過程」を仮定として導入する} ― Einstein 自身の
(1916/1917) 論文（短報が 1916、Physikalische Zeitschrift の本論文が
1917）に従い、自発放出・誘導吸収・誘導放出を\textbf{現象として要請}し、その帰結（プランク分布の再現）を見るのが本章のアプローチである。

しかし、それぞれの過程の\textbf{物理的起源}は本章では説明しない。これらは第VIII部（\textbf{QED
の見通し}）で扱われる：

\begin{itemize}
\tightlist
\item
  \textbf{自発放出}は、なぜ「外部の光なし」で起きるのか ―
  これは真空電磁場の零点ゆらぎ（節~\ref{sec-quantum-foundations-vacuum}）が励起原子を揺さぶる帰結。詳しくは
  節~\ref{sec-qed-minimum-natural-width} で扱う
\item
  \textbf{誘導放出・誘導吸収}は、量子化された電磁場の生成消滅演算子（第12章
  §12.3）と原子の電気双極子モーメントの結合から自然に出る。詳しくは
  節~\ref{sec-em-field-quantization-mode-oscillator}
\item
  \textbf{Einstein 関係式} \(A_{ul}/B_{ul} = 2h\nu^3/c^2\)
  そのものも、QED の枠組みで真空モード密度と直接結びつく ―
  節~\ref{sec-qed-minimum-einstein}
  で本書全体の「\textbf{プランク関数地図と水素線地図を結ぶ要}」として扱う
\end{itemize}

つまり本章は「Einstein
が辿った歴史的・現象論的ルート」を再現し、その\textbf{物理的基礎}は第VIII部で逆引きされる、という構造になっている。

\end{tcolorbox}

\section{11.3 Einstein 係数と詳細釣り合い ―
プランク分布の再導出}\label{sec-einstein-method-detailed-balance}

{[}本文目安：B4｜導出：M1{]}

Einstein の見事な議論は、\textbf{熱平衡における詳細釣り合い}（第6章
§6.6）だけからプランク分布を再導出できるという点である。

\textbf{設定}：温度 \(T\) の空洞に、二準位原子の集団と放射場 \(J_\nu\)
が熱平衡にある。下準位の占有数を \(n_l\)、上準位を \(n_u\)、外部放射場を
\(J_\nu\) とする。

\begin{tcolorbox}[enhanced jigsaw, arc=.35mm, bottomrule=.15mm, bottomtitle=1mm, breakable, colback=white, colbacktitle=quarto-callout-note-color!10!white, colframe=quarto-callout-note-color-frame, coltitle=black, left=2mm, leftrule=.75mm, opacityback=0, opacitybacktitle=0.6, rightrule=.15mm, title=\textcolor{quarto-callout-note-color}{\faInfo}\hspace{0.5em}{ノート}, titlerule=0mm, toprule=.15mm, toptitle=1mm]

\textbf{Einstein 係数の規約について}（重要）：放射場のどの量で \(B\)
係数を表すかで、係数の値・次元・関係式の前因子が変わる。本書は
\textbf{平均強度 \(J_\nu\)
規約}（\textbf{\(J_\nu\)-convention}）を採る：

\[
n_l B_{lu} J_\nu \quad \text{[s}^{-1}\text{]}, \qquad A_{ul}/B_{ul} = \frac{2h\nu^3}{c^2}
\]

これは Rybicki \& Lightman, Hubeny \& Mihalas 等の天体物理の標準。

別規約として、\textbf{エネルギー密度 \(u_\nu\) 規約} では
\(u_\nu = (4\pi/c) J_\nu\)（等方場）を用いて

\[
n_l B_{lu}^{(u)} u_\nu, \qquad A_{ul}/B_{ul}^{(u)} = \frac{8\pi h\nu^3}{c^3}
\]

となる（Mihalas (1978) や量子光学の文献に多い）。両規約の \(B\)
係数の関係は
\(B_{lu}^{(u)} = (c/4\pi) B_{lu}\)。論文・他書を読むときは、どちらの規約か必ず確認すること。

\end{tcolorbox}

\textbf{詳細釣り合い}：上下準位間の遷移率が両向きで等しい：

\begin{equation}\protect\phantomsection\label{eq-einstein-detailed-balance}{
n_l B_{lu} J_\nu \;=\; n_u (A_{ul} + B_{ul} J_\nu)
}\end{equation}

左辺は \(l \to u\)（吸収）、右辺は
\(u \to l\)（自発放出＋誘導放出）の総和。これを \(J_\nu\) について解くと

\begin{equation}\protect\phantomsection\label{eq-einstein-jnu-form}{
J_\nu = \frac{A_{ul}/B_{ul}}{(n_l/n_u)(B_{lu}/B_{ul}) - 1}
}\end{equation}

ここで \(n_u/n_l\) は熱平衡なら Boltzmann 分布で

\begin{equation}\protect\phantomsection\label{eq-einstein-boltzmann}{
\frac{n_u}{n_l} = \frac{g_u}{g_l} e^{-h\nu/k_B T}
}\end{equation}

これを 式~\ref{eq-einstein-jnu-form} に代入して整理すると

\begin{equation}\protect\phantomsection\label{eq-einstein-jnu-final}{
J_\nu = \frac{A_{ul}/B_{ul}}{(g_l B_{lu}/g_u B_{ul}) e^{h\nu/k_B T} - 1}
}\end{equation}

\textbf{観測されるプランク分布 \(J_\nu = B_\nu(T)\)
と一致させる}ためには、Einstein 係数の間に次の関係が必要となる：

\begin{equation}\protect\phantomsection\label{eq-einstein-B-ratio}{
g_l B_{lu} = g_u B_{ul}
}\end{equation}

\begin{equation}\protect\phantomsection\label{eq-einstein-A-B-ratio}{
\frac{A_{ul}}{B_{ul}} = \frac{2h\nu^3}{c^2}
}\end{equation}

これらが \textbf{Einstein の関係式}（\textbf{Einstein
relations}）である。詳細釣り合いという原理だけから、これら係数の比が自然界に必然として出る。

\begin{tcolorbox}[enhanced jigsaw, arc=.35mm, bottomrule=.15mm, breakable, colback=white, colframe=quarto-callout-tip-color-frame, left=2mm, leftrule=.75mm, opacityback=0, rightrule=.15mm, toprule=.15mm]
\begin{minipage}[t]{5.5mm}
\textcolor{quarto-callout-tip-color}{\faLightbulb}
\end{minipage}%
\begin{minipage}[t]{\textwidth - 5.5mm}

\textbf{問い}：なぜ Einstein
はわざわざ「自発放出」を別の過程として導入する必要があったのか？

\textbf{短答}：誘導吸収・誘導放出だけだと、放射場のない時にも準位間遷移が起こらないことになり、観測される現象（線スペクトルの放出）が説明できない。自発放出
\(A_{ul}\) がないと熱平衡に達するメカニズムが消える。

\textbf{もう一歩}：自発放出の物理的起源は、真空中の電磁場ゆらぎ ―
量子真空の零点振動 ― である（第VIII部・第26章）。Einstein
は古典的描像にこだわらず、必要な過程を仮定として導入することで物理を進めた。これは観測の必然から理論を駆動する好例である。

\end{minipage}%
\end{tcolorbox}

\section{11.4 観測とのつながり ―
連続と線の橋}\label{sec-einstein-method-bridge}

{[}本文目安：B3-B4{]}

Einstein 係数は、本書の二つの中心地図 ―
\textbf{プランク関数地図}と\textbf{水素線地図}（第VII部）―
を結ぶ要となる量である。

\textbf{連続スペクトル側との関係}：式~\ref{eq-einstein-A-B-ratio}
で出てくる \(2h\nu^3/c^2\)
は、\textbf{プランク関数の前因子と同じ}である。これは偶然ではなく、両者ともが「電磁場の同じモード構造」を見ている結果（第VIII部・第26章で詳しく扱う）。

\textbf{線スペクトル側との関係}：

\begin{itemize}
\tightlist
\item
  \textbf{自発放出率 \(A_{ul}\)}
  は、観測される\textbf{線の自然幅}（\textbf{natural
  width}）と直結する。励起寿命 \(\tau = 1/A_{ul}\)、自然幅
  \(\Delta\nu_{\mathrm{nat}} = A_{ul}/(2\pi)\)（第23章で本格的に扱う）。
\item
  \textbf{B 係数}は、観測される\textbf{線吸収係数
  \(\alpha_\nu^{\mathrm{line}}\) の振動子強度}と直接つながる（第22章）。
\item
  線のソース関数（第5章）：LTE では 式~\ref{eq-line-S-explicit} に従い
  \(S_\nu = (2h\nu^3/c^2)/[(g_u n_l)/(g_l n_u) - 1]\) となり、Boltzmann
  比 \(n_u/n_l = (g_u/g_l) e^{-h\nu/k_B T}\) を代入すると分母が
  \(e^{h\nu/k_B T} - 1\) に整理され、\(S_\nu = B_\nu(T)\)（第21章 §21.4
  で再導出）。
\end{itemize}

\begin{tcolorbox}[enhanced jigsaw, arc=.35mm, bottomrule=.15mm, bottomtitle=1mm, breakable, colback=white, colbacktitle=quarto-callout-note-color!10!white, colframe=quarto-callout-note-color-frame, coltitle=black, left=2mm, leftrule=.75mm, opacityback=0, opacitybacktitle=0.6, rightrule=.15mm, title=\textcolor{quarto-callout-note-color}{\faInfo}\hspace{0.5em}{ノート}, titlerule=0mm, toprule=.15mm, toptitle=1mm]

\textbf{観測との対応}：第VII部の線スペクトルの計算は、すべて Einstein
係数を介して行われる。「連続スペクトル」（第I〜V部）と「線スペクトル」（第VII部）が本書の二つの主柱と言ったが、両者は
Einstein
係数を介して\textbf{同じ電磁場--物質相互作用の二側面}として統合される。第VIII部
§25 でこの統合を本格的に閉じる。

\end{tcolorbox}

\section{11.5
分光学・天体物理への接続}\label{sec-einstein-method-astrophysics}

{[}本文目安：B3-B4{]}

Einstein 係数は天体物理の実務で常用される量である。代表的な応用：

\begin{itemize}
\tightlist
\item
  \textbf{線の強度}：観測される線スペクトル強度から、上準位の占有数
  \(n_u\) と自発放出率 \(A_{ul}\)
  の積を読む。これが第VII部の「線形成」の中心話題。
\item
  \textbf{励起温度}：観測される線比から \(n_u/n_l\) を求め、Boltzmann
  分布で温度を推定（第24章）。
\item
  \textbf{臨界密度}：衝突遷移率 \(C_{ul}\) と自発放出率 \(A_{ul}\)
  が等しくなる密度。これより低密度ではその線は non-LTE
  で観測されない（第21章）。
\item
  \textbf{メーザー}：誘導放出が自発放出を上回る非熱的状態 ―
  宇宙メーザー（OH, H\(_2\)O, SiO 等）の物理。
\end{itemize}

Einstein の (1916/1917)
論文は、現代の天体物理学・分光学・量子光学のすべてに基礎を提供している。

\begin{tcolorbox}[enhanced jigsaw, arc=.35mm, bottomrule=.15mm, bottomtitle=1mm, breakable, colback=white, colbacktitle=quarto-callout-note-color!10!white, colframe=quarto-callout-note-color-frame, coltitle=black, left=2mm, leftrule=.75mm, opacityback=0, opacitybacktitle=0.6, rightrule=.15mm, title=\textcolor{quarto-callout-note-color}{\faInfo}\hspace{0.5em}{ノート}, titlerule=0mm, toprule=.15mm, toptitle=1mm]

\textbf{使えるようになった道具}（§11.1〜11.5）：

\begin{itemize}
\tightlist
\item
  自発放出・誘導放出・誘導吸収という素過程に分解する見方
\item
  詳細釣り合い 式~\ref{eq-einstein-detailed-balance} から Einstein 関係
  式~\ref{eq-einstein-B-ratio}, 式~\ref{eq-einstein-A-B-ratio}
  を導く方法
\item
  \(A_{ul}\) から線の自然幅、\(B_{lu}\) から振動子強度 ―
  線スペクトルの計算手段全般
\end{itemize}

\end{tcolorbox}

\begin{center}\rule{0.5\linewidth}{0.5pt}\end{center}

\section*{この章で何がわかったか}\label{sec-einstein-method-summary}
\addcontentsline{toc}{section}{この章で何がわかったか}

\markright{この章で何がわかったか}

\textbf{中心地図に戻る}

本章で、中心地図プランク関数の構造が \textbf{詳細釣り合い}と
\textbf{三つの遷移過程}（自発・誘導吸収・誘導放出）から再構成された。Einstein
関係式 式~\ref{eq-einstein-A-B-ratio} の右辺に現れる
\textbf{\(2h\nu^3/c^2\)} は、プランク関数の前因子そのものである。

本章で導入された \textbf{Einstein 係数}は、本書の二つの中心地図 ―
プランク関数地図と水素線地図 ―
を結ぶ要となる量である。第VII部で本格的に活用するが、その下地はここで整っている。

\textbf{次章へ}：第12章では、本書の三度目のプランク分布導出ルート ―
\textbf{量子力学の基礎概念}（状態と観測量、調和振動子、生成消滅演算子、零点エネルギー）―
へ進む。これは第V部の電磁モード密度の議論と第VIII部の電磁場量子化への準備でもある。Planck
の振動子と Einstein の光量子が、量子力学の枠組みで一つに統合される。

\section*{演習問題}\label{sec-einstein-method-exercises}
\addcontentsline{toc}{section}{演習問題}

\markright{演習問題}

\textbf{模範解答}
は折りたたみを展開して確認できる（オンライン版）。まず自力で解いてから開くこと。

\subsection*{問題 11-1　詳細釣り合いから Einstein 関係式}\label{ex-11-1}
\addcontentsline{toc}{subsection}{問題 11-1　詳細釣り合いから Einstein
関係式}

{[}★ 難易度：☆☆ {]}

§11.3 は詳細釣り合いの式から Einstein
関係式を「整理すると」得た。途中代数を完全に埋めよ。出発点：

\[
n_l B_{lu}J_\nu=n_u(A_{ul}+B_{ul}J_\nu),\qquad \frac{n_u}{n_l}=\frac{g_u}{g_l}e^{-h\nu/k_BT}.
\]

\begin{enumerate}
\def\labelenumi{\arabic{enumi}.}
\tightlist
\item
  \(J_\nu\)
  について解き、\(J_\nu=\dfrac{A_{ul}/B_{ul}}{(g_l B_{lu}/g_u B_{ul})e^{h\nu/k_BT}-1}\)
  を導け。
\item
  これが任意の \(T\) で観測されるプランク分布
  \(B_\nu=\dfrac{2h\nu^3/c^2}{e^{h\nu/k_BT}-1}\)
  に一致する条件として、2つの Einstein 関係式 \(g_lB_{lu}=g_uB_{ul}\) と
  \(A_{ul}/B_{ul}=2h\nu^3/c^2\) を導け。
\item
  なぜ「任意の \(T\)
  で」一致させることが2式を同時に要求するのか説明せよ。
\end{enumerate}

\textbf{関連}：\hyperref[sec-einstein-method-detailed-balance]{§11.3
詳細釣り合い}／詳細釣り合いの原理は\hyperref[ux7b2c6ux7ae0-ux71b1ux5e73ux8861ux3068ux306fux4f55ux304b]{第6章
§6.6}、同じ占有因子の別ルートは\hyperref[ux7b2c8ux7ae0-ux5149ux5b50ux306eux7d71ux8a08]{第8章}と\hyperref[ex-10-1]{演習
10-1}。

\section{模範解答（問題
11-1）}\label{ux6a21ux7bc4ux89e3ux7b54ux554fux984c-11-1}

\textbf{(1)} \(J_\nu\) の項を集める：

\[
n_l B_{lu}J_\nu-n_u B_{ul}J_\nu=n_u A_{ul}\ \Rightarrow\ J_\nu=\frac{n_u A_{ul}}{n_l B_{lu}-n_u B_{ul}}=\frac{A_{ul}/B_{ul}}{\dfrac{n_l}{n_u}\dfrac{B_{lu}}{B_{ul}}-1}.
\]

\(n_l/n_u=(g_l/g_u)e^{+h\nu/k_BT}\) を代入：

\[
J_\nu=\frac{A_{ul}/B_{ul}}{\dfrac{g_lB_{lu}}{g_uB_{ul}}e^{h\nu/k_BT}-1}.
\]

\textbf{(2)} プランク分布 \(B_\nu=\dfrac{2h\nu^3/c^2}{e^{h\nu/k_BT}-1}\)
と項ごとに比較する。分母の指数の係数が \(1\) であるためには

\[
\frac{g_lB_{lu}}{g_uB_{ul}}=1\ \Longrightarrow\ g_lB_{lu}=g_uB_{ul}.
\]

これを使うと分子は \(A_{ul}/B_{ul}\)、これがプランクの前因子に等しい：

\[
\frac{A_{ul}}{B_{ul}}=\frac{2h\nu^3}{c^2}.
\]

\textbf{(3)} Einstein 係数は \(T\) に依らない原子の固有量。一方
\(J_\nu=B_\nu(T)\) はすべての \(T\) で成り立たねばならない。指数
\(e^{h\nu/k_BT}\) は \(T\)
で連続的に変わるので、「指数の係数＝1」（\(T\)
依存項の一致）と「全体の前因子の一致」（\(T\)
非依存項の一致）は独立な条件となり、2式を同時に要求する。

\textbf{答え}：\(g_lB_{lu}=g_uB_{ul}\) かつ
\(A_{ul}/B_{ul}=2h\nu^3/c^2\)。\(T\)
依存・非依存の両部分を一致させるため2式が必要。

\subsection*{問題 11-2　自発放出と誘導放出の比}\label{ex-11-2}
\addcontentsline{toc}{subsection}{問題 11-2　自発放出と誘導放出の比}

{[}★ 難易度：☆ {]}

熱放射場の中で、1個の上準位原子が起こす自発放出率と誘導放出率の比を求め、どの波長帯でどちらが優勢かを論じる。

\begin{enumerate}
\def\labelenumi{\arabic{enumi}.}
\tightlist
\item
  比 \(R=A_{ul}/(B_{ul}J_\nu)\) を、熱平衡 \(J_\nu=B_\nu(T)\) と
  Einstein 関係 \(A_{ul}/B_{ul}=2h\nu^3/c^2\)
  を使って簡単化し、\(R=e^{h\nu/k_BT}-1\) を示せ。
\item
  電波領域（\(h\nu\ll k_BT\)）と可視・紫外（\(h\nu\gg k_BT\)）で \(R\)
  の大小を評価し、それぞれどちらの放出が優勢かを述べよ。
\item
  \begin{enumerate}
  \def\labelenumii{(\arabic{enumii})}
  \setcounter{enumii}{1}
  \tightlist
  \item
    の結果が、宇宙メーザー（誘導放出の卓越）が電波領域に多い理由とどう関係するか簡潔に述べよ。
  \end{enumerate}
\end{enumerate}

\textbf{関連}：\hyperref[sec-einstein-method-three-processes]{§11.2
三つの素過程}、\hyperref[sec-einstein-method-astrophysics]{§11.5
メーザー}。

\section{模範解答（問題
11-2）}\label{ux6a21ux7bc4ux89e3ux7b54ux554fux984c-11-2}

\textbf{(1)} 熱場で
\(J_\nu=B_\nu=\dfrac{2h\nu^3/c^2}{e^{h\nu/k_BT}-1}\)。よって

\[
R=\frac{A_{ul}}{B_{ul}J_\nu}=\frac{A_{ul}/B_{ul}}{B_\nu}=\frac{2h\nu^3/c^2}{\dfrac{2h\nu^3/c^2}{e^{h\nu/k_BT}-1}}=e^{h\nu/k_BT}-1.
\]

\textbf{(2)}
電波（\(h\nu\ll k_BT\)）：\(R\approx h\nu/k_BT\ll1\)、誘導放出が自発放出を圧倒。可視・紫外（\(h\nu\gg k_BT\)）：\(R\approx e^{h\nu/k_BT}\gg1\)、自発放出が圧倒的に優勢。

\textbf{(3)}
メーザーは誘導放出が自発放出を上回ることで増幅が成立する。\(R\)
が小さい電波領域では、もともと熱場でも誘導放出が優勢になりやすく、反転分布が少し作られれば誘導放出が雪崩的に効く。一方、可視・紫外では
\(R\gg1\)（自発放出優勢）のため、増幅（レーザー）には強い反転分布と共振器が必要になる。よって自然界の「メーザー」は電波領域に多い。

\textbf{答え}：\(R=e^{h\nu/k_BT}-1\)。電波で
\(R\ll1\)（誘導優勢）、可視・紫外で
\(R\gg1\)（自発優勢）。宇宙メーザーが電波に多いことと整合。

\subsection*{問題 11-3　B 係数の規約変換（Jν 規約と uν
規約）}\label{ex-11-3}
\addcontentsline{toc}{subsection}{問題 11-3　B 係数の規約変換（Jν 規約と
uν 規約）}

{[}★ 難易度：☆☆ {]}

§11.3 のコラムは2つの規約を併記し \(B_{lu}^{(u)}=(c/4\pi)B_{lu}\)
とだけ述べた。これを導く。等方場では \(u_\nu=(4\pi/c)J_\nu\)。

\begin{enumerate}
\def\labelenumi{\arabic{enumi}.}
\tightlist
\item
  遷移率は規約に依らず物理的に同一であること、すなわち
  \(B_{lu}J_\nu=B_{lu}^{(u)}u_\nu\)
  を出発点に、\(B_{lu}^{(u)}=(c/4\pi)B_{lu}\) を導け。
\item
  \(J_\nu\) 規約での \(A_{ul}/B_{ul}=2h\nu^3/c^2\) から、\(u_\nu\)
  規約での比 \(A_{ul}/B_{ul}^{(u)}=8\pi h\nu^3/c^3\) を導け。
\item
  論文を読むとき規約の取り違えがどんな誤りを生むか（係数の数値が
  \(4\pi/c\) 倍ずれる等）を一言で述べよ。
\end{enumerate}

\textbf{関連}：\hyperref[sec-einstein-method-detailed-balance]{§11.3
係数の規約コラム}／平均強度 \(J_\nu\) とエネルギー密度 \(u_\nu\)
の定義は\hyperref[ux7b2c4ux7ae0-ux653eux5c04ux5834ux306eux57faux672cux91cf]{第4章
放射場}。

\section{模範解答（問題
11-3）}\label{ux6a21ux7bc4ux89e3ux7b54ux554fux984c-11-3}

\textbf{(1)} 同じ物理過程（吸収率
{[}s\(^{-1}\){]}）を2つの規約で書いた量は等しい：

\[
B_{lu}J_\nu=B_{lu}^{(u)}u_\nu.
\]

\(u_\nu=(4\pi/c)J_\nu\) を代入して \(J_\nu\) を消すと

\[
B_{lu}J_\nu=B_{lu}^{(u)}\frac{4\pi}{c}J_\nu\ \Rightarrow\ B_{lu}^{(u)}=\frac{c}{4\pi}B_{lu}.
\]

\textbf{(2)} \(B_{ul}\) も同様に \(B_{ul}^{(u)}=(c/4\pi)B_{ul}\)。よって

\[
\frac{A_{ul}}{B_{ul}^{(u)}}=\frac{A_{ul}}{(c/4\pi)B_{ul}}=\frac{4\pi}{c}\cdot\frac{A_{ul}}{B_{ul}}=\frac{4\pi}{c}\cdot\frac{2h\nu^3}{c^2}=\frac{8\pi h\nu^3}{c^3}.
\]

これは \(u_\nu\) 規約での標準形（量子光学でよく見る
\(8\pi h\nu^3/c^3\)）と一致する。

\textbf{(3)} どちらの規約で書かれた \(B\) かを取り違えると、係数の数値が
\(c/4\pi\)（または逆数）倍ずれ、線吸収係数・振動子強度・遷移率の換算がすべて狂う。論文・教科書を読むときは「\(B\)
が \(J_\nu\) に掛かるか \(u_\nu\) に掛かるか」を必ず確認する。

\textbf{答え}：\(B_{lu}^{(u)}=\tfrac{c}{4\pi}B_{lu}\)、\(A_{ul}/B_{ul}^{(u)}=8\pi h\nu^3/c^3\)。規約取り違えは係数を
\(4\pi/c\) 倍ずらす。

\subsection*{問題 11-4　Aul から励起寿命・自然幅・波長幅}\label{ex-11-4}
\addcontentsline{toc}{subsection}{問題 11-4　Aul
から励起寿命・自然幅・波長幅}

{[}★ 難易度：☆☆☆ {]}

§11.4 は「\(A_{ul}\)
は線の自然幅と直結、\(\tau=1/A_{ul}\)、\(\Delta\nu_\mathrm{nat}=A_{ul}/2\pi\)」とだけ述べた。典型的な許容遷移
\(A_{ul}=1\times10^8\ \mathrm{s^{-1}}\)、波長 \(\lambda=500\) nm
の線について数値評価せよ。

\begin{enumerate}
\def\labelenumi{\arabic{enumi}.}
\tightlist
\item
  励起寿命 \(\tau\) を求めよ。
\item
  自然幅 \(\Delta\nu_\mathrm{nat}=A_{ul}/2\pi\) を Hz で求めよ。
\item
  対応する波長幅
  \(\Delta\lambda_\mathrm{nat}=(\lambda^2/c)\Delta\nu_\mathrm{nat}\)
  を求め、線中心波長との相対幅 \(\Delta\lambda/\lambda\)
  を評価せよ。なぜ自然幅は実際の観測線幅より普通ずっと狭いのか（ドップラー幅との対比）を一言で述べよ。
\end{enumerate}

\textbf{関連}：\hyperref[sec-einstein-method-bridge]{§11.4
連続と線の橋}／自然幅の本格的扱いは\href{../part7/23-line-shape.qmd}{第23章
線の形}、その量子的起源（零点振動）は\hyperref[sec-quantum-foundations-vacuum]{§12.4}と\hyperref[ex-12-3]{演習
12-3}。

\section{模範解答（問題
11-4）}\label{ux6a21ux7bc4ux89e3ux7b54ux554fux984c-11-4}

\textbf{(1)} \(\tau=1/A_{ul}=1/10^8=1\times10^{-8}\) s \(=10\) ns。

\textbf{(2)}
\(\Delta\nu_\mathrm{nat}=A_{ul}/2\pi=10^8/(2\pi)\approx1.6\times10^7\)
Hz \(\approx16\) MHz。

\textbf{(3)} \(\lambda=500\)
nm、\(\lambda^2/c=(5.0\times10^{-7})^2/(3.0\times10^8)=2.5\times10^{-13}/3.0\times10^8\approx8.3\times10^{-22}\)
s·m。よって

\[
\Delta\lambda_\mathrm{nat}=\frac{\lambda^2}{c}\Delta\nu_\mathrm{nat}\approx8.3\times10^{-22}\times1.6\times10^7\approx1.3\times10^{-14}\ \mathrm{m}=1.3\times10^{-5}\ \mathrm{nm}.
\]

相対幅
\(\Delta\lambda/\lambda\approx1.3\times10^{-14}/5.0\times10^{-7}\approx2.7\times10^{-8}\)。きわめて狭い。\textbf{理由}：実際の星間・恒星大気の線は熱運動による\textbf{ドップラー幅}が支配的で、\(\Delta\nu_D/\nu\sim v_\mathrm{th}/c\sim10^{-5}\)
程度と自然幅より3桁ほど広い。自然幅が見えるのは線の遠い翼（ローレンツ翼）や低温・低速度分散の特殊な状況に限られる。

\textbf{答え}：\(\tau=10\) ns、\(\Delta\nu_\mathrm{nat}\approx16\)
MHz、\(\Delta\lambda_\mathrm{nat}\approx1.3\times10^{-5}\)
nm（\(\Delta\lambda/\lambda\sim3\times10^{-8}\)）。通常はドップラー幅が支配的で自然幅は埋もれる。

\section*{さらに学ぶための参考文献}\label{sec-einstein-method-further}
\addcontentsline{toc}{section}{さらに学ぶための参考文献}

\markright{さらに学ぶための参考文献}

\begin{itemize}
\tightlist
\item
  Einstein, ``Zur Quantentheorie der Strahlung'', \emph{Physikalische
  Zeitschrift} 18 (1917) --- Einstein 係数の原論文（独語）
\item
  Rybicki \& Lightman, \emph{Radiative Processes in Astrophysics}
  (Wiley, 1979) --- 第1章「Einstein Coefficients」と線輸送
\item
  Loudon, \emph{The Quantum Theory of Light} (Oxford, 3rd ed., 2000) ---
  Einstein 係数の量子光学的扱い
\item
  Hilborn, ``Einstein coefficients, cross sections, f values, dipole
  moments, and all that'', \emph{American Journal of Physics} 50, 982
  (1982) --- 係数間の換算公式の決定版
\end{itemize}

\chapter{第12章　量子力学の基礎へ戻る}\label{ux7b2c12ux7ae0-ux91cfux5b50ux529bux5b66ux306eux57faux790eux3078ux623bux308b}

\textbf{この章の主題}：第10章では Planck の量子仮説、第11章では Einstein
の係数法でプランク分布を導いた。これらは量子論の入口を別々に切り開く議論だったが、両者を統一的に整理する枠組みが
\textbf{量子力学}
である。本章では、観測スペクトルから第10・11章を経て要請された量子論的描像
―「量子化された調和振動子」「生成消滅演算子」「零点エネルギー」―
を量子力学の言葉で再定式化する。これは第V部の電磁モード密度の議論と、第VIII部の電磁場量子化（QED
の入口）への直接的な準備でもある。

\section*{この章の中心地図}\label{sec-quantum-foundations-map}
\addcontentsline{toc}{section}{この章の中心地図}

\markright{この章の中心地図}

\textbf{中心地図：プランク関数 ― 本部で扱う因子}

\[
B_\nu(T) = \frac{2\, \boxed{h}\, \nu^3}{c^2} \cdot \frac{1}{e^{\boxed{h}\nu/kT} - 1}
\]

強調されている \textbf{Planck 定数
\(h\)}（前因子の中と、指数の中の二箇所）が、第IV部で逆引きする対象である。エネルギーが連続でなく
\(h\nu\) という離散単位でやりとりされる ―
これが量子論の核心で、第10章（Planck の量子仮説）と第11章（Einstein
の方法）でその必然性を二つの異なるルートから示す。

\begin{tcolorbox}[enhanced jigsaw, arc=.35mm, bottomrule=.15mm, bottomtitle=1mm, breakable, colback=white, colbacktitle=quarto-callout-note-color!10!white, colframe=quarto-callout-note-color-frame, coltitle=black, left=2mm, leftrule=.75mm, opacityback=0, opacitybacktitle=0.6, rightrule=.15mm, title=\textcolor{quarto-callout-note-color}{\faInfo}\hspace{0.5em}{ノート}, titlerule=0mm, toprule=.15mm, toptitle=1mm]

\textbf{方針}：本章は第IV部の総まとめである。量子力学の数学的形式 ―
ヒルベルト空間、エルミート演算子、固有値問題 ― を最低限導入し、それが
Planck の振動子と Einstein
の光量子を統一的に扱える枠組みであることを示す。本書は量子力学の体系的教科書ではないので、必要最小限に絞る。

\end{tcolorbox}

\section*{この章で答える問い}\label{sec-quantum-foundations-questions}
\addcontentsline{toc}{section}{この章で答える問い}

\markright{この章で答える問い}

\begin{itemize}
\tightlist
\item
  量子力学の「状態」と「観測量」は、観測スペクトルの何に対応するのか
\item
  調和振動子の量子化は、なぜ電磁場の量子化と同型なのか
\item
  生成消滅演算子はどんな物理操作に対応するのか
\item
  零点エネルギーは観測可能なのか（→ Casimir 効果・自発放出）
\end{itemize}

\section*{到達目標}\label{sec-quantum-foundations-goals}
\addcontentsline{toc}{section}{到達目標}

\markright{到達目標}

この章を読み終えると、読者は次のことができるようになる：

\begin{itemize}
\tightlist
\item
  量子力学の基本枠組み（状態・観測量・固有値問題）を、本書の流れに位置づけて理解できる
\item
  調和振動子の量子化を、電磁場各モードの量子化へ自然に拡張できる準備を整える
\item
  生成消滅演算子の言葉で、光子の生成・消滅を語れる
\end{itemize}

\begin{center}\rule{0.5\linewidth}{0.5pt}\end{center}

\section{12.1 状態と観測量 ―
量子力学の最小限の言葉}\label{sec-quantum-foundations-states}

{[}本文目安：B3{]}

量子力学を、観測スペクトルとの関係に絞って最小限の言葉で整理する。

\textbf{状態}（\textbf{state}）：系の状態を、複素ベクトル空間（\textbf{Hilbert
space}）の要素 \(|\psi\rangle\)
で表す。古典力学では位相空間の点（座標と運動量の組）が状態だったが、量子力学ではより抽象的なベクトル。

\textbf{観測量}（\textbf{observable}）：観測される物理量は、Hilbert
空間上の \textbf{エルミート演算子}（\textbf{Hermitian operator},
\textbf{self-adjoint operator}）\(\hat{O}\)
で表される。古典的なエネルギー、運動量、位置などは、すべて対応する演算子（\(\hat{H}, \hat{p}, \hat{x}\)）として再定義される。

\textbf{固有値問題}：観測量を測定すると、対応する演算子の\textbf{固有値}（\textbf{eigenvalue}）のどれかが結果として得られる。

\begin{equation}\protect\phantomsection\label{eq-eigenvalue-problem}{
\hat{O} |o\rangle = o\, |o\rangle
}\end{equation}

ここで \(|o\rangle\) は固有ベクトル（\textbf{eigenstate}）、\(o\)
は固有値。たとえばエネルギー演算子 \(\hat{H}\) の固有値 \(E_n\)
がエネルギー準位を与える。

\textbf{確率振幅}：系が状態 \(|\psi\rangle\) にあるとき、観測値 \(o\)
が出る確率は
\(|\langle o | \psi \rangle|^2\)。これが量子論の確率解釈である。

\begin{tcolorbox}[enhanced jigsaw, arc=.35mm, bottomrule=.15mm, bottomtitle=1mm, breakable, colback=white, colbacktitle=quarto-callout-note-color!10!white, colframe=quarto-callout-note-color-frame, coltitle=black, left=2mm, leftrule=.75mm, opacityback=0, opacitybacktitle=0.6, rightrule=.15mm, title=\textcolor{quarto-callout-note-color}{\faInfo}\hspace{0.5em}{ノート}, titlerule=0mm, toprule=.15mm, toptitle=1mm]

\textbf{観測との対応}：本書のこれまでの議論で、観測スペクトル →
原子・分子のエネルギー準位 → 量子化
という流れがあった。「準位」とは具体的にエネルギー演算子 \(\hat{H}\)
の固有値である。観測される線スペクトルの振動数 \(\nu_{ul}\)
は、上下準位の固有値差 \(\nu_{ul} = (E_u - E_l)/h\)
で決まる。これが第22章で本格的に扱う水素原子のリュードベリ系列の起源である。

\end{tcolorbox}

\section{12.2 調和振動子の量子化 ―
構造の核心}\label{sec-quantum-foundations-oscillator}

{[}本文目安：B3{]}

第10章で見た「振動子のエネルギーが \(h\nu\)
の整数倍に量子化される」という現象を、量子力学の言葉で再定式化する。

質量 \(m\)、固有振動数 \(\omega = 2\pi\nu\) の調和振動子の Hamiltonian
は

\begin{equation}\protect\phantomsection\label{eq-harmonic-hamiltonian}{
\hat{H} = \frac{\hat{p}^2}{2m} + \frac{1}{2} m\omega^2 \hat{x}^2
}\end{equation}

ここで \(\hat{x}\) と \(\hat{p}\) は位置と運動量の演算子で、交換関係
\([\hat{x}, \hat{p}] = i\hbar\)（\(\hbar = h/2\pi\)）を満たす。

式~\ref{eq-harmonic-hamiltonian} の固有値問題を解くと、エネルギー準位は

\begin{equation}\protect\phantomsection\label{eq-harmonic-eigenvalues}{
E_n = \hbar\omega \left( n + \frac{1}{2} \right), \qquad n = 0, 1, 2, \ldots
}\end{equation}

となる。これが Planck の量子化に対応する離散エネルギーである。

\textbf{等間隔の準位}：隣接準位の差は
\(E_{n+1} - E_n = \hbar\omega = h\nu\)。これが Planck の「エネルギー量子
\(h\nu\)」の正体である。

\textbf{零点エネルギー}（\textbf{zero-point
energy}）：最低エネルギー（\(n = 0\)）でも
\(E_0 = \hbar\omega/2 \neq 0\)。古典では振動子が完全に止まる状態がエネルギーゼロだったが、量子では\textbf{不確定性原理}（\(\Delta x \cdot \Delta p \geq \hbar/2\)）から完全な静止が許されない。これが「真空でもエネルギーを持つ」という量子の特徴で、第VIII部・第25章で本格的に扱う。

\begin{tcolorbox}[enhanced jigsaw, arc=.35mm, bottomrule=.15mm, breakable, colback=white, colframe=quarto-callout-tip-color-frame, left=2mm, leftrule=.75mm, opacityback=0, rightrule=.15mm, toprule=.15mm]
\begin{minipage}[t]{5.5mm}
\textcolor{quarto-callout-tip-color}{\faLightbulb}
\end{minipage}%
\begin{minipage}[t]{\textwidth - 5.5mm}

\textbf{問い}：零点エネルギーは観測できるのか？

\textbf{短答}：絶対値（Hamiltonian
の零点）は観測不能だが、\textbf{差}は観測可能。たとえば二つの調和振動子の零点エネルギー差や、真空のモード構造変化（\textbf{Casimir
効果}）として現れる。

\textbf{もう一歩}：自発放出も零点エネルギーの観測的帰結である ―
真空場との結合により励起原子の自発放出率が非ゼロになる。第VIII部・第26章で扱う。CMB
の存在 ― そのものは古典電磁場として観測されるが ―
も、究極的には零点エネルギーから出発した宇宙史を映している。

\end{minipage}%
\end{tcolorbox}

\section{12.3 生成消滅演算子 ―
量子の最小公倍数}\label{sec-quantum-foundations-operators}

{[}本文目安：B4｜導出：M1{]}

式~\ref{eq-harmonic-hamiltonian}
は次のエレガントな形に書き直せる。\textbf{生成消滅演算子}（\textbf{creation
and annihilation operators}）を

\begin{equation}\protect\phantomsection\label{eq-creation-annihilation-def}{
\hat{a} = \sqrt{\frac{m\omega}{2\hbar}} \left( \hat{x} + \frac{i}{m\omega} \hat{p} \right), \qquad \hat{a}^\dagger = \sqrt{\frac{m\omega}{2\hbar}} \left( \hat{x} - \frac{i}{m\omega} \hat{p} \right)
}\end{equation}

と定義すると、Hamiltonian は

\begin{equation}\protect\phantomsection\label{eq-hamiltonian-number}{
\hat{H} = \hbar\omega \left( \hat{a}^\dagger \hat{a} + \frac{1}{2} \right) = \hbar\omega \left( \hat{n} + \frac{1}{2} \right)
}\end{equation}

ここで \(\hat{n} = \hat{a}^\dagger \hat{a}\) は
\textbf{数演算子}（\textbf{number operator}）で、固有値が
\(n = 0, 1, 2, \ldots\)。

\(\hat{a}^\dagger\) は数を一つ増やし（\textbf{creation},
生成）、\(\hat{a}\) は一つ減らす（\textbf{annihilation}, 消滅）：

\begin{equation}\protect\phantomsection\label{eq-creation-annihilation-action}{
\hat{a}^\dagger |n\rangle = \sqrt{n+1}\, |n+1\rangle, \qquad \hat{a} |n\rangle = \sqrt{n}\, |n-1\rangle
}\end{equation}

これらの交換関係 \([\hat{a}, \hat{a}^\dagger] = 1\)
から、すべてが綺麗に整う。

\textbf{物理的解釈}：生成消滅演算子は、振動子のエネルギー準位を一段階上げ下げする。電磁場の文脈では、\(\hat{a}^\dagger\)
が\textbf{光子を一つ生成}し、\(\hat{a}\)
が\textbf{光子を一つ消滅}させる。これが第VIII部の電磁場量子化の核心である。

\begin{tcolorbox}[enhanced jigsaw, arc=.35mm, bottomrule=.15mm, bottomtitle=1mm, breakable, colback=white, colbacktitle=quarto-callout-note-color!10!white, colframe=quarto-callout-note-color-frame, coltitle=black, left=2mm, leftrule=.75mm, opacityback=0, opacitybacktitle=0.6, rightrule=.15mm, title=\textcolor{quarto-callout-note-color}{\faInfo}\hspace{0.5em}{ノート}, titlerule=0mm, toprule=.15mm, toptitle=1mm]

\textbf{用語}：「生成」「消滅」は粒子の数を増やしたり減らしたりする操作の名前である。文献によっては「上昇/下降演算子」（ladder
operators）とも呼ばれる。本書では「生成消滅演算子」を主に用いる。

\end{tcolorbox}

\section{12.4 真空状態と数状態 ―
量子場の入口}\label{sec-quantum-foundations-vacuum}

{[}本文目安：B4｜導出：M1{]}

数演算子の固有状態 \(|n\rangle\) は \textbf{数状態}（\textbf{number
state}, \textbf{Fock
state}）と呼ばれ、粒子数が確定した状態を表す。最低エネルギー状態
\(|0\rangle\) を \textbf{真空}（\textbf{vacuum}）と呼ぶ。

真空 \(|0\rangle\) の特性：

\begin{itemize}
\tightlist
\item
  \(\hat{a} |0\rangle = 0\)（消滅は無効）
\item
  \(\hat{n} |0\rangle = 0\)（粒子数ゼロ）
\item
  \(\hat{H} |0\rangle = \frac{1}{2}\hbar\omega |0\rangle\)（零点エネルギー）
\end{itemize}

真空は「粒子がいない」状態だが、\textbf{ゼロ点振動}を持つ。物理的には：

\begin{itemize}
\tightlist
\item
  位置・運動量がともに不確定（\(\Delta x \neq 0\), \(\Delta p \neq 0\)）
\item
  真空のエネルギー揺らぎが、励起状態にある原子に\textbf{自発放出}を引き起こす（第11章
  §11.3 で導入した \(A_{ul}\) の起源）
\end{itemize}

電磁場の場合、各モードがこの調和振動子に対応する。すべてのモードが真空にある状態
\(|0\rangle_\text{EM}\)
が「電磁真空」で、\(\hat{a}^\dagger_{\nu, \mathbf{k}, \sigma}\) が振動数
\(\nu\)、波数 \(\mathbf{k}\)、偏光 \(\sigma\)
の光子を一つ生成する。これが第VIII部・第25章で本格的に扱う「電磁場の量子化」である。

\begin{tcolorbox}[enhanced jigsaw, arc=.35mm, bottomrule=.15mm, bottomtitle=1mm, breakable, colback=white, colbacktitle=quarto-callout-note-color!10!white, colframe=quarto-callout-note-color-frame, coltitle=black, left=2mm, leftrule=.75mm, opacityback=0, opacitybacktitle=0.6, rightrule=.15mm, title=\textcolor{quarto-callout-note-color}{\faInfo}\hspace{0.5em}{ノート}, titlerule=0mm, toprule=.15mm, toptitle=1mm]

\textbf{観測との対応}：真空エネルギーは観測スペクトルにも顔を出す。第11章で見た自発放出率
\(A_{ul}\)
は、真空場との結合（真空電磁場揺らぎとの相互作用）が遷移を誘起する効果として理解される。観測される線の自然幅
\(\Gamma_{\text{nat}} = A_{ul}/(2\pi)\)（第23章）は、究極的には真空ゼロ点振動の観測的痕跡である。CMB
の黒体スペクトルも、光子という量子場の励起状態が宇宙膨張で温度を下げてきた歴史の表れである。

\end{tcolorbox}

\section{12.5
第V部・第VIII部への準備}\label{sec-quantum-foundations-preview}

{[}本文目安：B3{]}

本章で導入した道具は、第V部以降の議論の前提になる：

\textbf{第V部（電磁モード）}：第V部で電磁場のモード密度
\(g(\nu) = 8\pi\nu^2/c^3\)
を本格的に導出する。各モードは本章の調和振動子に対応し、量子化されると光子の集まりとして扱える。

\textbf{第VII部（線スペクトル）}：第22章で水素原子のエネルギー準位を
Schrödinger 方程式から導く。これは本章 §12.1 で導入した固有値問題
式~\ref{eq-eigenvalue-problem}
の具体例である。第23章で線の自然幅を扱うが、その起源は本章で導入した零点エネルギー（→
自発放出）にある。

\textbf{第VIII部（QED）}：第25章で電磁場の量子化を扱う。これは本章 §12.3
の生成消滅演算子を、空間の各点での電磁モードに対して並べたものに等しい。第26章で
Einstein 係数の関係式 式~\ref{eq-einstein-A-B-ratio} を QED
の言葉で再構成する。

\begin{tcolorbox}[enhanced jigsaw, arc=.35mm, bottomrule=.15mm, bottomtitle=1mm, breakable, colback=white, colbacktitle=quarto-callout-note-color!10!white, colframe=quarto-callout-note-color-frame, coltitle=black, left=2mm, leftrule=.75mm, opacityback=0, opacitybacktitle=0.6, rightrule=.15mm, title=\textcolor{quarto-callout-note-color}{\faInfo}\hspace{0.5em}{ノート}, titlerule=0mm, toprule=.15mm, toptitle=1mm]

\textbf{使えるようになった道具}（§12.1〜12.5）：

\begin{itemize}
\tightlist
\item
  量子力学の最小限の枠組み（状態・観測量・固有値問題）
\item
  調和振動子の固有値 \(E_n = (n+1/2)\hbar\omega\) と等間隔準位 ―
  これが第10章で Planck が要請した「エネルギー量子
  \(h\nu\)」の量子力学的根拠
\item
  生成消滅演算子と数演算子 ― これが第VIII部の電磁場量子化の準備
\item
  真空の零点エネルギーと、それが自発放出・線の自然幅・Casimir
  効果として観測される構造
\end{itemize}

\end{tcolorbox}

\begin{center}\rule{0.5\linewidth}{0.5pt}\end{center}

\section*{この章で何がわかったか}\label{sec-quantum-foundations-summary}
\addcontentsline{toc}{section}{この章で何がわかったか}

\markright{この章で何がわかったか}

\textbf{中心地図に戻る}

第IV部全体（第9〜12章）で、中心地図プランク関数の \textbf{Planck 定数
\(h\)}
が、観測される黒体スペクトルから三つの異なるルートで必然として要請されることを示した：

\begin{enumerate}
\def\labelenumi{\arabic{enumi}.}
\tightlist
\item
  古典の破綻（第9章）から、エネルギーが離散単位であることが必要
\item
  Planck の量子仮説（第10章）から、その単位が \(h\nu\)
  であり、振動子の量子化が等分配を破る
\item
  Einstein の方法（第11章）から、詳細釣り合いだけで
  \(A_{ul}/B_{ul} = 2h\nu^3/c^2\) が出る
\end{enumerate}

そして本章（第12章）で、これらをすべて統一する量子力学の枠組み ―
調和振動子の固有値、生成消滅演算子、零点エネルギー ― を整えた。

中心地図プランク関数の \textbf{指数の中の \(h\)} と \textbf{前因子の
\(h\)} が、これでほぼ説明される。残るは前因子の
\(\nu^3/c^2\)（電磁モード密度の振動数依存）であり、これが第V部の主題になる。

\textbf{次部へ}：第V部（第13〜14章）では、本章で借用していた電磁モード密度
\(g(\nu) = 8\pi\nu^2/c^3\) を、Maxwell
方程式と空洞の境界条件から本格的に導出する。これでプランク分布の全因子が「観測
→ 物理」の逆引きで完全に説明される。

\section*{演習問題}\label{sec-quantum-foundations-exercises}
\addcontentsline{toc}{section}{演習問題}

\markright{演習問題}

\textbf{模範解答}
は折りたたみを展開して確認できる（オンライン版）。まず自力で解いてから開くこと。

\subsection*{問題
12-1　生成消滅演算子でハミルトニアンを対角化}\label{ex-12-1}
\addcontentsline{toc}{subsection}{問題
12-1　生成消滅演算子でハミルトニアンを対角化}

{[}★ 難易度：☆☆ {]}

§12.3 は \(\hat H=\hbar\omega(\hat a^\dagger\hat a+\tfrac12)\)
を結果として与えた。定義

\[
\hat a=\sqrt{\tfrac{m\omega}{2\hbar}}\left(\hat x+\tfrac{i}{m\omega}\hat p\right),\quad \hat a^\dagger=\sqrt{\tfrac{m\omega}{2\hbar}}\left(\hat x-\tfrac{i}{m\omega}\hat p\right)
\]

と \([\hat x,\hat p]=i\hbar\)
から、\(\hat H=\dfrac{\hat p^2}{2m}+\dfrac12 m\omega^2\hat x^2\)
がこの形になることを示せ。さらに \([\hat a,\hat a^\dagger]=1\)
を確かめよ。

\textbf{関連}：\hyperref[sec-quantum-foundations-operators]{§12.3
生成消滅演算子}、\hyperref[sec-quantum-foundations-oscillator]{§12.2
調和振動子}／準位 \(E_n=(n+\tfrac12)\hbar\omega\)
は第10章のエネルギー量子 \(h\nu\)
の量子力学的根拠（\hyperref[ex-10-1]{演習 10-1}）。

\section{模範解答（問題
12-1）}\label{ux6a21ux7bc4ux89e3ux7b54ux554fux984c-12-1}

\(\hat a^\dagger\hat a\) を展開する：

\[
\hat a^\dagger\hat a=\frac{m\omega}{2\hbar}\left(\hat x-\frac{i}{m\omega}\hat p\right)\left(\hat x+\frac{i}{m\omega}\hat p\right)=\frac{m\omega}{2\hbar}\left[\hat x^2+\frac{1}{m^2\omega^2}\hat p^2+\frac{i}{m\omega}(\hat x\hat p-\hat p\hat x)\right].
\]

\(\hat x\hat p-\hat p\hat x=[\hat x,\hat p]=i\hbar\) なので
\(\dfrac{i}{m\omega}(i\hbar)=-\dfrac{\hbar}{m\omega}\)。代入：

\[
\hat a^\dagger\hat a=\frac{m\omega}{2\hbar}\hat x^2+\frac{1}{2\hbar m\omega}\hat p^2-\frac12.
\]

両辺に \(\hbar\omega\) を掛ける：

\[
\hbar\omega\,\hat a^\dagger\hat a=\frac12 m\omega^2\hat x^2+\frac{\hat p^2}{2m}-\frac12\hbar\omega=\hat H-\frac12\hbar\omega.
\]

\[
\therefore\ \hat H=\hbar\omega\left(\hat a^\dagger\hat a+\frac12\right).
\]

\textbf{交換関係}：同様に
\(\hat a\hat a^\dagger=\dfrac{m\omega}{2\hbar}\hat x^2+\dfrac{1}{2\hbar m\omega}\hat p^2+\dfrac12\)。差をとると

\[
[\hat a,\hat a^\dagger]=\hat a\hat a^\dagger-\hat a^\dagger\hat a=\frac12-\left(-\frac12\right)=1.
\]

\textbf{答え}：\(\hat H=\hbar\omega(\hat a^\dagger\hat a+\tfrac12)\)。零点
\(\tfrac12\hbar\omega\) は交換子 \([\hat x,\hat p]=i\hbar\)
から出る。\([\hat a,\hat a^\dagger]=1\)。

\subsection*{問題
12-2　零点エネルギーを不確定性原理から導く}\label{ex-12-2}
\addcontentsline{toc}{subsection}{問題
12-2　零点エネルギーを不確定性原理から導く}

{[}★ 難易度：☆ {]}

§12.2 は「不確定性原理から完全静止が許されず
\(E_0=\tfrac12\hbar\omega\)」と述べた。これを最小化問題として導け。基底状態では
\(\langle\hat x\rangle=\langle\hat p\rangle=0\) なので
\(\langle\hat x^2\rangle=(\Delta x)^2,\ \langle\hat p^2\rangle=(\Delta p)^2\)。

\begin{enumerate}
\def\labelenumi{\arabic{enumi}.}
\tightlist
\item
  エネルギー期待値
  \(E=\dfrac{(\Delta p)^2}{2m}+\dfrac12 m\omega^2(\Delta x)^2\)
  に、最小不確定 \(\Delta p=\dfrac{\hbar}{2\Delta x}\) を代入せよ。
\item
  \(\Delta x\) について最小化し、\(E_{\min}=\tfrac12\hbar\omega\)
  と最適な \(\Delta x\) を求めよ。
\end{enumerate}

\textbf{関連}：\hyperref[sec-quantum-foundations-oscillator]{§12.2
零点エネルギー}／零点振動の観測的帰結（自発放出・Casimir）は\hyperref[sec-quantum-foundations-vacuum]{§12.4}、\href{../part8/25-em-field-quantization.qmd}{第25章}。

\section{模範解答（問題
12-2）}\label{ux6a21ux7bc4ux89e3ux7b54ux554fux984c-12-2}

\textbf{(1)} \(\Delta p=\hbar/(2\Delta x)\) を代入：

\[
E(\Delta x)=\frac{\hbar^2}{8m(\Delta x)^2}+\frac12 m\omega^2(\Delta x)^2.
\]

\textbf{(2)} \(u=(\Delta x)^2\) とおいて \(dE/du=0\)：

\[
-\frac{\hbar^2}{8m u^2}+\frac12 m\omega^2=0\ \Rightarrow\ u^2=\frac{\hbar^2}{4m^2\omega^2}\ \Rightarrow\ u=(\Delta x)^2=\frac{\hbar}{2m\omega}.
\]

代入：

\[
E_{\min}=\frac{\hbar^2}{8m}\cdot\frac{2m\omega}{\hbar}+\frac12 m\omega^2\cdot\frac{\hbar}{2m\omega}=\frac{\hbar\omega}{4}+\frac{\hbar\omega}{4}=\frac12\hbar\omega.
\]

運動エネルギーとポテンシャルが各 \(\tfrac14\hbar\omega\)
ずつ等分される点が最小。

\textbf{答え}：\((\Delta x)^2=\hbar/2m\omega\) で
\(E_{\min}=\tfrac12\hbar\omega\)。完全静止（\(\Delta x=0\)）は不確定性が禁じる。

\subsection*{問題
12-3　零点エネルギーはなぜプランク分布に現れないのか}\label{ex-12-3}
\addcontentsline{toc}{subsection}{問題
12-3　零点エネルギーはなぜプランク分布に現れないのか}

{[}★ 難易度：☆☆ {]}

第10章の量子化振動子では \(\langle E\rangle=h\nu/(e^x-1)\)
を使ったが、完全な量子力学では \(E_n=(n+\tfrac12)\hbar\omega\) で零点
\(\tfrac12\hbar\omega\) がある。両者の整合を考える。

\begin{enumerate}
\def\labelenumi{\arabic{enumi}.}
\tightlist
\item
  準位 \(E_n=(n+\tfrac12)\hbar\omega\) の熱平均が
  \(\langle E\rangle=\tfrac12\hbar\omega+\dfrac{\hbar\omega}{e^{\hbar\omega/k_BT}-1}\)
  となることを、占有数 \(\langle n\rangle=1/(e^{\hbar\omega/k_BT}-1)\)
  を使って示せ。
\item
  なぜ観測される放射エネルギー密度 \(u_\nu\)（プランク分布）に零点項
  \(\tfrac12\hbar\omega\)
  が現れないのか。物理的理由を述べよ（ヒント：定数オフセット、エネルギーの基準）。
\item
  では零点エネルギーは「観測できない」のか。観測可能な帰結の例を §12.4
  から2つ挙げよ。
\end{enumerate}

\textbf{関連}：\hyperref[sec-quantum-foundations-vacuum]{§12.4
真空状態}／プランク導出は\hyperref[sec-planck-quantum-derivation]{§10.3}（\hyperref[ex-10-1]{演習
10-1}）、零点の本格的扱いは\href{../part8/25-em-field-quantization.qmd}{第25章}、\href{../part8/26-qed-minimum.qmd}{第26章}。

\section{模範解答（問題
12-3）}\label{ux6a21ux7bc4ux89e3ux7b54ux554fux984c-12-3}

\textbf{(1)} Boltzmann 重み
\(P_n=e^{-\beta E_n}/Z\)。\(E_n=(n+\tfrac12)\hbar\omega\) なので

\[
\langle E\rangle=\sum_n\left(n+\tfrac12\right)\hbar\omega\,P_n=\tfrac12\hbar\omega+\hbar\omega\langle n\rangle=\frac12\hbar\omega+\frac{\hbar\omega}{e^{\hbar\omega/k_BT}-1}.
\]

（\(\langle n\rangle=\sum n P_n=1/(e^{\beta\hbar\omega}-1)\) は問題 10-1
の計算と同じ。）

\textbf{(2)} 零点項 \(\tfrac12\hbar\omega\)
は\textbf{温度に依らない定数オフセット}。放射場のエネルギー密度や物質との熱的なやりとり（吸収・放出）に効くのは、真空（\(n=0\)）より上の\textbf{励起分}
\(\hbar\omega\langle n\rangle\)
だけである。エネルギーの絶対基準（ハミルトニアンの零点）は物理に効かないので、観測される
\(u_\nu\)
は励起分すなわちプランク分布になる。零点をすべてのモードで足すと発散するが、これは絶対値として観測不能。

\textbf{(3)}
零点の\textbf{差}や\textbf{変化}は観測可能。例：①境界条件でモード構造を変えると真空エネルギーが変化する
\textbf{Casimir 効果}。②真空電磁場の揺らぎが励起原子を揺さぶって起こす
\textbf{自発放出}（\(A_{ul}\)
の起源、第11章）。これらは線の自然幅（第23章）として観測スペクトルにも痕跡を残す。

\textbf{答え}：\(\langle E\rangle=\tfrac12\hbar\omega+\hbar\omega/(e^{\hbar\omega/k_BT}-1)\)。零点は
\(T\) 非依存の定数オフセットで放射の励起分に効かないため \(u_\nu\)
に出ない。だが差は Casimir 効果・自発放出として観測可能。

\subsection*{問題 12-4　生成消滅演算子の作用の係数}\label{ex-12-4}
\addcontentsline{toc}{subsection}{問題 12-4　生成消滅演算子の作用の係数}

{[}★ 難易度：☆☆ {]}

§12.3 は係数を結果として与えた。\([\hat a,\hat a^\dagger]=1\) と数演算子
\(\hat n=\hat a^\dagger\hat a\)（\(\hat n|n\rangle=n|n\rangle\)、\(|n\rangle\)
は規格化）から、上昇・下降の係数を導け。

\begin{enumerate}
\def\labelenumi{\arabic{enumi}.}
\tightlist
\item
  \(\hat a^\dagger|n\rangle\) のノルムの2乗
  \(\langle n|\hat a\hat a^\dagger|n\rangle\)
  を計算し、\(\hat a^\dagger|n\rangle=\sqrt{n+1}\,|n+1\rangle\) を示せ。
\item
  同様に \(\hat a|n\rangle=\sqrt{n}\,|n-1\rangle\)
  を示し、\(\hat a|0\rangle=0\)（真空は消滅で潰れる）を確認せよ。
\item
  これらが電磁場の文脈で「光子の生成・消滅」に対応することを一言で述べよ。
\end{enumerate}

\textbf{関連}：\hyperref[sec-quantum-foundations-operators]{§12.3
生成消滅演算子}、\hyperref[sec-quantum-foundations-vacuum]{§12.4
数状態・真空}／電磁場各モードへの拡張は\href{../part8/25-em-field-quantization.qmd}{第25章
電磁場の量子化}。

\section{模範解答（問題
12-4）}\label{ux6a21ux7bc4ux89e3ux7b54ux554fux984c-12-4}

\textbf{(1)}
\(\hat a\hat a^\dagger=\hat a^\dagger\hat a+[\hat a,\hat a^\dagger]=\hat n+1\)
ゆえ

\[
\big\|\hat a^\dagger|n\rangle\big\|^2=\langle n|\hat a\hat a^\dagger|n\rangle=\langle n|(\hat n+1)|n\rangle=n+1.
\]

また
\(\hat n(\hat a^\dagger|n\rangle)=(\hat a^\dagger\hat n+\hat a^\dagger)|n\rangle=(n+1)\hat a^\dagger|n\rangle\)
より \(\hat a^\dagger|n\rangle\) は固有値 \(n+1\)
の状態。規格化して位相を実にとれば

\[
\hat a^\dagger|n\rangle=\sqrt{n+1}\,|n+1\rangle.
\]

\textbf{(2)} 同様に
\(\|\hat a|n\rangle\|^2=\langle n|\hat a^\dagger\hat a|n\rangle=\langle n|\hat n|n\rangle=n\)
ゆえ \(\hat a|n\rangle=\sqrt{n}\,|n-1\rangle\)。\(n=0\) では係数
\(\sqrt0=0\) となり \(\hat a|0\rangle=0\)。真空はこれ以上下げられない。

\textbf{(3)}
電磁場の各モードをこの調和振動子と同一視すると、\(\hat a^\dagger\)
はそのモードに光子を1個\textbf{生成}し、\(\hat a\)
は1個\textbf{消滅}させる。\(|n\rangle\) は「光子 \(n\)
個の状態」、\(|0\rangle\) は光子のない電磁真空に対応する。

\textbf{答え}：\(\|\hat a^\dagger|n\rangle\|^2=n+1\Rightarrow\hat a^\dagger|n\rangle=\sqrt{n+1}|n+1\rangle\)、\(\hat a|n\rangle=\sqrt n|n-1\rangle\)、\(\hat a|0\rangle=0\)。電磁場では光子の生成・消滅。

\section*{さらに学ぶための参考文献}\label{sec-quantum-foundations-further}
\addcontentsline{toc}{section}{さらに学ぶための参考文献}

\markright{さらに学ぶための参考文献}

\begin{itemize}
\tightlist
\item
  Sakurai \& Napolitano, \emph{Modern Quantum Mechanics} (Cambridge, 3rd
  ed., 2020) ---
  標準的な量子力学教科書、調和振動子と生成消滅演算子の体系的扱い
\item
  Griffiths, \emph{Introduction to Quantum Mechanics} (Cambridge, 3rd
  ed., 2018) --- 学部向け平易な入門
\item
  猪木慶治・川合光『量子力学 I・II』（講談社, 1994）---
  日本語の標準的な量子力学教科書
\item
  Cohen-Tannoudji et al., \emph{Quantum Mechanics} (Wiley, 1977) ---
  詳細な計算とアプリケーション豊富
\end{itemize}

\part{第V部　逆引き3：電磁気学へ}

\chapter{第13章　電磁波としての光}\label{ux7b2c13ux7ae0-ux96fbux78c1ux6ce2ux3068ux3057ux3066ux306eux5149}

\textbf{この章の主題}：第I〜IV部までで、観測される黒体スペクトルが
\(B_\nu(T) = (2h\nu^3/c^2)\cdot 1/(e^{h\nu/kT}-1)\)
の形であることを確認し、そのうち\textbf{指数関数因子}（占有因子）が光子の
Bose 統計とエネルギー量子化から出ることを示した。

残る課題は \textbf{前因子 \(2h\nu^3/c^2\) の起源}である。\(h\nu\)
はエネルギー量子（第IV部）、では \(\nu^3/c^2\)
はどこから来るのか。第8章・第9章でモード密度 \(g(\nu) = 8\pi\nu^2/c^3\)
を「借りた」が、その本格的な逆引きが第V部の仕事である。本章はその準備として、\textbf{光が電磁波であること}を
Maxwell 方程式から確認し、第14章での空洞モードの数え上げの基礎を整える。

\section*{この章の中心地図}\label{sec-electromagnetic-waves-map}
\addcontentsline{toc}{section}{この章の中心地図}

\markright{この章の中心地図}

\textbf{中心地図：プランク関数 ― 本部で扱う因子}

\[
B_\nu(T) = \frac{2 h\, \boxed{\nu^3}}{c^2} \cdot \frac{1}{e^{h\nu/kT} - 1}
\]

強調されている \textbf{\(\nu^3\)}（より精確には \(\nu^2/c^3\)
の振動数依存と、エネルギー量子 \(h\nu\)
の積）が、第V部で逆引きする対象である。これは空洞中の電磁場が三次元空間で離散的な定在波モードを持ち、振動数
\(\nu\) あたりのモード密度が \(\nu^2\)
に比例するという、電磁気学＋幾何の帰結である。第14章で本格的に導く。

\begin{tcolorbox}[enhanced jigsaw, arc=.35mm, bottomrule=.15mm, bottomtitle=1mm, breakable, colback=white, colbacktitle=quarto-callout-note-color!10!white, colframe=quarto-callout-note-color-frame, coltitle=black, left=2mm, leftrule=.75mm, opacityback=0, opacitybacktitle=0.6, rightrule=.15mm, title=\textcolor{quarto-callout-note-color}{\faInfo}\hspace{0.5em}{ノート}, titlerule=0mm, toprule=.15mm, toptitle=1mm]

\textbf{方針}：本章で扱うのはすべて\textbf{古典電磁気学}であり、量子は登場しない。第IV部で導入した
\(h\nu\)
と組み合わせるのは第14章末で。本章の射程は「光は電磁波である」「空洞中の境界条件の準備」までである。

\end{tcolorbox}

\section*{この章で答える問い}\label{sec-electromagnetic-waves-questions}
\addcontentsline{toc}{section}{この章で答える問い}

\markright{この章で答える問い}

\begin{itemize}
\tightlist
\item
  Maxwell
  方程式から波動方程式が出てくることは、観測される光のどの性質を保証するのか
\item
  偏光は観測の言葉では何を意味し、どんな天文学的情報を持つのか
\item
  Poynting ベクトルとフラックスの関係は何か
\item
  境界条件と空洞放射のスペクトルは、どこで結びつくのか
\end{itemize}

\section*{到達目標}\label{sec-electromagnetic-waves-goals}
\addcontentsline{toc}{section}{到達目標}

\markright{到達目標}

この章を読み終えると、読者は次のことができるようになる：

\begin{itemize}
\tightlist
\item
  Maxwell 方程式から平面波解を導き、振幅と偏光の自由度を説明できる
\item
  Poynting ベクトルを、観測量としてのフラックスに対応づけられる
\item
  境界条件の指定が、次章のモード密度議論の出発点になることを把握できる
\end{itemize}

\begin{center}\rule{0.5\linewidth}{0.5pt}\end{center}

\section{13.1 Maxwell 方程式 ―
古典電磁気の出発点}\label{sec-electromagnetic-waves-maxwell}

{[}本文目安：B2{]}

電磁気学の基本方程式 ― \textbf{Maxwell 方程式}（\textbf{Maxwell's
equations}） ― を出発点として受け入れる。真空中（電荷・電流なし）では
\textbf{SI 単位系で}

\begin{equation}\protect\phantomsection\label{eq-maxwell-divergence}{
\nabla \cdot \mathbf{E} = 0, \quad
\nabla \cdot \mathbf{B} = 0
}\end{equation}

\begin{equation}\protect\phantomsection\label{eq-maxwell-curl}{
\nabla \times \mathbf{E} = -\frac{\partial \mathbf{B}}{\partial t}, \quad
\nabla \times \mathbf{B} = \mu_0\epsilon_0\frac{\partial \mathbf{E}}{\partial t} = \frac{1}{c^2}\frac{\partial \mathbf{E}}{\partial t}
}\end{equation}

\(\mathbf{E}\) は電場（\textbf{electric field}）、\(\mathbf{B}\)
は磁場（\textbf{magnetic field}）、\(c = 1/\sqrt{\mu_0\epsilon_0}\)
は真空中の光速度。

\begin{tcolorbox}[enhanced jigsaw, arc=.35mm, bottomrule=.15mm, bottomtitle=1mm, breakable, colback=white, colbacktitle=quarto-callout-note-color!10!white, colframe=quarto-callout-note-color-frame, coltitle=black, left=2mm, leftrule=.75mm, opacityback=0, opacitybacktitle=0.6, rightrule=.15mm, title=\textcolor{quarto-callout-note-color}{\faInfo}\hspace{0.5em}{ノート}, titlerule=0mm, toprule=.15mm, toptitle=1mm]

\textbf{単位系について}：本書は SI を基本とする（付録 B §B.4
参照）。Gauss 単位系で書かれた古典文献（Jackson, Rybicki \& Lightman
など）と比較する際は、上式の各項に対し
\(\mathbf{B}_\text{SI} \leftrightarrow \mathbf{B}_\text{Gauss}/c\)、\(\epsilon_0 \to 1/(4\pi)\)
の置き換えで Gauss 形
\(\nabla \times \mathbf{B} = (1/c)\partial \mathbf{E}/\partial t\)
に対応する。

\end{tcolorbox}

\begin{tcolorbox}[enhanced jigsaw, arc=.35mm, bottomrule=.15mm, bottomtitle=1mm, breakable, colback=white, colbacktitle=quarto-callout-note-color!10!white, colframe=quarto-callout-note-color-frame, coltitle=black, left=2mm, leftrule=.75mm, opacityback=0, opacitybacktitle=0.6, rightrule=.15mm, title=\textcolor{quarto-callout-note-color}{\faInfo}\hspace{0.5em}{ノート}, titlerule=0mm, toprule=.15mm, toptitle=1mm]

\textbf{方針}：Maxwell
方程式そのものの導出は本書の射程外（標準的な学部電磁気学教科書を参照）。本書ではこれを与えられた出発点として、\textbf{光がここから出てくる}ことを示すことに集中する。

\end{tcolorbox}

\begin{tcolorbox}[enhanced jigsaw, arc=.35mm, bottomrule=.15mm, breakable, colback=white, colframe=quarto-callout-tip-color-frame, left=2mm, leftrule=.75mm, opacityback=0, rightrule=.15mm, toprule=.15mm]
\begin{minipage}[t]{5.5mm}
\textcolor{quarto-callout-tip-color}{\faLightbulb}
\end{minipage}%
\begin{minipage}[t]{\textwidth - 5.5mm}

\textbf{問い}：なぜ「真空中」の Maxwell 方程式から出発するのか？

\textbf{短答}：第14章で扱う空洞放射では、空洞内部は真空（または希薄）であり、放射場は電荷・電流のない領域での電磁場として扱える。プランク分布の前因子は、まさにこの真空中の電磁モードから出る。

\textbf{もう一歩}：物質中の Maxwell
方程式（誘電率・透磁率を含む）は、観測される連続放射のうち「ダストの修正黒体」（第15章）などに必要になる。本章は中心地図プランク分布の理想的前因子に集中する。

\end{minipage}%
\end{tcolorbox}

\section{\texorpdfstring{13.2 波動方程式 ― 光速度 \(c\)
の起源}{13.2 波動方程式 ― 光速度 c の起源}}\label{sec-electromagnetic-waves-wave-equation}

{[}本文目安：B2-B3{]}

式~\ref{eq-maxwell-curl} の二つの式から \(\mathbf{B}\)
を消去できる。一方の式に時間微分をとり、もう一方を空間 rot
で展開してベクトル恒等式
\(\nabla \times (\nabla \times \mathbf{E}) = \nabla(\nabla \cdot \mathbf{E}) - \nabla^2 \mathbf{E}\)
と 式~\ref{eq-maxwell-divergence} を使うと、

\begin{equation}\protect\phantomsection\label{eq-wave-equation}{
\nabla^2 \mathbf{E} - \frac{1}{c^2}\frac{\partial^2 \mathbf{E}}{\partial t^2} = 0
}\end{equation}

\(\mathbf{B}\) についても同じ形の方程式が出る。これが
\textbf{波動方程式}（\textbf{wave equation}）。

\textbf{重要な帰結}：Maxwell 方程式は、その中の定数 \(c\)
を\textbf{光の伝播速度}として持つ波動方程式を生む。電磁波の速度が真空中で
\(c\) であることは、Maxwell
方程式そのものに組み込まれている。これは観測される光の速さがどの基準系でも
\(c\) であることの古典的な根拠でもある（→
特殊相対性理論の出発点、本書の射程外）。

\begin{tcolorbox}[enhanced jigsaw, arc=.35mm, bottomrule=.15mm, bottomtitle=1mm, breakable, colback=white, colbacktitle=quarto-callout-note-color!10!white, colframe=quarto-callout-note-color-frame, coltitle=black, left=2mm, leftrule=.75mm, opacityback=0, opacitybacktitle=0.6, rightrule=.15mm, title=\textcolor{quarto-callout-note-color}{\faInfo}\hspace{0.5em}{ノート}, titlerule=0mm, toprule=.15mm, toptitle=1mm]

\textbf{観測との対応}：観測される黒体スペクトルの前因子 \(\nu^3/c^2\) の
\(c^2\) は、まさに \textbf{Maxwell 方程式に含まれる光速 \(c\)}
に直接由来する。光速が電磁気の基本定数として現れる ―
このことが、後で見るモード密度に \(c\) が現れる理由でもある。

\end{tcolorbox}

\section{13.3 平面波解 ―
振動数と波数}\label{sec-electromagnetic-waves-plane-wave}

{[}本文目安：B2-B3{]}

式~\ref{eq-wave-equation} の解の代表が \textbf{平面波}（\textbf{plane
wave}）である：

\begin{equation}\protect\phantomsection\label{eq-plane-wave}{
\mathbf{E}(\mathbf{r}, t) = \mathbf{E}_0 \exp\left[i(\mathbf{k}\cdot\mathbf{r} - \omega t)\right]
}\end{equation}

ここで \(\mathbf{k}\) は \textbf{波数ベクトル}（\textbf{wave
vector}）、\(\omega = 2\pi\nu\) は \textbf{角振動数}（\textbf{angular
frequency}）。式~\ref{eq-plane-wave} を 式~\ref{eq-wave-equation}
に代入すると、\textbf{分散関係}（\textbf{dispersion relation}）

\begin{equation}\protect\phantomsection\label{eq-dispersion-relation}{
\omega = c |\mathbf{k}|, \quad \text{すなわち} \quad \nu = c/\lambda
}\end{equation}

が出る。ここで \(\lambda = 2\pi/|\mathbf{k}|\) は
\textbf{波長}。これは本書を通じて使ってきた \(\nu = c/\lambda\)
の関係そのもの。

平面波には三つの「自由度」がある：

\begin{enumerate}
\def\labelenumi{\arabic{enumi}.}
\tightlist
\item
  \textbf{方向} \(\hat{\mathbf{k}}\)（単位ベクトル、空間角度2自由度）
\item
  \textbf{波数の大きさ} \(|\mathbf{k}|\)（あるいは振動数 \(\nu\)）
\item
  \textbf{偏光} \(\mathbf{E}_0\) の向き（次節）
\end{enumerate}

これらが空洞中で離散化されると、第14章のモードの数え上げになる。

\section{13.4 偏光 ―
横波の二つの自由度}\label{sec-electromagnetic-waves-polarization}

{[}本文目安：B2-B3{]}

式~\ref{eq-plane-wave} を 式~\ref{eq-maxwell-divergence} に代入すると、

\begin{equation}\protect\phantomsection\label{eq-transverse-condition}{
\mathbf{k} \cdot \mathbf{E}_0 = 0
}\end{equation}

つまり\textbf{電場は波数ベクトルに垂直}である ―
\textbf{横波}（\textbf{transverse wave}）の性質。

波数ベクトル \(\mathbf{k}\) が決まると、\(\mathbf{E}_0\) は
\(\mathbf{k}\)
に垂直な平面内で任意に選べる。この平面は二次元なので、\textbf{独立な偏光自由度は二つ}である：

\begin{itemize}
\tightlist
\item
  \textbf{直線偏光}（\textbf{linear polarization}）：\(\mathbf{E}_0\)
  が固定方向
\item
  \textbf{円偏光}（\textbf{circular polarization}）：時間とともに
  \(\mathbf{E}_0\) が回転（右円・左円）
\end{itemize}

任意の偏光は、二つの基底偏光の線形結合で書ける。

\begin{tcolorbox}[enhanced jigsaw, arc=.35mm, bottomrule=.15mm, bottomtitle=1mm, breakable, colback=white, colbacktitle=quarto-callout-tip-color!10!white, colframe=quarto-callout-tip-color-frame, coltitle=black, left=2mm, leftrule=.75mm, opacityback=0, opacitybacktitle=0.6, rightrule=.15mm, title=\textcolor{quarto-callout-tip-color}{\faLightbulb}\hspace{0.5em}{ヒント}, titlerule=0mm, toprule=.15mm, toptitle=1mm]

\textbf{観測との対応}：この \textbf{「2
つの偏光自由度」が、プランク分布の前因子に現れる係数 2
の起源}である。すなわち \(2h\nu^3/c^2\) の先頭の 2
は、電磁波が三次元空間で横波であることの直接的な帰結。本章の核となる事実。

\end{tcolorbox}

\begin{tcolorbox}[enhanced jigsaw, arc=.35mm, bottomrule=.15mm, bottomtitle=1mm, breakable, colback=white, colbacktitle=quarto-callout-note-color!10!white, colframe=quarto-callout-note-color-frame, coltitle=black, left=2mm, leftrule=.75mm, opacityback=0, opacitybacktitle=0.6, rightrule=.15mm, title=\textcolor{quarto-callout-note-color}{\faInfo}\hspace{0.5em}{ノート}, titlerule=0mm, toprule=.15mm, toptitle=1mm]

\textbf{観測との対応}：偏光は天文学では重要な情報源である：

\begin{itemize}
\tightlist
\item
  \textbf{シンクロトロン放射}（非熱的）の偏光 → 磁場の方向診断
\item
  \textbf{散乱光}の偏光 → 散乱体の幾何
\item
  \textbf{Zeeman 効果}（第23章）→ 線が偏光して分裂、磁場強度を測定
\item
  \textbf{CMB の B モード偏光} → 初期宇宙の重力波の指紋
\end{itemize}

ただし、本書の中心地図プランク分布そのものは\textbf{等方・無偏光}である（多数の方向と偏光の平均）。観測される偏光現象は、本書では第15章以降や第23章で扱う。

\end{tcolorbox}

\section{13.5 Poynting ベクトル ―
観測量への接続}\label{sec-electromagnetic-waves-poynting}

{[}本文目安：B3{]}

電磁場のエネルギーの流れを表すのが \textbf{Poynting
ベクトル}（\textbf{Poynting vector}）である。本書は SI
を主とするので、SI 形を主式とする：

\begin{equation}\protect\phantomsection\label{eq-poynting-vector}{
\mathbf{S} = \frac{1}{\mu_0} \mathbf{E} \times \mathbf{B}
}\end{equation}

Gauss 単位系では
\(\mathbf{S} = (c/4\pi)\, \mathbf{E} \times \mathbf{B}\)
と書く（換算の詳細は付録 B §B.4
参照）。以下の平面波の計算は、古典文献との対照のため Gauss 形で進める。

物理的意味：単位面積・単位時間あたりに通過する電磁エネルギー。すなわち\textbf{エネルギー流密度}である。第4章で扱った観測量としてのフラックス
\(F_\nu\) は、この Poynting
ベクトルの時間平均を振動数で分解したものに対応する。

平面波 式~\ref{eq-plane-wave} の場合、時間平均された Poynting ベクトルは

\begin{equation}\protect\phantomsection\label{eq-poynting-time-averaged}{
\langle \mathbf{S} \rangle = \frac{c}{8\pi} |\mathbf{E}_0|^2 \hat{\mathbf{k}}
}\end{equation}

これは波の方向 \(\hat{\mathbf{k}}\)
に向かう。電磁波がエネルギーを運ぶ方向は、波数ベクトルの方向と一致する。

\begin{tcolorbox}[enhanced jigsaw, arc=.35mm, bottomrule=.15mm, breakable, colback=white, colframe=quarto-callout-tip-color-frame, left=2mm, leftrule=.75mm, opacityback=0, rightrule=.15mm, toprule=.15mm]
\begin{minipage}[t]{5.5mm}
\textcolor{quarto-callout-tip-color}{\faLightbulb}
\end{minipage}%
\begin{minipage}[t]{\textwidth - 5.5mm}

\textbf{問い}：観測装置（望遠鏡＋検出器）が実際に検出するのは、\(\mathbf{E}\)
か \(\mathbf{B}\) か Poynting ベクトルか？

\textbf{短答}：本質的には Poynting ベクトル ―
単位面積・単位時間あたりのエネルギー ―
が検出器で熱や電子に変換される。波長分解（分光）すれば
\(F_\nu\)（あるいは比強度 \(I_\nu\)）になる。

\textbf{もう一歩}：電場や磁場そのものは、高周波（光）では振動が速すぎて直接測れない。しかし電波領域では位相情報を含めた検出が可能（電波干渉計など）。観測される量は、装置の感度帯域と分解能に依存する。

\end{minipage}%
\end{tcolorbox}

\section{13.6 境界条件 ―
第14章への入口}\label{sec-electromagnetic-waves-boundary}

{[}本文目安：B3{]}

第14章で扱う空洞放射の準備として、電磁場の
\textbf{境界条件}（\textbf{boundary conditions}）を確認する。

\textbf{完全導体の表面}（\textbf{perfect conductor}）：

\begin{itemize}
\tightlist
\item
  \textbf{電場の接線成分は連続でゼロ}：\(E_\parallel = 0\)
  （物理的には、導体内部の自由電子が瞬時に外部電場を打ち消す）
\item
  \textbf{磁場の法線成分はゼロ}：\(B_\perp = 0\) （表面電流の対応から）
\end{itemize}

これらが\textbf{空洞壁での反射}を導き、内部の電磁場が\textbf{定在波}（\textbf{standing
wave}）の重ね合わせとして表される根拠になる。

\begin{tcolorbox}[enhanced jigsaw, arc=.35mm, bottomrule=.15mm, bottomtitle=1mm, breakable, colback=white, colbacktitle=quarto-callout-note-color!10!white, colframe=quarto-callout-note-color-frame, coltitle=black, left=2mm, leftrule=.75mm, opacityback=0, opacitybacktitle=0.6, rightrule=.15mm, title=\textcolor{quarto-callout-note-color}{\faInfo}\hspace{0.5em}{ノート}, titlerule=0mm, toprule=.15mm, toptitle=1mm]

\textbf{方針}：実際の空洞は完全導体ではないが、空洞放射の\textbf{スペクトル形は壁の材質によらない}ことを第3章で示した。したがって、計算の便宜上「完全導体壁」を仮定して構わない。これは黒体放射の普遍性を逆手に取った計算戦略である。

\end{tcolorbox}

式~\ref{eq-plane-wave}
のような自由な平面波は空洞の有限境界では存在できず、境界条件
\(E_\parallel = 0\)
を満たす\textbf{離散的な定在波}だけが許される。この離散化が第14章のモードの数え上げ
― そしてプランク分布の前因子 \(\nu^3/c^2\) の起源 ― である。

\begin{tcolorbox}[enhanced jigsaw, arc=.35mm, bottomrule=.15mm, bottomtitle=1mm, breakable, colback=white, colbacktitle=quarto-callout-note-color!10!white, colframe=quarto-callout-note-color-frame, coltitle=black, left=2mm, leftrule=.75mm, opacityback=0, opacitybacktitle=0.6, rightrule=.15mm, title=\textcolor{quarto-callout-note-color}{\faInfo}\hspace{0.5em}{ノート}, titlerule=0mm, toprule=.15mm, toptitle=1mm]

\textbf{使えるようになった道具}（§13.1〜13.6）：

\begin{itemize}
\tightlist
\item
  Maxwell 方程式から波動方程式 式~\ref{eq-wave-equation} を導き、光速
  \(c\) がそこに組み込まれていることを確認
\item
  平面波解 式~\ref{eq-plane-wave} と分散関係 \(\omega = c|\mathbf{k}|\)
\item
  横波条件 式~\ref{eq-transverse-condition} から、偏光自由度が
  \textbf{2} であることを確認 ― これがプランク分布の前因子の係数 2
  の正体
\item
  Poynting ベクトル 式~\ref{eq-poynting-vector}
  が観測量フラックスに対応すること
\item
  完全導体壁での境界条件 (\(E_\parallel = 0\)) ―
  第14章で離散モードを数えるための準備
\end{itemize}

\end{tcolorbox}

\begin{center}\rule{0.5\linewidth}{0.5pt}\end{center}

\section*{この章で何がわかったか}\label{sec-electromagnetic-waves-summary}
\addcontentsline{toc}{section}{この章で何がわかったか}

\markright{この章で何がわかったか}

\textbf{中心地図に戻る}

本章で、中心地図プランク関数の前因子 \(2h\nu^3/c^2\) のうち、

\begin{itemize}
\tightlist
\item
  \textbf{光速 \(c\) の出どころ} ― Maxwell
  方程式の波動方程式に組み込まれている
\item
  \textbf{偏光自由度 2} ― 横波条件から出る
\end{itemize}

の二つが明らかになった。残るは \(\nu^2\)
という振動数依存性で、これは第14章で空洞中の電磁モードの数え上げから出てくる。

\textbf{次章へ}：第14章では、完全導体壁の空洞中での電磁定在波を数え上げ、\textbf{モード密度
\(g(\nu) = 8\pi\nu^2/c^3\)}
を本格的に導出する。第8章・第9章で「借用」していたこの式が、Maxwell
方程式と三次元幾何の帰結として現れる。これで中心地図プランク分布の\textbf{全因子}が、観測から逆向きに完全に解き明かされる。

\section*{演習問題}\label{sec-electromagnetic-waves-exercises}
\addcontentsline{toc}{section}{演習問題}

\markright{演習問題}

\textbf{模範解答}
は折りたたみを展開して確認できる（オンライン版）。まず自力で解いてから開くこと。

\subsection*{問題 13-1　Maxwell
方程式から波動方程式を導く}\label{ex-13-1}
\addcontentsline{toc}{subsection}{問題 13-1　Maxwell
方程式から波動方程式を導く}

{[}★ 難易度：☆☆ {]}

§13.2 はベクトル恒等式を使うとだけ述べた。真空 Maxwell 方程式から
\(\mathbf E\)
の波動方程式を導く。\(\nabla\cdot\mathbf E=0\)、\(\nabla\times\mathbf E=-\partial_t\mathbf B\)、\(\nabla\times\mathbf B=\dfrac{1}{c^2}\partial_t\mathbf E\)。

\begin{enumerate}
\def\labelenumi{\arabic{enumi}.}
\tightlist
\item
  \(\nabla\times\mathbf E=-\partial_t\mathbf B\) の両辺に
  \(\nabla\times\) を作用させ、恒等式
  \(\nabla\times(\nabla\times\mathbf E)=\nabla(\nabla\cdot\mathbf E)-\nabla^2\mathbf E\)
  を使って \(-\nabla^2\mathbf E=-\partial_t(\nabla\times\mathbf B)\)
  を示せ。
\item
  右辺に \(\nabla\times\mathbf B=\dfrac{1}{c^2}\partial_t\mathbf E\)
  を代入し、\(\nabla^2\mathbf E-\dfrac{1}{c^2}\partial_t^2\mathbf E=0\)
  を導け。
\item
  \(c=1/\sqrt{\mu_0\epsilon_0}\)
  が波の伝播速度としてこの式に現れることの意味（光速が電磁気の基本定数から決まる）を一言で述べよ。
\end{enumerate}

\textbf{関連}：\hyperref[sec-electromagnetic-waves-wave-equation]{§13.2
波動方程式}／前因子の \(1/c^2\)
の起源は\hyperref[sec-cavity-modes-frequency]{§14.4}（\hyperref[ex-14-1]{演習
14-1}）。

\section{模範解答（問題
13-1）}\label{ux6a21ux7bc4ux89e3ux7b54ux554fux984c-13-1}

\textbf{(1)}
\(\nabla\times(\nabla\times\mathbf E)=-\partial_t(\nabla\times\mathbf B)\)。左辺に恒等式を使い、\(\nabla\cdot\mathbf E=0\)（真空）より
\(\nabla(\nabla\cdot\mathbf E)=0\) ゆえ

\[
\nabla(\nabla\cdot\mathbf E)-\nabla^2\mathbf E=-\nabla^2\mathbf E=-\partial_t(\nabla\times\mathbf B).
\]

\textbf{(2)} \(\nabla\times\mathbf B=\dfrac{1}{c^2}\partial_t\mathbf E\)
を代入：

\[
-\nabla^2\mathbf E=-\partial_t\left(\frac{1}{c^2}\partial_t\mathbf E\right)=-\frac{1}{c^2}\partial_t^2\mathbf E\ \Rightarrow\ \nabla^2\mathbf E-\frac{1}{c^2}\partial_t^2\mathbf E=0.
\]

\textbf{(3)} この波動方程式の波の速さは方程式中の定数
\(c=1/\sqrt{\mu_0\epsilon_0}\)。電気・磁気の基本定数だけから光速が決まり、光が電磁波であること、その速度が真空中で
\(c\) であることが Maxwell
方程式に組み込まれている。これがプランク前因子の \(c\) の出どころ。

\textbf{答え}：\(\nabla^2\mathbf E-\frac{1}{c^2}\partial_t^2\mathbf E=0\)。波速
\(c=1/\sqrt{\mu_0\epsilon_0}\) が電磁気定数から決まる。

\subsection*{問題 13-2　分散関係と横波条件 ― 偏光自由度2}\label{ex-13-2}
\addcontentsline{toc}{subsection}{問題 13-2　分散関係と横波条件 ―
偏光自由度2}

{[}★ 難易度：☆☆ {]}

平面波
\(\mathbf E=\mathbf E_0\,e^{i(\mathbf k\cdot\mathbf r-\omega t)}\)
を使う。

\begin{enumerate}
\def\labelenumi{\arabic{enumi}.}
\tightlist
\item
  波動方程式に代入して分散関係
  \(\omega=c|\mathbf k|\)（\(\nu=c/\lambda\)）を導け。
\item
  \(\nabla\cdot\mathbf E=0\) に代入して
  \(\mathbf k\cdot\mathbf E_0=0\)（横波）を示せ。
\item
  横波条件から、与えられた \(\mathbf k\)
  に対し独立な偏光が2つであることを示し、これがプランク前因子
  \(2h\nu^3/c^2\) の先頭の「2」の起源であることを述べよ。
\end{enumerate}

\textbf{関連}：\hyperref[sec-electromagnetic-waves-plane-wave]{§13.3
平面波}、\hyperref[sec-electromagnetic-waves-polarization]{§13.4
偏光}／モード数え上げでの2倍は\hyperref[sec-cavity-modes-polarization]{§14.3}。

\section{模範解答（問題
13-2）}\label{ux6a21ux7bc4ux89e3ux7b54ux554fux984c-13-2}

\textbf{(1)} \(\partial_t\to-i\omega\)、\(\nabla\to i\mathbf k\) ゆえ
\(\nabla^2\to-|\mathbf k|^2\)、\(\partial_t^2\to-\omega^2\)。波動方程式は

\[
\left(-|\mathbf k|^2+\frac{\omega^2}{c^2}\right)\mathbf E_0=0\ \Rightarrow\ \omega^2=c^2|\mathbf k|^2\ \Rightarrow\ \omega=c|\mathbf k|.
\]

\(\omega=2\pi\nu\)、\(|\mathbf k|=2\pi/\lambda\) より
\(\nu=c/\lambda\)。

\textbf{(2)}
\(\nabla\cdot\mathbf E=i\mathbf k\cdot\mathbf E_0\,e^{i(\cdots)}=0\Rightarrow\mathbf k\cdot\mathbf E_0=0\)。電場は伝播方向に垂直＝横波。

\textbf{(3)} \(\mathbf E_0\) は \(\mathbf k\)
に垂直な平面（2次元）内で自由。よって独立な基底偏光は2つ（例：直線偏光2方向、または右・左円偏光）。任意の偏光はこの2つの線形結合。この「2」がモード数え上げで各
\(\mathbf k\) に2を掛ける根拠となり、プランク前因子の先頭の2になる。

\textbf{答え}：\(\omega=c|\mathbf k|\)、\(\mathbf k\cdot\mathbf E_0=0\)。横波ゆえ偏光は2自由度＝前因子の「2」。

\subsection*{問題 13-3　Poynting ベクトルの時間平均}\label{ex-13-3}
\addcontentsline{toc}{subsection}{問題 13-3　Poynting
ベクトルの時間平均}

{[}★ 難易度：☆☆ {]}

§13.5 は時間平均 Poynting ベクトル
\(\langle\mathbf S\rangle=\dfrac{c}{8\pi}|\mathbf E_0|^2\hat{\mathbf k}\)（Gauss
単位）を結果として与えた。導く。実数の平面波
\(\mathbf E=\mathbf E_0\cos(\mathbf k\cdot\mathbf r-\omega t)\)、電磁波では
\(|\mathbf B|=|\mathbf E|\) かつ
\(\mathbf B=\hat{\mathbf k}\times\mathbf E\)（Gauss）。

\begin{enumerate}
\def\labelenumi{\arabic{enumi}.}
\tightlist
\item
  \(\mathbf S=\dfrac{c}{4\pi}\mathbf E\times\mathbf B\)
  に代入し、\(\mathbf S=\dfrac{c}{4\pi}|\mathbf E_0|^2\cos^2(\mathbf k\cdot\mathbf r-\omega t)\,\hat{\mathbf k}\)
  を示せ。
\item
  \(\cos^2\) の時間平均が \(1/2\)
  であることを使い、\(\langle\mathbf S\rangle=\dfrac{c}{8\pi}|\mathbf E_0|^2\hat{\mathbf k}\)
  を導け。
\item
  エネルギーが波数ベクトル方向 \(\hat{\mathbf k}\)
  に運ばれること、これが観測量フラックス \(F_\nu\)
  に対応することを一言で述べよ。
\end{enumerate}

\textbf{関連}：\hyperref[sec-electromagnetic-waves-poynting]{§13.5
Poynting
ベクトル}／フラックスの定義は\hyperref[sec-radiation-field-flux]{第4章
§4.4}（\hyperref[sec-radiation-field-exercises]{第II部 演習}）。

\section{模範解答（問題
13-3）}\label{ux6a21ux7bc4ux89e3ux7b54ux554fux984c-13-3}

\textbf{(1)}
\(\mathbf B=\hat{\mathbf k}\times\mathbf E\)、\(\mathbf E\perp\hat{\mathbf k}\)
ゆえ
\(\mathbf E\times\mathbf B=\mathbf E\times(\hat{\mathbf k}\times\mathbf E)=\hat{\mathbf k}|\mathbf E|^2-\mathbf E(\mathbf E\cdot\hat{\mathbf k})=|\mathbf E|^2\hat{\mathbf k}\)（\(\mathbf E\cdot\hat{\mathbf k}=0\)）。\(|\mathbf E|^2=|\mathbf E_0|^2\cos^2(\cdots)\)。よって

\[
\mathbf S=\frac{c}{4\pi}|\mathbf E_0|^2\cos^2(\mathbf k\cdot\mathbf r-\omega t)\,\hat{\mathbf k}.
\]

\textbf{(2)} \(\langle\cos^2\rangle=1/2\)（1周期平均）ゆえ

\[
\langle\mathbf S\rangle=\frac{c}{4\pi}|\mathbf E_0|^2\cdot\frac12\,\hat{\mathbf k}=\frac{c}{8\pi}|\mathbf E_0|^2\hat{\mathbf k}.
\]

\textbf{(3)} エネルギー流は波数方向
\(\hat{\mathbf k}\)。これは単位面積・単位時間あたりのエネルギー流密度で、振動数分解すれば観測量フラックス
\(F_\nu\)（第4章）に対応する。望遠鏡が検出するのは本質的にこの量。

\textbf{答え}：\(\mathbf S=\frac{c}{4\pi}|\mathbf E_0|^2\cos^2(\cdots)\hat{\mathbf k}\)、時間平均で
\(\langle\mathbf S\rangle=\frac{c}{8\pi}|\mathbf E_0|^2\hat{\mathbf k}\)。\(\hat{\mathbf k}\)
方向のエネルギー流＝フラックス。

\subsection*{問題 13-4　偏光が運ぶ天文情報 ―
なぜ黒体は無偏光か}\label{ex-13-4}
\addcontentsline{toc}{subsection}{問題 13-4　偏光が運ぶ天文情報 ―
なぜ黒体は無偏光か}

{[}★ 難易度：☆☆ {]}

§13.4 は偏光の天文的意義を列挙した。整理する。

\begin{enumerate}
\def\labelenumi{\arabic{enumi}.}
\tightlist
\item
  偏光自由度が「ちょうど2」になるのは、電磁波が3次元空間で横波だからである。もし縦波成分が許されたら自由度は何個になるか。なぜ実際は2なのか。
\item
  シンクロトロン放射・散乱光・Zeeman
  効果が、それぞれ偏光からどんな物理量（磁場方向・散乱幾何・磁場強度）を引き出すか述べよ。
\item
  中心地図のプランク分布そのものは等方・無偏光である。なぜ熱平衡放射は無偏光なのか（多数の方向と偏光の平均）を述べよ。
\end{enumerate}

\textbf{関連}：\hyperref[sec-electromagnetic-waves-polarization]{§13.4
偏光}／散乱と偏光は\hyperref[sec-non-blackbody-scattering]{§15.5}、Zeeman
効果は\href{../part7/23-line-shape.qmd}{第23章}。

\section{模範解答（問題
13-4）}\label{ux6a21ux7bc4ux89e3ux7b54ux554fux984c-13-4}

\textbf{(1)} 3次元のベクトル場 \(\mathbf E_0\) は本来3成分。横波条件
\(\mathbf k\cdot\mathbf E_0=0\) が1成分（\(\mathbf k\)
方向＝縦波成分）を禁じるので、残り2成分（\(\mathbf k\)
垂直面内）が独立。もし縦波が許されれば3自由度だが、Maxwell
方程式（\(\nabla\cdot\mathbf E=0\)）が縦波を禁じるため実際は2。

\textbf{(2)}
シンクロトロン：放射の偏光面から磁場の方向（偏光は磁場に垂直）。散乱光：偏光の向き・度合いから散乱体の幾何（散乱角90°で最大偏光）。Zeeman
効果：線が偏光成分に分裂し、その分裂幅から磁場強度。

\textbf{(3)}
熱平衡放射は多数の独立な放射体（原子・モード）から、あらゆる方向・あらゆる偏光で等確率に放出される。2つの偏光状態が同じ強度で混ざり、方向平均すると正味の偏光は打ち消されてゼロになる。だから黒体放射は等方・無偏光。観測される偏光は、散乱や磁場など等方性を破る機構の存在を示す。

\textbf{答え}：Maxwell
が縦波を禁じ自由度2。偏光→磁場方向(シンクロトロン)/散乱幾何/磁場強度(Zeeman)。熱平衡は両偏光が等確率で混ざり無偏光。

\section*{さらに学ぶための参考文献}\label{sec-electromagnetic-waves-further}
\addcontentsline{toc}{section}{さらに学ぶための参考文献}

\markright{さらに学ぶための参考文献}

\begin{itemize}
\tightlist
\item
  Jackson, \emph{Classical Electrodynamics} (Wiley, 3rd ed., 1998) ---
  標準的な電磁気学の決定版
\item
  Griffiths, \emph{Introduction to Electrodynamics} (Cambridge, 4th ed.,
  2017) --- 学部向けの平易な入門
\item
  砂川重信『理論電磁気学』（紀伊國屋書店, 第3版, 1999）---
  日本語の標準教科書
\item
  Rybicki \& Lightman, \emph{Radiative Processes in Astrophysics}
  (Wiley, 1979) --- 第2章「Electromagnetic Theory」
\end{itemize}

\chapter{第14章　空洞中の電磁モード}\label{ux7b2c14ux7ae0-ux7a7aux6d1eux4e2dux306eux96fbux78c1ux30e2ux30fcux30c9}

\textbf{この章の主題}：第I〜IV部で構築してきた逆引きの
\textbf{最後のピース}が本章である。第13章で準備した Maxwell
方程式と完全導体壁の境界条件から、空洞中の電磁モードを数え上げ、\textbf{モード密度
\(g(\nu) = 8\pi\nu^2/c^3\)}
を本格的に導出する。これが第8〜10章で「借用」してきた式の起源である。導出が終われば、プランク分布

\[B_\nu(T) = \frac{2h\nu^3}{c^2} \cdot \frac{1}{e^{h\nu/k_BT}-1}\]

の\textbf{前因子と指数因子の両方}が、観測から逆向きに完全に説明されたことになる。本章は第V部の到達点であると同時に、第I〜V部全体（観測
→ 物理）の逆引き総まとめでもある。

\section*{この章の中心地図}\label{sec-cavity-modes-map}
\addcontentsline{toc}{section}{この章の中心地図}

\markright{この章の中心地図}

\textbf{中心地図：プランク関数 ― 本部で扱う因子}

\[
B_\nu(T) = \frac{2 h\, \boxed{\nu^3}}{c^2} \cdot \frac{1}{e^{h\nu/kT} - 1}
\]

強調されている \textbf{\(\nu^3\)}（より精確には \(\nu^2/c^3\)
の振動数依存と、エネルギー量子 \(h\nu\)
の積）が、第V部で逆引きする対象である。これは空洞中の電磁場が三次元空間で離散的な定在波モードを持ち、振動数
\(\nu\) あたりのモード密度が \(\nu^2\)
に比例するという、電磁気学＋幾何の帰結である。第14章で本格的に導く。

\begin{tcolorbox}[enhanced jigsaw, arc=.35mm, bottomrule=.15mm, bottomtitle=1mm, breakable, colback=white, colbacktitle=quarto-callout-note-color!10!white, colframe=quarto-callout-note-color-frame, coltitle=black, left=2mm, leftrule=.75mm, opacityback=0, opacitybacktitle=0.6, rightrule=.15mm, title=\textcolor{quarto-callout-note-color}{\faInfo}\hspace{0.5em}{ノート}, titlerule=0mm, toprule=.15mm, toptitle=1mm]

\textbf{方針}：本章の結末で、強調された \(\nu^3\) が「\textbf{モード密度
\(\nu^2\)（本章で導出）× 光子のエネルギー
\(h\nu\)（第IV部）}」の積として完全に逆引きされる。第I〜V部の物理的逆引きはここで一度閉じる。

\end{tcolorbox}

\section*{この章で答える問い}\label{sec-cavity-modes-questions}
\addcontentsline{toc}{section}{この章で答える問い}

\markright{この章で答える問い}

\begin{itemize}
\tightlist
\item
  空洞中の電磁波には、どんな「許される」状態が何個あるのか
\item
  なぜ三次元空間ではモード密度が \(\nu^2\) に比例するのか
\item
  プランク分布の前因子 \(2h\nu^3/c^2\)
  は、モード密度と何の積から来るのか
\item
  偏光自由度がモードの数え上げにどう効くのか
\end{itemize}

\section*{到達目標}\label{sec-cavity-modes-goals}
\addcontentsline{toc}{section}{到達目標}

\markright{到達目標}

この章を読み終えると、読者は次のことができるようになる：

\begin{itemize}
\tightlist
\item
  空洞中の電磁モードを境界条件のもとで列挙し、状態密度を導出できる
\item
  モード密度 \(\propto \nu^2\) が三次元の幾何から来ることを示せる
\item
  プランク分布の三層構造（熱平衡・量子化・モード密度）の最後のピースを埋められる
\end{itemize}

\begin{center}\rule{0.5\linewidth}{0.5pt}\end{center}

\section{14.1 空洞中の定在波 ―
境界条件と量子化}\label{sec-cavity-modes-standing-waves}

{[}本文目安：B2-B3{]}

辺長 \(L\) の立方体空洞を考える。壁は完全導体（第13章
§13.6）。三軸方向を \(x, y, z\) とし、原点を空洞の隅にとる。

第13章の平面波解 式~\ref{eq-plane-wave}
をそのまま入れると壁の境界条件を満たさない。代わりに、\textbf{定在波}（\textbf{standing
wave}）の重ね合わせとして電場を書く。電場の \(x\) 成分は、壁
\(y = 0, L\) と \(z = 0, L\) で接線方向（\(E_\parallel = 0\)）の条件から

\begin{equation}\protect\phantomsection\label{eq-Ex-standing-wave}{
E_x(\mathbf{r}, t) \propto \cos(k_x x) \sin(k_y y) \sin(k_z z) \, e^{-i\omega t}
}\end{equation}

の形になる（\(E_y, E_z\) も同様に対称な形）。境界条件
\(\sin(k_i L) = 0\) から、許される波数成分は

\begin{equation}\protect\phantomsection\label{eq-allowed-k}{
k_x = \frac{n_x \pi}{L}, \quad k_y = \frac{n_y \pi}{L}, \quad k_z = \frac{n_z \pi}{L}
}\end{equation}

ここで \(n_x, n_y, n_z\)
は\textbf{正の整数}（\(n_i = 1, 2, 3, \ldots\)）。\(n_x = n_y = n_z = 0\)
では全成分が消えて自明だが、軸成分の一部だけが 0
になる場合（例：\(n_x = 0\),
\(n_y, n_z \ge 1\)）は厳密には残る偏光自由度が一つで生き残る。ただし全モード数の中でこれらは「境界の格子点」上の有限個に過ぎず、熱力学極限（\(V \to \infty\)）でモード密度には寄与しない（次節以降の積分で測度ゼロ）。本書では
§14.2 以降の数え上げで \(n_i \ge 1\) のみを扱う。

\textbf{重要なポイント}：境界条件が、本来連続だった波数 \(\mathbf{k}\)
を\textbf{離散的な格子点}に量子化した。これは Planck
の量子化（第10章）とは別の意味の「量子化」 ―
古典電磁気と境界条件だけから出る幾何学的離散化である。

\begin{tcolorbox}[enhanced jigsaw, arc=.35mm, bottomrule=.15mm, bottomtitle=1mm, breakable, colback=white, colbacktitle=quarto-callout-note-color!10!white, colframe=quarto-callout-note-color-frame, coltitle=black, left=2mm, leftrule=.75mm, opacityback=0, opacitybacktitle=0.6, rightrule=.15mm, title=\textcolor{quarto-callout-note-color}{\faInfo}\hspace{0.5em}{ノート}, titlerule=0mm, toprule=.15mm, toptitle=1mm]

\textbf{用語}：「量子化」という言葉は、本書ではここまで二度の意味で使われている：

\begin{enumerate}
\def\labelenumi{\arabic{enumi}.}
\tightlist
\item
  \textbf{エネルギー量子化}（第10〜12章）：エネルギーが \(h\nu\)
  の整数倍に離散化される。\textbf{量子論固有}の現象
\item
  \textbf{モード量子化}（本節）：波数 \(\mathbf{k}\)
  が境界条件で離散化される。\textbf{古典的幾何}の現象
\end{enumerate}

両者は別物だが、最終的に第VIII部の電磁場の量子化で統合される（各「離散的モード」が「量子化された調和振動子」となり、その励起準位が光子数）。

\end{tcolorbox}

\section{\texorpdfstring{14.2 許される波数の数え上げ ― \(k\)
空間の格子}{14.2 許される波数の数え上げ ― k 空間の格子}}\label{sec-cavity-modes-k-counting}

{[}本文目安：B3{]}

式~\ref{eq-allowed-k} で許される \((n_x, n_y, n_z)\)
の組は、\textbf{正の整数の三つ組}である。これらを
\textbf{波数空間}（\textbf{k-space}）にプロットすると、原点付近に格子点が並ぶ。

\begin{figure}[H]

\centering{

\includegraphics[width=0.88\linewidth,height=\textheight,keepaspectratio]{index_files/mediabag/part5/../images/k-space-counting.pdf}

}

\caption{\label{fig-k-space}\(k\) 空間における格子点の並び（格子間隔
\(\pi/L\)）と、波数 \(k\) 未満の数え上げに使う \textbf{第一象限（1/8
球）の球殻}。プランク分布の前因子 \(8\pi\nu^2/c^3\)
を生む計算の幾何を、本書最大の数え上げの核として示す。}

\end{figure}%

格子間隔は各軸方向に
\(\pi/L\)。したがって、波数空間における\textbf{単位体積あたりの格子点密度}は

\begin{equation}\protect\phantomsection\label{eq-k-space-density}{
\rho_{k\text{-space}} = \frac{1}{(\pi/L)^3} = \frac{L^3}{\pi^3} = \frac{V}{\pi^3}
}\end{equation}

ここで \(V = L^3\) は空洞の体積。

\textbf{波数の大きさ}
\(k = |\mathbf{k}| = \sqrt{k_x^2 + k_y^2 + k_z^2}\) が \(k\)
未満であるモードの個数は、波数空間で\textbf{半径 \(k\)
の球の体積}を考えればよい。ただし格子点は正の整数のみなので、\textbf{第一象限（octant,
1/8 のみ）}を取る：

\begin{equation}\protect\phantomsection\label{eq-mode-count-k-onepol}{
N_\text{mode}^{(1)}(<k) = \rho_{k\text{-space}} \cdot \frac{1}{8} \cdot \frac{4\pi k^3}{3} = \frac{V k^3}{6\pi^2}
}\end{equation}

ここで指数 (1)
は「偏光自由度を考慮していない」ことを表す（次節で考慮）。

\begin{tcolorbox}[enhanced jigsaw, arc=.35mm, bottomrule=.15mm, breakable, colback=white, colframe=quarto-callout-tip-color-frame, left=2mm, leftrule=.75mm, opacityback=0, rightrule=.15mm, toprule=.15mm]
\begin{minipage}[t]{5.5mm}
\textcolor{quarto-callout-tip-color}{\faLightbulb}
\end{minipage}%
\begin{minipage}[t]{\textwidth - 5.5mm}

\textbf{問い}：なぜ「1/8 倍」が必要なのか？ 全球面 \(4\pi k^3/3\)
で何が間違うのか？

\textbf{短答}：許される \(\mathbf{k}\) は \(n_x, n_y, n_z > 0\)
のみ。\(\mathbf{k}\) と \(-\mathbf{k}\)
は同じ定在波を生む（\(\cos(k_x x) = \cos(-k_x x)\)）ので、両方を独立に数えると過大評価になる。波数空間の第一象限だけが独立なモードを表す。

\textbf{もう一歩}：もし周期境界条件（壁ではなく端と端をつなぐ）にすると、\(\mathbf{k}\)
は正負両方とり、格子間隔は \(2\pi/L\)
に変わる。両方法は最終的に同じモード密度を与える ―
「壁にどう接続するか」は最終結果に効かない。これは第3章の「黒体放射の普遍性」を逆手にとった計算戦略である。

\end{minipage}%
\end{tcolorbox}

\section{14.3 偏光自由度 ―
横波の2倍}\label{sec-cavity-modes-polarization}

{[}本文目安：B2-B3{]}

第13章 §13.4 で見たように、与えられた \(\mathbf{k}\)
ベクトルに対し、独立な偏光方向は\textbf{二つ}ある（横波条件）。

式~\ref{eq-mode-count-k-onepol} に偏光自由度を掛けて：

\begin{equation}\protect\phantomsection\label{eq-mode-count-k}{
N_\text{mode}(<k) = 2 \cdot \frac{V k^3}{6\pi^2} = \frac{V k^3}{3\pi^2}
}\end{equation}

これが波数 \(k\) 未満のモードの総数。\(V\)
に比例するのは、大きな空洞ほどモード数が多いという当然の結果。

\begin{tcolorbox}[enhanced jigsaw, arc=.35mm, bottomrule=.15mm, bottomtitle=1mm, breakable, colback=white, colbacktitle=quarto-callout-note-color!10!white, colframe=quarto-callout-note-color-frame, coltitle=black, left=2mm, leftrule=.75mm, opacityback=0, opacitybacktitle=0.6, rightrule=.15mm, title=\textcolor{quarto-callout-note-color}{\faInfo}\hspace{0.5em}{ノート}, titlerule=0mm, toprule=.15mm, toptitle=1mm]

\textbf{観測との対応}：この因子 2 が、プランク分布の前因子
\(2h\nu^3/c^2\) の先頭の 2 に直接対応する。観測される黒体スペクトル ―
CMB であれ恒星であれ ―
は、すべて電磁波が横波であるという\textbf{幾何的事実}を反映している。

\end{tcolorbox}

\section{14.4 振動数空間に変換 ―
モード密度の正体}\label{sec-cavity-modes-frequency}

{[}本文目安：B3{]}

波数 \(k\) と振動数 \(\nu\) は第13章の分散関係
式~\ref{eq-dispersion-relation}
で結ばれる：\(k = 2\pi\nu/c\)。式~\ref{eq-mode-count-k} に代入して

\begin{equation}\protect\phantomsection\label{eq-mode-count-nu}{
N_\text{mode}(<\nu) = \frac{V}{3\pi^2} \left(\frac{2\pi\nu}{c}\right)^3 = \frac{8\pi V \nu^3}{3 c^3}
}\end{equation}

これが振動数 \(\nu\) 未満のモードの総数。微分して
\textbf{振動数あたり・単位体積あたりのモード密度} を得る：

\begin{equation}\protect\phantomsection\label{eq-mode-density}{
g(\nu) = \frac{1}{V} \frac{dN_\text{mode}}{d\nu} = \frac{8\pi \nu^2}{c^3}
}\end{equation}

これが第8章 式~\ref{eq-photon-energy-density} と第9章
式~\ref{eq-mode-density-borrowed}
で\textbf{借用}していた式である。本章で本格的に逆引きが完了した。

\textbf{\(\nu^2\) の起源を辿る}：

{\def\LTcaptype{none} % do not increment counter
\begin{longtable}[]{@{}
  >{\raggedright\arraybackslash}p{(\linewidth - 2\tabcolsep) * \real{0.5000}}
  >{\raggedright\arraybackslash}p{(\linewidth - 2\tabcolsep) * \real{0.5000}}@{}}
\toprule\noalign{}
\begin{minipage}[b]{\linewidth}\raggedright
因子
\end{minipage} & \begin{minipage}[b]{\linewidth}\raggedright
起源
\end{minipage} \\
\midrule\noalign{}
\endhead
\bottomrule\noalign{}
\endlastfoot
\(\nu^2\) & \(k\) 空間の球の体積 \(k^3\) を微分（三次元空間） \\
\(8\pi\) & \(4\pi/3\)（球の体積） × \(1/8\)（第一象限） ×
\(2\)（偏光自由度） × \((2\pi)^3/\pi^3\)（\(k = 2\pi\nu/c\)
の三乗が格子密度 \(\pi^{-3}\) と打ち消し合い、\(2^3\) が残る） \\
\(1/c^3\) & 分散関係 \(k = 2\pi\nu/c\) から \\
\end{longtable}
}

\begin{tcolorbox}[enhanced jigsaw, arc=.35mm, bottomrule=.15mm, breakable, colback=white, colframe=quarto-callout-tip-color-frame, left=2mm, leftrule=.75mm, opacityback=0, rightrule=.15mm, toprule=.15mm]
\begin{minipage}[t]{5.5mm}
\textcolor{quarto-callout-tip-color}{\faLightbulb}
\end{minipage}%
\begin{minipage}[t]{\textwidth - 5.5mm}

\textbf{問い}：なぜ三次元空間でモード密度が \(\nu^2\) なのか？
他の次元では？

\textbf{短答}：\(d\) 次元空間では、\(k\) 空間の球の表面積は
\(\propto k^{d-1}\)。これに対応してモード密度は
\(\nu^{d-1}\)。三次元なら \(\nu^2\)、二次元なら
\(\nu\)、一次元なら定数。

\textbf{もう一歩}：もし宇宙が二次元空間だったら、モード密度は
\(g(\nu) \propto \nu/c^2\) となり（\(d\) 次元では
\(\nu^{d-1}/c^d\)）、プランク分布の前因子は \(2h\nu^3/c^2\) ではなく
\(\propto h\nu^2/c^2\)
の形をとる。黒体スペクトル全体の形（Rayleigh-Jeans
側のべき・ピーク位置・Wien
側の減衰）は完全に違うものになる。\textbf{観測される黒体スペクトルの形には、私たちが三次元空間に住んでいるという情報が組み込まれている}。

\end{minipage}%
\end{tcolorbox}

\section{14.5 プランク分布の前因子への合流 ―
第I〜V部の総まとめ}\label{sec-cavity-modes-bnu-derivation}

{[}本文目安：B3｜導出：M1{]}

モード密度 式~\ref{eq-mode-density}
を、第8章で立ち上げた\textbf{振動数あたりエネルギー密度}の表式に代入する。第8章
式~\ref{eq-photon-energy-density} から：

\begin{equation}\protect\phantomsection\label{eq-uvalues}{
u_\nu(T) = g(\nu) \cdot h\nu \cdot \bar{n}(\nu, T) = \frac{8\pi\nu^2}{c^3} \cdot h\nu \cdot \frac{1}{e^{h\nu/k_BT}-1}
}\end{equation}

整理すると

\begin{equation}\protect\phantomsection\label{eq-planck-energy-final}{
u_\nu(T) = \frac{8\pi h\nu^3}{c^3} \cdot \frac{1}{e^{h\nu/k_BT}-1}
}\end{equation}

比強度に変換する（等方放射場では \(B_\nu = c\, u_\nu / 4\pi\)、第4章
§4.3）：

\begin{equation}\protect\phantomsection\label{eq-planck-fully-derived}{
\boxed{
B_\nu(T) = \frac{2 h \nu^3}{c^2} \cdot \frac{1}{e^{h\nu/k_BT}-1}
}
}\end{equation}

\textbf{プランク分布が、観測から逆向きに完全に説明された。}
各因子の起源を一覧にまとめる：

{\def\LTcaptype{none} % do not increment counter
\begin{longtable}[]{@{}
  >{\raggedright\arraybackslash}p{(\linewidth - 4\tabcolsep) * \real{0.3333}}
  >{\raggedright\arraybackslash}p{(\linewidth - 4\tabcolsep) * \real{0.3333}}
  >{\raggedright\arraybackslash}p{(\linewidth - 4\tabcolsep) * \real{0.3333}}@{}}
\toprule\noalign{}
\begin{minipage}[b]{\linewidth}\raggedright
因子
\end{minipage} & \begin{minipage}[b]{\linewidth}\raggedright
起源
\end{minipage} & \begin{minipage}[b]{\linewidth}\raggedright
章
\end{minipage} \\
\midrule\noalign{}
\endhead
\bottomrule\noalign{}
\endlastfoot
\(1/(e^{h\nu/k_BT}-1)\) & 光子の Bose 統計＋ \(\mu = 0\) & 第8章 \\
\(h\nu\)（指数中） & エネルギー量子化（Planck の仮説） & 第10章 \\
\(h\nu\)（前因子中） & 各モードの光子エネルギー & 第10章 \\
\(\nu^2\)（前因子中） & 三次元空洞のモード密度 & \textbf{本章} \\
\(1/c^3\)（モード密度から） & Maxwell 方程式の分散関係 & 第13章 \\
\(2\)（前因子の先頭） & 偏光自由度（横波） & 第13章 \\
全体 \(1/c^2\)（\(1/c^3 \cdot c\) から） & 光速度 ＋ \(u_\nu \to B_\nu\)
変換 & 本章 \\
\end{longtable}
}

これで第I〜V部の逆引きが完結する。

\textbf{中心地図の各因子 ―
初出・借用・導出の往復}：本書では、いくつかの因子を\textbf{先に結果だけ借りて}おき、後の章で\textbf{本格的に導出して借用を返す}という往復構造をとった（第8章・第9章で予告した一覧がこれである）。

{\def\LTcaptype{none} % do not increment counter
\begin{longtable}[]{@{}
  >{\raggedright\arraybackslash}p{(\linewidth - 6\tabcolsep) * \real{0.2500}}
  >{\raggedright\arraybackslash}p{(\linewidth - 6\tabcolsep) * \real{0.2500}}
  >{\raggedright\arraybackslash}p{(\linewidth - 6\tabcolsep) * \real{0.2500}}
  >{\raggedright\arraybackslash}p{(\linewidth - 6\tabcolsep) * \real{0.2500}}@{}}
\toprule\noalign{}
\begin{minipage}[b]{\linewidth}\raggedright
因子
\end{minipage} & \begin{minipage}[b]{\linewidth}\raggedright
初出（仮に使う）
\end{minipage} & \begin{minipage}[b]{\linewidth}\raggedright
仮に借りる場所
\end{minipage} & \begin{minipage}[b]{\linewidth}\raggedright
本当に導出する場所
\end{minipage} \\
\midrule\noalign{}
\endhead
\bottomrule\noalign{}
\endlastfoot
占有因子 \(1/(e^{h\nu/k_BT}-1)\) & 第8章 & ―（借用なし） & 第8章（Bose
統計＋\(\mu=0\)） \\
エネルギー量子 \(h\nu\) & 第IV部 & 第8章（量子化を仮定して使用） &
第10章（Planck の量子仮説） \\
モード密度 \(8\pi\nu^2/c^3\) & 第8章 & 第8章
式~\ref{eq-photon-energy-density}・第9章
式~\ref{eq-mode-density-borrowed} & \textbf{本章（第14章）} \\
偏光因子 \(2\) & 第13章 & ―（借用なし） & 第13章（横波の2自由度） \\
\(u_\nu \to B_\nu\) 変換（\(c/4\pi\)） & 第4章 & ―（借用なし） & 第4章
§4.3 \\
\end{longtable}
}

この往復こそが、本書が「天下りを避ける」と述べた逆引きの実装である。モード密度だけが第8・9章で借用され、本章で返済された。

\begin{tcolorbox}[enhanced jigsaw, arc=.35mm, bottomrule=.15mm, bottomtitle=1mm, breakable, colback=white, colbacktitle=quarto-callout-tip-color!10!white, colframe=quarto-callout-tip-color-frame, coltitle=black, left=2mm, leftrule=.75mm, opacityback=0, opacitybacktitle=0.6, rightrule=.15mm, title=\textcolor{quarto-callout-tip-color}{\faLightbulb}\hspace{0.5em}{ヒント}, titlerule=0mm, toprule=.15mm, toptitle=1mm]

\textbf{観測との対応}：観測される黒体スペクトル ― CMB の 2.725 K
の完璧な黒体形、太陽の \(T \simeq 5800\) K
の連続光、降着円盤の多温度合成、ダストの \(30\) K 修正黒体 ―
これらすべてが、上の表のすべての物理（電磁気・量子・統計）の\textbf{同時の結果}であることが、本章までの逆引きで明らかになった。\textbf{観測スペクトルの一本の曲線}から、現代物理学の主要分野が逆向きに再構成された。これが本書の方法論の到達点である。

\end{tcolorbox}

\section{14.6 空洞放射からの「展望」 ―
第VI〜VIII部への橋}\label{sec-cavity-modes-bridge}

{[}本文目安：B3-B4{]}

第I〜V部で確立した中心地図プランク分布を出発点として、本書は次の三方向に進む。

\textbf{第VI部（現実の宇宙へ戻る）}：プランク分布は理想的な完全熱平衡の式。現実の天体スペクトルがどう・なぜずれるかを扱う。修正黒体（ダスト）、灰色大気近似、多温度成分、コンプトン化、非熱的放射の混入
― 観測スペクトルの解読に必要なすべての「ずれ」の物理。

\textbf{第VII部（線スペクトル）}：本書のもう一つの主柱。観測される線スペクトルの逆引き
― 原子物理、線形成、線形、観測応用 ―
を、プランク分布と平行な構造で展開する。Einstein
係数（第11章）が二つの地図を結ぶ要として再登場する。

\textbf{第VIII部（QED の見通し）}：本章で行った「モード密度 ×
量子振動子」を本格的に再定式化する。各電磁モードが\textbf{量子化された調和振動子}であり、励起準位が光子数。これにより本書の二つの中心地図（プランク関数と水素線吸収係数）が、同じ電磁場--物質相互作用の二側面として統合される。本章で借用していた
\textbf{「ν² モード密度 × \(h\nu\) エネルギー量子」}
という組合せが、第25章で「電磁場の量子化」として整理し直される。

\begin{tcolorbox}[enhanced jigsaw, arc=.35mm, bottomrule=.15mm, bottomtitle=1mm, breakable, colback=white, colbacktitle=quarto-callout-note-color!10!white, colframe=quarto-callout-note-color-frame, coltitle=black, left=2mm, leftrule=.75mm, opacityback=0, opacitybacktitle=0.6, rightrule=.15mm, title=\textcolor{quarto-callout-note-color}{\faInfo}\hspace{0.5em}{ノート}, titlerule=0mm, toprule=.15mm, toptitle=1mm]

\textbf{使えるようになった道具}（§14.1〜14.6）：

\begin{itemize}
\tightlist
\item
  空洞内の電磁定在波を境界条件のもとで数え上げる方法（k 空間の格子）
\item
  モード密度 \(g(\nu) = 8\pi\nu^2/c^3\) ―
  第8〜10章で借用していた式が完全に説明された
\item
  プランク分布の\textbf{すべての因子}が、観測から逆向きに既知の物理（電磁気＋統計力学＋量子論）で説明されたという事実
\item
  中心地図プランク関数を「観測の言葉」と「物理の言葉」のあいだの橋として完全に理解する
\end{itemize}

\end{tcolorbox}

\begin{center}\rule{0.5\linewidth}{0.5pt}\end{center}

\section*{この章で何がわかったか}\label{sec-cavity-modes-summary}
\addcontentsline{toc}{section}{この章で何がわかったか}

\markright{この章で何がわかったか}

\textbf{中心地図に戻る・第I〜V部の総まとめ}

本章で、プランク分布の前因子 \(2h\nu^3/c^2\) の最後のピース ― モード密度
\(\nu^2\) ―
が逆引きされた。これで中心地図プランク分布の\textbf{すべての因子}が、観測から逆向きに既知の物理で説明された。

第I部で「観測される黒体スペクトル」として現象を確認し、第II部で放射場の記述を整え、第III部で熱平衡と
Bose 統計から指数関数因子を、第IV部でエネルギー量子化から \(h\nu\)
を、第V部で電磁モード密度から \(\nu^2/c^3\)
を導出した。観測の一本の曲線から、現代物理学の主要分野が逆向きに再構成されたことになる。

\textbf{次部以降}：これで第I〜V部の「連続スペクトル（黒体放射）の逆引き」は完結する。

\begin{itemize}
\tightlist
\item
  \textbf{第VI部} では、現実の天体でこの理想からどうずれるかを扱う
\item
  \textbf{第VII部}
  では、もう一つの主柱\textbf{線スペクトル}の逆引きを展開する
\item
  \textbf{第VIII部} では、QED で連続と線の両方の地図を統一する
\end{itemize}

\section*{演習問題}\label{sec-cavity-modes-exercises}
\addcontentsline{toc}{section}{演習問題}

\markright{演習問題}

\textbf{模範解答}
は折りたたみを展開して確認できる（オンライン版）。まず自力で解いてから開くこと。

\subsection*{問題 14-1　定在波と境界条件からモード密度
8πν²/c³}\label{ex-14-1}
\addcontentsline{toc}{subsection}{問題
14-1　定在波と境界条件からモード密度 8πν²/c³}

{[}★ 難易度：☆☆☆ {]}

§14.1〜14.4 の数え上げを通しで再現する。辺長 \(L\)
の完全導体立方体。定在波 \(E_x\propto\cos(k_xx)\sin(k_yy)\sin(k_zz)\)
など。

\begin{enumerate}
\def\labelenumi{\arabic{enumi}.}
\tightlist
\item
  境界条件 \(E_\parallel=0\) から許される波数が
  \(k_i=n_i\pi/L\)（\(n_i=1,2,\dots\)）になることを示せ。
\item
  \(k\) 空間の格子点密度 \(V/\pi^3\)、第一象限（1/8
  球）、偏光2を使い、\(|\mathbf k|<k\) のモード数
  \(N(<k)=\dfrac{Vk^3}{3\pi^2}\) を導け。
\item
  分散関係 \(k=2\pi\nu/c\) を代入し
  \(N(<\nu)=\dfrac{8\pi V\nu^3}{3c^3}\)、微分して
  \(g(\nu)=\dfrac{1}{V}\dfrac{dN}{d\nu}=\dfrac{8\pi\nu^2}{c^3}\)
  を導け。
\end{enumerate}

\textbf{関連}：\hyperref[sec-cavity-modes-standing-waves]{§14.1
定在波}、\hyperref[sec-cavity-modes-frequency]{§14.4
モード密度}／同じ式を第9章側から導いたのが\hyperref[ex-9-1]{演習 9-1}。

\section{模範解答（問題
14-1）}\label{ux6a21ux7bc4ux89e3ux7b54ux554fux984c-14-1}

\textbf{(1)} 壁 \(y=0,L\) と \(z=0,L\) で \(E_x\)
の接線成分がゼロ。\(E_x\propto\sin(k_yy)\sin(k_zz)\) が \(y,z=0\)
で自動的にゼロ、\(y,z=L\) でゼロには
\(\sin(k_iL)=0\Rightarrow k_iL=n_i\pi\)。よって
\(k_i=n_i\pi/L\)（\(n_i=1,2,\dots\)、各軸独立）。

\textbf{(2)} 格子間隔 \(\pi/L\) ゆえ \(k\) 空間の格子点密度
\(1/(\pi/L)^3=V/\pi^3\)（\(V=L^3\)）。\(n_i>0\) なので半径 \(k\)
の球の第一象限（1/8）のみ：\(\frac18\cdot\frac{4\pi k^3}{3}\)。偏光2を掛けて

\[
N(<k)=2\cdot\frac{V}{\pi^3}\cdot\frac18\cdot\frac{4\pi k^3}{3}=\frac{Vk^3}{3\pi^2}.
\]

\textbf{(3)} \(k=2\pi\nu/c\)：

\[
N(<\nu)=\frac{V}{3\pi^2}\left(\frac{2\pi\nu}{c}\right)^3=\frac{8\pi V\nu^3}{3c^3},\qquad g(\nu)=\frac1V\frac{dN}{d\nu}=\frac{8\pi\nu^2}{c^3}.
\]

\textbf{答え}：\(k_i=n_i\pi/L\)、\(N(<k)=Vk^3/3\pi^2\)、\(g(\nu)=8\pi\nu^2/c^3\)。第9章の数え上げ（演習
9-1）と一致。

\subsection*{問題 14-2　なぜ「1/8 倍」か ― ±k の同一視}\label{ex-14-2}
\addcontentsline{toc}{subsection}{問題 14-2　なぜ「1/8 倍」か ― ±k
の同一視}

{[}★ 難易度：☆☆ {]}

§14.2 の問いを掘る。全球 \(4\pi k^3/3\) をそのまま使うと過大評価になる。

\begin{enumerate}
\def\labelenumi{\arabic{enumi}.}
\tightlist
\item
  なぜ許される \(\mathbf k\)
  は第一象限（\(n_x,n_y,n_z>0\)）だけなのか。\(\mathbf k\) と
  \(-\mathbf k\)
  が同じ定在波を表すこと（\(\cos(k_xx)=\cos(-k_xx)\)）を使って説明せよ。
\item
  周期境界条件（壁でなく端をつなぐ）にすると格子間隔が
  \(2\pi/L\)、\(\mathbf k\) は正負両方をとる。この場合も同じ \(g(\nu)\)
  になることを、両効果（間隔 \(2\) 倍＝密度
  \(1/8\)、全球を使う＝\(\times8\)）が打ち消すことから示せ。
\item
  「壁にどう接続するか」が最終結果に効かないことが、第3章の黒体放射の普遍性とどう整合するか述べよ。
\end{enumerate}

\textbf{関連}：\hyperref[sec-cavity-modes-k-counting]{§14.2
波数の数え上げ}／普遍性は\hyperref[sec-universality-cavity]{§3.3}（\hyperref[ex-3-1]{演習
3-1}）。

\section{模範解答（問題
14-2）}\label{ux6a21ux7bc4ux89e3ux7b54ux554fux984c-14-2}

\textbf{(1)} 定在波は \(\cos(k_xx)\) や \(\sin(k_xx)\) の形で、\(k_x\)
と \(-k_x\)
は同じ関数（偶／奇の符号差は振幅に吸収）を与え、物理的に同一のモード。\(n_i\)
を正負両方数えると同じモードを重複して数えることになる。独立なモードは
\(n_i>0\) の第一象限のみ。よって全球の 1/8。

\textbf{(2)} 周期境界条件では許される
\(k_i=2\pi n_i/L\)（\(n_i=0,\pm1,\pm2,\dots\)）で格子間隔
\(2\pi/L\)（壁の場合の2倍）。格子点密度は \((L/2\pi)^3=V/8\pi^3\)（壁の
\(V/\pi^3\) の 1/8）。一方 \(\mathbf k\) は正負両方とるので全球
\(4\pi k^3/3\)（1/8 球の8倍）を使う。密度の \(1/8\) と全球の \(\times8\)
が打ち消し、\(N(<k)=2\cdot\frac{V}{8\pi^3}\cdot\frac{4\pi k^3}{3}=\frac{Vk^3}{3\pi^2}\)
と同じ。

\textbf{(3)}
第3章では「空洞放射のスペクトルは壁の材質によらず一意」（熱力学第二法則）と示した。境界条件の選び方（完全導体壁
vs 周期境界）が \(g(\nu)\)
を変えないのは、まさにこの普遍性の現れ。計算の便宜でどちらを選んでもよく、これは普遍性を逆手にとった計算戦略。

\textbf{答え}：\(\pm\mathbf k\) が同一モードゆえ 1/8
球。周期境界では密度 1/8・全球 ×8 が打ち消し同じ
\(g(\nu)\)。境界の選び方に依らないのは黒体の普遍性の現れ。

\subsection*{問題 14-3　モード密度の次元依存性（d 次元）}\label{ex-14-3}
\addcontentsline{toc}{subsection}{問題 14-3　モード密度の次元依存性（d
次元）}

{[}★ 難易度：☆☆ {]}

§14.4 は3次元で \(g(\nu)\propto\nu^2\)
になる理由が幾何にあると述べた。次元を一般化する。

\begin{enumerate}
\def\labelenumi{\arabic{enumi}.}
\tightlist
\item
  \(d\) 次元で \(k\) 空間の球の体積 \(\propto k^d\)
  から、\(|\mathbf k|<k\) のモード数 \(N(<k)\propto k^d\)、よって
  \(g(k)\propto k^{d-1}\)、\(g(\nu)\propto\nu^{d-1}/c^d\) を示せ。
\item
  3次元・2次元・1次元でモード密度の振動数依存を書き、それぞれ
  \(\nu^2\)・\(\nu\)・定数になることを確かめよ。
\item
  もし宇宙が2次元なら、プランク分布のエネルギー密度
  \(u_\nu=g(\nu)h\nu\bar n\) は \(\nu\)
  の何乗の前因子を持つか。「観測される黒体スペクトルの形に3次元という情報が組み込まれている」という主張を説明せよ。
\end{enumerate}

\textbf{関連}：\hyperref[sec-cavity-modes-frequency]{§14.4
モード密度の正体}／モード密度の幾何は\hyperref[ex-9-1]{演習 9-1}。

\section{模範解答（問題
14-3）}\label{ux6a21ux7bc4ux89e3ux7b54ux554fux984c-14-3}

\textbf{(1)} \(d\) 次元の球の体積は \(\propto k^d\)。格子密度
\(\propto V/\pi^d\)、象限因子・偏光因子は定数なので
\(N(<k)\propto Vk^d\)。微分して
\(dN/dk\propto k^{d-1}\)。\(k=2\pi\nu/c\) より
\(g(\nu)\propto\nu^{d-1}/c^d\)。

\textbf{(2)}
3次元：\(g\propto\nu^2/c^3\)。2次元：\(g\propto\nu/c^2\)。1次元：\(g\propto\nu^0=\)
定数（\(\propto1/c\)）。これは「\(k\) 空間の球殻
\(k^{d-1}\)」の幾何そのもの。

\textbf{(3)} 2次元では
\(u_\nu=g(\nu)h\nu\bar n\propto\nu\cdot h\nu/(e^{h\nu/kT}-1)=h\nu^2/(e^{h\nu/kT}-1)\)、前因子は
\(\nu^2\)（3次元の \(\nu^3\)
ではない）。スペクトル全体の形（低振動数べき・ピーク位置・全エネルギーの温度依存）が次元で変わる。逆に言えば、観測される黒体スペクトルが
\(\nu^3\)
前因子を持つこと自体が、我々が3次元空間に住む証拠を含んでいる。

\textbf{答え}：\(g(\nu)\propto\nu^{d-1}/c^d\)。3D→\(\nu^2\)、2D→\(\nu\)、1D→定数。2次元なら
\(u_\nu\) 前因子は \(\nu^2\)。スペクトル形が空間次元を反映。

\subsection*{問題 14-4　プランク分布の全因子の合流 ―
第I〜V部の総まとめ}\label{ex-14-4}
\addcontentsline{toc}{subsection}{問題 14-4　プランク分布の全因子の合流
― 第I〜V部の総まとめ}

{[}★ 難易度：☆☆☆ {]}

§14.5 の総まとめを自分で組み立てる。モード密度
\(g(\nu)=8\pi\nu^2/c^3\)、各モードの光子エネルギー \(h\nu\)、平均光子数
\(\bar n=1/(e^{h\nu/kT}-1)\)。

\begin{enumerate}
\def\labelenumi{\arabic{enumi}.}
\tightlist
\item
  エネルギー密度 \(u_\nu=g(\nu)\cdot h\nu\cdot\bar n\)
  を組み立て、\(u_\nu=\dfrac{8\pi h\nu^3}{c^3}\dfrac{1}{e^{h\nu/kT}-1}\)
  を示せ。
\item
  等方場の関係
  \(B_\nu=\dfrac{c}{4\pi}u_\nu\)（第4章）を使い、\(B_\nu=\dfrac{2h\nu^3}{c^2}\dfrac{1}{e^{h\nu/kT}-1}\)
  を導け。
\item
  前因子 \(2h\nu^3/c^2\) と指数因子 \(1/(e^{h\nu/kT}-1)\)
  の各部品（2、\(\nu^2\)、\(1/c^3\)、\(h\nu\)、Bose＋\(\mu=0\)）が、それぞれどの章・どの物理に由来するかを一覧にせよ。
\end{enumerate}

\textbf{関連}：\hyperref[sec-cavity-modes-bnu-derivation]{§14.5
前因子への合流}／占有因子は\hyperref[ex-8-2]{演習
8-2}・\hyperref[ex-8-3]{8-3}、\(h\nu\) 量子化は\hyperref[ex-10-1]{演習
10-1}、\(B_\nu=\frac{c}{4\pi}u_\nu\) は\hyperref[ex-4-1]{演習 4-1}。

\section{模範解答（問題
14-4）}\label{ux6a21ux7bc4ux89e3ux7b54ux554fux984c-14-4}

\textbf{(1)}
\(u_\nu=g(\nu)\cdot h\nu\cdot\bar n=\dfrac{8\pi\nu^2}{c^3}\cdot h\nu\cdot\dfrac{1}{e^{h\nu/kT}-1}=\dfrac{8\pi h\nu^3}{c^3}\dfrac{1}{e^{h\nu/kT}-1}\)。

\textbf{(2)}
\(B_\nu=\dfrac{c}{4\pi}u_\nu=\dfrac{c}{4\pi}\cdot\dfrac{8\pi h\nu^3}{c^3}\dfrac{1}{e^{h\nu/kT}-1}=\dfrac{2h\nu^3}{c^2}\dfrac{1}{e^{h\nu/kT}-1}\)。（\(\frac{c}{4\pi}\cdot8\pi=2c\)、\(/c^3=1/c^2\)。）

\textbf{(3)} 各因子の起源：

{\def\LTcaptype{none} % do not increment counter
\begin{longtable}[]{@{}
  >{\raggedright\arraybackslash}p{(\linewidth - 4\tabcolsep) * \real{0.3333}}
  >{\raggedright\arraybackslash}p{(\linewidth - 4\tabcolsep) * \real{0.3333}}
  >{\raggedright\arraybackslash}p{(\linewidth - 4\tabcolsep) * \real{0.3333}}@{}}
\toprule\noalign{}
\begin{minipage}[b]{\linewidth}\raggedright
因子
\end{minipage} & \begin{minipage}[b]{\linewidth}\raggedright
起源
\end{minipage} & \begin{minipage}[b]{\linewidth}\raggedright
章/演習
\end{minipage} \\
\midrule\noalign{}
\endhead
\bottomrule\noalign{}
\endlastfoot
\(1/(e^{h\nu/kT}-1)\) & 光子の Bose 統計＋\(\mu=0\) &
第8章（8-2,8-3） \\
\(h\nu\)（指数中・前因子中） & エネルギー量子化 & 第10章（10-1） \\
\(\nu^2\) & 3次元空洞のモード密度 & 第14章（14-1） \\
\(1/c^3\) & 分散関係 \(k=2\pi\nu/c\) & 第13章（13-2） \\
\(2\) & 偏光自由度（横波） & 第13章（13-2,13-4） \\
全体 \(1/c^2\) & \(1/c^3\times c\)（\(u_\nu\to B_\nu\) 変換） &
第14章/第4章（4-1） \\
\end{longtable}
}

観測される一本の黒体曲線から、統計力学・量子論・電磁気学が逆向きに再構成された。

\textbf{答え}：\(u_\nu=\frac{8\pi h\nu^3}{c^3}\frac{1}{e^{h\nu/kT}-1}\)、\(B_\nu=\frac{2h\nu^3}{c^2}\frac{1}{e^{h\nu/kT}-1}\)。全因子が第8・10・13・14章の物理に分解される。

\section*{さらに学ぶための参考文献}\label{sec-cavity-modes-further}
\addcontentsline{toc}{section}{さらに学ぶための参考文献}

\markright{さらに学ぶための参考文献}

\begin{itemize}
\tightlist
\item
  Reif, \emph{Fundamentals of Statistical and Thermal Physics}
  (McGraw-Hill, 1965) --- 空洞放射と状態数のクラシック
\item
  Pathria \& Beale, \emph{Statistical Mechanics} (Academic Press, 3rd
  ed., 2011) --- モード密度の体系的扱い
\item
  Mandl, \emph{Statistical Physics} (Wiley, 2nd ed., 1988) ---
  学部向け、空洞モードと黒体放射の標準的扱い
\item
  Loudon, \emph{The Quantum Theory of Light} (Oxford, 3rd ed., 2000) ---
  第1〜2章で電磁モードと光子の対応を体系的に
\end{itemize}

\part{第VI部　現実の宇宙へ戻る}

\chapter{第15章　現実の天体は完全な黒体ではない}\label{ux7b2c15ux7ae0-ux73feux5b9fux306eux5929ux4f53ux306fux5b8cux5168ux306aux9ed2ux4f53ux3067ux306fux306aux3044}

\textbf{この章の主題}：第I〜V部で逆引きしてきたプランク分布 \(B_\nu(T)\)
は、\textbf{完全熱平衡}という理想状態の結果である。現実の天体は、熱平衡から少しずつ・あるいは大きくずれている。本章は第VI部の総論として、\textbf{「ずれを引き起こす物理機構」を体系的に分類}
する。第16〜18章で個別の天体に踏み込む前に、すべての応用で使う共通の道具立てをここで揃える。

\section*{この章の中心地図}\label{sec-non-blackbody-map}
\addcontentsline{toc}{section}{この章の中心地図}

\markright{この章の中心地図}

\textbf{中心地図：プランク関数 ― 本部の主題は「ずれ」}

\[
B_\nu(T) = \frac{2 h \nu^3}{c^2} \cdot \frac{1}{e^{h\nu/kT} - 1} \quad \overset{\text{現実}}{\longrightarrow} \quad \text{?}
\]

第VI部（現実の宇宙へ戻る）は、観測される現実の天体スペクトルが、この理想的プランク関数からどのように、そしてなぜ
\textbf{ずれる}かを扱う。ずれの主な起源：

\begin{itemize}
\tightlist
\item
  \textbf{振動数依存不透明度} → 修正黒体（第15章）
\item
  \textbf{多温度成分の重ね合わせ} → 多温度黒体スペクトル（第17・18章）
\item
  \textbf{散乱効果}（コンプトン化など）→ 高エネルギー側の変形（第18章）
\item
  \textbf{非熱的放射の混入} → シンクロトロン・ベキ則の重なり
\end{itemize}

本部は、観測スペクトルを「理想的黒体＋ずれ」として読み解き、ずれの定量化から物理状態を引き出す視点を養う。

\begin{tcolorbox}[enhanced jigsaw, arc=.35mm, bottomrule=.15mm, bottomtitle=1mm, breakable, colback=white, colbacktitle=quarto-callout-note-color!10!white, colframe=quarto-callout-note-color-frame, coltitle=black, left=2mm, leftrule=.75mm, opacityback=0, opacitybacktitle=0.6, rightrule=.15mm, title=\textcolor{quarto-callout-note-color}{\faInfo}\hspace{0.5em}{ノート}, titlerule=0mm, toprule=.15mm, toptitle=1mm]

\textbf{方針}：本部の中心地図は「ずれ」を主題化したプランク分布。本章で
\texttt{(deviation)}
因子の中身を整理し、第16〜18章では具体的な天体ごとに、それぞれの
\(f_\text{dev}\) がどう書けるかを見る。

\end{tcolorbox}

\section*{この章で答える問い}\label{sec-non-blackbody-questions}
\addcontentsline{toc}{section}{この章で答える問い}

\markright{この章で答える問い}

\begin{itemize}
\tightlist
\item
  なぜ恒星のスペクトルには深い吸収線が刻まれているのに、全体としては黒体に近いのか
\item
  灰色大気近似はどのくらい良い近似で、どこから破綻するのか
\item
  振動数依存の不透明度は、観測スペクトルの形をどう歪めるのか
\item
  「修正黒体」とは何で、なぜダストのスペクトルがそう呼ばれるのか
\item
  非熱的放射（シンクロトロンなど）は熱的放射からどう区別するのか
\end{itemize}

\section*{到達目標}\label{sec-non-blackbody-goals}
\addcontentsline{toc}{section}{到達目標}

\markright{到達目標}

この章を読み終えると、読者は次のことができるようになる：

\begin{itemize}
\tightlist
\item
  黒体スペクトルからのずれを引き起こす主要な物理機構を分類できる
\item
  灰色大気近似と振動数依存不透明度の違いを、観測との対応で説明できる
\item
  修正黒体スペクトルの形を、ダストの放射効率から導ける
\item
  熱的放射と非熱的放射を、スペクトル形状から見分ける手段を理解する
\end{itemize}

\begin{center}\rule{0.5\linewidth}{0.5pt}\end{center}

\section{15.1 ずれの全体像 ―
完全熱平衡からの三段階}\label{sec-non-blackbody-overview}

{[}本文目安：B3{]}

第I〜V部で示したように、プランク分布 \(B_\nu(T)\) が成り立つには次の
\textbf{三つの条件} すべてが必要だった：

\begin{enumerate}
\def\labelenumi{\arabic{enumi}.}
\tightlist
\item
  \textbf{物質と放射が完全熱平衡} にある（共通の温度 \(T\)）
\item
  \textbf{光子の化学ポテンシャルがゼロ}（光子数が自由に変動できる、第8章）
\item
  \textbf{すべての振動数で十分な不透明度}（外と切れていて熱化が完了している）
\end{enumerate}

現実の天体は、ほぼ常にこの三つのうち何かが緩い。緩み方には次の階層がある：

\begin{tcolorbox}[enhanced jigsaw, arc=.35mm, bottomrule=.15mm, bottomtitle=1mm, breakable, colback=white, colbacktitle=quarto-callout-note-color!10!white, colframe=quarto-callout-note-color-frame, coltitle=black, left=2mm, leftrule=.75mm, opacityback=0, opacitybacktitle=0.6, rightrule=.15mm, title=\textcolor{quarto-callout-note-color}{\faInfo}\hspace{0.5em}{ノート}, titlerule=0mm, toprule=.15mm, toptitle=1mm]

\textbf{用語：完全熱平衡 → LTE → non-LTE → 非熱的}（再掲）

\begin{itemize}
\tightlist
\item
  \textbf{完全熱平衡（CTE; complete thermal
  equilibrium）}：上記三条件すべて。CMB のような特殊な系のみ
\item
  \textbf{局所熱平衡（LTE; local thermal
  equilibrium）}：物質側だけが共通温度。放射場は自由に流れる（第6章
  §6.5）
\item
  \textbf{non-LTE}：物質の占有数が Boltzmann/Saha
  でも書けない。放射場が物質と独立に流れる
\item
  \textbf{非熱的（non-thermal）}：粒子の \textbf{エネルギー分布} が
  Maxwell-Boltzmann ですらない（べき則など）
\end{itemize}

本書の主たる射程は CTE〜LTE。non-LTE
は第VII部で個別に扱い、非熱的放射は本章 §15.6 で区別の方法だけを述べる

\end{tcolorbox}

これら四段階に対応して、観測スペクトルは

\begin{equation}\protect\phantomsection\label{eq-deviation-factor}{
I_\nu^\text{obs} = B_\nu(T) \cdot f_\text{dev}(\nu, T, \text{物理量})
}\end{equation}

と書けることが多い。\(f_\text{dev}\) が
\textbf{何でずれるか・どれくらいずれるか}
を読み解くのが本章以降の主題である。

本章および第VI部の応用章を貫く逆引きの流れを、あらかじめ一表にまとめておく：

{\def\LTcaptype{none} % do not increment counter
\begin{longtable}[]{@{}
  >{\raggedright\arraybackslash}p{(\linewidth - 6\tabcolsep) * \real{0.2500}}
  >{\raggedright\arraybackslash}p{(\linewidth - 6\tabcolsep) * \real{0.2500}}
  >{\raggedright\arraybackslash}p{(\linewidth - 6\tabcolsep) * \real{0.2500}}
  >{\raggedright\arraybackslash}p{(\linewidth - 6\tabcolsep) * \real{0.2500}}@{}}
\toprule\noalign{}
\begin{minipage}[b]{\linewidth}\raggedright
観測で見える特徴
\end{minipage} & \begin{minipage}[b]{\linewidth}\raggedright
最小モデル
\end{minipage} & \begin{minipage}[b]{\linewidth}\raggedright
逆引きされる物理量
\end{minipage} & \begin{minipage}[b]{\linewidth}\raggedright
破れる仮定
\end{minipage} \\
\midrule\noalign{}
\endhead
\bottomrule\noalign{}
\endlastfoot
連続吸収端・段差 & 振動数依存不透明度 ＋ 温度勾配 & \(\kappa_\nu\)
の振動数依存・\(T(\tau)\) & 灰色不透明度 \\
Rayleigh-Jeans 側の急な勾配 & 修正黒体 \(Q_\nu B_\nu(T)\) &
ダスト温度・放射率指数 \(\beta\) & \(Q_\nu = 1\)（理想黒体） \\
偏光の存在 & 散乱を含むソース関数 & 単一散乱アルベド・散乱体幾何 &
純吸収（LTE 局所性） \\
べき則スペクトル & 非熱的粒子分布 & スペクトル指数 \(\alpha\)・粒子分布
\(p\) & Maxwell-Boltzmann 分布 \\
\end{longtable}
}

\begin{tcolorbox}[enhanced jigsaw, arc=.35mm, bottomrule=.15mm, breakable, colback=white, colframe=quarto-callout-tip-color-frame, left=2mm, leftrule=.75mm, opacityback=0, rightrule=.15mm, toprule=.15mm]
\begin{minipage}[t]{5.5mm}
\textcolor{quarto-callout-tip-color}{\faLightbulb}
\end{minipage}%
\begin{minipage}[t]{\textwidth - 5.5mm}

\textbf{問い}：「黒体スペクトルからのずれ」と聞いて、最初に何を思い浮かべるか？

\textbf{短答}：おそらく吸収線・輝線。だがそれは「ずれ」の一種類にすぎない。連続光の振動数依存・全体的傾き・指数的減衰の鋭さ・低/高振動数側のべき指数
―
これらすべてが「ずれ」の情報を持っている。本章ではまず連続光のずれに焦点を当て、線スペクトルは第VII部で扱う

\end{minipage}%
\end{tcolorbox}

\section{15.2 灰色大気近似 ―
振動数を平均する}\label{sec-non-blackbody-grey}

{[}本文目安：B3{]}

恒星大気から放射輸送を解いて連続スペクトルを得る最も簡単な近似が
\textbf{灰色大気}（\textbf{grey
atmosphere}）。「灰色」とは「振動数によらず吸収率が一定」の意味（\textbf{\(\kappa_\nu = \kappa\)
for all \(\nu\)}）。

放射輸送方程式（第5章）から、灰色大気で得られる \textbf{Eddington 近似}
の解は

\begin{equation}\protect\phantomsection\label{eq-eddington-grey}{
T^4(\tau) = \frac{3}{4} T_\text{eff}^4 \left(\tau + \frac{2}{3}\right)
}\end{equation}

である。ここで \(\tau\) は光学的深さ、\(T_\text{eff}\) は
\textbf{有効温度}（\textbf{effective temperature}）。表面 \(\tau = 0\)
では \(T(0) = (1/2)^{1/4} T_\text{eff} \simeq 0.84\, T_\text{eff}\)
となる。

\textbf{導出の骨格}：放射平衡（各層で吸収＝放出、正味のエネルギー流束
\(F\) が深さによらず一定で
\(F = \sigma T_\text{eff}^4\)）のもとで、放射輸送方程式を立体角の 0
次・1 次モーメントで積分すると、平均強度 \(J\)・フラックス \(H\)・放射圧
\(K\)
に関するモーメント方程式が得られる。灰色（\(\kappa_\nu = \kappa\)）かつ
LTE では \(S = B = (\sigma/\pi) T^4\)
なので、振動数積分した方程式が閉じる。

ここで閉包として \textbf{Eddington 近似}（放射場がほぼ等方で
\(K = J/3\)、すなわち Eddington 因子
\(f = K/J = 1/3\)）を仮定し、表面の流出境界条件を加えると
\(J(\tau) = (3/4\pi)\,F\,(\tau + 2/3)\) が得られる。LTE で
\(J = B = (\sigma/\pi)T^4\) かつ \(F = \sigma T_\text{eff}^4\)
を用いれば、ただちに 式~\ref{eq-eddington-grey} に至る。表面項 \(2/3\)
は流出境界条件（Hopf 関数の Eddington 近似値）から、係数 \(3/4\) は
Eddington 因子 \(1/3\) から来る。詳しい導出は\hyperref[ex-15-1]{演習
15-1}で再現する（より厳密な扱いは Mihalas 1978 §3、Hubeny \& Mihalas
2014 第 17 章を参照）。

\begin{tcolorbox}[enhanced jigsaw, arc=.35mm, bottomrule=.15mm, bottomtitle=1mm, breakable, colback=white, colbacktitle=quarto-callout-note-color!10!white, colframe=quarto-callout-note-color-frame, coltitle=black, left=2mm, leftrule=.75mm, opacityback=0, opacitybacktitle=0.6, rightrule=.15mm, title=\textcolor{quarto-callout-note-color}{\faInfo}\hspace{0.5em}{ノート}, titlerule=0mm, toprule=.15mm, toptitle=1mm]

\textbf{この導出が置く仮定}：

\begin{itemize}
\tightlist
\item
  \textbf{平面平行}：物理量が鉛直 1 次元の光学的深さ \(\tau\)
  のみに依存する。
\item
  \textbf{灰色不透明度}：\(\kappa_\nu = \kappa\)（振動数によらない）。これにより振動数積分が単純化される。
\item
  \textbf{放射平衡}：各層で吸収と放出が釣り合い、正味フラックス
  \(F = \sigma T_\text{eff}^4\) が深さによらず一定。
\item
  \textbf{LTE}：ソース関数が局所のプランク関数に等しい（\(S_\nu = B_\nu(T)\)）。
\item
  \textbf{Eddington 近似}：放射場がほぼ等方で Eddington 因子
  \(K/J = 1/3\)。
\end{itemize}

\end{tcolorbox}

\textbf{灰色大気でのスペクトル}：振動数によらず吸収率が同じため、振動数依存性によるスペクトル形の歪み（次節）を導入しない。LTE
の仮定の下では各層のスペクトルはプランク関数で、観測されるフラックスは振動数ごとに
\(\tau \approx 2/3\)
付近の層の温度のプランク関数で近似される。\textbf{形の歪みは作らないが、現実の恒星にある吸収線・吸収端は当然再現できない}。

\begin{tcolorbox}[enhanced jigsaw, arc=.35mm, bottomrule=.15mm, breakable, colback=white, colframe=quarto-callout-warning-color-frame, left=2mm, leftrule=.75mm, opacityback=0, rightrule=.15mm, toprule=.15mm]
\begin{minipage}[t]{5.5mm}
\textcolor{quarto-callout-warning-color}{\faExclamationTriangle}
\end{minipage}%
\begin{minipage}[t]{\textwidth - 5.5mm}

\textbf{注意（観測）}：灰色大気近似は「形は黒体」と予想するが、実際の恒星スペクトルは形そのものが歪んでいる（特に紫外側で水素のバルマー連続吸収による段差）。これは振動数依存不透明度のため。灰色大気は
\textbf{全光度の正しい目安は与えるが、スペクトル形の正しい記述ではない}

\end{minipage}%
\end{tcolorbox}

\section{15.3 振動数依存不透明度 ―
形が歪む}\label{sec-non-blackbody-frequency-opacity}

{[}本文目安：B3-B4{]}

実際の物質は、振動数によって吸収率が大きく違う。水素原子なら、

\begin{itemize}
\tightlist
\item
  紫外側で Lyman 連続吸収（\(\lambda < 912\) Å）
\item
  可視〜赤外で Balmer・Paschen・Brackett
  連続吸収（順次しきい値が低くなる）
\item
  さらに離散的な吸収線（第VII部の主題）
\end{itemize}

不透明度が振動数で違うとき、観測者が見ているのは
\textbf{「振動数ごとに別の深さ」} での放射である。

\begin{tcolorbox}[enhanced jigsaw, arc=.35mm, bottomrule=.15mm, bottomtitle=1mm, breakable, colback=white, colbacktitle=quarto-callout-note-color!10!white, colframe=quarto-callout-note-color-frame, coltitle=black, left=2mm, leftrule=.75mm, opacityback=0, opacitybacktitle=0.6, rightrule=.15mm, title=\textcolor{quarto-callout-note-color}{\faInfo}\hspace{0.5em}{ノート}, titlerule=0mm, toprule=.15mm, toptitle=1mm]

\textbf{観測との対応}：ある \(\nu\) で吸収率が小さい → 深く透けて見える
→ 温度が高い領域の放射 →
その振動数で観測強度が大きい。逆に吸収率が大きい振動数では浅い表面しか見えない
→ 温度が低い → 観測強度が小さい。

つまり \textbf{振動数依存不透明度 + 大気の温度勾配}
が組み合わさって、観測スペクトルの形を歪める。これが恒星スペクトルの
\textbf{段差（連続吸収端）}・\textbf{ハンプ}・\textbf{くぼみ}
の起源（連続吸収端を作る束縛-自由過程の微視的な起源 ― 光電離断面積
\(\sigma_{\mathrm{bf}}\propto\nu^{-3}\)
としきい値、可視連続不透明度を担う H⁻ ―
は\hyperref[sec-continuum-bound-free]{第19章 §19.3}）

\end{tcolorbox}

形式的に書くと、振動数依存不透明度の下で、観測者方向（鉛直から
\(\theta = \arccos\mu\) の方向）への射出強度は（平面平行・LTE で）

\begin{equation}\protect\phantomsection\label{eq-eddington-barbier-intensity}{
I_\nu(0, \mu) = \int_0^\infty S_\nu(\tau_\nu) e^{-\tau_\nu/\mu} \frac{d\tau_\nu}{\mu} \approx S_\nu(\tau_\nu = \mu)
}\end{equation}

これが \textbf{Eddington-Barbier 近似（射出強度版）}：観測される振動数
\(\nu\) の方向 \(\mu\) への強度は、その振動数で光学的深さが
\(\tau_\nu = \mu\) になる層のソース関数値に近い。

半球積分した \textbf{射出フラックス}（emergent
flux、鉛直外向きフラックス）では別の流儀の Eddington-Barbier
近似が成り立つ：

\begin{equation}\protect\phantomsection\label{eq-eddington-barbier-flux}{
F_\nu(0) = \int_0^\infty 2\pi S_\nu(\tau_\nu) E_2(\tau_\nu) d\tau_\nu \approx \pi S_\nu(\tau_\nu = 2/3)
}\end{equation}

LTE では \(S_\nu = B_\nu(T)\) なので、両式の右辺の \(S_\nu\) は
\(B_\nu(T)\) に置き換わる。

\begin{tcolorbox}[enhanced jigsaw, arc=.35mm, bottomrule=.15mm, breakable, colback=white, colframe=quarto-callout-warning-color-frame, left=2mm, leftrule=.75mm, opacityback=0, rightrule=.15mm, toprule=.15mm]
\begin{minipage}[t]{5.5mm}
\textcolor{quarto-callout-warning-color}{\faExclamationTriangle}
\end{minipage}%
\begin{minipage}[t]{\textwidth - 5.5mm}

\textbf{注意（観測）}：Eddington-Barbier 近似には \textbf{二つの流儀}
がある：

\begin{itemize}
\tightlist
\item
  \textbf{射出強度} \(I_\nu(0, \mu)\) ：\(\tau_\nu = \mu\)（鉛直方向では
  \(\tau_\nu = 1\)）
\item
  \textbf{射出フラックス}（emergent flux）\(F_\nu\) ：\(\tau_\nu = 2/3\)
\end{itemize}

文献によってどちらを「Eddington-Barbier
近似」と呼ぶかが違うので注意。本書は両者を明示して使い分ける

\end{minipage}%
\end{tcolorbox}

\textbf{重要な含意}：観測スペクトルの「形のゆがみ」から、\(\kappa_\nu\)
の振動数依存と \(T(\tau)\)
の構造の両方が逆引きできる。これが恒星大気の現代的解析の核心。

\section{15.4 修正黒体 ―
ダストの粒度効果}\label{sec-non-blackbody-modified}

{[}本文目安：B3{]}

星間ダストや原始惑星系円盤のダストは、おおむね \(T \simeq 10\)〜\(1000\)
K
の温度範囲で熱的放射を放つ。観測されるスペクトルは黒体に似ているが、\textbf{Rayleigh-Jeans
側（低振動数側。以下 RJ 側）で予想より弱い} という特徴がある。

その原因は、粒子サイズが波長より小さくなる帯域での
\textbf{放射効率の低下}。放射効率 \(Q_\nu\) をプランク関数に掛けた形

\begin{equation}\protect\phantomsection\label{eq-modified-blackbody}{
I_\nu^\text{dust} = Q_\nu(a) \cdot B_\nu(T)
}\end{equation}

が \textbf{修正黒体}（\textbf{modified blackbody}; 文脈により
\textbf{greybody} とも呼ぶ）と呼ばれる。以後は「修正黒体」を用いる。

長波長極限では多くの物質で
\(Q_\nu \propto \nu^\beta\)（\(\beta \approx 1\)〜\(2\)）の振る舞いを示し、これを
\textbf{ダスト放射率指数}（emissivity
index）と呼ぶ。観測スペクトルからの \(\beta\)
の決定は、ダスト粒子の組成・大きさ分布の手がかりとなる（第17章 §17.3）。

\begin{tcolorbox}[enhanced jigsaw, arc=.35mm, bottomrule=.15mm, breakable, colback=white, colframe=quarto-callout-tip-color-frame, left=2mm, leftrule=.75mm, opacityback=0, rightrule=.15mm, toprule=.15mm]
\begin{minipage}[t]{5.5mm}
\textcolor{quarto-callout-tip-color}{\faLightbulb}
\end{minipage}%
\begin{minipage}[t]{\textwidth - 5.5mm}

\textbf{問い}：\(\beta = 2\) のとき、Rayleigh-Jeans 側で観測される
\(I_\nu\) は \(\nu\) の何乗に比例するか？

\textbf{短答}：黒体の Rayleigh-Jeans 側は \(B_\nu \propto \nu^2\)（第2章
§2.5）。これに \(Q_\nu \propto \nu^2\)
を掛けるので、\(I_\nu \propto \nu^4\)。

\textbf{もう一歩}：Wien 側（指数で減衰）では、\(Q_\nu\)
の振動数依存はほぼ効かず、依然として指数的減衰が支配する。\textbf{Rayleigh-Jeans
側でのみ \(\beta\) が見える} という非対称性が、観測解析の鍵となる

\end{minipage}%
\end{tcolorbox}

\section{15.5 散乱の効果 ―
再分布と偏光}\label{sec-non-blackbody-scattering}

{[}本文目安：B3-B4{]}

不透明度には吸収と散乱の二種類がある（第5章）。\textbf{散乱}（\textbf{scattering}）は光子を消さず、方向を変えるだけ。

\begin{itemize}
\tightlist
\item
  \textbf{吸収では局所熱平衡}：物質が光を熱に変え、再びプランク関数で放出する
\item
  \textbf{散乱では非局所}：光子が遠くで作られて、ここで方向を変えるだけ
\end{itemize}

散乱が支配的な系では、観測される光は \textbf{「ソース関数（第5章
§5.4）が局所のプランク関数とは大きく違う」}
状態になる。具体的には、ソース関数 \(S_\nu\) は

\begin{equation}\protect\phantomsection\label{eq-source-with-scattering}{
S_\nu = \frac{\kappa_\nu^\text{abs} B_\nu(T) + \sigma_\nu \langle J_\nu \rangle}{\kappa_\nu^\text{abs} + \sigma_\nu}
}\end{equation}

の形になり、散乱割合 \(\sigma_\nu/(\kappa_\nu^\text{abs} + \sigma_\nu)\)
が大きいと局所のプランク関数からの大きな乖離が生じる。

\begin{tcolorbox}[enhanced jigsaw, arc=.35mm, bottomrule=.15mm, bottomtitle=1mm, breakable, colback=white, colbacktitle=quarto-callout-note-color!10!white, colframe=quarto-callout-note-color-frame, coltitle=black, left=2mm, leftrule=.75mm, opacityback=0, opacitybacktitle=0.6, rightrule=.15mm, title=\textcolor{quarto-callout-note-color}{\faInfo}\hspace{0.5em}{ノート}, titlerule=0mm, toprule=.15mm, toptitle=1mm]

\textbf{観測との対応}：散乱の代表的指標は
\textbf{偏光}。熱平衡の放射場は無偏光だが、散乱は偏光を生む。よって、観測される偏光の存在は「散乱が無視できない」ことの直接の証拠。例：CMB
の偏光（再結合期の散乱が源）、太陽コロナの偏光（Thomson
散乱）、惑星間ダストの偏光散乱

\end{tcolorbox}

\section{15.6 非熱的放射との混合 ―
形で見分ける}\label{sec-non-blackbody-nonthermal}

{[}本文目安：B3{]}

ここまでは粒子のエネルギー分布が
Maxwell-Boltzmann（熱的分布）であることを前提していた。しかし宇宙では、粒子のエネルギー分布が
\textbf{べき則}（power-law）になっている場合がある（\textbf{非熱的}）：

\begin{equation}\protect\phantomsection\label{eq-nonthermal-distribution}{
\frac{dn}{dE} \propto E^{-p}
}\end{equation}

このような非熱的電子からの \textbf{シンクロトロン放射} や
\textbf{逆コンプトン放射}
は、観測スペクトルもべき則となる：\(I_\nu \propto \nu^{-\alpha}\)。これらの本格的な扱いは本書の主軸（熱的放射）の射程外であり、Rybicki
\& Lightman、Longair 等の専門書を参照してほしい（第18章 §18.3
でも降着円盤と共存する非熱的成分との関連で軽く触れる）。

\textbf{熱的 vs 非熱的の見分け方}：

{\def\LTcaptype{none} % do not increment counter
\begin{longtable}[]{@{}
  >{\raggedright\arraybackslash}p{(\linewidth - 4\tabcolsep) * \real{0.3333}}
  >{\raggedright\arraybackslash}p{(\linewidth - 4\tabcolsep) * \real{0.3333}}
  >{\raggedright\arraybackslash}p{(\linewidth - 4\tabcolsep) * \real{0.3333}}@{}}
\toprule\noalign{}
\begin{minipage}[b]{\linewidth}\raggedright
観測的特徴
\end{minipage} & \begin{minipage}[b]{\linewidth}\raggedright
熱的（黒体・修正黒体）
\end{minipage} & \begin{minipage}[b]{\linewidth}\raggedright
非熱的（シンクロトロン・逆コンプトン）
\end{minipage} \\
\midrule\noalign{}
\endhead
\bottomrule\noalign{}
\endlastfoot
高振動数側の減衰 & 指数的（Wien 則） & べき則 \\
低振動数側 & \(\nu^2\)（Rayleigh-Jeans） or \(\nu^{2+\beta}\) & べき則
\(\nu^{-\alpha}\) または特殊なターンオーバー \\
偏光 & ふつう無偏光 or 散乱偏光 & 強く偏光（シンクロトロン） \\
振動数帯 & 温度に対応する典型振動数 &
振動数帯は広範囲、しばしば多重スペクトル \\
\end{longtable}
}

\begin{tcolorbox}[enhanced jigsaw, arc=.35mm, bottomrule=.15mm, bottomtitle=1mm, breakable, colback=white, colbacktitle=quarto-callout-tip-color!10!white, colframe=quarto-callout-tip-color-frame, coltitle=black, left=2mm, leftrule=.75mm, opacityback=0, opacitybacktitle=0.6, rightrule=.15mm, title=\textcolor{quarto-callout-tip-color}{\faLightbulb}\hspace{0.5em}{ヒント}, titlerule=0mm, toprule=.15mm, toptitle=1mm]

\textbf{観測との対応}：高エネルギー天体（AGN, GRB, パルサー風星雲）では
\textbf{熱的成分と非熱的成分が共存}
することが多い。両者を分離する観測戦略：

\begin{itemize}
\tightlist
\item
  広帯域 SED（X 線〜電波）で形を見る
\item
  偏光観測
\item
  時間変動（非熱的成分はしばしば早く変動）
\end{itemize}

第18章では降着円盤の標準モデルが多温度黒体スペクトル（熱的）を生み、その上にコロナによるコンプトン化（非熱的に近い）が乗る具体例を扱う

\end{tcolorbox}

\begin{center}\rule{0.5\linewidth}{0.5pt}\end{center}

\section*{この章で何がわかったか}\label{sec-non-blackbody-summary}
\addcontentsline{toc}{section}{この章で何がわかったか}

\markright{この章で何がわかったか}

\textbf{中心地図に戻る}

本章で「中心地図プランク関数 ×
ずれ因子」という構造を導入した。ずれを引き起こす機構は：

\begin{enumerate}
\def\labelenumi{\arabic{enumi}.}
\tightlist
\item
  \textbf{振動数依存不透明度} ―
  スペクトル形の歪み（恒星大気の段差、Eddington-Barbier）
\item
  \textbf{修正黒体} ― ダスト粒度による Rayleigh-Jeans
  側の弱化（\(Q_\nu \propto \nu^\beta\)）
\item
  \textbf{散乱} ― ソース関数の非局所化と偏光の発生
\item
  \textbf{非熱的放射の混入} ― べき則スペクトルとの共存
\end{enumerate}

次章以降では、これらの道具を具体的天体に適用する：第16章 CMB、第17章
恒星・惑星・ダスト、第18章 降着円盤と高エネルギー天体

\section*{使えるようになった道具}\label{sec-non-blackbody-tools}
\addcontentsline{toc}{section}{使えるようになった道具}

\markright{使えるようになった道具}

\begin{tcolorbox}[enhanced jigsaw, arc=.35mm, bottomrule=.15mm, bottomtitle=1mm, breakable, colback=white, colbacktitle=quarto-callout-note-color!10!white, colframe=quarto-callout-note-color-frame, coltitle=black, left=2mm, leftrule=.75mm, opacityback=0, opacitybacktitle=0.6, rightrule=.15mm, title=\textcolor{quarto-callout-note-color}{\faInfo}\hspace{0.5em}{ノート}, titlerule=0mm, toprule=.15mm, toptitle=1mm]

\begin{itemize}
\tightlist
\item
  灰色大気近似と振動数依存不透明度の使い分け
\item
  Eddington-Barbier 近似 ― 観測スペクトルから温度構造を逆引き
\item
  修正黒体スペクトル \(Q_\nu B_\nu(T)\) ― ダスト解析の基本
\item
  散乱を含むソース関数 ― 偏光が「散乱卓越」の直接指標
\item
  熱的・非熱的の SED 上での見分け方
\end{itemize}

\end{tcolorbox}

\section*{演習問題}\label{sec-non-blackbody-exercises}
\addcontentsline{toc}{section}{演習問題}

\markright{演習問題}

\textbf{模範解答}
は折りたたみを展開して確認できる（オンライン版）。まず自力で解いてから開くこと。

\subsection*{問題
15-1　灰色大気の温度構造と「見える深さ」}\label{ex-15-1}
\addcontentsline{toc}{subsection}{問題
15-1　灰色大気の温度構造と「見える深さ」}

{[}★ 難易度：☆☆ {]}

§15.2 は灰色大気の Eddington 近似解
\(T^4(\tau)=\dfrac{3}{4}T_\mathrm{eff}^4(\tau+\dfrac23)\)
を結果として引いた。これを読み解く。

\begin{enumerate}
\def\labelenumi{\arabic{enumi}.}
\tightlist
\item
  表面 \(\tau=0\) での温度 \(T(0)\) を \(T_\mathrm{eff}\)
  で表し、\(T(0)\simeq0.84\,T_\mathrm{eff}\) を確かめよ。
\item
  \(T(\tau)=T_\mathrm{eff}\) となる光学的深さ \(\tau\) を求めよ。
\item
  \begin{enumerate}
  \def\labelenumii{(\arabic{enumii})}
  \setcounter{enumii}{1}
  \tightlist
  \item
    の結果が Eddington-Barbier 近似（射出フラックスは \(\tau_\nu=2/3\)
    層の \(S_\nu\)
    で近似）と整合することを述べ、観測される連続光が「\(\tau\simeq2/3\)
    の層の温度」を映す理由を説明せよ。
  \end{enumerate}
\end{enumerate}

\textbf{関連}：\hyperref[sec-non-blackbody-grey]{§15.2
灰色大気}、\hyperref[sec-non-blackbody-frequency-opacity]{§15.3
Eddington-Barbier} ／ \(\tau\sim1\) と光球は\hyperref[ex-5-4]{第II部
演習 5-4}、有効温度は\hyperref[sec-stars-effective-temperature]{§17.1}。

\section{模範解答（問題
15-1）}\label{ux6a21ux7bc4ux89e3ux7b54ux554fux984c-15-1}

\textbf{(1)}
\(\tau=0\)：\(T^4(0)=\dfrac34 T_\mathrm{eff}^4\cdot\dfrac23=\dfrac12 T_\mathrm{eff}^4\)。よって
\(T(0)=(1/2)^{1/4}T_\mathrm{eff}=0.841\,T_\mathrm{eff}\)。表面は有効温度より約16\%低い。

\textbf{(2)} \(T^4(\tau)=T_\mathrm{eff}^4\)
とおく：\(\dfrac34(\tau+\dfrac23)=1\Rightarrow\tau+\dfrac23=\dfrac43\Rightarrow\tau=\dfrac23\)。

\textbf{(3)} 射出フラックスの Eddington-Barbier 近似
\(F_\nu\approx\pi S_\nu(\tau_\nu=2/3)\) は「\(\tau=2/3\)
層のソース関数が観測フラックスを代表する」。灰色大気ではちょうどこの層で
\(T=T_\mathrm{eff}\)
なので、LTE（\(S_\nu=B_\nu\)）の下で観測される連続光は
\(B_\nu(T_\mathrm{eff})\) に近い。観測スペクトルは「\(\tau\simeq2/3\)
の層」を映している。

\textbf{答え}：\(T(0)=(1/2)^{1/4}T_\mathrm{eff}\simeq0.84T_\mathrm{eff}\)、\(T=T_\mathrm{eff}\)
は \(\tau=2/3\)。これが Eddington-Barbier の見える深さ。

\subsection*{問題 15-2　修正黒体のべき ― なぜ β は Rayleigh-Jeans
側でしか見えないか}\label{ex-15-2}
\addcontentsline{toc}{subsection}{問題 15-2　修正黒体のべき ― なぜ β は
Rayleigh-Jeans 側でしか見えないか}

{[}★ 難易度：☆☆ {]}

§15.4 の修正黒体
\(I_\nu=Q_\nu B_\nu(T)\)、\(Q_\nu\propto\nu^\beta\)（長波長極限）。

\begin{enumerate}
\def\labelenumi{\arabic{enumi}.}
\tightlist
\item
  Rayleigh-Jeans
  側（\(h\nu\ll kT\)、\(B_\nu\propto\nu^2\)）で、\(\beta=2\)
  のとき観測される \(I_\nu\) は \(\nu\) の何乗に比例するか。\(\beta=1\)
  なら？
\item
  Wien 側（\(h\nu\gg kT\)）では \(Q_\nu\propto\nu^\beta\)
  の冪は指数減衰に対してほぼ効かないことを説明せよ。
\item
  「\(\beta\) は Rayleigh-Jeans
  側でのみ見える」という非対称性が、観測でダストの \(\beta\)
  を決めるとき長波長（サブミリ・ミリ）観測を必要とする理由を述べよ。
\end{enumerate}

\textbf{関連}：\hyperref[sec-non-blackbody-modified]{§15.4 修正黒体} ／
両極限は\hyperref[sec-reading-spectra-limits]{§2.5}（\hyperref[ex-10-2]{第IV部
演習
10-2}）、ダスト解析は\hyperref[sec-dust-modified-blackbody]{§17.3}（\hyperref[ex-17-3]{演習
17-3}）。

\section{模範解答（問題
15-2）}\label{ux6a21ux7bc4ux89e3ux7b54ux554fux984c-15-2}

\textbf{(1)} RJ 側
\(B_\nu\propto\nu^2\)。\(I_\nu=Q_\nu B_\nu\propto\nu^\beta\cdot\nu^2=\nu^{2+\beta}\)。\(\beta=2\)
なら \(I_\nu\propto\nu^4\)、\(\beta=1\) なら
\(\nu^3\)。低振動数側の勾配が黒体（\(\nu^2\)）より急になる。

\textbf{(2)} Wien 側 \(B_\nu\propto\nu^3e^{-h\nu/kT}\)
で指数減衰が支配的。\(Q_\nu\propto\nu^\beta\)
の冪は指数関数に比べて緩やかで、スペクトルの落ち方をほとんど変えない。したがって
Wien 側の形からは \(\beta\) をほぼ読めない。

\textbf{(3)} \(\beta\) は RJ 側の勾配 \(\nu^{2+\beta}\)
にしか現れない。ダスト（\(T\sim10\)〜\(30\) K）の RJ
側は遠赤外〜サブミリ・ミリ波に来る。よって \(\beta\)
を測るには、ピークより長波長側のサブミリ・ミリ（Herschel/ALMA）観測点が必要。Wien
側（近赤外）だけでは \(\beta\) が決まらない。

\textbf{答え}：RJ 側 \(I_\nu\propto\nu^{2+\beta}\)（\(\beta=2\) で
\(\nu^4\)）。Wien 側は指数減衰が支配で \(\beta\) が見えない。ゆえに
\(\beta\) 決定には長波長観測が必須。

\subsection*{問題 15-3　散乱を含むソース関数と偏光}\label{ex-15-3}
\addcontentsline{toc}{subsection}{問題 15-3　散乱を含むソース関数と偏光}

{[}★ 難易度：☆☆ {]}

§15.5 のソース関数
\(S_\nu=\dfrac{\kappa_\nu^\mathrm{abs}B_\nu(T)+\sigma_\nu\langle J_\nu\rangle}{\kappa_\nu^\mathrm{abs}+\sigma_\nu}\)
を読み解く。単一散乱アルベド
\(\tilde\omega_\nu\equiv\dfrac{\sigma_\nu}{\kappa_\nu^\mathrm{abs}+\sigma_\nu}\)。

\begin{enumerate}
\def\labelenumi{\arabic{enumi}.}
\tightlist
\item
  \(S_\nu\) を \(\tilde\omega_\nu\)
  で書き直し、\(S_\nu=(1-\tilde\omega_\nu)B_\nu+\tilde\omega_\nu\langle J_\nu\rangle\)
  を示せ。
\item
  純吸収（\(\tilde\omega_\nu=0\)）と純散乱（\(\tilde\omega_\nu=1\)）の極限で
  \(S_\nu\)
  がそれぞれ何になるか述べ、後者で局所プランク関数から外れる理由を述べよ。
\item
  熱平衡放射は無偏光なのに、散乱が偏光を生む。観測される偏光が「散乱が無視できない」ことの直接の証拠になる理由を述べ、例を1つ挙げよ。
\end{enumerate}

\textbf{関連}：\hyperref[sec-non-blackbody-scattering]{§15.5 散乱の効果}
／
ソース関数の基礎は\hyperref[sec-radiative-transfer-source]{§5.4}（\hyperref[sec-radiative-transfer-exercises]{第II部
演習}）、偏光は\hyperref[ex-13-4]{第V部 演習 13-4}。

\section{模範解答（問題
15-3）}\label{ux6a21ux7bc4ux89e3ux7b54ux554fux984c-15-3}

\textbf{(1)} 分子分母を \(\kappa^\mathrm{abs}+\sigma\)
で割る：\(S_\nu=\dfrac{\kappa^\mathrm{abs}}{\kappa^\mathrm{abs}+\sigma}B_\nu+\dfrac{\sigma}{\kappa^\mathrm{abs}+\sigma}\langle J_\nu\rangle=(1-\tilde\omega_\nu)B_\nu+\tilde\omega_\nu\langle J_\nu\rangle\)。

\textbf{(2)} 純吸収
\(\tilde\omega=0\)：\(S_\nu=B_\nu(T)\)（LTE、局所プランク）。純散乱
\(\tilde\omega=1\)：\(S_\nu=\langle J_\nu\rangle\)（平均強度そのもの）。後者は局所の温度
\(T\) ではなく、媒質全体から散乱されてきた放射場
\(\langle J_\nu\rangle\) で決まるので、局所プランク関数
\(B_\nu(T_\mathrm{local})\) から外れる（非局所）。

\textbf{(3)} 熱平衡放射は両偏光が等確率で混ざり無偏光。散乱（特に
Thomson・Rayleigh
散乱）は散乱角に依存して特定偏光を優先的に散乱し、偏光を生む。よって観測される偏光は、その光が散乱を経たこと＝散乱が無視できないことの直接証拠。例：CMB
の偏光（再結合期の Thomson
散乱）、太陽コロナの偏光、惑星間ダストの散乱光。

\textbf{答え}：\(S_\nu=(1-\tilde\omega)B_\nu+\tilde\omega\langle J_\nu\rangle\)。純散乱で
\(S_\nu=\langle J_\nu\rangle\)（非局所）。偏光の存在＝散乱卓越の証拠（CMB
偏光など）。

\section*{さらに学ぶための参考文献}\label{sec-non-blackbody-further}
\addcontentsline{toc}{section}{さらに学ぶための参考文献}

\markright{さらに学ぶための参考文献}

\begin{itemize}
\tightlist
\item
  Mihalas, \emph{Stellar Atmospheres} (W.H. Freeman, 2nd ed., 1978) ---
  灰色大気・振動数依存不透明度の古典
\item
  Hubeny \& Mihalas, \emph{Theory of Stellar Atmospheres} (Princeton,
  2014) --- 現代的な放射輸送と恒星大気
\item
  Rybicki \& Lightman, \emph{Radiative Processes in Astrophysics}
  (Wiley, 1979) --- 修正黒体・散乱・非熱的放射の基礎
\item
  Draine, \emph{Physics of the Interstellar and Intergalactic Medium}
  (Princeton, 2011) --- ダストの放射と修正黒体
\end{itemize}

\chapter{第16章　宇宙マイクロ波背景放射}\label{ux7b2c16ux7ae0-ux5b87ux5b99ux30deux30a4ux30afux30edux6ce2ux80ccux666fux653eux5c04}

\textbf{この章の主題}：宇宙マイクロ波背景放射（CMB）は、\textbf{自然界で観測されるもっとも完璧な黒体スペクトル}である（COBE/FIRAS
で rms \(\sim 5 \times 10^{-5}\)、ピーク値
\(|f_\text{dev}| < 10^{-4}\)）。本章では「なぜこれほど完璧なのか」「いつ・どこで・どう熱化したのか」「宇宙膨張下でなぜ黒体形が保たれるのか」「観測されつつあるスペクトル歪みは何を見ているのか」を、第I〜V部で逆引きしたプランク分布の理解を直接適用して解読する。

\section*{この章の中心地図}\label{sec-cmb-map}
\addcontentsline{toc}{section}{この章の中心地図}

\markright{この章の中心地図}

\textbf{中心地図：プランク関数 ― 本部の主題は「ずれ」}

\[
B_\nu(T) = \frac{2 h \nu^3}{c^2} \cdot \frac{1}{e^{h\nu/kT} - 1} \quad \overset{\text{現実}}{\longrightarrow} \quad \text{?}
\]

第VI部（現実の宇宙へ戻る）は、観測される現実の天体スペクトルが、この理想的プランク関数からどのように、そしてなぜ
\textbf{ずれる}かを扱う。ずれの主な起源：

\begin{itemize}
\tightlist
\item
  \textbf{振動数依存不透明度} → 修正黒体（第15章）
\item
  \textbf{多温度成分の重ね合わせ} → 多温度黒体スペクトル（第17・18章）
\item
  \textbf{散乱効果}（コンプトン化など）→ 高エネルギー側の変形（第18章）
\item
  \textbf{非熱的放射の混入} → シンクロトロン・ベキ則の重なり
\end{itemize}

本部は、観測スペクトルを「理想的黒体＋ずれ」として読み解き、ずれの定量化から物理状態を引き出す視点を養う。

\begin{tcolorbox}[enhanced jigsaw, arc=.35mm, bottomrule=.15mm, bottomtitle=1mm, breakable, colback=white, colbacktitle=quarto-callout-note-color!10!white, colframe=quarto-callout-note-color-frame, coltitle=black, left=2mm, leftrule=.75mm, opacityback=0, opacitybacktitle=0.6, rightrule=.15mm, title=\textcolor{quarto-callout-note-color}{\faInfo}\hspace{0.5em}{ノート}, titlerule=0mm, toprule=.15mm, toptitle=1mm]

\textbf{方針}：CMB は「ずれ因子 \(f_\text{dev}\)
が事実上ゼロ」という極限ケース。中心地図のずれ因子の項を\textbf{裏返して理解する}
―
なぜ宇宙はこれほど完璧にプランク分布を実現できたか、を逆向きに問う章である。

\end{tcolorbox}

\section*{この章で答える問い}\label{sec-cmb-questions}
\addcontentsline{toc}{section}{この章で答える問い}

\markright{この章で答える問い}

\begin{itemize}
\tightlist
\item
  なぜ CMB はこれほど完璧な黒体スペクトルを示すのか
\item
  CMB の温度 2.725 K は、宇宙の何を測っているのか
\item
  宇宙膨張で温度が下がるのに、なぜ黒体スペクトルの形を保ち続けるのか
\item
  スペクトル歪み（\(\mu\), \(y\)
  歪み）は、宇宙史のどの時期の何を見ているのか
\end{itemize}

\section*{到達目標}\label{sec-cmb-goals}
\addcontentsline{toc}{section}{到達目標}

\markright{到達目標}

この章を読み終えると、読者は次のことができるようになる：

\begin{itemize}
\tightlist
\item
  CMB
  の黒体性が、宇宙が熱平衡を経たことの直接的証拠であることを説明できる
\item
  宇宙膨張下でも黒体スペクトルが保たれる理由を、断熱不変量の観点から論じられる
\item
  スペクトル歪みが宇宙論的観測のどの段階で意味を持つかを把握できる
\end{itemize}

\begin{center}\rule{0.5\linewidth}{0.5pt}\end{center}

\section{16.1 CMB スペクトルの観測 ―
完璧な黒体}\label{sec-cmb-observation}

{[}本文目安：B3{]}

1965 年に Penzias \& Wilson
が偶然発見した宇宙からの一様マイクロ波放射は、その後の精密観測により、\textbf{観測史上もっとも完璧な黒体スペクトル}であることが確かめられた。

決定的だったのは COBE 衛星の \textbf{FIRAS}（Far-InfraRed Absolute
Spectrophotometer）による観測（1989-1990）。FIRAS は CMB
を内部の極低温黒体源と比較するという絶妙な実験設計で、

\begin{itemize}
\tightlist
\item
  温度 \(T_\text{CMB} = 2.7255 \pm 0.0006\) K
\item
  黒体からの \textbf{ずれ} rms \(\sim 5 \times 10^{-5}\)、ピーク値
  \(|f_\text{dev}| < 10^{-4}\)（第1章 §1.5 と同じ値）
\end{itemize}

を確立した。これは第1章 §1.5 で観測動機として挙げた CMB
の「完璧な黒体スペクトル」の出典である。

\begin{figure}[H]

\centering{

\includegraphics[width=0.92\linewidth,height=\textheight,keepaspectratio]{index_files/mediabag/part6/../images/firas-spectrum.pdf}

}

\caption{\label{fig-firas-spectrum}\textbf{自然界でもっとも完璧な黒体スペクトル。}
COBE/FIRAS が測定した CMB モノポールスペクトル（赤点）と、\(T = 2.725\)
K
の黒体曲線（紺）。観測点は黒体曲線に完全に重なり、第I〜V部で逆引きしたプランク関数が現実の宇宙で厳密に成り立つことを示す。誤差棒は実際の
\(1\sigma\) 不確かさ（中央値 \(\sim 21\) kJy/sr、ピーク値の
\(\sim 10^{-4}\)）を \textbf{400 倍に拡大}
して初めて見える大きさである。下軸は周波数、上軸は波数。データは Fixsen
et al.~(1996, ApJ 473, 576)
のモノポールスペクトル（パブリックドメイン、NASA/COBE）。}

\end{figure}%

\begin{tcolorbox}[enhanced jigsaw, arc=.35mm, bottomrule=.15mm, breakable, colback=white, colframe=quarto-callout-warning-color-frame, left=2mm, leftrule=.75mm, opacityback=0, rightrule=.15mm, toprule=.15mm]
\begin{minipage}[t]{5.5mm}
\textcolor{quarto-callout-warning-color}{\faExclamationTriangle}
\end{minipage}%
\begin{minipage}[t]{\textwidth - 5.5mm}

\textbf{注意（観測）}：「黒体に見える」と「黒体である」の違い。観測精度
\(10^{-4}\) は十分高いが、\(10^{-5}\)
以下のスペクトル歪みは将来観測（PIXIE 等の提案）の標的。後述する \(\mu\)
歪み・\(y\)
歪みは、現在の観測限界をわずかに下回るレベルで存在することが理論的に予言されている。

なお、本書が扱うのは CMB の
\textbf{スペクトル歪み}（周波数方向の黒体からのずれ。FIRAS 残差や
\(\mu\)／\(y\) 歪み）であり、\textbf{空間温度ゆらぎ
\(\Delta T/T \sim 10^{-5}\)（異方性）}
とは別軸の量である。後者は宇宙の初期密度ゆらぎを担う観測量で、本書の射程外（§16.6
で位置づけのみ述べる）。

\end{minipage}%
\end{tcolorbox}

\section{16.2 なぜこれほど完璧な黒体なのか ―
熱化の三段階}\label{sec-cmb-thermalization}

{[}本文目安：B3-B4{]}

第8章で見たように、プランク分布が出るには
\textbf{光子数が自由に変動できる}（化学ポテンシャル
\(\mu = 0\)）必要があった。これを達成する初期宇宙の物理プロセスは以下の三つ。重要度の高い順に：

\begin{enumerate}
\def\labelenumi{\arabic{enumi}.}
\tightlist
\item
  \textbf{二重コンプトン散乱}（\textbf{double Compton scattering},
  \(e^- + \gamma \to e^- + 2\gamma\)）：光子数を増減できる。\(z \gtrsim 2 \times 10^6\)
  で効率的
\item
  \textbf{熱的制動放射}（\textbf{thermal
  bremsstrahlung}）：イオン-電子の衝突で光子を生む/吸う。\(z \gtrsim 10^5\)
  で効率的
\item
  \textbf{コンプトン散乱}（\textbf{Compton scattering},
  \(e^- + \gamma \to e^- + \gamma\)）：光子数は変えないが、エネルギーを再分配する。\(z \gtrsim 10^4\)
  で効率的
\end{enumerate}

\begin{tcolorbox}[enhanced jigsaw, arc=.35mm, bottomrule=.15mm, bottomtitle=1mm, breakable, colback=white, colbacktitle=quarto-callout-note-color!10!white, colframe=quarto-callout-note-color-frame, coltitle=black, left=2mm, leftrule=.75mm, opacityback=0, opacitybacktitle=0.6, rightrule=.15mm, title=\textcolor{quarto-callout-note-color}{\faInfo}\hspace{0.5em}{ノート}, titlerule=0mm, toprule=.15mm, toptitle=1mm]

\textbf{用語：再結合と熱化の区別}

「CMB の起源 = \(z \simeq 1100\) の
\textbf{再結合}」と覚えがちだが、これは不正確。再結合
(\(p + e^- \to H\)) は \textbf{光子と物質の結合が切れた時} で、CMB
が「最後に散乱された時」を決める ― いわば撮影時刻。

一方、\textbf{熱化} (\(\mu \to 0\))
は遥かに早期、\(z \gtrsim 2 \times 10^6\)
までに完了している。熱化が完了したのち、CMB は何も乱されずに
\textbf{凍結された完璧な黒体}
として宇宙膨張に乗って現在まで届く（次節）。

第1章 §1.5 で言及した「CMB
の完璧さは再結合より早期の物理の証拠」とは、この熱化過程のこと

\end{tcolorbox}

\section{16.3 熱史と再結合 ― 撮影時刻}\label{sec-cmb-recombination}

{[}本文目安：B3{]}

宇宙が膨張すると温度が下がる（§16.4）。\(T \simeq 3000\)
K（\(z \simeq 1100\)）で、自由電子と陽子が中性水素を作り始める：\textbf{再結合}（\textbf{recombination}）。

それまで宇宙はプラズマ状態で \textbf{Thomson 散乱}
により光子の平均自由行程が短く、光は不透明だった。再結合で自由電子が消えると、宇宙は急に透明になる。この瞬間に放たれた光が、現在私たちが
CMB として観測しているもの。

このため CMB は \textbf{「最終散乱面}」（last scattering
surface）からの光と呼ばれる。それより前の電磁波情報は手に入らない（重力波とニュートリノは別途）。

\begin{tcolorbox}[enhanced jigsaw, arc=.35mm, bottomrule=.15mm, breakable, colback=white, colframe=quarto-callout-tip-color-frame, left=2mm, leftrule=.75mm, opacityback=0, rightrule=.15mm, toprule=.15mm]
\begin{minipage}[t]{5.5mm}
\textcolor{quarto-callout-tip-color}{\faLightbulb}
\end{minipage}%
\begin{minipage}[t]{\textwidth - 5.5mm}

\textbf{問い}：もし宇宙の元素組成が違っていて（例：水素より重い元素ばかり）、再結合温度が違っていたら、現在観測される
CMB 温度 \(T_\text{CMB}^\text{obs}\) は変わるか？

\textbf{短答}：\textbf{変わらない}。今日の CMB 温度 \(T_0 \simeq 2.725\)
K
は、宇宙の光子のエントロピー密度（と膨張史）で決まるもので、再結合温度には依存しない。\(T \propto 1/(1+z)\)
は \textbf{恒等式}（§16.4）であり、「再結合温度 ÷
\((1+z_\text{rec})\)」と書いたときの両者は \textbf{恒等的に対応}
している。

\textbf{もう一歩}：仮に元素組成が違って再結合温度 \(T_\text{rec}\)
が違うとすれば、\textbf{\(z_\text{rec}\) がそれに応じて変わる}
だけで、\(T_0 = T_\text{rec}/(1+z_\text{rec})\)
は恒等的に成り立つ。\(T_0\)
を実際に決めているのは、宇宙の光子数密度の歴史（電子・陽電子対消滅前後の自由度の変化、バリオン-光子比
\(\eta\)
など）であり、再結合の詳細は「最終散乱面の位置」のみを左右する。

\end{minipage}%
\end{tcolorbox}

\section{16.4 宇宙膨張と温度 ―
なぜ形が崩れないか}\label{sec-cmb-expansion}

{[}本文目安：B3-B4{]}

膨張する宇宙では波長は伸びる（\textbf{redshift}）：\(\lambda \to a(t) \lambda\)、ここで
\(a(t)\) はスケール因子。よって振動数は \(\nu \to \nu/a\)、エネルギーは
\(h\nu \to h\nu/a\)。

プランク分布の指数因子は \(h\nu/k_B T\) を含むので、

\begin{equation}\protect\phantomsection\label{eq-redshift-temperature}{
\frac{h\nu}{k_B T}\quad \text{不変が成り立つには}\quad T \propto \frac{1}{a}
}\end{equation}

すなわち \textbf{温度がスケール因子の逆数で減る}
ならば、スペクトルの形は変わらない（温度だけが変わる）。これが宇宙膨張下で
\textbf{黒体形が保たれる} 理由。

\textbf{第14章の道具を使う}：モード密度 \(g(\nu) = 8\pi\nu^2/c^3\) は
\textbf{共動座標}
で考えると一定。波長が伸びれば、各モードの励起数も自然に変換される。光子数密度は
\(a^{-3}\) で薄まり、各光子のエネルギーは \(a^{-1}\)
で減るので、エネルギー密度は \(a^{-4}\) で減る ― これは
\(T \propto a^{-1}\) と \(u \propto T^4\)（Stefan-Boltzmann）の組合せ。

\begin{tcolorbox}[enhanced jigsaw, arc=.35mm, bottomrule=.15mm, bottomtitle=1mm, breakable, colback=white, colbacktitle=quarto-callout-note-color!10!white, colframe=quarto-callout-note-color-frame, coltitle=black, left=2mm, leftrule=.75mm, opacityback=0, opacitybacktitle=0.6, rightrule=.15mm, title=\textcolor{quarto-callout-note-color}{\faInfo}\hspace{0.5em}{ノート}, titlerule=0mm, toprule=.15mm, toptitle=1mm]

\textbf{観測との対応}：高赤方偏移銀河 (\(z \sim 3\)) で観測される CMB
の温度を、衝突励起遷移の温度計から測定すると、\(T_\text{CMB}(z) = (1+z) T_\text{CMB}(z=0) = 4(2.725) = 10.9\)
K
と一致する。これは「過去ほど宇宙は熱かった」ことの直接観測的確認である。

\end{tcolorbox}

\section{\texorpdfstring{16.5 スペクトル歪み ― \(\mu\) 歪みと \(y\)
歪み}{16.5 スペクトル歪み ― \textbackslash mu 歪みと y 歪み}}\label{sec-cmb-distortions}

{[}本文目安：B3-B4{]}

§16.2 の三つの熱化プロセスは、宇宙年齢が進むと一つずつ効率を失う：

\begin{itemize}
\tightlist
\item
  \(z \lesssim 2\times 10^6\)：二重コンプトン散乱が無効化 →
  \textbf{光子数を変える能力を失う}
\item
  \(z \lesssim 10^4\)：コンプトン散乱が無効化 →
  \textbf{エネルギーを再分配する能力も失う}
\end{itemize}

それぞれの時期に「もしエネルギーが宇宙に投入されると」、その「歪み」はその後消えずに残る：

\begin{itemize}
\tightlist
\item
  \textbf{\(\mu\)
  歪み}（\textbf{\(\mu\)-distortion}）：\(z \sim 5 \times 10^4 \to 2 \times 10^6\)
  の時期に投入されたエネルギーの記録。コンプトン散乱でエネルギー再分配は完了するが、光子数増減が不可能なため、\textbf{Bose-Einstein
  分布}（化学ポテンシャル \(\mu \neq 0\)）の形になる
\item
  \textbf{\(y\) 歪み}（\textbf{\(y\)-distortion},
  \textbf{Sunyaev-Zel'dovich type}）：\(z \lesssim 10^4\)
  の時期、または現在まで含む遅い時期に高温電子と CMB
  光子が散乱したときの記録。スペクトル全体が高振動数側へ引きずられる形（\textbf{コンプトン化}、第18章
  §18.3 で詳述）
\end{itemize}

標準的予言は \(|\mu|, |y| \sim 10^{-8}\)
レベル。これは現在の観測限界（FIRAS：\(|\mu| < 9 \times 10^{-5}\)、\(|y| < 1.5 \times 10^{-5}\)）より
\textbf{3〜4 桁} 下にあり、PIXIE、PRISM
等の将来計画の主標的となっている。

\begin{tcolorbox}[enhanced jigsaw, arc=.35mm, bottomrule=.15mm, bottomtitle=1mm, breakable, colback=white, colbacktitle=quarto-callout-note-color!10!white, colframe=quarto-callout-note-color-frame, coltitle=black, left=2mm, leftrule=.75mm, opacityback=0, opacitybacktitle=0.6, rightrule=.15mm, title=\textcolor{quarto-callout-note-color}{\faInfo}\hspace{0.5em}{ノート}, titlerule=0mm, toprule=.15mm, toptitle=1mm]

\textbf{観測との対応}：銀河団中の高温電子と CMB 光子の散乱による \(y\)
歪みは、\textbf{個別の銀河団}
で観測されている（\textbf{Sunyaev-Zel'dovich
効果}）。スペクトル全体が高振動数側へ引きずられるため、CMB
黒体スペクトルから差し引くと \textbf{\(\sim 217\) GHz
より低振動数側では銀河団方向が暗く（ディップ）、高振動数側では明るく（バンプ）}
見える ― 217 GHz
が「ヌル点」と呼ばれる境界。銀河団のガス圧の地図を作る道具として確立している。第18章
§18.3 のコンプトン化と同じ物理を、宇宙論的なスケールで見ている。

\end{tcolorbox}

\section{16.6 CMB の何を測っているか ―
4つの側面}\label{sec-cmb-what-we-measure}

{[}本文目安：B3{]}

CMB は \textbf{観測量}
として、現代宇宙論で最重要のひとつ。本書の射程である \textbf{スペクトル}
の側面に絞ると、CMB から読み取れることは：

\begin{enumerate}
\def\labelenumi{\arabic{enumi}.}
\tightlist
\item
  \textbf{温度 \(T_\text{CMB}\)} ―
  宇宙の現在の温度（\(z = 0\)）。\textbf{完璧に揃った温度}が、宇宙初期の熱化を直接示す
\item
  \textbf{スペクトル歪み} ―
  熱化が完了する前のエネルギー注入履歴。粒子物理（暗黒物質崩壊・原始ブラックホール蒸発など）の宇宙論的制約
\item
  \textbf{空間温度ゆらぎ} ―
  \(\Delta T/T \sim 10^{-5}\)。宇宙の初期密度ゆらぎ。本書の射程外（宇宙論の教科書を参照）
\item
  \textbf{偏光} ―
  再結合期の散乱の残り香。インフレーション理論検証の手段。本書の射程外
\end{enumerate}

\begin{tcolorbox}[enhanced jigsaw, arc=.35mm, bottomrule=.15mm, bottomtitle=1mm, breakable, colback=white, colbacktitle=quarto-callout-note-color!10!white, colframe=quarto-callout-note-color-frame, coltitle=black, left=2mm, leftrule=.75mm, opacityback=0, opacitybacktitle=0.6, rightrule=.15mm, title=\textcolor{quarto-callout-note-color}{\faInfo}\hspace{0.5em}{ノート}, titlerule=0mm, toprule=.15mm, toptitle=1mm]

\textbf{使えるようになった道具}（§16.1〜§16.6）：

\begin{itemize}
\tightlist
\item
  「黒体スペクトル」と「黒体に見えるスペクトル」を区別する観測精度の物差し
\item
  熱化と再結合の区別 ― 早期物理と後期物理を別軸として扱う
\item
  宇宙膨張下の黒体形保持 ― \(T \propto a^{-1}\) の意味
\item
  スペクトル歪み（\(\mu\),
  \(y\)）の物理的意味と、それが宇宙史のどの時期と対応するか
\end{itemize}

\end{tcolorbox}

\begin{center}\rule{0.5\linewidth}{0.5pt}\end{center}

\section*{この章で何がわかったか}\label{sec-cmb-summary}
\addcontentsline{toc}{section}{この章で何がわかったか}

\markright{この章で何がわかったか}

\textbf{中心地図に戻る}

CMB は、中心地図プランク分布の \textbf{ずれ因子 \(f_\text{dev}\)
が事実上ゼロ} という極限ケースだった。逆向きに読むと、これは
\textbf{宇宙が熱化を経た直接の証拠} である。

\begin{itemize}
\tightlist
\item
  熱化は \(z \gtrsim 2 \times 10^6\)
  で完了している（二重コンプトン散乱、第8章の光子数自由生成に対応）
\item
  再結合 \(z \simeq 1100\) で「最後の散乱」が起きて CMB は凍結
\item
  宇宙膨張下では \(T \propto a^{-1}\)
  で形が保たれる（モード密度の共動不変性）
\item
  スペクトル歪みは将来観測の主標的、熱化境界より早い宇宙物理を読む手段
\end{itemize}

次章では同じ道具を、もう少し身近な熱放射体 ― 恒星・惑星・ダスト ―
に適用する。

\section*{演習問題}\label{sec-cmb-exercises}
\addcontentsline{toc}{section}{演習問題}

\markright{演習問題}

\textbf{模範解答}
は折りたたみを展開して確認できる（オンライン版）。まず自力で解いてから開くこと。

\subsection*{問題 16-1　なぜ宇宙膨張で黒体形が崩れないか ― T は 1/a
に比例}\label{ex-16-1}
\addcontentsline{toc}{subsection}{問題
16-1　なぜ宇宙膨張で黒体形が崩れないか ― T は 1/a に比例}

{[}★ 難易度：☆☆ {]}

§16.4 は膨張下で黒体形が保たれる条件を述べた。確かめる。スケール因子
\(a(t)\)、膨張で波長 \(\lambda\to a\lambda\)、振動数 \(\nu\to\nu/a\)。

\begin{enumerate}
\def\labelenumi{\arabic{enumi}.}
\tightlist
\item
  プランク分布の指数 \(h\nu/k_BT\) が膨張下で不変であるためには
  \(T\propto1/a\) が必要であることを示せ。
\item
  \(T\propto a^{-1}\) と Stefan-Boltzmann \(u\propto T^4\)
  から、エネルギー密度が \(u\propto a^{-4}\) で減ることを示せ。
\item
  \(u\propto a^{-4}\) を「光子数密度
  \(\propto a^{-3}\)」と「1光子あたりのエネルギー
  \(\propto a^{-1}\)」の積として説明せよ。
\end{enumerate}

\textbf{関連}：\hyperref[sec-cmb-expansion]{§16.4
宇宙膨張と温度}／Stefan-Boltzmann は\hyperref[ex-10-3]{演習
10-3}、光子数密度 \(\propto T^3\) は\hyperref[ex-8-3]{演習 8-3}。

\section{模範解答（問題
16-1）}\label{ux6a21ux7bc4ux89e3ux7b54ux554fux984c-16-1}

\textbf{(1)} 膨張で各モードの振動数は \(\nu\to\nu/a\)。指数
\(h\nu/k_BT\) が不変であるには、\(\nu\) が \(1/a\) で減るのと同じ割合で
\(T\) も減らねばならない：\(T\propto1/a\)。このとき各振動数で占有因子
\(1/(e^{h\nu/k_BT}-1)\)
が保たれ、スペクトル形（プランク分布）は温度だけが下がる形で保存される。

\textbf{(2)} \(u\propto T^4\propto(a^{-1})^4=a^{-4}\)。

\textbf{(3)} 光子数密度
\(n_\gamma\propto T^3\propto a^{-3}\)（体積膨張で薄まる）。1光子あたりのエネルギー
\(h\nu\propto a^{-1}\)（赤方偏移）。積
\(u=n_\gamma\cdot\langle h\nu\rangle\propto a^{-3}\cdot a^{-1}=a^{-4}\)。物質（\(\rho\propto a^{-3}\)）より1段速く薄まるのが放射の特徴。

\textbf{答え}：指数不変には \(T\propto1/a\)。すると
\(u\propto T^4\propto a^{-4}=a^{-3}\)(数密度)×\(a^{-1}\)(光子エネルギー)。形は保たれ温度だけ下がる。

\subsection*{問題 16-2　高赤方偏移の CMB 温度}\label{ex-16-2}
\addcontentsline{toc}{subsection}{問題 16-2　高赤方偏移の CMB 温度}

{[}★ 難易度：☆☆ {]}

§16.4
は過去ほど宇宙が熱いことを述べた。\(T_\mathrm{CMB}(z)=(1+z)T_0\)、\(T_0=2.725\)
K。

\begin{enumerate}
\def\labelenumi{\arabic{enumi}.}
\tightlist
\item
  赤方偏移 \(z=3\) での CMB 温度を求めよ。
\item
  その温度の黒体の Wien ピーク波長
  \(\lambda_\mathrm{peak}=b/T\)（\(b=2.898\times10^{-3}\)
  m·K）を、\(z=3\) と現在（\(z=0\)）で計算して比べよ。
\item
  高赤方偏移銀河の吸収線温度計でこの \(T_\mathrm{CMB}(z)\)
  が確認されていることが、何の直接観測的証拠になるか述べよ。
\end{enumerate}

\textbf{関連}：\hyperref[sec-cmb-expansion]{§16.4 宇宙膨張と温度}／Wien
則は\hyperref[ex-2-2]{演習 2-2}。

\section{模範解答（問題
16-2）}\label{ux6a21ux7bc4ux89e3ux7b54ux554fux984c-16-2}

\textbf{(1)} \(T_\mathrm{CMB}(3)=(1+3)\times2.725=4\times2.725=10.9\)
K。

\textbf{(2)} 現在
\(z=0\)：\(\lambda_\mathrm{peak}=2.898\times10^{-3}/2.725=1.06\times10^{-3}\)
m \(=1.06\)
mm（マイクロ波）。\(z=3\)：\(\lambda_\mathrm{peak}=2.898\times10^{-3}/10.9=2.66\times10^{-4}\)
m \(=266\ \mu\)m（遠赤外）。過去はピークが短波長側。比はちょうど
\((1+z)=4\) 倍。

\textbf{(3)} 高赤方偏移銀河中の原子・分子の励起温度（CMB
光子で励起される準位比から測る温度計）が \((1+z)T_0\)
と一致することは、「過去ほど宇宙が熱かった」こと、すなわち
\(T\propto1/a\)（膨張で冷える黒体）という標準宇宙論の予言の直接観測的確認になる。

\textbf{答え}：\(T(z=3)=10.9\) K。\(\lambda_\mathrm{peak}\) は 1.06
mm→266 μm（\(1+z=4\) 倍短く）。\(T\propto1/a\) の直接確認。

\subsection*{問題 16-3　熱化と再結合の区別 ―
「撮影時刻」と「黒体化」}\label{ex-16-3}
\addcontentsline{toc}{subsection}{問題 16-3　熱化と再結合の区別 ―
「撮影時刻」と「黒体化」}

{[}★ 難易度：☆☆ {]}

§16.2・16.3 は「CMB
が再結合で黒体になった」という言い方が不正確だと述べた。整理する。

\begin{enumerate}
\def\labelenumi{\arabic{enumi}.}
\tightlist
\item
  「熱化（\(\mu\to0\)）」と「再結合（最終散乱）」は、それぞれ宇宙史のいつ起き、何を担う過程か区別して述べよ（赤方偏移の目安も）。
\item
  熱化を担う主過程（二重コンプトン散乱・制動放射・コンプトン散乱）のうち、光子数を変えられるのはどれか。なぜ
  \(\mu=0\) にこれが必要か。
\item
  もし元素組成が違って再結合温度が変わっても、現在の \(T_0=2.725\) K
  は変わらない。なぜか（\(T_0=T_\mathrm{rec}/(1+z_\mathrm{rec})\)
  が恒等式である理由）。
\end{enumerate}

\textbf{関連}：\hyperref[sec-cmb-thermalization]{§16.2
熱化の三段階}、\hyperref[sec-cmb-recombination]{§16.3 再結合}／\(\mu=0\)
の物理は\hyperref[ex-8-3]{演習
8-3}、第1章の動機は\hyperref[sec-where-blackbody-cmb]{§1.5}（\hyperref[ex-1-4]{演習
1-4}）。

\section{模範解答（問題
16-3）}\label{ux6a21ux7bc4ux89e3ux7b54ux554fux984c-16-3}

\textbf{(1)} 熱化：\(z\gtrsim2\times10^6\)
までに完了。放射のスペクトルを Planck
形（\(\mu\to0\)）に整える過程。再結合：\(z\simeq1100\)。電子と陽子が中性水素になり光が自由伝搬を始める（最終散乱＝CMB
の「撮影時刻」）。前者が黒体スペクトルを作り、後者がそれを解放する。

\textbf{(2)}
二重コンプトン散乱（\(e^-+\gamma\to e^-+2\gamma\)）と制動放射が光子数を変えられる。コンプトン散乱はエネルギー再分配のみで数を変えない。\(\mu=0\)（化学ポテンシャルゼロ＝光子数が自由に調整される）には光子数を増減できる過程が必要なので、二重コンプトン散乱が本質的。

\textbf{(3)} \(T_0=T_\mathrm{rec}/(1+z_\mathrm{rec})\) は
\(T\propto1/a\)（\(1+z=1/a\)）という恒等式。元素組成が変わって再結合温度
\(T_\mathrm{rec}\) が変われば、それに応じて \(z_\mathrm{rec}\)
も変わり、比 \(T_\mathrm{rec}/(1+z_\mathrm{rec})=T_0\)
は恒等的に保たれる。\(T_0\)
を実際に決めるのは光子数密度の歴史（バリオン-光子比など）で、再結合の詳細は最終散乱面の位置だけを左右する。

\textbf{答え}：熱化（\(z\gtrsim2\times10^6\)、黒体化）と再結合（\(z\simeq1100\)、撮影）は別過程。二重コンプトンが光子数を変え
\(\mu=0\) を作る。\(T_0\) は \(T\propto1/a\)
の恒等式ゆえ再結合温度に依らない。

\subsection*{問題 16-4　μ 歪みと y 歪み ―
いつの何を見るか}\label{ex-16-4}
\addcontentsline{toc}{subsection}{問題 16-4　μ 歪みと y 歪み ―
いつの何を見るか}

{[}★ 難易度：☆☆ {]}

§16.5 はスペクトル歪みを2種に分けた。整理する。

\begin{enumerate}
\def\labelenumi{\arabic{enumi}.}
\tightlist
\item
  \(\mu\) 歪みと \(y\)
  歪みが、それぞれ宇宙史のどの時期のエネルギー注入を記録するか述べよ（熱化過程の効率喪失の順序と結びつけて）。
\item
  \(\mu\) 歪みが「Bose-Einstein
  分布（\(\mu\neq0\)）」の形になる理由を、コンプトン散乱は働くが光子数を変えられない状況から説明せよ。
\item
  \(y\) 歪みが銀河団で観測される Sunyaev-Zel'dovich
  効果と同じ物理であること、約 217 GHz が「ヌル点」になることに触れよ。
\end{enumerate}

\textbf{関連}：\hyperref[sec-cmb-distortions]{§16.5
スペクトル歪み}／コンプトン化の同じ物理は\hyperref[sec-comptonization]{§18.3}（\hyperref[ex-18-2]{演習
18-2}）。

\section{模範解答（問題
16-4）}\label{ux6a21ux7bc4ux89e3ux7b54ux554fux984c-16-4}

\textbf{(1)} \(\mu\) 歪み：\(z\sim5\times10^4\)〜\(2\times10^6\)
のエネルギー注入を記録。この時期は二重コンプトンが無効化（光子数を変えられない）だがコンプトン散乱は有効（エネルギー再分配可能）。\(y\)
歪み：\(z\lesssim10^4\)（コンプトン散乱も非効率化して以降、現在まで含む）の高温電子との散乱を記録。熱化過程が効率を失う順序（光子数生成→エネルギー再分配）に対応する。

\textbf{(2)} \(\mu\)
歪みの時期は、コンプトン散乱がエネルギーを完全に再分配して平衡分布に達するが、光子数を変えられない。粒子数が保存される
Bose 気体の平衡は \(\mu\neq0\) の Bose-Einstein
分布（第8章）。だからプランク（\(\mu=0\)）からずれた \(\mu\neq0\)
の形になる。

\textbf{(3)} \(y\) 歪みは高温電子と CMB
光子の逆コンプトン散乱で、スペクトル全体が高振動数側へ引きずられる。これは銀河団中の高温ガスで起こる
Sunyaev-Zel'dovich
効果と同じ物理（スケールが宇宙論的か銀河団かの違いのみ）。低振動数側が暗く（ディップ）、高振動数側が明るく（バンプ）、その境界の約
217 GHz がヌル点。

\textbf{答え}：\(\mu\)
歪み＝\(z\sim5\times10^4\)〜\(2\times10^6\)（光子数凍結、BE
分布）、\(y\) 歪み＝\(z\lesssim10^4\)（高温電子散乱）。\(y\) は SZ
効果と同物理、217 GHz がヌル点。

\section*{さらに学ぶための参考文献}\label{sec-cmb-further}
\addcontentsline{toc}{section}{さらに学ぶための参考文献}

\markright{さらに学ぶための参考文献}

\begin{itemize}
\tightlist
\item
  Dodelson \& Schmidt, \emph{Modern Cosmology} (Academic Press, 2nd ed.,
  2020) --- CMB と宇宙論の標準教科書
\item
  Peebles, \emph{Principles of Physical Cosmology} (Princeton, 1993) ---
  古典、再結合と熱化の物理を丁寧に
\item
  Chluba \& Sunyaev, ``The evolution of CMB spectral distortions in the
  early Universe,'' MNRAS 419 (2012) 1294 --- スペクトル歪みの体系的計算
\item
  Mather et al., ``Measurement of the CMB Spectrum by the COBE FIRAS
  Instrument,'' ApJ 420 (1994) 439 --- FIRAS の原論文
\end{itemize}

\chapter{第17章　恒星・惑星・ダスト・電離ガスの熱放射}\label{ux7b2c17ux7ae0-ux6052ux661fux60d1ux661fux30c0ux30b9ux30c8ux96fbux96e2ux30acux30b9ux306eux71b1ux653eux5c04}

\textbf{この章の主題}：第15章でずれの体系を整理し、第16章で完全黒体の極限（CMB）を見た。本章では、もっとも身近な熱放射体
― 恒星・惑星・ダスト ― の
\textbf{観測スペクトルが、プランク分布を出発点としてどう書けるか}
を扱う。各天体のスペクトルから読み取れる物理量（温度・サイズ・距離・組成）と、それを引き出す典型的手順を、逆引きの立場で整理する。最後に、\textbf{面をもたない希薄な電離ガス}（HII
領域など）の自由--自由放射が、低振動数側で光学的に厚くなって黒体に漸近するという重要な観測例（§17.6）を加える。

\section*{この章の中心地図}\label{sec-stars-planets-dust-map}
\addcontentsline{toc}{section}{この章の中心地図}

\markright{この章の中心地図}

\textbf{中心地図：プランク関数 ― 本部の主題は「ずれ」}

\[
B_\nu(T) = \frac{2 h \nu^3}{c^2} \cdot \frac{1}{e^{h\nu/kT} - 1} \quad \overset{\text{現実}}{\longrightarrow} \quad \text{?}
\]

第VI部（現実の宇宙へ戻る）は、観測される現実の天体スペクトルが、この理想的プランク関数からどのように、そしてなぜ
\textbf{ずれる}かを扱う。ずれの主な起源：

\begin{itemize}
\tightlist
\item
  \textbf{振動数依存不透明度} → 修正黒体（第15章）
\item
  \textbf{多温度成分の重ね合わせ} → 多温度黒体スペクトル（第17・18章）
\item
  \textbf{散乱効果}（コンプトン化など）→ 高エネルギー側の変形（第18章）
\item
  \textbf{非熱的放射の混入} → シンクロトロン・ベキ則の重なり
\end{itemize}

本部は、観測スペクトルを「理想的黒体＋ずれ」として読み解き、ずれの定量化から物理状態を引き出す視点を養う。

\begin{tcolorbox}[enhanced jigsaw, arc=.35mm, bottomrule=.15mm, bottomtitle=1mm, breakable, colback=white, colbacktitle=quarto-callout-note-color!10!white, colframe=quarto-callout-note-color-frame, coltitle=black, left=2mm, leftrule=.75mm, opacityback=0, opacitybacktitle=0.6, rightrule=.15mm, title=\textcolor{quarto-callout-note-color}{\faInfo}\hspace{0.5em}{ノート}, titlerule=0mm, toprule=.15mm, toptitle=1mm]

\textbf{方針}：本章は最も「実務」に近い章。各節で「観測される SED →
仮定する物理モデル → 読み取れる物理量」の流れを意識して進める。

\end{tcolorbox}

本章で扱う各熱放射体の逆引きを、あらかじめ一表にまとめておく：

{\def\LTcaptype{none} % do not increment counter
\begin{longtable}[]{@{}
  >{\raggedright\arraybackslash}p{(\linewidth - 6\tabcolsep) * \real{0.2500}}
  >{\raggedright\arraybackslash}p{(\linewidth - 6\tabcolsep) * \real{0.2500}}
  >{\raggedright\arraybackslash}p{(\linewidth - 6\tabcolsep) * \real{0.2500}}
  >{\raggedright\arraybackslash}p{(\linewidth - 6\tabcolsep) * \real{0.2500}}@{}}
\toprule\noalign{}
\begin{minipage}[b]{\linewidth}\raggedright
観測で見える特徴
\end{minipage} & \begin{minipage}[b]{\linewidth}\raggedright
最小モデル
\end{minipage} & \begin{minipage}[b]{\linewidth}\raggedright
逆引きされる物理量
\end{minipage} & \begin{minipage}[b]{\linewidth}\raggedright
破れる仮定
\end{minipage} \\
\midrule\noalign{}
\endhead
\bottomrule\noalign{}
\endlastfoot
SED 全体の形・全光度 & 黒体 \(B_\nu(T_\text{eff})\) & 有効温度・半径 &
単一温度（大気の温度勾配） \\
惑星の赤外放射 & 入出力バランス & 平衡温度・アルベド &
大気なし（温室効果） \\
遠赤外〜サブミリのピークと勾配 & 修正黒体 \(Q_\nu B_\nu(T_d)\) &
ダスト温度・\(\beta\)・質量 & 理想黒体・単一温度 \\
電波連続光の折れ曲がり & 自由--自由放射（等温スラブ） & 電子温度
\(T_e\)・発光量度 EM & 光学的に薄い／一様 \\
\end{longtable}
}

\section*{この章で答える問い}\label{sec-stars-planets-dust-questions}
\addcontentsline{toc}{section}{この章で答える問い}

\markright{この章で答える問い}

\begin{itemize}
\tightlist
\item
  恒星のスペクトル分類（OBAFGKM）は何を反映しているのか
\item
  惑星の平衡温度を求めるとき、何を仮定し、何を無視しているのか
\item
  ダストの放射スペクトルが赤外で「黒体に似ているがずれている」のはなぜか
\item
  多温度成分（恒星表層と内部、HII
  領域と中性領域）は観測スペクトルにどう現れるのか
\item
  希薄な電離ガス（HII
  領域）の連続光が、低振動数側で「黒体的」に見えるのはなぜか
\end{itemize}

\section*{到達目標}\label{sec-stars-planets-dust-goals}
\addcontentsline{toc}{section}{到達目標}

\markright{到達目標}

この章を読み終えると、読者は次のことができるようになる：

\begin{itemize}
\tightlist
\item
  恒星の有効温度・惑星の平衡温度・ダストの色温度を、適切な近似のもとで計算できる
\item
  修正黒体スペクトルからダストの性質（放射効率・温度・量）を読み取れる
\item
  多温度合成スペクトルから、主要な温度成分を見分けられる
\item
  電離ガスの自由--自由放射が低振動数で黒体（\(B_\nu(T_e)\)）に漸近する仕組みを説明できる
\end{itemize}

\begin{center}\rule{0.5\linewidth}{0.5pt}\end{center}

\section{\texorpdfstring{17.1 恒星の有効温度 ―
\(L = 4\pi R^2 \sigma T_\text{eff}^4\)}{17.1 恒星の有効温度 ― L = 4\textbackslash pi R\^{}2 \textbackslash sigma T\_\textbackslash text\{eff\}\^{}4}}\label{sec-stars-effective-temperature}

{[}本文目安：B3{]}

恒星のスペクトルは、第15章 §15.2-§15.3 で扱った
\textbf{「振動数依存不透明度を持つ恒星大気」}
の典型例。完全な黒体ではないが、十分な近似で次のように書ける：

\begin{equation}\protect\phantomsection\label{eq-stefan-boltzmann-star}{
L_\star = 4\pi R_\star^2 \sigma T_\text{eff}^4
}\end{equation}

ここで \(L_\star\) は全光度、\(R_\star\) は恒星半径、\(T_\text{eff}\) は
\textbf{有効温度}。これは Stefan-Boltzmann 則（第2章
§2.4）の観測的応用：「もし恒星表面が温度 \(T_\text{eff}\)
の理想黒体だったとしたら、観測される全光度が同じになるような温度」と定義する。

\begin{tcolorbox}[enhanced jigsaw, arc=.35mm, bottomrule=.15mm, breakable, colback=white, colframe=quarto-callout-warning-color-frame, left=2mm, leftrule=.75mm, opacityback=0, rightrule=.15mm, toprule=.15mm]
\begin{minipage}[t]{5.5mm}
\textcolor{quarto-callout-warning-color}{\faExclamationTriangle}
\end{minipage}%
\begin{minipage}[t]{\textwidth - 5.5mm}

\textbf{注意（観測）}：「有効温度」は \textbf{観測上の便利な定義量}
であって、恒星の物理的温度ではない。表層は温度勾配を持っており、紫外側で「より深く」、赤外側で「より浅く」見えている（Eddington-Barbier、第15章
§15.3）。有効温度は
\textbf{積分量}（全光度）からの逆算値で、振動数ごとに違う「色温度」とは別概念

\end{minipage}%
\end{tcolorbox}

スペクトル分類 OBAFGKM は本質的に有効温度の分類：O型で約 \(40{,}000\)
K、M 型で約 \(3{,}000\) K。これに付随する違い ―
吸収線パターン、紫外/可視/赤外の連続光比、彩層活動 ―
はすべて温度から派生する物理である。

\begin{tcolorbox}[enhanced jigsaw, arc=.35mm, bottomrule=.15mm, breakable, colback=white, colframe=quarto-callout-tip-color-frame, left=2mm, leftrule=.75mm, opacityback=0, rightrule=.15mm, toprule=.15mm]
\begin{minipage}[t]{5.5mm}
\textcolor{quarto-callout-tip-color}{\faLightbulb}
\end{minipage}%
\begin{minipage}[t]{\textwidth - 5.5mm}

\textbf{問い}：太陽（\(T_\text{eff} = 5778\) K）と
Vega（\(T_\text{eff} = 9602\) K）の SED を比べたとき、両者の
\textbf{同一波長} での明るさ比（半径と距離を同じと仮定）はどれくらいか？
紫外、可視、赤外それぞれで違うか？

\textbf{短答}：Wien 則 \(\lambda_\text{peak} \propto 1/T\) から、Vega
のピークは太陽より短波長側。同一波長での比 \(B_\nu(T_1)/B_\nu(T_2)\) は

\begin{itemize}
\tightlist
\item
  可視（500 nm）：Vega が 約 7.6 倍明るい
\item
  紫外（250 nm）：Vega が 約 53 倍明るい（Wien 側の指数差が支配）
\item
  赤外（5 µm）：Vega が 約 1.85 倍明るい（Rayleigh-Jeans 側〔以下 RJ
  側〕、温度比のみ効く）
\end{itemize}

紫外側ほど温度差が劇的に効くのが「Wien 側の急峻性」の現れ

\end{minipage}%
\end{tcolorbox}

\section{17.2 惑星の平衡温度 ―
入出力バランス}\label{sec-planet-equilibrium}

{[}本文目安：B3{]}

惑星は内部の核融合エネルギーを持たない（地球型岩石惑星では地殻活動の貢献は弱い）。観測される熱放射は、ほぼすべて
\textbf{恒星からの入射光を吸収・再放射} したもの。

エネルギー収支：

\begin{equation}\protect\phantomsection\label{eq-planet-balance}{
\underbrace{(1-A) \cdot \frac{L_\star}{4\pi d^2} \cdot \pi R_p^2}_\text{吸収光} = \underbrace{4\pi R_p^2 \cdot \sigma T_\text{eq}^4}_\text{再放射}
}\end{equation}

ここで \(A\) はアルベド（反射率）、\(d\) は恒星-惑星距離、\(R_p\)
は惑星半径。解いて

\begin{equation}\protect\phantomsection\label{eq-planet-temperature}{
T_\text{eq} = T_\star \left(\frac{R_\star}{2d}\right)^{1/2} (1-A)^{1/4}
}\end{equation}

地球の場合（\(T_\star = 5778\) K, \(R_\star/d = 1/214\),
\(A = 0.3\)）：\(T_\text{eq} \simeq 255\) K \(\simeq -18\) °C。

\begin{tcolorbox}[enhanced jigsaw, arc=.35mm, bottomrule=.15mm, breakable, colback=white, colframe=quarto-callout-warning-color-frame, left=2mm, leftrule=.75mm, opacityback=0, rightrule=.15mm, toprule=.15mm]
\begin{minipage}[t]{5.5mm}
\textcolor{quarto-callout-warning-color}{\faExclamationTriangle}
\end{minipage}%
\begin{minipage}[t]{\textwidth - 5.5mm}

\textbf{注意（観測）}：実際の地球表面温度は約 +15 °C で、平衡温度より 33
K 高い。これは
\textbf{温室効果}（大気の赤外不透明度）による。式~\ref{eq-planet-temperature}
は \textbf{「大気がない」近似}
で、地球のように大気のある惑星では大きく外れる。

しかし「平衡温度」は次のような実用的価値を持つ：

\begin{itemize}
\tightlist
\item
  多くの太陽系外惑星では大気がよくわからない → 平衡温度が
  \textbf{最初の温度推定値} として使われる
\item
  「ハビタブルゾーン」（液体水が安定な軌道）は平衡温度から定義される
\end{itemize}

\end{minipage}%
\end{tcolorbox}

\section{17.3 ダストの修正黒体 ―
放射効率と粒子サイズ}\label{sec-dust-modified-blackbody}

{[}本文目安：B3{]}

星間ダスト、原始惑星系円盤の塵、星形成領域の冷たい塵 ―
これらの熱放射は、第15章 §15.4 の \textbf{修正黒体スペクトル}

\begin{equation}\protect\phantomsection\label{eq-modified-blackbody-applied}{
I_\nu^\text{dust} = Q_\nu \cdot B_\nu(T_d), \quad Q_\nu \propto \nu^\beta \text{ (long wavelength)}
}\end{equation}

で記述される。観測される \(I_\nu\) から、三つの物理量が読み取れる：

\begin{enumerate}
\def\labelenumi{\arabic{enumi}.}
\tightlist
\item
  \textbf{ダスト温度 \(T_d\)} ― Wien 側の指数的減衰の鋭さ、または SED
  のピーク位置から
\item
  \textbf{ダスト放射率指数 \(\beta\)}（emissivity index）―
  Rayleigh-Jeans 側の勾配の急さから
\item
  \textbf{ダスト質量} ― 全光度を \(B_\nu(T_d)\)
  で割ったときの規格化定数から（不透明度を仮定する必要あり）
\end{enumerate}

\(\beta\) の値はダストの物理を反映する：

\begin{itemize}
\tightlist
\item
  \textbf{大きな粒子（\(\gtrsim 100\) µm）}：\(\beta \to 0\)（粒子サイズ
  ≈ 波長の極限）
\item
  \textbf{小さな結晶粒子（シリケート、\(\lesssim 1\)
  µm）}：\(\beta \to 2\)
\item
  \textbf{非晶質シリケート、グラファイト}：\(\beta \simeq 1.5\)〜\(2\)
\end{itemize}

\begin{tcolorbox}[enhanced jigsaw, arc=.35mm, bottomrule=.15mm, bottomtitle=1mm, breakable, colback=white, colbacktitle=quarto-callout-note-color!10!white, colframe=quarto-callout-note-color-frame, coltitle=black, left=2mm, leftrule=.75mm, opacityback=0, opacitybacktitle=0.6, rightrule=.15mm, title=\textcolor{quarto-callout-note-color}{\faInfo}\hspace{0.5em}{ノート}, titlerule=0mm, toprule=.15mm, toptitle=1mm]

\textbf{観測との対応}：ALMA で観測される原始惑星系円盤のミリ波 SED を
\(\beta\) で分析すると、\textbf{円盤の中心部ほど粒子が大きい}（\(\beta\)
が小さい）ことが発見されている。これは円盤内でダストが成長して惑星形成へ向かう過程の直接観測的証拠

\end{tcolorbox}

\section{17.4 赤外線天文学 ―
多温度天体の現代観測}\label{sec-infrared-astronomy}

{[}本文目安：B3{]}

近代の宇宙論的観測では、赤外線・サブミリ波の天文学が決定的に重要。なぜなら：

\begin{itemize}
\tightlist
\item
  銀河の重元素の半分が星間ダストに閉じ込められており、その熱放射は
  \textbf{遠赤外（FIR）に大きなピーク} を持つ
\item
  高赤方偏移銀河は、紫外/可視光が宇宙膨張で赤方偏移し、観測波長帯が赤外/サブミリへ移動する
\end{itemize}

主要観測施設の波長帯：

{\def\LTcaptype{none} % do not increment counter
\begin{longtable}[]{@{}lll@{}}
\toprule\noalign{}
観測施設 & 波長帯 & 主な対象 \\
\midrule\noalign{}
\endhead
\bottomrule\noalign{}
\endlastfoot
Spitzer & 3.6〜160 µm & 中間赤外〜遠赤外、銀河系・近傍銀河 \\
Herschel & 60〜600 µm & 遠赤外〜サブミリ、星形成 \\
ALMA & 0.3〜3.6 mm & (サブ)ミリ波、原始惑星系円盤・高赤方偏移銀河 \\
JWST & 0.6〜28 µm & 近赤外〜中間赤外、高分解能分光・恒星形成 \\
\end{longtable}
}

これら異なる波長帯の同一天体観測を組み合わせた \textbf{多波長 SED}
を、複数温度のダスト成分でフィットすることで、温度・質量・組成・空間構造を逆引きする。

\section{17.5 多温度合成 ― SED
フィッティング}\label{sec-multi-temperature}

{[}本文目安：B3{]}

実際の銀河や星間雲は \textbf{複数の温度成分の混合} であることが普通：

\begin{itemize}
\tightlist
\item
  銀河の SED：恒星の光（\(\sim 10^4\) K の準黒体）＋
  暖かいダスト（\(\sim 100\) K）＋ 冷たいダスト（\(\sim 20\) K）＋ AGN
  や活発な星形成があれば中間赤外の高温成分
\item
  星形成領域：HII 領域（\(\sim 10^4\) K のプラズマ自由-自由放射）＋
  高温ダスト（\(\sim 100\) K）＋ 冷たい分子雲のダスト（\(\sim 15\) K）＋
  線スペクトル
\end{itemize}

全 SED は

\begin{equation}\protect\phantomsection\label{eq-multi-temperature-sed}{
I_\nu^\text{total} = \sum_i I_{\nu, \text{star}, i} + \sum_j Q_\nu(\beta_j)\, B_\nu(T_j)\, M_j + I_\nu^\text{lines} + I_\nu^\text{ff} + \ldots
}\end{equation}

の形で組み立てる。フィッティングは温度・質量・\(\beta\)
の組合せを観測波長点と適合させ、\(\chi^2\) 最小化や Bayesian inference
で求める。

\begin{tcolorbox}[enhanced jigsaw, arc=.35mm, bottomrule=.15mm, breakable, colback=white, colframe=quarto-callout-warning-color-frame, left=2mm, leftrule=.75mm, opacityback=0, rightrule=.15mm, toprule=.15mm]
\begin{minipage}[t]{5.5mm}
\textcolor{quarto-callout-warning-color}{\faExclamationTriangle}
\end{minipage}%
\begin{minipage}[t]{\textwidth - 5.5mm}

\textbf{注意（観測）}：「冷たい成分の質量」を多温度フィットから推定するとき、温度との
\textbf{強い縮退} に注意。同じピーク位置でも、(高温・少量) と
(低温・大量) で同程度の SED を作れる場合がある。長波長側（Rayleigh-Jeans
側）の観測点が、この縮退を解く鍵となる

\end{minipage}%
\end{tcolorbox}

\begin{tcolorbox}[enhanced jigsaw, arc=.35mm, bottomrule=.15mm, breakable, colback=white, colframe=quarto-callout-tip-color-frame, left=2mm, leftrule=.75mm, opacityback=0, rightrule=.15mm, toprule=.15mm]
\begin{minipage}[t]{5.5mm}
\textcolor{quarto-callout-tip-color}{\faLightbulb}
\end{minipage}%
\begin{minipage}[t]{\textwidth - 5.5mm}

\textbf{問い}：温度 \(T_d = 20\) K のダスト成分のピーク波長を Wien
則で求めよ。これは何の観測施設の波長帯か？

\textbf{短答}：\(\lambda_\text{peak} \approx 2898\,\text{µm·K} / 20\,\text{K} \simeq 145\)
µm。これは Herschel SPIRE / PACS の主帯域。冷たい星間ダストの観測が
Herschel で飛躍的に進んだのはこのため

\end{minipage}%
\end{tcolorbox}

\section{17.6 電離ガスの自由--自由放射 ―
電波で見る黒体への漸近}\label{sec-free-free-ionized}

{[}本文目安：B3{]}

恒星・惑星・ダストはいずれも「面をもつ高密度の熱放射体」だった。しかし
\textbf{希薄な電離ガス}（HII
領域・惑星状星雲・銀河中心の電離ガスなど）も、熱的な連続光 ―
\textbf{自由--自由放射}（\textbf{free-free
emission}、\textbf{熱的制動放射}, thermal bremsstrahlung）―
を放つ。自由電子がイオンの Coulomb
場で曲げられるときに光子を出す過程で、電子が Maxwell 分布（温度
\(T_e\)）にあれば \textbf{熱的放射}（第1章 §1.9）であり、LTE
のもとでソース関数は \(S_\nu=B_\nu(T_e)\) になる（第3章・第5章）。

鍵は、\textbf{自由--自由吸収係数が低振動数ほど大きい}ことである：

\begin{equation}\protect\phantomsection\label{eq-ff-absorption}{
\alpha_\nu^{\mathrm{ff}} \propto n_e\, n_i\, T_e^{-3/2}\, \nu^{-2}\, \bar g_{\mathrm{ff}}(\nu,T_e)
}\end{equation}

ここで \(n_e, n_i\) は電子・イオン数密度、\(\bar g_{\mathrm{ff}}\) は
Gaunt 因子（\(\nu,T_e\)
にゆるやかに依存）。自由--自由の放出係数・吸収係数の微視的な導出と、両者を結ぶ
Kirchhoff の関係 \(S_\nu=B_\nu(T_e)\)、Gaunt
因子の起源は\hyperref[sec-continuum-free-free]{第19章 §19.2}で扱う。

これを視線方向に積分したものが光学的厚さ
\(\tau_\nu=\int\alpha_\nu^{\mathrm{ff}}\,ds\) である。密度の寄与は
\textbf{発光量度}（\textbf{emission measure}）

\begin{equation}\protect\phantomsection\label{eq-emission-measure}{
\mathrm{EM}\equiv\int n_e^2\,ds \qquad(\text{水素プラズマでは } n_i\simeq n_e)
}\end{equation}

にまとまる。残る振動数・温度依存は、式~\ref{eq-ff-absorption} の
\(\nu^{-2}\,T_e^{-3/2}\) に Gaunt
因子のゆるやかな依存（\(\bar g_{\mathrm{ff}}\propto\nu^{-0.1}\,T_e^{0.15}\)
程度）が重なったもので、実用的な近似形は

\begin{equation}\protect\phantomsection\label{eq-ff-optical-depth}{
\tau_\nu \;\propto\; \mathrm{EM}\,\nu^{-2.1}\,T_e^{-1.35}
}\end{equation}

となる。\(\mathrm{EM}\) は吸収の強さ（密度の寄与）を、指数
\(\nu^{-2.1}\)
はその振動数依存を与える。後者が支配的で、\(\tau_\nu\propto\nu^{-2.1}\)
のため \textbf{低振動数で光学的に厚く、高振動数で薄い}。

等温スラブ（温度 \(T_e\)、LTE）の射出強度は第5章の解
\(I_\nu=B_\nu(T_e)(1-e^{-\tau_\nu})\)。電波帯では \(h\nu\ll k_BT_e\)
なので Rayleigh--Jeans が使え、輝度温度
\(T_b\equiv c^2 I_\nu/(2\nu^2 k_B)\) で書くと

\begin{equation}\protect\phantomsection\label{eq-ff-brightness-temp}{
T_b(\nu)=T_e\left(1-e^{-\tau_\nu}\right)
}\end{equation}

\textbf{低振動数側（光学的に厚い、\(\tau_\nu\gg1\)）}：\(T_b\to T_e\)
に飽和し、\(I_\nu\to B_\nu(T_e)\)。これは第5章 §5.6
の「光学的に厚い極限で \(I_\nu\to B_\nu(T)\)」そのもので、観測されるのは
\textbf{電子温度 \(T_e\) の黒体（の Rayleigh--Jeans 裾）}
である。フラックスは \(F_\nu\propto\nu^2\)。

\textbf{高振動数側（光学的に薄い、\(\tau_\nu\ll1\)）}：\(T_b\to T_e\,\tau_\nu\propto\nu^{-2.1}\)
なので \(F_\nu\propto\nu^2\cdot\nu^{-2.1}=\nu^{-0.1}\) ― ほぼ平坦。

両者の境目（\(\tau_\nu\simeq1\) となる \textbf{ターンオーバー振動数}
\(\nu_t\)）は \(\mathrm{EM}\) と \(T_e\) で決まる。したがって電波 HII
領域の連続スペクトルは、\textbf{低 \(\nu\) 側
\(F_\nu\propto\nu^2\)（黒体的）→ 高 \(\nu\) 側
\(F_\nu\propto\nu^{-0.1}\)（薄い）} という特徴的な折れ曲がりを示す。

\begin{tcolorbox}[enhanced jigsaw, arc=.35mm, bottomrule=.15mm, bottomtitle=1mm, breakable, colback=white, colbacktitle=quarto-callout-note-color!10!white, colframe=quarto-callout-note-color-frame, coltitle=black, left=2mm, leftrule=.75mm, opacityback=0, opacitybacktitle=0.6, rightrule=.15mm, title=\textcolor{quarto-callout-note-color}{\faInfo}\hspace{0.5em}{ノート}, titlerule=0mm, toprule=.15mm, toptitle=1mm]

\textbf{観測との対応}：電波で観測する HII
領域（オリオン大星雲など）・超コンパクト HII
領域・惑星状星雲の連続スペクトルは、まさにこの折れ曲がりを示す。低振動数側で輝度温度
\(T_b\)
が一定値に飽和すること自体が「\textbf{ガスが光学的に厚くなり電子温度
\(T_e\) の黒体に達した}」直接の証拠であり、そこから \(T_e\)
を、ターンオーバーから \(\mathrm{EM}\)（密度 ×
経路長）を読み取れる。第1章 §1.2--1.6
の黒体例（恒星光球・CMB・降着円盤）は \textbf{面をもつ天体}
だったが、ここでは \textbf{希薄な気体} が低振動数で黒体に漸近する ―
第5章 §5.6「光学的に厚い → 黒体」の最も直接的な気体の実例である。

\end{tcolorbox}

\begin{tcolorbox}[enhanced jigsaw, arc=.35mm, bottomrule=.15mm, breakable, colback=white, colframe=quarto-callout-tip-color-frame, left=2mm, leftrule=.75mm, opacityback=0, rightrule=.15mm, toprule=.15mm]
\begin{minipage}[t]{5.5mm}
\textcolor{quarto-callout-tip-color}{\faLightbulb}
\end{minipage}%
\begin{minipage}[t]{\textwidth - 5.5mm}

\textbf{問い}：自由--自由放射は「熱的放射」なのに、なぜ全帯域で単一の黒体（プランク分布）には
\textbf{ならない} のか？

\textbf{短答}：黒体になるのは \textbf{光学的に厚い} 振動数帯（低 \(\nu\)
側）だけ。高 \(\nu\)
側では光学的に薄く、\(I_\nu\approx\tau_\nu B_\nu(T_e)\ll B_\nu(T_e)\)
となって黒体から下にずれる。「熱的（粒子が Maxwell
分布）」と「黒体（光学的に厚く熱平衡放射場）」は別条件であり、自由--自由放射は前者を常に満たすが、後者は低振動数でのみ満たす。

\textbf{もう一歩}：同じ熱的自由--自由放射でも、銀河団の高温プラズマ（\(T_e\sim10^7\)〜\(10^8\)
K）からの \textbf{X 線制動放射}
は全帯域で光学的に薄く、\(F_\nu\propto e^{-h\nu/k_BT_e}\)
の指数カットオフをもつ「薄い側」の典型例である（第18章）。低振動数で黒体に漸近するのは、\(\nu^{-2}\)
で吸収が伸びる電波帯ならではの現象である。

\end{minipage}%
\end{tcolorbox}

\section*{17.7
使えるようになった道具}\label{sec-stars-planets-dust-tools}
\addcontentsline{toc}{section}{17.7 使えるようになった道具}

\markright{17.7 使えるようになった道具}

\begin{tcolorbox}[enhanced jigsaw, arc=.35mm, bottomrule=.15mm, bottomtitle=1mm, breakable, colback=white, colbacktitle=quarto-callout-note-color!10!white, colframe=quarto-callout-note-color-frame, coltitle=black, left=2mm, leftrule=.75mm, opacityback=0, opacitybacktitle=0.6, rightrule=.15mm, title=\textcolor{quarto-callout-note-color}{\faInfo}\hspace{0.5em}{ノート}, titlerule=0mm, toprule=.15mm, toptitle=1mm]

\begin{itemize}
\tightlist
\item
  恒星の有効温度の定義と限界（積分量としての \(T_\text{eff}\)）
\item
  惑星の平衡温度と温室効果のずれ
\item
  修正黒体スペクトル \(Q_\nu B_\nu(T)\) からのダスト 3 物理量読み取り
\item
  多温度 SED フィッティングの構造と縮退性
\item
  観測波長帯（赤外・サブミリ・ミリ）と各観測施設の対応関係
\item
  電離ガスの自由--自由放射：低 \(\nu\) 側は \(T_e\)
  の黒体（Rayleigh--Jeans）、高 \(\nu\) 側は
  \(F_\nu\propto\nu^{-0.1}\)、ターンオーバーから \(\mathrm{EM}\) と
  \(T_e\)
\end{itemize}

\end{tcolorbox}

\begin{center}\rule{0.5\linewidth}{0.5pt}\end{center}

\section*{この章で何がわかったか}\label{sec-stars-planets-dust-summary}
\addcontentsline{toc}{section}{この章で何がわかったか}

\markright{この章で何がわかったか}

\textbf{中心地図に戻る}

中心地図プランク分布 \(B_\nu(T)\)
に、ずれ因子として「振動数依存不透明度」「修正黒体係数
\(Q_\nu\)」「複数温度成分の和」を加えるだけで、宇宙にある熱放射体の大半を統一的に記述できることを見た。

\begin{itemize}
\tightlist
\item
  恒星：\(T_\text{eff}\) で代表される（厳密には大気の温度構造の積分量）
\item
  惑星：恒星光の再放射、\(T_\text{eq}\) ＋ 温室効果のずれ
\item
  ダスト：修正黒体 \(Q_\nu B_\nu(T_d)\)、\(\beta\) が粒子物理を反映
\item
  銀河・星形成領域：多温度合成、波長帯ごとに異なる成分が主役
\item
  電離ガス（HII 領域）：自由--自由放射が低振動数で光学的に厚くなり
  \(B_\nu(T_e)\) に漸近 ―
  面をもたない希薄な気体が黒体に達する観測例（§17.6）
\end{itemize}

次章では、もっと極端な物理条件 ― 強い重力場・強磁場・相対論的速度 ―
のもとでも、同じ枠組み（プランク関数 ＋
ずれ因子）がどこまで通用するかを見る

\section*{演習問題}\label{sec-stars-planets-dust-exercises}
\addcontentsline{toc}{section}{演習問題}

\markright{演習問題}

\textbf{模範解答}
は折りたたみを展開して確認できる（オンライン版）。まず自力で解いてから開くこと。

\subsection*{問題
17-1　惑星の平衡温度を入出力バランスから}\label{ex-17-1}
\addcontentsline{toc}{subsection}{問題
17-1　惑星の平衡温度を入出力バランスから}

{[}★ 難易度：☆☆ {]}

§17.2 のエネルギー収支から惑星平衡温度を導く。アルベド
\(A\)、恒星-惑星距離 \(d\)、恒星半径 \(R_\star\)、温度
\(T_\star\)、惑星半径 \(R_p\)。

\begin{enumerate}
\def\labelenumi{\arabic{enumi}.}
\tightlist
\item
  吸収光 \((1-A)\dfrac{L_\star}{4\pi d^2}\pi R_p^2\) と再放射
  \(4\pi R_p^2\sigma T_\mathrm{eq}^4\)
  を等置し、\(L_\star=4\pi R_\star^2\sigma T_\star^4\) を使って
  \(T_\mathrm{eq}=T_\star\left(\dfrac{R_\star}{2d}\right)^{1/2}(1-A)^{1/4}\)
  を導け。
\item
  断面積 \(\pi R_p^2\)（吸収）と表面積 \(4\pi R_p^2\)（放射）の比 4
  が、式のどこに効くか説明せよ。
\item
  地球（\(T_\star=5778\) K、\(R_\star/d=1/214\)、\(A=0.3\)）で
  \(T_\mathrm{eq}\) を計算し、実際の地表温度（+15
  °C）とのずれの原因を述べよ。
\end{enumerate}

\textbf{関連}：\hyperref[sec-planet-equilibrium]{§17.2
惑星の平衡温度}／Stefan-Boltzmann
と光度は\hyperref[sec-reading-spectra-stefan-boltzmann]{§2.4}（\hyperref[ex-2-5]{演習
2-5}）。

\section{模範解答（問題
17-1）}\label{ux6a21ux7bc4ux89e3ux7b54ux554fux984c-17-1}

\textbf{(1)}
等置：\((1-A)\dfrac{L_\star}{4\pi d^2}\pi R_p^2=4\pi R_p^2\sigma T_\mathrm{eq}^4\)。\(\pi R_p^2\)
を消し、\(L_\star=4\pi R_\star^2\sigma T_\star^4\) を代入：

\[(1-A)\frac{4\pi R_\star^2\sigma T_\star^4}{4\pi d^2}\cdot\frac14=\sigma T_\mathrm{eq}^4\ \Rightarrow\ T_\mathrm{eq}^4=(1-A)\frac{R_\star^2}{4d^2}T_\star^4.\]

\[T_\mathrm{eq}=T_\star\left(\frac{R_\star}{2d}\right)^{1/2}(1-A)^{1/4}.\]

\textbf{(2)} 惑星は断面 \(\pi R_p^2\) で恒星光を受け、全表面
\(4\pi R_p^2\) から再放射する。比 \(4\pi R_p^2/\pi R_p^2=4\) が右辺の
\(1/4\)（すなわち \(\sqrt{1/4}=1/2\) が \(R_\star/2d\) の
2）として効く。自転が速く全球で均一再放射する近似。

\textbf{(3)}
\(T_\mathrm{eq}=5778\times(1/(2\times214))^{1/2}\times(0.7)^{1/4}=5778\times(2.34\times10^{-3})^{1/2}\times0.915=5778\times0.0483\times0.915\simeq255\)
K（\(-18\)°C）。実際の +15°C との 33 K
の差は温室効果（大気の赤外不透明度）。この式は大気なし近似。

\textbf{答え}：\(T_\mathrm{eq}=T_\star\sqrt{R_\star/2d}\,(1-A)^{1/4}\)。地球で
255 K。実測との 33 K 差は温室効果。

\subsection*{問題 17-2　二温度天体の同一波長明るさ比}\label{ex-17-2}
\addcontentsline{toc}{subsection}{問題
17-2　二温度天体の同一波長明るさ比}

{[}★ 難易度：☆☆☆ {]}

§17.1 の問いを数値で解く。太陽（\(T_1=5778\) K）と Vega（\(T_2=9602\)
K）を半径・距離同じと仮定。同一振動数での比強度比
\(\dfrac{B_\nu(T_2)}{B_\nu(T_1)}=\dfrac{e^{h\nu/kT_1}-1}{e^{h\nu/kT_2}-1}\)
を使う。便利な関係 \(h\nu/k=hc/\lambda k\)、\(hc/k=1.439\times10^{-2}\)
m·K。

\begin{enumerate}
\def\labelenumi{\arabic{enumi}.}
\tightlist
\item
  可視 \(\lambda=500\) nm で \(x_i=hc/\lambda kT_i\)
  を求め、明るさ比を計算せよ（約7.6倍）。
\item
  紫外 \(\lambda=250\) nm
  で同様に計算せよ（約53倍）。なぜ紫外ほど比が大きいか。
\item
  赤外 \(\lambda=5\ \mu\)m での比を計算し（約1.85倍）、Rayleigh-Jeans
  側では温度比のみが効くことを確認せよ。
\end{enumerate}

\textbf{関連}：\hyperref[sec-stars-effective-temperature]{§17.1
恒星の有効温度}／両極限は\hyperref[sec-reading-spectra-limits]{§2.5}、恒星の色は\hyperref[sec-where-blackbody-stars]{§1.3}。

\section{模範解答（問題
17-2）}\label{ux6a21ux7bc4ux89e3ux7b54ux554fux984c-17-2}

\(x=hc/\lambda kT\)、\(hc/k=1.439\times10^{-2}\) m·K。比
\(=\dfrac{e^{x_1}-1}{e^{x_2}-1}\)（\(x_1\)＝太陽、\(x_2\)＝Vega、\(x_1>x_2\)）。

\textbf{(1)} \(\lambda=500\)
nm：\(hc/\lambda k=1.439\times10^{-2}/5.0\times10^{-7}=2.88\times10^4\)
K。\(x_1=28800/5778=4.98\)、\(x_2=28800/9602=3.00\)。比
\(=\dfrac{e^{4.98}-1}{e^{3.00}-1}=\dfrac{144.4}{19.1}\simeq7.6\)。Vega
が約7.6倍明るい。

\textbf{(2)} \(\lambda=250\) nm：\(hc/\lambda k=5.76\times10^4\)
K。\(x_1=9.96\)、\(x_2=5.99\)。比
\(=\dfrac{e^{9.96}-1}{e^{5.99}-1}=\dfrac{21100}{399}\simeq53\)。紫外は両者とも
Wien 側で指数 \(e^{x}\)
が支配し、温度差が指数的に増幅されるため比が大きい。

\textbf{(3)} \(\lambda=5\ \mu\)m：\(hc/\lambda k=2.88\times10^3\)
K。\(x_1=0.498\)、\(x_2=0.300\)。比
\(=\dfrac{e^{0.498}-1}{e^{0.300}-1}=\dfrac{0.645}{0.350}\simeq1.85\)。\(x\ll1\)
で \(e^x-1\simeq x\) ゆえ比 \(\simeq x_1/x_2=T_2/T_1=9602/5778=1.66\)
に近く、Rayleigh-Jeans 側では温度比のみが効く（指数増幅なし）。

\textbf{答え}：比は可視 7.6、紫外 53、赤外 1.85。Wien
側（紫外）ほど温度差が指数的に効き、RJ 側（赤外）は温度比のみ。

\subsection*{問題 17-3　ダスト修正黒体から3物理量を読む}\label{ex-17-3}
\addcontentsline{toc}{subsection}{問題
17-3　ダスト修正黒体から3物理量を読む}

{[}★ 難易度：☆☆ {]}

§17.3 の修正黒体
\(I_\nu=Q_\nu B_\nu(T_d)\)、\(Q_\nu\propto\nu^\beta\)。観測 SED
から3つの物理量を読む手順を整理する。

\begin{enumerate}
\def\labelenumi{\arabic{enumi}.}
\tightlist
\item
  ダスト温度 \(T_d\)、ダスト放射率指数 \(\beta\)、ダスト質量 がそれぞれ
  SED のどの特徴から読めるか述べよ。
\item
  \(\beta\) の値が粒子サイズを反映する（大粒子 \(\beta\to0\)、小結晶
  \(\beta\to2\)）。ALMA で原始惑星系円盤の中心部ほど \(\beta\)
  が小さい観測が、何を示すか述べよ。
\item
  多温度フィットで「冷たい成分の質量」を推定するとき、温度との縮退に注意が要る。長波長観測点がこの縮退を解く理由を、修正黒体の
  RJ 側（\hyperref[ex-15-2]{演習 15-2}）と結びつけて述べよ。
\end{enumerate}

\textbf{関連}：\hyperref[sec-dust-modified-blackbody]{§17.3
ダストの修正黒体}、\hyperref[sec-multi-temperature]{§17.5
多温度合成}／修正黒体のべきは\hyperref[ex-15-2]{演習
15-2}、修正黒体の定義は\hyperref[sec-non-blackbody-modified]{§15.4}。

\section{模範解答（問題
17-3）}\label{ux6a21ux7bc4ux89e3ux7b54ux554fux984c-17-3}

\textbf{(1)} \(T_d\)：Wien 側の指数減衰の鋭さ、または SED
のピーク位置（Wien 則）から。\(\beta\)：Rayleigh-Jeans
側（長波長）の勾配の急さ（\(I_\nu\propto\nu^{2+\beta}\)）から。ダスト質量：全光度を
\(B_\nu(T_d)\) で割ったときの規格化定数から（不透明度の仮定が必要）。

\textbf{(2)} \(\beta\) が小さいほど粒子が大きい（粒子サイズ≈波長の極限で
\(\beta\to0\)）。円盤中心部ほど \(\beta\)
が小さい＝粒子が大きく成長している。これは円盤内でダストが集積・成長して惑星形成へ向かう過程の直接観測的証拠。

\textbf{(3)} 同じピーク位置でも（高温・少量）と（低温・大量）で似た SED
を作れる縮退がある。RJ 側の勾配 \(\nu^{2+\beta}\)
と絶対強度は温度・質量・\(\beta\)
に異なる依存をするので、ピークより長波長（サブミリ・ミリ）の観測点を加えると、RJ
側の形が温度と質量を分離して縮退を解く。だから長波長観測が鍵。

\textbf{答え}：\(T_d\)＝Wien側/ピーク、\(\beta\)＝RJ側勾配、質量＝規格化。中心部の小
\(\beta\)＝粒子成長＝惑星形成の証拠。長波長点が温度-質量縮退を解く。

\subsection*{問題 17-4　ダストのピーク波長と観測施設}\label{ex-17-4}
\addcontentsline{toc}{subsection}{問題
17-4　ダストのピーク波長と観測施設}

{[}★ 難易度：☆ {]}

§17.5 の問いを一般化する。Wien 則
\(\lambda_\mathrm{peak}=b/T\)（\(b=2898\ \mu\)m·K）。

\begin{enumerate}
\def\labelenumi{\arabic{enumi}.}
\tightlist
\item
  \(T_d=20\) K
  の冷たいダストのピーク波長を求め、どの観測施設の帯域か述べよ。
\item
  暖かいダスト \(T_d=100\) K、ごく冷たいダスト \(T_d=15\) K
  のピーク波長も求めよ。
\item
  銀河の重元素の半分がダストに閉じ込められ、その熱放射が遠赤外〜サブミリにピークを持つ。なぜ赤外・サブミリ天文学（Herschel/ALMA）が宇宙の物質収支に決定的なのか述べよ。
\end{enumerate}

\textbf{関連}：\hyperref[sec-multi-temperature]{§17.5
多温度合成}、\hyperref[sec-infrared-astronomy]{§17.4 赤外線天文学}／Wien
変位則は\hyperref[ex-2-2]{演習 2-2}。

\section{模範解答（問題
17-4）}\label{ux6a21ux7bc4ux89e3ux7b54ux554fux984c-17-4}

\textbf{(1)} \(\lambda_\mathrm{peak}=2898/20=145\ \mu\)m。これは
Herschel（SPIRE/PACS、60〜600 μm）の主帯域。冷たい星間ダストの観測が
Herschel で飛躍したのはこのため。

\textbf{(2)} \(T_d=100\)
K：\(\lambda_\mathrm{peak}=2898/100=29\ \mu\)m（中間赤外、Spitzer/JWST）。\(T_d=15\)
K：\(\lambda_\mathrm{peak}=2898/15=193\ \mu\)m（サブミリ、Herschel/ALMA）。

\textbf{(3)} ダスト（\(T\sim15\)〜\(30\)
K）の熱放射ピークは遠赤外〜サブミリ。銀河の重元素の半分がダストに閉じ込められているので、可視・紫外だけ見ると物質の半分を見落とす。遠赤外・サブミリ（Herschel/ALMA）でこの放射を捉えて初めて、銀河の星間物質・ダスト質量・星形成の全体像（物質収支）が分かる。

\textbf{答え}：20 K→145 μm（Herschel）、100 K→29 μm、15 K→193
μm。重元素の半分がダスト＝遠赤外/サブミリ観測が物質収支に必須。

\subsection*{問題
17-5　電波HII領域の自由--自由スペクトルのターンオーバー}\label{ex-17-5}
\addcontentsline{toc}{subsection}{問題
17-5　電波HII領域の自由--自由スペクトルのターンオーバー}

{[}★ 難易度：☆☆ {]}

§17.6 の電離ガスの自由--自由放射を扱う。等温スラブ（電子温度
\(T_e\)、LTE）の輝度温度
\(T_b(\nu)=T_e(1-e^{-\tau_\nu})\)、自由--自由の光学的厚さ
\(\tau_\nu\propto\nu^{-2.1}\)（電波帯）とする。

\begin{enumerate}
\def\labelenumi{\arabic{enumi}.}
\tightlist
\item
  低振動数（光学的に厚い、\(\tau_\nu\gg1\)）で
  \(T_b\to T_e\)、\(I_\nu\to B_\nu(T_e)\)、フラックスが
  \(F_\nu\propto\nu^2\) になることを示せ。これが「電子温度 \(T_e\)
  の黒体（Rayleigh--Jeans 裾）」である理由を、第5章 §5.6
  と結びつけて述べよ。
\item
  高振動数（光学的に薄い、\(\tau_\nu\ll1\)）で
  \(T_b\to T_e\tau_\nu\)、\(F_\nu\propto\nu^{-0.1}\)
  になることを示せ（Rayleigh--Jeans で \(F_\nu\propto\nu^2 T_b\)
  を使う）。
\item
  「自由--自由放射は熱的なのに、全帯域で単一の黒体にはならない」のはなぜか。低
  \(\nu\) 側と高 \(\nu\) 側で何が違うかを一言で述べよ。
\end{enumerate}

\textbf{関連}：\hyperref[sec-free-free-ionized]{§17.6
電離ガスの自由--自由放射}／光学的に厚い極限は\hyperref[sec-radiative-transfer-thick]{§5.6}（\hyperref[ex-5-2]{演習
5-2}）、輝度温度は\hyperref[sec-reading-spectra-temperatures]{§2.6}、熱的/非熱的は\hyperref[sec-where-blackbody-thermal]{§1.9}、X
線制動放射は\href{../part6/18-accretion-high-energy.qmd}{第18章}。

\section{模範解答（問題
17-5）}\label{ux6a21ux7bc4ux89e3ux7b54ux554fux984c-17-5}

\textbf{(1)} \(\tau_\nu\gg1\) で \(1-e^{-\tau_\nu}\to1\) ゆえ
\(T_b\to T_e\)、すなわち
\(I_\nu=B_\nu(T_e)(1-e^{-\tau_\nu})\to B_\nu(T_e)\)。これは第5章 §5.6
の「光学的に厚い極限で \(I_\nu\to S_\nu=B_\nu(T)\)」そのもの（LTE で
\(S_\nu=B_\nu(T_e)\)）。電波帯は \(h\nu\ll k_BT_e\) なので
\(B_\nu(T_e)\simeq2\nu^2k_BT_e/c^2\)（Rayleigh--Jeans）、よって
\(F_\nu\propto\nu^2\)。観測されるのは電子温度 \(T_e\) の黒体の RJ 裾。

\textbf{(2)} \(\tau_\nu\ll1\) で \(1-e^{-\tau_\nu}\simeq\tau_\nu\) ゆえ
\(T_b\to T_e\tau_\nu\propto\nu^{-2.1}\)。フラックスは RJ で
\(F_\nu\propto\nu^2 T_b\propto\nu^2\cdot\nu^{-2.1}=\nu^{-0.1}\)（ほぼ平坦）。

\textbf{(3)} 黒体になるのは光学的に厚い低 \(\nu\) 側だけ。高 \(\nu\)
側は光学的に薄く \(I_\nu\approx\tau_\nu B_\nu(T_e)\ll B_\nu(T_e)\)
で黒体から下にずれる。「熱的（粒子が Maxwell
分布）」と「黒体（光学的に厚い）」は別条件で、自由--自由放射は前者を常に、後者を低
\(\nu\) 側でのみ満たす。

\textbf{答え}：低 \(\nu\)（厚い）→
\(T_b\to T_e\)、\(I_\nu\to B_\nu(T_e)\)、\(F_\nu\propto\nu^2\)（\(T_e\)
の黒体）。高 \(\nu\)（薄い）→
\(F_\nu\propto\nu^{-0.1}\)。黒体化は光学的に厚い低 \(\nu\) 側のみ。

\section*{さらに学ぶための参考文献}\label{sec-stars-planets-dust-further}
\addcontentsline{toc}{section}{さらに学ぶための参考文献}

\markright{さらに学ぶための参考文献}

\begin{itemize}
\tightlist
\item
  Carroll \& Ostlie, \emph{An Introduction to Modern Astrophysics}
  (Cambridge, 2nd ed., 2017) --- 恒星・惑星・銀河の標準的扱い
\item
  Pierrehumbert, \emph{Principles of Planetary Climate} (Cambridge,
  2010) --- 惑星平衡温度と温室効果の物理
\item
  Draine, \emph{Physics of the Interstellar and Intergalactic Medium}
  (Princeton, 2011) --- ダストの放射、\(\beta\)、修正黒体
\item
  Galametz et al.~の SED フィッティングのレビュー（A\&A, 2012）---
  多温度 SED の実用解析
\item
  Wilson, Rohlfs \& Hüttemeister, \emph{Tools of Radio Astronomy}
  (Springer, 6th ed., 2013) --- 電波連続光・自由--自由放射と HII
  領域のスペクトル
\item
  Rybicki \& Lightman, \emph{Radiative Processes in Astrophysics}
  (Wiley, 1979) ---
  第5章「Bremsstrahlung」（自由--自由放射の断面積と吸収）
\end{itemize}

\chapter{第18章　降着円盤と高エネルギー天体}\label{ux7b2c18ux7ae0-ux964dux7740ux5186ux76e4ux3068ux9ad8ux30a8ux30cdux30ebux30aeux30fcux5929ux4f53}

\textbf{この章の主題}：第VI部の最終章として、\textbf{プランク分布の枠組みが極端な物理条件
― 強重力場・強磁場・相対論的速度 ― の下でどこまで通用するか}
を、高エネルギー天体（X
線連星、AGN、中性子星、ブラックホール）を題材に検討する。本章で第I〜V部の道具立てを最大に拡張し、第VII部（線スペクトル）と第VIII部（QED）への準備が整う。

\section*{この章の中心地図}\label{sec-accretion-high-energy-map}
\addcontentsline{toc}{section}{この章の中心地図}

\markright{この章の中心地図}

\textbf{中心地図：プランク関数 ― 本部の主題は「ずれ」}

\[
B_\nu(T) = \frac{2 h \nu^3}{c^2} \cdot \frac{1}{e^{h\nu/kT} - 1} \quad \overset{\text{現実}}{\longrightarrow} \quad \text{?}
\]

第VI部（現実の宇宙へ戻る）は、観測される現実の天体スペクトルが、この理想的プランク関数からどのように、そしてなぜ
\textbf{ずれる}かを扱う。ずれの主な起源：

\begin{itemize}
\tightlist
\item
  \textbf{振動数依存不透明度} → 修正黒体（第15章）
\item
  \textbf{多温度成分の重ね合わせ} → 多温度黒体スペクトル（第17・18章）
\item
  \textbf{散乱効果}（コンプトン化など）→ 高エネルギー側の変形（第18章）
\item
  \textbf{非熱的放射の混入} → シンクロトロン・ベキ則の重なり
\end{itemize}

本部は、観測スペクトルを「理想的黒体＋ずれ」として読み解き、ずれの定量化から物理状態を引き出す視点を養う。

\begin{tcolorbox}[enhanced jigsaw, arc=.35mm, bottomrule=.15mm, bottomtitle=1mm, breakable, colback=white, colbacktitle=quarto-callout-note-color!10!white, colframe=quarto-callout-note-color-frame, coltitle=black, left=2mm, leftrule=.75mm, opacityback=0, opacitybacktitle=0.6, rightrule=.15mm, title=\textcolor{quarto-callout-note-color}{\faInfo}\hspace{0.5em}{ノート}, titlerule=0mm, toprule=.15mm, toptitle=1mm]

\textbf{方針}：本章では「局所的にプランク分布 ＋
大域的な変形」という構造が繰り返し現れる。各天体について、\textbf{「どこが局所黒体で、どこから外れるか」}
を明示的に区別して進む。

\end{tcolorbox}

本章の高エネルギー天体の逆引きを、あらかじめ一表にまとめておく：

{\def\LTcaptype{none} % do not increment counter
\begin{longtable}[]{@{}
  >{\raggedright\arraybackslash}p{(\linewidth - 6\tabcolsep) * \real{0.2500}}
  >{\raggedright\arraybackslash}p{(\linewidth - 6\tabcolsep) * \real{0.2500}}
  >{\raggedright\arraybackslash}p{(\linewidth - 6\tabcolsep) * \real{0.2500}}
  >{\raggedright\arraybackslash}p{(\linewidth - 6\tabcolsep) * \real{0.2500}}@{}}
\toprule\noalign{}
\begin{minipage}[b]{\linewidth}\raggedright
観測で見える特徴
\end{minipage} & \begin{minipage}[b]{\linewidth}\raggedright
最小モデル
\end{minipage} & \begin{minipage}[b]{\linewidth}\raggedright
逆引きされる物理量
\end{minipage} & \begin{minipage}[b]{\linewidth}\raggedright
破れる仮定
\end{minipage} \\
\midrule\noalign{}
\endhead
\bottomrule\noalign{}
\endlastfoot
連続光の \(\nu^{1/3}\) べき則 & 多温度円盤 \(T \propto R^{-3/4}\) &
降着率・内縁半径 & 局所黒体・薄い円盤 \\
高エネルギー側のべき則テイル & コンプトン化（\(y\) パラメータ） &
コロナ温度 \(T_e\)・光学的深さ \(\tau_e\) & 単純な局所黒体 \\
硬 X 線スペクトルと状態遷移 & コロナ＋円盤の幾何変化 &
コロナ・円盤の幾何 & 定常・単一成分 \\
非対称な鉄輝線輪郭 & 相対論的円盤線 & ブラックホールスピン・内縁半径 &
非相対論的・静的 \\
\end{longtable}
}

\section*{この章で答える問い}\label{sec-accretion-high-energy-questions}
\addcontentsline{toc}{section}{この章で答える問い}

\markright{この章で答える問い}

\begin{itemize}
\tightlist
\item
  降着円盤の各半径が「局所的に」黒体になるのは、どんな仮定のもとでか
\item
  多温度円盤の合成スペクトルがべき関数的に見えるのはなぜか
\item
  コンプトン化は黒体スペクトルをどう変形するのか
\item
  中性子星表面やブラックホール近傍では、何が黒体放射の理想からずれるのか
\end{itemize}

\section*{到達目標}\label{sec-accretion-high-energy-goals}
\addcontentsline{toc}{section}{到達目標}

\markright{到達目標}

この章を読み終えると、読者は次のことができるようになる：

\begin{itemize}
\tightlist
\item
  標準降着円盤モデルから、多温度黒体スペクトルを構成できる
\item
  コンプトン化による高エネルギー側のスペクトル変形を質的に予想できる
\item
  強い重力場や強い磁場が、観測スペクトルにどう影響するかを述べられる
\end{itemize}

\begin{center}\rule{0.5\linewidth}{0.5pt}\end{center}

\section{18.1 局所黒体近似 ―
各半径に温度を割り当てる}\label{sec-disk-local-blackbody}

{[}本文目安：B3-B4{]}

降着円盤（\textbf{accretion
disk}）の物理的設定：中心天体（白色矮星・中性子星・ブラックホール）の周りで、重力ポテンシャルエネルギーを解放しながら降着するガス。エネルギー解放率（光度）は降着率
\(\dot{M}\) と中心天体の重力 (\(GM/R\)) で決まる。

\textbf{標準降着円盤}（Shakura-Sunyaev 1973）の核心仮定：

\begin{enumerate}
\def\labelenumi{\arabic{enumi}.}
\tightlist
\item
  \textbf{薄い円盤}：高さ \(H \ll R\)（軌道半径）
\item
  \textbf{回転は Kepler 速度}：粘性が小さく軌道は摂動程度
\item
  \textbf{局所熱平衡（LTE）}：各半径 \(R\)
  で、ガスが局所的に熱化し、その温度 \(T(R)\) のプランク分布で放射
\end{enumerate}

この仮定 3 が \textbf{局所黒体近似}。各半径から \textbf{温度 \(T(R)\)
の黒体スペクトル} が放たれる：

\begin{equation}\protect\phantomsection\label{eq-disk-local-blackbody}{
I_\nu(R) = B_\nu(T(R))
}\end{equation}

ただし、\textbf{\(T(R)\)
は半径によって変わる}。これが多温度合成スペクトル（次節）を生む。

\begin{tcolorbox}[enhanced jigsaw, arc=.35mm, bottomrule=.15mm, breakable, colback=white, colframe=quarto-callout-warning-color-frame, left=2mm, leftrule=.75mm, opacityback=0, rightrule=.15mm, toprule=.15mm]
\begin{minipage}[t]{5.5mm}
\textcolor{quarto-callout-warning-color}{\faExclamationTriangle}
\end{minipage}%
\begin{minipage}[t]{\textwidth - 5.5mm}

\textbf{注意（観測）}：「局所的に黒体」は\textbf{強い仮定}。実際の降着円盤では：

\begin{itemize}
\tightlist
\item
  高温・希薄なコロナ層（次節後半）
\item
  不透明度が振動数依存（修正黒体的なずれ）
\item
  散乱が支配的（電子散乱が吸収より卓越する高温体）
\end{itemize}

これらは局所黒体からのずれとして現れる。\(T(R)\) から実際の
\textbf{色温度} \(T_\text{col}(R)\) への補正係数
\(f_\text{col} \simeq 1.5\)〜\(2\) が経験的に必要

\end{minipage}%
\end{tcolorbox}

\section{\texorpdfstring{18.2 多温度円盤 ― \(\nu^{1/3}\)
べき則の出現}{18.2 多温度円盤 ― \textbackslash nu\^{}\{1/3\} べき則の出現}}\label{sec-disk-multi-temperature}

{[}本文目安：B4{]}

標準降着円盤では、エネルギー収支から温度の半径依存が

\begin{equation}\protect\phantomsection\label{eq-disk-temperature-profile}{
T(R) \propto R^{-3/4}
}\end{equation}

となる（薄い円盤の Stefan-Boltzmann
バランス）。内側ほど高温、外側ほど低温。

全円盤からの合成スペクトルは、円盤面の各半径からの放射を面積分した形で

\begin{equation}\protect\phantomsection\label{eq-disk-total-flux}{
F_\nu \propto \int_{R_\text{in}}^{R_\text{out}} 2\pi R \cdot B_\nu(T(R))\, dR
}\end{equation}

と書ける。ここで \(F_\nu\) は \textbf{比例関係}
として理解すべき量で、右辺には観測者までの距離 \(d\)、円盤の傾斜角
\(i\)（投影
\(\cos i\)）、両面放射の有無といった幾何因子を含めていない。実際に観測されるフラックス
\(F_\nu^\text{obs}\)
は、面フラックス（円盤単位面積あたりの放射）にこれらの幾何因子を掛けたもので、たとえば
\(F_\nu^\text{obs} \simeq (\cos i / d^2)\int 2\pi R\, B_\nu(T(R))\, dR\)
のように書ける。本節で着目するのは絶対値ではなく振動数依存（スペクトルの形）なので、以下では比例関係で扱う。

ここで重要なのは \textbf{中間振動数帯}
の振る舞い。\(k_B T(R_\text{in}) \gg h\nu \gg k_B T(R_\text{out})\)（あるいは同値に
\(T(R_\text{in}) \gg h\nu/k_B \gg T(R_\text{out})\)）となる振動数では、温度分布が積分を支配する。

式~\ref{eq-disk-temperature-profile}
を使って変数変換すると（導出概略は下記の問い、丁寧な計算は Frank, King
\& Raine 第 5 章参照）

\begin{equation}\protect\phantomsection\label{eq-disk-13-powerlaw}{
F_\nu \propto \nu^{1/3} \quad \text{(中間帯)}
}\end{equation}

この \textbf{\(\nu^{1/3}\) べき則} は、多温度円盤の
\textbf{特徴的シグネチャ}。観測される X 線連星や活動銀河中心核 (AGN)
の連続光に、この勾配が広い範囲で見られる。

\begin{tcolorbox}[enhanced jigsaw, arc=.35mm, bottomrule=.15mm, bottomtitle=1mm, breakable, colback=white, colbacktitle=quarto-callout-note-color!10!white, colframe=quarto-callout-note-color-frame, coltitle=black, left=2mm, leftrule=.75mm, opacityback=0, opacitybacktitle=0.6, rightrule=.15mm, title=\textcolor{quarto-callout-note-color}{\faInfo}\hspace{0.5em}{ノート}, titlerule=0mm, toprule=.15mm, toptitle=1mm]

\textbf{観測との対応}：「べき則
\(\nu^{1/3}\)」を見たら、\textbf{多数の異なる温度の黒体成分の合成}
を疑う ―
これは熱的物理の結果。同じく「べき則」と見えるが、非熱的シンクロトロンや逆コンプトン
(§18.3) では指数が違うので、区別がつく

\end{tcolorbox}

\begin{tcolorbox}[enhanced jigsaw, arc=.35mm, bottomrule=.15mm, breakable, colback=white, colframe=quarto-callout-tip-color-frame, left=2mm, leftrule=.75mm, opacityback=0, rightrule=.15mm, toprule=.15mm]
\begin{minipage}[t]{5.5mm}
\textcolor{quarto-callout-tip-color}{\faLightbulb}
\end{minipage}%
\begin{minipage}[t]{\textwidth - 5.5mm}

\textbf{問い}：\(\nu^{1/3}\) という指数 1/3 は、どこから来るか？

\textbf{短答}：\(T \propto R^{-3/4}\) より
\(R \propto T^{-4/3}\)、\(2\pi R\, dR \propto T^{-11/3} dT\)。中間振動数帯では
Rayleigh-Jeans 側（\(B_\nu \propto \nu^2 T\)。以下 RJ
側）が支配する範囲があり、\(\int T \cdot T^{-11/3} dT\)
で適度な温度範囲を積分すると、最終的に \(F_\nu \propto \nu^{1/3}\)
となる。

\textbf{もう一歩}：もし温度プロファイルが \(T \propto R^{-\alpha}\)
なら、\$F\_\nu \propto \nu\^{}\{(3\alpha - 2)/\alpha\} \$
という形になり、\(\alpha = 3/4\) で \(\nu^{1/3}\)、\(\alpha = 1\) で
\(\nu^1\) など、円盤の物理に応じて勾配が変わる

\end{minipage}%
\end{tcolorbox}

\section{18.3 コンプトン化 ―
高温電子からの逆風}\label{sec-comptonization}

{[}本文目安：B4{]}

降着円盤の中心近くには、しばしば
\textbf{高温電子コロナ}（\(T_e \sim 10^8\)〜\(10^9\) K、\(\sim 100\)
keV）が存在する。低エネルギーのプランク光子がこのコロナを通過すると、\textbf{逆コンプトン散乱}（\textbf{inverse
Compton scattering}）で電子からエネルギーを受け取る。

低温の黒体光子（種光子、\(h\nu_\text{seed}\)）が熱い電子（\(k_B T_e \gg h\nu_\text{seed}\)）と散乱した一回あたりのエネルギー利得：

\begin{equation}\protect\phantomsection\label{eq-compton-gain-per-scattering}{
\frac{\Delta \nu}{\nu} \simeq \frac{4 k_B T_e}{m_e c^2}
}\end{equation}

光子が脱出までに \(N\) 回散乱すると、累積利得
\({(1 + \Delta\nu/\nu)}^N\)。コロナの光学的深さ \(\tau_e\)
と温度から見積もる \textbf{Compton \(y\) パラメータ}：

\begin{equation}\protect\phantomsection\label{eq-compton-y}{
y \equiv \frac{4 k_B T_e}{m_e c^2} \max(\tau_e, \tau_e^2)
}\end{equation}

\begin{itemize}
\tightlist
\item
  \(y \ll 1\)：光子はわずかにエネルギー利得 ―
  元の黒体形のわずかな歪み（\textbf{Wien 側へのテイル}）
\item
  \(y \gtrsim 1\)：光子はコロナで完全に再分配される ― \textbf{Wien 分布}
  に近い高温化スペクトル
\item
  中間域：\textbf{べき則 \(F_\nu \propto \nu^{-\alpha}\)} が現れ、Wien
  カットオフで切れる
\end{itemize}

\begin{tcolorbox}[enhanced jigsaw, arc=.35mm, bottomrule=.15mm, bottomtitle=1mm, breakable, colback=white, colbacktitle=quarto-callout-note-color!10!white, colframe=quarto-callout-note-color-frame, coltitle=black, left=2mm, leftrule=.75mm, opacityback=0, opacitybacktitle=0.6, rightrule=.15mm, title=\textcolor{quarto-callout-note-color}{\faInfo}\hspace{0.5em}{ノート}, titlerule=0mm, toprule=.15mm, toptitle=1mm]

\textbf{観測との対応}：X
線連星のスペクトルは、二つの状態とその間の遷移として整理できる。

\begin{itemize}
\tightlist
\item
  \textbf{観測特徴}：\textbf{ハード状態}（hard
  state）は高エネルギー側（硬 X 線）が相対的に強い状態で、観測される指数
  \(\alpha \simeq 0.5\)〜\(1.5\) のべき則と \(\sim 100\) keV
  の急速減衰を示す。\textbf{ソフト状態}（soft
  state）は円盤黒体成分が卓越し、軟 X 線側が強い状態である。
\item
  \textbf{物理的解釈}：ハード状態は高温電子コロナによるコンプトン化が支配的で、\(T_e\)・\(\tau_e\)
  の組合せがスペクトルを決める。ソフト状態は多温度円盤（§18.2）の熱的放射が支配的である。
\item
  \textbf{状態遷移}：両状態の間の遷移は、コロナと円盤の幾何（内縁半径・コロナの量）の変化を反映すると解釈される。
\end{itemize}

\end{tcolorbox}

第16章 §16.5 の \textbf{\(y\) 歪み}（CMB）と本節の \(y\) パラメータは
\textbf{完全に同じ物理}。スケールが宇宙論的か、銀河近傍かの違いだけ。

\section{18.4 中性子星表面と強磁場}\label{sec-neutron-star-surface}

{[}本文目安：B4｜M1{]}

中性子星表面からの X
線熱放射は、原理的にはプランク分布。しかし強い物理条件が複合する：

\begin{itemize}
\tightlist
\item
  \textbf{温度} \(T \sim 10^6\)〜\(10^7\) K → ピークは X 線帯
\item
  \textbf{重力赤方偏移} \(z_g = (1-2GM/Rc^2)^{-1/2} - 1 \simeq 0.3\) →
  観測温度は \(T_\infty = T / (1+z_g) \simeq T/1.3\)
\item
  \textbf{強磁場} \(B \sim 10^{12}\) G（パルサー）〜\(10^{15}\)
  G（マグネター）：\textbf{サイクロトロン共鳴吸収線}、量子効果による不透明度変形
\item
  \textbf{大気組成} ：水素薄層、ヘリウム、重元素 → 振動数依存不透明度
\end{itemize}

観測される X 線スペクトルは、これら効果が重ねあわされた
\textbf{「修正黒体」} とも呼ぶべき形状。観測から逆引きされる物理量：

\begin{itemize}
\tightlist
\item
  表面有効温度 \(T_\infty\)
\item
  質量と半径の組合せ（重力赤方偏移と Stefan-Boltzmann を連立）
\item
  磁場の方向と大きさ（サイクロトロン線の振動数）
\item
  大気組成（吸収線パターン）
\end{itemize}

\begin{tcolorbox}[enhanced jigsaw, arc=.35mm, bottomrule=.15mm, breakable, colback=white, colframe=quarto-callout-warning-color-frame, left=2mm, leftrule=.75mm, opacityback=0, rightrule=.15mm, toprule=.15mm]
\begin{minipage}[t]{5.5mm}
\textcolor{quarto-callout-warning-color}{\faExclamationTriangle}
\end{minipage}%
\begin{minipage}[t]{\textwidth - 5.5mm}

\textbf{注意（観測）}：\textbf{中性子星半径の精密測定}
は核物質の状態方程式を制約する重要な物理。NICER
衛星（2017〜）はパルサーの X 線スペクトルと時間プロファイルから \(R\) と
\(M\)
を独立に測定する観測戦略を採用している。本書の射程である「スペクトルから物理を逆引き」の最も精密な現代例

\end{minipage}%
\end{tcolorbox}

\section{18.5 ブラックホール近傍 ―
一般相対論的修正}\label{sec-black-hole-effects}

{[}本文目安：M1-M2{]}

ブラックホール降着円盤の最も内側（\(R \sim 10\,GM/c^2\)）では、一般相対論的効果が無視できない：

\begin{enumerate}
\def\labelenumi{\arabic{enumi}.}
\tightlist
\item
  \textbf{重力赤方偏移}：光子のエネルギーが減る
  \(\nu \to \nu(1 - 2GM/Rc^2)^{1/2}\)
\item
  \textbf{特殊相対論的ビーミング}：円盤のケプラー回転速度 \(v \to c\)
  で、近づく側が明るく・遠ざかる側が暗くなる
\item
  \textbf{光線の曲がり}：円盤の裏側が見える、Einstein リング
\item
  \textbf{時間遅れ}：円盤内側ほど時間がゆっくり進む
\end{enumerate}

これらの効果は、第13章 §13.4 の偏光自由度や §18.1 の局所黒体仮定の
\textbf{上に重ねて} 効く。観測される SED は

\begin{equation}\protect\phantomsection\label{eq-grdisk-flux}{
F_\nu^\text{obs} = \int B_\nu(T(R, \phi)) \cdot g^3(R, \phi) \cdot dA_\text{proper}
}\end{equation}

の形（\(g\) は赤方偏移因子、\(dA_\text{proper}\)
は固有面積）。一般相対論的計算が必要。

特に重要な観測量は \textbf{広線輪郭} ― 鉄輝線 (6.4 keV)
が、円盤の高速回転と重力赤方偏移により、特徴的な非対称形状（青翼が高く、赤翼が長く伸びる）に変形される。これがブラックホールスピンの測定に使われる。

\begin{tcolorbox}[enhanced jigsaw, arc=.35mm, bottomrule=.15mm, bottomtitle=1mm, breakable, colback=white, colbacktitle=quarto-callout-note-color!10!white, colframe=quarto-callout-note-color-frame, coltitle=black, left=2mm, leftrule=.75mm, opacityback=0, opacitybacktitle=0.6, rightrule=.15mm, title=\textcolor{quarto-callout-note-color}{\faInfo}\hspace{0.5em}{ノート}, titlerule=0mm, toprule=.15mm, toptitle=1mm]

\textbf{観測との対応}：第VII部の主題である \textbf{線スペクトル}
が、ここで連続スペクトルと再会する。鉄 K\(\alpha\) 線（中性鉄 6.4
keV）の輪郭は、原子物理（線のエネルギーと振動子強度）＋
円盤運動学（Doppler 効果）＋ 一般相対論（重力赤方偏移）の
\textbf{三層の物理}
が刻まれた観測量。第VII部では原子物理側を、本章では円盤物理側を別々に整理しておく

\end{tcolorbox}

\section{18.6
第V部までの道具で見える限界と次への橋}\label{sec-accretion-bridge}

{[}本文目安：B3{]}

本章で扱った高エネルギー天体は、本書の第I〜V部の枠組み（プランク分布＋ずれ因子）で観測の
\textbf{大半が} 記述できることを示した。しかし：

\begin{itemize}
\tightlist
\item
  \textbf{コロナでのコンプトン化} ―
  電子分布が熱的かほぼ熱的だが、放射過程は QED 的（第VIII部で本格議論）
\item
  \textbf{シンクロトロン放射} ― AGN
  ジェットやマグネター磁気圏では、非熱的電子分布と強磁場が組合さって、本書の枠組み外の物理が支配的
\item
  \textbf{線スペクトルの広がり} ―
  第VII部の主題と、本章の運動学的・重力的効果が組合さる
\end{itemize}

\begin{tcolorbox}[enhanced jigsaw, arc=.35mm, bottomrule=.15mm, bottomtitle=1mm, breakable, colback=white, colbacktitle=quarto-callout-note-color!10!white, colframe=quarto-callout-note-color-frame, coltitle=black, left=2mm, leftrule=.75mm, opacityback=0, opacitybacktitle=0.6, rightrule=.15mm, title=\textcolor{quarto-callout-note-color}{\faInfo}\hspace{0.5em}{ノート}, titlerule=0mm, toprule=.15mm, toptitle=1mm]

\textbf{使えるようになった道具}（§18.1〜§18.5）：

\begin{itemize}
\tightlist
\item
  局所黒体仮定とその限界
\item
  多温度円盤の \(\nu^{1/3}\) べき則（合成スペクトルの特徴シグネチャ）
\item
  Compton 化と \(y\) パラメータ（X 線連星のハード状態スペクトル）
\item
  中性子星表面物理（重力赤方偏移、強磁場、大気組成）
\item
  ブラックホール近傍の一般相対論的修正
\end{itemize}

\end{tcolorbox}

\begin{center}\rule{0.5\linewidth}{0.5pt}\end{center}

\section*{この章で何がわかったか}\label{sec-accretion-high-energy-summary}
\addcontentsline{toc}{section}{この章で何がわかったか}

\markright{この章で何がわかったか}

\textbf{中心地図に戻る・第I〜VI部の総まとめ}

本章で、プランク分布の枠組みが、もっとも極端な物理条件（強重力場・強磁場・相対論的速度）の下でも、\textbf{局所黒体
＋ 大域的修正} の組合せとして大きく適用できることを見た。

第I〜V部で逆引きしてきたプランク分布の理解 ―
統計的源泉、量子的源泉、電磁気的源泉 ―
が、第VI部の応用章群で次々と観測スペクトルの解読に動員された：

\begin{itemize}
\tightlist
\item
  CMB（第16章）：完璧な黒体、熱化の証拠
\item
  恒星・惑星・ダスト（第17章）：振動数依存不透明度、修正黒体、平衡温度
\item
  降着円盤（第18章）：局所黒体、多温度合成、コンプトン化、相対論的修正
\end{itemize}

これで第I〜VI部は完結。\textbf{第VII部} では、本書のもう一つの主柱 ―
\textbf{線スペクトル} ― の逆引きが始まる。Einstein
係数（第11章）と水素線吸収係数の地図が再登場する

\section*{演習問題}\label{sec-accretion-high-energy-exercises}
\addcontentsline{toc}{section}{演習問題}

\markright{演習問題}

\textbf{模範解答}
は折りたたみを展開して確認できる（オンライン版）。まず自力で解いてから開くこと。

\subsection*{問題 18-1　多温度円盤の 1/3 乗べき則}\label{ex-18-1}
\addcontentsline{toc}{subsection}{問題 18-1　多温度円盤の 1/3 乗べき則}

{[}★ 難易度：☆☆☆ {]}

§18.2 は標準降着円盤の中間振動数帯で \(F_\nu\propto\nu^{1/3}\)
になると述べた。導出の骨子を埋める。温度プロファイル
\(T(R)\propto R^{-3/4}\)、合成フラックス
\(F_\nu=\displaystyle\int_{R_\mathrm{in}}^{R_\mathrm{out}}2\pi R\,B_\nu(T(R))\,dR\)。

\begin{enumerate}
\def\labelenumi{\arabic{enumi}.}
\tightlist
\item
  \(T\propto R^{-3/4}\) から
  \(R\propto T^{-4/3}\)、\(2\pi R\,dR\propto T^{-11/3}\,dT\) を示せ。
\item
  中間帯（\(kT(R_\mathrm{in})\gg h\nu\gg kT(R_\mathrm{out})\)）では各半径の
  Rayleigh-Jeans 側 \(B_\nu\propto\nu^2 T\)
  が効く。\(F_\nu\propto\nu^2\displaystyle\int T\cdot T^{-11/3}\,dT=\nu^2\int T^{-8/3}dT\)
  を変数 \(x=h\nu/kT\) に変換し、\(F_\nu\propto\nu^{1/3}\) を導け。
\item
  温度プロファイルが一般に \(T\propto R^{-\alpha}\) なら
  \(F_\nu\propto\nu^{(3\alpha-2)/\alpha}\)
  となることを確かめ、\(\alpha=3/4\) で \(1/3\) になることを示せ。
\end{enumerate}

\textbf{関連}：\hyperref[sec-disk-multi-temperature]{§18.2
多温度円盤}／局所黒体は\hyperref[sec-disk-local-blackbody]{§18.1}（\hyperref[ex-18-3]{演習
18-3}）、RJ 極限は\hyperref[sec-reading-spectra-limits]{§2.5}。

\section{模範解答（問題
18-1）}\label{ux6a21ux7bc4ux89e3ux7b54ux554fux984c-18-1}

\textbf{(1)}
\(T\propto R^{-3/4}\Rightarrow R\propto T^{-4/3}\)。\(dR\propto T^{-7/3}dT\)。\(2\pi R\,dR\propto T^{-4/3}\cdot T^{-7/3}dT=T^{-11/3}dT\)。

\textbf{(2)} 中間帯で各半径の RJ 側 \(B_\nu\propto\nu^2 T\)。よって

\[F_\nu\propto\int\nu^2 T\cdot T^{-11/3}dT=\nu^2\int T^{-8/3}dT.\]

\(x=h\nu/kT\) で
\(T=h\nu/kx\)、\(dT\propto-\nu\,x^{-2}dx\)、\(T^{-8/3}\propto\nu^{-8/3}x^{8/3}\)。よって
\(\int T^{-8/3}dT\propto\nu^{-8/3}\cdot\nu\int x^{8/3-2}dx=\nu^{-5/3}\int x^{2/3}dx\)。中間帯では積分は
\(\nu\) に依らない定数（端の寄与が効かない範囲）。したがって
\(F_\nu\propto\nu^2\cdot\nu^{-5/3}=\nu^{1/3}\)。

\textbf{(3)} 一般に
\(T\propto R^{-\alpha}\Rightarrow R\propto T^{-1/\alpha}\)、\(2\pi R\,dR\propto T^{-2/\alpha}dT\)。\(F_\nu\propto\nu^2\int T\cdot T^{-2/\alpha}dT=\nu^2\int T^{1-2/\alpha}dT\)。\(x=h\nu/kT\)
変換で
\(F_\nu\propto\nu^2\cdot\nu^{(1-2/\alpha)+1}\cdot(\text{const})\cdot\nu^{-1}\)\ldots{}
整理すると指数は \((3\alpha-2)/\alpha\)。\(\alpha=3/4\) で
\((9/4-2)/(3/4)=(1/4)/(3/4)=1/3\)。

\textbf{答え}：\(2\pi R\,dR\propto T^{-11/3}dT\)、中間帯で
\(F_\nu\propto\nu^{1/3}\)。一般に \(T\propto R^{-\alpha}\) で
\(F_\nu\propto\nu^{(3\alpha-2)/\alpha}\)、\(\alpha=3/4\) で \(1/3\)。

\subsection*{問題 18-2　コンプトン y パラメータ}\label{ex-18-2}
\addcontentsline{toc}{subsection}{問題 18-2　コンプトン y パラメータ}

{[}★ 難易度：☆☆ {]}

§18.3
のコンプトン化を読み解く。種光子（低温黒体）が高温電子コロナ（\(T_e\)）を通過する。1回散乱あたりのエネルギー利得
\(\dfrac{\Delta\nu}{\nu}\simeq\dfrac{4k_BT_e}{m_ec^2}\)、\(y\equiv\dfrac{4k_BT_e}{m_ec^2}\max(\tau_e,\tau_e^2)\)。\(m_ec^2=511\)
keV。

\begin{enumerate}
\def\labelenumi{\arabic{enumi}.}
\tightlist
\item
  \(T_e=10^9\) K（\(k_BT_e\simeq86\) keV）で1回あたりのエネルギー利得
  \(\Delta\nu/\nu\) を求めよ。
\item
  \(y\ll1\)、\(y\gtrsim1\)、中間域で観測スペクトルがどう変わるか（弱い
  Wien テイル／Wien 分布化／べき則）述べよ。
\item
  この \(y\) パラメータが第16章 CMB の \(y\) 歪み（Sunyaev-Zel'dovich
  効果）と同じ物理であること、スケールの違いだけであることを述べよ。
\end{enumerate}

\textbf{関連}：\hyperref[sec-comptonization]{§18.3 コンプトン化}／CMB の
\(y\)
歪みは\hyperref[sec-cmb-distortions]{§16.5}（\hyperref[ex-16-4]{演習
16-4}）。

\section{模範解答（問題
18-2）}\label{ux6a21ux7bc4ux89e3ux7b54ux554fux984c-18-2}

\textbf{(1)}
\(\dfrac{\Delta\nu}{\nu}=\dfrac{4k_BT_e}{m_ec^2}=\dfrac{4\times86\ \text{keV}}{511\ \text{keV}}=\dfrac{344}{511}\simeq0.67\)。1回の散乱で振動数が約
67\% 増える。（\(k_BT_e=8.617\times10^{-5}\times10^9=8.6\times10^4\) eV
\(=86\) keV。）

\textbf{(2)}
\(y\ll1\)：光子はわずかにエネルギーを得るだけ＝元の黒体形に弱い Wien
側テイルが付く。\(y\gtrsim1\)：光子がコロナで完全に再分配され、Wien
分布に近い高温化スペクトル。中間域：べき則 \(F_\nu\propto\nu^{-\alpha}\)
が現れ、高エネルギー側は Wien カットオフで切れる。

\textbf{(3)}
どちらも高温電子と光子の逆コンプトン散乱で、\(y=\dfrac{4k_BT_e}{m_ec^2}\max(\tau,\tau^2)\)
という同じパラメータでスペクトル変形を測る。X
線連星のコロナ（銀河近傍スケール）も、銀河団ガス（宇宙論スケール）も、同じ物理。違うのはスケール（\(T_e\)、\(\tau\)
の値と対象の大きさ）だけ。

\textbf{答え}：\(\Delta\nu/\nu\simeq0.67\)（\(T_e=10^9\) K）。\(y\ll1\)
弱テイル、\(y\gtrsim1\) Wien 化、中間でべき則。CMB の \(y\) 歪み（SZ
効果）と同物理。

\subsection*{問題 18-3　局所黒体仮定とその限界}\label{ex-18-3}
\addcontentsline{toc}{subsection}{問題 18-3　局所黒体仮定とその限界}

{[}★ 難易度：☆☆ {]}

§18.1
は標準降着円盤が「各半径で局所的に黒体」と仮定する。この仮定が成り立つ条件と破れ方を整理する。

\begin{enumerate}
\def\labelenumi{\arabic{enumi}.}
\tightlist
\item
  各半径が温度 \(T(R)\)
  の黒体として放射するには、第3章・第5章のどんな条件（熱化時間、光学的厚さ、単一温度）が局所的に満たされる必要があるか述べよ。
\item
  実際の円盤で局所黒体から外れる要因を2つ挙げよ（高温コロナ、振動数依存不透明度、電子散乱卓越など）。
\item
  電子散乱が吸収より卓越する高温体では、観測される色温度
  \(T_\mathrm{col}\) が実効温度 \(T(R)\) より高くなり、補正係数
  \(f_\mathrm{col}\simeq1.5\)〜\(2\)
  が経験的に要る。なぜ散乱卓越が色温度を上げるのか、\hyperref[ex-15-3]{演習
  15-3} のソース関数と結びつけて一言で述べよ。
\end{enumerate}

\textbf{関連}：\hyperref[sec-disk-local-blackbody]{§18.1
局所黒体近似}／黒体化条件は\hyperref[sec-universality-real-objects]{§3.6}（\hyperref[ex-3-4]{演習
3-4}）、LTE
は\hyperref[sec-radiative-transfer-lte]{§5.5}、散乱ソース関数は\hyperref[ex-15-3]{演習
15-3}。

\section{模範解答（問題
18-3）}\label{ux6a21ux7bc4ux89e3ux7b54ux554fux984c-18-3}

\textbf{(1)} ①
熱化時間が円盤の力学的時間より短い（局所的に熱平衡に達する）、②
各半径で鉛直方向に光学的に厚い（\(\tau\gg1\)、吸収・再放出が十分）、③
その半径内では単一温度 \(T(R)\)
で記述できる。これらが満たされれば各半径が \(B_\nu(T(R))\)
で放射する（第3章 §3.6、第5章 §5.6）。

\textbf{(2)}
例：高温・希薄なコロナ層（局所熱平衡を破る）、不透明度の振動数依存（修正黒体的なずれ）、電子散乱が吸収より卓越（純黒体から外れる）。

\textbf{(3)} 電子散乱卓越ではソース関数
\(S_\nu=(1-\tilde\omega)B_\nu+\tilde\omega\langle J_\nu\rangle\)（\hyperref[ex-15-3]{演習
15-3}）が局所プランクから外れ、光子は吸収で熱化されにくく深部の高温を反映したまま散乱で抜ける。結果、見かけの色温度が実効温度より高くなり、\(f_\mathrm{col}\simeq1.5\)〜\(2\)
の補正が要る。

\textbf{答え}：局所熱平衡・光学的厚・単一温度が必要。破れ＝コロナ／振動数依存不透明度／電子散乱卓越。散乱卓越は熱化を妨げ色温度を上げる（\(f_\mathrm{col}\sim1.5\)〜\(2\)）。

\subsection*{問題 18-4　中性子星表面の重力赤方偏移}\label{ex-18-4}
\addcontentsline{toc}{subsection}{問題 18-4　中性子星表面の重力赤方偏移}

{[}★ 難易度：☆☆ {]}

§18.4
は中性子星の観測温度が重力赤方偏移を受けると述べた。\(z_g=(1-2GM/Rc^2)^{-1/2}-1\)、観測温度
\(T_\infty=T/(1+z_g)\)。\(M=1.4M_\odot=2.78\times10^{30}\) kg、\(R=10\)
km、\(G=6.674\times10^{-11}\)、\(c=3\times10^8\) m/s。

\begin{enumerate}
\def\labelenumi{\arabic{enumi}.}
\tightlist
\item
  コンパクトネス \(2GM/Rc^2\) を計算せよ。
\item
  重力赤方偏移 \(z_g\) を求め、\(z_g\simeq0.3\) を確かめよ。
\item
  表面温度 \(T=10^7\) K の中性子星の観測温度 \(T_\infty\)
  を求めよ。重力赤方偏移と Stefan-Boltzmann
  を連立すると質量と半径の組合せが制約できる（NICER
  の戦略）ことに一言触れよ。
\end{enumerate}

\textbf{関連}：\hyperref[sec-neutron-star-surface]{§18.4
中性子星表面}／Fermi
縮退は\hyperref[sec-photon-statistics-bose-fermi]{§8.2}、相対論的修正は\hyperref[sec-black-hole-effects]{§18.5}。

\section{模範解答（問題
18-4）}\label{ux6a21ux7bc4ux89e3ux7b54ux554fux984c-18-4}

\textbf{(1)}
\(\dfrac{2GM}{Rc^2}=\dfrac{2\times6.674\times10^{-11}\times2.78\times10^{30}}{1.0\times10^4\times(3\times10^8)^2}=\dfrac{3.71\times10^{20}}{9.0\times10^{20}}=0.412\)。

\textbf{(2)}
\(z_g=(1-0.412)^{-1/2}-1=(0.588)^{-1/2}-1=1.304-1=0.304\simeq0.30\)。

\textbf{(3)} \(T_\infty=T/(1+z_g)=10^7/1.30=7.7\times10^6\)
K。観測される温度は表面温度より約23\%低い。重力赤方偏移 \(z_g\) は
\(M/R\) に依存し、Stefan-Boltzmann で得る放射半径は \(M,R\)
の別の組合せに依存するので、両者を連立すると \(M\) と \(R\)
を独立に制約できる。NICER
はこれにパルス時間プロファイルを加えて核物質の状態方程式を制約している。

\textbf{答え}：\(2GM/Rc^2=0.41\)、\(z_g\simeq0.30\)、\(T_\infty=7.7\times10^6\)
K。\(z_g\) と Stefan-Boltzmann の連立で \(M,R\) を制約（NICER）。

\section*{さらに学ぶための参考文献}\label{sec-accretion-high-energy-further}
\addcontentsline{toc}{section}{さらに学ぶための参考文献}

\markright{さらに学ぶための参考文献}

\begin{itemize}
\tightlist
\item
  Frank, King \& Raine, \emph{Accretion Power in Astrophysics}
  (Cambridge, 3rd ed., 2002) --- 降着円盤と高エネルギー天体の標準教科書
\item
  Shakura \& Sunyaev, A\&A 24 (1973) 337 --- 標準円盤モデルの原論文
\item
  Rybicki \& Lightman, \emph{Radiative Processes in Astrophysics}
  (Wiley, 1979) --- コンプトン化、シンクロトロン放射
\item
  Kompaneets, JETP 4 (1957) 730 --- コンプトン化方程式の原論文
\item
  Bambi, \emph{Black Holes: A Laboratory for Testing Strong Gravity}
  (Springer, 2017) --- BH近傍の相対論的放射過程
\end{itemize}

\chapter{第19章　自由電子が作る連続光 ―
自由-自由・束縛-自由・再結合}\label{ux7b2c19ux7ae0-ux81eaux7531ux96fbux5b50ux304cux4f5cux308bux9023ux7d9aux5149-ux81eaux7531-ux81eaux7531ux675fux7e1b-ux81eaux7531ux518dux7d50ux5408}

\textbf{この章の主題}：第15章で「振動数依存不透明度」を、第17章で電離ガスの自由-自由放射を、いずれも吸収係数
\(\alpha_\nu\) と放射率 \(j_\nu\) を \textbf{与えられたもの}
として使った。本章はその出どころに降りる。希薄プラズマや恒星大気の
\textbf{連続スペクトル} を作るのは、自由電子が絡む三つの素過程 ―
\textbf{自由-自由}（制動放射）、\textbf{束縛-自由}（光電離）、その逆過程である
\textbf{自由-束縛}（再結合） ―
である。これらを放出と吸収の両面で、Kirchhoff の法則と Milne
関係、そして Saha
の電離平衡で結び、観測される連続吸収端・電波ターンオーバー・X
線の指数カットオフ・再結合連続光を逆引きする。最後に、再結合が連続光だけでなく
\textbf{線スペクトル}（再結合線）も生むことを示し、第VII部への橋とする。

\section*{この章の中心地図}\label{sec-continuum-micro-map}
\addcontentsline{toc}{section}{この章の中心地図}

\markright{この章の中心地図}

\textbf{中心地図：プランク関数 ― 本部の主題は「ずれ」}

\[
B_\nu(T) = \frac{2 h \nu^3}{c^2} \cdot \frac{1}{e^{h\nu/kT} - 1} \quad \overset{\text{現実}}{\longrightarrow} \quad \text{?}
\]

第VI部（現実の宇宙へ戻る）は、観測される現実の天体スペクトルが、この理想的プランク関数からどのように、そしてなぜ
\textbf{ずれる}かを扱う。ずれの主な起源：

\begin{itemize}
\tightlist
\item
  \textbf{振動数依存不透明度} → 修正黒体（第15章）
\item
  \textbf{多温度成分の重ね合わせ} → 多温度黒体スペクトル（第17・18章）
\item
  \textbf{散乱効果}（コンプトン化など）→ 高エネルギー側の変形（第18章）
\item
  \textbf{非熱的放射の混入} → シンクロトロン・ベキ則の重なり
\end{itemize}

本部は、観測スペクトルを「理想的黒体＋ずれ」として読み解き、ずれの定量化から物理状態を引き出す視点を養う。

\begin{tcolorbox}[enhanced jigsaw, arc=.35mm, bottomrule=.15mm, bottomtitle=1mm, breakable, colback=white, colbacktitle=quarto-callout-note-color!10!white, colframe=quarto-callout-note-color-frame, coltitle=black, left=2mm, leftrule=.75mm, opacityback=0, opacitybacktitle=0.6, rightrule=.15mm, title=\textcolor{quarto-callout-note-color}{\faInfo}\hspace{0.5em}{ノート}, titlerule=0mm, toprule=.15mm, toptitle=1mm]

\textbf{方針}：本章の中心地図は、連続の吸収・放射係数を素過程に分解した形

\[
\alpha_\nu^{\mathrm{cont}} = \underbrace{n_e n_i\,\bar\sigma_{\mathrm{ff}}(\nu,T)}_{\text{自由-自由}} \;+\; \underbrace{\textstyle\sum_i n_i\,\sigma_{\mathrm{bf}}^{(i)}(\nu)}_{\text{束縛-自由}}, \qquad S_\nu \overset{\mathrm{LTE}}{=} B_\nu(T)
\]

である。占有数 \(n_e, n_i, n_X\) は
\textbf{Saha--Boltzmann}（§19.5）が決め、放出 \(j_\nu\) と吸収
\(\alpha_\nu\) は \textbf{Kirchhoff（LTE）／詳細つり合い（Milne）}
で結ばれる。「中心地図プランク関数のどの因子が、どの素過程から来るか」を素過程の側から裏打ちするのが本章である。

\end{tcolorbox}

\section*{この章で答える問い}\label{sec-continuum-micro-questions}
\addcontentsline{toc}{section}{この章で答える問い}

\markright{この章で答える問い}

\begin{itemize}
\tightlist
\item
  恒星スペクトルの紫外側の段差（Balmer jump
  など）は、何の素過程が作るのか。なぜ「端」が立つのか。
\item
  HII
  領域の電波連続光が、低振動数で黒体に漸近し高振動数で平坦になるのはなぜか（第17章
  §17.6 の微視的根拠）。
\item
  銀河団の X 線連続光が \(e^{-h\nu/k_BT}\) で切れるのはなぜか。
\item
  太陽光球の可視連続光を、なぜ豊富な中性水素そのものではなく
  \textbf{微量の H⁻} が担うのか。
\item
  同じ再結合が、なぜ連続光（自由-束縛）と輝線（束縛-束縛カスケード）の
  \textbf{両方} を生むのか。
\end{itemize}

\section*{到達目標}\label{sec-continuum-micro-goals}
\addcontentsline{toc}{section}{到達目標}

\markright{到達目標}

この章を読み終えると、読者は次のことができるようになる：

\begin{itemize}
\tightlist
\item
  自由-自由・束縛-自由・再結合（自由-束縛）の三過程を、エネルギー準位図の上で区別して説明できる。
\item
  光電離断面積 \(\sigma_{\mathrm{bf}}(\nu)\)
  のしきい値と連続吸収端の関係、\(\sim\nu^{-3}\) 減衰を述べられる。
\item
  Milne 関係（詳細つり合い）で \(\sigma_{\mathrm{bf}}\)
  と再結合断面積を結び、再結合係数 \(\alpha_{\mathrm{rec}}(T)\)
  の意味と温度依存を説明できる。
\item
  Saha 式を使って電離度から連続の \(\alpha_\nu, j_\nu\)
  を自己無撞着に決める道筋を述べられる。
\item
  連続吸収端・電波ターンオーバー・X
  線カットオフ・再結合連続光を、観測から逆引きできる。
\end{itemize}

\begin{center}\rule{0.5\linewidth}{0.5pt}\end{center}

\section{19.1 連続を作る三素過程 ―
全体像}\label{sec-continuum-three-processes}

{[}本文目安：B2{]}

第VII部で扱う線スペクトルは、原子の \textbf{束縛準位の間}
の遷移（\textbf{束縛-束縛}、bound-bound）である。離散的な準位差
\(E_u-E_l\) に対応する振動数 \(\nu_0\)
にだけ光が出入りするので、スペクトルは \textbf{線} になる。

これに対し、電子の片側または両側が \textbf{自由状態}（連続エネルギー
\(E>0\)）にある遷移は、連続したエネルギー範囲を取れるので
\textbf{連続スペクトル} を作る。自由電子が関与する遷移は三種類ある：

\begin{itemize}
\tightlist
\item
  \textbf{束縛-自由}（bound-free,
  bf）：束縛電子が光子を吸って電離する＝\textbf{光電離}（photoionization）。しきい値
  \(h\nu \ge \chi\) で起こる。逆過程（§19.4）は
  \textbf{自由-束縛}（free-bound）＝\textbf{再結合}で、光子を放出する。
\item
  \textbf{自由-自由}（free-free, ff）：自由電子がイオンの Coulomb
  場で曲げられる際に光子を吸収・放出する＝\textbf{制動放射}（bremsstrahlung）。束縛準位を経由しないので、任意の振動数で起こる。
\item
  （\textbf{束縛-束縛}：線。第VII部。）
\end{itemize}

\begin{figure}[H]

\centering{

\includegraphics[width=0.88\linewidth,height=\textheight,keepaspectratio]{index_files/mediabag/part6/../images/continuum-processes.pdf}

}

\caption{\label{fig-continuum-processes}連続を作る三素過程の準位図。横線が束縛準位、上部の網掛け帯が連続状態（\(E>0\)）。束縛-束縛（線、第VII部）と対比して、束縛-自由（光電離↔再結合）と自由-自由（制動放射）が連続帯を使うことを示す。第11章
§11.2／第21章 §21.2 の二準位＋三過程図の姉妹図。}

\end{figure}%

\begin{tcolorbox}[enhanced jigsaw, arc=.35mm, bottomrule=.15mm, bottomtitle=1mm, breakable, colback=white, colbacktitle=quarto-callout-note-color!10!white, colframe=quarto-callout-note-color-frame, coltitle=black, left=2mm, leftrule=.75mm, opacityback=0, opacitybacktitle=0.6, rightrule=.15mm, title=\textcolor{quarto-callout-note-color}{\faInfo}\hspace{0.5em}{ノート}, titlerule=0mm, toprule=.15mm, toptitle=1mm]

\textbf{用語}：本章の三過程は、電子が \textbf{Maxwell 分布}（温度
\(T_e\)）にある限りすべて \textbf{熱的放射}（第1章
§1.9）である。ここで二つの条件を明確に区別しておく：

\begin{itemize}
\tightlist
\item
  \textbf{熱的}＝粒子（電子）のエネルギー分布が Maxwell-Boltzmann
  であること。これは放射を出す側の \textbf{粒子分布} の性質。
\item
  \textbf{黒体的}＝放射場が局所の温度の黒体 \(B_\nu(T)\)
  に達していること。これには \textbf{光学的に厚く、放射が熱化している}
  ことが必要。
\end{itemize}

すなわち「熱的」であっても、光学的に薄ければ黒体にはならない。この区別は第17章
§17.6
で自由-自由放射について見た通りで、本章の三過程すべてに共通する（§19.6）。

\end{tcolorbox}

\section{19.2 自由-自由放射と吸収 ―
熱的制動放射}\label{sec-continuum-free-free}

{[}本文目安：B3{]}

自由電子がイオンの Coulomb
場で加速度を受けると、古典電磁気（第V部）の通り電磁波を放つ。これが
\textbf{自由-自由放射}（制動放射）である。

ここで放出の量を表す係数を二つ区別しておく。\textbf{放出係数}（emission
coefficient）\(j_\nu\)
は単位体積・単位振動数・\textbf{単位立体角あたり}に放出されるエネルギー（次元は
erg s⁻¹ cm⁻³ Hz⁻¹
sr⁻¹）であり、放射輸送方程式にソース項として現れる。これに対し、\textbf{体積放射率}
\(\varepsilon_\nu\) は全立体角に放出されるエネルギー（次元は erg s⁻¹
cm⁻³ Hz⁻¹）で、等方放射では両者は立体角積分
\(\varepsilon_\nu = 4\pi\, j_\nu\) で結ばれる。

電子が Maxwell 分布にあるとき、自由-自由の体積放射率は

\begin{equation}\protect\phantomsection\label{eq-ff-emissivity}{
\varepsilon_\nu^{\mathrm{ff}} \;\propto\; n_e\,n_i\,Z^2\,T_e^{-1/2}\,e^{-h\nu/k_BT_e}\,\bar g_{\mathrm{ff}}(\nu,T_e)
}\end{equation}

の形をとる（自由-自由放射は等方なので、放出係数は
\(j_\nu^{\mathrm{ff}} = \varepsilon_\nu^{\mathrm{ff}}/4\pi\)）。\(Z\)
はイオンの電荷、\(\bar g_{\mathrm{ff}}\) は \textbf{Gaunt
因子}（量子補正、\(\nu,T_e\) にゆるやかに依存し \(O(1)\)）。指数因子
\(e^{-h\nu/k_BT_e}\) は、振動数 \(\nu\) の光子を出すには電子が少なくとも
\(h\nu\) の運動エネルギーを持たねばならず、その割合が Maxwell
分布で抑えられることを表す。

対応する \textbf{吸収係数} は、第17章 §17.6 で見たように

\begin{equation}\protect\phantomsection\label{eq-ff-absorption-micro}{
\alpha_\nu^{\mathrm{ff}} \;\propto\; n_e\,n_i\,Z^2\,T_e^{-3/2}\,\nu^{-2}\,\bar g_{\mathrm{ff}}\,\bigl(1-e^{-h\nu/k_BT_e}\bigr)
}\end{equation}

括弧の \((1-e^{-h\nu/k_BT_e})\) は \textbf{誘導放出補正}（第11章・第21章
§21.2
と同じ因子）である。ソース関数は放出係数（単位立体角あたり）と吸収係数の比
\(S_\nu = j_\nu/\alpha_\nu\) で定義されるので、等方放射で
\(j_\nu^{\mathrm{ff}} = \varepsilon_\nu^{\mathrm{ff}}/4\pi\)
を用いる。式~\ref{eq-ff-emissivity} と 式~\ref{eq-ff-absorption-micro}
の比をとると、\(n_e,n_i\) も \(\bar g_{\mathrm{ff}}\) も消えて

\begin{equation}\protect\phantomsection\label{eq-ff-kirchhoff}{
S_\nu^{\mathrm{ff}} = \frac{j_\nu^{\mathrm{ff}}}{\alpha_\nu^{\mathrm{ff}}} = \frac{1}{4\pi}\frac{\varepsilon_\nu^{\mathrm{ff}}}{\alpha_\nu^{\mathrm{ff}}} = \frac{2h\nu^3}{c^2}\frac{1}{e^{h\nu/k_BT_e}-1} = B_\nu(T_e)
}\end{equation}

となり、\textbf{Kirchhoff の法則} \(S_\nu=B_\nu(T_e)\)
が素過程のレベルで確かめられる（第3章・第5章）。立体角因子 \(4\pi\) は
\(j_\nu\) と \(\varepsilon_\nu\)
の換算に現れるが、係数の振動数・温度依存はそのまま \(B_\nu(T_e)\)
を再現する。すなわち自由-自由放射のソース関数は、つねに電子温度の黒体である。

観測スペクトルの形を決めるのは、ここから先の \textbf{光学的厚さ}
\(\tau_\nu=\int\alpha_\nu^{\mathrm{ff}}\,ds\)
である。式~\ref{eq-ff-absorption-micro} の \(\nu^{-2}\) により、

\begin{itemize}
\tightlist
\item
  \textbf{低振動数（電波）}：\(\alpha_\nu^{\mathrm{ff}}\)
  が大きく光学的に厚い → \(I_\nu\to B_\nu(T_e)\)。HII
  領域の電波連続光が黒体に漸近する（第17章 §17.6 のターンオーバー）。
\item
  \textbf{高振動数・高温・希薄（X 線）}：光学的に薄く、観測される
  \(F_\nu\propto e^{-h\nu/k_BT_e}\) の
  \textbf{指数カットオフ}。銀河団の高温ガス（\(T_e\sim10^7\)〜\(10^8\)
  K）が典型（第18章）。
\end{itemize}

\begin{figure}[H]

\centering{

\includegraphics[width=0.92\linewidth,height=\textheight,keepaspectratio]{index_files/mediabag/part6/../images/continuum-ff-spectrum.pdf}

}

\caption{\label{fig-continuum-ff-spectrum}自由-自由放射のスペクトル。低振動数では光学的に厚く
\(F_\nu\propto\nu^2\)（電子温度 \(T_e\) の黒体の Rayleigh--Jeans
裾）、ターンオーバー \(\nu_t\)
を越えると光学的に薄くほぼ平坦（\(\propto\nu^{-0.1}\)）、\(h\nu\sim k_BT_e\)
で指数カットオフ。電波 HII 領域と銀河団 X
線は同じ熱的制動放射の両端である。}

\end{figure}%

\begin{tcolorbox}[enhanced jigsaw, arc=.35mm, bottomrule=.15mm, breakable, colback=white, colframe=quarto-callout-tip-color-frame, left=2mm, leftrule=.75mm, opacityback=0, rightrule=.15mm, toprule=.15mm]
\begin{minipage}[t]{5.5mm}
\textcolor{quarto-callout-tip-color}{\faLightbulb}
\end{minipage}%
\begin{minipage}[t]{\textwidth - 5.5mm}

\textbf{問い}：自由-自由放射は全帯域で「熱的」なのに、なぜ全帯域で単一の黒体スペクトルには
\textbf{ならない} のか？

\textbf{短答}：黒体になるのは光学的に厚い帯（低 \(\nu\)）だけ。高
\(\nu\) では光学的に薄く \(I_\nu\approx\tau_\nu B_\nu \ll B_\nu\)
となって黒体から下にずれ、\(e^{-h\nu/k_BT_e}\) で切れる。「熱的（Maxwell
分布）」と「黒体（光学的に厚い熱平衡放射場）」は別条件である（第17章
§17.6）。

\end{minipage}%
\end{tcolorbox}

\section{19.3 束縛-自由 ― 光電離断面積と連続吸収端（H⁻
を含む）}\label{sec-continuum-bound-free}

{[}本文目安：B3{]}

束縛電子が \(h\nu\ge\chi\) の光子を吸って自由になる過程が
\textbf{光電離}（束縛-自由吸収）である。下準位 \(n\) の電離エネルギーを
\(\chi_n\) とすると、吸収はしきい振動数

\begin{equation}\protect\phantomsection\label{eq-bf-threshold}{
\nu_{\mathrm{th}} = \chi_n/h
}\end{equation}

で \textbf{急に立ち上がり}、それ以上の振動数では水素様原子で

\begin{equation}\protect\phantomsection\label{eq-bf-cross-section}{
\sigma_{\mathrm{bf}}(\nu) \;\propto\; \frac{Z^4}{n^5}\,\nu^{-3}\,g_{\mathrm{bf}}(\nu,n) \qquad (\nu\ge\nu_{\mathrm{th}})
}\end{equation}

と \(\nu^{-3}\) で減衰する（Kramers の断面積、\(g_{\mathrm{bf}}\)
は束縛-自由 Gaunt 因子）。しきい値での不連続な立ち上がりと \(\nu^{-3}\)
の減衰が、観測される \textbf{連続吸収端}（continuum
edge）を作る。水素では

\begin{itemize}
\tightlist
\item
  \textbf{Lyman 端}（\(n=1\)、\(\lambda=912\) Å、紫外）
\item
  \textbf{Balmer 端}（\(n=2\)、\(\lambda=3646\) Å、可視と紫外の境）
\item
  \textbf{Paschen 端}（\(n=3\)、\(\lambda=8204\) Å、近赤外）
\end{itemize}

と、下準位が浅くなるほど端の波長が長くなる。これが第15章 §15.3
で「振動数依存不透明度」として与えた連続吸収端の \textbf{微視的な起源}
である。各準位の占有数
\(n_i\)（Boltzmann／Saha、§19.5）が、どの端がどれだけ立つかを決める。

\subsection{H⁻ ―
低温星の可視連続不透明度の主役}\label{h-ux4f4eux6e29ux661fux306eux53efux8996ux9023ux7d9aux4e0dux900fux660eux5ea6ux306eux4e3bux5f79}

水素の束縛-自由端はいずれも紫外〜近赤外にあり、\textbf{可視のど真ん中には水素中性原子の端がない}。それなら太陽（\(T_{\mathrm{eff}}\simeq5800\)
K）の光球は可視で透明なはずだが、実際は可視全域で不透明である。担い手は中性水素そのものではなく、\textbf{負水素イオン
H⁻} ― 水素原子に電子がもう一つゆるく束縛したイオン ― である。

H⁻ の束縛エネルギーはわずか \(0.754\) eV
しかない。したがって束縛-自由しきい波長は

\begin{equation}\protect\phantomsection\label{eq-hminus-threshold}{
\lambda_{\mathrm{th}}^{\mathrm{H}^-} = \frac{hc}{0.754\ \mathrm{eV}} \simeq 1.64\ \mu\mathrm{m}
}\end{equation}

で、\(1.64\ \mu\)m より短い \textbf{可視〜近赤外の全域}
で光電離吸収が効く（吸収のピークは
\(\sim0.85\ \mu\)m）。さらに長波長側では H⁻ の \textbf{自由-自由}
吸収（§19.2 と同型）が連続して効く。

H⁻ が低温星で卓越するのには、Saha 平衡（§19.5）が効いている：

\begin{itemize}
\tightlist
\item
  \textbf{中性水素が豊富} であること（H⁻
  の素材）。高温星では水素が電離して中性 H が枯れ、H⁻ も消える。
\item
  \textbf{自由電子の供給}
  があること。低温星では水素はほとんど電離しないので、自由電子は
  \textbf{電離ポテンシャルの低い金属}（Na, K, Ca, Mg, Fe
  など）の電離が供給する。したがって H⁻ 不透明度は電子密度
  \(n_e\)（≒金属量と Saha）に比例して増える。
\end{itemize}

高温星（O・B 型）では水素・ヘリウムの束縛-自由／自由-自由と
\textbf{電子散乱}（第15章
§15.5）に主役が交代する。「中性水素が主成分なのに、可視の連続不透明度は微量イオン
H⁻ が担う」という逆転は、連続スペクトルの逆引きが Saha
平衡と分かちがたく結びついていることを示す好例である。

\begin{tcolorbox}[enhanced jigsaw, arc=.35mm, bottomrule=.15mm, bottomtitle=1mm, breakable, colback=white, colbacktitle=quarto-callout-note-color!10!white, colframe=quarto-callout-note-color-frame, coltitle=black, left=2mm, leftrule=.75mm, opacityback=0, opacitybacktitle=0.6, rightrule=.15mm, title=\textcolor{quarto-callout-note-color}{\faInfo}\hspace{0.5em}{ノート}, titlerule=0mm, toprule=.15mm, toptitle=1mm]

\textbf{観測との対応}：太陽スペクトルの連続光のなだらかな形（可視で最大、紫外・赤外へ向けて落ちる）は、H⁻
の吸収係数の波長依存（\(\sim0.85\ \mu\)m
にピーク）でよく説明される。恒星大気モデル（第15章 §15.3、Mihalas
1978）では H⁻ が G・K
型星の連続不透明度の標準的な第一成分として組み込まれている。

\end{tcolorbox}

\section{19.4 再結合 ― 束縛-自由の逆過程と Milne
関係}\label{sec-continuum-recombination}

{[}本文目安：B3{]}

光電離（束縛-自由）の \textbf{逆過程} が
\textbf{再結合}（recombination）である：自由電子がイオンに捕まって準位
\(n\) に落ち、余ったエネルギーを光子として放つ。

\begin{equation}\protect\phantomsection\label{eq-recombination-reaction}{
e^- + X^+ \;\longrightarrow\; X(n) + h\nu
}\end{equation}

捕まる前の電子の運動エネルギーを \(\tfrac12 m_e v^2\)
とすると、放出光子のエネルギーは

\begin{equation}\protect\phantomsection\label{eq-fb-photon-energy}{
h\nu = \tfrac12 m_e v^2 + \chi_n
}\end{equation}

である。電子は連続的な速度分布（Maxwell）を持つので、放出される光子も
\textbf{連続スペクトル} になる ― これが
\textbf{自由-束縛放射}（free-bound
emission）＝\textbf{再結合連続光}（recombination
continuum）である。ただし \(h\nu\) は最小でも
\(\chi_n\)（\(v\to0\)）なので、再結合連続光は各準位のしきい \(\chi_n\)
で \textbf{下端（放出端）}
を持ち、そこから高振動数へ裾を引く。吸収側の連続吸収端（§19.3）と同じ位置に、今度は
\textbf{放出の端} が立つ。

\begin{figure}[H]

\centering{

\includegraphics[width=0.94\linewidth,height=\textheight,keepaspectratio]{index_files/mediabag/part6/../images/continuum-edges-recomb.pdf}

}

\caption{\label{fig-continuum-edges-recomb}同じしきい値（Lyman 912
Å・Balmer 3646 Å・Paschen 8204 Å）が、恒星大気では
\textbf{束縛-自由の吸収端}（上：短波長へ向け急に立ち上がり
\(\propto\nu^{-3}\) で減衰、可視は H⁻ が担う）として、星雲では
\textbf{自由-束縛の放出端＝再結合連続光}（下：各しきいで放出が跳ね上がり短波長へ
Wien 減衰、Balmer
ジャンプが顕著）として現れる。背景光の有無で吸収にも放出にも見える（§19.6）。}

\end{figure}%

\subsection{Milne 関係 ―
詳細つり合いが結ぶ吸収と放出}\label{milne-ux95a2ux4fc2-ux8a73ux7d30ux3064ux308aux5408ux3044ux304cux7d50ux3076ux5438ux53ceux3068ux653eux51fa}

熱平衡では、光電離（吸収）と再結合（放出）が振動数ごとに釣り合う（\textbf{詳細つり合い}、第6章）。この釣り合いから、再結合断面積
\(\sigma_{\mathrm{rec}}(v)\) と光電離断面積
\(\sigma_{\mathrm{bf}}(\nu)\) が直接結ばれる ― \textbf{Milne 関係}：

\begin{equation}\protect\phantomsection\label{eq-milne-relation}{
\sigma_{\mathrm{rec}}(v) = \frac{g_n}{2\,g_{+}}\,\frac{(h\nu)^2}{m_e^2 c^2 v^2}\,\sigma_{\mathrm{bf}}(\nu)
}\end{equation}

\(g_n, g_{+}\) は再結合後の準位とイオンの統計重率、因子 \(2\)
は自由電子のスピン状態数。これは第11章で Einstein 係数 \(A,B\)
を詳細つり合いで結んだのと \textbf{同じ論法}
の連続版である：放出の量（\(A\) や
\(\sigma_{\mathrm{rec}}\)）を、吸収の量（\(B\) や
\(\sigma_{\mathrm{bf}}\)）から熱平衡の要請だけで決める。

準位 \(n\) への再結合率を Maxwell 分布で平均したものが
\textbf{再結合係数}

\begin{equation}\protect\phantomsection\label{eq-recombination-coefficient}{
\alpha_n(T) = \langle \sigma_{\mathrm{rec}} v\rangle = \int_0^\infty \sigma_{\mathrm{rec}}(v)\,v\,f_{\mathrm{Maxwell}}(v,T)\,dv
}\end{equation}

で、単位は \(\mathrm{cm^3\,s^{-1}}\)。Milne 関係から
\(\sigma_{\mathrm{rec}}\propto v^{-2}\sigma_{\mathrm{bf}}\)
であり、Maxwell 平均をとると \(\alpha_n(T)\propto T^{-1/2}\)
程度（係数の精密値は Gaunt 因子と準位和に依存し、全準位和の場合
\(\alpha_A\propto T^{-0.7\sim-0.8}\)）の \textbf{温度の減少関数}
になる。高温ほど電子が速く通り過ぎて捕まりにくい、という直観に合う。

\begin{tcolorbox}[enhanced jigsaw, arc=.35mm, bottomrule=.15mm, bottomtitle=1mm, breakable, colback=white, colbacktitle=quarto-callout-note-color!10!white, colframe=quarto-callout-note-color-frame, coltitle=black, left=2mm, leftrule=.75mm, opacityback=0, opacitybacktitle=0.6, rightrule=.15mm, title=\textcolor{quarto-callout-note-color}{\faInfo}\hspace{0.5em}{ノート}, titlerule=0mm, toprule=.15mm, toptitle=1mm]

\textbf{用語}：全準位（\(n\ge1\)）への再結合を足したものを \textbf{case
A}、基底状態（\(n=1\)）への直接再結合で出る Lyman
連続光が周囲に再吸収されて実質効かない極限を \textbf{case B}
と呼ぶ（§19.7 で線スペクトルとして再登場）。星雲では多くの場合 case B
が適切である。

\end{tcolorbox}

\section{19.5 Saha 平衡と連続放射の自己無撞着}\label{sec-continuum-saha}

{[}本文目安：B3{]}

§19.2--19.4 の各係数は、いずれも
\(n_e\)（電子）、\(n_i\)（イオン）、\(n_X\)（中性原子の各準位）という
\textbf{占有数} に比例する。これらを決めるのが、電離反応
\(X \rightleftharpoons X^+ + e^-\) の熱平衡条件 ― \textbf{Saha の式}
である：

\begin{equation}\protect\phantomsection\label{eq-saha-maintext}{
\frac{n_{X^+}\,n_e}{n_X} = \left(\frac{2\pi m_e k_B T}{h^2}\right)^{3/2} 2\,\frac{g_{X^+}}{g_X}\,e^{-\chi/k_BT}
}\end{equation}

（導出と記号は付録 A.3.3 を参照。\(\chi\) は電離エネルギー、右辺の \(2\)
は自由電子のスピン。）Saha の式は \textbf{Boltzmann 分布}（第7章
§7.3、準位間の占有比）の電離版であり、両者を合わせた
\textbf{Saha--Boltzmann} が、

\[
\text{温度 }T \;\longrightarrow\; \text{電離度・各準位占有 }n_e,n_i,n_X \;\longrightarrow\; \alpha_\nu^{\mathrm{ff}},\ \sigma_{\mathrm{bf}}\text{ の重み},\ \varepsilon_\nu^{\mathrm{rec}}
\]

という連鎖で、連続の吸収・放射係数を \textbf{自己無撞着} に与える。LTE
が成り立つ限り、この連鎖の出力はつねに
\(S_\nu=B_\nu(T)\)（Kirchhoff、式~\ref{eq-ff-kirchhoff}）に帰着する ―
素過程が何であってもソース関数は局所温度の黒体になる、というのが LTE
連続放射の要点である。

\begin{tcolorbox}[enhanced jigsaw, arc=.35mm, bottomrule=.15mm, bottomtitle=1mm, breakable, colback=white, colbacktitle=quarto-callout-note-color!10!white, colframe=quarto-callout-note-color-frame, coltitle=black, left=2mm, leftrule=.75mm, opacityback=0, opacitybacktitle=0.6, rightrule=.15mm, title=\textcolor{quarto-callout-note-color}{\faInfo}\hspace{0.5em}{ノート}, titlerule=0mm, toprule=.15mm, toptitle=1mm]

\textbf{注意（観測）}：恒星大気では H⁻ 不透明度（§19.3）が \(n_e\)
に比例し、その \(n_e\) は金属の Saha
電離で決まる。したがって連続スペクトルの形そのものが、間接的に金属量・表面重力（\(\log g\)）の情報を運ぶ。連続と線の双方を
Saha--Boltzmann で一貫して解くのが現代の恒星大気モデルである（第15章
§15.3、第21章 §21.4）。

\end{tcolorbox}

\section{19.6 観測への逆引き}\label{sec-continuum-inverse}

{[}本文目安：B2{]}

本章の三過程は、連続スペクトルの特徴的な形から物理量を読み出す道具になる：

\begin{itemize}
\tightlist
\item
  \textbf{連続吸収端}（束縛-自由、§19.3）：恒星大気の \textbf{Balmer
  jump}
  の大きさから有効温度・表面重力を、端の有無から組成・電離段階を逆引きする。
\item
  \textbf{電波の折れ曲がり}（自由-自由、§19.2）：低振動数の
  \(F_\nu\propto\nu^2\) への飽和から \textbf{電子温度 \(T_e\)}
  を、ターンオーバー振動数から \textbf{発光量度
  \(\mathrm{EM}=\int n_e^2\,ds\)} を読む（第17章 §17.6、演習 17-5）。
\item
  \textbf{X
  線の指数カットオフ}（自由-自由、§19.2）：\(F_\nu\propto e^{-h\nu/k_BT_e}\)
  の傾きから銀河団高温ガスの \(T_e\sim10^7\)〜\(10^8\) K
  を測る（第18章）。
\item
  \textbf{再結合連続光のジャンプ}（自由-束縛、§19.4）：星雲の
  \textbf{Balmer 連続光ジャンプ} の大きさは \(T_e\)
  に依存し、輝線とは独立な電子温度診断を与える。
\end{itemize}

\begin{tcolorbox}[enhanced jigsaw, arc=.35mm, bottomrule=.15mm, bottomtitle=1mm, breakable, colback=white, colbacktitle=quarto-callout-note-color!10!white, colframe=quarto-callout-note-color-frame, coltitle=black, left=2mm, leftrule=.75mm, opacityback=0, opacitybacktitle=0.6, rightrule=.15mm, title=\textcolor{quarto-callout-note-color}{\faInfo}\hspace{0.5em}{ノート}, titlerule=0mm, toprule=.15mm, toptitle=1mm]

\textbf{観測との対応}：恒星大気では Balmer 端は
\textbf{吸収側の段差}（端より青で不透明度増）として、星雲では同じ
\(3646\) Å に \textbf{再結合連続光の放出ジャンプ}
として現れる。同じしきい値が、幾何条件（背景光の有無）で吸収にも放出にも見える
―
これは線スペクトルで同じ遷移が吸収線にも輝線にもなる対比（\hyperref[sec-where-lines-emission]{第20章
§20.3}）と同じ論理である。

\end{tcolorbox}

\section{19.7 再結合は線も生む ―
連続と線をつなぐ電離}\label{sec-continuum-bridge}

{[}本文目安：B2{]}

再結合（§19.4）で電子がいったん高い準位 \(n\) に捕まると、そこからは
\textbf{束縛-束縛遷移}
を次々と下って基底状態へ落ちる（\textbf{再結合カスケード}）。各ステップが線光子を出すので、再結合は連続光（自由-束縛）だけでなく
\textbf{再結合線}（recombination lines）も生む。HII 領域で輝く Hα など
Balmer 輝線は、この再結合カスケードの産物である。

\[
e^- + \mathrm{H}^+ \;\longrightarrow\; \mathrm{H}(n)\ \underbrace{+\,h\nu_{\mathrm{fb}}}_{\text{連続（§19.4）}} \;\overset{\text{カスケード}}{\longrightarrow}\; \text{Balmer・Paschen 輝線（束縛-束縛）}
\]

すなわち
\textbf{電離と再結合は、連続スペクトル（自由-自由・束縛-自由）と線スペクトル（束縛-束縛）の両方を同時に統べる}。case
A／case B（§19.4）の区別は、Lyman 線光子が再吸収されるかどうかという
\textbf{線} の問題が \textbf{連続}
の再結合率に跳ね返る、両者の結合の一例である。

ここから先 ― 束縛-束縛遷移の振動数・強度・選択則・線形 ―
が第VII部の主題である。本章で整えた電離・再結合の枠組みは、第21章（線形成の
Einstein
係数）・第22章（原子準位と選択則）・第24章（電離度診断と再結合線）で線スペクトルの言葉に翻訳され、最終的に第VIII部（第26章）の
QED で連続と線が同じ Hamiltonian の二側面として統合される。

\section*{19.8 使えるようになった道具}\label{sec-continuum-micro-tools}
\addcontentsline{toc}{section}{19.8 使えるようになった道具}

\markright{19.8 使えるようになった道具}

\begin{tcolorbox}[enhanced jigsaw, arc=.35mm, bottomrule=.15mm, bottomtitle=1mm, breakable, colback=white, colbacktitle=quarto-callout-note-color!10!white, colframe=quarto-callout-note-color-frame, coltitle=black, left=2mm, leftrule=.75mm, opacityback=0, opacitybacktitle=0.6, rightrule=.15mm, title=\textcolor{quarto-callout-note-color}{\faInfo}\hspace{0.5em}{ノート}, titlerule=0mm, toprule=.15mm, toptitle=1mm]

\begin{itemize}
\tightlist
\item
  連続を作る三素過程（自由-自由・束縛-自由・自由-束縛）の準位図による区別
\item
  自由-自由の放出 式~\ref{eq-ff-emissivity} と吸収
  式~\ref{eq-ff-absorption-micro}、Kirchhoff
  \(S_\nu=B_\nu(T_e)\)、電波ターンオーバーと X 線カットオフの起源
\item
  光電離断面積 \(\sigma_{\mathrm{bf}}\propto\nu^{-3}\)
  としきい値が作る連続吸収端（Lyman/Balmer/Paschen）
\item
  H⁻ が低温星の可視連続不透明度を担う仕組みと Saha 依存
\item
  Milne 関係による
  \(\sigma_{\mathrm{bf}}\leftrightarrow\sigma_{\mathrm{rec}}\)、再結合係数
  \(\alpha_n(T)\propto T^{-1/2}\) 級
\item
  Saha--Boltzmann による占有数 → 連続係数の自己無撞着な決定
\item
  再結合連続光と再結合線が同一過程から出ること（連続↔線の橋）
\end{itemize}

\end{tcolorbox}

\begin{center}\rule{0.5\linewidth}{0.5pt}\end{center}

\section*{この章で何がわかったか}\label{sec-continuum-micro-summary}
\addcontentsline{toc}{section}{この章で何がわかったか}

\markright{この章で何がわかったか}

\textbf{中心地図に戻る}

中心地図プランク関数 \(B_\nu(T)\)
を作る連続の吸収・放射係数は、自由電子が絡む三つの素過程に分解できる：

\begin{itemize}
\tightlist
\item
  \textbf{自由-自由}（制動放射）：\(\alpha_\nu\propto n_e n_i\,\nu^{-2}T^{-3/2}\)。低
  \(\nu\) で黒体（電波）、高 \(\nu\) で指数カットオフ（X 線）。
\item
  \textbf{束縛-自由}（光電離）：\(\sigma_{\mathrm{bf}}\propto\nu^{-3}\)
  としきい値が連続吸収端を作る。可視は H⁻ が主役。
\item
  \textbf{自由-束縛}（再結合）：束縛-自由の逆過程。Milne
  関係で結ばれ、再結合連続光と再結合線を生む。
\end{itemize}

これらを束ねるのが Saha--Boltzmann（占有数）と
Kirchhoff／Milne（放出と吸収の関係）で、LTE ではソース関数はつねに
\(S_\nu=B_\nu(T)\)
に帰着する。電離・再結合は連続と線の両方を統べ、次の第VII部（線スペクトル）への橋となる。

\section*{演習問題}\label{sec-continuum-micro-exercises}
\addcontentsline{toc}{section}{演習問題}

\markright{演習問題}

\textbf{模範解答}
は折りたたみを展開して確認できる（オンライン版）。まず自力で解いてから開くこと。

\subsection*{問題 19-1　電波 HII
領域のターンオーバー（自由-自由の逆引き）}\label{ex-19-1}
\addcontentsline{toc}{subsection}{問題 19-1　電波 HII
領域のターンオーバー（自由-自由の逆引き）}

{[}★ 難易度：☆☆ {]}

§19.2・第17章 §17.6 の自由-自由吸収
\(\alpha_\nu^{\mathrm{ff}}\propto n_e n_i T_e^{-3/2}\nu^{-2}\)
を使う。等温スラブ（電子温度 \(T_e\)、LTE）の射出強度は
\(I_\nu=B_\nu(T_e)(1-e^{-\tau_\nu})\)。

\begin{enumerate}
\def\labelenumi{\arabic{enumi}.}
\tightlist
\item
  光学的に厚い極限（低 \(\nu\)）で \(I_\nu\to B_\nu(T_e)\)
  となること、電波帯（Rayleigh--Jeans）で \(F_\nu\propto\nu^2\)
  になることを示せ。
\item
  光学的に薄い極限（高 \(\nu\)）で \(F_\nu\)
  がほぼ平坦（\(\propto\nu^{-0.1}\)）になることを、\(\tau_\nu\propto\nu^{-2.1}\)
  から示せ。
\item
  ターンオーバー振動数
  \(\nu_t\)（\(\tau_{\nu_t}\simeq1\)）から何が読めるか述べよ。
\end{enumerate}

\textbf{関連}：\hyperref[sec-continuum-free-free]{§19.2
自由-自由}／電波の漸近は\hyperref[sec-free-free-ionized]{第17章
§17.6}（\hyperref[ex-17-5]{演習
17-5}）、光学的に厚い極限は\hyperref[sec-radiative-transfer-thick]{§5.6}。

\section{模範解答（問題
19-1）}\label{ux6a21ux7bc4ux89e3ux7b54ux554fux984c-19-1}

\textbf{(1)} \(\tau_\nu\gg1\) で \(1-e^{-\tau_\nu}\to1\) ゆえ
\(I_\nu\to B_\nu(T_e)\)。電波帯 \(h\nu\ll k_BT_e\) では
\(B_\nu\simeq2\nu^2k_BT_e/c^2\propto\nu^2\)、よって
\(F_\nu\propto\nu^2\)。観測される輝度温度は \(T_b\to T_e\) に飽和する。

\textbf{(2)} \(\tau_\nu\ll1\) で
\(1-e^{-\tau_\nu}\simeq\tau_\nu\propto\nu^{-2.1}\)。\(F_\nu\propto\nu^2\cdot\nu^{-2.1}=\nu^{-0.1}\)
でほぼ平坦。

\textbf{(3)} \(\nu_t\) は \(\tau_{\nu_t}\simeq1\)
で決まり、\(\tau_\nu\propto\mathrm{EM}\,\nu^{-2.1}T_e^{-1.35}\)
なので、飽和値から \(T_e\)、\(\nu_t\) から発光量度
\(\mathrm{EM}=\int n_e^2 ds\)（密度 × 経路長）が読める。

\textbf{答え}：低 \(\nu\) で \(F_\nu\propto\nu^2\)（\(T_b\to T_e\)）、高
\(\nu\) で \(\propto\nu^{-0.1}\)。飽和から \(T_e\)、\(\nu_t\) から
\(\mathrm{EM}\)。

\subsection*{問題 19-2　Balmer 連続光ジャンプ ―
吸収端と放出端}\label{ex-19-2}
\addcontentsline{toc}{subsection}{問題 19-2　Balmer 連続光ジャンプ ―
吸収端と放出端}

{[}★ 難易度：☆☆ {]}

§19.3・§19.4 の束縛-自由。水素 \(n=2\) の電離エネルギーは
\(\chi_2=13.6/2^2=3.40\) eV。

\begin{enumerate}
\def\labelenumi{\arabic{enumi}.}
\tightlist
\item
  Balmer 端の波長 \(\lambda=hc/\chi_2\) を求め、\(3646\) Å
  に一致することを確かめよ。
\item
  恒星大気では Balmer 端より
  \textbf{短波長側で連続光が暗くなる}（吸収段差）のに、星雲では同じ波長に
  \textbf{明るいジャンプ}（放出）が出る。なぜ符号が逆に見えるか、背景光の有無で説明せよ。
\item
  星雲の Balmer 連続光ジャンプが電子温度 \(T_e\)
  の診断になる理由を一言で述べよ。
\end{enumerate}

\textbf{関連}：\hyperref[sec-continuum-bound-free]{§19.3
束縛-自由}、\hyperref[sec-continuum-recombination]{§19.4
再結合}／吸収線と輝線の幾何は\hyperref[sec-where-lines-emission]{第20章
§20.3}、線のソース関数は\hyperref[sec-line-source-function]{第21章
§21.3}。

\section{模範解答（問題
19-2）}\label{ux6a21ux7bc4ux89e3ux7b54ux554fux984c-19-2}

\textbf{(1)}
\(\lambda=hc/\chi_2=(1.24\times10^4\ \mathrm{eV\,\AA})/3.40\ \mathrm{eV}=3647\)
Å。Balmer 端（\(3646\) Å）に一致。

\textbf{(2)} 恒星大気は背後に明るい連続光源（深層）があり、Balmer
端より青では \(n=2\)
からの光電離吸収（§19.3）で不透明度が跳ね上がるため、より浅い（冷たい）層しか見えず暗くなる＝吸収段差。星雲は背景連続光がなく、\(n=2\)
への再結合（§19.4）で出る自由-束縛放射そのものを見るため、しきい
\(3646\) Å
の高振動数側に放出のジャンプが立つ。同じしきい値が幾何条件で吸収にも放出にも見える。

\textbf{(3)} 再結合連続光の形（ジャンプ高さと裾の傾き）は電子の Maxwell
分布の温度を反映するので、\(T_e\)
に依存する。輝線比とは独立な温度診断になる。

\textbf{答え}：\(3647\)
Å。背景光ありで吸収段差、なしで再結合の放出ジャンプ。ジャンプ形が
\(T_e\) 依存。

\subsection*{問題 19-3　Milne 関係と再結合係数の温度依存}\label{ex-19-3}
\addcontentsline{toc}{subsection}{問題 19-3　Milne
関係と再結合係数の温度依存}

{[}★ 難易度：☆☆☆ {]}

§19.4 の Milne 関係
\(\sigma_{\mathrm{rec}}(v)\propto (h\nu)^2/(v^2)\,\sigma_{\mathrm{bf}}(\nu)\)
と \(h\nu=\tfrac12 m_e v^2+\chi_n\) を使う。

\begin{enumerate}
\def\labelenumi{\arabic{enumi}.}
\tightlist
\item
  しきい近傍（\(\tfrac12 m_e v^2\ll\chi_n\)）で \(h\nu\simeq\chi_n\)
  とみなすと、\(\sigma_{\mathrm{rec}}\propto v^{-2}\sigma_{\mathrm{bf}}(\nu_{\mathrm{th}})\)
  となることを示せ。
\item
  再結合係数 \(\alpha_n(T)=\langle\sigma_{\mathrm{rec}}v\rangle\) を
  Maxwell 分布で平均すると、\(\alpha_n(T)\propto T^{-1/2}\)
  級になることを、\(\sigma_{\mathrm{rec}}v\propto v^{-1}\) と
  \(\langle v^{-1}\rangle\propto T^{-1/2}\) から示せ。
\item
  高温ほど再結合係数が小さい物理的理由を一言で述べよ。
\end{enumerate}

\textbf{関連}：\hyperref[sec-continuum-recombination]{§19.4
再結合}／詳細つり合いの論法は\hyperref[sec-einstein-method-detailed-balance]{第11章
§11.3}、Maxwell
分布は\hyperref[sec-statistical-mechanics-boltzmann]{第7章 §7.3}。

\section{模範解答（問題
19-3）}\label{ux6a21ux7bc4ux89e3ux7b54ux554fux984c-19-3}

\textbf{(1)} Milne 関係
\(\sigma_{\mathrm{rec}}\propto (h\nu)^2 v^{-2}\sigma_{\mathrm{bf}}\)
で、しきい近傍 \(h\nu\simeq\chi_n\)（定数）とすれば
\(\sigma_{\mathrm{rec}}\propto v^{-2}\sigma_{\mathrm{bf}}(\nu_{\mathrm{th}})\)。

\textbf{(2)}
\(\sigma_{\mathrm{rec}}v\propto v^{-2}\cdot v=v^{-1}\)。Maxwell 平均
\(\langle v^{-1}\rangle=\int v^{-1}f(v)dv\) は
\(f(v)\propto v^2 e^{-m_ev^2/2k_BT}\) で、変数を \(x=v/\sqrt{k_BT/m_e}\)
にとると \(\langle v^{-1}\rangle\propto(k_BT)^{-1/2}\)。よって
\(\alpha_n(T)\propto T^{-1/2}\)。

\textbf{(3)}
高温では電子の典型速度が大きく、イオンの近くを速く通り過ぎて捕まりにくい（相互作用時間が短い）ため、再結合係数が小さくなる。

\textbf{答え}：\(\sigma_{\mathrm{rec}}\propto v^{-2}\sigma_{\mathrm{bf}}\)、\(\alpha_n\propto T^{-1/2}\)。高温ほど速く通り過ぎ捕まりにくい。

\subsection*{問題 19-4　X 線制動放射で銀河団温度}\label{ex-19-4}
\addcontentsline{toc}{subsection}{問題 19-4　X 線制動放射で銀河団温度}

{[}★ 難易度：☆☆ {]}

§19.2 の高振動数（光学的に薄い）極限 \(F_\nu\propto e^{-h\nu/k_BT_e}\)
を使う。\(k_B=8.617\times10^{-5}\) eV/K。

\begin{enumerate}
\def\labelenumi{\arabic{enumi}.}
\tightlist
\item
  銀河団 X 線スペクトルの指数カットオフが \(h\nu\sim8\) keV
  付近で効くとき、電子温度 \(T_e\) を概算せよ。
\item
  この \(T_e\) が \(10^7\)〜\(10^8\) K
  の範囲にあることを確かめ、なぜ銀河団ガスが「X 線で光る」のかを述べよ。
\item
  同じ自由-自由放射なのに、電波 HII 領域（問題
  19-1）とは見え方が全く違う理由を、光学的厚さの観点から一言で述べよ。
\end{enumerate}

\textbf{関連}：\hyperref[sec-continuum-free-free]{§19.2
自由-自由}／高温降着・銀河団は\hyperref[ux7b2c18ux7ae0-ux964dux7740ux5186ux76e4ux3068ux9ad8ux30a8ux30cdux30ebux30aeux30fcux5929ux4f53]{第18章}。

\section{模範解答（問題
19-4）}\label{ux6a21ux7bc4ux89e3ux7b54ux554fux984c-19-4}

\textbf{(1)} カットオフは \(h\nu\sim k_BT_e\) で効く。\(k_BT_e\sim8\)
keV
\(\Rightarrow T_e\sim8\times10^3\,\mathrm{eV}/(8.617\times10^{-5}\,\mathrm{eV/K})\simeq9\times10^7\)
K。

\textbf{(2)} \(9\times10^7\) K は \(10^7\)〜\(10^8\) K
の範囲。これだけ高温だと熱的制動放射のピークが X
線帯に来るので、銀河団ガスは X 線で明るく光る。

\textbf{(3)} 電波 HII 領域は低 \(\nu\) で \textbf{光学的に厚く}
黒体に漸近するが、銀河団 X 線は高温・希薄で全帯域 \textbf{光学的に薄い}
ため、放射率の形（指数カットオフ）がそのまま観測される。同じ熱的自由-自由でも光学的厚さで見え方が変わる。

\textbf{答え}：\(T_e\sim9\times10^7\) K（\(10^7\)〜\(10^8\)
K）。高温ゆえ X 線で光る。電波は厚い／X 線は薄いの違い。

\subsection*{問題 19-5　H⁻
のしきい値と太陽光球の連続不透明度}\label{ex-19-5}
\addcontentsline{toc}{subsection}{問題 19-5　H⁻
のしきい値と太陽光球の連続不透明度}

{[}★ 難易度：☆☆ {]}

§19.3 の H⁻。束縛エネルギー \(0.754\) eV。\(hc=1.24\times10^4\) eV·Å。

\begin{enumerate}
\def\labelenumi{\arabic{enumi}.}
\tightlist
\item
  H⁻ の束縛-自由しきい波長
  \(\lambda_{\mathrm{th}}=hc/0.754\ \mathrm{eV}\) を求め、\(1.64\ \mu\)m
  になることを確かめよ。
\item
  中性水素（\(n=1\) の Lyman 端 \(912\) Å、\(n=2\) の Balmer 端 \(3646\)
  Å）の束縛-自由端はいずれも可視中央にない。それでも太陽光球が可視全域で不透明なのはなぜか、H⁻
  で説明せよ。
\item
  H⁻ が低温星でのみ効き、高温星（O・B 型）では効かない理由を、Saha
  平衡（中性 H と自由電子の供給）から述べよ。
\end{enumerate}

\textbf{関連}：\hyperref[sec-continuum-bound-free]{§19.3
H⁻}、\hyperref[sec-continuum-saha]{§19.5
Saha}／恒星大気の振動数依存不透明度は\hyperref[sec-non-blackbody-frequency-opacity]{第15章
§15.3}。

\section{模範解答（問題
19-5）}\label{ux6a21ux7bc4ux89e3ux7b54ux554fux984c-19-5}

\textbf{(1)}
\(\lambda_{\mathrm{th}}=1.24\times10^4\ \mathrm{eV\,\AA}/0.754\ \mathrm{eV}=1.64\times10^4\)
Å \(=1.64\ \mu\)m。

\textbf{(2)} \(1.64\ \mu\)m より短い可視〜近赤外の全域で H⁻
の光電離吸収が効く（ピーク \(\sim0.85\ \mu\)m）。中性 H
の束縛-自由端は紫外（Lyman/Balmer）にしかないので可視を説明できないが、H⁻
が可視連続不透明度を担うため光球は可視で不透明になる。さらに長波長では
H⁻ 自由-自由が続く。

\textbf{(3)} H⁻ には (i) 素材の中性 H と (ii)
余分な電子の供給源（自由電子）が要る。低温星では水素はほとんど中性で豊富、自由電子は電離ポテンシャルの低い金属の電離が供給する。高温星では水素が電離して中性
H が枯れ、H⁻ も作れない（自由電子は豊富でも素材がない）。よって主役は
H・He の束縛-自由／自由-自由と電子散乱に移る。

\textbf{答え}：\(1.64\ \mu\)m。H⁻ が可視連続不透明度を担う。低温星は中性
H 豊富＋金属由来電子で H⁻ が効くが、高温星は H 電離で素材が枯れる。

\section*{さらに学ぶための参考文献}\label{sec-continuum-micro-further}
\addcontentsline{toc}{section}{さらに学ぶための参考文献}

\markright{さらに学ぶための参考文献}

\begin{itemize}
\tightlist
\item
  Rybicki \& Lightman, \emph{Radiative Processes in Astrophysics}
  (Wiley, 1979) --- 自由-自由・束縛-自由・Milne 関係の標準的導出。
\item
  Osterbrock \& Ferland, \emph{Astrophysics of Gaseous Nebulae and
  Active Galactic Nuclei} (2nd ed., 2006) --- 再結合・case
  A/B・再結合連続光と再結合線。
\item
  Mihalas, \emph{Stellar Atmospheres} (W.H. Freeman, 2nd ed., 1978) ---
  H⁻ を含む恒星大気の連続不透明度。
\item
  Gray, \emph{The Observation and Analysis of Stellar Photospheres} (3rd
  ed., 2005) --- 連続不透明度の実務的な扱い。
\end{itemize}

\part{第VII部　線スペクトル ― もう一つの逆引き}

\chapter{第20章　線スペクトルはどこに現れるか}\label{ux7b2c20ux7ae0-ux7ddaux30b9ux30daux30afux30c8ux30ebux306fux3069ux3053ux306bux73feux308cux308bux304b}

\textbf{この章の主題}：第VII部は本書のもう一つの主柱 ―
\textbf{線スペクトル} ―
の逆引きを始める。第1章が連続スペクトル（黒体放射）の観測動機を集めたのと対をなす本章では、太陽・恒星・星雲・電波源・AGN
といった\textbf{線スペクトルが現れる天体}の実例を一覧する。線スペクトルが「宇宙物理の二本目の主柱」であることを、観測の側から確立する。

\section*{この章の中心地図}\label{sec-where-lines-map}
\addcontentsline{toc}{section}{この章の中心地図}

\markright{この章の中心地図}

\begin{figure}[H]

\centering{

\includegraphics[width=0.9\linewidth,height=\textheight,keepaspectratio]{index_files/mediabag/part7/../images/hydrogen-map.pdf}

}

\caption{\label{fig-hydrogen-map}水素線吸収係数 \(\alpha_\nu\)
の各因子と、それを逆引きする章群の対応。\textbf{因子の色＝担当する章/部の色}で結ばれており、矢印に頼らず読み取れる設計。第VII部の§19-23と、線スペクトルが連続側と統合される第VIII部までの流れを一望できる。}

\end{figure}%

\textbf{水素線スペクトルの地図 ― 二層構造}

\textbf{第一層：イコニック・アンカー（リュードベリ公式）}

\begin{equation}\protect\phantomsection\label{eq-rydberg}{
\frac{1}{\lambda} = R_\mathrm{H} \left( \frac{1}{n_1^2} - \frac{1}{n_2^2} \right)
}\end{equation}

線スペクトルの「骨格」を一目で示す。\(R_\mathrm{H}\)
は基本定数の組合せ、\(1/n^2\)
は束縛状態のエネルギー、離散性が「線が存在する理由」を語る。

\textbf{第二層：構造的マップ（Hα を標準例とした線吸収係数）}

\begin{equation}\protect\phantomsection\label{eq-line-absorption}{
\alpha_\nu = \frac{\pi e^2}{m_e c} \, f_{lu} \, n_l \, \phi(\nu - \nu_0)
}\end{equation}

{\def\LTcaptype{none} % do not increment counter
\begin{longtable}[]{@{}lll@{}}
\toprule\noalign{}
因子 & 物理的意味 & 対応する章 \\
\midrule\noalign{}
\endhead
\bottomrule\noalign{}
\endlastfoot
\(\alpha_\nu\) & 観測量（線吸収係数） & 第21章 \\
\(\pi e^2 / m_e c\) & 古典極限の線強度 & 第22章＋第VIII部 \\
\(f_{lu}\) & 振動子強度 & 第22章 \\
\(n_l\) & 下準位の占有数 & 第21章＋Part III \\
\(\phi(\nu - \nu_0)\) & 線輪郭（自然・Doppler・Voigt） & 第23章 \\
\(\nu_0\) & 線位置（リュードベリ） & 第22章 \\
\end{longtable}
}

\textbf{適用範囲（重要）}：この式の前因子 \(\pi e^2/m_e c\) と振動子強度
\(f_{lu}\) は、古典 Lorentz 振動子に対応する
\textbf{電気双極子（E1）許容線} に対するものである。21 cm
線（基底状態超微細構造の \textbf{磁気双極子 M1 遷移}）や
\textbf{禁制線（E2・M1 など）} では、この前因子・\(f_{lu}\)
をそのまま使えない。これらの線では、測定された自発放出率 \(A_{ul}\) から
Einstein
関係を介して有効遷移強度（\(B_{lu}\)）に読み替える（式~\ref{eq-line-absorption}
の構造は保ったまま、強度因子だけを差し替える）。

\begin{tcolorbox}[enhanced jigsaw, arc=.35mm, bottomrule=.15mm, breakable, colback=white, colframe=quarto-callout-tip-color-frame, left=2mm, leftrule=.75mm, opacityback=0, rightrule=.15mm, toprule=.15mm]

\textbf{地図の使い方}：第VII部の各章はこの式の異なる因子を解き明かす。プランク関数地図と並べて読めば、本書を貫く二つの逆引きが同時に把握できる。第VIII部の
QED で、両者は前因子 \(2h\nu^3/c^2\)
を介して統一される（式~\ref{eq-einstein-A} 参照）。

\end{tcolorbox}

\begin{tcolorbox}[enhanced jigsaw, arc=.35mm, bottomrule=.15mm, bottomtitle=1mm, breakable, colback=white, colbacktitle=quarto-callout-note-color!10!white, colframe=quarto-callout-note-color-frame, coltitle=black, left=2mm, leftrule=.75mm, opacityback=0, opacitybacktitle=0.6, rightrule=.15mm, title=\textcolor{quarto-callout-note-color}{\faInfo}\hspace{0.5em}{ノート}, titlerule=0mm, toprule=.15mm, toptitle=1mm]

\textbf{方針}：第VII部の中心地図は \textbf{水素線吸収係数}
\(\alpha_\nu = (\pi e^2/m_e c) f_{lu} n_l \phi(\nu-\nu_0)\)。本章ではまだ式の中身に踏み込まず、観測される線スペクトルの
\textbf{「ファクト」} を提示する。第21-23章で各因子（\(f_{lu}\),
\(n_l\), \(\phi\)）を逆引きする。

\end{tcolorbox}

\section*{この章で答える問い}\label{sec-where-lines-questions}
\addcontentsline{toc}{section}{この章で答える問い}

\markright{この章で答える問い}

\begin{itemize}
\tightlist
\item
  なぜ太陽光のスペクトルには細い暗線（Fraunhofer 線）が見えるのか（→
  第21章、§20.1）
\item
  同じ水素から、ある天体では吸収線、別の天体では輝線が見えるのはなぜか（→
  §21.4 線ソース関数）
\item
  21 cm 線は何の遷移で、なぜ宇宙物理で特別な役割を持つのか（→ §22.5
  超微細構造の磁気双極子遷移、§24.5）
\item
  分子線は原子線とどう違い、何を診断するのに使えるのか（→ §22.6
  分子の回転・振動準位）
\item
  AGN・クェーサーの幅広い輝線は、何を語っているのか（→ §23.5
  速度場、§18.5 ブラックホール近傍）
\end{itemize}

\section*{到達目標}\label{sec-where-lines-goals}
\addcontentsline{toc}{section}{到達目標}

\markright{到達目標}

この章を読み終えると、読者は次のことができるようになる：

\begin{itemize}
\tightlist
\item
  宇宙のさまざまな天体に現れる線スペクトル（Fraunhofer 線・輝線・21
  cm・分子線・AGN 線）を例示できる
\item
  同じ遷移から吸収線にも輝線にもなる条件を、放射輸送の観点から説明できる
\item
  線スペクトルが連続スペクトルと並んで、宇宙物理の二本目の主柱であることを位置づけられる
\end{itemize}

\begin{center}\rule{0.5\linewidth}{0.5pt}\end{center}

\section{20.1 太陽の Fraunhofer 線 ―
線スペクトル発見の起点}\label{sec-where-lines-fraunhofer}

{[}本文目安：高校-B1{]}

1814 年、Joseph von Fraunhofer
は太陽光をプリズムで分解し、連続スペクトルの中に
\textbf{数百本の細い暗線} を発見した。彼はそれぞれの線に A, B, C,
\ldots{}
と名付け、後にこれらが地球上の元素の吸収線と一致することが判明する：

\begin{itemize}
\tightlist
\item
  \textbf{D 線} (\(\lambda \simeq 589\) nm)：ナトリウム
\item
  \textbf{C 線} (\(\lambda = 656\) nm, H\(\alpha\))：水素
\item
  \textbf{K 線、H 線} (\(\lambda \simeq 393, 397\)
  nm)：イオン化カルシウム
\end{itemize}

このとき同時に、線スペクトルが \textbf{宇宙物理の中核観測量}
になることが運命づけられた：「太陽が
\textbf{何で出来ているか}」を直接観測する手段が、連続スペクトルとは別の経路で確立されたからである。

\begin{tcolorbox}[enhanced jigsaw, arc=.35mm, bottomrule=.15mm, bottomtitle=1mm, breakable, colback=white, colbacktitle=quarto-callout-note-color!10!white, colframe=quarto-callout-note-color-frame, coltitle=black, left=2mm, leftrule=.75mm, opacityback=0, opacitybacktitle=0.6, rightrule=.15mm, title=\textcolor{quarto-callout-note-color}{\faInfo}\hspace{0.5em}{ノート}, titlerule=0mm, toprule=.15mm, toptitle=1mm]

\textbf{観測との対応}：Fraunhofer 線は
\textbf{太陽の表層（光球と彩層）の物質組成}
を直接告げる。地上の元素と一致するということは、太陽も同じ元素から出来ているという、宇宙物理の根幹的事実の発見だった。同じ手法で恒星・銀河・QSO
に対して \textbf{元素組成}
を測ることが、本部のクライマックス（第24章）の応用に繋がる。

\end{tcolorbox}

\section{20.2
恒星の吸収線とスペクトル分類}\label{sec-where-lines-stellar-classification}

{[}本文目安：B2{]}

太陽以外の恒星も、それぞれの温度に応じて異なる吸収線パターンを示す。これがスペクトル分類
\textbf{OBAFGKM} の物理的基盤：

{\def\LTcaptype{none} % do not increment counter
\begin{longtable}[]{@{}
  >{\raggedright\arraybackslash}p{(\linewidth - 6\tabcolsep) * \real{0.2500}}
  >{\raggedright\arraybackslash}p{(\linewidth - 6\tabcolsep) * \real{0.2500}}
  >{\raggedright\arraybackslash}p{(\linewidth - 6\tabcolsep) * \real{0.2500}}
  >{\raggedright\arraybackslash}p{(\linewidth - 6\tabcolsep) * \real{0.2500}}@{}}
\toprule\noalign{}
\begin{minipage}[b]{\linewidth}\raggedright
型
\end{minipage} & \begin{minipage}[b]{\linewidth}\raggedright
\(T_\text{eff}\) (K)
\end{minipage} & \begin{minipage}[b]{\linewidth}\raggedright
主な強い線
\end{minipage} & \begin{minipage}[b]{\linewidth}\raggedright
物理
\end{minipage} \\
\midrule\noalign{}
\endhead
\bottomrule\noalign{}
\endlastfoot
O & \(\gtrsim 30{,}000\) & He II 線、N III 輝線 & He\(^+\)（He
II）線が現れるほど高温 \\
B & \(10{,}000\)〜\(30{,}000\) & He I、強い水素線 &
水素は大部分電離せず励起 \\
A & \(7{,}500\)〜\(10{,}000\) & H 線がピーク & 水素励起が最大効率 \\
F & \(6{,}000\)〜\(7{,}500\) & Ca II H,K 線が現れる &
水素励起がやや弱まる \\
G & \(5{,}200\)〜\(6{,}000\) & Ca II 強く、金属線 & 太陽はここ \\
K & \(3{,}700\)〜\(5{,}200\) & 金属線、分子帯（CN 等） &
低温で分子も出始める \\
M & \(\lesssim 3{,}700\) & TiO バンド、強い分子線 & 分子線が支配 \\
\end{longtable}
}

水素のバルマー線（H\(\alpha\), H\(\beta\), \ldots）の \textbf{強さは A
型でピーク} になる ― これが Saha-Boltzmann 分布の効果（第VII部の中心地図
\(n_l\)
因子）。低温では水素励起準位の占有が低く、高温では水素が電離してしまう。

\begin{tcolorbox}[enhanced jigsaw, arc=.35mm, bottomrule=.15mm, breakable, colback=white, colframe=quarto-callout-tip-color-frame, left=2mm, leftrule=.75mm, opacityback=0, rightrule=.15mm, toprule=.15mm]
\begin{minipage}[t]{5.5mm}
\textcolor{quarto-callout-tip-color}{\faLightbulb}
\end{minipage}%
\begin{minipage}[t]{\textwidth - 5.5mm}

\textbf{問い}：「水素は宇宙で最も多い元素」なのに、なぜ水素吸収線は
\textbf{A 型恒星でだけ最強} になり、太陽（G 型）では H\(\alpha\)
が比較的弱いのか？

\textbf{短答}：水素線が見えるためには、基底状態（\(n=1\)）ではなく
\textbf{\(n=2\)（励起状態）に電子がいる必要がある}（バルマー線は
\(n=2 \to n \ge 3\) の遷移）。占有確率
\(\propto e^{-13.6\,\text{eV} \cdot 3/4 / kT}\)
が大きいには高温が必要だが、温度が高すぎると水素が電離して水素原子そのものが消える。両者のバランスから、水素線強度は
\(T \simeq 10{,}000\) K（A 型）で最大になる。

\textbf{もう一歩}：これは第7章 §7.3（Boltzmann
分布）で予告した「中心地図の \(n_l\)
因子が観測スペクトルに直接効く」の具体例。Saha-Boltzmann
から定量化される

\end{minipage}%
\end{tcolorbox}

\section{20.3 星雲の輝線スペクトル ―
吸収線とは何が違うのか}\label{sec-where-lines-emission}

{[}本文目安：B2-B3{]}

\begin{figure}[H]

\centering{

\includegraphics[width=0.92\linewidth,height=\textheight,keepaspectratio]{index_files/mediabag/part7/../images/line-formation-geometry.pdf}

}

\caption{\label{fig-line-geometry}同じ原子遷移が、観測の幾何条件によって
\textbf{吸収線} にも \textbf{輝線} にも見える仕組み。(a)
背後に明るい連続光源があってその前を冷たいガスが覆う →
吸収線（恒星大気）、(b) 背景に連続光がなく、励起されたガス自体を観測 →
輝線（星雲）、(c) 温度勾配のある大気 →
表面付近では冷たく見えるため吸収線。Kirchhoff
の三法則の現代的視覚化、§21.4 でも再登場する。}

\end{figure}%

恒星のスペクトルが \textbf{連続光に重なった吸収線}
であるのに対し、星雲（HII 領域、惑星状星雲、超新星残骸など）では
\textbf{連続光がほとんどなく、輝線が主役} になる。

代表例：

\begin{itemize}
\tightlist
\item
  \textbf{HII 領域}（オリオン大星雲など）：H\(\alpha\), H\(\beta\), {[}O
  III{]} \$\lambda\$5007, {[}N II{]}, {[}S II{]} 等の
  \textbf{輝線スペクトル}
\item
  \textbf{惑星状星雲}（環状星雲など）：似た輝線＋{[}O III{]}
  が極端に強い
\item
  \textbf{超新星残骸}：水素・酸素・硫黄等の輝線、衝撃波と関連した特殊な輪郭
\end{itemize}

これらの天体が \textbf{輝線スペクトル} を示すのは、次の二条件の組合せ：

\begin{enumerate}
\def\labelenumi{\arabic{enumi}.}
\tightlist
\item
  \textbf{励起源}（多くは隣の高温星の紫外光、または衝撃波）が原子を励起・電離させる
\item
  励起された原子が \textbf{自発放出}
  で下準位に落ちるとき、その線振動数だけが選択的に強く出る
\end{enumerate}

連続光は、星雲ガス自体が比較的低密度・低温で熱平衡的な連続放射を作れないため、相対的に弱い。

\begin{tcolorbox}[enhanced jigsaw, arc=.35mm, bottomrule=.15mm, bottomtitle=1mm, breakable, colback=white, colbacktitle=quarto-callout-note-color!10!white, colframe=quarto-callout-note-color-frame, coltitle=black, left=2mm, leftrule=.75mm, opacityback=0, opacitybacktitle=0.6, rightrule=.15mm, title=\textcolor{quarto-callout-note-color}{\faInfo}\hspace{0.5em}{ノート}, titlerule=0mm, toprule=.15mm, toptitle=1mm]

\textbf{用語}：\textbf{禁制線}（\textbf{forbidden line}）― {[}O III{]}
\$\lambda\$5007 のように、四角括弧で囲んだ線。原子物理の選択則（第22章
§22.5）では本来「禁じられた」遷移だが、地上では衝突で先に脱励起されるため見えない一方、\textbf{極低密度}
の星雲環境では時間が許されて自発放出が起こる。HII 領域の代表線である
{[}O III{]}
が「禁制線」を冠するのは、\textbf{星雲が地上では実現できない希薄環境であることの直接観測的証拠}

\textbf{観測との対応}：禁制線は
\textbf{通常の実験室条件（高密度）では衝突脱励起で潰されて見えにくい遷移}。ただし極低密度の系では地球上でも観測される
― 例えば \textbf{オーロラの \([\text{O I}]\)
\$\lambda\$5577}（緑光）は地球超高層大気の典型的な禁制線で、EBIT
等の希薄プラズマ実験室でも測られる。むしろ「禁制線が出る
\(\Leftrightarrow\)
低密度」という関係を読みとる手段として、星雲・上層大気・宇宙環境の共通の道具になる。第24章
§24.3 で電子密度・温度の診断に活躍する

\end{tcolorbox}

\section{20.4 21 cm 線 ―
中性水素の宇宙論的トレーサー}\label{sec-where-lines-21cm}

{[}本文目安：B3{]}

水素原子の基底状態 (\(n=1\)) には、電子スピン ↑ と陽子スピン ↑ が
\textbf{平行か反平行か} で微小なエネルギー差がある（\textbf{超微細構造},
\textbf{hyperfine structure}）。エネルギー差はわずか
\(\Delta E \simeq 5.9 \times 10^{-6}\) eV、対応する振動数は

\begin{equation}\protect\phantomsection\label{eq-21cm-frequency}{
\nu_{21\text{cm}} = 1420.405\,\text{MHz} \quad (\lambda = 21.106\,\text{cm})
}\end{equation}

これが \textbf{21 cm 線}（\textbf{21 cm line}, \textbf{HI line}）。

宇宙物理での 21 cm 線の特別な役割：

\begin{itemize}
\tightlist
\item
  \textbf{広がっている}：宇宙の中性水素はどこにでもある（銀河系内の星間ガス、銀河ハロー、銀河間媒質）
\item
  \textbf{電波で観測}：可視光と違って塵に吸収されない →
  銀河系の遠くまで見える
\item
  \textbf{回転曲線}：銀河の中性水素分布を 21 cm
  線で測れば、\textbf{ダークマターの存在の歴史的証拠}となった回転曲線が得られる
\item
  \textbf{宇宙再電離期}：高赤方偏移 \(z \sim 6\)〜\(20\)
  の中性水素分布が、宇宙再電離期の探査ツール（SKA 等が目指す）
\end{itemize}

\begin{tcolorbox}[enhanced jigsaw, arc=.35mm, bottomrule=.15mm, breakable, colback=white, colframe=quarto-callout-tip-color-frame, left=2mm, leftrule=.75mm, opacityback=0, rightrule=.15mm, toprule=.15mm]
\begin{minipage}[t]{5.5mm}
\textcolor{quarto-callout-tip-color}{\faLightbulb}
\end{minipage}%
\begin{minipage}[t]{\textwidth - 5.5mm}

\textbf{問い}：21 cm 線の自発放出確率
\(A_{ul} \simeq 2.9 \times 10^{-15}\) s\(^{-1}\)
は、可視光の水素線（例：Ly\(\alpha\) の \(A \simeq 6.3 \times 10^8\)
s\(^{-1}\)）より \(\sim 10^{23}\)
倍も小さい。それでも観測できるのはなぜか？

\textbf{短答}：圧倒的に多くの \textbf{中性水素が宇宙に存在する}
から。第VII部の中心地図
\(\alpha_\nu = (\pi e^2/m_e c) f_{lu} n_l \phi(\nu-\nu_0)\) の \(n_l\)
が非常に大きいので、実効的な吸収強度が極端に小さくても観測可能な吸収（または輝線）として現れる。

\textbf{もう一歩}：個別の水素原子では一千万年に一度しか脱励起しない遷移も、宇宙的スケールの大量の中性水素が積算すれば、地上で測れる電波信号になる。「弱い遷移＋膨大な物質量」の組合せが
21 cm 線を宇宙物理ツールにしている。

\textbf{注意（前因子の扱い）}：中心地図の前因子 \(\pi e^2/m_e c\)
は、古典 Lorentz 振動子に対応する \textbf{電気双極子（E1）遷移}
の積分吸収断面積である（§22.4b）。21 cm 線は基底状態超微細構造の
\textbf{磁気双極子（M1）遷移} なので、この前因子と通常の振動子強度
\(f_{lu}\) をそのまま当てはめることはできない。21 cm
線の吸収・放出強度は、測定された自発放出率
\(A_{ul} \simeq 2.9 \times 10^{-15}\) s\(^{-1}\) から Einstein
関係を介して読み替える（式~\ref{eq-A-f-conversion} の関係を E1
に限らず一般化し、\(A_{ul}\) ベースで \(B_{lu}\) を求める）必要がある。

\end{minipage}%
\end{tcolorbox}

\section{20.5 分子線 ―
電波・サブミリの主役}\label{sec-where-lines-molecular}

{[}本文目安：B3{]}

分子は原子と違って \textbf{回転・振動の準位}
を持つので、線スペクトルがはるかに豊か。代表例：

\begin{itemize}
\tightlist
\item
  \textbf{CO (一酸化炭素)} \(J=1 \to 0\) at 115
  GHz：分子雲のトレーサー、星間で最重要
\item
  \textbf{H\(_2\)O} maser 22 GHz：星形成領域・銀河中心
\item
  \textbf{OH} maser 1.6 GHz：星間ガス・恒星質量放出
\item
  \textbf{NH\(_3\), HCN, HCO+, CS, \ldots{}}
  ：暗黒分子雲の温度・密度診断
\end{itemize}

分子の特徴：

\begin{itemize}
\tightlist
\item
  \textbf{低エネルギー遷移}：回転準位の間隔は典型的に
  \(h\nu \sim k_B \cdot 10\) K 程度 → 電波・サブミリ
\item
  \textbf{多数の線}：選択則と回転状態数の組合せで、複雑な線スペクトル
\item
  \textbf{化学のトレーサー}：どの分子が見えるかが化学組成と物理条件を反映
\end{itemize}

\begin{tcolorbox}[enhanced jigsaw, arc=.35mm, bottomrule=.15mm, bottomtitle=1mm, breakable, colback=white, colbacktitle=quarto-callout-note-color!10!white, colframe=quarto-callout-note-color-frame, coltitle=black, left=2mm, leftrule=.75mm, opacityback=0, opacitybacktitle=0.6, rightrule=.15mm, title=\textcolor{quarto-callout-note-color}{\faInfo}\hspace{0.5em}{ノート}, titlerule=0mm, toprule=.15mm, toptitle=1mm]

\textbf{観測との対応}：星間分子線は
\textbf{暗黒分子雲（星形成の場）の物理条件}
を測る現代天文学の主要観測対象。ALMA、IRAM、NOEMA、JCMT
等のサブミリ・ミリ波望遠鏡の主要ターゲット。線スペクトルの「もう一つの主柱」性が明示される領域である。

\end{tcolorbox}

\section{20.6 AGN・クェーサーの線 ―
広線と狭線の物語}\label{sec-where-lines-agn}

{[}本文目安：B3{]}

活動銀河中心核（AGN）やクェーサー（QSO）のスペクトルには、印象的な
\textbf{二種類の輝線} が現れる：

\begin{itemize}
\tightlist
\item
  \textbf{広線（broad emission line）}：FWHM \(\sim 1000\)〜\(10{,}000\)
  km/s。Hα, Hβ, C IV \$\lambda\$1549, Mg II \$\lambda\$2798 等
\item
  \textbf{狭線（narrow emission line）}：FWHM \(\sim 100\)〜\(1000\)
  km/s。{[}O III{]} \$\lambda\$5007, {[}N II{]}, {[}S II{]} 等
\end{itemize}

両者の \textbf{線幅} から、放射領域の速度場（運動学）が逆引きできる：

\begin{itemize}
\tightlist
\item
  広線領域（BLR）：中心ブラックホールの近く（\(\sim 0.01\)〜\(1\)
  pc）、高速ケプラー回転または乱流
\item
  狭線領域（NLR）：もっと外側（\(\sim 100\)〜\(1000\)
  pc）、ホスト銀河スケールの速度
\end{itemize}

\begin{tcolorbox}[enhanced jigsaw, arc=.35mm, bottomrule=.15mm, bottomtitle=1mm, breakable, colback=white, colbacktitle=quarto-callout-note-color!10!white, colframe=quarto-callout-note-color-frame, coltitle=black, left=2mm, leftrule=.75mm, opacityback=0, opacitybacktitle=0.6, rightrule=.15mm, title=\textcolor{quarto-callout-note-color}{\faInfo}\hspace{0.5em}{ノート}, titlerule=0mm, toprule=.15mm, toptitle=1mm]

\textbf{観測との対応}：QSO の \textbf{広線}
から、ブラックホール質量を測る（\textbf{reverberation
mapping}：連続光変動と広線変動の遅れから BLR サイズを測る →
ケプラー速度と組合せて中心質量）。これは現代宇宙論の標準ツール。線スペクトルの幅を観測すれば、\(10^9 M_\odot\)
の超大質量ブラックホールの質量が決まる、というのは線幅から力学量を推定できる代表例である。

\end{tcolorbox}

\section{20.7
線スペクトルの逆引きへ向けて}\label{sec-where-lines-toward-inverse}

{[}本文目安：B2{]}

第20章で見た事実をまとめると、宇宙の線スペクトルは次の
\textbf{多様な情報} を運んでいる：

{\def\LTcaptype{none} % do not increment counter
\begin{longtable}[]{@{}lll@{}}
\toprule\noalign{}
観測量 & 物理量（逆引き先） & 章 \\
\midrule\noalign{}
\endhead
\bottomrule\noalign{}
\endlastfoot
線が出ている振動数 & 原子・分子のエネルギー準位 & 第22章 \\
吸収線 or 輝線 & 線形成（ソース関数と局所放射場） & 第21章 \\
線の強度・等価幅 & 占有数 \(n_l\)、振動子強度 \(f_{lu}\) & 第21-22章 \\
線の輪郭・幅 & 温度・速度場・密度・磁場 & 第23章 \\
線比 & 電離度・電子密度・化学組成 & 第24章 \\
\end{longtable}
}

第21-24章で、これら \textbf{観測量 → 物理量}
の逆引き経路を一つずつ整備する。これは第I-VI部で連続スペクトルに対して行ったのと
\textbf{対応する構造の作業} である。

\begin{tcolorbox}[enhanced jigsaw, arc=.35mm, bottomrule=.15mm, bottomtitle=1mm, breakable, colback=white, colbacktitle=quarto-callout-note-color!10!white, colframe=quarto-callout-note-color-frame, coltitle=black, left=2mm, leftrule=.75mm, opacityback=0, opacitybacktitle=0.6, rightrule=.15mm, title=\textcolor{quarto-callout-note-color}{\faInfo}\hspace{0.5em}{ノート}, titlerule=0mm, toprule=.15mm, toptitle=1mm]

\textbf{使えるようになった道具}（本章では観測例の把握、本格的道具立ては第21-24章）：

\begin{itemize}
\tightlist
\item
  線スペクトルが宇宙のあらゆる場所で観測されるという基本ファクト
\item
  吸収線・輝線・禁制線・21 cm 線・分子線・AGN 線という観測カテゴリ
\item
  線スペクトルから読み取れる物理量カタログ（逆引き先一覧）
\item
  第VII部の中心地図 \(\alpha_\nu\)
  の各因子（\(f_{lu}, n_l, \phi\)）の役割の予告
\end{itemize}

\end{tcolorbox}

\begin{center}\rule{0.5\linewidth}{0.5pt}\end{center}

\section*{この章で何がわかったか}\label{sec-where-lines-summary}
\addcontentsline{toc}{section}{この章で何がわかったか}

\markright{この章で何がわかったか}

\textbf{中心地図に戻る}

本章は第VII部の
\textbf{観測動機の章}。連続スペクトルが第1章で「宇宙のあちこちで見える熱的連続光」だったように、線スペクトルも宇宙のあちこちで観測される
― 太陽・恒星・星雲・電波源・AGN。

第VII部の中心地図
\(\alpha_\nu = (\pi e^2/m_e c) f_{lu} n_l \phi(\nu-\nu_0)\)
には次の章で踏み込む：

\begin{itemize}
\tightlist
\item
  第21章：線形成の物理 ― 吸収線と輝線の出る条件
\item
  第22章：原子物理 ― 線振動数と \(f_{lu}\) の起源
\item
  第23章：線輪郭 ― \(\phi(\nu-\nu_0)\) の物理機構
\item
  第24章：観測応用 ― 線スペクトルからの物理量読み取りの実践
\end{itemize}

\section*{演習問題}\label{sec-where-lines-exercises}
\addcontentsline{toc}{section}{演習問題}

\markright{演習問題}

\textbf{模範解答}
は折りたたみを展開して確認できる（オンライン版）。まず自力で解いてから開くこと。

\subsection*{問題 20-1　21 cm 線 ―
極端に弱い遷移がなぜ観測できるか}\label{ex-20-1}
\addcontentsline{toc}{subsection}{問題 20-1　21 cm 線 ―
極端に弱い遷移がなぜ観測できるか}

{[}★ 難易度：☆☆ {]}

§20.4 は 21 cm 線の自発放出率 \(A_{ul}\simeq2.9\times10^{-15}\)
s\(^{-1}\) が Ly\(\alpha\)（\(A\simeq6.3\times10^8\) s\(^{-1}\)）より約
\(10^{23}\) 倍小さいと述べた。中心地図
\(\alpha_\nu=(\pi e^2/m_ec)f_{lu}n_l\phi\) で考える。

\begin{enumerate}
\def\labelenumi{\arabic{enumi}.}
\tightlist
\item
  個々の原子では一千万年に一度しか脱励起しない遷移が、観測可能な信号になる理由を、\(\alpha_\nu\)
  のどの因子で説明するか述べよ。
\item
  励起寿命 \(\tau=1/A_{ul}\) を年単位で求めよ。
\item
  21 cm
  線が宇宙物理で特別な役割を持つ理由を3つ挙げよ（観測波長帯・ダスト・対象天体の観点から）。
\end{enumerate}

\textbf{関連}：\hyperref[sec-where-lines-21cm]{§20.4 21 cm
線}／中心地図の \(n_l\)
は\hyperref[sec-line-source-function]{§21.3}、超微細構造の M1
遷移は\hyperref[sec-qed-minimum-selection]{§26.5}、宇宙論応用は\hyperref[sec-line-applications-cosmology]{§24.5}。

\section{模範解答（問題
20-1）}\label{ux6a21ux7bc4ux89e3ux7b54ux554fux984c-20-1}

\textbf{(1)} 中心地図 \(\alpha_\nu=(\pi e^2/m_ec)f_{lu}\,n_l\,\phi\) の
\textbf{\(n_l\)（下準位占有数＝柱密度）}。21 cm の振動子強度
\(f_{lu}\)（\(\propto A_{ul}\)）は極端に小さいが、宇宙には中性水素が膨大に存在し
\(n_l\)
が桁外れに大きい。両者の積として観測可能な吸収・輝線になる。「弱い遷移＋膨大な物質量」の組合せ。

\textbf{(2)} \(\tau=1/A_{ul}=1/(2.9\times10^{-15})=3.4\times10^{14}\)
s。年に直すと \(3.4\times10^{14}/(3.16\times10^7)\simeq1.1\times10^7\)
年（約一千万年）。

\textbf{(3)} ① 電波（21
cm）なのでダストに吸収されず銀河系の遠くまで見通せる。②
中性水素は宇宙のどこにでもある（星間ガス・ハロー・銀河間媒質）。③
銀河回転曲線（ダークマターの証拠）や高赤方偏移の宇宙再電離期（SKA
の標的）を探れる。

\textbf{答え}：\(n_l\)（膨大な中性水素）が小さな \(f_{lu}\)
を補う。\(\tau\simeq1.1\times10^7\)
年。ダスト透過・遍在・回転曲線/再電離で特別。

\subsection*{問題 20-2　禁制線は「低密度の証拠」}\label{ex-20-2}
\addcontentsline{toc}{subsection}{問題 20-2　禁制線は「低密度の証拠」}

{[}★ 難易度：☆☆ {]}

§20.3 は {[}O III{]} \$\lambda\$5007
等の禁制線が星雲で強く見える一方、地上実験室では見えにくいと述べた。

\begin{enumerate}
\def\labelenumi{\arabic{enumi}.}
\tightlist
\item
  禁制線が地上の高密度環境で見えにくい理由を、衝突脱励起と自発放出の競合から述べよ。
\item
  極低密度の星雲（\(n_e\sim10^2\)--\(10^4\)
  cm\(^{-3}\)）でなぜ自発放出が間に合うのか。臨界密度の概念に触れて述べよ。
\item
  「禁制線が出ている」という観測事実が、その天体について何を語るか一言で述べよ。
\end{enumerate}

\textbf{関連}：\hyperref[sec-where-lines-emission]{§20.3
星雲の輝線}／選択則と禁制線の次数は\hyperref[sec-atomic-selection-rules]{§22.5}（\hyperref[ex-22-3]{演習
22-3}）、密度診断は\hyperref[sec-line-applications-diagnostics]{§24.3}（\hyperref[ex-24-2]{演習
24-2}）。

\section{模範解答（問題
20-2）}\label{ux6a21ux7bc4ux89e3ux7b54ux554fux984c-20-2}

\textbf{(1)} 禁制線の上準位は自発放出率 \(A_{ul}\)
が極端に小さい（長寿命）。高密度では、上準位の原子が自発放出で光子を出す前に、他粒子との衝突で脱励起されてしまう（衝突脱励起）。その結果、線として放出されず見えない。

\textbf{(2)} 衝突脱励起率は密度 \(n_e\) に比例。自発放出率 \(A_{ul}\)
と衝突脱励起率が拮抗する密度が\textbf{臨界密度}
\(n_\mathrm{crit}\sim A_{ul}/\langle\sigma v\rangle\)。星雲は \(n_e\)
が臨界密度より低いので、衝突より先に自発放出が起こり、禁制線が放出される。

\textbf{(3)}
禁制線が観測される＝その天体の密度が臨界密度より低い、すなわち\textbf{極低密度環境である}ことの直接の証拠。地上では実現できない希薄さの指標になる。

\textbf{答え}：高密度では衝突脱励起が自発放出に勝ち禁制線が潰れる。星雲は臨界密度以下で自発放出が間に合う。禁制線の存在＝低密度の証拠。

\subsection*{問題 20-3　AGN 広線からブラックホール質量}\label{ex-20-3}
\addcontentsline{toc}{subsection}{問題 20-3　AGN
広線からブラックホール質量}

{[}★ 難易度：☆☆ {]}

§20.6 は AGN
の広線幅から中心ブラックホール質量が測れると述べた（reverberation
mapping）。広線領域（BLR）の運動をケプラー回転（ビリアル）とみなす：\(M_\mathrm{BH}\simeq\dfrac{v^2 R}{G}\)。

\begin{enumerate}
\def\labelenumi{\arabic{enumi}.}
\tightlist
\item
  広線の FWHM から速度 \(v\) を、reverberation mapping から BLR 半径
  \(R\) を得る手順を一言で述べよ。
\item
  \(v=5000\) km/s、\(R=20\) 光日として \(M_\mathrm{BH}\)
  を太陽質量単位で概算せよ（1 光日 \(=2.59\times10^{13}\)
  m、\(G=6.67\times10^{-11}\)、\(M_\odot=1.99\times10^{30}\) kg）。
\item
  狭線（FWHM \(\sim\) 数百
  km/s）は同じ手法で中心質量を測れない。なぜか述べよ。
\end{enumerate}

\textbf{関連}：\hyperref[sec-where-lines-agn]{§20.6
AGN・クェーサーの線}／ブラックホール近傍の相対論的効果は\hyperref[sec-black-hole-effects]{§18.5}（\hyperref[ux7b2c18ux7ae0-ux964dux7740ux5186ux76e4ux3068ux9ad8ux30a8ux30cdux30ebux30aeux30fcux5929ux4f53]{第VI部
演習}）、線幅は\hyperref[sec-line-shape-macro]{§23.5}。

\section{模範解答（問題
20-3）}\label{ux6a21ux7bc4ux89e3ux7b54ux554fux984c-20-3}

\textbf{(1)} 広線の FWHM（ドップラー幅）から BLR ガスの速度分散 \(v\)
を得る。reverberation mapping は連続光の変動と広線の変動の時間遅れ
\(\Delta t\) から \(R=c\,\Delta t\)（光路差）で BLR 半径を測る。両者を
\(M=v^2R/G\) に入れる。

\textbf{(2)} \(v=5\times10^6\)
m/s、\(R=20\times2.59\times10^{13}=5.18\times10^{14}\)
m。\(M=\dfrac{(5\times10^6)^2\times5.18\times10^{14}}{6.67\times10^{-11}}=\dfrac{2.5\times10^{13}\times5.18\times10^{14}}{6.67\times10^{-11}}=\dfrac{1.30\times10^{28}}{6.67\times10^{-11}}=1.94\times10^{38}\)
kg。太陽質量で
\(1.94\times10^{38}/1.99\times10^{30}\simeq9.7\times10^7\ M_\odot\)（約
\(10^8\ M_\odot\)）。

\textbf{(3)} 狭線領域（NLR）は中心から \(\sim100\)--\(1000\) pc
も外側にあり、その速度はホスト銀河スケールの運動（中心ブラックホールの重力が支配しない）を反映する。よって
NLR 速度に \(M=v^2R/G\)
を当てても中心ブラックホール質量にはならない。中心質量を測るには中心近傍の
BLR が必要。

\textbf{答え}：FWHM→\(v\)、時間遅れ→\(R=c\Delta t\)。\(M_\mathrm{BH}\simeq10^8\,M_\odot\)。狭線は中心重力圏外なので使えない。

\section*{さらに学ぶための参考文献}\label{sec-where-lines-further}
\addcontentsline{toc}{section}{さらに学ぶための参考文献}

\markright{さらに学ぶための参考文献}

\begin{itemize}
\tightlist
\item
  Osterbrock \& Ferland, \emph{Astrophysics of Gaseous Nebulae and
  Active Galactic Nuclei} (University Science Books, 2nd ed., 2006) ---
  星雲と AGN の線スペクトル
\item
  Gray, \emph{The Observation and Analysis of Stellar Photospheres}
  (Cambridge, 3rd ed., 2005) --- 恒星吸収線の標準書
\item
  Draine, \emph{Physics of the Interstellar and Intergalactic Medium}
  (Princeton, 2011) --- 21 cm 線・分子線・ISM の線スペクトル
\item
  Netzer, \emph{The Physics and Evolution of Active Galactic Nuclei}
  (Cambridge, 2013) --- AGN の広線・狭線、reverberation mapping
\end{itemize}

\chapter{第21章　線形成 ―
吸収線と輝線の物理}\label{ux7b2c21ux7ae0-ux7ddaux5f62ux6210-ux5438ux53ceux7ddaux3068ux8f1dux7ddaux306eux7269ux7406}

\textbf{この章の主題}：第20章で見たさまざまな線スペクトルは、すべて「\textbf{原子・分子の離散的遷移と、その場の放射場との相互作用}」の結果として理解できる。本章では、第11章で導入した
\textbf{Einstein 係数} \(A_{ul}, B_{lu}, B_{ul}\)
を線スペクトルの主役として再登場させ、\textbf{線吸収・線輝線・線ソース関数}
を逆引きで体系化する。これは連続スペクトルの第4-5章（放射輸送とソース関数）の
\textbf{線スペクトル版} にあたる。

\section*{この章の中心地図}\label{sec-line-formation-map}
\addcontentsline{toc}{section}{この章の中心地図}

\markright{この章の中心地図}

\textbf{水素線スペクトルの地図 ― 二層構造}

\textbf{第一層：イコニック・アンカー（リュードベリ公式）}

\begin{equation}\protect\phantomsection\label{eq-rydberg}{
\frac{1}{\lambda} = R_\mathrm{H} \left( \frac{1}{n_1^2} - \frac{1}{n_2^2} \right)
}\end{equation}

線スペクトルの「骨格」を一目で示す。\(R_\mathrm{H}\)
は基本定数の組合せ、\(1/n^2\)
は束縛状態のエネルギー、離散性が「線が存在する理由」を語る。

\textbf{第二層：構造的マップ（Hα を標準例とした線吸収係数）}

\begin{equation}\protect\phantomsection\label{eq-line-absorption}{
\alpha_\nu = \frac{\pi e^2}{m_e c} \, f_{lu} \, n_l \, \phi(\nu - \nu_0)
}\end{equation}

{\def\LTcaptype{none} % do not increment counter
\begin{longtable}[]{@{}lll@{}}
\toprule\noalign{}
因子 & 物理的意味 & 対応する章 \\
\midrule\noalign{}
\endhead
\bottomrule\noalign{}
\endlastfoot
\(\alpha_\nu\) & 観測量（線吸収係数） & 第21章 \\
\(\pi e^2 / m_e c\) & 古典極限の線強度 & 第22章＋第VIII部 \\
\(f_{lu}\) & 振動子強度 & 第22章 \\
\(n_l\) & 下準位の占有数 & 第21章＋Part III \\
\(\phi(\nu - \nu_0)\) & 線輪郭（自然・Doppler・Voigt） & 第23章 \\
\(\nu_0\) & 線位置（リュードベリ） & 第22章 \\
\end{longtable}
}

\textbf{適用範囲（重要）}：この式の前因子 \(\pi e^2/m_e c\) と振動子強度
\(f_{lu}\) は、古典 Lorentz 振動子に対応する
\textbf{電気双極子（E1）許容線} に対するものである。21 cm
線（基底状態超微細構造の \textbf{磁気双極子 M1 遷移}）や
\textbf{禁制線（E2・M1 など）} では、この前因子・\(f_{lu}\)
をそのまま使えない。これらの線では、測定された自発放出率 \(A_{ul}\) から
Einstein
関係を介して有効遷移強度（\(B_{lu}\)）に読み替える（式~\ref{eq-line-absorption}
の構造は保ったまま、強度因子だけを差し替える）。

\begin{tcolorbox}[enhanced jigsaw, arc=.35mm, bottomrule=.15mm, breakable, colback=white, colframe=quarto-callout-tip-color-frame, left=2mm, leftrule=.75mm, opacityback=0, rightrule=.15mm, toprule=.15mm]

\textbf{地図の使い方}：第VII部の各章はこの式の異なる因子を解き明かす。プランク関数地図と並べて読めば、本書を貫く二つの逆引きが同時に把握できる。第VIII部の
QED で、両者は前因子 \(2h\nu^3/c^2\)
を介して統一される（式~\ref{eq-einstein-A} 参照）。

\end{tcolorbox}

\begin{tcolorbox}[enhanced jigsaw, arc=.35mm, bottomrule=.15mm, bottomtitle=1mm, breakable, colback=white, colbacktitle=quarto-callout-note-color!10!white, colframe=quarto-callout-note-color-frame, coltitle=black, left=2mm, leftrule=.75mm, opacityback=0, opacitybacktitle=0.6, rightrule=.15mm, title=\textcolor{quarto-callout-note-color}{\faInfo}\hspace{0.5em}{ノート}, titlerule=0mm, toprule=.15mm, toptitle=1mm]

\textbf{方針}：本章では中心地図
\(\alpha_\nu = (\pi e^2/m_e c) f_{lu} n_l \phi(\nu-\nu_0)\) の
\textbf{\(n_l\) 因子} と \textbf{線形成（吸収 or 輝線）の選択}
に焦点を当てる。\(f_{lu}\) は第22章、\(\phi(\nu-\nu_0)\)
は第23章で別途扱う。

\end{tcolorbox}

\section*{この章で答える問い}\label{sec-line-formation-questions}
\addcontentsline{toc}{section}{この章で答える問い}

\markright{この章で答える問い}

\begin{itemize}
\tightlist
\item
  同じ水素 H\(\alpha\) 線が、太陽では吸収線、HII
  領域では輝線になるのはなぜか
\item
  線の等価幅とは何で、なぜ観測量として有用なのか
\item
  LTE では何が単純で、non-LTE では何が複雑になるのか
\item
  成長曲線（curve of growth）は何を、どのように測る道具なのか
\end{itemize}

\section*{到達目標}\label{sec-line-formation-goals}
\addcontentsline{toc}{section}{到達目標}

\markright{到達目標}

この章を読み終えると、読者は次のことができるようになる：

\begin{itemize}
\tightlist
\item
  線の吸収・自発放出・誘導放出を Einstein
  係数で表現し、各々の物理的意味を述べられる
\item
  線ソース関数とプランク関数の関係を、LTE と non-LTE
  で区別して説明できる
\item
  線の等価幅と成長曲線（curve of
  growth）から物理量を読み取る基本手順を理解する
\end{itemize}

\begin{center}\rule{0.5\linewidth}{0.5pt}\end{center}

\section{21.1 線の観測量 ― 等価幅}\label{sec-line-equivalent-width}

{[}本文目安：B2-B3{]}

吸収線の \textbf{観測量} として最も標準的なのが
\textbf{等価幅}（\textbf{equivalent width}, \textbf{EW}）。

\begin{equation}\protect\phantomsection\label{eq-equivalent-width}{
W_\lambda \equiv \int_\text{line} \frac{I_c - I_\lambda}{I_c} d\lambda
}\end{equation}

\begin{figure}[H]

\centering{

\includegraphics[width=0.85\linewidth,height=\textheight,keepaspectratio]{index_files/mediabag/part7/../images/equivalent-width.pdf}

}

\caption{\label{fig-equivalent-width}等価幅 \(W_\lambda\) の幾何的な定義
― \textbf{吸収線が連続光から引き抜いた光の総量（青の面積）}
を、深さが連続光と同じ高さの \textbf{赤の矩形}
に置き換えたときの幅。輪郭の細かい形状によらない積分量なので、観測しやすく理論との比較も容易。}

\end{figure}%

ここで \(I_c\) は線がない場合に予想される連続光強度、\(I_\lambda\)
は実際の観測強度。\(W_\lambda\) は
\textbf{「線によって連続光からどれだけ光が引き抜かれたか」を波長単位で表した量}。物理的には：

\begin{itemize}
\tightlist
\item
  浅くて広い線と、深くて狭い線が \textbf{同じ \(W_\lambda\)}
  になる場合がある
\item
  線輪郭（line profile）の詳細によらない積分量
\item
  観測しやすく、理論との比較がしやすい
\end{itemize}

輝線についても同様に、連続光からの \textbf{余剰光}
の積分として等価幅を定義する。

\begin{tcolorbox}[enhanced jigsaw, arc=.35mm, bottomrule=.15mm, breakable, colback=white, colframe=quarto-callout-warning-color-frame, left=2mm, leftrule=.75mm, opacityback=0, rightrule=.15mm, toprule=.15mm]
\begin{minipage}[t]{5.5mm}
\textcolor{quarto-callout-warning-color}{\faExclamationTriangle}
\end{minipage}%
\begin{minipage}[t]{\textwidth - 5.5mm}

\textbf{注意（観測）}：等価幅は
\textbf{観測スペクトル分解能の限界以上の線}
には鋭敏である。一方、強く飽和した線（線芯が完全に黒）では、追加の物質を加えても等価幅はゆっくりとしか増えない（§21.6
の成長曲線の「飽和域」）。等価幅から物理量を読むには、線がどの「成長領域」にあるかを意識する必要がある。

\end{minipage}%
\end{tcolorbox}

\section{21.2 線の三つの素過程 ― Einstein
係数の再登場}\label{sec-line-three-processes}

{[}本文目安：B3{]}

第11章 §11.2 で導入した \textbf{Einstein
係数}（同じ図を再掲、図~\ref{fig-einstein-three-processes-line}）：

\begin{figure}[H]

\centering{

\includegraphics[width=0.88\linewidth,height=\textheight,keepaspectratio]{index_files/mediabag/part7/../images/einstein-processes.pdf}

}

\caption{\label{fig-einstein-three-processes-line}三つの素過程の準位図（再掲）。連続スペクトル側（第IV部）で導入した同じ図が、線スペクトル側でも線形成の出発点となる。両中心地図が共通の
Hamiltonian の二側面であることを示す最初の表れ。}

\end{figure}%

\begin{itemize}
\tightlist
\item
  \(A_{ul}\)：\textbf{自発放出}（\textbf{spontaneous emission}）―
  \(u \to l\) で光子放出。確率 {[}s\(^{-1}\){]}
\item
  \(B_{lu}\)：\textbf{誘導吸収}（\textbf{induced absorption}）―
  \(l \to u\) で光子吸収
\item
  \(B_{ul}\)：\textbf{誘導放出}（\textbf{stimulated emission}）―
  \(u \to l\)、入射光と同じモードへ放出
\end{itemize}

これらは原子物理（第22章で再導出）と熱平衡条件（第11章 §11.3 の Einstein
関係式）

\begin{equation}\protect\phantomsection\label{eq-einstein-relations}{
\frac{A_{ul}}{B_{ul}} = \frac{2h\nu^3}{c^2}, \quad \frac{g_l B_{lu}}{g_u B_{ul}} = 1
}\end{equation}

で結ばれる。ここで \(g_l, g_u\) は下/上準位の縮重度。これらは
\textbf{線形成の主役} である。

線吸収係数を Einstein 係数で書き直すと

\begin{equation}\protect\phantomsection\label{eq-line-absorption-einstein}{
\alpha_\nu = \frac{h\nu_0}{4\pi} n_l B_{lu} \phi(\nu - \nu_0) \left[1 - \frac{g_l n_u}{g_u n_l}\right]
}\end{equation}

右辺の括弧内が \textbf{誘導放出補正}。\(n_u > 0\)
なら吸収を弱める効果を表す。

\begin{tcolorbox}[enhanced jigsaw, arc=.35mm, bottomrule=.15mm, breakable, colback=white, colframe=quarto-callout-tip-color-frame, left=2mm, leftrule=.75mm, opacityback=0, rightrule=.15mm, toprule=.15mm]
\begin{minipage}[t]{5.5mm}
\textcolor{quarto-callout-tip-color}{\faLightbulb}
\end{minipage}%
\begin{minipage}[t]{\textwidth - 5.5mm}

\textbf{問い}：誘導放出補正 \([1 - g_l n_u/(g_u n_l)]\) が \textbf{ゼロ}
になるのはどんな条件か？ \textbf{負} になるのはどんなとき？
物理的意味は？

\textbf{短答}：

\begin{itemize}
\tightlist
\item
  ゼロ：\(g_l n_u/g_u = n_l\) ―
  上下準位の占有が同等。「これ以上吸収できない」
\item
  負：\(n_u/g_u > n_l/g_l\) ― \textbf{準位反転}（\textbf{population
  inversion}）。負の吸収係数 = \textbf{増幅}（メーザー、レーザー）
\end{itemize}

\textbf{もう一歩}：銀河中心の H\(_2\)O メーザー（22 GHz）、星形成領域の
OH メーザー（1.6 GHz）はこの \textbf{天然のレーザー}
現象。観測すれば、\(n_u/g_u > n_l/g_l\)
を実現する非平衡ポンプ機構（衝突、放射ポンプ）の存在が逆引きされる

\end{minipage}%
\end{tcolorbox}

\section{21.3 線ソース関数 ―
連続スペクトルから線スペクトルへ}\label{sec-line-source-function}

{[}本文目安：B3{]}

第5章 §5.4 で定義した \textbf{ソース関数} \(S_\nu\) は、線スペクトルでは

\begin{equation}\protect\phantomsection\label{eq-line-source-function}{
S_\nu^\text{line} = \frac{j_\nu^\text{line}}{\alpha_\nu^\text{line}}
}\end{equation}

の形で構成される。ここで \(j_\nu^\text{line}\) は線の
\textbf{放出係数}（\textbf{emission coefficient}）で、本書では
\textbf{自発放出のみ} をその寄与とする。一方、\(\alpha_\nu^\text{line}\)
は線吸収係数で、\textbf{誘導放出を「負の吸収」として吸収係数の中に繰り込む流儀}（式~\ref{eq-line-absorption-einstein}
の補正項
\([1 - g_l n_u/(g_u n_l)]\)）を採用する。すなわち本書では一貫して、誘導放出は吸収係数側の補正として扱い、放出係数
\(j_\nu^\text{line}\) には含めない。Einstein
係数を使うと、\(\phi(\nu-\nu_0)\) が両者で共通になる前提のもとで

\begin{equation}\protect\phantomsection\label{eq-line-S-explicit}{
S_\nu^\text{line} = \frac{A_{ul} n_u}{B_{lu} n_l - B_{ul} n_u} = \frac{2h\nu_0^3/c^2}{(g_u n_l)/(g_l n_u) - 1}
}\end{equation}

となる（分子が自発放出
\(A_{ul} n_u\)、分母が誘導放出を差し引いた正味の吸収
\(B_{lu} n_l - B_{ul} n_u\)）。これは形式的に
\textbf{プランク関数と同じ構造} だが、\(h\nu/k_BT\) の代わりに
\(\ln[(g_u n_l)/(g_l n_u)]\) が指数因子に入っている。これを
\textbf{「励起温度」}（\textbf{excitation temperature}）\(T_\text{ex}\)
を使って

\begin{equation}\protect\phantomsection\label{eq-excitation-temperature}{
\frac{g_u n_l}{g_l n_u} = e^{h\nu_0/k_B T_\text{ex}}
}\end{equation}

と書き直すと、\(S_\nu^\text{line} = B_{\nu_0}(T_\text{ex})\)。すなわち
\textbf{線ソース関数は、励起温度に対応する黒体スペクトル} である。

\begin{tcolorbox}[enhanced jigsaw, arc=.35mm, bottomrule=.15mm, bottomtitle=1mm, breakable, colback=white, colbacktitle=quarto-callout-note-color!10!white, colframe=quarto-callout-note-color-frame, coltitle=black, left=2mm, leftrule=.75mm, opacityback=0, opacitybacktitle=0.6, rightrule=.15mm, title=\textcolor{quarto-callout-note-color}{\faInfo}\hspace{0.5em}{ノート}, titlerule=0mm, toprule=.15mm, toptitle=1mm]

\textbf{用語}：\textbf{励起温度} \(T_\text{ex}\)
は、観測者が観測される線の強度から \textbf{逆引き}
で求める「上下準位の占有比を Boltzmann
分布で書いたときの温度」である。物理的に物質と熱平衡している温度とは限らない（LTE
/ non-LTE の区別、§21.4-§21.5）。

\textbf{観測との対応}：観測される線の輝度温度から \(T_\text{ex}\)
を測れる ―
これが線スペクトルから物理状態を逆引きする最初のステップである。

\end{tcolorbox}

\section{\texorpdfstring{21.4 LTE 線形成 ― 線ソース関数 ≈
\(B_\nu(T_\text{kin})\)}{21.4 LTE 線形成 ― 線ソース関数 ≈ B\_\textbackslash nu(T\_\textbackslash text\{kin\})}}\label{sec-LTE-line-formation}

{[}本文目安：B3{]}

\textbf{LTE}（局所熱平衡、第6章 §6.5）が成り立つ環境では、占有数
\(n_l, n_u\) は局所の \textbf{キネティック温度}
\(T_\text{kin}\)（粒子の運動エネルギー由来の温度）に対する Boltzmann
分布で書ける：

\begin{equation}\protect\phantomsection\label{eq-boltzmann}{
\frac{n_u}{g_u} = \frac{n_l}{g_l} e^{-h\nu_0/k_B T_\text{kin}}
}\end{equation}

すると \(T_\text{ex} = T_\text{kin}\) となり、線ソース関数は
\textbf{キネティック温度のプランク関数} になる：

\begin{equation}\protect\phantomsection\label{eq-LTE-source-function}{
S_\nu^\text{line, LTE} = B_\nu(T_\text{kin})
}\end{equation}

これが「LTE の良い」点 ― \textbf{線スペクトルが連続スペクトルと同じ温度}
で記述できる。例えば、恒星大気の Eddington-Barbier 近似（第15章
§15.3）が線スペクトルにも適用できる：観測される線輝度 = 線で光学的深さ
2/3 になる層の温度 \(T_\text{kin}\) のプランク関数値。

\begin{tcolorbox}[enhanced jigsaw, arc=.35mm, bottomrule=.15mm, breakable, colback=white, colframe=quarto-callout-tip-color-frame, left=2mm, leftrule=.75mm, opacityback=0, rightrule=.15mm, toprule=.15mm]
\begin{minipage}[t]{5.5mm}
\textcolor{quarto-callout-tip-color}{\faLightbulb}
\end{minipage}%
\begin{minipage}[t]{\textwidth - 5.5mm}

\textbf{問い}：LTE で、線が \textbf{吸収線として} 観測される条件は何か？

\textbf{短答}：観測する視線方向に、\textbf{温度が深さとともに上がる}
温度勾配がある場合（恒星大気の通常）。線の中心では浅い表面しか見えず（\(T\)
が低い）、連続光は深い層から（\(T\) が高い）。差し引き、線が
\textbf{連続光より暗く} 見える → 吸収線。

逆に、温度が深さとともに \textbf{下がる}
場合（彩層・コロナ）や、観測対象がガス雲そのもの（背景の連続光がない）では、輝線として現れる。\textbf{吸収線
↔ 輝線は、温度構造と背景光の組合せ} で決まる

\end{minipage}%
\end{tcolorbox}

\begin{figure}[H]

\centering{

\includegraphics[width=0.88\linewidth,height=\textheight,keepaspectratio]{index_files/mediabag/part7/../images/line-formation-geometry.pdf}

}

\caption{\label{fig-line-geometry-lte}同じ原子遷移が観測の幾何条件で吸収線にも輝線にもなる三つのケース（第20章
§20.3 で導入した図を、§21.4
でソース関数の文脈で再掲）。温度勾配と背景光の組合せがどう線形態を決めるかが視覚的に分かる。}

\end{figure}%

\section{21.5 non-LTE 線形成 ― いつ・なぜ LTE
が破綻するか}\label{sec-non-LTE-line-formation}

{[}本文目安：B3-B4{]}

実際には LTE が破綻する状況が多い：

\begin{itemize}
\tightlist
\item
  希薄ガス（衝突率が低い）
\item
  強い放射場（放射過程が衝突過程より速い）
\item
  急激な温度勾配（局所平衡が追いつかない）
\end{itemize}

このとき、占有数 \(n_l, n_u\) は Boltzmann
分布から外れ、\textbf{励起温度 \(T_\text{ex}\) がキネティック温度
\(T_\text{kin}\) と乖離する}。線ソース関数は \(B_\nu(T_\text{ex})\)
のままだが、\(T_\text{ex}\) が物理温度を反映しない。

non-LTE 線形成を扱う数値的枠組み：

\begin{enumerate}
\def\labelenumi{\arabic{enumi}.}
\tightlist
\item
  \textbf{占有数の方程式}：各準位への入力（励起・吸収）＝各準位からの出力（脱励起・放出）の収支
\item
  \textbf{放射輸送方程式}：占有数で決まる \(S_\nu\) と \(\alpha_\nu\)
  から、放射場を解く
\item
  \textbf{両者の反復}：放射場と占有数が \textbf{自己無撞着}
  になるまで反復
\end{enumerate}

これが non-LTE radiative transfer の標準枠組み。

\begin{tcolorbox}[enhanced jigsaw, arc=.35mm, bottomrule=.15mm, breakable, colback=white, colframe=quarto-callout-warning-color-frame, left=2mm, leftrule=.75mm, opacityback=0, rightrule=.15mm, toprule=.15mm]
\begin{minipage}[t]{5.5mm}
\textcolor{quarto-callout-warning-color}{\faExclamationTriangle}
\end{minipage}%
\begin{minipage}[t]{\textwidth - 5.5mm}

\textbf{注意（観測）}：「LTE
解析からの偏差」は、観測者が誤った結論に至る源。一例として、\textbf{恒星大気の
Fe I 線} を LTE で解析すると、紫外側で \textbf{Fe の
over-ionization（過電離）} により占有数が Saha 予測を下回り、組成決定で
0.1〜0.4 dex のずれが生じることが知られている（Asplund 2005
のレビュー）。non-LTE 補正は現代の線スペクトル解析の必須道具となっている

\end{minipage}%
\end{tcolorbox}

\begin{tcolorbox}[enhanced jigsaw, arc=.35mm, bottomrule=.15mm, bottomtitle=1mm, breakable, colback=white, colbacktitle=quarto-callout-note-color!10!white, colframe=quarto-callout-note-color-frame, coltitle=black, left=2mm, leftrule=.75mm, opacityback=0, opacitybacktitle=0.6, rightrule=.15mm, title=\textcolor{quarto-callout-note-color}{\faInfo}\hspace{0.5em}{ノート}, titlerule=0mm, toprule=.15mm, toptitle=1mm]

\textbf{用語の整理}（第6章 §6.5 と整合）：

\begin{itemize}
\tightlist
\item
  \textbf{LTE}：物質側の占有数が局所キネティック温度の Boltzmann/Saha
  分布 → 線ソース関数 = \(B_\nu(T_\text{kin})\)
\item
  \textbf{non-LTE}：物質側の占有数が局所温度の Boltzmann から外れる →
  線ソース関数 ≠ \(B_\nu(T_\text{kin})\)、しかし依然として
  \(= B_\nu(T_\text{ex})\)
\item
  \textbf{熱的 vs 非熱的}：粒子のエネルギー分布が Maxwell-Boltzmann
  か否か。LTE / non-LTE とは別の軸（第15章 §15.1）
\end{itemize}

\end{tcolorbox}

\section{21.6 成長曲線（curve of growth） ―
等価幅から柱密度を測る}\label{sec-curve-of-growth}

{[}本文目安：B3{]}

弱い吸収線から強い吸収線まで、観測される \textbf{等価幅} \(W\) と
\textbf{柱密度} \(N \equiv \int n_l\, ds\)
の関係は、\textbf{成長曲線}（\textbf{curve of
growth}、等価幅曲線とも）として知られる：

\begin{equation}\protect\phantomsection\label{eq-curve-of-growth}{
W(N) = \begin{cases}
\propto N & \text{(線形領域)} \\
\propto \sqrt{\ln N} & \text{(飽和領域)} \\
\propto \sqrt{N} & \text{(減衰翼領域)}
\end{cases}
}\end{equation}

物理的解釈：

\begin{itemize}
\tightlist
\item
  \textbf{線形領域}：吸収線の光学的深さ
  \(\tau < 1\)。等価幅は柱密度に比例
\item
  \textbf{飽和領域}：線芯が完全に黒（\(\tau \gg 1\)）。中心が「これ以上吸収できない」、ただ広がるだけ
\item
  \textbf{減衰翼領域}：自然幅・衝突幅が支配（第23章）。Lorentz 翼の
  \(\propto \sqrt{N}\) 依存
\end{itemize}

\begin{tcolorbox}[enhanced jigsaw, arc=.35mm, bottomrule=.15mm, bottomtitle=1mm, breakable, colback=white, colbacktitle=quarto-callout-note-color!10!white, colframe=quarto-callout-note-color-frame, coltitle=black, left=2mm, leftrule=.75mm, opacityback=0, opacitybacktitle=0.6, rightrule=.15mm, title=\textcolor{quarto-callout-note-color}{\faInfo}\hspace{0.5em}{ノート}, titlerule=0mm, toprule=.15mm, toptitle=1mm]

\textbf{観測との対応}：恒星・QSO 吸収線スペクトルの解析で、成長曲線は
\textbf{観測等価幅 → 柱密度} を求める基本道具である。化学組成、QSO
視線方向の星間吸収系の柱密度測定などに使われる。第24章 §24.2
で本格的に応用する。

\end{tcolorbox}

\section*{21.7 使えるようになった道具}\label{sec-line-formation-tools}
\addcontentsline{toc}{section}{21.7 使えるようになった道具}

\markright{21.7 使えるようになった道具}

\begin{tcolorbox}[enhanced jigsaw, arc=.35mm, bottomrule=.15mm, bottomtitle=1mm, breakable, colback=white, colbacktitle=quarto-callout-note-color!10!white, colframe=quarto-callout-note-color-frame, coltitle=black, left=2mm, leftrule=.75mm, opacityback=0, opacitybacktitle=0.6, rightrule=.15mm, title=\textcolor{quarto-callout-note-color}{\faInfo}\hspace{0.5em}{ノート}, titlerule=0mm, toprule=.15mm, toptitle=1mm]

\begin{itemize}
\tightlist
\item
  等価幅 \(W_\lambda\) ― 線の最も基本的観測量
\item
  Einstein 係数による吸収・輝線・誘導放出の統一記述
\item
  線ソース関数 \(S_\nu^\text{line} = B_\nu(T_\text{ex})\)
  と励起温度の概念
\item
  LTE と non-LTE の区別、いつどちらが使えるか
\item
  成長曲線（curve of growth）― \(W \to N\) の換算手段
\end{itemize}

\end{tcolorbox}

\begin{center}\rule{0.5\linewidth}{0.5pt}\end{center}

\section*{この章で何がわかったか}\label{sec-line-formation-summary}
\addcontentsline{toc}{section}{この章で何がわかったか}

\markright{この章で何がわかったか}

\textbf{中心地図に戻る}

本章で、中心地図 \(\alpha_\nu\) の \textbf{\(n_l\) 因子} と
\textbf{線形成の選択（吸収 or 輝線）} が物理的に解読された：

\begin{itemize}
\tightlist
\item
  \(n_l\) は LTE では Boltzmann/Saha、non-LTE では statistical
  equilibrium で決まる
\item
  線が吸収か輝線かは、線ソース関数
  \(S_\nu^\text{line} = B_\nu(T_\text{ex})\)
  と背景の連続光輝度の比較で決まる
\item
  等価幅 \(W\) と柱密度 \(N\) は成長曲線（curve of growth）で結ばれる
\end{itemize}

次章では中心地図のもう一つの因子 ― \textbf{振動子強度 \(f_{lu}\)} ―
を、原子物理から逆引きする。

\section*{演習問題}\label{sec-line-formation-exercises}
\addcontentsline{toc}{section}{演習問題}

\markright{演習問題}

\textbf{模範解答}
は折りたたみを展開して確認できる（オンライン版）。まず自力で解いてから開くこと。

\subsection*{問題 21-1　線ソース関数が Bν(Tex) になる}\label{ex-21-1}
\addcontentsline{toc}{subsection}{問題 21-1　線ソース関数が Bν(Tex)
になる}

{[}★ 難易度：☆☆ {]}

§21.3 は線ソース関数を励起温度 \(T_\mathrm{ex}\)
のプランク関数に書き換えた。導出を埋める。\(S_\nu^\mathrm{line}=\dfrac{A_{ul}n_u}{B_{lu}n_l-B_{ul}n_u}\)、Einstein
関係 \(A_{ul}/B_{ul}=2h\nu_0^3/c^2\)、\(g_lB_{lu}=g_uB_{ul}\)。

\begin{enumerate}
\def\labelenumi{\arabic{enumi}.}
\tightlist
\item
  Einstein 関係を使って
  \(S_\nu^\mathrm{line}=\dfrac{2h\nu_0^3/c^2}{(g_un_l)/(g_ln_u)-1}\)
  を導け。
\item
  励起温度を \(\dfrac{g_un_l}{g_ln_u}=e^{h\nu_0/k_BT_\mathrm{ex}}\)
  で定義し、\(S_\nu^\mathrm{line}=B_{\nu_0}(T_\mathrm{ex})\) を示せ。
\item
  \(T_\mathrm{ex}\)
  が物理的なキネティック温度と一致するのはどんな条件か（LTE）。一致しない場合（non-LTE）に
  \(S_\nu\) がプランク関数でなくなるわけではない理由を述べよ。
\end{enumerate}

\textbf{関連}：\hyperref[sec-line-source-function]{§21.3
線ソース関数}／Einstein 関係の導出は\hyperref[ex-11-1]{演習
11-1}、ソース関数の基礎は\hyperref[ex-5-5]{演習 5-5}、LTE
は\hyperref[sec-LTE-line-formation]{§21.4}。

\section{模範解答（問題
21-1）}\label{ux6a21ux7bc4ux89e3ux7b54ux554fux984c-21-1}

\textbf{(1)} 分子分母を \(B_{ul}n_u\) で割る。\(g_lB_{lu}=g_uB_{ul}\)
より \(B_{lu}=（g_u/g_l)B_{ul}\)：

\[S_\nu^\mathrm{line}=\frac{A_{ul}n_u}{B_{lu}n_l-B_{ul}n_u}=\frac{A_{ul}/B_{ul}}{\dfrac{B_{lu}n_l}{B_{ul}n_u}-1}=\frac{2h\nu_0^3/c^2}{\dfrac{g_un_l}{g_ln_u}-1}.\]

\textbf{(2)} \(\dfrac{g_un_l}{g_ln_u}=e^{h\nu_0/k_BT_\mathrm{ex}}\)
を代入すると分母が \(e^{h\nu_0/k_BT_\mathrm{ex}}-1\) となり

\[S_\nu^\mathrm{line}=\frac{2h\nu_0^3/c^2}{e^{h\nu_0/k_BT_\mathrm{ex}}-1}=B_{\nu_0}(T_\mathrm{ex}).\]

\textbf{(3)} 占有数が局所キネティック温度 \(T_\mathrm{kin}\) の
Boltzmann 分布で書ける（LTE）とき
\(T_\mathrm{ex}=T_\mathrm{kin}\)。non-LTE では占有数が Boltzmann
から外れて \(T_\mathrm{ex}\ne T_\mathrm{kin}\)
になるが、\(T_\mathrm{ex}\) を「占有比を Boltzmann
形で書いたときの温度」と\textbf{定義}する限り、\(S_\nu=B_\nu(T_\mathrm{ex})\)
は形式的に常に成り立つ。ただし \(T_\mathrm{ex}\)
は物理温度を反映しない。

\textbf{答え}：\(S_\nu^\mathrm{line}=\dfrac{2h\nu_0^3/c^2}{(g_un_l)/(g_ln_u)-1}=B_{\nu_0}(T_\mathrm{ex})\)。LTE
で \(T_\mathrm{ex}=T_\mathrm{kin}\)、non-LTE でも定義上
\(B_\nu(T_\mathrm{ex})\) の形は保たれる。

\subsection*{問題
21-2　誘導放出補正と準位反転（メーザー）}\label{ex-21-2}
\addcontentsline{toc}{subsection}{問題
21-2　誘導放出補正と準位反転（メーザー）}

{[}★ 難易度：☆☆ {]}

§21.2 の線吸収係数
\(\alpha_\nu=\dfrac{h\nu_0}{4\pi}n_lB_{lu}\phi\left[1-\dfrac{g_ln_u}{g_un_l}\right]\)
の誘導放出補正 \([1-g_ln_u/(g_un_l)]\) を読み解く。

\begin{enumerate}
\def\labelenumi{\arabic{enumi}.}
\tightlist
\item
  補正項がゼロになる条件（\(n_u/g_u=n_l/g_l\)）の物理的意味を述べよ。
\item
  補正項が負になる条件（\(n_u/g_u>n_l/g_l\)）を求め、そのとき
  \(\alpha_\nu<0\)（負の吸収＝増幅）になることを示せ。これが何の現象か。
\item
  宇宙メーザー（H\(_2\)O 22 GHz、OH 1.6
  GHz）が観測されることから、何が逆引きできるか述べよ。
\end{enumerate}

\textbf{関連}：\hyperref[sec-line-three-processes]{§21.2
線の三つの素過程}／誘導放出は\hyperref[sec-einstein-method-three-processes]{§11.2}、自発/誘導放出の比は\hyperref[ex-11-2]{演習
11-2}。

\section{模範解答（問題
21-2）}\label{ux6a21ux7bc4ux89e3ux7b54ux554fux984c-21-2}

\textbf{(1)} \(\dfrac{g_ln_u}{g_un_l}=1\) すなわち
\(\dfrac{n_u}{g_u}=\dfrac{n_l}{g_l}\)
で補正がゼロ。上下準位の縮重度あたりの占有が等しく、誘導放出が誘導吸収を完全に打ち消す。「これ以上正味で吸収できない」状態。

\textbf{(2)} \(\dfrac{g_ln_u}{g_un_l}>1\) すなわち
\(\dfrac{n_u}{g_u}>\dfrac{n_l}{g_l}\)（\textbf{準位反転}）のとき補正項が負になり
\(\alpha_\nu<0\)。負の吸収係数は、光が媒質を通るほど強められる＝\textbf{増幅}。これがメーザー／レーザーの原理。

\textbf{(3)} 宇宙メーザーが見える＝準位反転 \(n_u/g_u>n_l/g_l\)
が実現している＝熱平衡では起こらない\textbf{非平衡なポンプ機構}（衝突ポンプ・放射ポンプ）が働いていることが逆引きできる。星形成領域・恒星質量放出・銀河中心の特殊な物理環境を示す。

\textbf{答え}：補正ゼロ＝占有が等しい。負＝準位反転（\(n_u/g_u>n_l/g_l\)）→\(\alpha_\nu<0\)＝増幅＝メーザー。観測は非平衡ポンプ機構の証拠。

\subsection*{問題 21-3　等価幅と成長曲線の三領域}\label{ex-21-3}
\addcontentsline{toc}{subsection}{問題 21-3　等価幅と成長曲線の三領域}

{[}★ 難易度：☆☆ {]}

§21.1・20.6 の等価幅 \(W_\lambda=\int\frac{I_c-I_\lambda}{I_c}d\lambda\)
と成長曲線（curve of growth）を読む。

\begin{enumerate}
\def\labelenumi{\arabic{enumi}.}
\tightlist
\item
  等価幅が「線の輪郭の詳細によらない積分量」である利点を述べよ。浅く広い線と深く狭い線が同じ
  \(W\) になりうる理由は。
\item
  成長曲線の三領域（線形 \(W\propto N\)、飽和
  \(W\propto\sqrt{\ln N}\)、減衰翼
  \(W\propto\sqrt{N}\)）が、それぞれ光学的深さ \(\tau\)
  のどんな状態に対応するか述べよ。
\item
  飽和領域の線から柱密度 \(N\)
  を精度よく決めるのが難しい理由を述べよ。組成解析で複数の強さの線を使う理由と結びつけよ。
\end{enumerate}

\textbf{関連}：\hyperref[sec-line-equivalent-width]{§21.1
等価幅}、\hyperref[sec-curve-of-growth]{§21.6
成長曲線}／組成解析は\hyperref[sec-line-applications-abundance]{§24.2}、減衰翼は\hyperref[sec-line-shape-natural]{§23.1}。

\section{模範解答（問題
21-3）}\label{ux6a21ux7bc4ux89e3ux7b54ux554fux984c-21-3}

\textbf{(1)}
等価幅は「線が連続光から引き抜いた光の総量」を波長幅で表した積分量で、装置の分解能で輪郭がぼやけても保存される。同じ総吸収量なら、浅く広い線も深く狭い線も面積が等しく同じ
\(W\) になる。観測しやすく理論と比べやすい。

\textbf{(2)} 線形領域：\(\tau<1\)（光学的に薄い）、吸収量が柱密度に比例
\(W\propto N\)。飽和領域：線芯
\(\tau\gg1\)（完全に黒）で中心はこれ以上吸収できず、\(W\) は
\(\sqrt{\ln N}\) でしか増えない。減衰翼領域：Lorentz
翼（自然幅・衝突幅）が支配し \(W\propto\sqrt{N}\)。

\textbf{(3)} 飽和領域では \(W\) が \(N\)
にほとんど鈍感（\(\sqrt{\ln N}\)）なので、\(W\) の小さな測定誤差が \(N\)
の大きな誤差になる。だから組成解析では弱い線（線形領域）・中間・強い線（減衰翼）を組み合わせ、成長曲線全体を観測点で押さえて
\(N\) を精度よく決める。

\textbf{答え}：\(W\)
は輪郭によらぬ積分量。線形(\(\tau<1\))→飽和(\(\tau\gg1\))→減衰翼。飽和域は
\(N\) に鈍感ゆえ複数の強さの線が必要。

\subsection*{問題 21-4　LTE で吸収線になる条件 ―
温度勾配と背景光}\label{ex-21-4}
\addcontentsline{toc}{subsection}{問題 21-4　LTE で吸収線になる条件 ―
温度勾配と背景光}

{[}★ 難易度：☆☆ {]}

§21.4
は線が吸収か輝線かが温度構造と背景光で決まると述べた。LTE（\(S_\nu=B_\nu(T)\)）で考える。

\begin{enumerate}
\def\labelenumi{\arabic{enumi}.}
\tightlist
\item
  恒星大気のように温度が深さとともに上がる場合、線中心と連続光がそれぞれどの深さ（温度）を見ているか述べ、なぜ吸収線になるか説明せよ。
\item
  彩層・コロナのように温度が外側で上がる、あるいは背景連続光がないガス雲では、なぜ輝線になるか述べよ。
\item
  線の Eddington-Barbier（線で \(\tau_\nu=2/3\) になる層の
  \(B_\nu\)）が、線中心と連続でなぜ違う温度を与えるか、線中心の方が
  \(\alpha_\nu\) が大きいことと結びつけて述べよ。
\end{enumerate}

\textbf{関連}：\hyperref[sec-LTE-line-formation]{§21.4 LTE
線形成}／吸収線か輝線かは\hyperref[ex-5-3]{演習 5-3}、Eddington-Barbier
は\hyperref[ex-15-1]{演習 15-1}。

\section{模範解答（問題
21-4）}\label{ux6a21ux7bc4ux89e3ux7b54ux554fux984c-21-4}

\textbf{(1)} 線中心は \(\alpha_\nu\)
が大きく光学的厚さがすぐ1に達するので\textbf{浅い（外側＝低温）層}を見る。連続光は
\(\alpha_\nu\) が小さく\textbf{深い（高温）層}まで見通せる。LTE で
\(S_\nu=B_\nu(T)\)
ゆえ、線中心は低温のプランク関数値（暗い）、連続は高温（明るい）。差し引き線中心が連続より暗く＝吸収線。

\textbf{(2)}
温度が外側で上がる（彩層・コロナ）と、線中心が見る浅い層の方が高温になり、線中心が連続より明るく＝輝線。あるいは背景連続光がないガス雲では、ガス自身の放出
\(\tau S_\nu\) だけが見え、線中心で強い輝線になる。

\textbf{(3)} Eddington-Barbier 近似では、各振動数で \(\tau_\nu=2/3\)
になる層の \(B_\nu(T)\) が出てくる。線中心は \(\alpha_\nu\) が大きいので
\(\tau_\nu=2/3\)
に達する幾何的深さが浅い（連続より外側＝別の温度の層）。よって線中心と連続が異なる温度を映し、温度勾配の向きで吸収/輝線が決まる。

\textbf{答え}：線中心＝浅い層（\(\alpha_\nu\)
大）、連続＝深い層。温度が深さで上がれば吸収線、外側で上がる/背景光なしなら輝線。Eddington-Barbier
で別温度を映す。

\section*{さらに学ぶための参考文献}\label{sec-line-formation-further}
\addcontentsline{toc}{section}{さらに学ぶための参考文献}

\markright{さらに学ぶための参考文献}

\begin{itemize}
\tightlist
\item
  Mihalas, \emph{Stellar Atmospheres} (Freeman, 2nd ed., 1978) ---
  線形成と non-LTE の古典
\item
  Hubeny \& Mihalas, \emph{Theory of Stellar Atmospheres} (Princeton,
  2014) --- 現代版、non-LTE radiative transfer の標準書
\item
  Rutten, \emph{Radiative Transfer in Stellar Atmospheres} (Lecture
  notes, Utrecht) --- 線形成の入門的講義ノート
\item
  Spitzer, \emph{Physical Processes in the Interstellar Medium} (Wiley,
  1978) --- 成長曲線（curve of growth）の標準扱い
\end{itemize}

\chapter{第22章　線スペクトルから原子物理へ}\label{ux7b2c22ux7ae0-ux7ddaux30b9ux30daux30afux30c8ux30ebux304bux3089ux539fux5b50ux7269ux7406ux3078}

\textbf{この章の主題}：第20章で観測される線スペクトルを「ファクト」として確認し、第21章で線形成の枠組み（Einstein
係数とソース関数）を整えた。本章では中心地図
\(\alpha_\nu = (\pi e^2/m_e c) f_{lu} n_l \phi(\nu-\nu_0)\) の
\textbf{\(f_{lu}\)（振動子強度）} と \textbf{\(\nu_0\)（線振動数）}
の起源を、量子力学から逆引きする。これは連続スペクトルにおける「第IV部の量子論」に対応する、線スペクトル側の
\textbf{量子力学への逆引き} である。

\section*{この章の中心地図}\label{sec-atomic-physics-map}
\addcontentsline{toc}{section}{この章の中心地図}

\markright{この章の中心地図}

\textbf{水素線スペクトルの地図 ― 二層構造}

\textbf{第一層：イコニック・アンカー（リュードベリ公式）}

\begin{equation}\protect\phantomsection\label{eq-rydberg}{
\frac{1}{\lambda} = R_\mathrm{H} \left( \frac{1}{n_1^2} - \frac{1}{n_2^2} \right)
}\end{equation}

線スペクトルの「骨格」を一目で示す。\(R_\mathrm{H}\)
は基本定数の組合せ、\(1/n^2\)
は束縛状態のエネルギー、離散性が「線が存在する理由」を語る。

\textbf{第二層：構造的マップ（Hα を標準例とした線吸収係数）}

\begin{equation}\protect\phantomsection\label{eq-line-absorption}{
\alpha_\nu = \frac{\pi e^2}{m_e c} \, f_{lu} \, n_l \, \phi(\nu - \nu_0)
}\end{equation}

{\def\LTcaptype{none} % do not increment counter
\begin{longtable}[]{@{}lll@{}}
\toprule\noalign{}
因子 & 物理的意味 & 対応する章 \\
\midrule\noalign{}
\endhead
\bottomrule\noalign{}
\endlastfoot
\(\alpha_\nu\) & 観測量（線吸収係数） & 第21章 \\
\(\pi e^2 / m_e c\) & 古典極限の線強度 & 第22章＋第VIII部 \\
\(f_{lu}\) & 振動子強度 & 第22章 \\
\(n_l\) & 下準位の占有数 & 第21章＋Part III \\
\(\phi(\nu - \nu_0)\) & 線輪郭（自然・Doppler・Voigt） & 第23章 \\
\(\nu_0\) & 線位置（リュードベリ） & 第22章 \\
\end{longtable}
}

\textbf{適用範囲（重要）}：この式の前因子 \(\pi e^2/m_e c\) と振動子強度
\(f_{lu}\) は、古典 Lorentz 振動子に対応する
\textbf{電気双極子（E1）許容線} に対するものである。21 cm
線（基底状態超微細構造の \textbf{磁気双極子 M1 遷移}）や
\textbf{禁制線（E2・M1 など）} では、この前因子・\(f_{lu}\)
をそのまま使えない。これらの線では、測定された自発放出率 \(A_{ul}\) から
Einstein
関係を介して有効遷移強度（\(B_{lu}\)）に読み替える（式~\ref{eq-line-absorption}
の構造は保ったまま、強度因子だけを差し替える）。

\begin{tcolorbox}[enhanced jigsaw, arc=.35mm, bottomrule=.15mm, breakable, colback=white, colframe=quarto-callout-tip-color-frame, left=2mm, leftrule=.75mm, opacityback=0, rightrule=.15mm, toprule=.15mm]

\textbf{地図の使い方}：第VII部の各章はこの式の異なる因子を解き明かす。プランク関数地図と並べて読めば、本書を貫く二つの逆引きが同時に把握できる。第VIII部の
QED で、両者は前因子 \(2h\nu^3/c^2\)
を介して統一される（式~\ref{eq-einstein-A} 参照）。

\end{tcolorbox}

\begin{tcolorbox}[enhanced jigsaw, arc=.35mm, bottomrule=.15mm, bottomtitle=1mm, breakable, colback=white, colbacktitle=quarto-callout-note-color!10!white, colframe=quarto-callout-note-color-frame, coltitle=black, left=2mm, leftrule=.75mm, opacityback=0, opacitybacktitle=0.6, rightrule=.15mm, title=\textcolor{quarto-callout-note-color}{\faInfo}\hspace{0.5em}{ノート}, titlerule=0mm, toprule=.15mm, toptitle=1mm]

\textbf{方針}：本章では水素原子を中心に置き、線振動数 \(\nu_0\)
と振動子強度 \(f_{lu}\)
を量子力学から逆引きする。多電子原子と分子は概要に留め、第VIII部（QED）で電磁場と物質場の量子化として統一する。

\end{tcolorbox}

\section*{この章で答える問い}\label{sec-atomic-physics-questions}
\addcontentsline{toc}{section}{この章で答える問い}

\markright{この章で答える問い}

\begin{itemize}
\tightlist
\item
  水素原子のスペクトル（Lyman・Balmer・Paschen 系列）は、なぜ Rydberg
  公式で書けるのか
\item
  振動子強度 \(f_{lu}\) は何を、どの近似で表しているのか
\item
  選択則と禁制線の物理的根拠は何か
\item
  水素原子の枠組みは、多電子原子・分子にどう拡張されるのか
\end{itemize}

\section*{到達目標}\label{sec-atomic-physics-goals}
\addcontentsline{toc}{section}{到達目標}

\markright{到達目標}

この章を読み終えると、読者は次のことができるようになる：

\begin{itemize}
\tightlist
\item
  Rydberg 公式と水素原子のエネルギー準位を量子力学から再導出できる
\item
  振動子強度 \(f_{lu}\) の物理的意味と Schrödinger
  方程式との関係を述べられる
\item
  選択則の起源と、それを破る現象（禁制線）の役割を説明できる
\end{itemize}

\begin{center}\rule{0.5\linewidth}{0.5pt}\end{center}

\section{22.1 Rydberg 公式 ―
水素線の振動数の起源}\label{sec-atomic-rydberg}

{[}本文目安：B1-B2{]}

第20章 §20.1 で見た太陽の Fraunhofer C 線（\(\lambda = 656\) nm,
H\(\alpha\)）と、それに続く水素線系列の振動数は、すべて \textbf{Rydberg
公式} で書ける：

\begin{equation}\protect\phantomsection\label{eq-rydberg-formula}{
\frac{1}{\lambda} = R_H \left(\frac{1}{n_l^2} - \frac{1}{n_u^2}\right)
}\end{equation}

ここで \(R_H \simeq 1.097 \times 10^7\) m\(^{-1}\) は Rydberg
定数、\(n_l, n_u\) は主量子数。

\begin{itemize}
\tightlist
\item
  \textbf{Lyman 系列}（\(n_l = 1\)）：紫外、\(\lambda < 122\) nm
\item
  \textbf{Balmer 系列}（\(n_l = 2\)）：可視、H\(\alpha\) = 656 nm,
  H\(\beta\) = 486 nm, \ldots{}
\item
  \textbf{Paschen 系列}（\(n_l = 3\)）：近赤外
\item
  \textbf{Brackett、Pfund 系列}（\(n_l = 4, 5\)）：赤外〜遠赤外
\end{itemize}

これは経験式として 1888 年に発見された。\textbf{なぜこの形になるか}
を逆引きするのが本節以降の主題。

\begin{tcolorbox}[enhanced jigsaw, arc=.35mm, bottomrule=.15mm, bottomtitle=1mm, breakable, colback=white, colbacktitle=quarto-callout-note-color!10!white, colframe=quarto-callout-note-color-frame, coltitle=black, left=2mm, leftrule=.75mm, opacityback=0, opacitybacktitle=0.6, rightrule=.15mm, title=\textcolor{quarto-callout-note-color}{\faInfo}\hspace{0.5em}{ノート}, titlerule=0mm, toprule=.15mm, toptitle=1mm]

\textbf{観測との対応}：観測可能なすべての水素線振動数が
\textbf{二つの整数 \(n_l, n_u\) で記述できる}
という事実は、本書の主題「観測スペクトルから量子物理を逆引きする」の典型例である。1888
年時点でこの規則性が認識されたが、それが何を意味するかは、Bohr
模型（1913）まで待たねばならなかった。

\end{tcolorbox}

\section{22.2 Bohr 模型 ― 線振動数の最初の説明}\label{sec-atomic-bohr}

{[}本文目安：B1-B2{]}

Bohr（1913）は、Rutherford の原子模型に \textbf{角運動量量子化}
を導入した：

\begin{equation}\protect\phantomsection\label{eq-bohr-quantization}{
L = n\hbar, \quad n = 1, 2, 3, \ldots
}\end{equation}

これと古典電磁気の Coulomb 力 + Newton
運動方程式を組み合わせると、エネルギー準位

\begin{equation}\protect\phantomsection\label{eq-bohr-energy}{
E_n = -\frac{1}{2} \frac{m_e e^4}{(4\pi\epsilon_0)^2 \hbar^2} \frac{1}{n^2} = -\frac{13.6\,\text{eV}}{n^2}
}\end{equation}

が出る。線振動数は \(h\nu = E_u - E_l\) から

\begin{equation}\protect\phantomsection\label{eq-bohr-rydberg}{
\frac{1}{\lambda} = \frac{m_e e^4}{(4\pi\epsilon_0)^2 \cdot 4\pi \hbar^3 c} \left(\frac{1}{n_l^2} - \frac{1}{n_u^2}\right) = R_H \left(\frac{1}{n_l^2} - \frac{1}{n_u^2}\right)
}\end{equation}

これで Rydberg 公式が量子物理から \textbf{逆引き} された。

\begin{tcolorbox}[enhanced jigsaw, arc=.35mm, bottomrule=.15mm, breakable, colback=white, colframe=quarto-callout-warning-color-frame, left=2mm, leftrule=.75mm, opacityback=0, rightrule=.15mm, toprule=.15mm]
\begin{minipage}[t]{5.5mm}
\textcolor{quarto-callout-warning-color}{\faExclamationTriangle}
\end{minipage}%
\begin{minipage}[t]{\textwidth - 5.5mm}

\textbf{注意（観測）}：Bohr 模型は \textbf{線振動数は完璧に予言する}
が、\textbf{振動子強度 \(f_{lu}\)
は予言できない}。なぜなら、模型は電子の位置を許される軌道の上に「正確に」置く描像で、波動関数を持たない。\(f_{lu}\)
は次節の Schrödinger 方程式と波動関数によって初めて計算できる

\end{minipage}%
\end{tcolorbox}

\section{22.3 Schrödinger
方程式と水素原子の波動関数}\label{sec-atomic-schrodinger}

{[}本文目安：B2-B3{]}

電子の波動的性質を取り入れた Schrödinger 方程式（1926）を Coulomb
ポテンシャル \(V(r) = -e^2/(4\pi\epsilon_0 r)\) の下で解くと、波動関数

\begin{equation}\protect\phantomsection\label{eq-hydrogen-wavefunction}{
\psi_{nlm}(r, \theta, \phi) = R_{nl}(r) Y_{lm}(\theta, \phi)
}\end{equation}

が得られる。ここで \(n, l, m\)
は主量子数・軌道角運動量量子数・磁気量子数。エネルギーは（非相対論的・スピン無視）

\begin{equation}\protect\phantomsection\label{eq-schrodinger-energy}{
E_n = -\frac{13.6\,\text{eV}}{n^2}
}\end{equation}

Bohr 模型と \textbf{同じエネルギー} だが、構造が深い：

\begin{enumerate}
\def\labelenumi{\arabic{enumi}.}
\tightlist
\item
  \textbf{縮退}：\(n\) が同じで \(l, m\)
  が違う状態はエネルギーが等しい（Coulomb ポテンシャル固有の縮退）
\item
  \textbf{波動関数}：電子の存在確率分布が空間的に広がっている
\item
  \textbf{角運動量量子化}：\(L^2 = l(l+1)\hbar^2\)（Bohr の
  \(L = n\hbar\) は不正確だった）
\end{enumerate}

\begin{tcolorbox}[enhanced jigsaw, arc=.35mm, bottomrule=.15mm, breakable, colback=white, colframe=quarto-callout-tip-color-frame, left=2mm, leftrule=.75mm, opacityback=0, rightrule=.15mm, toprule=.15mm]
\begin{minipage}[t]{5.5mm}
\textcolor{quarto-callout-tip-color}{\faLightbulb}
\end{minipage}%
\begin{minipage}[t]{\textwidth - 5.5mm}

\textbf{問い}：Bohr 模型と Schrödinger
方程式は同じエネルギーを与えるのに、なぜ Schrödinger
方程式が「正しい」量子力学なのか？

\textbf{短答}：Bohr 模型は
\textbf{エネルギーは正しいが波動関数を持たない}。次節で計算する
\textbf{振動子強度
\(f_{lu}\)}（観測される線の強度を決める）は波動関数の積分で決まる。観測される線強度を再現するには、波動関数つき
Schrödinger 方程式が必要。

\textbf{もう一歩}：相対論的補正（微細構造、Lamb シフト）も Schrödinger
方程式を Dirac 方程式 → QED
へ拡張することで導出される。これらの精密構造も観測される（次節以降）

\end{minipage}%
\end{tcolorbox}

\section{\texorpdfstring{22.4 振動子強度 \(f_{lu}\) ―
観測線強度の起源}{22.4 振動子強度 f\_\{lu\} ― 観測線強度の起源}}\label{sec-atomic-oscillator-strength}

{[}本文目安：B3{]}

中心地図 \(\alpha_\nu = (\pi e^2/m_e c) f_{lu} n_l \phi(\nu-\nu_0)\) の
\textbf{振動子強度} \(f_{lu}\) は、量子力学では \textbf{遷移行列要素}
から計算される。\(f_{lu}\) は無次元量であり、行列要素を「位置
\(\mathbf{r}\) で書く形」と「電気双極子モーメント \(\mathbf{d}\)
で書く形」の二通りがある。

\textbf{位置演算子で書く形}（\(\langle u|\mathbf{r}|l\rangle\)
は長さの次元。電荷 \(e\) を前因子に別出ししてあるため、SI・CGS
で同形）：

\begin{equation}\protect\phantomsection\label{eq-oscillator-strength}{
f_{lu} = \frac{2 m_e \omega_{ul}}{3 \hbar\, g_l} \,|\langle u | \mathbf{r} | l \rangle|^2
}\end{equation}

\textbf{電気双極子モーメントで書く形}（\(\mathbf{d} = e\mathbf{r}\)、SI
では C·m、CGS(Gauss) では esu·cm の次元。\(e^2\)
を分母に置くことで両単位系とも無次元になり、\(\epsilon_0\)
は現れない）：

\begin{equation}\protect\phantomsection\label{eq-oscillator-strength-dipole}{
f_{lu} = \frac{2 m_e \omega_{ul}}{3 \hbar\, g_l\, e^2} \,|\langle u | \mathbf{d} | l \rangle|^2
}\end{equation}

ここで \(\langle u | \mathbf{r} | l \rangle\) は
\textbf{電気双極子遷移行列要素}（位置表示）、\(\omega_{ul} = (E_u - E_l)/\hbar\)。いずれの形でも
\(f_{lu}\) は無次元である。物理的解釈：

\begin{itemize}
\tightlist
\item
  \(f_{lu}\) は古典的振動子の \textbf{強度を表す無次元量}
  で、\(\sum_u f_{lu} = N_\text{electrons}\)（\textbf{Thomas-Reiche-Kuhn
  和則}）
\item
  電気双極子の振動が「強い」遷移ほど \(f_{lu}\)
  が大きく、観測される線が強い
\item
  遷移行列要素は波動関数の重なりで決まる ― Bohr 模型では計算不可能
\end{itemize}

\begin{tcolorbox}[enhanced jigsaw, arc=.35mm, bottomrule=.15mm, bottomtitle=1mm, breakable, colback=white, colbacktitle=quarto-callout-note-color!10!white, colframe=quarto-callout-note-color-frame, coltitle=black, left=2mm, leftrule=.75mm, opacityback=0, opacitybacktitle=0.6, rightrule=.15mm, title=\textcolor{quarto-callout-note-color}{\faInfo}\hspace{0.5em}{ノート}, titlerule=0mm, toprule=.15mm, toptitle=1mm]

\textbf{用語}：\(\pi e^2 / m_e c\) という前因子は \textbf{古典 Lorentz
振動子の積分吸収断面積}（次元は cm\(^2 \cdot\)
Hz）に対応する量で、これは \textbf{CGS(Gauss) 形} である。\textbf{SI 形}
では \(\pi e^2/(4\pi\epsilon_0 m_e c) = e^2/(4\epsilon_0 m_e c)\)
となる。振動数積分した吸収係数は
\(\int \sigma\, d\nu = (\pi e^2/m_e c)\, f_{lu}\)（CGS）の形で現れる。線輪郭
\(\phi(\nu-\nu_0)\) を掛けて積分すると規格化条件
\(\int \phi\, d\nu = 1\) により全積分が \((\pi e^2/m_e c)\, f_{lu}\)
に戻る ― \textbf{古典 + 量子論の橋} の典型例である（Thomson 散乱断面積
\(\sigma_T = 8\pi r_e^2/3\) とは別物なので混同しないこと）。

\textbf{観測との対応}：\(f_{lu}\)
の値は理論計算・実験室で精密に求められており、原子物理の標準データ（NIST
Atomic Spectra Database 等）として整備されている。観測線強度 →
柱密度の換算（第21章 成長曲線）にはこのデータが必須である。

\end{tcolorbox}

\section{\texorpdfstring{22.4b \(\pi e^2/m_e c\) の逆引き ― 古典 Lorentz
振動子から}{22.4b \textbackslash pi e\^{}2/m\_e c の逆引き ― 古典 Lorentz 振動子から}}\label{sec-atomic-lorentz-prefactor}

{[}本文目安：B3／導出：M1{]}

中心地図 \(\alpha_\nu = (\pi e^2/m_e c) f_{lu} n_l \phi(\nu-\nu_0)\) の
\textbf{前因子 \(\pi e^2/m_e c\)}
は、第VII部の中心地図のうち、まだ逆引きを明示していなかった因子である。本節でその起源を、古典
Lorentz 振動子（damped harmonic oscillator）モデルから示し、Einstein
係数との関係を通じて量子論との橋を完成させる。

\subsection{古典 Lorentz
振動子の吸収断面積}\label{ux53e4ux5178-lorentz-ux632fux52d5ux5b50ux306eux5438ux53ceux65adux9762ux7a4d}

電場 \(\mathbf{E}(t) = \mathbf{E}_0 e^{-i\omega t}\)
に弱く減衰する束縛電子（質量 \(m_e\)、電荷 \(-e\)、固有角振動数
\(\omega_0\)、減衰率 \(\gamma\)）の運動方程式：

\begin{equation}\protect\phantomsection\label{eq-lorentz-oscillator}{
\ddot{\mathbf{r}} + \gamma \dot{\mathbf{r}} + \omega_0^2 \mathbf{r} = -\frac{e}{m_e}\mathbf{E}(t)
}\end{equation}

定常解の振幅から、振動電子による双極子モーメント \(\mathbf{p}(t)\)
を求め、これに対する
\textbf{吸収断面積}（時間平均パワー吸収率を入射エネルギーフラックスで割ったもの）を計算すると、振動数積分された値は

\begin{equation}\protect\phantomsection\label{eq-classical-integrated-cross-section}{
\int_0^\infty \sigma_\text{cl}(\nu)\, d\nu = \frac{\pi e^2}{m_e c} \quad (\text{CGS(Gauss) 形})
}\end{equation}

である（\textbf{SI 形} では
\(\dfrac{\pi e^2}{4\pi\epsilon_0 m_e c} = \dfrac{e^2}{4\epsilon_0 m_e c}\)）。ここで前因子
\(\pi e^2/m_e c\) は \textbf{減衰率 \(\gamma\)
にも線輪郭のピークの位置にもよらない普遍量}
であり、古典電磁気の力学パラメータのみで決まる。

\subsection{\texorpdfstring{量子論への橋 ― 振動子強度 \(f_{lu}\)
の登場}{量子論への橋 ― 振動子強度 f\_\{lu\} の登場}}\label{ux91cfux5b50ux8ad6ux3078ux306eux6a4b-ux632fux52d5ux5b50ux5f37ux5ea6-f_lu-ux306eux767bux5834}

実際の原子は古典振動子ではなく、量子的な多準位系である。各遷移
\(l \to u\) の \textbf{実効的な吸収強度}
は、古典的振動子の振動数積分断面積に \textbf{無次元の補正係数}
\(f_{lu}\)（振動子強度、第21.4節）を掛けたものとして書ける：

\begin{equation}\protect\phantomsection\label{eq-quantum-integrated-cross-section}{
\int_0^\infty \sigma_{lu}(\nu)\, d\nu = \frac{\pi e^2}{m_e c} \cdot f_{lu} \quad (\text{CGS(Gauss) 形})
}\end{equation}

\(f_{lu}\) は量子力学の遷移行列要素 式~\ref{eq-oscillator-strength}
から計算される（SI 形では前因子を \(e^2/(4\epsilon_0 m_e c)\)
に置き換える）。Thomas-Reiche-Kuhn 和則
\(\sum_u f_{lu} = N_\text{electrons}\)
は、「すべての遷移を足し合わせると、古典的に予想される総吸収量に戻る」という意味で、古典-量子対応を定量的に保証する。

線輪郭 \(\phi(\nu-\nu_0)\)（規格化
\(\int \phi\, d\nu = 1\)）を使って、吸収断面積を振動数依存の形に展開する：

\begin{equation}\protect\phantomsection\label{eq-cross-section-line}{
\sigma_{lu}(\nu) = \frac{\pi e^2}{m_e c} f_{lu}\, \phi(\nu - \nu_0) \quad (\text{CGS(Gauss) 形})
}\end{equation}

これを \textbf{柱密度 × 断面積 = 吸収係数} の関係に代入して、中心地図

\begin{equation}\protect\phantomsection\label{eq-line-absorption-coefficient-complete}{
\boxed{
\alpha_\nu = \frac{\pi e^2}{m_e c}\, f_{lu}\, n_l\, \phi(\nu - \nu_0)
}\quad (\text{CGS(Gauss) 形})
}\end{equation}

を再現する（第VII部冒頭で経験式として受け入れていたこの式が、いま完全に逆引きされた）。SI
形では前因子を \(e^2/(4\epsilon_0 m_e c)\) に置き換えればよい。

\subsection{\texorpdfstring{Einstein \(B\)
係数との換算}{Einstein B 係数との換算}}\label{einstein-b-ux4fc2ux6570ux3068ux306eux63dbux7b97}

第21章 §21.2 で導入した Einstein 関係
\(\alpha_\nu = (h\nu_0/4\pi) n_l B_{lu} \phi(\nu-\nu_0)\) と本節の
式~\ref{eq-line-absorption-coefficient-complete} を比べると、両者は
\textbf{同一の物理量を別の言葉で書いたもの} であり、

\begin{equation}\protect\phantomsection\label{eq-B-f-conversion}{
B_{lu} = \frac{4\pi^2 e^2}{m_e c h\nu_0}\, f_{lu} \quad (\text{CGS(Gauss) 形})
}\end{equation}

の換算で結ばれる（\(J_\nu\) 規約、§11.3。SI 形では
\(e^2 \to e^2/(4\pi\epsilon_0)\)）。さらに Einstein 関係
\(A_{ul} = (2h\nu^3/c^2) B_{ul} = (2h\nu^3/c^2)(g_l/g_u) B_{lu}\)
と組合せて、自発放出率 \(A_{ul}\) も振動子強度から得られる：

\begin{equation}\protect\phantomsection\label{eq-A-f-conversion}{
A_{ul} = \frac{8\pi^2 e^2 \nu_0^2}{m_e c^3}\, \frac{g_l}{g_u}\, f_{lu} \quad (\text{CGS(Gauss) 形})
}\end{equation}

である（SI 形では \(e^2 \to e^2/(4\pi\epsilon_0)\)）。この \(A_{ul}\) と
\(f_{lu}\) の関係は電気双極子（E1）遷移を前提とするが、\(A_{ul}\)
を一般の遷移（M1・E2 等）の測定値で与えれば、式~\ref{eq-B-f-conversion}
を逆に解いて実効的な \(B_{lu}\) を得る読み替えが可能である（21 cm
線の扱い、§20.4）。

\begin{tcolorbox}[enhanced jigsaw, arc=.35mm, bottomrule=.15mm, bottomtitle=1mm, breakable, colback=white, colbacktitle=quarto-callout-note-color!10!white, colframe=quarto-callout-note-color-frame, coltitle=black, left=2mm, leftrule=.75mm, opacityback=0, opacitybacktitle=0.6, rightrule=.15mm, title=\textcolor{quarto-callout-note-color}{\faInfo}\hspace{0.5em}{ノート}, titlerule=0mm, toprule=.15mm, toptitle=1mm]

\textbf{第VII部中心地図の完成}：これで線スペクトル中心地図の全因子（\(\pi e^2/m_e c\),
\(f_{lu}\), \(n_l\), \(\phi\), \(\nu_0\)）が逆引きされた：

{\def\LTcaptype{none} % do not increment counter
\begin{longtable}[]{@{}
  >{\raggedright\arraybackslash}p{(\linewidth - 4\tabcolsep) * \real{0.3333}}
  >{\raggedright\arraybackslash}p{(\linewidth - 4\tabcolsep) * \real{0.3333}}
  >{\raggedright\arraybackslash}p{(\linewidth - 4\tabcolsep) * \real{0.3333}}@{}}
\toprule\noalign{}
\begin{minipage}[b]{\linewidth}\raggedright
因子
\end{minipage} & \begin{minipage}[b]{\linewidth}\raggedright
起源
\end{minipage} & \begin{minipage}[b]{\linewidth}\raggedright
章
\end{minipage} \\
\midrule\noalign{}
\endhead
\bottomrule\noalign{}
\endlastfoot
\(\pi e^2 / m_e c\) & 古典 Lorentz 振動子の積分吸収断面積 &
\textbf{本節（§22.4b）} \\
\(f_{lu}\) & 遷移行列要素
\(|\langle u|\mathbf{r}|l\rangle|^2\)（量子論的補正） & §22.4 \\
\(n_l\) & LTE（Boltzmann/Saha）または non-LTE 統計平衡 & §21.4-§21.5 \\
\(\phi(\nu - \nu_0)\) & 自然幅・Doppler・衝突幅・Voigt & 第23章 \\
\(\nu_0\) & 原子・分子のエネルギー準位 & §22.1-§22.3 \\
\end{longtable}
}

連続スペクトル側（第14章で逆引き完了）と並んで、\textbf{第VII部の中心地図も完全に逆引きされた}。第VIII部の
QED で両中心地図が統一されるとき、\(\pi e^2/m_e c\) と \(2h\nu^3/c^2\)
の両前因子が、同じ電磁場--物質相互作用 Hamiltonian
から生じることが明らかになる

\end{tcolorbox}

\section{22.5 選択則と禁制線 ―
なぜ「許されない」遷移が見えるか}\label{sec-atomic-selection-rules}

{[}本文目安：B3{]}

電気双極子遷移は、波動関数のパリティと角運動量の選択則を持つ：

\begin{equation}\protect\phantomsection\label{eq-selection-rules-E1}{
\Delta l = \pm 1, \quad \Delta m = 0, \pm 1
}\end{equation}

これらを満たさない遷移は、電気双極子では起こらない ―
「\textbf{禁制}」（\textbf{forbidden}）。

しかし、より高次の遷移は弱いながら可能：

\begin{itemize}
\tightlist
\item
  \textbf{電気四重極遷移}（\textbf{E2}）：\(f \sim \alpha^2 f_{E1} \sim 10^{-4}\)
  倍
\item
  \textbf{磁気双極子遷移}（\textbf{M1}）：\(f \sim \alpha^2 f_{E1} \sim 10^{-4}\)
  倍
\end{itemize}

これらの弱い遷移は、地上の実験室では \textbf{衝突脱励起}
で先に消されて観測できないが、\textbf{極低密度の宇宙環境}
では時間が許されて自発放出が起こる。

第20章 §20.3 で触れた \textbf{禁制線}（{[}O III{]} \$\lambda\$5007
等）の正体はこれ。

\begin{tcolorbox}[enhanced jigsaw, arc=.35mm, bottomrule=.15mm, breakable, colback=white, colframe=quarto-callout-tip-color-frame, left=2mm, leftrule=.75mm, opacityback=0, rightrule=.15mm, toprule=.15mm]
\begin{minipage}[t]{5.5mm}
\textcolor{quarto-callout-tip-color}{\faLightbulb}
\end{minipage}%
\begin{minipage}[t]{\textwidth - 5.5mm}

\textbf{問い}：{[}O III{]}
\$\lambda\(5007 は「禁制」と呼ばれるが、HII 領域や惑星状星雲では H\)\beta\$
よりも強く観測されることがある。なぜか？

\textbf{短答}：HII 領域の電子密度（\(n_e \sim 10^2\)〜\(10^4\)
cm\(^{-3}\)）では衝突脱励起時間が長く、\(A_{ul} \sim 10^{-2}\)
s\(^{-1}\) の禁制線でも自発放出が間に合う。一方、{[}O III{]}
のエネルギー準位構造は
\textbf{衝突励起断面積が大きく}、励起される割合も高い。両者の組合せで、許容線と同程度の強度が出る。

\textbf{もう一歩}：{[}O III{]} \$\lambda\$5007 の強度は、星雲の
\textbf{電子温度} に強く依存し、§24.3
の温度診断に使われる。「禁制」という言葉は、地上実験室での測りにくさを反映するが、宇宙物理では主役級の道具

\end{minipage}%
\end{tcolorbox}

\section{22.6 多電子原子と分子 ―
水素原子の枠組みの拡張}\label{sec-atomic-multi-electron}

{[}本文目安：B3-B4{]}

水素以外の原子では、複数の電子が同時に存在するため、\(\psi\)
は単一粒子波動関数の積では書けない（Pauli
の排他律と電子間相互作用）。代わりに：

\begin{itemize}
\tightlist
\item
  \textbf{Hartree-Fock
  近似}：各電子の波動関数を、他電子の平均場の中で解く
\item
  \textbf{項記号}（\textbf{term symbol}）：原子の状態を \(^{2S+1}L_J\)
  で表す（例：太陽 D 線の上準位は \(^2P_{3/2}\) と \(^2P_{1/2}\)）
\item
  \textbf{微細構造}：スピン-軌道結合による準位分裂、線分裂の起源
\item
  \textbf{超微細構造}：核スピンとの結合、21 cm 線の起源（§20.4）
\end{itemize}

分子では、\textbf{回転・振動準位} が加わる：

\begin{itemize}
\tightlist
\item
  \textbf{回転準位}：\(E_J = \hbar^2 J(J+1) / 2I\)、\(I\)
  は慣性モーメント。電波・サブミリ波
\item
  \textbf{振動準位}：\(E_v = \hbar \omega (v + 1/2)\)、\(\omega\)
  は分子振動振動数。中間赤外
\item
  \textbf{電子準位}：原子と同類。紫外〜可視
\end{itemize}

これらの組合せで、分子線スペクトルの \textbf{巨大なライン群（rotational
band 等）} が生まれる。

\begin{tcolorbox}[enhanced jigsaw, arc=.35mm, bottomrule=.15mm, breakable, colback=white, colframe=quarto-callout-warning-color-frame, left=2mm, leftrule=.75mm, opacityback=0, rightrule=.15mm, toprule=.15mm]
\begin{minipage}[t]{5.5mm}
\textcolor{quarto-callout-warning-color}{\faExclamationTriangle}
\end{minipage}%
\begin{minipage}[t]{\textwidth - 5.5mm}

\textbf{注意（観測）}：分子線スペクトルは原子線より遥かに複雑だが、回転定数
\(B\) と振動定数 \(\omega_e\) の \textbf{二つの定数}
で大半が記述できる。NIST Molecular Spectroscopy Database
等で公開されており、観測者は分子線同定とパラメータ決定にこれを利用する

\end{minipage}%
\end{tcolorbox}

\section*{22.7 使えるようになった道具}\label{sec-atomic-tools}
\addcontentsline{toc}{section}{22.7 使えるようになった道具}

\markright{22.7 使えるようになった道具}

\begin{tcolorbox}[enhanced jigsaw, arc=.35mm, bottomrule=.15mm, bottomtitle=1mm, breakable, colback=white, colbacktitle=quarto-callout-note-color!10!white, colframe=quarto-callout-note-color-frame, coltitle=black, left=2mm, leftrule=.75mm, opacityback=0, opacitybacktitle=0.6, rightrule=.15mm, title=\textcolor{quarto-callout-note-color}{\faInfo}\hspace{0.5em}{ノート}, titlerule=0mm, toprule=.15mm, toptitle=1mm]

\begin{itemize}
\tightlist
\item
  Rydberg 公式とその Bohr 模型・Schrödinger 方程式からの逆引き導出
\item
  振動子強度 \(f_{lu}\) の量子力学的意味（遷移行列要素）
\item
  選択則と禁制線の起源、極低密度宇宙環境での観測可能性
\item
  多電子原子（項記号、微細構造）・分子（回転・振動準位）への拡張
\item
  中心地図 \(\alpha_\nu\) の \(f_{lu}\) と \(\nu_0\) 因子の完全な逆引き
\end{itemize}

\end{tcolorbox}

\begin{center}\rule{0.5\linewidth}{0.5pt}\end{center}

\section*{この章で何がわかったか}\label{sec-atomic-physics-summary}
\addcontentsline{toc}{section}{この章で何がわかったか}

\markright{この章で何がわかったか}

\textbf{中心地図に戻る}

本章で、中心地図
\(\alpha_\nu = (\pi e^2/m_e c) f_{lu} n_l \phi(\nu-\nu_0)\) の
\textbf{線振動数 \(\nu_0\)} と \textbf{振動子強度 \(f_{lu}\)}
が量子力学から逆引きされた：

\begin{itemize}
\tightlist
\item
  \(\nu_0\) は原子・分子のエネルギー準位差（Rydberg
  公式、項記号、回転・振動準位）
\item
  \(f_{lu}\) は遷移行列要素 \(|\langle u|\mathbf{r}|l\rangle|^2\)
  から計算される無次元量
\item
  選択則と禁制線が、観測される線スペクトルの構造（強い線・弱い線・宇宙環境固有の線）を決める
\end{itemize}

次章では中心地図の最後の因子 ― \textbf{線輪郭 \(\phi(\nu-\nu_0)\)} ―
の物理を探る。

\section*{演習問題}\label{sec-atomic-physics-exercises}
\addcontentsline{toc}{section}{演習問題}

\markright{演習問題}

\textbf{模範解答}
は折りたたみを展開して確認できる（オンライン版）。まず自力で解いてから開くこと。

\subsection*{問題 22-1　Rydberg 公式と Hα の波長}\label{ex-22-1}
\addcontentsline{toc}{subsection}{問題 22-1　Rydberg 公式と Hα の波長}

{[}★ 難易度：☆☆ {]}

§22.1・21.2 の Rydberg 公式
\(\dfrac{1}{\lambda}=R_H\left(\dfrac{1}{n_l^2}-\dfrac{1}{n_u^2}\right)\)、\(R_H=1.097\times10^7\)
m\(^{-1}\)、Bohr エネルギー \(E_n=-13.6/n^2\) eV。

\begin{enumerate}
\def\labelenumi{\arabic{enumi}.}
\tightlist
\item
  \(h\nu=E_u-E_l\) から Rydberg 公式の形が出ることを示せ。
\item
  Balmer 系列の H\(\alpha\)（\(n_l=2\to n_u=3\)）の波長を計算し、656 nm
  を確かめよ。
\item
  Lyman
  系列の系列極限（\(n_u\to\infty\)、\(n_l=1\)）の波長を求め、これが水素の電離エネルギー
  13.6 eV に対応することを確かめよ。
\end{enumerate}

\textbf{関連}：\hyperref[sec-atomic-rydberg]{§22.1 Rydberg
公式}、\hyperref[sec-atomic-bohr]{§22.2 Bohr
模型}／準位の量子力学的起源は\hyperref[sec-quantum-foundations-states]{§12.1}、Boltzmann
占有は\hyperref[ex-7-4]{演習 7-4}。

\section{模範解答（問題
22-1）}\label{ux6a21ux7bc4ux89e3ux7b54ux554fux984c-22-1}

\textbf{(1)}
\(h\nu=E_u-E_l=-\dfrac{13.6}{n_u^2}-\left(-\dfrac{13.6}{n_l^2}\right)=13.6\,\text{eV}\left(\dfrac{1}{n_l^2}-\dfrac{1}{n_u^2}\right)\)。\(\nu=c/\lambda\)
より
\(\dfrac{1}{\lambda}=\dfrac{13.6\,\text{eV}}{hc}\left(\dfrac{1}{n_l^2}-\dfrac{1}{n_u^2}\right)=R_H\left(\dfrac{1}{n_l^2}-\dfrac{1}{n_u^2}\right)\)。

\textbf{(2)}
\(\dfrac{1}{\lambda}=1.097\times10^7\left(\dfrac14-\dfrac19\right)=1.097\times10^7\times0.1389=1.524\times10^6\)
m\(^{-1}\)。\(\lambda=6.56\times10^{-7}\) m \(=656\) nm。✓（太陽
Fraunhofer C 線）

\textbf{(3)}
\(n_u\to\infty\)、\(n_l=1\)：\(\dfrac{1}{\lambda}=R_H\times1=1.097\times10^7\)
m\(^{-1}\)、\(\lambda=91.2\) nm。対応エネルギー
\(E=hc/\lambda=1.986\times10^{-25}/9.12\times10^{-8}=2.18\times10^{-18}\)
J \(=13.6\) eV。これは基底状態（\(n=1\)）からの電離エネルギーに一致。

\textbf{答え}：\(h\nu=E_u-E_l\) から Rydberg 公式。H\(\alpha\)＝656
nm。Lyman 極限 91.2 nm＝13.6 eV（電離エネルギー）。

\subsection*{問題 22-2　Einstein 係数と振動子強度の換算}\label{ex-22-2}
\addcontentsline{toc}{subsection}{問題 22-2　Einstein
係数と振動子強度の換算}

{[}★ 難易度：☆☆ {]}

§22.4b は中心地図の前因子 \(\pi e^2/m_ec\) と Einstein
係数を結ぶ換算を与えた。導出を埋める。中心地図
\(\alpha_\nu=\dfrac{\pi e^2}{m_ec}f_{lu}n_l\phi\) と Einstein 形
\(\alpha_\nu=\dfrac{h\nu_0}{4\pi}n_lB_{lu}\phi\)（誘導放出補正を当面無視）。

\begin{enumerate}
\def\labelenumi{\arabic{enumi}.}
\tightlist
\item
  両式を等置して
  \(B_{lu}=\dfrac{4\pi^2e^2}{m_ec\,h\nu_0}f_{lu}\)（\(J_\nu\)
  規約）を導け。
\item
  Einstein 関係
  \(A_{ul}=\dfrac{2h\nu_0^3}{c^2}B_{ul}=\dfrac{2h\nu_0^3}{c^2}\dfrac{g_l}{g_u}B_{lu}\)
  と組合せ、\(A_{ul}=\dfrac{8\pi^2e^2\nu_0^2}{m_ec^3}\dfrac{g_l}{g_u}f_{lu}\)
  を導け。
\item
  \(\pi e^2/m_ec\) が「古典 Lorentz
  振動子の積分吸収断面積」であり、\(f_{lu}\)
  がその量子論的補正（Thomas-Reiche-Kuhn 和則
  \(\sum_u f_{lu}=N_\mathrm{electrons}\)）であることの意味を一言で述べよ。
\end{enumerate}

\textbf{関連}：\hyperref[sec-atomic-lorentz-prefactor]{§22.4b
前因子の逆引き}、\hyperref[sec-atomic-oscillator-strength]{§22.4
振動子強度}／Einstein 関係は\hyperref[ex-11-1]{演習 11-1}、QED からの
\(A_{ul}\) は\hyperref[sec-qed-minimum-spontaneous]{§26.2}。

\section{模範解答（問題
22-2）}\label{ux6a21ux7bc4ux89e3ux7b54ux554fux984c-22-2}

\textbf{(1)} 両式の \(n_l\phi\)
を消して係数を等置：\(\dfrac{\pi e^2}{m_ec}f_{lu}=\dfrac{h\nu_0}{4\pi}B_{lu}\)。\(B_{lu}\)
について解く：

\[B_{lu}=\frac{4\pi}{h\nu_0}\cdot\frac{\pi e^2}{m_ec}f_{lu}=\frac{4\pi^2e^2}{m_ec\,h\nu_0}f_{lu}.\]

\textbf{(2)} \(B_{ul}=（g_l/g_u)B_{lu}\) を
\(A_{ul}=\dfrac{2h\nu_0^3}{c^2}B_{ul}\) に代入：

\[A_{ul}=\frac{2h\nu_0^3}{c^2}\frac{g_l}{g_u}\frac{4\pi^2e^2}{m_ec\,h\nu_0}f_{lu}=\frac{8\pi^2e^2\nu_0^2}{m_ec^3}\frac{g_l}{g_u}f_{lu}.\]

（\(h\) が消え、\(\nu_0^3/\nu_0=\nu_0^2\)。）

\textbf{(3)} \(\pi e^2/m_ec\) は古典の減衰調和振動子（Lorentz
振動子）が与える振動数積分吸収断面積で、減衰率にもピーク位置にもよらない普遍量。実際の原子はこれに無次元の量子補正
\(f_{lu}\) を掛けたもので、全遷移の和
\(\sum_u f_{lu}=N_\mathrm{electrons}\)（TRK
和則）が「すべて足すと古典総吸収量に戻る」ことを保証する。古典-量子の橋。

\textbf{答え}：\(B_{lu}=\dfrac{4\pi^2e^2}{m_ech\nu_0}f_{lu}\)、\(A_{ul}=\dfrac{8\pi^2e^2\nu_0^2}{m_ec^3}\dfrac{g_l}{g_u}f_{lu}\)。\(\pi e^2/m_ec\)＝古典積分断面積、\(f_{lu}\)＝量子補正（TRK
和則）。

\subsection*{問題 22-3　選択則と禁制線の次数}\label{ex-22-3}
\addcontentsline{toc}{subsection}{問題 22-3　選択則と禁制線の次数}

{[}★ 難易度：☆☆ {]}

§22.5 の電気双極子選択則 \(\Delta l=\pm1\)、\(\Delta m=0,\pm1\)
と禁制遷移を読む。

\begin{enumerate}
\def\labelenumi{\arabic{enumi}.}
\tightlist
\item
  電気双極子（E1）選択則を満たさない遷移が「禁制」と呼ばれる理由を、遷移行列要素
  \(\langle u|\mathbf r|l\rangle\) の観点から述べよ。
\item
  磁気双極子（M1）・電気四重極（E2）遷移の強度が E1 の約
  \(\alpha^2\sim10^{-4}\)
  倍であること（\(\alpha\simeq1/137\)）を確かめ、\(\alpha^2\)
  がどこから来るか一言で述べよ。
\item
  禁制線が地上では見えにくいのに星間で見える理由を、演習 20-2
  と結びつけて述べよ。21 cm 線が M1 禁制遷移であることにも触れよ。
\end{enumerate}

\textbf{関連}：\hyperref[sec-atomic-selection-rules]{§22.5
選択則と禁制線}／QED
からの選択則は\hyperref[sec-qed-minimum-selection]{§26.5}、禁制線と低密度は\hyperref[ex-20-2]{演習
20-2}、密度診断は\hyperref[sec-line-applications-diagnostics]{§24.3}。

\section{模範解答（問題
22-3）}\label{ux6a21ux7bc4ux89e3ux7b54ux554fux984c-22-3}

\textbf{(1)} E1 遷移強度は電気双極子行列要素
\(\langle u|\mathbf r|l\rangle\)
で決まる。パリティと角運動量の保存から、\(\Delta l=\pm1\)
等を満たさない遷移ではこの行列要素がゼロになり、最低次（E1）では遷移が起こらない。これが「禁制」。

\textbf{(2)} M1・E2 は次数が一つ高い過程で、典型的に振幅が微細構造定数
\(\alpha\) 程度小さく、強度（振幅の2乗）が
\(\alpha^2\simeq(1/137)^2\simeq5.3\times10^{-5}\sim10^{-4}\)
倍。\(\alpha\) は電磁相互作用の結合定数で、より高次の多重極展開の各段で
\(\alpha\) のべきが効く。

\textbf{(3)} 禁制線は \(A_{ul}\)
が極端に小さく長寿命。地上の高密度では自発放出より衝突脱励起が先に起こり消える（\hyperref[ex-20-2]{演習
20-2}）。星間の極低密度では衝突が稀で自発放出が間に合い観測できる。21 cm
線は水素基底状態の超微細構造の M1 禁制遷移で、\(A_{ul}\sim10^{-15}\)
s\(^{-1}\) と極小だが膨大な中性水素のおかげで観測される。

\textbf{答え}：E1 選択則違反＝行列要素ゼロ＝禁制。M1/E2 は
\(\alpha^2\sim10^{-4}\) 倍。低密度の星間でのみ観測（21 cm も M1 禁制）。

\subsection*{問題 22-4　なぜ Bohr
模型では線強度が出ないか}\label{ex-22-4}
\addcontentsline{toc}{subsection}{問題 22-4　なぜ Bohr
模型では線強度が出ないか}

{[}★ 難易度：☆ {]}

§22.2・21.3 は Bohr 模型と Schrödinger
方程式が同じエネルギーを与えるのに、後者が「正しい」量子力学だと述べた。

\begin{enumerate}
\def\labelenumi{\arabic{enumi}.}
\tightlist
\item
  Bohr 模型と Schrödinger
  方程式が一致する点（エネルギー準位）と、異なる点（波動関数・縮退・角運動量）を述べよ。
\item
  振動子強度 \(f_{lu}\propto|\langle u|\mathbf r|l\rangle|^2\) がなぜ
  Bohr 模型では計算できないか述べよ。
\item
  観測される線の\textbf{振動数}は Bohr
  模型で説明できるが、線の\textbf{強度}には Schrödinger
  方程式（波動関数）が必要、という構図を、中心地図 \(\alpha_\nu\) の
  \(\nu_0\) と \(f_{lu}\) に対応づけて述べよ。
\end{enumerate}

\textbf{関連}：\hyperref[sec-atomic-schrodinger]{§22.3 Schrödinger
方程式}／固有値問題は\hyperref[sec-quantum-foundations-states]{§12.1}、振動子強度の換算は\hyperref[ex-22-2]{演習
22-2}。

\section{模範解答（問題
22-4）}\label{ux6a21ux7bc4ux89e3ux7b54ux554fux984c-22-4}

\textbf{(1)} 一致：エネルギー準位 \(E_n=-13.6/n^2\)
eV（線振動数を正しく与える）。異なる：Bohr
は電子を許される軌道上に「正確に」置く描像で波動関数を持たず、角運動量を
\(L=n\hbar\) とする（Schrödinger は \(L^2=l(l+1)\hbar^2\)）。Schrödinger
は同じ \(n\) で \(l,m\)
の異なる縮退状態と、空間に広がる存在確率分布（波動関数）を持つ。

\textbf{(2)} \(f_{lu}\) は上下準位の波動関数の重なり
\(\langle u|\mathbf r|l\rangle=\int\psi_u^*\,\mathbf r\,\psi_l\,d^3r\)
から計算される。Bohr 模型は波動関数を持たないので、この積分を実行できず
\(f_{lu}\) を予言できない。

\textbf{(3)} 中心地図
\(\alpha_\nu=(\pi e^2/m_ec)f_{lu}n_l\phi(\nu-\nu_0)\)
で、\(\nu_0\)（線振動数＝準位差）は Bohr
模型でも出るが、\(f_{lu}\)（線強度を決める振動子強度）は波動関数の行列要素ゆえ
Schrödinger
方程式が要る。線の「位置」と「強さ」で必要な理論の深さが違う。

\textbf{答え}：エネルギー（\(\nu_0\)）は一致するが、Bohr
は波動関数を持たず \(f_{lu}=|\langle u|\mathbf r|l\rangle|^2\)
を計算できない。線強度には Schrödinger が必要。

\section*{さらに学ぶための参考文献}\label{sec-atomic-physics-further}
\addcontentsline{toc}{section}{さらに学ぶための参考文献}

\markright{さらに学ぶための参考文献}

\begin{itemize}
\tightlist
\item
  Bransden \& Joachain, \emph{Physics of Atoms and Molecules} (Pearson,
  2nd ed., 2003) --- 原子物理の標準教科書
\item
  Sobelman, \emph{Atomic Spectra and Radiative Transitions} (Springer,
  2nd ed., 1992) --- 原子線と振動子強度
\item
  Herzberg, \emph{Molecular Spectra and Molecular Structure} (Krieger,
  vol.~I-III, 1989-1991) --- 分子分光の標準書
\item
  NIST Atomic Spectra Database (\url{https://physics.nist.gov/asd}) ---
  原子線データの定本
\end{itemize}

\chapter{第23章　線輪郭（line
profile）の物理}\label{ux7b2c23ux7ae0-ux7ddaux8f2aux90edline-profileux306eux7269ux7406}

\textbf{この章の主題}：中心地図
\(\alpha_\nu = (\pi e^2/m_e c) f_{lu} n_l \phi(\nu-\nu_0)\) の最後の因子
― \textbf{線輪郭関数 \(\phi(\nu-\nu_0)\)}（line profile、line shape
ともいう）―
を逆引きする。観測される線は決して数学的なデルタ関数ではなく、有限の幅を持つ。\textbf{自然幅・Doppler
拡がり・衝突拡がり}、そしてそれらの合成である \textbf{Voigt
プロファイル} の物理を整理し、観測される
\textbf{線輪郭から物理条件}（温度・密度・速度場・磁場）を逆引きする道具を整える。

\section*{この章の中心地図}\label{sec-line-shape-map}
\addcontentsline{toc}{section}{この章の中心地図}

\markright{この章の中心地図}

\textbf{水素線スペクトルの地図 ― 二層構造}

\textbf{第一層：イコニック・アンカー（リュードベリ公式）}

\begin{equation}\protect\phantomsection\label{eq-rydberg}{
\frac{1}{\lambda} = R_\mathrm{H} \left( \frac{1}{n_1^2} - \frac{1}{n_2^2} \right)
}\end{equation}

線スペクトルの「骨格」を一目で示す。\(R_\mathrm{H}\)
は基本定数の組合せ、\(1/n^2\)
は束縛状態のエネルギー、離散性が「線が存在する理由」を語る。

\textbf{第二層：構造的マップ（Hα を標準例とした線吸収係数）}

\begin{equation}\protect\phantomsection\label{eq-line-absorption}{
\alpha_\nu = \frac{\pi e^2}{m_e c} \, f_{lu} \, n_l \, \phi(\nu - \nu_0)
}\end{equation}

{\def\LTcaptype{none} % do not increment counter
\begin{longtable}[]{@{}lll@{}}
\toprule\noalign{}
因子 & 物理的意味 & 対応する章 \\
\midrule\noalign{}
\endhead
\bottomrule\noalign{}
\endlastfoot
\(\alpha_\nu\) & 観測量（線吸収係数） & 第21章 \\
\(\pi e^2 / m_e c\) & 古典極限の線強度 & 第22章＋第VIII部 \\
\(f_{lu}\) & 振動子強度 & 第22章 \\
\(n_l\) & 下準位の占有数 & 第21章＋Part III \\
\(\phi(\nu - \nu_0)\) & 線輪郭（自然・Doppler・Voigt） & 第23章 \\
\(\nu_0\) & 線位置（リュードベリ） & 第22章 \\
\end{longtable}
}

\textbf{適用範囲（重要）}：この式の前因子 \(\pi e^2/m_e c\) と振動子強度
\(f_{lu}\) は、古典 Lorentz 振動子に対応する
\textbf{電気双極子（E1）許容線} に対するものである。21 cm
線（基底状態超微細構造の \textbf{磁気双極子 M1 遷移}）や
\textbf{禁制線（E2・M1 など）} では、この前因子・\(f_{lu}\)
をそのまま使えない。これらの線では、測定された自発放出率 \(A_{ul}\) から
Einstein
関係を介して有効遷移強度（\(B_{lu}\)）に読み替える（式~\ref{eq-line-absorption}
の構造は保ったまま、強度因子だけを差し替える）。

\begin{tcolorbox}[enhanced jigsaw, arc=.35mm, bottomrule=.15mm, breakable, colback=white, colframe=quarto-callout-tip-color-frame, left=2mm, leftrule=.75mm, opacityback=0, rightrule=.15mm, toprule=.15mm]

\textbf{地図の使い方}：第VII部の各章はこの式の異なる因子を解き明かす。プランク関数地図と並べて読めば、本書を貫く二つの逆引きが同時に把握できる。第VIII部の
QED で、両者は前因子 \(2h\nu^3/c^2\)
を介して統一される（式~\ref{eq-einstein-A} 参照）。

\end{tcolorbox}

\begin{tcolorbox}[enhanced jigsaw, arc=.35mm, bottomrule=.15mm, bottomtitle=1mm, breakable, colback=white, colbacktitle=quarto-callout-note-color!10!white, colframe=quarto-callout-note-color-frame, coltitle=black, left=2mm, leftrule=.75mm, opacityback=0, opacitybacktitle=0.6, rightrule=.15mm, title=\textcolor{quarto-callout-note-color}{\faInfo}\hspace{0.5em}{ノート}, titlerule=0mm, toprule=.15mm, toptitle=1mm]

\textbf{方針}：本章は中心地図の \(\phi(\nu-\nu_0)\)
の中身を見る。第22章で得た \(\nu_0, f_{lu}\)
と組み合わせて、観測される線輪郭が物理量の関数として書ける。観測 →
物理量の逆引きが完成する。

\end{tcolorbox}

\section*{この章で答える問い}\label{sec-line-shape-questions}
\addcontentsline{toc}{section}{この章で答える問い}

\markright{この章で答える問い}

\begin{itemize}
\tightlist
\item
  線の幅は何で決まるのか ― 量子論、運動学、衝突のどれが支配するか
\item
  なぜ線の中心は Gaussian、翼は Lorentzian という Voigt 形になるのか
\item
  観測される線輪郭から、温度・密度・速度場・磁場をどうやって読み取るのか
\item
  マクロな速度場（自転、対流、乱流）の効果は、ミクロな線幅とどう区別されるのか
\end{itemize}

\section*{到達目標}\label{sec-line-shape-goals}
\addcontentsline{toc}{section}{到達目標}

\markright{到達目標}

この章を読み終えると、読者は次のことができるようになる：

\begin{itemize}
\tightlist
\item
  自然幅・Doppler 幅・衝突幅の物理的起源と振動数依存を説明できる
\item
  Voigt プロファイルを Doppler（コア）と Lorentz（翼）の合成として描ける
\item
  観測される線輪郭から、温度・密度・速度場・磁場を逆引きする手順を組み立てられる
\end{itemize}

\begin{center}\rule{0.5\linewidth}{0.5pt}\end{center}

\section{23.1 自然幅 ― Lorentzian の起源}\label{sec-line-shape-natural}

{[}本文目安：B3-B4{]}

第11章 §11.2 と第12章で見たように、励起準位 \(u\)
には有限の寿命がある。重要なのは、この寿命を決めるのは \textbf{上準位
\(u\) から到達可能なすべての下準位 \(l'\) への自発放出率の総和}
である点である：

\begin{equation}\protect\phantomsection\label{eq-total-decay-rate}{
\Gamma_u = \sum_{l'} A_{ul'}
}\end{equation}

すなわち上準位の全崩壊率 \(\Gamma_u\) は、注目している遷移 \(u \to l\)
の単一の \(A_{ul}\) ではなく、\(u\)
から落ちうるすべての経路の和である。上準位の寿命は
\(\tau_u = 1/\Gamma_u\) で与えられる。励起状態の量子振幅が
\(e^{-\Gamma_u t/2}\) で減衰することから、エネルギー準位そのものが
Lorentzian の形でエネルギー幅を持つ。観測される線 \(u \to l\)
のうち、この遷移が全崩壊に占める割合は
\textbf{分岐比}（\textbf{branching ratio}）\(A_{ul}/\Gamma_u\)
で与えられ、線の強度を決める。

\begin{tcolorbox}[enhanced jigsaw, arc=.35mm, bottomrule=.15mm, bottomtitle=1mm, breakable, colback=white, colbacktitle=quarto-callout-note-color!10!white, colframe=quarto-callout-note-color-frame, coltitle=black, left=2mm, leftrule=.75mm, opacityback=0, opacitybacktitle=0.6, rightrule=.15mm, title=\textcolor{quarto-callout-note-color}{\faInfo}\hspace{0.5em}{ノート}, titlerule=0mm, toprule=.15mm, toptitle=1mm]

\textbf{記号の規約}（混同しやすいので注意）：

\begin{itemize}
\tightlist
\item
  \(A_{ul}\) ：遷移 \(u \to l\) の自発放出率。次元 {[}s\(^{-1}\){]}
\item
  \(\Gamma_u \equiv \sum_{l'} A_{ul'}\) ：上準位の
  \textbf{全崩壊率}（全下準位への和）。自然幅を決めるのはこれ。下準位
  \(l\) も励起準位なら、その崩壊率 \(\Gamma_l\) も加わり
  \(\Gamma_\text{nat} = \Gamma_u + \Gamma_l\) となる
\item
  \(\Gamma_\text{nat}\) ：\textbf{角振動数単位} の自然幅（\textbf{FWHM},
  \(\omega\) 軸）。基底状態への遷移など下準位が安定な場合は
  \(\Gamma_\text{nat} = \Gamma_u\)
\item
  \(\gamma_\text{nat} \equiv \Gamma_\text{nat}/(2\pi)\)
  ：\textbf{振動数単位} の自然幅（\textbf{FWHM}, \(\nu\) 軸、Hz）
\item
  HWHM（半値半幅）はこれらの半分：\(\Gamma_\text{nat}/2\)（\(\omega\)）、\(\gamma_\text{nat}/2 = \Gamma_\text{nat}/(4\pi)\)（\(\nu\)）
\end{itemize}

単一の下準位しか持たない遷移（例：基底状態への共鳴線）では
\(\Gamma_u = A_{ul}\) となり、\(\Gamma_\text{nat} = A_{ul}\)
と簡約される。これらは「\textbf{寿命 \(\tau_u\) → 振幅減衰 → エネルギー
Lorentzian
の幅}」という一連の関係から出るもので、\(\Delta E \cdot \tau_u \gtrsim \hbar\)
という不確定性原理の不等式は \textbf{オーダー推定}
にすぎず、厳密には等号にならない（古典的減衰振動子の Fourier 変換が
Lorentzian になることが厳密な起源、付録 §A.8.2 参照）。

\end{tcolorbox}

これに対応する規格化された線輪郭関数（\textbf{Lorentzian}）は

\begin{equation}\protect\phantomsection\label{eq-lorentzian}{
\phi_L(\nu - \nu_0) = \frac{1}{\pi} \frac{\gamma_\text{nat}/2}{(\nu - \nu_0)^2 + (\gamma_\text{nat}/2)^2}
}\end{equation}

特徴：

\begin{itemize}
\tightlist
\item
  中心 \(\nu = \nu_0\)
  で最大、\(|\nu - \nu_0| = \gamma_\text{nat}/2\)（HWHM）で半値
\item
  \textbf{翼}（wing）が遠くまでべき的に減衰（\(\propto 1/(\Delta\nu)^2\)）―
  これが Voigt 線の遠方を決める
\item
  \(\int \phi_L\, d\nu = 1\)（規格化）
\end{itemize}

第12章で扱った \textbf{古典的調和振動子の減衰モデル} が同じ Lorentzian
を与えることが、振動子強度モデルの古典-量子対応を示す。

\begin{tcolorbox}[enhanced jigsaw, arc=.35mm, bottomrule=.15mm, bottomtitle=1mm, breakable, colback=white, colbacktitle=quarto-callout-note-color!10!white, colframe=quarto-callout-note-color-frame, coltitle=black, left=2mm, leftrule=.75mm, opacityback=0, opacitybacktitle=0.6, rightrule=.15mm, title=\textcolor{quarto-callout-note-color}{\faInfo}\hspace{0.5em}{ノート}, titlerule=0mm, toprule=.15mm, toptitle=1mm]

\textbf{観測との対応}：水素 Ly\(\alpha\)（\(\lambda = 121.6\)
nm）は上準位 \(2p\) から基底状態 \(1s\) への単一経路で崩壊するため
\(\Gamma_u = A_{ul} = 6.27 \times 10^8\) s\(^{-1}\)
となり、自然幅は波長単位で
\(\Delta\lambda_\text{nat} = \lambda^2 \Gamma_\text{nat}/(2\pi c) \simeq 4.9 \times 10^{-5}\)
Å
である。多くの場合、自然幅は他の拡がり（Doppler、衝突）に比べて小さいが、\textbf{線の翼}（damping
wing）には自然幅の Lorentzian の寄与が現れる。QSO の Ly\(\alpha\)
強吸収系（damped Ly\(\alpha\) system）で観測される \textbf{減衰翼}
がその例である。

\end{tcolorbox}

\section{23.2 Doppler 拡がり ― Gaussian
の起源}\label{sec-line-shape-doppler}

{[}本文目安：B2-B3{]}

ガスの粒子は熱運動で \textbf{視線方向に速度}
を持つ。各粒子は自分の速度に応じて、観測者の系で線中心振動数が
\textbf{Doppler シフト} する：

\begin{equation}\protect\phantomsection\label{eq-doppler-shift}{
\nu = \nu_0 \left(1 + \frac{v_\parallel}{c}\right)
}\end{equation}

熱的 Maxwell-Boltzmann 速度分布のもとでは、視線速度成分 \(v_\parallel\)
も Gaussian 分布になる。これを線輪郭として積分すると、

\begin{equation}\protect\phantomsection\label{eq-gaussian-doppler}{
\phi_D(\nu - \nu_0) = \frac{1}{\Delta\nu_D \sqrt{\pi}} \exp\left[-\frac{(\nu - \nu_0)^2}{\Delta\nu_D^2}\right]
}\end{equation}

ここで \textbf{Doppler 幅}

\begin{equation}\protect\phantomsection\label{eq-doppler-width}{
\Delta\nu_D = \frac{\nu_0}{c} \sqrt{\frac{2k_B T}{m}}
}\end{equation}

\(m\) は原子・分子質量。

特徴：

\begin{itemize}
\tightlist
\item
  中心が Gaussian で \textbf{狭くて深い}
\item
  翼が \textbf{指数的に急減衰}（\(\propto e^{-(\Delta\nu)^2}\)）
\item
  質量に依存：軽い水素は重い鉄より広い Doppler 幅
\end{itemize}

\begin{tcolorbox}[enhanced jigsaw, arc=.35mm, bottomrule=.15mm, breakable, colback=white, colframe=quarto-callout-tip-color-frame, left=2mm, leftrule=.75mm, opacityback=0, rightrule=.15mm, toprule=.15mm]
\begin{minipage}[t]{5.5mm}
\textcolor{quarto-callout-tip-color}{\faLightbulb}
\end{minipage}%
\begin{minipage}[t]{\textwidth - 5.5mm}

\textbf{問い}：温度 \(T = 10^4\) K の HII 領域における、水素 H\(\alpha\)
の Doppler 幅（FWHM）を km/s で求めよ。

\textbf{短答}：\(\Delta v_D \simeq \sqrt{2k_B T/m_p} \simeq 12.8\)
km/s（標準偏差）、FWHM = \(2\sqrt{\ln 2} \cdot \Delta v_D \simeq 21\)
km/s。

\textbf{もう一歩}：同じ温度の鉄（\(m = 56 m_p\)）の線幅は
\(\sqrt{56} \simeq 7.5\) 倍狭い。観測される線幅の元素質量依存性から
\textbf{熱的拡がり vs 他の拡がり} を区別できる ― これは観測解析の鍵

\end{minipage}%
\end{tcolorbox}

\section{23.3 衝突拡がり ― Lorentzian
の追加}\label{sec-line-shape-pressure}

{[}本文目安：B3{]}

ガスが密になると、原子は \textbf{他の粒子との衝突}
で励起準位の位相が乱される。これは励起準位の有効寿命を短くし、自然幅と同形の
\textbf{Lorentzian 拡がり} を加える：

\begin{equation}\protect\phantomsection\label{eq-collision-width}{
\Gamma_\text{coll} = 2\sum_i n_i \langle \sigma_i v_i \rangle
}\end{equation}

ここで \(n_i\) は衝突相手の数密度、\(\sigma_i v_i\)
は衝突断面積×相対速度の平均。

\textbf{特徴}：

\begin{itemize}
\tightlist
\item
  \(\Gamma_\text{coll}\) は \textbf{密度に比例} →
  高密度天体（恒星深部、巨星大気）で大きい
\item
  Lorentzian の形は自然幅と同じ ― 両者は \textbf{合算} して扱える
\item
  圧力に依存することから「圧力 broadening」とも呼ばれる
\end{itemize}

\begin{tcolorbox}[enhanced jigsaw, arc=.35mm, bottomrule=.15mm, bottomtitle=1mm, breakable, colback=white, colbacktitle=quarto-callout-note-color!10!white, colframe=quarto-callout-note-color-frame, coltitle=black, left=2mm, leftrule=.75mm, opacityback=0, opacitybacktitle=0.6, rightrule=.15mm, title=\textcolor{quarto-callout-note-color}{\faInfo}\hspace{0.5em}{ノート}, titlerule=0mm, toprule=.15mm, toptitle=1mm]

\textbf{観測との対応}：恒星の \textbf{重力（圧力）の感度線} ―
例えば水素線の自然幅 + 衝突幅の合算が重力に強く依存する。Mg I b 線や Na
D 線の翼の形を解析することで、\textbf{恒星の表面重力 \(\log g\)}
が決まる。これは恒星の質量・半径推定の基本パラメータである。

\end{tcolorbox}

\section{23.4 Voigt プロファイル ―
二つの拡がりの合成}\label{sec-line-shape-voigt}

{[}本文目安：B3-B4{]}

\begin{figure}[H]

\centering{

\includegraphics[width=0.92\linewidth,height=\textheight,keepaspectratio]{index_files/mediabag/part7/../images/voigt-profile.pdf}

}

\caption{\label{fig-voigt-profile}Voigt
プロファイル（黒）と、その元となる
Gaussian（赤・点線）／Lorentzian（青・破線）の比較。\textbf{線形スケール}ではコアの形が、\textbf{対数スケール}では翼の形が比較しやすい。\textbf{コアは
Gaussian で急速に落ち、翼に行くと Lorentzian がべき則 \(1/\Delta\nu^2\)
で長く伸びる} ― これが Voigt
線の二面性で、観測者がコア（温度）と翼（密度・圧力）を別々のプローブとして使える理由。}

\end{figure}%

自然幅 + 衝突幅（合算した Lorentzian、幅 \(\gamma\)）と Doppler
拡がり（Gaussian、幅 \(\Delta\nu_D\)）の \textbf{両方}
が効くのが現実である。線輪郭は両者の \textbf{畳み込み} になる：

\begin{equation}\protect\phantomsection\label{eq-voigt-convolution}{
\phi_V(\nu - \nu_0) = \phi_L \otimes \phi_D = \int \phi_L(\nu - \nu' - \nu_0) \phi_D(\nu' - \nu_0) d\nu'
}\end{equation}

これが \textbf{Voigt
プロファイル}。閉形式では書けないが、無次元パラメータ

\[
a = \frac{\gamma}{2\Delta\nu_D}, \quad x = \frac{\nu - \nu_0}{\Delta\nu_D}
\]

を使って Voigt 関数 \(H(a, x)\) で表される。

\textbf{振る舞いの二相}：

{\def\LTcaptype{none} % do not increment counter
\begin{longtable}[]{@{}lll@{}}
\toprule\noalign{}
領域 & 支配 & 振る舞い \\
\midrule\noalign{}
\endhead
\bottomrule\noalign{}
\endlastfoot
\$ & x & \textless{} \$ 数 \\
\(|x| \gg 1\) & Lorentz 翼 & \(\propto 1/x^2\)（緩減衰） \\
\end{longtable}
}

\begin{tcolorbox}[enhanced jigsaw, arc=.35mm, bottomrule=.15mm, breakable, colback=white, colframe=quarto-callout-tip-color-frame, left=2mm, leftrule=.75mm, opacityback=0, rightrule=.15mm, toprule=.15mm]
\begin{minipage}[t]{5.5mm}
\textcolor{quarto-callout-tip-color}{\faLightbulb}
\end{minipage}%
\begin{minipage}[t]{\textwidth - 5.5mm}

\textbf{問い}：Voigt プロファイルで、\textbf{コア（Doppler）}
が支配する範囲を観測者はどう判別できるか？

\textbf{短答}：典型的に
\(|x| < 3\)（\(|\Delta\nu| < 3 \Delta\nu_D\)）でコアが支配。\(|x| > 5\)
になると、Doppler の Gaussian は \(e^{-25} \simeq 10^{-11}\)
で消え、Lorentz 翼が浮かび上がる。

\textbf{もう一歩}：観測者は line core と line wing
を別々の物理プローブとして使う：

\begin{itemize}
\tightlist
\item
  Core：温度（Doppler 幅）、視線速度（中心シフト）
\item
  Wing：自然幅 + 衝突幅 → 圧力、密度
\item
  両者を \textbf{同時にフィット}
  することで、温度と密度を独立に決められる
\end{itemize}

\end{minipage}%
\end{tcolorbox}

\section{23.5 マクロな速度場 ―
自転、対流、乱流}\label{sec-line-shape-macro}

{[}本文目安：B2-B3{]}

ミクロな熱運動（Doppler 拡がり）のほかに、\textbf{マクロな速度場}
も線幅に寄与する：

\begin{itemize}
\tightlist
\item
  \textbf{恒星自転}：表面の一方が近づき他方が遠ざかる → 線が
  \textbf{平頂化}
\item
  \textbf{対流}（granulation）：恒星表面の上下流 → 線の
  \textbf{非対称化}（線芯青方シフト）
\item
  \textbf{乱流}（micro/macroturbulence）：等方ランダム速度 → 線幅の
  \textbf{追加 Gaussian 成分}
\item
  \textbf{吹き出し風}（stellar wind）：青翼が吸収（P Cygni 形）
\end{itemize}

これらマクロ速度場は \textbf{熱的線幅と等価な Gaussian}
として加算的に効く：

\begin{equation}\protect\phantomsection\label{eq-velocity-quadrature-sum}{
\Delta\nu_\text{D,total}^2 = \Delta\nu_\text{D,thermal}^2 + \Delta\nu_\text{D,turb}^2 + \ldots
}\end{equation}

熱的成分と乱流成分の分離は、\textbf{元素質量に対する依存性}
から行える：熱は \(1/\sqrt{m}\) 依存、乱流は元素質量によらない。

\begin{tcolorbox}[enhanced jigsaw, arc=.35mm, bottomrule=.15mm, breakable, colback=white, colframe=quarto-callout-warning-color-frame, left=2mm, leftrule=.75mm, opacityback=0, rightrule=.15mm, toprule=.15mm]
\begin{minipage}[t]{5.5mm}
\textcolor{quarto-callout-warning-color}{\faExclamationTriangle}
\end{minipage}%
\begin{minipage}[t]{\textwidth - 5.5mm}

\textbf{注意}：上の「Gaussian
の二乗和」が正当化できるのは、速度場のスケールが光子の
\textbf{平均自由行程より小さい} 場合 ― すなわち
\textbf{マイクロ乱流}（\textbf{microturbulence}）。この場合は線吸収係数
\(\alpha_\nu\) 自体に乱流成分が織り込まれ、\(\Delta\nu_D\)
の単純な拡張で書ける。

一方、速度場のスケールが平均自由行程より大きい
\textbf{マクロ乱流}（\textbf{macroturbulence}）では、線形成は局所的に進み、最後に
\textbf{射出スペクトルへの畳み込み}
として速度場が効く。両者の数学的扱いは異なる（前者は線吸収係数の修正、後者は出てきたスペクトルの畳み込み）。Gray
(2005) 第 17-18 章の扱いを参照

\end{minipage}%
\end{tcolorbox}

\begin{tcolorbox}[enhanced jigsaw, arc=.35mm, bottomrule=.15mm, bottomtitle=1mm, breakable, colback=white, colbacktitle=quarto-callout-note-color!10!white, colframe=quarto-callout-note-color-frame, coltitle=black, left=2mm, leftrule=.75mm, opacityback=0, opacitybacktitle=0.6, rightrule=.15mm, title=\textcolor{quarto-callout-note-color}{\faInfo}\hspace{0.5em}{ノート}, titlerule=0mm, toprule=.15mm, toptitle=1mm]

\textbf{観測との対応}：太陽の吸収線解析では、\(\xi_\text{turb} \simeq 1\)〜\(2\)
km/s のマイクロ乱流と、\(\zeta_\text{macro} \simeq 3\)〜\(5\) km/s
のマクロ乱流が経験的に必要である。これは対流による速度場の積分量で、3D
数値流体シミュレーションと整合が取れている。

\end{tcolorbox}

\section{23.6 磁場の効果 ― Zeeman 効果}\label{sec-line-shape-zeeman}

{[}本文目安：B3{]}

磁場 \(\mathbf{B}\) の下では、原子のエネルギー準位が磁気量子数 \(m_J\)
に応じて分裂する：\textbf{Zeeman 効果}。

\begin{equation}\protect\phantomsection\label{eq-zeeman-energy}{
\Delta E_\text{Zeeman} = g_J \mu_B B m_J
}\end{equation}

ここで \(g_J\) は Landé \(g\) 因子、\(\mu_B = e\hbar/2m_e\) は Bohr
磁子。

観測される線輪郭への影響：

\begin{itemize}
\tightlist
\item
  弱磁場：単一線が \textbf{三本} に分裂（ノーマル
  Zeeman）または複数本に分裂（アノーマラス Zeeman）
\item
  線分裂は \textbf{円偏光} を伴う（縦磁場成分）または
  \textbf{直線偏光}（横磁場成分）
\item
  分解できない場合でも、線が \textbf{広がって見える}（Zeeman
  broadening）
\end{itemize}

\begin{tcolorbox}[enhanced jigsaw, arc=.35mm, bottomrule=.15mm, bottomtitle=1mm, breakable, colback=white, colbacktitle=quarto-callout-note-color!10!white, colframe=quarto-callout-note-color-frame, coltitle=black, left=2mm, leftrule=.75mm, opacityback=0, opacitybacktitle=0.6, rightrule=.15mm, title=\textcolor{quarto-callout-note-color}{\faInfo}\hspace{0.5em}{ノート}, titlerule=0mm, toprule=.15mm, toptitle=1mm]

\textbf{観測との対応}：太陽黒点の磁場は Fe I \$\lambda\$6173 線の Zeeman
分裂から測られる。星形成領域での磁場は OH メーザーや CN 線の Zeeman
効果で測られる ― \textbf{線スペクトルが磁場の数少ない直接観測手段}
であることは、第VII部の応用の重要側面である。

\end{tcolorbox}

\section*{23.7 使えるようになった道具}\label{sec-line-shape-tools}
\addcontentsline{toc}{section}{23.7 使えるようになった道具}

\markright{23.7 使えるようになった道具}

\begin{tcolorbox}[enhanced jigsaw, arc=.35mm, bottomrule=.15mm, bottomtitle=1mm, breakable, colback=white, colbacktitle=quarto-callout-note-color!10!white, colframe=quarto-callout-note-color-frame, coltitle=black, left=2mm, leftrule=.75mm, opacityback=0, opacitybacktitle=0.6, rightrule=.15mm, title=\textcolor{quarto-callout-note-color}{\faInfo}\hspace{0.5em}{ノート}, titlerule=0mm, toprule=.15mm, toptitle=1mm]

\begin{itemize}
\tightlist
\item
  自然幅 ―
  Lorentzian、\(\gamma_\text{nat} = \Gamma_\text{nat}/(2\pi)\)（\(\Gamma_\text{nat} = \sum_{l'} A_{ul'}\)
  ＝上準位の全崩壊率）
\item
  Doppler 幅 ― Gaussian、\(\Delta\nu_D = \nu_0 \sqrt{2k_BT/m}/c\)
\item
  衝突幅 ― Lorentzian、密度・圧力に比例
\item
  Voigt プロファイル ― コア（温度）と翼（密度）を独立にプローブ
\item
  マクロ速度場 ― 自転・対流・乱流、元素質量依存性で分離
\item
  Zeeman 効果 ― 線スペクトルから磁場を直接測る
\end{itemize}

\end{tcolorbox}

\begin{center}\rule{0.5\linewidth}{0.5pt}\end{center}

\section*{この章で何がわかったか}\label{sec-line-shape-summary}
\addcontentsline{toc}{section}{この章で何がわかったか}

\markright{この章で何がわかったか}

\textbf{中心地図に戻る}

本章で、中心地図
\(\alpha_\nu = (\pi e^2/m_e c) f_{lu} n_l \phi(\nu-\nu_0)\) の最後の因子
\(\phi(\nu-\nu_0)\) が物理的に解読された：

\begin{itemize}
\tightlist
\item
  自然幅は不確定性原理から → Lorentzian
\item
  Doppler 幅は熱運動の Maxwell-Boltzmann 分布から → Gaussian
\item
  衝突幅は密度・圧力から → 追加 Lorentzian
\item
  Voigt プロファイルが両者の畳み込み
\end{itemize}

これで中心地図のすべての因子（\(\pi e^2/m_e c\), \(f_{lu}\), \(n_l\),
\(\phi\)）が逆引きされた（前因子 \(\pi e^2/m_e c\) は §22.4b の古典
Lorentz 振動子からの導出で完結）。\textbf{線スペクトルから物理を読む}
ための道具立てが揃った。

次章では、これらの道具を \textbf{実際の観測解析} にどう動員するか ―
視線速度、化学組成、温度・密度診断、宇宙論的応用 ― を見る。

\section*{演習問題}\label{sec-line-shape-exercises}
\addcontentsline{toc}{section}{演習問題}

\markright{演習問題}

\textbf{模範解答}
は折りたたみを展開して確認できる（オンライン版）。まず自力で解いてから開くこと。

\subsection*{問題 23-1　Doppler 幅 ― 熱運動の Gaussian}\label{ex-23-1}
\addcontentsline{toc}{subsection}{問題 23-1　Doppler 幅 ― 熱運動の
Gaussian}

{[}★ 難易度：☆☆ {]}

§23.2 の Doppler 幅
\(\Delta\nu_D=\dfrac{\nu_0}{c}\sqrt{\dfrac{2k_BT}{m}}\)。HII
領域（\(T=10^4\) K）の水素 H\(\alpha\) について。

\begin{enumerate}
\def\labelenumi{\arabic{enumi}.}
\tightlist
\item
  速度の標準偏差に相当する \(\Delta v_D=\sqrt{2k_BT/m_p}\) を km/s
  で求めよ（\(m_p=1.67\times10^{-27}\) kg）。
\item
  FWHM \(=2\sqrt{\ln2}\,\Delta v_D\) を求めよ。
\item
  同じ温度の鉄（\(m=56\,m_p\)）の Doppler
  幅は水素の何倍か。観測線幅の元素質量依存性から熱的拡がりと他の拡がりをどう区別できるか述べよ。
\end{enumerate}

\textbf{関連}：\hyperref[sec-line-shape-doppler]{§23.2 Doppler
拡がり}／Maxwell
分布は\hyperref[ux7b2c7ux7ae0-ux7d71ux8a08ux529bux5b66ux306eux57faux672c]{第III部
演習}、マクロ速度場は\hyperref[sec-line-shape-macro]{§23.5}（\hyperref[ex-23-4]{演習
23-4}）。

\section{模範解答（問題
23-1）}\label{ux6a21ux7bc4ux89e3ux7b54ux554fux984c-23-1}

\textbf{(1)}
\(\Delta v_D=\sqrt{2k_BT/m_p}=\sqrt{2\times1.381\times10^{-23}\times10^4/1.67\times10^{-27}}=\sqrt{1.654\times10^8}=1.29\times10^4\)
m/s \(=12.9\) km/s。

\textbf{(2)} FWHM
\(=2\sqrt{\ln2}\times12.9=2\times0.833\times12.9=21.4\) km/s。

\textbf{(3)} \(\Delta\nu_D\propto1/\sqrt{m}\) ゆえ鉄は
\(\sqrt{56}\simeq7.5\)
倍狭い。観測で「重い元素の線ほど狭い」なら熱的拡がりが支配（質量依存）。一方、乱流・自転などマクロ速度場は元素質量に依らないので、異なる質量の元素の線幅を比べると、熱的成分と非熱的（質量無依存）成分を分離できる。

\textbf{答え}：\(\Delta v_D=12.9\) km/s、FWHM \(=21.4\) km/s。鉄は
\(\sqrt{56}\simeq7.5\) 倍狭い。質量依存性で熱的と乱流を分離。

\subsection*{問題 23-2　自然幅 ― 寿命から Lorentzian}\label{ex-23-2}
\addcontentsline{toc}{subsection}{問題 23-2　自然幅 ― 寿命から
Lorentzian}

{[}★ 難易度：☆☆ {]}

§23.1 の自然幅。励起寿命 \(\tau_u=1/A_{ul}\)、振動数単位の自然幅
\(\gamma_\mathrm{nat}=A_{ul}/(2\pi)\)。Ly\(\alpha\)（\(\lambda=121.6\)
nm、\(A_{ul}=6.27\times10^8\) s\(^{-1}\)）について。

\begin{enumerate}
\def\labelenumi{\arabic{enumi}.}
\tightlist
\item
  励起寿命 \(\tau_u\) と自然幅 \(\gamma_\mathrm{nat}\)（Hz）を求めよ。
\item
  波長単位の自然幅
  \(\Delta\lambda_\mathrm{nat}=\dfrac{\lambda^2 A_{ul}}{2\pi c}\) を Å
  単位で求め、\(4.9\times10^{-5}\) Å を確かめよ。
\item
  自然幅 Lorentzian の「翼」が \(1/(\Delta\nu)^2\)
  でべき的に伸びることが、QSO の damped Ly\(\alpha\)
  系の減衰翼として観測される。なぜ自然幅は線中心より翼で重要なのか述べよ。
\end{enumerate}

\textbf{関連}：\hyperref[sec-line-shape-natural]{§23.1
自然幅}／\(A_{ul}\) と自然幅は\hyperref[ex-11-4]{演習
11-4}、真空揺らぎとの関係は\hyperref[sec-qed-minimum-natural-width]{§26.4}（\href{../part8/26-qed-minimum.qmd}{第VIII部
演習}）。

\section{模範解答（問題
23-2）}\label{ux6a21ux7bc4ux89e3ux7b54ux554fux984c-23-2}

\textbf{(1)} \(\tau_u=1/A_{ul}=1/(6.27\times10^8)=1.6\times10^{-9}\)
s（1.6
ns）。\(\gamma_\mathrm{nat}=A_{ul}/(2\pi)=6.27\times10^8/6.283=9.98\times10^7\simeq10^8\)
Hz。

\textbf{(2)}
\(\Delta\lambda_\mathrm{nat}=\dfrac{\lambda^2 A_{ul}}{2\pi c}=\dfrac{(1.216\times10^{-7})^2\times6.27\times10^8}{2\pi\times3\times10^8}=\dfrac{1.479\times10^{-14}\times6.27\times10^8}{1.885\times10^9}=\dfrac{9.27\times10^{-6}}{1.885\times10^9}=4.9\times10^{-15}\)
m \(=4.9\times10^{-5}\) Å。✓

\textbf{(3)} 線中心は Doppler の Gaussian
コアが支配し（指数減衰で狭い）、自然幅 Lorentzian は埋もれる。しかし
Lorentzian の翼は \(1/(\Delta\nu)^2\)
でべき的に緩く伸びるため、中心から十分離れた翼では
Gaussian（\(e^{-(\Delta\nu)^2}\)）が先に消えて Lorentzian
が浮かび上がる。よって自然幅は線翼（damping wing）で重要。

\textbf{答え}：\(\tau_u=1.6\) ns、\(\gamma_\mathrm{nat}\simeq10^8\)
Hz、\(\Delta\lambda_\mathrm{nat}=4.9\times10^{-5}\) Å。Lorentzian 翼が
\(1/\Delta\nu^2\) で伸びるため翼で支配的。

\subsection*{問題 23-3　Voigt プロファイル ―
コアと翼を別プローブに}\label{ex-23-3}
\addcontentsline{toc}{subsection}{問題 23-3　Voigt プロファイル ―
コアと翼を別プローブに}

{[}★ 難易度：☆☆ {]}

§23.4 の Voigt プロファイルは Gaussian（Doppler）と
Lorentzian（自然幅＋衝突幅）の畳み込み。無次元パラメータ
\(a=\gamma/(2\Delta\nu_D)\)、\(x=(\nu-\nu_0)/\Delta\nu_D\)。

\begin{enumerate}
\def\labelenumi{\arabic{enumi}.}
\tightlist
\item
  コア（\(|x|<\) 数）と翼（\(|x|\gg1\)）で、それぞれ Gaussian と
  Lorentzian のどちらが支配するか述べよ。\(|x|>5\) で Doppler 成分が
  \(e^{-25}\sim10^{-11}\) に消えることを確かめよ。
\item
  コアと翼が、観測者にとってそれぞれどんな物理量（温度・速度／密度・圧力）のプローブになるか述べよ。
\item
  コアと翼を同時にフィットすることで、温度と密度を独立に決められる理由を述べよ。
\end{enumerate}

\textbf{関連}：\hyperref[sec-line-shape-voigt]{§23.4 Voigt
プロファイル}、\hyperref[sec-line-shape-pressure]{§23.3
衝突拡がり}／恒星の \(\log g\)
決定は\hyperref[sec-line-shape-pressure]{§23.3}、curve of growth
の減衰翼は\hyperref[ex-21-3]{演習 21-3}。

\section{模範解答（問題
23-3）}\label{ux6a21ux7bc4ux89e3ux7b54ux554fux984c-23-3}

\textbf{(1)} コア（\(|x|<3\) 程度）は
Gaussian（Doppler）が支配。翼（\(|x|\gg1\)）は Lorentzian
が支配。\(|x|=5\) で Gaussian は
\(e^{-x^2}=e^{-25}\simeq1.4\times10^{-11}\) とほぼ消え、\(1/x^2\)
で緩く減衰する Lorentz 翼だけが残る。

\textbf{(2)} コア：Doppler 幅 \(\Delta\nu_D\propto\sqrt{T/m}\)
から\textbf{温度}、中心シフトから\textbf{視線速度}。翼：Lorentzian 幅
\(\gamma=\)
自然幅＋衝突幅で、衝突幅が\textbf{密度・圧力}に比例するので翼の形から\textbf{密度・圧力}。

\textbf{(3)}
コアは温度に、翼は密度・圧力に主に依存し、両者は線輪郭の異なる部分に別々に効く。観測されたコアと翼を同時にフィットすると、温度（コア幅）と密度（翼の広がり）を分離して独立に決められる。恒星の表面重力
\(\log g\) は翼の解析から得られる。

\textbf{答え}：コア＝Gaussian（温度・速度）、翼＝Lorentzian（密度・圧力）。\(|x|>5\)
で Gaussian は \(10^{-11}\)
に消え翼が浮かぶ。同時フィットで温度と密度を分離。

\subsection*{問題 23-4　熱的幅と乱流の分離、Zeeman
で磁場}\label{ex-23-4}
\addcontentsline{toc}{subsection}{問題 23-4　熱的幅と乱流の分離、Zeeman
で磁場}

{[}★ 難易度：☆☆ {]}

§23.5・22.6 のマクロ速度場と磁場。Gaussian 成分の二乗和
\(\Delta\nu_{D,\mathrm{total}}^2=\Delta\nu_{D,\mathrm{thermal}}^2+\Delta\nu_{D,\mathrm{turb}}^2\)、Zeeman
分裂 \(\Delta E=g_J\mu_B B m_J\)。

\begin{enumerate}
\def\labelenumi{\arabic{enumi}.}
\tightlist
\item
  熱的成分（\(\propto1/\sqrt m\)）と乱流成分（元素質量に依らない）が、異なる質量の元素の線幅比較でどう分離できるか述べよ。
\item
  マイクロ乱流とマクロ乱流で数学的扱いが異なる理由（速度場のスケールと光子平均自由行程の大小）を一言で述べよ。
\item
  Zeeman
  分裂が線スペクトルから磁場を測る数少ない直接手段である理由を述べ、観測例（太陽黒点・星形成領域）を1つ挙げよ。
\end{enumerate}

\textbf{関連}：\hyperref[sec-line-shape-macro]{§23.5
マクロな速度場}、\hyperref[sec-line-shape-zeeman]{§23.6 Zeeman
効果}／偏光は\hyperref[ex-13-4]{演習 13-4}。

\section{模範解答（問題
23-4）}\label{ux6a21ux7bc4ux89e3ux7b54ux554fux984c-23-4}

\textbf{(1)} 熱的 Doppler 幅は \(\Delta v_\mathrm{th}\propto1/\sqrt m\)
で重い元素ほど狭い。乱流は等方ランダム速度で元素質量に依らず一定。複数の元素（軽い
H と重い Fe
など）の線幅を測り、\(\Delta v^2=\Delta v_\mathrm{th}^2(m)+\xi_\mathrm{turb}^2\)
を質量に対してプロットすると、質量依存部分（熱的）と一定部分（乱流
\(\xi_\mathrm{turb}\)）に分離できる。

\textbf{(2)}
速度場のスケールが光子の平均自由行程より\textbf{小さい}＝マイクロ乱流：線吸収係数
\(\alpha_\nu\) 自体に織り込まれ、\(\Delta\nu_D\)
の単純拡張で扱える。\textbf{大きい}＝マクロ乱流：線形成は局所的に進み、最後に射出スペクトルへの畳み込みとして効く。扱う段階が違う。

\textbf{(3)}
磁場は放射スペクトルに直接の連続的痕跡を残しにくいが、Zeeman
効果は準位を \(m_J\) に応じて分裂させ（分裂幅
\(\propto B\)）、円偏光・直線偏光を伴うので、線の分裂・偏光から磁場の強度と方向を直接測れる。例：太陽黒点（Fe
I \$\lambda\$6173）、星形成領域（OH メーザー・CN 線）。

\textbf{答え}：熱的（\(\propto1/\sqrt m\)）と乱流（質量無依存）は元素質量比較で分離。マイクロ/マクロは平均自由行程との大小で扱いが違う。Zeeman
分裂・偏光が磁場の直接手段（太陽黒点など）。

\section*{さらに学ぶための参考文献}\label{sec-line-shape-further}
\addcontentsline{toc}{section}{さらに学ぶための参考文献}

\markright{さらに学ぶための参考文献}

\begin{itemize}
\tightlist
\item
  Sobelman et al., \emph{Excitation of Atoms and Broadening of Spectral
  Lines} (Springer, 2nd ed., 1995) --- 線輪郭の系統的扱い
\item
  Gray, \emph{The Observation and Analysis of Stellar Photospheres}
  (Cambridge, 3rd ed., 2005) --- 恒星線輪郭と速度場
\item
  del Toro Iniesta, \emph{Introduction to Spectropolarimetry}
  (Cambridge, 2003) --- Zeeman 効果と偏光分光
\item
  Mihalas, \emph{Stellar Atmospheres} (Freeman, 2nd ed., 1978) --- Voigt
  プロファイルの解析応用
\end{itemize}

\chapter{第24章　線スペクトルから天体物理を読む}\label{ux7b2c24ux7ae0-ux7ddaux30b9ux30daux30afux30c8ux30ebux304bux3089ux5929ux4f53ux7269ux7406ux3092ux8aadux3080}

\textbf{この章の主題}：第20-23章で整備した線スペクトルの逆引き道具（Einstein
係数・占有数・振動子強度・線輪郭）を、\textbf{現実の観測解析}
に動員する。視線速度・化学組成・温度・密度・電離度・磁場・宇宙論的赤方偏移という、\textbf{線スペクトルが直接読ませてくれる物理量}
を体系的に整理し、第VII部を閉じる。これは第VI部（連続スペクトルの応用）と対をなす、線スペクトル側の実務章である。

\section*{この章の中心地図}\label{sec-line-applications-map}
\addcontentsline{toc}{section}{この章の中心地図}

\markright{この章の中心地図}

\textbf{水素線スペクトルの地図 ― 二層構造}

\textbf{第一層：イコニック・アンカー（リュードベリ公式）}

\begin{equation}\protect\phantomsection\label{eq-rydberg}{
\frac{1}{\lambda} = R_\mathrm{H} \left( \frac{1}{n_1^2} - \frac{1}{n_2^2} \right)
}\end{equation}

線スペクトルの「骨格」を一目で示す。\(R_\mathrm{H}\)
は基本定数の組合せ、\(1/n^2\)
は束縛状態のエネルギー、離散性が「線が存在する理由」を語る。

\textbf{第二層：構造的マップ（Hα を標準例とした線吸収係数）}

\begin{equation}\protect\phantomsection\label{eq-line-absorption}{
\alpha_\nu = \frac{\pi e^2}{m_e c} \, f_{lu} \, n_l \, \phi(\nu - \nu_0)
}\end{equation}

{\def\LTcaptype{none} % do not increment counter
\begin{longtable}[]{@{}lll@{}}
\toprule\noalign{}
因子 & 物理的意味 & 対応する章 \\
\midrule\noalign{}
\endhead
\bottomrule\noalign{}
\endlastfoot
\(\alpha_\nu\) & 観測量（線吸収係数） & 第21章 \\
\(\pi e^2 / m_e c\) & 古典極限の線強度 & 第22章＋第VIII部 \\
\(f_{lu}\) & 振動子強度 & 第22章 \\
\(n_l\) & 下準位の占有数 & 第21章＋Part III \\
\(\phi(\nu - \nu_0)\) & 線輪郭（自然・Doppler・Voigt） & 第23章 \\
\(\nu_0\) & 線位置（リュードベリ） & 第22章 \\
\end{longtable}
}

\textbf{適用範囲（重要）}：この式の前因子 \(\pi e^2/m_e c\) と振動子強度
\(f_{lu}\) は、古典 Lorentz 振動子に対応する
\textbf{電気双極子（E1）許容線} に対するものである。21 cm
線（基底状態超微細構造の \textbf{磁気双極子 M1 遷移}）や
\textbf{禁制線（E2・M1 など）} では、この前因子・\(f_{lu}\)
をそのまま使えない。これらの線では、測定された自発放出率 \(A_{ul}\) から
Einstein
関係を介して有効遷移強度（\(B_{lu}\)）に読み替える（式~\ref{eq-line-absorption}
の構造は保ったまま、強度因子だけを差し替える）。

\begin{tcolorbox}[enhanced jigsaw, arc=.35mm, bottomrule=.15mm, breakable, colback=white, colframe=quarto-callout-tip-color-frame, left=2mm, leftrule=.75mm, opacityback=0, rightrule=.15mm, toprule=.15mm]

\textbf{地図の使い方}：第VII部の各章はこの式の異なる因子を解き明かす。プランク関数地図と並べて読めば、本書を貫く二つの逆引きが同時に把握できる。第VIII部の
QED で、両者は前因子 \(2h\nu^3/c^2\)
を介して統一される（式~\ref{eq-einstein-A} 参照）。

\end{tcolorbox}

\begin{tcolorbox}[enhanced jigsaw, arc=.35mm, bottomrule=.15mm, bottomtitle=1mm, breakable, colback=white, colbacktitle=quarto-callout-note-color!10!white, colframe=quarto-callout-note-color-frame, coltitle=black, left=2mm, leftrule=.75mm, opacityback=0, opacitybacktitle=0.6, rightrule=.15mm, title=\textcolor{quarto-callout-note-color}{\faInfo}\hspace{0.5em}{ノート}, titlerule=0mm, toprule=.15mm, toptitle=1mm]

\textbf{方針}：本章は「\textbf{逆引きの結果としての観測道具}」の章。中心地図
\(\alpha_\nu = (\pi e^2/m_e c) f_{lu} n_l \phi(\nu-\nu_0)\)
の各因子から、どんな物理量が読めるかを再整理し、観測の具体例で確認する。

\end{tcolorbox}

本章の応用を俯瞰するため、「観測で見える特徴 → 最小モデル →
逆引きされる物理量 → 破れる仮定」の対応を最初にまとめておく。

{\def\LTcaptype{none} % do not increment counter
\begin{longtable}[]{@{}
  >{\raggedright\arraybackslash}p{(\linewidth - 6\tabcolsep) * \real{0.2500}}
  >{\raggedright\arraybackslash}p{(\linewidth - 6\tabcolsep) * \real{0.2500}}
  >{\raggedright\arraybackslash}p{(\linewidth - 6\tabcolsep) * \real{0.2500}}
  >{\raggedright\arraybackslash}p{(\linewidth - 6\tabcolsep) * \real{0.2500}}@{}}
\toprule\noalign{}
\begin{minipage}[b]{\linewidth}\raggedright
観測で見える特徴
\end{minipage} & \begin{minipage}[b]{\linewidth}\raggedright
最小モデル
\end{minipage} & \begin{minipage}[b]{\linewidth}\raggedright
逆引きされる物理量
\end{minipage} & \begin{minipage}[b]{\linewidth}\raggedright
破れる仮定
\end{minipage} \\
\midrule\noalign{}
\endhead
\bottomrule\noalign{}
\endlastfoot
線中心のシフト & 非相対論的 Doppler & 視線速度・赤方偏移 & 高
\(z\)（相対論・宇宙膨張）、線ブレンド \\
等価幅 & 成長曲線（curve of growth） & 柱密度・化学組成 &
飽和、non-LTE、未解像の微細構造 \\
禁制線の線比 & 二準位＋衝突の励起平衡 & 電子温度・電子密度 &
臨界密度近傍、温度非一様 \\
多重電離段階の線比 & 電離平衡（Saha・光電離） & 電離度・電離パラメータ &
非平衡電離、幾何の不定性 \\
線輪郭の幅・分裂 & Voigt＋Zeeman & 温度・乱流・圧力・磁場 &
速度場の重畳、視線方向の混合 \\
\end{longtable}
}

\section*{この章で答える問い}\label{sec-line-applications-questions}
\addcontentsline{toc}{section}{この章で答える問い}

\markright{この章で答える問い}

\begin{itemize}
\tightlist
\item
  線の中心振動数のシフトから、どんな速度情報が読めるか
\item
  等価幅と成長曲線（curve of growth）から元素組成を測る手順は
\item
  線比（line ratio）から温度・密度を診断する原理は
\item
  高赤方偏移 QSO のスペクトルからは、宇宙の何が読めるか
\item
  系外惑星大気の線スペクトルからは、何がわかるか
\end{itemize}

\section*{到達目標}\label{sec-line-applications-goals}
\addcontentsline{toc}{section}{到達目標}

\markright{到達目標}

この章を読み終えると、読者は次のことができるようになる：

\begin{itemize}
\tightlist
\item
  線の中心シフトから視線速度と宇宙論的赤方偏移を測れる
\item
  等価幅と成長曲線（curve of
  growth）から元素組成を求める手順を組み立てられる
\item
  線比診断（{[}O III{]}、{[}S II{]} 等）から電子温度と密度を分けて測れる
\item
  第VII部の中心地図が、観測 → 物理量の \textbf{網羅的逆引き地図}
  であることを実感できる
\end{itemize}

\begin{center}\rule{0.5\linewidth}{0.5pt}\end{center}

\section{24.1 視線速度 ―
線中心のシフト}\label{sec-line-applications-velocity}

{[}本文目安：B1-B2{]}

線スペクトルの最も基本的応用は \textbf{視線速度測定}：観測される線中心
\(\nu_\text{obs}\) と実験室での既知の値 \(\nu_0\) の比較から、Doppler
シフト

\begin{equation}\protect\phantomsection\label{eq-doppler-velocity}{
\frac{v_\parallel}{c} = \frac{\nu_0 - \nu_\text{obs}}{\nu_0} \quad (\text{非相対論的})
}\end{equation}

これだけで多くの応用が可能：

\begin{itemize}
\tightlist
\item
  \textbf{連星系の軌道}：成分星の視線速度変化 → 質量比、傾き → 質量
\item
  \textbf{銀河の回転曲線}：21 cm 線や H\(\alpha\) 線の視線速度マップ →
  ダークマター存在の証拠
\item
  \textbf{宇宙論的赤方偏移}：\(z \equiv (\nu_0 - \nu_\text{obs})/\nu_\text{obs}\)
  で定義され、銀河の距離尺度（Hubble 法則）
\item
  \textbf{系外惑星検出}：恒星の視線速度の周期的揺らぎ → 惑星の存在
\end{itemize}

\begin{tcolorbox}[enhanced jigsaw, arc=.35mm, bottomrule=.15mm, breakable, colback=white, colframe=quarto-callout-warning-color-frame, left=2mm, leftrule=.75mm, opacityback=0, rightrule=.15mm, toprule=.15mm]
\begin{minipage}[t]{5.5mm}
\textcolor{quarto-callout-warning-color}{\faExclamationTriangle}
\end{minipage}%
\begin{minipage}[t]{\textwidth - 5.5mm}

\textbf{注意（観測）}：宇宙論的赤方偏移 \(z\) を
\(v_\text{recession}/c\) と書くのは \textbf{\(z \ll 1\)
の低赤方偏移近似のみ}
で正当化される。一般相対論的には、宇宙論的赤方偏移は時空計量のスケール因子比

\[
1 + z = \frac{a(t_\text{obs})}{a(t_\text{emit})}
\]

で定義され、単純な Doppler 速度ではない。\(z \gtrsim 0.1\)
では特殊相対論的 Doppler
公式や宇宙論モデル（\(\Lambda\)CDM）に依存した距離--赤方偏移関係を使う必要がある。本書は線スペクトル応用を主眼とし、宇宙論の詳細は専門書（Dodelson
\& Schmidt 等）を参照

\end{minipage}%
\end{tcolorbox}

\begin{tcolorbox}[enhanced jigsaw, arc=.35mm, bottomrule=.15mm, breakable, colback=white, colframe=quarto-callout-tip-color-frame, left=2mm, leftrule=.75mm, opacityback=0, rightrule=.15mm, toprule=.15mm]
\begin{minipage}[t]{5.5mm}
\textcolor{quarto-callout-tip-color}{\faLightbulb}
\end{minipage}%
\begin{minipage}[t]{\textwidth - 5.5mm}

\textbf{問い}：太陽の周りを地球と同じ軌道半径で公転する木星質量惑星は、太陽にどのくらいの視線速度ゆらぎを起こすか？

\textbf{短答}：\(v_\star = (M_p/M_\star) v_p \simeq (10^{-3}) \cdot 30\)
km/s = 30 m/s。Doppler 法での系外惑星検出の典型感度（HARPS, ESPRESSO
等）はサブ m/s 〜数 m/s
で、これは十分検出可能。地球質量同等の惑星検出には \(\sim 10\) cm/s
感度が必要で、これは現在の最先端

\textbf{もう一歩}：視線速度測定の精度は、線中心の決定精度に依存する。Voigt
プロファイルの中心を \(\sim 10^{-3} \cdot \Delta\nu_D\)
レベルで決定できれば、\textasciitilde10 m/s レベルの速度精度が得られる

\end{minipage}%
\end{tcolorbox}

\section{24.2 化学組成 ― 等価幅と成長曲線（curve of
growth）}\label{sec-line-applications-abundance}

{[}本文目安：B3{]}

恒星・星間ガス・QSO 視線方向の \textbf{化学組成}
を測る古典的方法。手順：

\begin{enumerate}
\def\labelenumi{\arabic{enumi}.}
\tightlist
\item
  \textbf{観測}：吸収線スペクトルから等価幅 \(W\)
  を測る（複数の線、複数の元素）
\item
  \textbf{成長曲線（curve of growth）}：第21章 §21.6 から
  \(W \to N\)（柱密度）の換算
\item
  \textbf{正規化}：水素柱密度を基準とし、\textbf{相対組成 \(N_X/N_H\)}
  を求める
\item
  \textbf{太陽組成との比}：{[}X/H{]}
  \(= \log_{10}(N_X/N_H) - \log_{10}(N_X/N_H)_\odot\)
\end{enumerate}

応用例：

\begin{itemize}
\tightlist
\item
  \textbf{銀河化学進化}：金属量 {[}Fe/H{]} と \(\alpha\) 元素過剰
  {[}\(\alpha\)/Fe{]} から、銀河中で星形成と超新星爆発が起きた
  \textbf{タイムスケール} が読める
\item
  \textbf{QSO 吸収系}：高赤方偏移の \textbf{金属吸収系}（C IV, Mg II, Fe
  II 等）→ 宇宙初期銀河の化学組成
\item
  \textbf{系外惑星宿主星}：金属量と惑星形成の相関 →
  「金属に富む恒星は巨大惑星を持ちやすい」という観測法則
\end{itemize}

\begin{tcolorbox}[enhanced jigsaw, arc=.35mm, bottomrule=.15mm, breakable, colback=white, colframe=quarto-callout-warning-color-frame, left=2mm, leftrule=.75mm, opacityback=0, rightrule=.15mm, toprule=.15mm]
\begin{minipage}[t]{5.5mm}
\textcolor{quarto-callout-warning-color}{\faExclamationTriangle}
\end{minipage}%
\begin{minipage}[t]{\textwidth - 5.5mm}

\textbf{注意（観測）}：等価幅から組成を求めるとき、成長曲線のどの領域に線があるかで誤差源が変わる：

\begin{itemize}
\tightlist
\item
  \textbf{線形領域}：誤差は等価幅の測定精度（連続光のフィット精度に依存）
\item
  \textbf{飽和領域}：\(N\) への感度が落ちる、ほぼ無用
\item
  \textbf{減衰翼領域}：感度回復、しかし衝突幅の理論計算誤差が新たな誤差源
\end{itemize}

優れた組成解析は \textbf{複数の線（弱・中・強）}
を組み合わせて、成長曲線全体を観測点で押さえる。

\end{minipage}%
\end{tcolorbox}

\section{24.3 温度・密度診断 ―
線比の原理}\label{sec-line-applications-diagnostics}

{[}本文目安：B3{]}

二つの異なる線の \textbf{強度比} \(I_1/I_2\) は、しばしば
\textbf{温度のみ} または \textbf{密度のみ}
に強く依存し、もう一方には鈍感。これを利用した \textbf{diagnostic line
ratio}：

\textbf{温度診断の例（電子温度）}：

\begin{itemize}
\tightlist
\item
  \textbf{{[}O III{]} \$\lambda\$4363 /
  \$\lambda\$5007}：両者の上準位エネルギーが大きく違うので、温度に強く依存
\item
  観測 \$\lambda\(4363/\)\lambda\$5007 比 → \(T_e \simeq 10^4\) K
\end{itemize}

\textbf{密度診断の例（電子密度）}：

\begin{itemize}
\tightlist
\item
  **{[}S II{]} \$\lambda\$6716 /
  \$\lambda\(6731**：両者の上準位エネルギーは近いが、各上準位の **臨界密度**（**critical density** ― 衝突脱励起率と自発放出率が拮抗する電子密度）が異なるため、両線の強度比は密度に応じて変わる。低密度極限と高密度極限の中間 (\)n\_e
  \sim 10\textsuperscript{2\$〜\$10}4\$ cm\(^{-3}\)) で診断感度が高い
\item
  観測 \$\lambda\(6716/\)\lambda\$6731 比 → \(n_e \sim 10^2\)〜\(10^4\)
  cm\(^{-3}\)
\end{itemize}

\begin{tcolorbox}[enhanced jigsaw, arc=.35mm, bottomrule=.15mm, bottomtitle=1mm, breakable, colback=white, colbacktitle=quarto-callout-note-color!10!white, colframe=quarto-callout-note-color-frame, coltitle=black, left=2mm, leftrule=.75mm, opacityback=0, opacitybacktitle=0.6, rightrule=.15mm, title=\textcolor{quarto-callout-note-color}{\faInfo}\hspace{0.5em}{ノート}, titlerule=0mm, toprule=.15mm, toptitle=1mm]

\textbf{用語}：これらは \textbf{禁制線}（第20章 §20.3、第22章
§22.5）。禁制線が \textbf{密度に敏感}
なのは、衝突脱励起が許容線に比べて短時間で起こるため。極低密度では禁制線が強く、高密度では衝突脱励起で弱まる。\textbf{「禁制線の存在自体が低密度の証拠」}
とも言える

\textbf{観測との対応}：HII 領域・惑星状星雲・AGN
の電子温度・密度を測る現代的標準ツールである。観測の波長分解能と感度の進歩により、\(\delta T_e \sim 100\)
K、\(\delta \log n_e \sim 0.1\) レベルの精度が達成されている。

\end{tcolorbox}

\section{24.4 電離度と化学種の分布 ―
イオン化バランス}\label{sec-line-applications-ionization}

{[}本文目安：B3{]}

同じ元素の異なる電離段階（例：H I vs H II、C IV vs C II）の
\textbf{存在比} から、電離度が読める（電離平衡を定める Saha
式と、それが連続放射に効く仕組みは\hyperref[sec-continuum-saha]{第19章
§19.5}）：

\begin{itemize}
\tightlist
\item
  \textbf{太陽光球}：水素は中性が圧倒的（\(T \simeq 6000\) K）
\item
  \textbf{HII 領域}：水素は完全電離（\(T \simeq 10^4\)
  K、紫外光子で電離）
\item
  \textbf{コロナ}：鉄が Fe IX〜Fe XX まで電離（\(T \sim 10^6\)
  K、衝突電離）
\end{itemize}

複数の電離段階の線を観測すれば、\textbf{電子温度} と
\textbf{電離パラメータ}（\(U \equiv n_\text{ionizing photon}/n_e\)）の組合せが決まる。これが活動銀河中心核（AGN）の
\textbf{電離度解析}（\textbf{BPT diagram}：{[}O III{]}/Hβ vs {[}N
II{]}/Hα の二次元プロット）の基礎。

\begin{tcolorbox}[enhanced jigsaw, arc=.35mm, bottomrule=.15mm, breakable, colback=white, colframe=quarto-callout-tip-color-frame, left=2mm, leftrule=.75mm, opacityback=0, rightrule=.15mm, toprule=.15mm]
\begin{minipage}[t]{5.5mm}
\textcolor{quarto-callout-tip-color}{\faLightbulb}
\end{minipage}%
\begin{minipage}[t]{\textwidth - 5.5mm}

\textbf{問い}：同じ温度の HII 領域でも、励起源（中心星）が暖かい O
型星か、冷たい B 型星かで、観測される線スペクトルはどう違うか？

\textbf{短答}：O 型星は \textbf{より高エネルギーの紫外光子}
を出すので、より深い電離 → {[}O III{]} が強く、低電離度の {[}O II{]}
が弱い。B 型星では逆。この \textbf{励起構造（excitation sequence）}
が銀河の高解像分光で銀河系内の様々な星雲を分類するのに使われる

\end{minipage}%
\end{tcolorbox}

\section{24.5 宇宙論的応用 ― QSO 吸収系と 21 cm
線}\label{sec-line-applications-cosmology}

{[}本文目安：B3-B4{]}

\textbf{QSO 吸収系}：高赤方偏移 QSO のスペクトルには、視線方向の
\textbf{多数の吸収系} が並ぶ：

\begin{itemize}
\tightlist
\item
  \textbf{Lyα 森（Lyman-alpha forest）}：低柱密度の中性水素 → 宇宙の
  \textbf{バリオン分布マップ}
\item
  \textbf{Damped Lyα system（DLA）}：高柱密度系 → 高赤方偏移銀河の前駆体
\item
  \textbf{金属吸収系}：C IV, Mg II 等 →
  宇宙初期の化学進化、銀河流出（galactic outflow）
\end{itemize}

これらの観測から、宇宙再電離の時期、銀河風の歴史、宇宙の重元素生成史が逆引きされる。

\textbf{21 cm 線の宇宙論}：再電離期 \(z \sim 6\)〜\(20\)
の中性水素分布を 21 cm 線（観測波長は \(0.21 \times (1+z)\) なので約
\(1.5\)〜\(4.4\) m）で見る。SKA（Square Kilometre
Array）の主要観測ターゲット。

\begin{tcolorbox}[enhanced jigsaw, arc=.35mm, bottomrule=.15mm, bottomtitle=1mm, breakable, colback=white, colbacktitle=quarto-callout-note-color!10!white, colframe=quarto-callout-note-color-frame, coltitle=black, left=2mm, leftrule=.75mm, opacityback=0, opacitybacktitle=0.6, rightrule=.15mm, title=\textcolor{quarto-callout-note-color}{\faInfo}\hspace{0.5em}{ノート}, titlerule=0mm, toprule=.15mm, toptitle=1mm]

\textbf{観測との対応}：第16章で見た
CMB（連続スペクトル側の宇宙論プローブ）と本節の QSO 吸収・21
cm（線スペクトル側）が、\textbf{宇宙物理学の二つの主柱}
として並列に効くことが、第VI部 ＋
第VII部の最大の成果である。一方は宇宙の温度の歴史、他方は構造形成・電離・化学進化の歴史を語る。

\end{tcolorbox}

\section{24.6 系外惑星大気 ―
透過分光と高分散分光}\label{sec-line-applications-exoplanet}

{[}本文目安：B3-B4{]}

系外惑星が中心星の前を通過するとき（transit）、惑星大気を
\textbf{透過した光} に大気成分の吸収線が刻まれる：

\begin{itemize}
\tightlist
\item
  \textbf{HST、JWST}：可視〜赤外分光で水蒸気・メタン・CO\(_2\)・酸素等の検出
\item
  \textbf{高分散分光}（ESPRESSO, IGRINS 等、\(R \sim 10^5\)）：個別線の
  Doppler シフトから惑星の速度成分まで読める
\end{itemize}

応用：

\begin{itemize}
\tightlist
\item
  大気組成 → 惑星形成過程の手がかり
\item
  バイオシグナチャー（O\(_2\)、CH\(_4\) の共存）の探査
\item
  地球型惑星の居住可能性（habitability）評価
\end{itemize}

\begin{tcolorbox}[enhanced jigsaw, arc=.35mm, bottomrule=.15mm, bottomtitle=1mm, breakable, colback=white, colbacktitle=quarto-callout-note-color!10!white, colframe=quarto-callout-note-color-frame, coltitle=black, left=2mm, leftrule=.75mm, opacityback=0, opacitybacktitle=0.6, rightrule=.15mm, title=\textcolor{quarto-callout-note-color}{\faInfo}\hspace{0.5em}{ノート}, titlerule=0mm, toprule=.15mm, toptitle=1mm]

\textbf{観測との対応}：系外惑星の大気観測は
\textbf{線スペクトル道具の代表的な現代応用}
である。本書の第20-23章で整備した枠組み（Einstein 係数 → 振動子強度 →
線輪郭フィッティング）が、すべて動員される。

\end{tcolorbox}

\section{24.7 まとめ ―
線スペクトル逆引き地図の完成}\label{sec-line-applications-summary}

{[}本文目安：B2{]}

第VII部全体を振り返ると、\textbf{線スペクトル → 物理量}
の逆引きは以下のようにマップ化される：

{\def\LTcaptype{none} % do not increment counter
\begin{longtable}[]{@{}llll@{}}
\toprule\noalign{}
中心地図の因子 & 観測量 & 逆引き先の物理量 & 章 \\
\midrule\noalign{}
\endhead
\bottomrule\noalign{}
\endlastfoot
\(\nu_0\) & 線振動数 & 原子・分子のエネルギー準位 & 21 \\
\(\nu_0\) シフト & 中心シフト & 視線速度、宇宙論赤方偏移 & 23 \\
\(f_{lu}\) & 線強度（と \(n_l\) から） & 遷移行列要素（量子力学） &
21 \\
\(n_l\) & 等価幅 & 柱密度、化学組成 & 20, 23 \\
\(n_l\) ratio & 線比 & 励起温度、電子温度・密度、電離度 & 20, 23 \\
\(\phi(\nu-\nu_0)\) 幅 & Doppler 幅 & キネティック温度、乱流速度 & 22 \\
\(\phi(\nu-\nu_0)\) 翼 & 翼の形 & 密度、圧力 & 22 \\
\(\phi(\nu-\nu_0)\) 分裂 & Zeeman 分裂 & 磁場 & 22 \\
\end{longtable}
}

これが第VII部の到達点 ―
\textbf{線スペクトル一本から、原子物理から銀河化学・宇宙論まで、多層の物理が逆引きできる地図}
が完成した。

\begin{tcolorbox}[enhanced jigsaw, arc=.35mm, bottomrule=.15mm, bottomtitle=1mm, breakable, colback=white, colbacktitle=quarto-callout-note-color!10!white, colframe=quarto-callout-note-color-frame, coltitle=black, left=2mm, leftrule=.75mm, opacityback=0, opacitybacktitle=0.6, rightrule=.15mm, title=\textcolor{quarto-callout-note-color}{\faInfo}\hspace{0.5em}{ノート}, titlerule=0mm, toprule=.15mm, toptitle=1mm]

\textbf{使えるようになった道具}（§24.1〜§24.6）：

\begin{itemize}
\tightlist
\item
  視線速度測定 ― 連星・回転曲線・宇宙論・系外惑星検出
\item
  化学組成決定 ― 成長曲線（curve of growth）+ 太陽組成基準
\item
  温度・密度診断 ― 線比（特に禁制線）
\item
  電離度解析 ― 多重電離段階の線比較
\item
  高赤方偏移宇宙の探査 ― QSO 吸収系、21 cm
\item
  系外惑星大気観測 ― 透過分光、高分散分光
\end{itemize}

\end{tcolorbox}

\begin{center}\rule{0.5\linewidth}{0.5pt}\end{center}

\section*{この章で何がわかったか}\label{sec-line-applications-summary-final}
\addcontentsline{toc}{section}{この章で何がわかったか}

\markright{この章で何がわかったか}

\textbf{中心地図に戻る・第VII部の総まとめ}

第VII部で、中心地図
\(\alpha_\nu = (\pi e^2/m_e c) f_{lu} n_l \phi(\nu-\nu_0)\)
の各因子が完全に逆引きされた：

\begin{itemize}
\tightlist
\item
  \(\nu_0\)（第22章）：原子・分子のエネルギー準位 ← 量子力学
\item
  \(f_{lu}\)（第22章）：振動子強度 ← 遷移行列要素
\item
  \(n_l\)（第21章）：占有数 ← LTE/non-LTE 統計力学
\item
  \(\phi(\nu-\nu_0)\)（第23章）：線輪郭 ←
  不確定性原理＋熱運動＋衝突＋磁場
\item
  観測応用（本章）：全部組み合わせて、物理量の網羅的読み取り
\end{itemize}

第I-VI部の連続スペクトルと並ぶ、\textbf{線スペクトル側の逆引き}
が完成した。第VIII部では、この二つの主柱（プランク分布と水素線吸収係数）を、より深い物理
― \textbf{電磁場の量子化と QED} ― から統一的に俯瞰する。

\section*{演習問題}\label{sec-line-applications-exercises}
\addcontentsline{toc}{section}{演習問題}

\markright{演習問題}

\textbf{模範解答}
は折りたたみを展開して確認できる（オンライン版）。まず自力で解いてから開くこと。

\subsection*{問題 24-1　視線速度と系外惑星の検出}\label{ex-24-1}
\addcontentsline{toc}{subsection}{問題 24-1　視線速度と系外惑星の検出}

{[}★ 難易度：☆☆ {]}

§24.1 の視線速度
\(\dfrac{v_\parallel}{c}=\dfrac{\nu_0-\nu_\mathrm{obs}}{\nu_0}\)
を系外惑星検出に使う。中心星と惑星は共通重心を回り、\(M_\star v_\star=M_p v_p\)。

\begin{enumerate}
\def\labelenumi{\arabic{enumi}.}
\tightlist
\item
  地球と同じ軌道半径（1 AU、公転速度 \(v_p\simeq30\)
  km/s）を回る木星質量惑星（\(M_p/M_\star\simeq10^{-3}\)）が中心星に起こす視線速度ゆらぎ
  \(v_\star\) を求めよ。
\item
  その \(v_\star\) が現在の Doppler 法（HARPS/ESPRESSO、感度サブ m/s〜数
  m/s）で検出可能か述べよ。地球質量の検出には何 cm/s
  感度が要るか概算せよ。
\item
  宇宙論的赤方偏移 \(z\) を単純な \(v/c\)
  と書けるのは低赤方偏移近似に限ることを、\(1+z=a(t_\mathrm{obs})/a(t_\mathrm{emit})\)
  に触れて述べよ。
\end{enumerate}

\textbf{関連}：\hyperref[sec-line-applications-velocity]{§24.1
視線速度}／Doppler
シフトは\hyperref[sec-line-shape-doppler]{§23.2}、宇宙膨張と赤方偏移は\hyperref[ex-16-2]{演習
16-2}。

\section{模範解答（問題
24-1）}\label{ux6a21ux7bc4ux89e3ux7b54ux554fux984c-24-1}

\textbf{(1)} \(v_\star=（M_p/M_\star)v_p=10^{-3}\times30\) km/s
\(=0.03\) km/s \(=30\) m/s。

\textbf{(2)} 30 m/s は HARPS/ESPRESSO（サブ m/s〜数
m/s）で十分検出可能。地球質量惑星なら
\(M_p/M_\star\simeq3\times10^{-6}\) ゆえ
\(v_\star\simeq3\times10^{-6}\times30\) km/s \(\simeq0.09\) m/s
\(\simeq10\) cm/s 感度が必要で、これは現在の最先端。

\textbf{(3)} \(z\simeq v/c\) は \(z\ll1\)
でのみ成立。一般には宇宙論的赤方偏移は時空のスケール因子比
\(1+z=a(t_\mathrm{obs})/a(t_\mathrm{emit})\)
で定義され、単純な後退速度ではない。\(z\gtrsim0.1\) では特殊相対論的
Doppler 公式や宇宙論モデル依存の距離-赤方偏移関係が要る。

\textbf{答え}：\(v_\star=30\) m/s（検出可能）。地球質量には \(\sim10\)
cm/s 感度。\(z=v/c\) は低赤方偏移近似のみ、本来は
\(1+z=a_\mathrm{obs}/a_\mathrm{emit}\)。

\subsection*{問題 24-2　線比による温度・密度診断}\label{ex-24-2}
\addcontentsline{toc}{subsection}{問題 24-2　線比による温度・密度診断}

{[}★ 難易度：☆☆ {]}

§24.3 の禁制線比による診断。{[}O III{]}
\$\lambda\(4363/\)\lambda\$5007（温度）と {[}S II{]}
\$\lambda\(6716/\)\lambda\$6731（密度）。

\begin{enumerate}
\def\labelenumi{\arabic{enumi}.}
\tightlist
\item
  {[}O III{]} \$\lambda\(4363/\)\lambda\$5007
  比が電子温度に強く依存する理由を、両線の上準位エネルギー差から述べよ。
\item
  {[}S II{]} \$\lambda\(6716/\)\lambda\$6731
  比が電子密度に依存する理由を、両上準位の臨界密度の違いから述べよ。
\item
  一つの線比が「温度のみ」または「密度のみ」に依存し、もう一方に鈍感だと、なぜ温度と密度を分離して測れるのか述べよ。
\end{enumerate}

\textbf{関連}：\hyperref[sec-line-applications-diagnostics]{§24.3
温度・密度診断}／禁制線と臨界密度は\hyperref[ex-20-2]{演習
20-2}、Boltzmann 占有は\hyperref[ex-7-4]{演習 7-4}。

\section{模範解答（問題
24-2）}\label{ux6a21ux7bc4ux89e3ux7b54ux554fux984c-24-2}

\textbf{(1)} {[}O III{]} \$\lambda\$4363 と \$\lambda\$5007
は上準位のエネルギーが大きく異なる。上準位の占有比は Boltzmann 因子
\(e^{-\Delta E/k_BT_e}\) で決まり、エネルギー差 \(\Delta E\)
が大きいほど温度に敏感。両線の強度比はこの占有比を反映し、電子温度
\(T_e\) に強く依存する。

\textbf{(2)} {[}S II{]} \$\lambda\$6716 と
\$\lambda\(6731 は上準位エネルギーが近い（温度に鈍感）が、各上準位の臨界密度が異なる。電子密度が両者の臨界密度の中間（\)n\_e\sim10\textsuperscript{2\$--\$10}4\$
cm\(^{-3}\)）にあると、密度に応じて衝突脱励起の効き方が変わり、強度比が密度の関数になる。

\textbf{(3)} {[}O III{]} 比は温度のみ、{[}S II{]}
比は密度のみに主に依存し互いに鈍感なので、二つの線比を独立な方程式として連立すれば、温度と密度を別々に解ける。これが星雲の電子温度・密度を分離して測る標準法。

\textbf{答え}：{[}O III{]} は上準位エネルギー差大→温度敏感、{[}S II{]}
は臨界密度差→密度敏感。互いに鈍感ゆえ温度と密度を分離して測れる。

\subsection*{問題 24-3　高赤方偏移宇宙 ― Lyα 森と 21 cm}\label{ex-24-3}
\addcontentsline{toc}{subsection}{問題 24-3　高赤方偏移宇宙 ― Lyα 森と
21 cm}

{[}★ 難易度：☆☆ {]}

§24.5 の宇宙論的応用。QSO 吸収系（Ly\(\alpha\)
森・DLA・金属吸収系）と再電離期の 21 cm 線。

\begin{enumerate}
\def\labelenumi{\arabic{enumi}.}
\tightlist
\item
  Ly\(\alpha\)
  森・DLA・金属吸収系が、それぞれ宇宙の何（バリオン分布・銀河前駆体・化学進化）を語るか述べよ。
\item
  再電離期 \(z=7\) の中性水素 21 cm 線（静止波長 0.211
  m）の観測波長を求めよ。\(z=6\) と \(z=20\) では何 m か。
\item
  第16章の CMB（連続スペクトル側の宇宙論プローブ）と本問の
  Ly\(\alpha\)/21
  cm（線スペクトル側）が、宇宙物理の二つの主柱として何の歴史を相補的に語るか述べよ。
\end{enumerate}

\textbf{関連}：\hyperref[sec-line-applications-cosmology]{§24.5
宇宙論的応用}／CMB
は\hyperref[ux7b2c16ux7ae0-ux5b87ux5b99ux30deux30a4ux30afux30edux6ce2ux80ccux666fux653eux5c04]{第16章}（\hyperref[ux7b2c16ux7ae0-ux5b87ux5b99ux30deux30a4ux30afux30edux6ce2ux80ccux666fux653eux5c04]{第VI部
演習}）、21 cm
の起源は\hyperref[sec-where-lines-21cm]{§20.4}（\hyperref[ex-20-1]{演習
20-1}）。

\section{模範解答（問題
24-3）}\label{ux6a21ux7bc4ux89e3ux7b54ux554fux984c-24-3}

\textbf{(1)} Ly\(\alpha\)
森：低柱密度の中性水素による多数の吸収線＝宇宙のバリオン分布マップ。DLA（damped
Ly\(\alpha\)）：高柱密度系＝高赤方偏移銀河の前駆体。金属吸収系（C IV, Mg
II 等）：宇宙初期の化学進化・銀河流出の記録。

\textbf{(2)} \(\lambda_\mathrm{obs}=(1+z)\times0.211\)
m。\(z=7\)：\(8\times0.211=1.69\) m。\(z=6\)：\(7\times0.211=1.48\)
m。\(z=20\)：\(21\times0.211=4.43\) m。再電離期は約 1.5〜4.4 m（SKA
の標的帯）。

\textbf{(3)} CMB
は連続スペクトル（黒体）で宇宙の\textbf{温度の歴史}（熱化・膨張）を語る。Ly\(\alpha\)/21
cm
は線スペクトルで\textbf{構造形成・電離・化学進化の歴史}を語る。両者は宇宙物理学の二つの主柱として相補的に、初期宇宙から現在までの異なる側面を逆引きさせる。

\textbf{答え}：Ly\(\alpha\)
森＝バリオン分布、DLA＝銀河前駆体、金属系＝化学進化。21
cm(\(z=7\))＝1.69 m（\(z=6\):1.48、\(z=20\):4.43
m）。CMB＝温度史、線＝構造・電離・化学史で相補的。

\section*{さらに学ぶための参考文献}\label{sec-line-applications-further}
\addcontentsline{toc}{section}{さらに学ぶための参考文献}

\markright{さらに学ぶための参考文献}

\begin{itemize}
\tightlist
\item
  Asplund, ``New Light on Stellar Abundance Analyses,'' ARA\&A 43 (2005)
  481 --- 恒星化学組成の現代手法
\item
  Osterbrock \& Ferland, \emph{Astrophysics of Gaseous Nebulae and
  Active Galactic Nuclei} (Univ. Science Books, 2nd ed., 2006) ---
  線比診断の標準書
\item
  Prochaska, ``Quasar Absorption Line Spectroscopy'' レビュー --- QSO
  吸収系の網羅的解説
\item
  Madhusudhan, ``Exoplanetary Atmospheres,'' ARA\&A 57 (2019) 617 ---
  系外惑星大気観測のレビュー
\item
  Dalgarno \& Layzer (eds.), \emph{Spectroscopy of Astrophysical
  Plasmas} (Cambridge, 1987) --- 線比診断の体系的記述
\end{itemize}

\part{第VIII部　QEDからの俯瞰 ― 連続と線を結ぶ見通し}

\chapter{第25章　電磁場を量子化するとは何か ―
直観的描像}\label{ux7b2c25ux7ae0-ux96fbux78c1ux5834ux3092ux91cfux5b50ux5316ux3059ux308bux3068ux306fux4f55ux304b-ux76f4ux89b3ux7684ux63cfux50cf}

\textbf{この章の主題}：電磁場の各モードを調和振動子として扱う描像、光子の生成消滅、真空状態と零点エネルギー、黒体放射を「電磁場の熱状態」として書き直す。本章は
QED
への入口として、\textbf{厳密な場の理論の構築ではなく、観測スペクトルを統一的に理解するための直観的描像}
を獲得することを目的とする。

\begin{tcolorbox}[enhanced jigsaw, arc=.35mm, bottomrule=.15mm, bottomtitle=1mm, breakable, colback=white, colbacktitle=quarto-callout-important-color!10!white, colframe=quarto-callout-important-color-frame, coltitle=black, left=2mm, leftrule=.75mm, opacityback=0, opacitybacktitle=0.6, rightrule=.15mm, title=\textcolor{quarto-callout-important-color}{\faExclamation}\hspace{0.5em}{重要}, titlerule=0mm, toprule=.15mm, toptitle=1mm]

\textbf{本部の位置づけ ― 「見通し」レベル}

第VIII部は本書の最終章ではあるが、第I〜VII
部までで到達できた理解を「もう一段高い視点から眺める」ためのものである。場の量子論を厳密に構築するのではなく、

\begin{itemize}
\tightlist
\item
  Einstein 係数の関係 (\(2h\nu^3/c^2\)
  がプランク関数と共通の前因子になる理由)
\item
  自然幅
\item
  真空揺らぎ
\item
  選択則
\end{itemize}

の四点を \textbf{見通し}
として整理することを目的とする。厳密な議論は専門書に譲り、本書では「観測される連続スペクトルと線スペクトルが、同じ電磁場--物質相互作用の二側面である」という統一像を読者が持てる程度で十分とする。

\end{tcolorbox}

\section*{この章の中心地図}\label{sec-em-field-quantization-map}
\addcontentsline{toc}{section}{この章の中心地図}

\markright{この章の中心地図}

\textbf{プランク関数（連続スペクトルの地図）}

\begin{equation}\protect\phantomsection\label{eq-planck-function}{
B_\nu(T) = \frac{2 h \nu^3}{c^2} \cdot \frac{1}{e^{h\nu/kT} - 1}
}\end{equation}

{\def\LTcaptype{none} % do not increment counter
\begin{longtable}[]{@{}lll@{}}
\toprule\noalign{}
因子 & 物理的意味 & 結びつく部 \\
\midrule\noalign{}
\endhead
\bottomrule\noalign{}
\endlastfoot
\(B_\nu\) & 観測量（比強度） & 第II部 \\
\(1/(e^{h\nu/kT}-1)\) & 熱平衡・光子統計 & 第III部 \\
\(h\nu\) & エネルギー量子 & 第IV部 \\
\(\nu^3/c^2\) & 電磁場のモード密度 & 第V部 \\
黒体からのずれ & 現実の天体 & 第VI部 \\
\end{longtable}
}

\begin{tcolorbox}[enhanced jigsaw, arc=.35mm, bottomrule=.15mm, breakable, colback=white, colframe=quarto-callout-tip-color-frame, left=2mm, leftrule=.75mm, opacityback=0, rightrule=.15mm, toprule=.15mm]

\textbf{地図の使い方}：各部の冒頭でこの式が再掲され、その部で扱う因子が強調表示される。読者は自分が今どこを学んでいるかをこの式の上で確認できる。

\end{tcolorbox}

\textbf{水素線スペクトルの地図 ― 二層構造}

\textbf{第一層：イコニック・アンカー（リュードベリ公式）}

\begin{equation}\protect\phantomsection\label{eq-rydberg}{
\frac{1}{\lambda} = R_\mathrm{H} \left( \frac{1}{n_1^2} - \frac{1}{n_2^2} \right)
}\end{equation}

線スペクトルの「骨格」を一目で示す。\(R_\mathrm{H}\)
は基本定数の組合せ、\(1/n^2\)
は束縛状態のエネルギー、離散性が「線が存在する理由」を語る。

\textbf{第二層：構造的マップ（Hα を標準例とした線吸収係数）}

\begin{equation}\protect\phantomsection\label{eq-line-absorption}{
\alpha_\nu = \frac{\pi e^2}{m_e c} \, f_{lu} \, n_l \, \phi(\nu - \nu_0)
}\end{equation}

{\def\LTcaptype{none} % do not increment counter
\begin{longtable}[]{@{}lll@{}}
\toprule\noalign{}
因子 & 物理的意味 & 対応する章 \\
\midrule\noalign{}
\endhead
\bottomrule\noalign{}
\endlastfoot
\(\alpha_\nu\) & 観測量（線吸収係数） & 第21章 \\
\(\pi e^2 / m_e c\) & 古典極限の線強度 & 第22章＋第VIII部 \\
\(f_{lu}\) & 振動子強度 & 第22章 \\
\(n_l\) & 下準位の占有数 & 第21章＋Part III \\
\(\phi(\nu - \nu_0)\) & 線輪郭（自然・Doppler・Voigt） & 第23章 \\
\(\nu_0\) & 線位置（リュードベリ） & 第22章 \\
\end{longtable}
}

\textbf{適用範囲（重要）}：この式の前因子 \(\pi e^2/m_e c\) と振動子強度
\(f_{lu}\) は、古典 Lorentz 振動子に対応する
\textbf{電気双極子（E1）許容線} に対するものである。21 cm
線（基底状態超微細構造の \textbf{磁気双極子 M1 遷移}）や
\textbf{禁制線（E2・M1 など）} では、この前因子・\(f_{lu}\)
をそのまま使えない。これらの線では、測定された自発放出率 \(A_{ul}\) から
Einstein
関係を介して有効遷移強度（\(B_{lu}\)）に読み替える（式~\ref{eq-line-absorption}
の構造は保ったまま、強度因子だけを差し替える）。

\begin{tcolorbox}[enhanced jigsaw, arc=.35mm, bottomrule=.15mm, breakable, colback=white, colframe=quarto-callout-tip-color-frame, left=2mm, leftrule=.75mm, opacityback=0, rightrule=.15mm, toprule=.15mm]

\textbf{地図の使い方}：第VII部の各章はこの式の異なる因子を解き明かす。プランク関数地図と並べて読めば、本書を貫く二つの逆引きが同時に把握できる。第VIII部の
QED で、両者は前因子 \(2h\nu^3/c^2\)
を介して統一される（式~\ref{eq-einstein-A} 参照）。

\end{tcolorbox}

\begin{tcolorbox}[enhanced jigsaw, arc=.35mm, bottomrule=.15mm, bottomtitle=1mm, breakable, colback=white, colbacktitle=quarto-callout-note-color!10!white, colframe=quarto-callout-note-color-frame, coltitle=black, left=2mm, leftrule=.75mm, opacityback=0, opacitybacktitle=0.6, rightrule=.15mm, title=\textcolor{quarto-callout-note-color}{\faInfo}\hspace{0.5em}{ノート}, titlerule=0mm, toprule=.15mm, toptitle=1mm]

\textbf{方針}：本章で初めて、二つの中心地図を \textbf{同時に}
掲げる。第VIII部の目的が、両者が同じ枠組み（電磁場と物質の量子化）から現れることを示すことにあるためである。

\end{tcolorbox}

\section*{この章で答える問い}\label{sec-em-field-quantization-questions}
\addcontentsline{toc}{section}{この章で答える問い}

\markright{この章で答える問い}

\begin{itemize}
\tightlist
\item
  なぜ電磁場を量子化する必要があるのか（古典論との不一致は何だったか）
\item
  各モードを「調和振動子」と見なすと、何が新しく見えてくるのか
\item
  「光子」とは、量子化された電磁場の中で具体的に何を指すのか
\item
  真空状態が「空っぽ」ではないというのは、観測の言葉で何を意味するのか
\item
  黒体放射は、量子化された電磁場の言葉でどう書き直されるのか
\end{itemize}

\section*{到達目標}\label{sec-em-field-quantization-goals}
\addcontentsline{toc}{section}{到達目標}

\markright{到達目標}

この章を読み終えると、読者は次のことができるようになる：

\begin{itemize}
\tightlist
\item
  電磁場の各モードを調和振動子とみなす描像を、自分の言葉で説明できる
\item
  光子の生成・消滅・数演算子の意味を、観測スペクトルとの対応で理解できる
\item
  真空揺らぎの存在を、観測可能な帰結（自発放出・Casimir
  効果）と結びつけられる
\end{itemize}

\begin{center}\rule{0.5\linewidth}{0.5pt}\end{center}

\section{25.1
なぜ電磁場の量子化が必要だったか}\label{sec-em-field-quantization-why}

{[}本文目安：M1{]}

ここまで本書では、電磁場と量子論を \textbf{別々に} 扱ってきた：

\begin{itemize}
\tightlist
\item
  \textbf{電磁場（古典）}：第V部で Maxwell 方程式 → 平面波・モード密度
\item
  \textbf{量子論}：第IV部でエネルギー量子化（プランク仮説）、第VII部で原子の波動関数
\end{itemize}

しかし、第IV-VII部のあちこちで「\textbf{古典電磁気の枠を一歩出ないと説明できない}」現象に出会ってきた：

{\def\LTcaptype{none} % do not increment counter
\begin{longtable}[]{@{}
  >{\raggedright\arraybackslash}p{(\linewidth - 4\tabcolsep) * \real{0.3333}}
  >{\raggedright\arraybackslash}p{(\linewidth - 4\tabcolsep) * \real{0.3333}}
  >{\raggedright\arraybackslash}p{(\linewidth - 4\tabcolsep) * \real{0.3333}}@{}}
\toprule\noalign{}
\begin{minipage}[b]{\linewidth}\raggedright
現象
\end{minipage} & \begin{minipage}[b]{\linewidth}\raggedright
出てきた章
\end{minipage} & \begin{minipage}[b]{\linewidth}\raggedright
古典電磁気では説明できない理由
\end{minipage} \\
\midrule\noalign{}
\endhead
\bottomrule\noalign{}
\endlastfoot
\textbf{エネルギー量子化} \(h\nu\) & 第10章 & 古典波には離散性がない \\
\textbf{自発放出} \(A_{ul}\) & 第11章 & 古典電磁気では「ゼロ放射場 →
ゼロ放射」、減衰なし \\
\textbf{自然幅} & 第23章 &
励起準位の有限寿命を古典電子モデルで「決め打ち」していた \\
\textbf{光子の Bose 統計} & 第8章 & 古典場は数えられる粒子ではない \\
\textbf{モードあたり光子数の Planck 分布} & 第8章 & 「モード =
振動子」の量子化が暗黙の前提 \\
\end{longtable}
}

これらが \textbf{同じ枠組み}
から出るためには、\textbf{電磁場そのものを量子化}
する必要がある。これが第VIII部の主題である。

\begin{tcolorbox}[enhanced jigsaw, arc=.35mm, bottomrule=.15mm, bottomtitle=1mm, breakable, colback=white, colbacktitle=quarto-callout-note-color!10!white, colframe=quarto-callout-note-color-frame, coltitle=black, left=2mm, leftrule=.75mm, opacityback=0, opacitybacktitle=0.6, rightrule=.15mm, title=\textcolor{quarto-callout-note-color}{\faInfo}\hspace{0.5em}{ノート}, titlerule=0mm, toprule=.15mm, toptitle=1mm]

\textbf{方針}：本章では、量子化の
\textbf{厳密な手続き}（正準量子化、生成消滅演算子の代数、Fock
空間構成）は最小限に留め、\textbf{観測スペクトルとの対応関係}
に注力する。厳密な構成は Loudon、Mandel \& Wolf、Peskin \& Schroeder
等の専門書を参照してほしい。

\end{tcolorbox}

\section{25.2 各モードは調和振動子 ―
直観的描像}\label{sec-em-field-quantization-mode-oscillator}

{[}本文目安：M1{]}

第14章 §14.1 で見たように、空洞中の電磁場は \textbf{離散的モード}
\(\{(\mathbf{k}, \sigma)\}\) の集合として表せる（\(\sigma\)
は偏光自由度）。

第VIII部の核心は次の \textbf{一行の置き換え}：

\begin{quote}
\textbf{各モードを、量子化された調和振動子に置き換える}
\end{quote}

この一手で、第I〜VII部のすべての断片が一つの構造に収まる。各モードの
Hamiltonian は

\begin{equation}\protect\phantomsection\label{eq-mode-hamiltonian}{
\hat{H}_\text{mode}(\mathbf{k}, \sigma) = \hbar\omega_\mathbf{k} \left(\hat{a}^\dagger_{\mathbf{k}\sigma} \hat{a}_{\mathbf{k}\sigma} + \frac{1}{2}\right)
}\end{equation}

ここで
\(\hat{a}^\dagger_{\mathbf{k}\sigma}, \hat{a}_{\mathbf{k}\sigma}\) は
\textbf{生成・消滅演算子}、\(\omega_\mathbf{k} = c|\mathbf{k}|\)
は古典分散関係（第13章 §13.3）。

{\def\LTcaptype{none} % do not increment counter
\begin{longtable}[]{@{}
  >{\raggedright\arraybackslash}p{(\linewidth - 4\tabcolsep) * \real{0.3333}}
  >{\raggedright\arraybackslash}p{(\linewidth - 4\tabcolsep) * \real{0.3333}}
  >{\raggedright\arraybackslash}p{(\linewidth - 4\tabcolsep) * \real{0.3333}}@{}}
\toprule\noalign{}
\begin{minipage}[b]{\linewidth}\raggedright
演算子
\end{minipage} & \begin{minipage}[b]{\linewidth}\raggedright
役割
\end{minipage} & \begin{minipage}[b]{\linewidth}\raggedright
物理
\end{minipage} \\
\midrule\noalign{}
\endhead
\bottomrule\noalign{}
\endlastfoot
\(\hat{a}^\dagger_{\mathbf{k}\sigma}\) & 「モード
\((\mathbf{k}, \sigma)\) に光子を1個加える」 & 放出 \\
\(\hat{a}_{\mathbf{k}\sigma}\) & 「モード \((\mathbf{k}, \sigma)\)
から光子を1個取り去る」 & 吸収 \\
\(\hat{n}_{\mathbf{k}\sigma} = \hat{a}^\dagger \hat{a}\) &
「そのモードにいる光子数」 & 観測量 \\
\end{longtable}
}

調和振動子の励起準位 \(|n\rangle\) がそのまま \textbf{「そのモードに
\(n\) 個の光子が存在する状態」} と読み替えられる。

\begin{tcolorbox}[enhanced jigsaw, arc=.35mm, bottomrule=.15mm, bottomtitle=1mm, breakable, colback=white, colbacktitle=quarto-callout-note-color!10!white, colframe=quarto-callout-note-color-frame, coltitle=black, left=2mm, leftrule=.75mm, opacityback=0, opacitybacktitle=0.6, rightrule=.15mm, title=\textcolor{quarto-callout-note-color}{\faInfo}\hspace{0.5em}{ノート}, titlerule=0mm, toprule=.15mm, toptitle=1mm]

\textbf{用語の対応関係}：

\begin{itemize}
\tightlist
\item
  第10章「エネルギーが \(h\nu\) の整数倍に量子化される」 →
  「モード励起準位が \(n h\nu\)」
\item
  第8章「光子の Bose 統計」 → 「各モードの調和振動子の熱状態」
\item
  第14章「モード密度 \(g(\nu)\)」 → 「調和振動子の振動数ごとの個数」
\end{itemize}

第I〜VII部で別々の道具と思っていたものが、すべて \textbf{「各モード ＝
調和振動子」} の一行で接続される。これが第VIII部の \textbf{最大の効用}
である。

\end{tcolorbox}

\section{25.3 真空状態と零点エネルギー ―
「空ではない」真空}\label{sec-em-field-quantization-vacuum}

{[}本文目安：M1{]}

すべてのモードが基底状態 \(|0\rangle\) にある状態を
\textbf{真空}（\textbf{vacuum}）と呼ぶ。古典電磁気では「電磁場 =
0」と言うところだが、量子論ではそうではない。

式~\ref{eq-mode-hamiltonian} に \(|0\rangle\)
を作用させると、エネルギーは

\begin{equation}\protect\phantomsection\label{eq-zero-point-energy}{
E_0 = \sum_{\mathbf{k}, \sigma} \frac{1}{2} \hbar \omega_\mathbf{k}
}\end{equation}

これが \textbf{零点エネルギー}（\textbf{zero-point
energy}）。形式的には発散するが、観測可能な物理量の差として現れる帰結は
\textbf{有限} であり、実際に観測されている。

真空が「空ではない」ことの \textbf{観測可能な帰結}：

\begin{enumerate}
\def\labelenumi{\arabic{enumi}.}
\tightlist
\item
  \textbf{自発放出}（\textbf{spontaneous emission}, 第11章 §11.2）

  \begin{itemize}
  \tightlist
  \item
    古典論では「ゼロ放射場 → 放出ゼロ」のはず
  \item
    量子論では真空揺らぎとの相互作用が遷移を誘起し、放出が生じる
  \item
    \textbf{次章 §26.2 で QED から \(A_{ul}\) を導出}
  \end{itemize}
\item
  \textbf{線の自然幅}（第23章 §23.1）

  \begin{itemize}
  \tightlist
  \item
    励起準位の有限寿命 → 不確定性関係でエネルギー幅 \(\Gamma = A_{ul}\)
  \item
    真空揺らぎの存在の \textbf{直接的観測結果}
  \end{itemize}
\item
  \textbf{Casimir 効果}（\textbf{Casimir effect}）

  \begin{itemize}
  \tightlist
  \item
    二枚の金属板の間で許されるモードが減る →
    板間の零点エネルギーが板外より小さい
  \item
    結果として板間に \textbf{引力} \(F \propto 1/L^4\)
  \item
    1948 年 Casimir 予言、1958 年 Sparnaay の初観測、1997 年 Lamoreaux
    の高精度検証
  \end{itemize}
\end{enumerate}

\begin{tcolorbox}[enhanced jigsaw, arc=.35mm, bottomrule=.15mm, breakable, colback=white, colframe=quarto-callout-tip-color-frame, left=2mm, leftrule=.75mm, opacityback=0, rightrule=.15mm, toprule=.15mm]
\begin{minipage}[t]{5.5mm}
\textcolor{quarto-callout-tip-color}{\faLightbulb}
\end{minipage}%
\begin{minipage}[t]{\textwidth - 5.5mm}

\textbf{問い}：「真空に揺らぎがある」と「真空は粒子で満ちている」は、同じことか違うことか？

\textbf{短答}：違う。真空状態 \(|0\rangle\) は厳密に光子 0
個の状態。観測される「揺らぎ」は、\(\hat{a}\) と \(\hat{a}^\dagger\)
の非可換性 \([\hat{a}, \hat{a}^\dagger] = 1\)
から来る、\textbf{電場や磁場の期待値の分散}
\(\langle (\Delta E)^2 \rangle \neq 0\) である。

\textbf{もう一歩}：粒子描像 \(|n\rangle\) と場の描像
\(\hat{E}, \hat{B}\)
は、ハイゼンベルクの不確定性関係で双対的に結ばれている。粒子数を定めると場の値が揺らぎ、場の値を定めると粒子数が揺らぐ。

\end{minipage}%
\end{tcolorbox}

\begin{tcolorbox}[enhanced jigsaw, arc=.35mm, bottomrule=.15mm, bottomtitle=1mm, breakable, colback=white, colbacktitle=quarto-callout-note-color!10!white, colframe=quarto-callout-note-color-frame, coltitle=black, left=2mm, leftrule=.75mm, opacityback=0, opacitybacktitle=0.6, rightrule=.15mm, title=\textcolor{quarto-callout-note-color}{\faInfo}\hspace{0.5em}{ノート}, titlerule=0mm, toprule=.15mm, toptitle=1mm]

\textbf{観測との対応}：宇宙論との交差点 ―
真空のエネルギー密度を素直に評価すると、観測される
\textbf{暗黒エネルギー密度} よりおよそ \(10^{120}\)
倍も大きくなってしまう。これは
\textbf{宇宙定数問題}（\textbf{cosmological constant
problem}）と呼ばれ、現代物理学最大の未解決問題のひとつである。本書の射程は超えるが、第VIII部で扱った真空のエネルギーが宇宙論的観測とも接続する例として重要である。

\end{tcolorbox}

\section{25.4
黒体放射を電磁場の熱状態として書き直す}\label{sec-em-field-quantization-blackbody}

{[}本文目安：M1{]}

プランク分布の最終解釈を、第VIII部の言葉で書き直す。

第8章では、プランク分布を「光子気体の Bose
統計」として導いた。第10章では「振動子の熱平均」として導いた。第14章では「モードあたり熱励起
×
モード密度」として導いた。\textbf{これらはすべて、同じ一つの構造の異なる側面である}：

\textbf{電磁場の各モード（調和振動子）の熱密度行列の表現}

すなわち、温度 \(T\) の熱浴と平衡にある電磁場の \textbf{熱密度行列}

\begin{equation}\protect\phantomsection\label{eq-thermal-density-matrix}{
\hat{\rho}_\text{th} = \frac{1}{Z} \exp(-\hat{H}/k_B T), \quad \hat{H} = \sum_{\mathbf{k}, \sigma} \hbar\omega_\mathbf{k} \hat{a}^\dagger \hat{a}
}\end{equation}

の下で、各モードの \textbf{平均光子数} は

\begin{equation}\protect\phantomsection\label{eq-thermal-photon-number}{
\bar{n}(\nu, T) = \text{tr}(\hat{\rho}_\text{th} \hat{n}_{\mathbf{k}\sigma}) = \frac{1}{e^{h\nu/k_B T} - 1}
}\end{equation}

これに \textbf{モード密度}
\(g(\nu) = 8\pi\nu^2/c^3\)（第14章で逆引き）と \textbf{光子エネルギー}
\(h\nu\) を掛けて、

\begin{equation}\protect\phantomsection\label{eq-thermal-energy-density}{
u_\nu(T) = g(\nu) \cdot h\nu \cdot \bar{n}(\nu, T) = \frac{8\pi h\nu^3}{c^3} \cdot \frac{1}{e^{h\nu/k_B T} - 1}
}\end{equation}

そして比強度 \(B_\nu = c u_\nu / 4\pi\)（第4章 §4.3 と第14章
§14.5）から、

\begin{equation}\protect\phantomsection\label{eq-planck-quantum-field}{
\boxed{
B_\nu(T) = \frac{2h\nu^3}{c^2} \cdot \frac{1}{e^{h\nu/k_B T} - 1}
}
}\end{equation}

これは第14章 式~\ref{eq-planck-fully-derived}
と同じ式だが、\textbf{「電磁場の熱状態の表現」}
という単一の物理的言葉でまとまっている。

\subsection{プランク関数地図の最終解釈}\label{ux30d7ux30e9ux30f3ux30afux95a2ux6570ux5730ux56f3ux306eux6700ux7d42ux89e3ux91c8}

中心地図の各因子が、第VIII部の言葉で：

{\def\LTcaptype{none} % do not increment counter
\begin{longtable}[]{@{}
  >{\raggedright\arraybackslash}p{(\linewidth - 2\tabcolsep) * \real{0.5000}}
  >{\raggedright\arraybackslash}p{(\linewidth - 2\tabcolsep) * \real{0.5000}}@{}}
\toprule\noalign{}
\begin{minipage}[b]{\linewidth}\raggedright
因子
\end{minipage} & \begin{minipage}[b]{\linewidth}\raggedright
第VIII部での意味
\end{minipage} \\
\midrule\noalign{}
\endhead
\bottomrule\noalign{}
\endlastfoot
\(2h\nu^3/c^2\) & モード密度 \(\nu^2/c^3\) × 偏光自由度 2 ×
エネルギー量子 \(h\nu\) \\
\(1/(e^{h\nu/k_BT}-1)\) &
各モードを調和振動子の熱状態とみなしたときの平均光子数 \\
全体 \(B_\nu(T)\) & \textbf{「電磁場の熱状態」のスペクトル密度} \\
\end{longtable}
}

これで連続スペクトル側の中心地図は、\textbf{「電磁場の熱状態」というひとつの構造}
に整理された。

\begin{tcolorbox}[enhanced jigsaw, arc=.35mm, bottomrule=.15mm, bottomtitle=1mm, breakable, colback=white, colbacktitle=quarto-callout-note-color!10!white, colframe=quarto-callout-note-color-frame, coltitle=black, left=2mm, leftrule=.75mm, opacityback=0, opacitybacktitle=0.6, rightrule=.15mm, title=\textcolor{quarto-callout-note-color}{\faInfo}\hspace{0.5em}{ノート}, titlerule=0mm, toprule=.15mm, toptitle=1mm]

\textbf{観測との対応}：第16章で見た CMB の完璧な黒体性は、宇宙が
\textbf{「電磁場の熱状態」を経験した}
ことの直接的観測である。これが第I-VIII部全体の到達点であり、\textbf{観測スペクトル
→ 電磁場の量子化された熱状態} という逆引きが完成したことを意味する。

\end{tcolorbox}

\section{25.5
線スペクトルへの橋}\label{sec-em-field-quantization-bridge}

{[}本文目安：M1{]}

連続スペクトル側の地図は本章で「電磁場の熱状態」として完成した。線スペクトル側の地図はどうか？

第VII部の中心地図
\(\alpha_\nu = (\pi e^2/m_e c) f_{lu} n_l \phi(\nu-\nu_0)\)
は、本章まででは原子物理（\(\nu_0, f_{lu}\)）、占有数（\(n_l\)）、線輪郭（\(\phi\)）の
\textbf{半古典的記述} ― 物質を量子化したが電磁場は古典場として扱う ―
にとどまっていた。

第26章では、本章で導入した \textbf{電磁場の量子化} と物質側の量子化を
\textbf{同時に} 扱うことで：

\begin{itemize}
\tightlist
\item
  自発放出率 \(A_{ul}\) が真空揺らぎから自然に出る（\textbf{Fermi
  の黄金律}）
\item
  Einstein 関係式 \(A_{ul} = (2h\nu^3/c^2) B_{ul}\) が QED
  の構造から直接従う ― \textbf{二つの地図の共通前因子 \(2h\nu^3/c^2\)
  の起源}
\item
  線の自然幅・選択則・禁制線が一貫した枠組みで現れる
\item
  連続スペクトル（プランク関数）と線スペクトル（水素線吸収係数）が、\textbf{同じ電磁場--物質相互作用
  Hamiltonian} の二側面として統一される
\end{itemize}

\begin{tcolorbox}[enhanced jigsaw, arc=.35mm, bottomrule=.15mm, bottomtitle=1mm, breakable, colback=white, colbacktitle=quarto-callout-note-color!10!white, colframe=quarto-callout-note-color-frame, coltitle=black, left=2mm, leftrule=.75mm, opacityback=0, opacitybacktitle=0.6, rightrule=.15mm, title=\textcolor{quarto-callout-note-color}{\faInfo}\hspace{0.5em}{ノート}, titlerule=0mm, toprule=.15mm, toptitle=1mm]

\textbf{使えるようになった道具}（§25.1〜§25.5）：

\begin{itemize}
\tightlist
\item
  「各モード = 量子調和振動子」の直観的描像
\item
  生成消滅演算子 \(\hat{a}, \hat{a}^\dagger\) と光子数演算子 \(\hat{n}\)
\item
  真空状態 \(|0\rangle\)
  と零点エネルギー、観測可能な帰結（自発放出、Casimir、自然幅）
\item
  黒体放射の最終解釈：「電磁場の熱状態」
\item
  連続スペクトル中心地図の最終整理（モード密度 + 偏光 + エネルギー量子 +
  熱光子数）
\end{itemize}

\end{tcolorbox}

\begin{center}\rule{0.5\linewidth}{0.5pt}\end{center}

\section*{この章で何がわかったか}\label{sec-em-field-quantization-summary}
\addcontentsline{toc}{section}{この章で何がわかったか}

\markright{この章で何がわかったか}

\textbf{中心地図に戻る}

本章で、プランク関数地図の最終解釈が定まった。各因子は別々の物理から来ているのではなく、\textbf{「電磁場を量子化したときの熱状態の表現」というひとつの構造の中で同時に決まる}：

\begin{itemize}
\tightlist
\item
  モード密度（第14章）
\item
  偏光自由度（第13章）
\item
  エネルギー量子（第10章）
\item
  Bose 統計（第8章）
\end{itemize}

これらすべてが「\textbf{各モード ＝
量子調和振動子}」という一行から、自動的に揃う。

次章では、線スペクトル地図にも同じ視点（電磁場の量子化 ＋
物質との相互作用）を適用し、Einstein
係数・自然幅・選択則を見通す。二つの中心地図が \textbf{同じ Hamiltonian
の二側面} であることが明らかになる。

\section*{演習問題}\label{sec-em-field-quantization-exercises}
\addcontentsline{toc}{section}{演習問題}

\markright{演習問題}

\textbf{模範解答}
は折りたたみを展開して確認できる（オンライン版）。まず自力で解いてから開くこと。

\subsection*{問題
25-1　各モード＝量子調和振動子という置き換え}\label{ex-25-1}
\addcontentsline{toc}{subsection}{問題
25-1　各モード＝量子調和振動子という置き換え}

{[}★ 難易度：☆☆ {]}

§25.2
は「各モードを量子調和振動子に置き換える」一行で第I〜VII部が接続されると述べた。モード
Hamiltonian
\(\hat H_\mathrm{mode}=\hbar\omega_{\mathbf k}(\hat a^\dagger_{\mathbf k\sigma}\hat a_{\mathbf k\sigma}+\tfrac12)\)。

\begin{enumerate}
\def\labelenumi{\arabic{enumi}.}
\tightlist
\item
  \(\hat a^\dagger\)、\(\hat a\)、\(\hat n=\hat a^\dagger\hat a\)
  がそれぞれ電磁場の文脈で何を表すか（放出・吸収・光子数）述べよ。
\item
  調和振動子の励起準位 \(|n\rangle\) が「そのモードに \(n\)
  個の光子がある状態」と読み替えられることを、第10章のエネルギー量子
  \(h\nu\)、第8章の Bose 統計、第14章のモード密度と対応づけよ。
\item
  「第I〜VII部で別々の道具と思っていたものが一行で接続される」とはどういう意味か、自分の言葉でまとめよ。
\end{enumerate}

\textbf{関連}：\hyperref[sec-em-field-quantization-mode-oscillator]{§25.2
各モードは調和振動子}／生成消滅演算子は\hyperref[ex-12-1]{演習
12-1}・\hyperref[ex-12-4]{演習
12-4}、モード密度は\hyperref[ex-14-1]{演習 14-1}、Bose
統計は\hyperref[ex-8-2]{演習 8-2}。

\section{模範解答（問題
25-1）}\label{ux6a21ux7bc4ux89e3ux7b54ux554fux984c-25-1}

\textbf{(1)} \(\hat a^\dagger_{\mathbf k\sigma}\)：モード
\((\mathbf k,\sigma)\)
に光子を1個生成＝放出。\(\hat a_{\mathbf k\sigma}\)：光子を1個消滅＝吸収。\(\hat n=\hat a^\dagger\hat a\)：そのモードの光子数（固有値
\(n=0,1,2,\dots\)）＝観測量。

\textbf{(2)} 励起準位 \(|n\rangle\) のエネルギーは
\(E_n=(n+\tfrac12)\hbar\omega\) で、隣接準位差
\(\hbar\omega=h\nu\)＝第10章のエネルギー量子。各モードを熱平衡の調和振動子とみた平均励起数
\(\langle n\rangle=1/(e^{h\nu/kT}-1)\)＝第8章の Bose 占有数。振動数
\(\nu\) 近傍のモードの個数＝第14章のモード密度 \(g(\nu)\)。

\textbf{(3)} エネルギー量子化（第10章）、光子の Bose
統計（第8章）、モード密度（第14章）、自発放出（第11章）―
これらは別々の物理に見えたが、すべて「各モード＝量子化された調和振動子」という一つの構造の異なる側面である。一行の置き換えで、観測される連続スペクトルの全因子が単一の枠組みに収まる。

\textbf{答え}：\(\hat a^\dagger\)＝放出、\(\hat a\)＝吸収、\(\hat n\)＝光子数。\(|n\rangle\)＝光子
\(n\) 個。量子化・Bose
統計・モード密度・自発放出がすべて「モード＝量子振動子」の一構造の側面。

\subsection*{問題
25-2　黒体放射を電磁場の熱状態として書き直す}\label{ex-25-2}
\addcontentsline{toc}{subsection}{問題
25-2　黒体放射を電磁場の熱状態として書き直す}

{[}★ 難易度：☆☆ {]}

§25.4
はプランク分布を「電磁場の熱状態の表現」として再構成した。熱密度行列
\(\hat\rho_\mathrm{th}=\dfrac1Z e^{-\hat H/k_BT}\) から組み立てる。

\begin{enumerate}
\def\labelenumi{\arabic{enumi}.}
\tightlist
\item
  各モードの平均光子数
  \(\bar n(\nu,T)=\mathrm{tr}(\hat\rho_\mathrm{th}\hat n)=\dfrac{1}{e^{h\nu/k_BT}-1}\)
  が、第III部で導いた Bose
  占有数と一致することを確認せよ（化学ポテンシャル \(\mu=0\)
  の理由も）。
\item
  モード密度 \(g(\nu)=8\pi\nu^2/c^3\) と光子エネルギー \(h\nu\) を掛けて
  \(u_\nu=\dfrac{8\pi h\nu^3}{c^3}\dfrac{1}{e^{h\nu/kT}-1}\)、\(B_\nu=\dfrac{c}{4\pi}u_\nu=\dfrac{2h\nu^3}{c^2}\dfrac{1}{e^{h\nu/kT}-1}\)
  を導け。
\item
  第8章（Bose
  統計）・第10章（振動子の熱平均）・第14章（モード密度×熱励起）の三つの導出が「同じ一つの構造」であるとはどういう意味か述べよ。
\end{enumerate}

\textbf{関連}：\hyperref[sec-em-field-quantization-blackbody]{§25.4
電磁場の熱状態}／Bose 占有数は\hyperref[ex-8-2]{演習
8-2}・\hyperref[ex-8-3]{演習 8-3}、全因子の合流は\hyperref[ex-14-4]{演習
14-4}。

\section{模範解答（問題
25-2）}\label{ux6a21ux7bc4ux89e3ux7b54ux554fux984c-25-2}

\textbf{(1)} 単一モードの熱平均
\(\bar n=\sum_n n\,e^{-n\hbar\omega/kT}/\sum_n e^{-n\hbar\omega/kT}=1/(e^{\hbar\omega/kT}-1)=1/(e^{h\nu/kT}-1)\)。これは第III部の
Bose
占有数（\(\mu=0\)）と同一。光子は数が保存しない（吸収・放出で生成消滅）ので
\(\mu=0\)。

\textbf{(2)}
\(u_\nu=g(\nu)\cdot h\nu\cdot\bar n=\dfrac{8\pi\nu^2}{c^3}\cdot h\nu\cdot\dfrac{1}{e^{h\nu/kT}-1}=\dfrac{8\pi h\nu^3}{c^3}\dfrac{1}{e^{h\nu/kT}-1}\)。等方場で
\(B_\nu=\dfrac{c}{4\pi}u_\nu=\dfrac{2h\nu^3}{c^2}\dfrac{1}{e^{h\nu/kT}-1}\)。

\textbf{(3)} 「光子気体の Bose
統計」（第8章）、「量子化振動子の熱平均」（第10章）、「モード密度×熱励起」（第14章）は出発点も言葉も違うが、すべて「温度
\(T\)
の熱浴と平衡にある電磁場（各モード＝調和振動子）の熱状態」を別の角度から記述したもの。同じ密度行列
\(\hat\rho_\mathrm{th}\) から同じプランク分布が出る。

\textbf{答え}：\(\bar n=1/(e^{h\nu/kT}-1)\)（\(\mu=0\)）。\(B_\nu=\frac{2h\nu^3}{c^2}\frac{1}{e^{h\nu/kT}-1}\)。三つの導出は電磁場の熱状態の同一構造の異なる側面。

\subsection*{問題
25-3　真空の零点エネルギーと観測可能な帰結}\label{ex-25-3}
\addcontentsline{toc}{subsection}{問題
25-3　真空の零点エネルギーと観測可能な帰結}

{[}★ 難易度：☆☆ {]}

§25.3 は真空 \(|0\rangle\)
が「空ではない」ことの観測的帰結を挙げた。零点エネルギー
\(E_0=\sum_{\mathbf k,\sigma}\tfrac12\hbar\omega_{\mathbf k}\)。

\begin{enumerate}
\def\labelenumi{\arabic{enumi}.}
\tightlist
\item
  零点エネルギーは形式的に発散するのに、なぜ観測される物理量（の差）は有限なのか述べよ。プランク分布に零点項が現れない理由（演習
  12-3）と結びつけよ。
\item
  真空が「空ではない」ことの観測可能な帰結を3つ挙げよ（自発放出・自然幅・Casimir
  効果）。Casimir 力が \(F\propto1/L^4\) になることに触れよ。
\item
  「真空に揺らぎがある」と「真空は粒子で満ちている」は同じか違うか。交換関係
  \([\hat a,\hat a^\dagger]=1\) と場の分散
  \(\langle(\Delta E)^2\rangle\neq0\) に触れて述べよ。
\end{enumerate}

\textbf{関連}：\hyperref[sec-em-field-quantization-vacuum]{§25.3
真空状態と零点エネルギー}／零点が観測量に効かないことは\hyperref[ex-7-2]{演習
7-2}・\hyperref[ex-12-3]{演習
12-3}、自然幅は\hyperref[sec-qed-minimum-natural-width]{§26.4}。

\section{模範解答（問題
25-3）}\label{ux6a21ux7bc4ux89e3ux7b54ux554fux984c-25-3}

\textbf{(1)}
零点エネルギーはエネルギーの絶対基準（定数オフセット）で、その絶対値は観測できない。観測できるのは零点エネルギーの\textbf{差や変化}（境界条件の変化など）で、これは有限。プランク分布に零点項
\(\tfrac12\hbar\omega\)
が現れないのも、それが温度非依存の定数で放射の励起分（物質とやりとりする部分）に効かないから（演習
12-3）。

\textbf{(2)} ① 自発放出（真空揺らぎが励起原子を揺さぶる、§26.2）。②
線の自然幅（励起準位の有限寿命、\(\Gamma=A_{ul}\)）。③ Casimir
効果（二枚の金属板間で許されるモードが減り、板間の零点エネルギーが板外より小さくなり引力
\(F\propto1/L^4\) が生じる）。

\textbf{(3)} 違う。真空 \(|0\rangle\)
は厳密に光子0個の状態（\(\hat n|0\rangle=0\)）。観測される「揺らぎ」は粒子がいることではなく、\([\hat a,\hat a^\dagger]=1\)
の非可換性から来る電場・磁場の期待値の分散
\(\langle(\Delta E)^2\rangle\neq0\)。粒子数を定めると場の値が揺らぎ、場の値を定めると粒子数が揺らぐ（不確定性関係の双対）。

\textbf{答え}：零点の絶対値は観測不能、差は有限。帰結＝自発放出・自然幅・Casimir(\(F\propto1/L^4\))。揺らぎ＝場の分散(\([\hat a,\hat a^\dagger]=1\))で、粒子で満ちているのとは別。

\subsection*{問題 25-4　なぜ電磁場の量子化が必要だったか}\label{ex-25-4}
\addcontentsline{toc}{subsection}{問題
25-4　なぜ電磁場の量子化が必要だったか}

{[}★ 難易度：☆ {]}

§25.1
は、古典電磁気＋（別々の）量子論では説明しきれず、電磁場そのものの量子化が要る現象を列挙した。

\begin{enumerate}
\def\labelenumi{\arabic{enumi}.}
\tightlist
\item
  次の各現象が古典電磁気では説明できない理由を一言ずつ述べよ：エネルギー量子化
  \(h\nu\)、自発放出 \(A_{ul}\)、線の自然幅、光子の Bose 統計。
\item
  これらが「同じ枠組み（電磁場の量子化）」から出ることが、なぜ本書の逆引きの到達点なのか述べよ。
\item
  本書では電磁場と量子論を第V部と第IV部で別々に扱ってきた。第VIII部でそれらを「同時に量子化する」ことの意味を述べよ。
\end{enumerate}

\textbf{関連}：\hyperref[sec-em-field-quantization-why]{§25.1
なぜ量子化が必要か}／自発放出は\hyperref[ex-11-1]{演習 11-1}、Bose
統計は\hyperref[ex-8-1]{演習 8-1}。

\section{模範解答（問題
25-4）}\label{ux6a21ux7bc4ux89e3ux7b54ux554fux984c-25-4}

\textbf{(1)}
エネルギー量子化：古典波には離散性がない。自発放出：古典では「ゼロ放射場→放出ゼロ」で減衰の起源がない。自然幅：励起準位の有限寿命を古典電子モデルでは「決め打ち」するしかない。Bose
統計：古典場は数えられる粒子ではなく、識別不能粒子の統計が出てこない。

\textbf{(2)}
第I〜VII部で別々に逆引きしてきた断片（量子化・統計・自発放出・自然幅）が、「電磁場を量子化する」という一つの枠組みからすべて出る。観測スペクトルから出発した逆引きが、最終的に単一の物理基盤（QED）に収束する
― これが到達点。

\textbf{(3)}
第IV部は物質（原子）を量子化したが電磁場は古典場、第V部は電磁場を古典的に扱った。第VIII部で電磁場と物質を\textbf{同時に}量子化することで、両者の相互作用（自発放出・Einstein
関係・選択則）が一貫した枠組みで現れ、連続スペクトルと線スペクトルが統一される。

\textbf{答え}：量子化・自発放出・自然幅・Bose
統計は古典場では出ない。電磁場量子化で全断片が一基盤に収束。第VIII部は電磁場と物質を同時量子化して両スペクトルを統一。

\section*{さらに学ぶための参考文献}\label{sec-em-field-quantization-further}
\addcontentsline{toc}{section}{さらに学ぶための参考文献}

\markright{さらに学ぶための参考文献}

\begin{itemize}
\tightlist
\item
  Loudon, \emph{The Quantum Theory of Light} (Oxford, 3rd ed., 2000) ---
  直観重視の電磁場量子化、本章のレベルの標準書
\item
  Mandel \& Wolf, \emph{Optical Coherence and Quantum Optics}
  (Cambridge, 1995) --- より体系的に進めたい読者向け
\item
  Milonni, \emph{The Quantum Vacuum} (Academic Press, 1994) ---
  真空揺らぎ・Casimir 効果の専門書
\item
  Peskin \& Schroeder, \emph{An Introduction to Quantum Field Theory}
  (Westview, 1995) --- 場の量子論の本格教科書、最初の数章
\item
  Weinberg, ``The Cosmological Constant Problem,'' Rev.~Mod. Phys. 61
  (1989) 1 --- 宇宙定数問題の古典レビュー
\end{itemize}

\chapter{第26章　Einstein係数・自然幅・選択則を QED
で見通す}\label{ux7b2c26ux7ae0-einsteinux4fc2ux6570ux81eaux7136ux5e45ux9078ux629eux5247ux3092-qed-ux3067ux898bux901aux3059}

\textbf{この章の主題}：量子電磁力学（QED）の最小限の枠組みで、本書を貫いてきた四つのテーマ
―（i）Einstein 係数の関係、（ii）自然幅、（iii）真空揺らぎ、（iv）選択則
―
を統一的に見通す。最終章として、連続スペクトル（プランク関数地図）と線スペクトル（水素線地図）が、同じ電磁場--物質相互作用の二側面であることを総括する。

\begin{tcolorbox}[enhanced jigsaw, arc=.35mm, bottomrule=.15mm, bottomtitle=1mm, breakable, colback=white, colbacktitle=quarto-callout-important-color!10!white, colframe=quarto-callout-important-color-frame, coltitle=black, left=2mm, leftrule=.75mm, opacityback=0, opacitybacktitle=0.6, rightrule=.15mm, title=\textcolor{quarto-callout-important-color}{\faExclamation}\hspace{0.5em}{重要}, titlerule=0mm, toprule=.15mm, toptitle=1mm]

\textbf{本章の方針 ― 厳密な QED ではなく「見通し」}

本章では Feynman
図やくり込みといった本格的な場の理論の道具は扱わない。代わりに、第VII部までで顔を出した次の四つの主題を、QED
の言葉でひとつにまとめ直す：

\begin{enumerate}
\def\labelenumi{\arabic{enumi}.}
\tightlist
\item
  \textbf{Einstein 係数}：\(A_{ul} = (2h\nu^3/c^2)(g_l/g_u) B_{lu}\)
  がなぜ成立するか
\item
  \textbf{自然幅}：線の有限の幅が、自発放出率 \(A_{ul}\)
  の逆数として現れる理由
\item
  \textbf{真空揺らぎ}：自発放出を引き起こす「真空のゆらぎ」の意味
\item
  \textbf{選択則}：許容遷移と禁制遷移が、電気双極子行列要素から決まる仕組み
\end{enumerate}

これらを通じて、読者は厳密な QED
に進む前段階として、観測天体物理学を理解するのに必要な「見通し」を持つ。

\end{tcolorbox}

\section*{この章の中心地図}\label{sec-qed-minimum-map}
\addcontentsline{toc}{section}{この章の中心地図}

\markright{この章の中心地図}

\textbf{プランク関数（連続スペクトルの地図）}

\begin{equation}\protect\phantomsection\label{eq-planck-function}{
B_\nu(T) = \frac{2 h \nu^3}{c^2} \cdot \frac{1}{e^{h\nu/kT} - 1}
}\end{equation}

{\def\LTcaptype{none} % do not increment counter
\begin{longtable}[]{@{}lll@{}}
\toprule\noalign{}
因子 & 物理的意味 & 結びつく部 \\
\midrule\noalign{}
\endhead
\bottomrule\noalign{}
\endlastfoot
\(B_\nu\) & 観測量（比強度） & 第II部 \\
\(1/(e^{h\nu/kT}-1)\) & 熱平衡・光子統計 & 第III部 \\
\(h\nu\) & エネルギー量子 & 第IV部 \\
\(\nu^3/c^2\) & 電磁場のモード密度 & 第V部 \\
黒体からのずれ & 現実の天体 & 第VI部 \\
\end{longtable}
}

\begin{tcolorbox}[enhanced jigsaw, arc=.35mm, bottomrule=.15mm, breakable, colback=white, colframe=quarto-callout-tip-color-frame, left=2mm, leftrule=.75mm, opacityback=0, rightrule=.15mm, toprule=.15mm]

\textbf{地図の使い方}：各部の冒頭でこの式が再掲され、その部で扱う因子が強調表示される。読者は自分が今どこを学んでいるかをこの式の上で確認できる。

\end{tcolorbox}

\textbf{水素線スペクトルの地図 ― 二層構造}

\textbf{第一層：イコニック・アンカー（リュードベリ公式）}

\begin{equation}\protect\phantomsection\label{eq-rydberg}{
\frac{1}{\lambda} = R_\mathrm{H} \left( \frac{1}{n_1^2} - \frac{1}{n_2^2} \right)
}\end{equation}

線スペクトルの「骨格」を一目で示す。\(R_\mathrm{H}\)
は基本定数の組合せ、\(1/n^2\)
は束縛状態のエネルギー、離散性が「線が存在する理由」を語る。

\textbf{第二層：構造的マップ（Hα を標準例とした線吸収係数）}

\begin{equation}\protect\phantomsection\label{eq-line-absorption}{
\alpha_\nu = \frac{\pi e^2}{m_e c} \, f_{lu} \, n_l \, \phi(\nu - \nu_0)
}\end{equation}

{\def\LTcaptype{none} % do not increment counter
\begin{longtable}[]{@{}lll@{}}
\toprule\noalign{}
因子 & 物理的意味 & 対応する章 \\
\midrule\noalign{}
\endhead
\bottomrule\noalign{}
\endlastfoot
\(\alpha_\nu\) & 観測量（線吸収係数） & 第21章 \\
\(\pi e^2 / m_e c\) & 古典極限の線強度 & 第22章＋第VIII部 \\
\(f_{lu}\) & 振動子強度 & 第22章 \\
\(n_l\) & 下準位の占有数 & 第21章＋Part III \\
\(\phi(\nu - \nu_0)\) & 線輪郭（自然・Doppler・Voigt） & 第23章 \\
\(\nu_0\) & 線位置（リュードベリ） & 第22章 \\
\end{longtable}
}

\textbf{適用範囲（重要）}：この式の前因子 \(\pi e^2/m_e c\) と振動子強度
\(f_{lu}\) は、古典 Lorentz 振動子に対応する
\textbf{電気双極子（E1）許容線} に対するものである。21 cm
線（基底状態超微細構造の \textbf{磁気双極子 M1 遷移}）や
\textbf{禁制線（E2・M1 など）} では、この前因子・\(f_{lu}\)
をそのまま使えない。これらの線では、測定された自発放出率 \(A_{ul}\) から
Einstein
関係を介して有効遷移強度（\(B_{lu}\)）に読み替える（式~\ref{eq-line-absorption}
の構造は保ったまま、強度因子だけを差し替える）。

\begin{tcolorbox}[enhanced jigsaw, arc=.35mm, bottomrule=.15mm, breakable, colback=white, colframe=quarto-callout-tip-color-frame, left=2mm, leftrule=.75mm, opacityback=0, rightrule=.15mm, toprule=.15mm]

\textbf{地図の使い方}：第VII部の各章はこの式の異なる因子を解き明かす。プランク関数地図と並べて読めば、本書を貫く二つの逆引きが同時に把握できる。第VIII部の
QED で、両者は前因子 \(2h\nu^3/c^2\)
を介して統一される（式~\ref{eq-einstein-A} 参照）。

\end{tcolorbox}

そして本章の主役：

\textbf{Einstein 係数の関係 ― 連続と線を結ぶ式}

\begin{equation}\protect\phantomsection\label{eq-einstein-A}{
A_{ul} = \frac{2 h \nu^3}{c^2} \cdot \frac{g_l}{g_u} \cdot B_{lu}
}\end{equation}

（添字 \(u\) は upper、\(l\) は
lower。本書の本文表記と統一。古典文献では \(A_{21}, B_{12}, g_1/g_2\)
と書かれることが多い。）

強調した前因子 \(2h\nu^3/c^2\)
は、プランク関数（式~\ref{eq-planck-function}）と\textbf{まったく同じ形}で現れる。これは連続放射と線スペクトルが、同じ電磁場--物質相互作用の二側面であることを示す数学的証拠である。本書全体が、この式によって
QED で締めくくられる。

\section*{この章で答える問い}\label{sec-qed-minimum-questions}
\addcontentsline{toc}{section}{この章で答える問い}

\markright{この章で答える問い}

\begin{itemize}
\tightlist
\item
  なぜ Einstein A 係数の前因子 \(2h\nu^3/c^2\)
  が、プランク関数の前因子と一致するのか
\item
  自発放出は何が原因で起きるのか（真空揺らぎとの関係）
\item
  線の自然幅は、観測量としては何を測っているのか
\item
  選択則はなぜ「許される」と「禁制」を区別するのか
\item
  連続スペクトルと線スペクトルは、QED の中でどう一つに収束するのか
\end{itemize}

\section*{到達目標}\label{sec-qed-minimum-goals}
\addcontentsline{toc}{section}{到達目標}

\markright{到達目標}

この章を読み終えると、読者は次のことができるようになる：

\begin{itemize}
\tightlist
\item
  Einstein 係数の関係式 \(A_{ul} = (2h\nu^3/c^2)(g_l/g_u) B_{lu}\)
  を、QED 的な「見通し」で再導出できる
\item
  自然幅 ↔ 自発放出率 ↔ 真空揺らぎの三角関係を整理できる
\item
  連続スペクトル（プランク関数地図）と線スペクトル（水素線地図）が共通の電磁場--物質相互作用の二側面であることを、本書全体の総括として説明できる
\end{itemize}

\begin{center}\rule{0.5\linewidth}{0.5pt}\end{center}

\section{26.1 相互作用 Hamiltonian ―
物質と電磁場をつなぐ}\label{sec-qed-minimum-interaction}

{[}本文目安：M1{]}

\begin{tcolorbox}[enhanced jigsaw, arc=.35mm, bottomrule=.15mm, bottomtitle=1mm, breakable, colback=white, colbacktitle=quarto-callout-warning-color!10!white, colframe=quarto-callout-warning-color-frame, coltitle=black, left=2mm, leftrule=.75mm, opacityback=0, opacitybacktitle=0.6, rightrule=.15mm, title=\textcolor{quarto-callout-warning-color}{\faExclamationTriangle}\hspace{0.5em}{警告}, titlerule=0mm, toprule=.15mm, toptitle=1mm]

\textbf{本章の単位系について}：本章は \textbf{QED
の構造を見通すことが主目的} であるが、本章の主式
式~\ref{eq-interaction-hamiltonian}、式~\ref{eq-vector-potential-quantized}、式~\ref{eq-A-coefficient-QED}
は \textbf{本書の方針どおり SI で厳密に記す}（前因子
\(4\pi\epsilon_0\)、\(1/c\) などを略さずに含める）。一方、QED
の標準文献は Gauss/Heaviside-Lorentz
単位系を用いるものも多く、同じ式でも前因子の見かけが異なる。\textbf{SI
と Gauss/HL の式形の対応は付録 B §B.4 に集約した}
ので、論文と比較する際はそちらと Loudon 第 4
章を参照し、必ず単位系を確認してほしい。各係数の詳細な導出は専門書（Loudon、Sakurai、Cohen-Tannoudji）に譲る。

\end{tcolorbox}

QED の言葉で、原子と電磁場の相互作用は次の Hamiltonian
で記述される（最低次、非相対論的、SI 系の代表形）：

\begin{equation}\protect\phantomsection\label{eq-interaction-hamiltonian}{
\hat{H}_\text{int} = -\frac{e}{m_e} \hat{\mathbf{A}} \cdot \hat{\mathbf{p}} + (\text{higher orders})
}\end{equation}

ここで：

\begin{itemize}
\tightlist
\item
  \(\hat{\mathbf{p}}\) は
  \textbf{電子の運動量演算子}（第22章で扱った原子の波動関数に作用）
\item
  \(\hat{\mathbf{A}}\) は
  \textbf{電磁場のベクトルポテンシャル演算子}（第25章で量子化されたもの）
\end{itemize}

\(\hat{\mathbf{A}}\) は生成消滅演算子で書き下せる：

\begin{equation}\protect\phantomsection\label{eq-vector-potential-quantized}{
\hat{\mathbf{A}}(\mathbf{r}, t) = \sum_{\mathbf{k}, \sigma} \sqrt{\frac{\hbar}{2\epsilon_0 \omega_\mathbf{k} V}} \left[\hat{a}_{\mathbf{k}\sigma} \boldsymbol{\epsilon}_{\mathbf{k}\sigma} e^{i\mathbf{k}\cdot\mathbf{r} - i\omega t} + \text{h.c.}\right]
}\end{equation}

\(\boldsymbol{\epsilon}_{\mathbf{k}\sigma}\) は偏光ベクトル、\(V\)
は規格化体積、「h.c.」は Hermitian 共役（生成項）。

\begin{tcolorbox}[enhanced jigsaw, arc=.35mm, bottomrule=.15mm, bottomtitle=1mm, breakable, colback=white, colbacktitle=quarto-callout-note-color!10!white, colframe=quarto-callout-note-color-frame, coltitle=black, left=2mm, leftrule=.75mm, opacityback=0, opacitybacktitle=0.6, rightrule=.15mm, title=\textcolor{quarto-callout-note-color}{\faInfo}\hspace{0.5em}{ノート}, titlerule=0mm, toprule=.15mm, toptitle=1mm]

\textbf{方針}：本章では \(\hat{H}_\text{int}\) から自発放出率
\(A_{ul}\)、Einstein 関係、線輪郭を \textbf{同時に}
導く流れを示す。各項の厳密な計算は専門書（Loudon, Sakurai,
Cohen-Tannoudji）に譲り、本書では \textbf{どこから何が出るか}
の見通しを与える。

\end{tcolorbox}

\begin{tcolorbox}[enhanced jigsaw, arc=.35mm, bottomrule=.15mm, bottomtitle=1mm, breakable, colback=white, colbacktitle=quarto-callout-note-color!10!white, colframe=quarto-callout-note-color-frame, coltitle=black, left=2mm, leftrule=.75mm, opacityback=0, opacitybacktitle=0.6, rightrule=.15mm, title=\textcolor{quarto-callout-note-color}{\faInfo}\hspace{0.5em}{ノート}, titlerule=0mm, toprule=.15mm, toptitle=1mm]

\textbf{双極子近似とゲージ}：相互作用は上のように速度形（\(\hat{\mathbf{A}}\cdot\hat{\mathbf{p}}\)、Coulomb
ゲージ）でも、長さ形（\(-\hat{\mathbf{d}}\cdot\hat{\mathbf{E}}\)、\(\hat{\mathbf{d}}=-e\hat{\mathbf{r}}\)）でも書ける。両者はゲージ変換（Power-Zienau-Woolley
変換）で結ばれた \textbf{同値な記述}
であり、エネルギー保存則を満たす遷移（実遷移）に対しては同じ遷移率を与える。本章で
\(A_{ul}\) を電気双極子モーメント \(\hat{\mathbf{d}}\)
で書く（式~\ref{eq-A-coefficient-QED}）のは、原子の大きさ \(a_0\)
が放出光の波長 \(\lambda\)
より十分小さく（\(a_0 \ll \lambda\)）、\(e^{i\mathbf{k}\cdot\mathbf{r}}\simeq 1\)
と置ける \textbf{長波長（電気双極子）近似}
が成り立つためである。この近似が破れる高次の項が、後の磁気双極子（M1）・電気四重極（E2）遷移（§26.5）に対応する。

\end{tcolorbox}

\section{\texorpdfstring{26.2 自発放出と Fermi の黄金律 ― \(A_{ul}\)
の導出}{26.2 自発放出と Fermi の黄金律 ― A\_\{ul\} の導出}}\label{sec-qed-minimum-spontaneous}

{[}本文目安：M1｜発展：M2+{]}

自発放出は次のように起こる：

\begin{itemize}
\tightlist
\item
  初期状態：原子は上準位 \(|u\rangle\)、電磁場は真空 \(|0\rangle\) ―
  全体状態 \(|u; 0\rangle\)
\item
  終状態：原子は下準位 \(|l\rangle\)、電磁場に光子 1 個 ― 全体状態
  \(|l; 1_{\mathbf{k}\sigma}\rangle\)
\end{itemize}

遷移率は \textbf{Fermi の黄金律}：

\begin{equation}\protect\phantomsection\label{eq-fermi-golden-rule}{
\Gamma_{u \to l} = \frac{2\pi}{\hbar} |\langle l; 1_{\mathbf{k}\sigma} | \hat{H}_\text{int} | u; 0\rangle|^2 \cdot \rho(\omega)
}\end{equation}

ここで \(\rho(\omega)\)
は終状態のモード密度。第14章で逆引きしたモード密度を使うと、

\begin{equation}\protect\phantomsection\label{eq-A-coefficient-QED}{
A_{ul} = \frac{\omega^3}{3\pi\epsilon_0 \hbar c^3} |\langle l | \hat{\mathbf{d}} | u \rangle|^2
}\end{equation}

\(\hat{\mathbf{d}} = -e\hat{\mathbf{r}}\) は \textbf{電気双極子演算子}。

ここで本書全体の流れが閉じる：

\begin{itemize}
\tightlist
\item
  真空 \(|0\rangle\) から光子を生成する \(\hat{a}^\dagger\) の
  \textbf{真空揺らぎ} が、自発放出の \textbf{物理的駆動力} である
\item
  古典電磁気では「ゼロ放射場 → 放出ゼロ」だったが、QED
  では真空場との結合により自発放出率が非ゼロになり、\(A_{ul} \neq 0\)
  を生む
\item
  式~\ref{eq-A-coefficient-QED} の前因子 \(\omega^3\)
  は、\textbf{モード密度 \(\propto \nu^2\) × エネルギー量子 \(h\nu\)}
  の積から来ており、これは第25章でプランク関数の前因子に出てきた
  \textbf{\(2h\nu^3/c^2\) と同じ構造}
\end{itemize}

\begin{tcolorbox}[enhanced jigsaw, arc=.35mm, bottomrule=.15mm, breakable, colback=white, colframe=quarto-callout-tip-color-frame, left=2mm, leftrule=.75mm, opacityback=0, rightrule=.15mm, toprule=.15mm]
\begin{minipage}[t]{5.5mm}
\textcolor{quarto-callout-tip-color}{\faLightbulb}
\end{minipage}%
\begin{minipage}[t]{\textwidth - 5.5mm}

\textbf{問い}：「自発放出」は、なぜ「自発」と呼ばれるのか？
本当に「自発」なのか？

\textbf{短答}：古典的な観点では「外部刺激なしに勝手に起こる」ように見えるので「自発」と呼ぶが、QED
では \textbf{真空揺らぎという形の電磁場の外部刺激}
が常にあり、それが誘起している。「真空揺らぎで誘起された自発放出」とでも言うべき現象。

\textbf{もう一歩}：教科書的に分類されてきた「自発放出
\(A_{ul}\)」と「誘導放出 \(B_{ul}\)」は、QED の立場では \textbf{同じ
Fermi の黄金律の二側面}：

\begin{itemize}
\tightlist
\item
  \(B_{ul}\)：終状態のモードが既に光子を持つ場合 → \(|n + 1\rangle\)
  への遷移、Bose 増強因子 \(\sqrt{n+1}\)
\item
  \(A_{ul}\)：終状態のモードが真空の場合 →
  \(|0\rangle \to |1\rangle\)、増強因子 1
\end{itemize}

両者の比
\(A_{ul}/B_{ul} = 2h\nu^3/c^2\)（モード密度＋エネルギー量子）が、第11章で詳細釣り合いから導いた
\textbf{Einstein 関係式} と一致する。

\end{minipage}%
\end{tcolorbox}

\section{26.3 Einstein 関係式の QED 起源 ―
二つの地図の交点}\label{sec-qed-minimum-einstein}

{[}本文目安：M1{]}

前節で出た \(A_{ul} = (2h\nu^3/c^2) B_{ul}\) の前因子 \(2h\nu^3/c^2\)
を、二つの中心地図の \textbf{共通の起源} として整理する：

{\def\LTcaptype{none} % do not increment counter
\begin{longtable}[]{@{}lll@{}}
\toprule\noalign{}
因子 & プランク関数（第14章） & 自発放出率（本章） \\
\midrule\noalign{}
\endhead
\bottomrule\noalign{}
\endlastfoot
\textbf{\(\nu^3\)} & モード密度 \(\nu^2\) × エネルギー量子 \(h\nu\) &
同じ \\
\textbf{\(1/c^2\)} & 分散関係 \(k = 2\pi\nu/c\) ＋ \(u_\nu \to B_\nu\) &
同じ \\
\textbf{2} & 偏光自由度（横波の二つの自由度） & 同じ \\
\textbf{\(h\)} & エネルギー量子 & 同じ \\
\end{longtable}
}

つまり \(2h\nu^3/c^2\) は
\textbf{「電磁場の同じモード構造を、それぞれの観点から見たもの」}
に他ならない：

\begin{itemize}
\tightlist
\item
  プランク関数で見れば：単位体積・単位振動数あたりのモード密度 ×
  エネルギー量子
\item
  自発放出で見れば：原子が放出するときの「許される終状態モード」の密度 ×
  エネルギー量子
\end{itemize}

これが、本書を
\textbf{逆引きで貫いてきた到達点}：プランク関数と水素線吸収係数は別々の物理ではなく、\textbf{QED
という共通の土台の上で連なっている}。

\begin{tcolorbox}[enhanced jigsaw, arc=.35mm, bottomrule=.15mm, bottomtitle=1mm, breakable, colback=white, colbacktitle=quarto-callout-note-color!10!white, colframe=quarto-callout-note-color-frame, coltitle=black, left=2mm, leftrule=.75mm, opacityback=0, opacitybacktitle=0.6, rightrule=.15mm, title=\textcolor{quarto-callout-note-color}{\faInfo}\hspace{0.5em}{ノート}, titlerule=0mm, toprule=.15mm, toptitle=1mm]

\textbf{観測との対応}：観測されるさまざまなスペクトル ― CMB
の完璧な黒体（第16章）、恒星 H\(\alpha\) 線の Voigt
プロファイル（第23章）、QSO
吸収系（第24章）、降着円盤多温度合成スペクトル（第18章）―
は、いずれも電磁場の量子化されたモード構造を背景に持つ。連続放射では前因子
\(2h\nu^3/c^2\) がプランク関数に直接現れ、線スペクトルでは同じ因子が
Einstein 係数比 \(A_{ul}/B_{ul}\)
を介して輝線・吸収線の強度関係に効いてくる（線吸収係数の中心地図そのものにこの前因子が直接入るわけではない点に注意）。したがって、日常的な分光観測にも、QED
で記述される電磁場のモード構造が反映されている。

\end{tcolorbox}

\section{26.4 自然幅と真空揺らぎ ―
観測量としての真空のゆらぎ}\label{sec-qed-minimum-natural-width}

{[}本文目安：M1{]}

第23章で「線輪郭」の機構として導入した自然幅
\(\Gamma_\text{nat} = A_{ul}/(2\pi)\) を、本節で QED 的に意味づける。

QED の直観：

\begin{itemize}
\tightlist
\item
  古典電磁気では、励起原子は「自分で勝手に光を出す」理由がない
\item
  量子化された電磁場の真空には、各モードに \textbf{零点揺らぎ} がある
\item
  この
  \textbf{真空揺らぎとの相互作用が励起原子の遷移を絶えず誘起し}、それが自発放出として観測される
\item
  励起状態の寿命 \(\tau = 1/A_{ul}\) は有限となり、不確定性関係
  \(\Delta E \cdot \tau \gtrsim \hbar\) から線にはエネルギー幅
  \(\Delta E = \hbar A_{ul}\)、すなわち振動数幅 \(A_{ul}/(2\pi)\)
  が現れる
\end{itemize}

つまり
\textbf{線の自然幅は真空揺らぎの存在を観測している}。第23章で「線輪郭機構の一つ」として挙げた自然幅は、QED
の核心 ― 真空が「空ではない」こと ― と直結する。

\begin{tcolorbox}[enhanced jigsaw, arc=.35mm, bottomrule=.15mm, bottomtitle=1mm, breakable, colback=white, colbacktitle=quarto-callout-tip-color!10!white, colframe=quarto-callout-tip-color-frame, coltitle=black, left=2mm, leftrule=.75mm, opacityback=0, opacitybacktitle=0.6, rightrule=.15mm, title=\textcolor{quarto-callout-tip-color}{\faLightbulb}\hspace{0.5em}{ヒント}, titlerule=0mm, toprule=.15mm, toptitle=1mm]

\textbf{例（数値）}：水素 Ly\(\alpha\) 線（\(\lambda = 121.6\) nm,
\(\nu = 2.47 \times 10^{15}\) Hz）の自発放出率
\(A_{ul} = 6.27 \times 10^8\) s\(^{-1}\) から、自然幅は
\(\gamma_\text{nat} = A_{ul}/(2\pi) \simeq 10^8\) Hz、波長単位で
\(\Delta\lambda_\text{nat} = \lambda^2 \gamma_\text{nat}/c \simeq 4.9 \times 10^{-5}\)
Å（第23章 §23.1 と同値）。

これは
\textbf{「水素原子が真空揺らぎに揺さぶられる強さ」を観測量にしたもの}
である。観測される QSO の damped Lyα
系の減衰翼は、まさにこの自然幅起源の Lorentzian
翼を、地球上の分光計で測っている。

\end{tcolorbox}

\begin{tcolorbox}[enhanced jigsaw, arc=.35mm, bottomrule=.15mm, bottomtitle=1mm, breakable, colback=white, colbacktitle=quarto-callout-note-color!10!white, colframe=quarto-callout-note-color-frame, coltitle=black, left=2mm, leftrule=.75mm, opacityback=0, opacitybacktitle=0.6, rightrule=.15mm, title=\textcolor{quarto-callout-note-color}{\faInfo}\hspace{0.5em}{ノート}, titlerule=0mm, toprule=.15mm, toptitle=1mm]

\textbf{観測との対応}：21 cm 線（第20章 §20.4、第24章 §24.5）の
\(A_{ul} \simeq 2.9 \times 10^{-15}\) s\(^{-1}\)
は、人類が観測する最も弱い遷移のひとつである。それでも
\textbf{宇宙に存在する膨大な中性水素のおかげで}
観測可能となる。真空揺らぎの効果がいかに弱くても、宇宙的スケールでは無視できない、というのが線スペクトル天文学の根底にある事実である。

\end{tcolorbox}

\section{26.5 選択則を行列要素で見通す}\label{sec-qed-minimum-selection}

{[}本文目安：M1｜発展：M2+{]}

第22章 §22.5 で「許される遷移」と「禁制遷移」を扱った。本節ではこれを
QED の \textbf{電気双極子行列要素}

\begin{equation}\protect\phantomsection\label{eq-dipole-matrix-element}{
\mathbf{d}_{ul} \equiv \langle u | \hat{\mathbf{d}} | l \rangle = -e \langle u | \hat{\mathbf{r}} | l \rangle
}\end{equation}

から見通す。

QED の構造から：

\begin{itemize}
\tightlist
\item
  電場と原子の相互作用は \textbf{最低次で電気双極子}
  \(\hat{\mathbf{d}}\) を介する（式~\ref{eq-interaction-hamiltonian}
  の長波長極限）
\item
  \(\langle u | \hat{\mathbf{d}} | l \rangle = 0\) の遷移は最低次で禁止
  → \textbf{電気双極子禁制（E1 forbidden）遷移}
\item
  パリティと角運動量の保存から、\(\Delta l = \pm 1\)、\(\Delta m = 0, \pm 1\)
  の選択則が自然に導かれる
\item
  禁制遷移は完全に消えるわけではなく、\textbf{磁気双極子（M1）・電気四重極（E2）など高次過程}
  として残り、観測上は弱いがゼロではない
\end{itemize}

天文学的意味：

\begin{itemize}
\tightlist
\item
  禁制遷移は地上では衝突で潰れて見えないが、\textbf{星間ガス（低密度・低温）では観測できる}
  → ISM の物理状態診断に重要（{[}O III{]} \$\lambda\$5007
  等の輝線、第24章 §24.3）
\item
  \textbf{21 cm 線} は水素の超微細構造の磁気双極子遷移 →
  強い禁制だが宇宙には水素原子が大量にあるため観測できる（§26.4
  末尾の例も参照）
\end{itemize}

\begin{tcolorbox}[enhanced jigsaw, arc=.35mm, bottomrule=.15mm, bottomtitle=1mm, breakable, colback=white, colbacktitle=quarto-callout-note-color!10!white, colframe=quarto-callout-note-color-frame, coltitle=black, left=2mm, leftrule=.75mm, opacityback=0, opacitybacktitle=0.6, rightrule=.15mm, title=\textcolor{quarto-callout-note-color}{\faInfo}\hspace{0.5em}{ノート}, titlerule=0mm, toprule=.15mm, toptitle=1mm]

\textbf{第22章との対応}：第22章 §22.5
で「選択則は経験則として」述べた規則が、本章では
\textbf{電気双極子相互作用の最低次} という QED
構造から自然に導かれる。これは「天下り的に与えない」という本書の原則 ③
を、最終章で線スペクトル側に適用した結果である。

\end{tcolorbox}

\section{26.6 連続と線の統一 ―
本書全体の総括}\label{sec-qed-minimum-unification}

{[}本文目安：M1{]}

連続スペクトル（プランク関数）と線スペクトル（水素線吸収係数）は、それぞれ別の地図を持ち、別々の章群（第I〜VI部
vs 第VII部）で逆引きしてきた。しかし両者は QED で次のように統合される：

{\def\LTcaptype{none} % do not increment counter
\begin{longtable}[]{@{}
  >{\raggedright\arraybackslash}p{(\linewidth - 4\tabcolsep) * \real{0.3333}}
  >{\raggedright\arraybackslash}p{(\linewidth - 4\tabcolsep) * \real{0.3333}}
  >{\raggedright\arraybackslash}p{(\linewidth - 4\tabcolsep) * \real{0.3333}}@{}}
\toprule\noalign{}
\begin{minipage}[b]{\linewidth}\raggedright
\end{minipage} & \begin{minipage}[b]{\linewidth}\raggedright
連続（黒体放射）
\end{minipage} & \begin{minipage}[b]{\linewidth}\raggedright
線スペクトル
\end{minipage} \\
\midrule\noalign{}
\endhead
\bottomrule\noalign{}
\endlastfoot
観測量 & プランク分布 \(B_\nu(T)\) & 線吸収係数 \(\alpha_\nu\) \\
統計 & 各モードの熱平均（光子の Bose 統計） &
原子準位の占有数（Boltzmann/Saha） \\
量子化 & エネルギー量子 \(h\nu\) & 原子の束縛状態のエネルギー差 \\
モード密度 & \(\nu^2/c^3\)（電磁場側） & \(\nu^2/c^3\)（同じ！）→
\(A_{ul}\) の前因子 \\
偏光自由度 & 2 & 2（同じ） \\
QED 解釈 & 電磁場の熱状態 &
物質と電磁場の遷移（真空揺らぎが自発放出を駆動） \\
相互作用 & 平衡確立後は \(\hat{H}_\text{int}\) を切って
\textbf{自由場の熱状態}
として記述（実際の熱化過程では物質との相互作用が不可欠、第16章 §16.2） &
\(\hat{H}_\text{int} \neq 0\)（双極子遷移を明示的に扱う） \\
\end{longtable}
}

最終的に、\(2h\nu^3/c^2\)
という前因子が両者に共通して現れることが、この統合の
\textbf{数学的な印}：

\begin{equation}\protect\phantomsection\label{eq-prefactor-unification}{
\boxed{
\frac{A_{ul}}{B_{ul}} = \frac{2h\nu^3}{c^2} = \text{(モード密度)} \times \text{(偏光自由度)} \times \text{(エネルギー量子)}
}
}\end{equation}

\textbf{プランク関数と水素線スペクトルは、同じ電磁場の同じモード構造を、それぞれ「電磁場側から」「物質側から」眺めたものに他ならない。}

これが、本書を「逆引き」で進めてきた到達点である。連続も線も、宇宙のスペクトルとして観測される「結果」から出発して、それぞれの背後にある物理を逐次明らかにし、最終的に
QED で \textbf{ひとつに収束} する。

\subsection{八部構成の総括}\label{ux516bux90e8ux69cbux6210ux306eux7dcfux62ec}

本書の旅を一覧する：

{\def\LTcaptype{none} % do not increment counter
\begin{longtable}[]{@{}lll@{}}
\toprule\noalign{}
部 & 主題 & 到達点 \\
\midrule\noalign{}
\endhead
\bottomrule\noalign{}
\endlastfoot
I & 観測される黒体放射 & プランク関数を観測の言葉で確認 \\
II & 放射場の記述 & 比強度・放射輸送方程式・ソース関数 \\
III & 統計力学への逆引き & 光子の Bose 統計、\(\mu = 0\) の意味 \\
IV & 量子論への逆引き & エネルギー量子 \(h\nu\)、Einstein 係数 \\
V & 電磁気学への逆引き & 偏光、モード密度 \(8\pi\nu^2/c^3\) \\
VI & 現実の天体への戻り & CMB・恒星・降着円盤への適用 \\
VII & 線スペクトルの逆引き & 原子物理・線形成・線輪郭・観測応用 \\
VIII & QED からの俯瞰 & 連続と線が共通の Hamiltonian の二側面と判明 \\
\end{longtable}
}

\begin{tcolorbox}[enhanced jigsaw, arc=.35mm, bottomrule=.15mm, bottomtitle=1mm, breakable, colback=white, colbacktitle=quarto-callout-note-color!10!white, colframe=quarto-callout-note-color-frame, coltitle=black, left=2mm, leftrule=.75mm, opacityback=0, opacitybacktitle=0.6, rightrule=.15mm, title=\textcolor{quarto-callout-note-color}{\faInfo}\hspace{0.5em}{ノート}, titlerule=0mm, toprule=.15mm, toptitle=1mm]

\textbf{使えるようになった道具}（第VIII部 §25.1〜§26.6）：

\begin{itemize}
\tightlist
\item
  電磁場の量子化（モード ＝ 量子調和振動子）
\item
  生成消滅演算子と光子数演算子
\item
  真空状態と零点エネルギー
\item
  Fermi の黄金律から \(A_{ul}\) を導く流れ
\item
  Einstein 関係 \(A_{ul}/B_{ul} = 2h\nu^3/c^2\) の QED 起源
\item
  自発放出 ↔ 真空揺らぎ ↔ 自然幅の三位一体
\item
  電気双極子行列要素から選択則
\item
  連続スペクトルと線スペクトルが \textbf{同じ Hamiltonian}
  から出るという統一的視点
\end{itemize}

\end{tcolorbox}

\begin{center}\rule{0.5\linewidth}{0.5pt}\end{center}

\section*{この章で何がわかったか}\label{sec-qed-minimum-summary}
\addcontentsline{toc}{section}{この章で何がわかったか}

\markright{この章で何がわかったか}

\textbf{本書全体の地図に戻る}

本章で、プランク関数地図と水素線地図が QED でつながった。\(2h\nu^3/c^2\)
という前因子は、両者の地図の \textbf{交点} であり、本書全体が
\textbf{一本に閉じる場所} である。

読者は、観測されるスペクトル ― 連続も線も ―
を見たとき、その背後でどの物理がどう積み重なって式を立ち上げているかを、\textbf{自分の言葉で説明できる}
ようになっているはずである：

\begin{itemize}
\tightlist
\item
  CMB の黒体スペクトル → 電磁場の熱状態 → 宇宙の熱化の証拠
\item
  恒星の Fraunhofer 線 → 原子の束縛状態 + 真空揺らぎ駆動の遷移 →
  元素同定
\item
  AGN の鉄線輪郭 → Doppler + 重力赤方偏移 + 自然幅 →
  中心ブラックホール質量とスピン
\item
  21 cm 線 → 超微細構造（M1 禁制）× 膨大な中性水素 → 宇宙再電離史
\end{itemize}

すべてが共通の物理基盤の上に乗っている。

ここで本書は終わる。しかし、本書を超えてさらに学びたい読者には、付録の参考文献から、より厳密な
QED、より高度な放射輸送、そして個別の天体物理応用へと進む道が用意されている。

\section*{あとがき ― 逆引きの終わりに}\label{sec-qed-minimum-afterword}
\addcontentsline{toc}{section}{あとがき ― 逆引きの終わりに}

\markright{あとがき ― 逆引きの終わりに}

本書は \textbf{「観測される宇宙の熱的放射」} ―
連続放射と線スペクトルの両方 ―
から出発し、その背後にある物理を逆向きにたどることで、現代物理学の主要分野（放射輸送・統計力学・量子論・電磁気学・原子物理・QED）を一筆書きでつないだ。

二つの中心地図 ― プランク関数と水素線吸収係数 ― は、本書の旅の
\textbf{道しるべ}
だった。各章の終わりで地図に戻り、各因子がどこから来ているかを確認することで、読者は理論の細部に迷い込まずに済んだ。

最終章で両地図が QED で交わったとき、本書の三原則 ―（① 逆引き、②
背景物理を曖昧にしない、③ 天下り的に式を与えない）―
が同時に満たされる。観測から出発し、各段階の物理を明示しつつ、最終的に式の必然性が
QED から見通される。

宇宙のスペクトルは、これからも観測装置の進歩とともに新たな現象を見せ続けるだろう。\textbf{読者がそのスペクトルを前にしたとき、本書で得た「逆引きの目」が役に立つ}
― これが著者の願いである。

\section*{演習問題}\label{sec-qed-minimum-exercises}
\addcontentsline{toc}{section}{演習問題}

\markright{演習問題}

\textbf{模範解答}
は折りたたみを展開して確認できる（オンライン版）。まず自力で解いてから開くこと。

\subsection*{問題 26-1　Fermi の黄金律から自発放出率 Aul}\label{ex-26-1}
\addcontentsline{toc}{subsection}{問題 26-1　Fermi
の黄金律から自発放出率 Aul}

{[}★ 難易度：☆☆ {]}

§26.2 は Fermi の黄金律
\(\Gamma_{u\to l}=\dfrac{2\pi}{\hbar}|\langle l;1_{\mathbf k\sigma}|\hat H_\mathrm{int}|u;0\rangle|^2\rho(\omega)\)
から
\(A_{ul}=\dfrac{\omega^3}{3\pi\epsilon_0\hbar c^3}|\langle l|\hat{\mathbf d}|u\rangle|^2\)
が出るとした。構造を読む。

\begin{enumerate}
\def\labelenumi{\arabic{enumi}.}
\tightlist
\item
  自発放出の始状態 \(|u;0\rangle\)（原子上準位＋電磁場真空）と終状態
  \(|l;1_{\mathbf k\sigma}\rangle\)（原子下準位＋光子1個）を述べ、なぜ「真空から光子を生成する」過程なのか説明せよ。
\item
  \(A_{ul}\) の前因子 \(\omega^3\) が「モード密度 \(\propto\nu^2\) ×
  エネルギー量子
  \(h\nu\)」の構造から来ることを述べ、これが第25章でプランク前因子に出た
  \(2h\nu^3/c^2\) と同じ構造であることを指摘せよ。
\item
  古典電磁気では自発放出が起こらないのに QED では \(A_{ul}\neq0\)
  になる物理的理由（真空揺らぎ）を述べよ。
\end{enumerate}

\textbf{関連}：\hyperref[sec-qed-minimum-spontaneous]{§26.2 Fermi
の黄金律}／振動子強度との関係は\hyperref[ex-22-2]{演習 22-2}、Einstein
係数は\hyperref[ex-11-1]{演習 11-1}。

\section{模範解答（問題
26-1）}\label{ux6a21ux7bc4ux89e3ux7b54ux554fux984c-26-1}

\textbf{(1)} 始状態 \(|u;0\rangle\)：原子は上準位
\(u\)、電磁場は真空（光子0個）。終状態
\(|l;1_{\mathbf k\sigma}\rangle\)：原子は下準位 \(l\)、電磁場にモード
\((\mathbf k,\sigma)\) の光子が1個。相互作用 \(\hat H_\mathrm{int}\)
の生成演算子 \(\hat a^\dagger\)
が真空から光子を1個作る。外部光がない真空からの遷移なので「自発」放出。

\textbf{(2)} Fermi の黄金律の終状態モード密度
\(\rho(\omega)\propto\omega^2\)（モード密度
\(\propto\nu^2\)、第14章）と、放出光子のエネルギー \(\hbar\omega=h\nu\)
が掛かって \(\omega^3\) が出る。これはプランク前因子 \(2h\nu^3/c^2\) の
\(\nu^3\)（＝モード密度 \(\nu^2\)×エネルギー量子 \(h\nu\)）と同じ構造。

\textbf{(3)}
古典電磁気では「ゼロ放射場→放出ゼロ」で励起原子が落ちる理由がない。QED
では真空に零点揺らぎ（\(\hat a^\dagger\)
が真空から光子を生成しうること）があり、この真空場との結合が励起原子の遷移を誘起する。だから真空でも
\(A_{ul}\neq0\)。

\textbf{答え}：\(|u;0\rangle\to|l;1\rangle\)＝真空から光子生成。\(\omega^3\)＝モード密度
\(\nu^2\)×\(h\nu\)＝プランク前因子と同構造。真空揺らぎが自発放出を駆動。

\subsection*{問題 26-2　Aul/Bul=2hν³/c² ―
二つの地図の交点}\label{ex-26-2}
\addcontentsline{toc}{subsection}{問題 26-2　Aul/Bul=2hν³/c² ―
二つの地図の交点}

{[}★ 難易度：☆☆☆ {]}

§26.3 は前因子 \(2h\nu^3/c^2\)
がプランク関数地図と水素線地図の共通起源だと総括した。

\begin{enumerate}
\def\labelenumi{\arabic{enumi}.}
\tightlist
\item
  \(2h\nu^3/c^2\)
  の各因子（\(\nu^3\)、\(1/c^2\)、\(2\)、\(h\)）が、プランク関数と自発放出率の両方で同じ物理（モード密度・分散関係・偏光・エネルギー量子）から来ることを表で整理せよ。
\item
  自発放出 \(A_{ul}\) と誘導放出 \(B_{ul}\) が「同じ Fermi
  の黄金律の二側面」である（終状態モードが真空か既に光子を持つか、Bose
  増強因子 \(\sqrt{n+1}\)）ことを述べよ。
\item
  「プランク関数と水素線吸収係数は、同じ電磁場の同じモード構造を電磁場側・物質側から眺めたもの」という総括を、自分の言葉でまとめよ。
\end{enumerate}

\textbf{関連}：\hyperref[sec-qed-minimum-einstein]{§26.3 Einstein
関係式の QED 起源}／プランク前因子の分解は\hyperref[ex-14-4]{演習
14-4}、Einstein 関係は\hyperref[ex-11-1]{演習
11-1}・\hyperref[ex-11-2]{演習 11-2}。

\section{模範解答（問題
26-2）}\label{ux6a21ux7bc4ux89e3ux7b54ux554fux984c-26-2}

\textbf{(1)}

{\def\LTcaptype{none} % do not increment counter
\begin{longtable}[]{@{}
  >{\raggedright\arraybackslash}p{(\linewidth - 4\tabcolsep) * \real{0.3333}}
  >{\raggedright\arraybackslash}p{(\linewidth - 4\tabcolsep) * \real{0.3333}}
  >{\raggedright\arraybackslash}p{(\linewidth - 4\tabcolsep) * \real{0.3333}}@{}}
\toprule\noalign{}
\begin{minipage}[b]{\linewidth}\raggedright
因子
\end{minipage} & \begin{minipage}[b]{\linewidth}\raggedright
プランク関数
\end{minipage} & \begin{minipage}[b]{\linewidth}\raggedright
自発放出率
\end{minipage} \\
\midrule\noalign{}
\endhead
\bottomrule\noalign{}
\endlastfoot
\(\nu^3\) & モード密度 \(\nu^2\) × \(h\nu\) & 終状態モード密度 \(\nu^2\)
× \(h\nu\) \\
\(1/c^2\) & 分散関係 \(k=2\pi\nu/c\) ＋ \(u_\nu\to B_\nu\) & 同じ \\
\(2\) & 偏光自由度（横波） & 同じ \\
\(h\) & エネルギー量子 & 同じ \\
\end{longtable}
}

どちらも「電磁場の同じモード構造」を見ている。

\textbf{(2)} Fermi の黄金律で、終状態のモードが既に \(n\)
個の光子を持つ場合は \(|n\rangle\to|n+1\rangle\) の遷移で Bose 増強因子
\(\sqrt{n+1}\)（→ 誘導放出 \(B_{ul}\)、\(\propto n\)
部分）と自発項（\(+1\) 部分）が同時に出る。終状態が真空なら \(n=0\)
で増強因子1＝自発放出 \(A_{ul}\)。両者の比
\(A_{ul}/B_{ul}=2h\nu^3/c^2\)。

\textbf{(3)}
プランク関数は「電磁場の各モードの熱励起」を電磁場側から見たスペクトル、水素線吸収係数は「原子が電磁場のモードと遷移する」のを物質側から見た量。両者に共通する
\(2h\nu^3/c^2\)
は、同じ電磁場のモード構造（モード密度×偏光×エネルギー量子）の現れ。連続も線も、同じ電磁場--物質相互作用
Hamiltonian の二側面。

\textbf{答え}：\(2h\nu^3/c^2\)
の各因子は両地図で同一物理（モード密度・分散・偏光・量子）。\(A_{ul}\)/\(B_{ul}\)
は Fermi
黄金律の真空/光子有りの二側面（\(\sqrt{n+1}\)）。連続と線は同じモード構造の二つの眺め。

\subsection*{問題 26-3　自発放出 ↔ 真空揺らぎ ↔
自然幅の三位一体}\label{ex-26-3}
\addcontentsline{toc}{subsection}{問題 26-3　自発放出 ↔ 真空揺らぎ ↔
自然幅の三位一体}

{[}★ 難易度：☆☆ {]}

§26.4 は自然幅を QED
的に意味づけた。Ly\(\alpha\)（\(A_{ul}=6.27\times10^8\)
s\(^{-1}\)）を例に三者の関係を整理する。

\begin{enumerate}
\def\labelenumi{\arabic{enumi}.}
\tightlist
\item
  真空揺らぎ → 自発放出 → 励起寿命 \(\tau=1/A_{ul}\) → 自然幅
  \(\gamma_\mathrm{nat}=A_{ul}/(2\pi)\) という連鎖を述べよ。
\item
  「線の自然幅は真空揺らぎの存在を観測している」という主張の意味を述べよ。観測例（QSO
  の damped Ly\(\alpha\) 系の減衰翼）に触れよ。
\item
  21 cm 線（\(A_{ul}\simeq2.9\times10^{-15}\)
  s\(^{-1}\)）の自然幅は極端に小さいのに観測される。その理由（膨大な中性水素）を演習
  20-1 と結びつけて述べよ。
\end{enumerate}

\textbf{関連}：\hyperref[sec-qed-minimum-natural-width]{§26.4
自然幅と真空揺らぎ}／自然幅の計算は\hyperref[ex-23-2]{演習 23-2}、21 cm
は\hyperref[ex-20-1]{演習 20-1}。

\section{模範解答（問題
26-3）}\label{ux6a21ux7bc4ux89e3ux7b54ux554fux984c-26-3}

\textbf{(1)} 量子化された電磁場の真空には各モードに零点揺らぎがある →
この真空場との相互作用が励起原子の遷移を誘起して自発放出を生む（\(A_{ul}\)）→
励起準位の寿命が有限 \(\tau=1/A_{ul}\) → 不確定性関係でエネルギー幅
\(\Delta E\sim\hbar A_{ul}\)、振動数幅
\(\gamma_\mathrm{nat}=A_{ul}/(2\pi)\) が線に現れる。

\textbf{(2)}
自然幅の物理的起源は真空揺らぎが励起原子を揺さぶること。よって観測される自然幅（Lorentzian
の翼）は、真空が「空ではない」ことの直接の観測的痕跡。QSO の damped
Ly\(\alpha\) 系の減衰翼は、まさにこの自然幅起源の Lorentzian
翼を地上の分光計で測ったもの。

\textbf{(3)} 21 cm 線の \(A_{ul}\) は極小（自然幅も極小）だが、中心地図
\(\alpha_\nu=(\pi e^2/m_ec)f_{lu}n_l\phi\) の
\(n_l\)（中性水素の柱密度）が宇宙的スケールで膨大なので、積として観測可能な信号になる（演習
20-1）。真空揺らぎの効果がいかに弱くても、宇宙的物質量では無視できない。

\textbf{答え}：真空揺らぎ→自発放出 \(A_{ul}\)→寿命
\(\tau=1/A_{ul}\)→自然幅
\(A_{ul}/2\pi\)。自然幅＝真空揺らぎの観測（damped Ly\(\alpha\) 翼）。21
cm は膨大な \(n_l\) で観測可能。

\subsection*{問題
26-4　選択則を行列要素で、そして連続と線の統一}\label{ex-26-4}
\addcontentsline{toc}{subsection}{問題
26-4　選択則を行列要素で、そして連続と線の統一}

{[}★ 難易度：☆☆ {]}

§26.5・25.6 の総括。電気双極子行列要素
\(\mathbf d_{ul}=\langle u|\hat{\mathbf d}|l\rangle=-e\langle u|\hat{\mathbf r}|l\rangle\)
と本書全体の統一。

\begin{enumerate}
\def\labelenumi{\arabic{enumi}.}
\tightlist
\item
  電気双極子（E1）選択則 \(\Delta l=\pm1\)、\(\Delta m=0,\pm1\)
  が、相互作用の最低次（双極子）とパリティ・角運動量保存から自然に出ること、禁制遷移が高次（M1/E2）として残ることを述べよ。
\item
  第22章で「経験則」として述べた選択則が、本章で QED
  構造から導かれることが、本書の原則③「天下り的に式を与えない」とどう関係するか述べよ。
\item
  本書の総括として、連続スペクトル（プランク関数）と線スペクトル（水素線吸収係数）が「同じ電磁場--物質相互作用
  Hamiltonian
  の二側面」であることを、これまでの全演習を踏まえて自分の言葉でまとめよ。
\end{enumerate}

\textbf{関連}：\hyperref[sec-qed-minimum-selection]{§26.5
選択則}、\hyperref[sec-qed-minimum-unification]{§26.6
連続と線の統一}／選択則の原子物理側は\hyperref[ex-22-3]{演習
22-3}、全因子の合流は\hyperref[ex-14-4]{演習 14-4}。

\section{模範解答（問題
26-4）}\label{ux6a21ux7bc4ux89e3ux7b54ux554fux984c-26-4}

\textbf{(1)} 相互作用 \(\hat H_\mathrm{int}\) の長波長極限が電気双極子
\(\hat{\mathbf d}\)
を介する（最低次）。\(\langle u|\hat{\mathbf d}|l\rangle\)
はパリティと角運動量の保存から \(\Delta l=\pm1\)、\(\Delta m=0,\pm1\)
で非ゼロ。これを満たさない遷移は E1
で行列要素ゼロ＝禁制だが、磁気双極子(M1)・電気四重極(E2)など高次過程として弱く残る（強度
\(\sim\alpha^2\) 倍）。

\textbf{(2)}
第22章では選択則を観測規則（経験則）として受け入れたが、本章で「電気双極子相互作用の最低次」という
QED
構造から必然として導いた。これは本書の原則③「天下り的に式を与えず、必然性として理解させる」を線スペクトル側で最終的に実行したもの。観測規則の背後にある物理が明示された。

\textbf{(3)}
プランク関数（連続）は電磁場の各モードの熱励起を電磁場側から、水素線吸収係数（線）は原子と電磁場モードの遷移を物質側から見た量。両者は同じ相互作用
Hamiltonian
\(\hat H_\mathrm{int}=-(e/m_e)\hat{\mathbf A}\cdot\hat{\mathbf p}\)
から出て、共通の前因子
\(2h\nu^3/c^2\)（モード密度×偏光×エネルギー量子）を持つ。観測されるスペクトル
― CMB の黒体、恒星の Fraunhofer 線、AGN の鉄線、21 cm ―
はすべて、量子化された電磁場のモード構造の現れであり、観測者は毎日 QED
を見ている。これが本書の逆引きの到達点。

\textbf{答え}：E1
選択則は双極子の最低次＋保存則から、禁制は高次として残る。経験則を QED
から導く＝原則③の完成。連続も線も同じ \(\hat H_\mathrm{int}\)
の二側面、共通前因子 \(2h\nu^3/c^2\) が統一の印。

\section*{さらに学ぶための参考文献}\label{sec-qed-minimum-further}
\addcontentsline{toc}{section}{さらに学ぶための参考文献}

\markright{さらに学ぶための参考文献}

\begin{itemize}
\tightlist
\item
  Sakurai \& Napolitano, \emph{Modern Quantum Mechanics} (Cambridge, 3rd
  ed., 2020) --- 放射場と物質の相互作用の章、本章のレベルから一段深く
\item
  Loudon, \emph{The Quantum Theory of Light} (Oxford, 3rd ed., 2000) ---
  自発放出・自然幅の QED 的扱い、本書と相性が良い
\item
  Cohen-Tannoudji, Dupont-Roc \& Grynberg, \emph{Photons and Atoms} /
  \emph{Atom-Photon Interactions} (Wiley, 1989/1992) ---
  光と原子の相互作用の体系的扱い
\item
  Rybicki \& Lightman, \emph{Radiative Processes in Astrophysics}
  (Wiley, 1979) --- 天体物理応用での放射理論の標準教科書
\item
  Peskin \& Schroeder, \emph{An Introduction to Quantum Field Theory}
  (Westview, 1995) --- 本格的な場の量子論
\item
  Weinberg, \emph{The Quantum Theory of Fields} vol.~I-III (Cambridge,
  1995-2000) --- さらに深く進みたい読者へ
\end{itemize}

\bookmarksetup{startatroot}

\chapter{参考文献}\label{ux53c2ux8003ux6587ux732e}

\bookmarksetup{startatroot}

\chapter*{参考文献}\label{sec-references}
\addcontentsline{toc}{chapter}{参考文献}

\markboth{参考文献}{参考文献}

\protect\phantomsection\label{refs}

\cleardoublepage
\phantomsection
\addcontentsline{toc}{part}{Appendices}
\appendix

\chapter{付録
A　数学補章}\label{ux4ed8ux9332-a-ux6570ux5b66ux88dcux7ae0}

\chapter*{付録 A　数学補章}\label{sec-math-notes}
\addcontentsline{toc}{chapter}{付録 A　数学補章}

\markboth{付録 A　数学補章}{付録 A　数学補章}

\textbf{この付録の主題}：本書で使う数学の最低限の整理。本文の論旨を追うのに必要な道具を、定義
→ 結果 →
物理的意味の順でまとめる。厳密な証明や一般論は数学の専門書（小平、Arfken
\& Weber、Reed \& Simon 等）に譲り、ここでは
\textbf{本書のどの章のどの式で使われるか} を併記する実用版とする。

\begin{center}\rule{0.5\linewidth}{0.5pt}\end{center}

\section{A.1 立体角と球面積分}\label{sec-math-solid-angle}

\subsection{A.1.1
立体角の定義}\label{a.1.1-ux7acbux4f53ux89d2ux306eux5b9aux7fa9}

\textbf{立体角}（\textbf{solid angle},
\(\Omega\)）は、点から見た面積の「広がり」を表す無次元量。単位は
\textbf{ステラジアン}（\textbf{steradian}, \textbf{sr}）。

半径 \(r\) の球面上で、面積 \(dA\) を覆う部分の立体角は

\begin{equation}\protect\phantomsection\label{eq-solid-angle-def}{
d\Omega = \frac{dA}{r^2}
}\end{equation}

球全体の立体角は \(4\pi\) sr。

球面座標 \((\theta, \phi)\) を用いると、

\begin{equation}\protect\phantomsection\label{eq-solid-angle-spherical}{
d\Omega = \sin\theta\, d\theta\, d\phi
}\end{equation}

\subsection{\texorpdfstring{A.1.2 半球での積分（Stefan-Boltzmann の
\(\pi\)
因子）}{A.1.2 半球での積分（Stefan-Boltzmann の \textbackslash pi 因子）}}\label{a.1.2-ux534aux7403ux3067ux306eux7a4dux5206stefan-boltzmann-ux306e-pi-ux56e0ux5b50}

第4章 §4.4 で出てくる \textbf{半球フラックス積分}：

\begin{equation}\protect\phantomsection\label{eq-hemisphere-flux}{
F = \int_\text{半球} I \cos\theta\, d\Omega = \int_0^{2\pi}\int_0^{\pi/2} I \cos\theta \sin\theta\, d\theta\, d\phi
}\end{equation}

等方放射場 \$I = \$ const の場合：

\begin{equation}\protect\phantomsection\label{eq-hemisphere-flux-result}{
F = 2\pi I \int_0^{\pi/2} \cos\theta \sin\theta\, d\theta = 2\pi I \cdot \frac{1}{2} = \pi I
}\end{equation}

この \(\pi\) が、Stefan-Boltzmann 則 \(F = \sigma T^4 = \pi B(T)\)
に現れる。

\textbf{本文での使用}：第4章 §4.4（Stefan-Boltzmann の \(\pi\)
因子）、第14章 §14.5（\(B_\nu = c u_\nu/4\pi\)）

\subsection{A.1.3
立体角と方向余弦}\label{a.1.3-ux7acbux4f53ux89d2ux3068ux65b9ux5411ux4f59ux5f26}

平面平行近似の放射輸送（第5章）では、\(z\) 軸方向の余弦
\(\mu \equiv \cos\theta\) を変数とし、

\[
d\mu = -\sin\theta\, d\theta, \quad d\Omega = -d\mu\, d\phi
\]

積分は \(\mu \in [-1, 1]\) に対して行う。等方放射場の平均強度

\begin{equation}\protect\phantomsection\label{eq-mean-intensity}{
J \equiv \frac{1}{4\pi}\int I\, d\Omega = \frac{1}{2}\int_{-1}^{1} I\, d\mu
}\end{equation}

これが第4章 §4.5 の \(J_\nu\) の定義。

\begin{center}\rule{0.5\linewidth}{0.5pt}\end{center}

\section{A.2 ガウス積分とガンマ関数}\label{sec-math-gaussian}

\subsection{A.2.1
ガウス積分}\label{a.2.1-ux30acux30a6ux30b9ux7a4dux5206}

\begin{equation}\protect\phantomsection\label{eq-gaussian-basic}{
\int_{-\infty}^{\infty} e^{-x^2}\, dx = \sqrt{\pi}
}\end{equation}

より一般に、

\begin{equation}\protect\phantomsection\label{eq-gaussian-moments}{
\int_{-\infty}^{\infty} e^{-ax^2}\, dx = \sqrt{\pi/a}, \quad \int_{-\infty}^{\infty} x^2 e^{-ax^2}\, dx = \frac{1}{2a}\sqrt{\pi/a}
}\end{equation}

\textbf{本文での使用}：第23章 §23.2（Doppler 線形 Gaussian
の規格化）、第3章（Maxwell-Boltzmann 速度分布）

\subsection{A.2.2
ガンマ関数}\label{a.2.2-ux30acux30f3ux30deux95a2ux6570}

ガンマ関数の定義：

\begin{equation}\protect\phantomsection\label{eq-gamma-function}{
\Gamma(s) \equiv \int_0^\infty x^{s-1} e^{-x}\, dx, \quad \text{Re}(s) > 0
}\end{equation}

性質：

\begin{itemize}
\tightlist
\item
  \(\Gamma(n) = (n-1)!\) for positive integer \(n\)
\item
  \(\Gamma(1/2) = \sqrt{\pi}\)
\item
  \(\Gamma(s+1) = s\Gamma(s)\)（再帰関係）
\end{itemize}

\subsection{A.2.3 Riemann
ゼータ関数}\label{a.2.3-riemann-ux30bcux30fcux30bfux95a2ux6570}

\begin{equation}\protect\phantomsection\label{eq-zeta-function}{
\zeta(s) \equiv \sum_{n=1}^\infty \frac{1}{n^s} = \frac{1}{\Gamma(s)}\int_0^\infty \frac{x^{s-1}}{e^x - 1}\, dx
}\end{equation}

代表値：

\begin{itemize}
\tightlist
\item
  \(\zeta(2) = \pi^2/6\)
\item
  \(\zeta(3) \simeq 1.20206\)
\item
  \(\zeta(4) = \pi^4/90\)
\end{itemize}

\subsection{A.2.4
プランク積分}\label{a.2.4-ux30d7ux30e9ux30f3ux30afux7a4dux5206}

Stefan-Boltzmann 則（第14章 §14.5）で使う積分：

\begin{equation}\protect\phantomsection\label{eq-planck-integral}{
\int_0^\infty \frac{x^3}{e^x - 1}\, dx = \Gamma(4)\zeta(4) = 6 \cdot \frac{\pi^4}{90} = \frac{\pi^4}{15}
}\end{equation}

これからプランク分布の全エネルギー密度

\begin{equation}\protect\phantomsection\label{eq-energy-density-T4}{
u(T) = \int_0^\infty u_\nu\, d\nu = \frac{8\pi^5 k_B^4}{15 h^3 c^3} T^4 = aT^4
}\end{equation}

放射定数 \(a \simeq 7.566 \times 10^{-16}\) J m\(^{-3}\)
K\(^{-4}\)。フラックス側では \(\sigma = ac/4\) で Stefan-Boltzmann
定数。

\textbf{本文での使用}：第4章 §4.4、第14章 §14.5、第17章
§17.1（恒星の有効温度）

\begin{center}\rule{0.5\linewidth}{0.5pt}\end{center}

\section{A.3 統計力学のための組合せ論}\label{sec-math-combinatorics}

\subsection{A.3.1
三つの分布}\label{a.3.1-ux4e09ux3064ux306eux5206ux5e03}

第7-8章で扱った三つの統計分布：

\textbf{Maxwell-Boltzmann}（\textbf{MB}, 古典）：

\begin{equation}\protect\phantomsection\label{eq-MB-distribution}{
\bar{n}_\text{MB}(\epsilon) = e^{-(\epsilon - \mu)/k_BT}
}\end{equation}

\textbf{Bose-Einstein}（\textbf{BE}, ボソン）：

\begin{equation}\protect\phantomsection\label{eq-BE-distribution}{
\bar{n}_\text{BE}(\epsilon) = \frac{1}{e^{(\epsilon - \mu)/k_BT} - 1}
}\end{equation}

\textbf{Fermi-Dirac}（\textbf{FD}, フェルミオン）：

\begin{equation}\protect\phantomsection\label{eq-FD-distribution}{
\bar{n}_\text{FD}(\epsilon) = \frac{1}{e^{(\epsilon - \mu)/k_BT} + 1}
}\end{equation}

光子（\(\mu = 0\)、ボソン）では 式~\ref{eq-BE-distribution}
が中心地図プランク関数の指数部分。

\subsection{A.3.2 Boltzmann
分布の導出（簡単版）}\label{a.3.2-boltzmann-ux5206ux5e03ux306eux5c0eux51faux7c21ux5358ux7248}

孤立系のエネルギー保存を考慮した状態数を最大化する手続き（Lagrange
未定乗数）から、

\begin{equation}\protect\phantomsection\label{eq-boltzmann-ratio}{
\frac{n_u}{n_l} = \frac{g_u}{g_l} e^{-(\epsilon_u - \epsilon_l)/k_BT}
}\end{equation}

これが第21章 §21.4 で LTE 線形成に使った Boltzmann 比。

\subsection{A.3.3 Saha
の式（電離平衡）}\label{a.3.3-saha-ux306eux5f0fux96fbux96e2ux5e73ux8861}

イオン化反応 \(X \leftrightarrow X^+ + e^-\) の電離平衡：

\begin{equation}\protect\phantomsection\label{eq-saha}{
\frac{n_{X^+} n_e}{n_X} = \left(\frac{2\pi m_e k_B T}{h^2}\right)^{3/2} \cdot 2 \cdot \frac{g_{X^+}}{g_X} e^{-\chi/k_BT}
}\end{equation}

\(\chi\) は電離エネルギー、\(g\) は各状態の縮重度。

\textbf{本文での使用}：第19章
§19.5（連続放射の電離平衡と自己無撞着）、第20章
§20.2（恒星スペクトル分類 OBAFGKM での水素線強度ピーク）、第24章
§24.4（電離度診断）

\begin{center}\rule{0.5\linewidth}{0.5pt}\end{center}

\section{A.4 Fourier 変換と Voigt 関数}\label{sec-math-fourier}

\subsection{A.4.1 Fourier
変換の規約}\label{a.4.1-fourier-ux5909ux63dbux306eux898fux7d04}

時間 ↔ 振動数の Fourier 変換（本書での規約）：

\begin{equation}\protect\phantomsection\label{eq-fourier-pair}{
\tilde{f}(\nu) = \int_{-\infty}^\infty f(t) e^{-2\pi i \nu t}\, dt, \quad f(t) = \int_{-\infty}^\infty \tilde{f}(\nu) e^{2\pi i \nu t}\, d\nu
}\end{equation}

\subsection{A.4.2 畳み込み定理（Voigt
プロファイルへの応用）}\label{a.4.2-ux7573ux307fux8fbcux307fux5b9aux7406voigt-ux30d7ux30edux30d5ux30a1ux30a4ux30ebux3078ux306eux5fdcux7528}

二つの関数の \textbf{畳み込み}：

\begin{equation}\protect\phantomsection\label{eq-convolution-def}{
(f \otimes g)(x) \equiv \int f(x') g(x - x')\, dx'
}\end{equation}

Fourier 変換は畳み込みを積に変換：

\begin{equation}\protect\phantomsection\label{eq-convolution-theorem}{
\widetilde{f \otimes g} = \tilde{f} \cdot \tilde{g}
}\end{equation}

\subsection{A.4.3 Voigt
プロファイル}\label{a.4.3-voigt-ux30d7ux30edux30d5ux30a1ux30a4ux30eb}

第23章 §23.4 で出てきた \textbf{Voigt プロファイル}
\(\phi_V = \phi_L \otimes \phi_D\) は、Lorentzian と Gaussian
の畳み込み：

\begin{equation}\protect\phantomsection\label{eq-voigt-convolution-A}{
\phi_V(\nu) = \int_{-\infty}^\infty \phi_L(\nu - \nu') \phi_D(\nu')\, d\nu'
}\end{equation}

無次元パラメータ

\[
a = \frac{\gamma}{2\Delta\nu_D}, \quad u = \frac{\nu - \nu_0}{\Delta\nu_D}
\]

を使うと、Voigt 関数 \(H(a, u)\) で表される：

\begin{equation}\protect\phantomsection\label{eq-voigt-H}{
H(a, u) = \frac{a}{\pi}\int_{-\infty}^\infty \frac{e^{-y^2}}{(u-y)^2 + a^2}\, dy
}\end{equation}

特殊極限：

\begin{itemize}
\tightlist
\item
  \(a \to 0\)：Doppler コア \(H(0, u) = e^{-u^2}\)
\item
  \(|u| \gg 1\)：Lorentz 翼 \(H(a, u) \to a/[\pi(u^2 + a^2)]\)
\end{itemize}

数値評価は Faddeeva 関数 \(w(z) = e^{-z^2}\text{erfc}(-iz)\)
を介して行うのが標準（SciPy \texttt{scipy.special.wofz}）。

\textbf{本文での使用}：第23章 §23.4（線形の合成）

\begin{center}\rule{0.5\linewidth}{0.5pt}\end{center}

\section{A.5 演算子の交換関係 ―
量子力学の基礎道具}\label{sec-math-commutators}

\subsection{A.5.1
基本交換関係}\label{a.5.1-ux57faux672cux4ea4ux63dbux95a2ux4fc2}

位置・運動量の正準交換関係：

\begin{equation}\protect\phantomsection\label{eq-canonical-commutator}{
[\hat{x}, \hat{p}] = i\hbar
}\end{equation}

これから生まれる不確定性関係：

\begin{equation}\protect\phantomsection\label{eq-uncertainty}{
\Delta x \cdot \Delta p \ge \frac{\hbar}{2}
}\end{equation}

\subsection{A.5.2
生成消滅演算子}\label{a.5.2-ux751fux6210ux6d88ux6ec5ux6f14ux7b97ux5b50}

調和振動子の生成・消滅演算子 \(\hat{a}^\dagger, \hat{a}\) の交換関係：

\begin{equation}\protect\phantomsection\label{eq-bosonic-commutator}{
[\hat{a}, \hat{a}^\dagger] = 1, \quad [\hat{a}, \hat{a}] = [\hat{a}^\dagger, \hat{a}^\dagger] = 0
}\end{equation}

これが Bose 統計の根本。光子は \textbf{電磁場モードの調和振動子の励起}
として現れる（第25章 §25.2）。

\subsection{A.5.3 数演算子}\label{a.5.3-ux6570ux6f14ux7b97ux5b50}

\(\hat{n} \equiv \hat{a}^\dagger \hat{a}\) を
\textbf{数演算子}（\textbf{number operator}）と呼ぶ。固有値は非負整数
\(0, 1, 2, \ldots\) で、各々が「そのモードに \(n\)
個の光子が存在する状態」を表す。

\begin{equation}\protect\phantomsection\label{eq-number-action}{
\hat{n}|n\rangle = n|n\rangle, \quad \hat{a}^\dagger|n\rangle = \sqrt{n+1}\,|n+1\rangle, \quad \hat{a}|n\rangle = \sqrt{n}\,|n-1\rangle
}\end{equation}

\subsection{A.5.4 Bose
増強因子}\label{a.5.4-bose-ux5897ux5f37ux56e0ux5b50}

式~\ref{eq-number-action} の係数 \(\sqrt{n+1}\) が、\textbf{Bose
増強}（\textbf{Bose enhancement}）の起源。誘導放出率 \(B_{ul}\)
がモード光子数 \(n\) に比例することの数学的表現：

\begin{equation}\protect\phantomsection\label{eq-bose-enhancement}{
|\langle n+1|\hat{a}^\dagger|n\rangle|^2 = n + 1
}\end{equation}

「\(+1\)」が自発放出（真空 \(|0\rangle\) からの \(|1\rangle\)
への遷移）に、「\(n\)」が誘導放出に対応する。

\textbf{本文での使用}：第8章（光子の Bose 統計）、第25章
§25.2-§25.3（モード ＝ 量子調和振動子、真空状態）、第26章
§26.2（自発放出と誘導放出の統一）

\begin{center}\rule{0.5\linewidth}{0.5pt}\end{center}

\section{A.6 Maxwell 方程式と波動方程式 ―
ベクトル解析の道具}\label{sec-math-vector-calc}

\subsection{A.6.1 真空中の Maxwell
方程式（SI）}\label{a.6.1-ux771fux7a7aux4e2dux306e-maxwell-ux65b9ux7a0bux5f0fsi}

第13章 §13.1 と同じ：

\[
\nabla \cdot \mathbf{E} = 0, \quad \nabla \cdot \mathbf{B} = 0
\] \begin{equation}\protect\phantomsection\label{eq-maxwell-vacuum-app}{
\nabla \times \mathbf{E} = -\frac{\partial \mathbf{B}}{\partial t}, \quad \nabla \times \mathbf{B} = \mu_0\epsilon_0 \frac{\partial \mathbf{E}}{\partial t}
}\end{equation}

\subsection{A.6.2
ベクトル恒等式}\label{a.6.2-ux30d9ux30afux30c8ux30ebux6052ux7b49ux5f0f}

\begin{equation}\protect\phantomsection\label{eq-curl-curl}{
\nabla \times (\nabla \times \mathbf{A}) = \nabla(\nabla \cdot \mathbf{A}) - \nabla^2 \mathbf{A}
}\end{equation}

これを 式~\ref{eq-maxwell-vacuum-app} の curl-of-curl
に使うと、波動方程式

\begin{equation}\protect\phantomsection\label{eq-wave-equation-app}{
\nabla^2 \mathbf{E} - \frac{1}{c^2}\frac{\partial^2 \mathbf{E}}{\partial t^2} = 0, \quad c = \frac{1}{\sqrt{\mu_0\epsilon_0}}
}\end{equation}

が出る。これが第13章 §13.2 の出発点。

\subsection{A.6.3
球面波と平面波}\label{a.6.3-ux7403ux9762ux6ce2ux3068ux5e73ux9762ux6ce2}

無限遠から見れば球面波は平面波で近似できる。平面波解

\begin{equation}\protect\phantomsection\label{eq-plane-wave-app}{
\mathbf{E}(\mathbf{r}, t) = \mathbf{E}_0 e^{i(\mathbf{k}\cdot\mathbf{r} - \omega t)}, \quad \omega = c|\mathbf{k}|
}\end{equation}

これが第13章 §13.3 の中心式。

\begin{center}\rule{0.5\linewidth}{0.5pt}\end{center}

\section{A.7 球面調和関数 ―
水素原子と回転状態}\label{sec-math-spherical-harmonics}

\subsection{A.7.1
球面調和関数の定義}\label{a.7.1-ux7403ux9762ux8abfux548cux95a2ux6570ux306eux5b9aux7fa9}

\begin{equation}\protect\phantomsection\label{eq-spherical-harmonics}{
Y_{lm}(\theta, \phi) = \sqrt{\frac{(2l+1)}{4\pi}\frac{(l-m)!}{(l+m)!}} P_l^m(\cos\theta) e^{im\phi}
}\end{equation}

\(P_l^m\) は Legendre 陪関数、\(l = 0, 1, 2, \ldots\)
は角運動量量子数、\(m = -l, \ldots, l\) は磁気量子数。

\subsection{A.7.2
直交性と完備性}\label{a.7.2-ux76f4ux4ea4ux6027ux3068ux5b8cux5099ux6027}

\[
\int Y_{l'm'}^* Y_{lm}\, d\Omega = \delta_{l'l} \delta_{m'm}, \quad \sum_{l, m} Y_{lm}(\theta, \phi) Y_{lm}^*(\theta', \phi') = \delta(\Omega - \Omega')
\]

\subsection{A.7.3
水素原子と回転分子への応用}\label{a.7.3-ux6c34ux7d20ux539fux5b50ux3068ux56deux8ee2ux5206ux5b50ux3078ux306eux5fdcux7528}

\begin{itemize}
\tightlist
\item
  \textbf{水素原子}（第22章 §22.3）：波動関数
  \(\psi_{nlm} = R_{nl}(r) Y_{lm}(\theta, \phi)\)
\item
  \textbf{二原子分子の回転}：エネルギー
  \(E_J = \hbar^2 J(J+1)/(2I)\)、\(J\) は回転量子数
\end{itemize}

選択則 \(\Delta l = \pm 1\)（第22章 §22.5）は、球面調和関数の積分
\(\int Y_{l'm'}^* \mathbf{r} Y_{lm}\, d\Omega\) がいつ非零かで決まる。

\textbf{本文での使用}：第22章
§22.3-§22.5（水素原子の波動関数と選択則）、第23章 §23.6（Zeeman 分裂）

\begin{center}\rule{0.5\linewidth}{0.5pt}\end{center}

\section{A.8 不確定性原理と寿命幅}\label{sec-math-uncertainty-lifetime}

\subsection{A.8.1
エネルギー-時間の不確定性}\label{a.8.1-ux30a8ux30cdux30ebux30aeux30fc-ux6642ux9593ux306eux4e0dux78baux5b9aux6027}

正準対 \([\hat{x}, \hat{p}] = i\hbar\) から
\(\Delta x\Delta p \ge \hbar/2\)
が出るのに対し、エネルギー-時間の関係はやや異なる文脈で：

\begin{equation}\protect\phantomsection\label{eq-energy-time-uncertainty}{
\Delta E \cdot \tau \gtrsim \hbar
}\end{equation}

ここで \(\tau\) は \textbf{励起状態の寿命}、\(\Delta E\)
はエネルギー幅。これが本文の自然幅（第23章 §23.1）：

\begin{equation}\protect\phantomsection\label{eq-natural-width-app}{
\Gamma_\text{nat} = \frac{1}{2\pi\tau} = \frac{A_{ul}}{2\pi}
}\end{equation}

\subsection{A.8.2 Lorentzian の Fourier
起源}\label{a.8.2-lorentzian-ux306e-fourier-ux8d77ux6e90}

寿命 \(\tau\) で指数減衰する古典場
\(E(t) \propto e^{-t/(2\tau)} \cos(\omega_0 t)\) の Fourier 変換は
\textbf{Lorentzian}：

\begin{equation}\protect\phantomsection\label{eq-lorentzian-fourier}{
|\tilde{E}(\omega)|^2 \propto \frac{1}{(\omega - \omega_0)^2 + (1/2\tau)^2}
}\end{equation}

これが第23章 §23.1 の自然幅 Lorentzian の \textbf{古典的起源}。QED
では真空揺らぎが減衰を駆動する（第26章 §26.4）。

\textbf{本文での使用}：第23章 §23.1（自然幅）、第26章
§26.4（真空揺らぎと自然幅）

\begin{center}\rule{0.5\linewidth}{0.5pt}\end{center}

\section{A.9 数値計算の常識}\label{sec-math-numerical}

本書の各章で出てくる計算を実際に行うときの実用メモ。

\subsection{A.9.1
プランク関数の数値評価}\label{a.9.1-ux30d7ux30e9ux30f3ux30afux95a2ux6570ux306eux6570ux5024ux8a55ux4fa1}

指数因子 \(h\nu/k_BT\) が大きいとき（Wien
側、\(h\nu/k_BT \gtrsim 50\)）はオーバーフローしないよう注意：

\begin{Shaded}
\begin{Highlighting}[]
\ImportTok{import}\NormalTok{ numpy }\ImportTok{as}\NormalTok{ np}
\KeywordTok{def}\NormalTok{ planck\_safe(nu, T, h}\OperatorTok{=}\FloatTok{6.626e{-}34}\NormalTok{, k}\OperatorTok{=}\FloatTok{1.381e{-}23}\NormalTok{, c}\OperatorTok{=}\FloatTok{2.998e8}\NormalTok{):}
\NormalTok{    x }\OperatorTok{=}\NormalTok{ h}\OperatorTok{*}\NormalTok{nu}\OperatorTok{/}\NormalTok{(k}\OperatorTok{*}\NormalTok{T)}
    \CommentTok{\# オーバーフロー防止}
\NormalTok{    x }\OperatorTok{=}\NormalTok{ np.where(x }\OperatorTok{\textgreater{}} \DecValTok{700}\NormalTok{, }\DecValTok{700}\NormalTok{, x)}
    \ControlFlowTok{return} \DecValTok{2}\OperatorTok{*}\NormalTok{h}\OperatorTok{*}\NormalTok{nu}\OperatorTok{**}\DecValTok{3}\OperatorTok{/}\NormalTok{c}\OperatorTok{**}\DecValTok{2} \OperatorTok{/}\NormalTok{ (np.exp(x) }\OperatorTok{{-}} \DecValTok{1}\NormalTok{)}
\end{Highlighting}
\end{Shaded}

\subsection{A.9.2 Voigt
プロファイルの数値評価}\label{a.9.2-voigt-ux30d7ux30edux30d5ux30a1ux30a4ux30ebux306eux6570ux5024ux8a55ux4fa1}

\texttt{scipy.special.wofz} を使う：

\begin{Shaded}
\begin{Highlighting}[]
\ImportTok{from}\NormalTok{ scipy.special }\ImportTok{import}\NormalTok{ wofz}
\KeywordTok{def}\NormalTok{ voigt\_profile(nu, nu0, sigma, gamma):}
\NormalTok{    z }\OperatorTok{=}\NormalTok{ ((nu }\OperatorTok{{-}}\NormalTok{ nu0) }\OperatorTok{+} \OtherTok{1j}\OperatorTok{*}\NormalTok{gamma) }\OperatorTok{/}\NormalTok{ (sigma}\OperatorTok{*}\NormalTok{np.sqrt(}\DecValTok{2}\NormalTok{))}
    \ControlFlowTok{return}\NormalTok{ np.real(wofz(z)) }\OperatorTok{/}\NormalTok{ (sigma}\OperatorTok{*}\NormalTok{np.sqrt(}\DecValTok{2}\OperatorTok{*}\NormalTok{np.pi))}
\end{Highlighting}
\end{Shaded}

\subsection{A.9.3
単位系の混在に注意}\label{a.9.3-ux5358ux4f4dux7cfbux306eux6df7ux5728ux306bux6ce8ux610f}

CGS と SI の混在は計算ミスの典型源。本書の数値例は \textbf{SI を基本}
にする（付録 B）。文献によっては
CGS（電磁気学では特に）が使われるので、論文との比較時には常に単位系を確認すること。

\textbf{実装例}：本書付属のコード
\texttt{code/planck\_curve.py}、\texttt{code/voigt\_profile.py} 等を参照

\begin{center}\rule{0.5\linewidth}{0.5pt}\end{center}

\section*{A.10 さらに進むための参考文献}\label{sec-math-references}
\addcontentsline{toc}{section}{A.10 さらに進むための参考文献}

\markright{A.10 さらに進むための参考文献}

\begin{itemize}
\tightlist
\item
  Arfken, Weber \& Harris, \emph{Mathematical Methods for Physicists}
  (Academic Press, 7th ed., 2012) --- 物理学者のための数学の標準書
\item
  小平邦彦『複素解析』『微分積分学』（岩波書店）--- 厳密な数学基礎
\item
  河村哲也『数値計算のためのFortran77/90/95』--- 数値計算の実装面
\item
  Press et al., \emph{Numerical Recipes} (Cambridge, 3rd ed., 2007) ---
  物理計算アルゴリズム集
\item
  Olver et al., \emph{NIST Handbook of Mathematical Functions}
  (Cambridge, 2010) --- Voigt
  関数等の特殊関数の定本（\url{https://dlmf.nist.gov}）
\end{itemize}

\chapter{付録
B　単位と基本定数}\label{ux4ed8ux9332-b-ux5358ux4f4dux3068ux57faux672cux5b9aux6570}

\chapter*{付録 B　単位と基本定数}\label{sec-units-constants}
\addcontentsline{toc}{chapter}{付録 B　単位と基本定数}

\markboth{付録 B　単位と基本定数}{付録 B　単位と基本定数}

\textbf{この付録の主題}：本書で使う物理基本定数、天文単位、エネルギー・振動数・波長の相互換算、CGS
と SI の対応をまとめる。値は CODATA 2018 と IAU 2015
を基準とする。本文中の数値例で「自分で計算したい」読者のために、桁感覚も含めて整理する。

\begin{center}\rule{0.5\linewidth}{0.5pt}\end{center}

\section{B.1 物理基本定数}\label{sec-fundamental-constants}

\subsection{B.1.1 SI 単位での値（CODATA
2018）}\label{b.1.1-si-ux5358ux4f4dux3067ux306eux5024codata-2018}

{\def\LTcaptype{none} % do not increment counter
\begin{longtable}[]{@{}llll@{}}
\toprule\noalign{}
定数 & 記号 & SI 値 & 単位 \\
\midrule\noalign{}
\endhead
\bottomrule\noalign{}
\endlastfoot
光速度 & \(c\) & \(2.998 \times 10^8\) & m s\(^{-1}\) \\
Planck 定数 & \(h\) & \(6.626 \times 10^{-34}\) & J s \\
換算 Planck 定数 & \(\hbar = h/2\pi\) & \(1.055 \times 10^{-34}\) & J
s \\
Boltzmann 定数 & \(k_B\) & \(1.381 \times 10^{-23}\) & J K\(^{-1}\) \\
電気素量 & \(e\) & \(1.602 \times 10^{-19}\) & C \\
電子質量 & \(m_e\) & \(9.109 \times 10^{-31}\) & kg \\
陽子質量 & \(m_p\) & \(1.673 \times 10^{-27}\) & kg \\
原子質量単位 & \(u\) & \(1.661 \times 10^{-27}\) & kg \\
Avogadro 定数 & \(N_A\) & \(6.022 \times 10^{23}\) & mol\(^{-1}\) \\
真空誘電率 & \(\epsilon_0\) & \(8.854 \times 10^{-12}\) & F
m\(^{-1}\) \\
真空透磁率 & \(\mu_0\) & \(1.257 \times 10^{-6}\) & H m\(^{-1}\) \\
\end{longtable}
}

\subsection{B.1.2
天体物理でよく使う組合せ定数}\label{b.1.2-ux5929ux4f53ux7269ux7406ux3067ux3088ux304fux4f7fux3046ux7d44ux5408ux305bux5b9aux6570}

{\def\LTcaptype{none} % do not increment counter
\begin{longtable}[]{@{}
  >{\raggedright\arraybackslash}p{(\linewidth - 6\tabcolsep) * \real{0.2500}}
  >{\raggedright\arraybackslash}p{(\linewidth - 6\tabcolsep) * \real{0.2500}}
  >{\raggedright\arraybackslash}p{(\linewidth - 6\tabcolsep) * \real{0.2500}}
  >{\raggedright\arraybackslash}p{(\linewidth - 6\tabcolsep) * \real{0.2500}}@{}}
\toprule\noalign{}
\begin{minipage}[b]{\linewidth}\raggedright
定数
\end{minipage} & \begin{minipage}[b]{\linewidth}\raggedright
記号
\end{minipage} & \begin{minipage}[b]{\linewidth}\raggedright
値
\end{minipage} & \begin{minipage}[b]{\linewidth}\raggedright
単位
\end{minipage} \\
\midrule\noalign{}
\endhead
\bottomrule\noalign{}
\endlastfoot
Stefan-Boltzmann 定数 & \(\sigma = \pi^2 k_B^4 / (60 \hbar^3 c^2)\) &
\(5.670 \times 10^{-8}\) & W m\(^{-2}\) K\(^{-4}\) \\
放射定数 & \(a = 4\sigma/c\) & \(7.566 \times 10^{-16}\) & J m\(^{-3}\)
K\(^{-4}\) \\
Wien 変位定数 & \(b = h c / (4.965 k_B)\) & \(2.898 \times 10^{-3}\) & m
K \\
Wien 変位（振動数版） & \(b_\nu = 2.821 k_B/h\) &
\(5.879 \times 10^{10}\) & Hz K\(^{-1}\) \\
微細構造定数 & \(\alpha = e^2/(4\pi\epsilon_0 \hbar c)\) &
\(7.297 \times 10^{-3}\) & (無次元、\(\simeq 1/137\)) \\
Bohr 半径 & \(a_0 = 4\pi\epsilon_0 \hbar^2/(m_e e^2)\) &
\(5.292 \times 10^{-11}\) & m \\
Rydberg エネルギー & \(\text{Ry} = m_e c^2 \alpha^2/2\) & \(13.606\) &
eV \\
Rydberg 定数 & \(R_\infty = \text{Ry}/(hc)\) & \(1.097 \times 10^7\) &
m\(^{-1}\) \\
Compton 波長（電子） & \(\lambda_C = h/(m_e c)\) &
\(2.426 \times 10^{-12}\) & m \\
古典電子半径 & \(r_e = e^2/(4\pi\epsilon_0 m_e c^2)\) &
\(2.818 \times 10^{-15}\) & m \\
Thomson 断面積 & \(\sigma_T = 8\pi r_e^2/3\) & \(6.652 \times 10^{-29}\)
& m\(^2\) \\
\end{longtable}
}

\textbf{本文での使用}：

\begin{itemize}
\tightlist
\item
  \(\sigma\), \(a\)：第2章 §2.4（Stefan-Boltzmann 則の初出）、第4章
  §4.4（半球フラックスでの \(\pi\) 因子）、第17章 §17.1（恒星有効温度）
\item
  \(b\), \(b_\nu\)：第1章、第2章 §2.3（Wien 変位則）
\item
  \(\alpha\)：第22章 §22.5（電気四重極/磁気双極子遷移強度の見積もり）
\item
  \(a_0\), Ry：第22章 §22.1-§22.3（Rydberg 公式と Bohr 模型）
\item
  \(r_e\), \(\sigma_T\)：第18章 §18.3（Compton 散乱）
\end{itemize}

\begin{center}\rule{0.5\linewidth}{0.5pt}\end{center}

\section{B.2 天文単位}\label{sec-astronomical-units}

\subsection{B.2.1 距離・時間}\label{b.2.1-ux8dddux96e2ux6642ux9593}

{\def\LTcaptype{none} % do not increment counter
\begin{longtable}[]{@{}
  >{\raggedright\arraybackslash}p{(\linewidth - 6\tabcolsep) * \real{0.2500}}
  >{\raggedright\arraybackslash}p{(\linewidth - 6\tabcolsep) * \real{0.2500}}
  >{\raggedright\arraybackslash}p{(\linewidth - 6\tabcolsep) * \real{0.2500}}
  >{\raggedright\arraybackslash}p{(\linewidth - 6\tabcolsep) * \real{0.2500}}@{}}
\toprule\noalign{}
\begin{minipage}[b]{\linewidth}\raggedright
単位
\end{minipage} & \begin{minipage}[b]{\linewidth}\raggedright
記号
\end{minipage} & \begin{minipage}[b]{\linewidth}\raggedright
値
\end{minipage} & \begin{minipage}[b]{\linewidth}\raggedright
由来
\end{minipage} \\
\midrule\noalign{}
\endhead
\bottomrule\noalign{}
\endlastfoot
天文単位 & au & \(1.496 \times 10^{11}\) m & 地球-太陽平均距離（IAU
定義値） \\
光年 & ly & \(9.461 \times 10^{15}\) m & 光が真空中で 1 年に進む距離 \\
パーセク & pc & \(3.086 \times 10^{16}\) m & 地球軌道半径が 1
秒角に見える距離 \\
キロパーセク & kpc & \(10^3\) pc & 銀河系スケール \\
メガパーセク & Mpc & \(10^6\) pc & 銀河間距離 \\
ギガパーセク & Gpc & \(10^9\) pc & 観測可能宇宙のスケール \\
ユリウス年 & yr & \(3.156 \times 10^7\) s & \(= 365.25 \cdot 86400\)
s \\
\end{longtable}
}

\subsection{B.2.2
質量・光度・半径}\label{b.2.2-ux8ceaux91cfux5149ux5ea6ux534aux5f84}

太陽値を基準にとる：

{\def\LTcaptype{none} % do not increment counter
\begin{longtable}[]{@{}lll@{}}
\toprule\noalign{}
量 & 記号 & 値 \\
\midrule\noalign{}
\endhead
\bottomrule\noalign{}
\endlastfoot
太陽質量 & \(M_\odot\) & \(1.989 \times 10^{30}\) kg \\
太陽光度 & \(L_\odot\) & \(3.828 \times 10^{26}\) W \\
太陽半径 & \(R_\odot\) & \(6.957 \times 10^8\) m \\
太陽有効温度 & \(T_{\text{eff},\odot}\) & \(5772\) K \\
太陽絶対等級 & \(M_{V,\odot}\) & \(+4.83\) \\
\end{longtable}
}

その他の天体スケール：

{\def\LTcaptype{none} % do not increment counter
\begin{longtable}[]{@{}
  >{\raggedright\arraybackslash}p{(\linewidth - 4\tabcolsep) * \real{0.3333}}
  >{\raggedright\arraybackslash}p{(\linewidth - 4\tabcolsep) * \real{0.3333}}
  >{\raggedright\arraybackslash}p{(\linewidth - 4\tabcolsep) * \real{0.3333}}@{}}
\toprule\noalign{}
\begin{minipage}[b]{\linewidth}\raggedright
天体
\end{minipage} & \begin{minipage}[b]{\linewidth}\raggedright
質量
\end{minipage} & \begin{minipage}[b]{\linewidth}\raggedright
半径
\end{minipage} \\
\midrule\noalign{}
\endhead
\bottomrule\noalign{}
\endlastfoot
木星 & \(M_J \simeq 1.9 \times 10^{27}\) kg \(\simeq 10^{-3} M_\odot\) &
\(7.15 \times 10^7\) m \\
地球 & \(M_\oplus \simeq 5.97 \times 10^{24}\) kg & \(6.37 \times 10^6\)
m \\
白色矮星 & \(\sim 0.6 M_\odot\) & \(\sim 10^7\) m \\
中性子星 & \(\sim 1.4 M_\odot\) & \(\sim 10^4\) m \\
超巨大ブラックホール & \(10^6\)〜\(10^{10} M_\odot\) &
\(r_s = 2GM/c^2\) \\
\end{longtable}
}

\subsection{B.2.3 宇宙論パラメータ（Planck
2018）}\label{b.2.3-ux5b87ux5b99ux8ad6ux30d1ux30e9ux30e1ux30fcux30bfplanck-2018}

{\def\LTcaptype{none} % do not increment counter
\begin{longtable}[]{@{}lll@{}}
\toprule\noalign{}
パラメータ & 記号 & 値 \\
\midrule\noalign{}
\endhead
\bottomrule\noalign{}
\endlastfoot
Hubble 定数 & \(H_0\) & \(67.4\) km s\(^{-1}\) Mpc\(^{-1}\) \\
物質密度比 & \(\Omega_m\) & \(0.315\) \\
暗黒エネルギー密度比 & \(\Omega_\Lambda\) & \(0.685\) \\
バリオン密度比 & \(\Omega_b\) & \(0.049\) \\
CMB 温度 & \(T_\text{CMB}\) & \(2.7255\) K \\
宇宙年齢 & \(t_0\) & \(1.38 \times 10^{10}\) yr \\
\end{longtable}
}

\textbf{本文での使用}：

\begin{itemize}
\tightlist
\item
  \(au, pc, M_\odot\)：第17章（恒星・惑星）、第18章（降着円盤）
\item
  \(L_\odot, R_\odot, T_{\text{eff},\odot}\)：第17章
  §17.1-§17.2（恒星と惑星平衡温度）
\item
  \(T_\text{CMB}\)：第16章（CMB）
\end{itemize}

\begin{center}\rule{0.5\linewidth}{0.5pt}\end{center}

\section{B.3 振動数・波長・エネルギーの相互換算}\label{sec-conversions}

電磁波の同じ「色」を、目的に応じて違う変数で書く。基本関係：

\begin{equation}\protect\phantomsection\label{eq-em-conversions}{
\nu = \frac{c}{\lambda}, \quad E = h\nu = \frac{hc}{\lambda} = k_B T_E
}\end{equation}

ここで \(T_E \equiv E/k_B\)
は「\textbf{特性温度}」（ある光子エネルギーに対応する温度）。

\subsection{B.3.1 換算定数}\label{b.3.1-ux63dbux7b97ux5b9aux6570}

{\def\LTcaptype{none} % do not increment counter
\begin{longtable}[]{@{}
  >{\raggedright\arraybackslash}p{(\linewidth - 4\tabcolsep) * \real{0.3333}}
  >{\raggedright\arraybackslash}p{(\linewidth - 4\tabcolsep) * \real{0.3333}}
  >{\raggedright\arraybackslash}p{(\linewidth - 4\tabcolsep) * \real{0.3333}}@{}}
\toprule\noalign{}
\begin{minipage}[b]{\linewidth}\raggedright
量の換算
\end{minipage} & \begin{minipage}[b]{\linewidth}\raggedright
計算式
\end{minipage} & \begin{minipage}[b]{\linewidth}\raggedright
数値
\end{minipage} \\
\midrule\noalign{}
\endhead
\bottomrule\noalign{}
\endlastfoot
\(hc\) & & \(1240\) eV nm \(= 1.24\) µm eV \\
\(h/k_B\) & & \(4.80 \times 10^{-11}\) K Hz\(^{-1}\) \\
\(k_B\) (温度 → エネルギー) & & \(1\) K \(= 8.617 \times 10^{-5}\) eV \\
\(hc/k_B\) (波長 ↔ 温度) & & \(\lambda \cdot T_E = hc/k_B \simeq 14388\)
µm·K（つまり「波長 {[}µm{]} × 特性温度 {[}K{]}」が一定） \\
1 eV \(\to\) 振動数 & \(E/h\) & \(2.418 \times 10^{14}\) Hz \\
1 eV \(\to\) 波長 & \(hc/E\) & \(1239.8\) nm \\
1 eV \(\to\) 温度 & \(E/k_B\) & \(1.160 \times 10^4\) K \\
\end{longtable}
}

\subsection{B.3.2
電磁波帯域の対応表}\label{b.3.2-ux96fbux78c1ux6ce2ux5e2fux57dfux306eux5bfeux5fdcux8868}

{\def\LTcaptype{none} % do not increment counter
\begin{longtable}[]{@{}
  >{\raggedright\arraybackslash}p{(\linewidth - 10\tabcolsep) * \real{0.1667}}
  >{\raggedright\arraybackslash}p{(\linewidth - 10\tabcolsep) * \real{0.1667}}
  >{\raggedright\arraybackslash}p{(\linewidth - 10\tabcolsep) * \real{0.1667}}
  >{\raggedright\arraybackslash}p{(\linewidth - 10\tabcolsep) * \real{0.1667}}
  >{\raggedright\arraybackslash}p{(\linewidth - 10\tabcolsep) * \real{0.1667}}
  >{\raggedright\arraybackslash}p{(\linewidth - 10\tabcolsep) * \real{0.1667}}@{}}
\toprule\noalign{}
\begin{minipage}[b]{\linewidth}\raggedright
帯域
\end{minipage} & \begin{minipage}[b]{\linewidth}\raggedright
波長
\end{minipage} & \begin{minipage}[b]{\linewidth}\raggedright
振動数
\end{minipage} & \begin{minipage}[b]{\linewidth}\raggedright
エネルギー
\end{minipage} & \begin{minipage}[b]{\linewidth}\raggedright
特性温度
\end{minipage} & \begin{minipage}[b]{\linewidth}\raggedright
主な天体
\end{minipage} \\
\midrule\noalign{}
\endhead
\bottomrule\noalign{}
\endlastfoot
電波（広義） & \(> 1\) mm（〜長波長まで） & \(< 300\) GHz &
\(< 10^{-3}\) eV & \(< 15\) K & 21cm
線、分子線、シンクロトロン（電波天文学全般） \\
―うち マイクロ波 & 1 mm -- 1 m & 300 MHz -- 300 GHz & \(10^{-6}\) --
\(10^{-3}\) eV & \(\sim 1\) K & \textbf{CMB}、低温ダスト \\
―うちサブミリ波 & 0.3 -- 1 mm & 300 GHz -- 1 THz & \(\sim 10^{-3}\) eV &
\(\sim 30\) K & ALMA、原始惑星系円盤 \\
遠赤外 & 30--300 µm & 1--10 THz & \(\sim 0.01\) eV & \(\sim 100\) K &
冷たいダスト、Herschel \\
中間赤外 & 3--30 µm & 10--100 THz & \(\sim 0.1\) eV & \(\sim 1000\) K &
暖かいダスト、JWST \\
近赤外 & 0.7--3 µm & \(\sim 10^{14}\) Hz & \(\sim 1\) eV & \(\sim 3000\)
K & 冷たい恒星、ダスト透過 \\
可視 & 380--750 nm & \(\sim 5 \times 10^{14}\) Hz & \(\sim 2\) eV &
\(\sim 6000\) K & 太陽・恒星 \\
紫外 & 10--380 nm & \(\sim 10^{15}\)--\(10^{16}\) Hz & 3--100 eV &
\(\sim 10^5\) K & OB 型恒星、HII 領域励起 \\
X 線 & 0.01--10 nm & \(10^{16}\)--\(10^{19}\) Hz & 0.1--100 keV &
\(10^6\)--\(10^9\) K & コロナ、降着円盤、中性子星 \\
ガンマ線 & \(< 0.01\) nm & \(> 10^{19}\) Hz & \(> 100\) keV & \(> 10^9\)
K & GRB、パルサー、ブラックホール \\
\end{longtable}
}

\textbf{本文での使用}：第17章
§17.4（赤外線天文学の波長帯）、第18章（X線天体）、第20章（CMB から GRB
まで）

\begin{center}\rule{0.5\linewidth}{0.5pt}\end{center}

\section{B.4 SI ↔ CGS の換算 ― 電磁気学の落とし穴}\label{sec-si-cgs}

電磁気学では、SI（国際単位系）と CGS（ガウス単位系）で
\textbf{式の形そのものが違う}。論文（特に天体物理の古典文献）では CGS
が標準なので、両者の対応を押さえておく。

\begin{tcolorbox}[enhanced jigsaw, arc=.35mm, bottomrule=.15mm, bottomtitle=1mm, breakable, colback=white, colbacktitle=quarto-callout-warning-color!10!white, colframe=quarto-callout-warning-color-frame, coltitle=black, left=2mm, leftrule=.75mm, opacityback=0, opacitybacktitle=0.6, rightrule=.15mm, title=\textcolor{quarto-callout-warning-color}{\faExclamationTriangle}\hspace{0.5em}{警告}, titlerule=0mm, toprule=.15mm, toptitle=1mm]

\textbf{本書の単位系方針（重要）}

本書は \textbf{SI を基本}
とするが、章・節の出自と慣用に応じて以下のように使い分ける：

{\def\LTcaptype{none} % do not increment counter
\begin{longtable}[]{@{}
  >{\raggedright\arraybackslash}p{(\linewidth - 4\tabcolsep) * \real{0.3333}}
  >{\raggedright\arraybackslash}p{(\linewidth - 4\tabcolsep) * \real{0.3333}}
  >{\raggedright\arraybackslash}p{(\linewidth - 4\tabcolsep) * \real{0.3333}}@{}}
\toprule\noalign{}
\begin{minipage}[b]{\linewidth}\raggedright
章
\end{minipage} & \begin{minipage}[b]{\linewidth}\raggedright
単位系
\end{minipage} & \begin{minipage}[b]{\linewidth}\raggedright
理由
\end{minipage} \\
\midrule\noalign{}
\endhead
\bottomrule\noalign{}
\endlastfoot
第1〜10章 & SI & 力学・統計力学・量子論の入門 \\
第11章 & SI（Einstein 係数） & 標準教科書との整合 \\
第13章（電磁気・Maxwell 方程式） & \textbf{SI} に統一（旧稿の Gauss
表記は SI に置換予定） & 学部電磁気の標準 \\
第14章（モード密度） & SI & 第13章との整合 \\
第22章（Bohr 模型・水素原子） & SI（\(4\pi\epsilon_0\) を露わに） &
線吸収係数 \(\alpha_\nu\) の SI 形 \\
第26章（QED） & \textbf{構造のみ提示}、係数は単位系依存と明示 &
見通し重視、Loudon／Sakurai 等の専門書を参照 \\
\end{longtable}
}

各節の冒頭または初出式で、単位系が SI から外れる場合は
\textbf{「(CGS)」「(Gauss)」}
を明記する。論文との比較時にも、必ず単位系を確認すること。

\end{tcolorbox}

\subsection{B.4.1
力学・熱の換算}\label{b.4.1-ux529bux5b66ux71b1ux306eux63dbux7b97}

{\def\LTcaptype{none} % do not increment counter
\begin{longtable}[]{@{}llll@{}}
\toprule\noalign{}
量 & SI & CGS & 換算 \\
\midrule\noalign{}
\endhead
\bottomrule\noalign{}
\endlastfoot
長さ & m & cm & \(1\,\text{m} = 10^2\) cm \\
質量 & kg & g & \(1\,\text{kg} = 10^3\) g \\
時間 & s & s & 同じ \\
力 & N & dyne & \(1\,\text{N} = 10^5\) dyne \\
エネルギー & J & erg & \(1\,\text{J} = 10^7\) erg \\
圧力 & Pa & dyne cm\(^{-2}\) & \(1\,\text{Pa} = 10\) dyne cm\(^{-2}\) \\
\end{longtable}
}

\subsection{B.4.2
電磁気の換算（注意）}\label{b.4.2-ux96fbux78c1ux6c17ux306eux63dbux7b97ux6ce8ux610f}

{\def\LTcaptype{none} % do not increment counter
\begin{longtable}[]{@{}
  >{\raggedright\arraybackslash}p{(\linewidth - 6\tabcolsep) * \real{0.2500}}
  >{\raggedright\arraybackslash}p{(\linewidth - 6\tabcolsep) * \real{0.2500}}
  >{\raggedright\arraybackslash}p{(\linewidth - 6\tabcolsep) * \real{0.2500}}
  >{\raggedright\arraybackslash}p{(\linewidth - 6\tabcolsep) * \real{0.2500}}@{}}
\toprule\noalign{}
\begin{minipage}[b]{\linewidth}\raggedright
量
\end{minipage} & \begin{minipage}[b]{\linewidth}\raggedright
SI
\end{minipage} & \begin{minipage}[b]{\linewidth}\raggedright
CGS (Gaussian)
\end{minipage} & \begin{minipage}[b]{\linewidth}\raggedright
換算
\end{minipage} \\
\midrule\noalign{}
\endhead
\bottomrule\noalign{}
\endlastfoot
電荷 & C & esu & \(1\) C \(= 3 \times 10^9\) esu \\
電場 & V m\(^{-1}\) & statvolt cm\(^{-1}\) & \(1\) V m\(^{-1}\)
\(\simeq 3.34 \times 10^{-5}\) statvolt cm\(^{-1}\) \\
磁場 & T (Tesla) & G (Gauss) & \(1\) T \(= 10^4\) G \\
\end{longtable}
}

\subsection{B.4.3
式の形の違い}\label{b.4.3-ux5f0fux306eux5f62ux306eux9055ux3044}

\textbf{Coulomb 法則}：

\[
\text{SI:}\quad F = \frac{1}{4\pi\epsilon_0}\frac{q_1 q_2}{r^2}, \quad \text{CGS:}\quad F = \frac{q_1 q_2}{r^2}
\]

\textbf{Maxwell 方程式（CGS、第13章 SI 版と比較）}：

\[
\nabla \cdot \mathbf{E} = 4\pi\rho, \quad \nabla \times \mathbf{B} = \frac{4\pi}{c}\mathbf{J} + \frac{1}{c}\frac{\partial \mathbf{E}}{\partial t}
\]

CGS では \(4\pi\) と \(1/c\) が露わに出る。本文の SI 表現では
\(\epsilon_0, \mu_0\) の中に隠れている。

\textbf{Rydberg エネルギー}（一例として）：

\[
\text{SI:}\quad \text{Ry} = \frac{m_e e^4}{2(4\pi\epsilon_0)^2 \hbar^2}, \quad \text{CGS:}\quad \text{Ry} = \frac{m_e e^4}{2\hbar^2}
\]

\subsection{B.4.4 SI ⇔ Gauss
式形対応表（本文の主要式）}\label{sec-si-gauss-form-table}

本書本文で前因子を含む式は、SI と
Gauss（CGS）で次のように対応する。Gauss 形では一般に
\(e^2 \to e^2/(4\pi\epsilon_0)\) の置き換えで SI
から移れる（電磁気の前因子が \(4\pi\epsilon_0\)
を介して隠れる）。論文と比較する際はこの表を基準にすること。

{\def\LTcaptype{none} % do not increment counter
\begin{longtable}[]{@{}
  >{\raggedright\arraybackslash}p{(\linewidth - 6\tabcolsep) * \real{0.2500}}
  >{\raggedright\arraybackslash}p{(\linewidth - 6\tabcolsep) * \real{0.2500}}
  >{\raggedright\arraybackslash}p{(\linewidth - 6\tabcolsep) * \real{0.2500}}
  >{\raggedright\arraybackslash}p{(\linewidth - 6\tabcolsep) * \real{0.2500}}@{}}
\toprule\noalign{}
\begin{minipage}[b]{\linewidth}\raggedright
量・式
\end{minipage} & \begin{minipage}[b]{\linewidth}\raggedright
SI 形
\end{minipage} & \begin{minipage}[b]{\linewidth}\raggedright
Gauss（CGS）形
\end{minipage} & \begin{minipage}[b]{\linewidth}\raggedright
本文
\end{minipage} \\
\midrule\noalign{}
\endhead
\bottomrule\noalign{}
\endlastfoot
Coulomb 力 & \(\dfrac{1}{4\pi\epsilon_0}\dfrac{q_1 q_2}{r^2}\) &
\(\dfrac{q_1 q_2}{r^2}\) & 第13章 \\
微細構造定数 \(\alpha\) & \(\dfrac{e^2}{4\pi\epsilon_0 \hbar c}\) &
\(\dfrac{e^2}{\hbar c}\) & 第22章 \\
Rydberg エネルギー & \(\dfrac{m_e e^4}{2(4\pi\epsilon_0)^2 \hbar^2}\) &
\(\dfrac{m_e e^4}{2\hbar^2}\) & 第22章 \\
古典電子半径 \(r_e\) & \(\dfrac{e^2}{4\pi\epsilon_0 m_e c^2}\) &
\(\dfrac{e^2}{m_e c^2}\) & 第18章 \\
Poynting ベクトル & \(\dfrac{1}{\mu_0}\mathbf{E}\times\mathbf{B}\) &
\(\dfrac{c}{4\pi}\mathbf{E}\times\mathbf{B}\) & 第13章 \\
線吸収の積分強度（Lorentz 振動子前因子） &
\(\dfrac{\pi e^2}{4\pi\epsilon_0\, m_e c}=\dfrac{e^2}{4\epsilon_0 m_e c}\)
& \(\dfrac{\pi e^2}{m_e c}\) & 第22章・水素線地図 \\
自発放出率 \(A_{ul}\) &
\(\dfrac{\omega^3}{3\pi\epsilon_0 \hbar c^3}\,|\mathbf{d}_{ul}|^2\) &
\(\dfrac{4\omega^3}{3\hbar c^3}\,|\mathbf{d}_{ul}|^2\) & 第26章 \\
\end{longtable}
}

\begin{tcolorbox}[enhanced jigsaw, arc=.35mm, bottomrule=.15mm, breakable, colback=white, colframe=quarto-callout-note-color-frame, left=2mm, leftrule=.75mm, opacityback=0, rightrule=.15mm, toprule=.15mm]
\begin{minipage}[t]{5.5mm}
\textcolor{quarto-callout-note-color}{\faInfo}
\end{minipage}%
\begin{minipage}[t]{\textwidth - 5.5mm}

\textbf{読み方}：天体物理の慣用式（特に線吸収係数の前因子
\(\pi e^2/m_e c\)、第22章・水素線地図）は Gauss
形である。本書本文ではこれを SI 形 \(e^2/(4\epsilon_0 m_e c)\)
で記し、慣用の Gauss 形を併記する方針をとる。第26章の QED
の主式（式~\ref{eq-A-coefficient-QED} など）は SI
形で厳密に記し、Gauss/HL 形との差はこの表で確認できる。

\end{minipage}%
\end{tcolorbox}

\begin{tcolorbox}[enhanced jigsaw, arc=.35mm, bottomrule=.15mm, breakable, colback=white, colframe=quarto-callout-warning-color-frame, left=2mm, leftrule=.75mm, opacityback=0, rightrule=.15mm, toprule=.15mm]
\begin{minipage}[t]{5.5mm}
\textcolor{quarto-callout-warning-color}{\faExclamationTriangle}
\end{minipage}%
\begin{minipage}[t]{\textwidth - 5.5mm}

\textbf{注意}：天文学の論文では断りなく CGS
が使われていることが多い。数値を比較するときは \textbf{必ず単位系を確認}
すること。特に電磁気が絡む量（断面積、磁場、Stefan-Boltzmann
定数、\ldots）は要注意である。

\end{minipage}%
\end{tcolorbox}

\begin{center}\rule{0.5\linewidth}{0.5pt}\end{center}

\section{B.5 数値の桁感覚 ―
オーダー推定のための定数}\label{sec-order-of-magnitude}

「ざっくり計算」のための桁感覚。

\subsection{B.5.1
エネルギー目盛り}\label{b.5.1-ux30a8ux30cdux30ebux30aeux30fcux76eeux76dbux308a}

{\def\LTcaptype{none} % do not increment counter
\begin{longtable}[]{@{}
  >{\raggedright\arraybackslash}p{(\linewidth - 6\tabcolsep) * \real{0.2500}}
  >{\raggedright\arraybackslash}p{(\linewidth - 6\tabcolsep) * \real{0.2500}}
  >{\raggedright\arraybackslash}p{(\linewidth - 6\tabcolsep) * \real{0.2500}}
  >{\raggedright\arraybackslash}p{(\linewidth - 6\tabcolsep) * \real{0.2500}}@{}}
\toprule\noalign{}
\begin{minipage}[b]{\linewidth}\raggedright
現象
\end{minipage} & \begin{minipage}[b]{\linewidth}\raggedright
エネルギー
\end{minipage} & \begin{minipage}[b]{\linewidth}\raggedright
特性温度
\end{minipage} & \begin{minipage}[b]{\linewidth}\raggedright
対応波長
\end{minipage} \\
\midrule\noalign{}
\endhead
\bottomrule\noalign{}
\endlastfoot
室温の熱運動 & \(0.025\) eV & 300 K & 50 µm \\
水素 21 cm 超微細構造 & \(5.9 \times 10^{-6}\) eV & \(0.07\) K & 21
cm \\
CMB 光子（典型） & \(6.6 \times 10^{-4}\) eV & 2.725 K & 1 mm \\
可視光（H\(\alpha\)） & \(1.89\) eV & \(2.2 \times 10^4\) K & 656 nm \\
水素 Lyα & \(10.2\) eV & \(1.2 \times 10^5\) K & 121.6 nm \\
水素電離 & \(13.6\) eV & \(1.6 \times 10^5\) K & 91.2 nm \\
鉄 K\(\alpha\)（中性鉄） & \(6.4\) keV & \(7.4 \times 10^7\) K & 0.19
nm \\
電子静止質量 & \(0.511\) MeV & \(5.9 \times 10^9\) K &
\(\lambda_C = 2.4\) pm \\
陽子静止質量 & \(938\) MeV & \(1.1 \times 10^{13}\) K & \(\sim 1.3\)
fm \\
\end{longtable}
}

\subsection{B.5.2
重力スケール}\label{b.5.2-ux91cdux529bux30b9ux30b1ux30fcux30eb}

{\def\LTcaptype{none} % do not increment counter
\begin{longtable}[]{@{}
  >{\raggedright\arraybackslash}p{(\linewidth - 4\tabcolsep) * \real{0.3333}}
  >{\raggedright\arraybackslash}p{(\linewidth - 4\tabcolsep) * \real{0.3333}}
  >{\raggedright\arraybackslash}p{(\linewidth - 4\tabcolsep) * \real{0.3333}}@{}}
\toprule\noalign{}
\begin{minipage}[b]{\linewidth}\raggedright
量
\end{minipage} & \begin{minipage}[b]{\linewidth}\raggedright
計算
\end{minipage} & \begin{minipage}[b]{\linewidth}\raggedright
例
\end{minipage} \\
\midrule\noalign{}
\endhead
\bottomrule\noalign{}
\endlastfoot
Schwarzschild 半径 & \(r_s = 2GM/c^2\) & \(r_{s,\odot} \simeq 3\) km,
\(r_{s,M87} \simeq 2 \times 10^{13}\) m \\
表面重力（恒星） & \(g = GM/R^2\) & 太陽 \(g_\odot = 274\) m s\(^{-2}\),
白色矮星 \(\sim 10^6\) m s\(^{-2}\) \\
自由落下時間 & \(t_\text{ff} = \sqrt{1/G\rho}\) & 太陽密度 \(1.4\) g/cc
→ \(t_\text{ff} \sim 30\) min \\
\end{longtable}
}

\subsection{B.5.3
観測的密度}\label{b.5.3-ux89b3ux6e2cux7684ux5bc6ux5ea6}

{\def\LTcaptype{none} % do not increment counter
\begin{longtable}[]{@{}lll@{}}
\toprule\noalign{}
環境 & 数密度 {[}cm\(^{-3}\){]} & 温度 {[}K{]} \\
\midrule\noalign{}
\endhead
\bottomrule\noalign{}
\endlastfoot
銀河間媒質 & \(10^{-6}\) & \(10^5\)--\(10^7\) \\
暖かい星間ガス（WIM） & \(10^{-1}\) & \(10^4\) \\
HII 領域 & \(10^2\)--\(10^4\) & \(10^4\) \\
暗黒分子雲 & \(10^3\)--\(10^5\) & \(10\)--\(20\) \\
原始惑星系円盤 & \(10^{10}\)--\(10^{12}\) & \(20\)--\(1000\) \\
恒星光球 & \(\sim 10^{17}\) & \(\sim 6000\) \\
太陽コロナ & \(10^8\)--\(10^{10}\) & \(10^6\)--\(10^7\) \\
\end{longtable}
}

\textbf{本文での使用}：第15章以降の応用章で「桁感覚」が必要なとき、本表を参照してほしい。

\begin{center}\rule{0.5\linewidth}{0.5pt}\end{center}

\section*{B.6
さらに進むための参考文献・データソース}\label{sec-data-references}
\addcontentsline{toc}{section}{B.6
さらに進むための参考文献・データソース}

\markright{B.6 さらに進むための参考文献・データソース}

\begin{itemize}
\tightlist
\item
  CODATA 2018 fundamental physical
  constants（\url{https://physics.nist.gov/cuu/Constants}）
\item
  IAU 2015 nominal values for solar and planetary properties（IAU
  Resolution B3）
\item
  NIST Atomic Spectra Database（\url{https://physics.nist.gov/asd}）
\item
  NIST Molecular Spectroscopy Database
\item
  \emph{Allen's Astrophysical Quantities} (4th ed., Springer, 2000) ---
  天文定数の定番一覧
\item
  \emph{Particle Data Group}, \emph{Review of Particle Physics} ---
  素粒子・宇宙論定数（\url{https://pdg.lbl.gov}）
\item
  Planck Collaboration 2018 results --- 宇宙論パラメータ
\end{itemize}

\chapter{用語索引（和英・英和）}\label{sec-glossary}

本書で用いた主要用語の和英・英和対照表。天文学に固有の用語は、日本天文学会『\href{https://astro-dic.jp/}{天文学辞典}』（astro-dic.jp）の標準表記に準拠した。物理用語のうち、揺れやすいものは本書全体で次の訳を統一して用いている：\textbf{frequency＝振動数}（×周波数）、\textbf{neutrino＝ニュートリノ}、\textbf{mid-infrared＝中間赤外}。

英語見出しの末尾に付した E1 / M1 / E2 などは多重極遷移の慣用記号、(LTE)
/ (CMB) などは本文で用いた略号である。

なお放射の「出る側」の係数については、\textbf{放出係数
\(j_\nu\)（emission coefficient）は単位立体角あたり}
の量、\textbf{体積放射率 \(\varepsilon_\nu\)（volume
emissivity）は全立体角 \(4\pi\)
について積分済み}（\(\varepsilon_\nu = 4\pi j_\nu\)、等方放射の場合）の量として区別する。両者を「放射率」と一括りにしない。

本索引の収録語数：157 語（ほかに別表記からの参照項目を含む）。

\section{C.1
和英索引（五十音順）}\label{c.1-ux548cux82f1ux7d22ux5f15ux4e94ux5341ux97f3ux9806}

{\def\LTcaptype{none} % do not increment counter
\begin{longtable}[]{@{}ll@{}}
\toprule\noalign{}
日本語 & English \\
\midrule\noalign{}
\endhead
\bottomrule\noalign{}
\endlastfoot
Einstein 係数 & Einstein coefficients \\
Einstein の関係式 & Einstein relations \\
色温度 & color temperature \\
Wien の変位則 & Wien's displacement law \\
宇宙マイクロ波背景放射 & cosmic microwave background (CMB) \\
エネルギー密度 & energy density \\
エネルギー量子 & energy quantum \\
エントロピー & entropy \\
円偏光 & circular polarization \\
灰色大気 & grey atmosphere \\
化学ポテンシャル & chemical potential \\
角振動数 & angular frequency \\
Casimir 効果 & Casimir effect \\
完全熱平衡 & complete thermodynamic equilibrium \\
観測量 & observable \\
Gauss 型 & Gaussian profile \\
輝線 & emission line \\
輝度温度 & brightness temperature \\
吸収 & absorption \\
吸収係数 & absorption coefficient \\
吸収線 & absorption line \\
境界条件 & boundary conditions \\
局所熱力学平衡 & local thermodynamic equilibrium (LTE) \\
巨視的乱流 & macroturbulence \\
Kirchhoff の法則 & Kirchhoff's law \\
禁制線 & forbidden line \\
逆コンプトン散乱 & inverse Compton scattering \\
空洞放射 & cavity radiation \\
空洞モード & cavity mode \\
減衰翼 & damping wing \\
源泉関数 → ソース関数 & （→ source function を見よ） \\
光学的厚さ & optical depth \\
光学的に厚い & optically thick \\
光学的に薄い & optically thin \\
項記号 & term symbol \\
光子 & photon \\
光子気体 & photon gas \\
恒星風 & stellar wind \\
降着円盤 & accretion disk \\
光電効果 & photoelectric effect \\
高分散分光 & high-resolution spectroscopy \\
光量子 & light quantum \\
黒体 & black body \\
黒体放射 & blackbody radiation \\
固有状態 & eigenstate \\
固有値 & eigenvalue \\
再結合 & recombination \\
最大エントロピー原理 & principle of maximum entropy \\
再電離 & reionization \\
Saha の式 & Saha equation \\
散乱 & scattering \\
散乱係数 & scattering coefficient \\
紫外線破綻 & ultraviolet catastrophe \\
自然幅 & natural width \\
修正黒体 & modified blackbody \\
縮退度 & degeneracy \\
Stefan--Boltzmann 定数 & Stefan--Boltzmann constant \\
主量子数 & principal quantum number \\
詳細釣り合い & detailed balance \\
衝突脱励起 & collisional de-excitation \\
真空状態 & vacuum state \\
真空のゆらぎ & vacuum fluctuations \\
シンクロトロン放射 & synchrotron radiation \\
振動子強度 & oscillator strength \\
振動数 & frequency \\
磁気双極子遷移 & magnetic dipole transition (M1) \\
自発放出 & spontaneous emission \\
数演算子 & number operator \\
数状態 & number state \\
スペクトル型 & spectral type \\
スペクトル線 & spectral line \\
スペクトル歪み & spectral distortion \\
生成消滅演算子 & creation and annihilation operators \\
成長曲線 & curve of growth \\
制動放射 & bremsstrahlung \\
赤方偏移 & redshift \\
遷移 & transition \\
選択則 & selection rule \\
占有因子 & occupation factor \\
占有数 & occupation number \\
占有数反転 & population inversion \\
線輪郭 & line profile \\
ソース関数 & source function \\
体積放射率 & volume emissivity \\
多温度円盤 & multi-temperature disk \\
ダスト & dust \\
中性水素 & neutral hydrogen \\
超微細構造 & hyperfine structure \\
調和振動子 & harmonic oscillator \\
直線偏光 & linear polarization \\
定在波 & standing wave \\
電気四重極遷移 & electric quadrupole transition (E2) \\
電気双極子遷移 & electric dipole transition (E1) \\
等価幅 & equivalent width \\
等分配定理 & equipartition theorem \\
ドップラー効果 & Doppler effect \\
ドップラー幅 & Doppler width \\
21 cm 線 & 21 cm line \\
ニュートリノ & neutrino \\
熱化 & thermalization \\
熱化深さ & thermalization depth \\
熱的制動放射 & thermal bremsstrahlung \\
熱的放射 & thermal radiation \\
熱平衡 & thermal equilibrium \\
熱力学的平衡 & thermodynamic equilibrium \\
波数空間 & k-space \\
波数ベクトル & wave vector \\
波長 & wavelength \\
波動方程式 & wave equation \\
場の量子化 & field quantization \\
比強度 & specific intensity \\
非熱的放射 & non-thermal radiation \\
標準降着円盤 & standard accretion disk \\
微細構造 & fine structure \\
微細構造定数 & fine-structure constant \\
微視的乱流 & microturbulence \\
Fermi の黄金律 & Fermi's golden rule \\
Voigt プロファイル & Voigt profile \\
不透明度 & opacity \\
フラックス & flux \\
分散関係 & dispersion relation \\
分配関数 & partition function \\
プランク関数 & Planck function \\
Planck 定数 & Planck constant \\
プランク分布 & Planck distribution \\
平均自由行程 & mean free path \\
平面波 & plane wave \\
偏光 & polarization \\
べき則 & power law \\
放射強度 & radiation intensity \\
放射効率 & emissivity \\
放射輸送方程式 & radiative transfer equation \\
放出 & emission \\
放出係数 & emission coefficient \\
Boltzmann 因子 & Boltzmann factor \\
Boltzmann 分布 & Boltzmann distribution \\
Bose--Einstein 分布 & Bose--Einstein distribution \\
Bose 増強 & Bose enhancement \\
Poynting ベクトル & Poynting vector \\
Maxwell 方程式 & Maxwell's equations \\
Maxwell--Boltzmann 統計 & Maxwell--Boltzmann statistics \\
マクロ状態 & macrostate \\
ミクロ状態 & microstate \\
モード密度 & mode density \\
有効温度 & effective temperature \\
誘導吸収 & induced absorption \\
誘導放出 & stimulated emission \\
横波 & transverse wave \\
立体角 & solid angle \\
リュードベリ定数 & Rydberg constant \\
量子化 & quantization \\
量子数 & quantum number \\
臨界密度 & critical density \\
励起温度 & excitation temperature \\
零点エネルギー & zero-point energy \\
Rayleigh--Jeans 則 & Rayleigh--Jeans law \\
連続スペクトル & continuous spectrum \\
Lorentz 型 & Lorentzian profile \\
\end{longtable}
}

\section{C.2
英和索引（アルファベット順）}\label{c.2-ux82f1ux548cux7d22ux5f15ux30a2ux30ebux30d5ux30a1ux30d9ux30c3ux30c8ux9806}

{\def\LTcaptype{none} % do not increment counter
\begin{longtable}[]{@{}ll@{}}
\toprule\noalign{}
English & 日本語 \\
\midrule\noalign{}
\endhead
\bottomrule\noalign{}
\endlastfoot
absorption & 吸収 \\
absorption coefficient & 吸収係数 \\
absorption line & 吸収線 \\
accretion disk & 降着円盤 \\
angular frequency & 角振動数 \\
black body & 黒体 \\
blackbody radiation & 黒体放射 \\
Boltzmann distribution & Boltzmann 分布 \\
Boltzmann factor & Boltzmann 因子 \\
Bose enhancement & Bose 増強 \\
Bose--Einstein distribution & Bose--Einstein 分布 \\
boundary conditions & 境界条件 \\
bremsstrahlung & 制動放射 \\
brightness temperature & 輝度温度 \\
Casimir effect & Casimir 効果 \\
cavity mode & 空洞モード \\
cavity radiation & 空洞放射 \\
chemical potential & 化学ポテンシャル \\
circular polarization & 円偏光 \\
21 cm line & 21 cm 線 \\
collisional de-excitation & 衝突脱励起 \\
color temperature & 色温度 \\
complete thermodynamic equilibrium & 完全熱平衡 \\
continuous spectrum & 連続スペクトル \\
cosmic microwave background (CMB) & 宇宙マイクロ波背景放射 \\
creation and annihilation operators & 生成消滅演算子 \\
critical density & 臨界密度 \\
curve of growth & 成長曲線 \\
damping wing & 減衰翼 \\
degeneracy & 縮退度 \\
detailed balance & 詳細釣り合い \\
dispersion relation & 分散関係 \\
Doppler effect & ドップラー効果 \\
Doppler width & ドップラー幅 \\
dust & ダスト \\
effective temperature & 有効温度 \\
eigenstate & 固有状態 \\
eigenvalue & 固有値 \\
Einstein coefficients & Einstein 係数 \\
Einstein relations & Einstein の関係式 \\
electric dipole transition (E1) & 電気双極子遷移 \\
electric quadrupole transition (E2) & 電気四重極遷移 \\
emission & 放出 \\
emission coefficient & 放出係数 \\
emission line & 輝線 \\
emissivity & 放射効率 \\
energy density & エネルギー密度 \\
energy quantum & エネルギー量子 \\
entropy & エントロピー \\
equipartition theorem & 等分配定理 \\
equivalent width & 等価幅 \\
excitation temperature & 励起温度 \\
Fermi's golden rule & Fermi の黄金律 \\
field quantization & 場の量子化 \\
fine structure & 微細構造 \\
fine-structure constant & 微細構造定数 \\
flux & フラックス \\
forbidden line & 禁制線 \\
frequency & 振動数 \\
Gaussian profile & Gauss 型 \\
grey atmosphere & 灰色大気 \\
harmonic oscillator & 調和振動子 \\
high-resolution spectroscopy & 高分散分光 \\
hyperfine structure & 超微細構造 \\
induced absorption & 誘導吸収 \\
inverse Compton scattering & 逆コンプトン散乱 \\
k-space & 波数空間 \\
Kirchhoff's law & Kirchhoff の法則 \\
light quantum & 光量子 \\
line profile & 線輪郭 \\
linear polarization & 直線偏光 \\
local thermodynamic equilibrium (LTE) & 局所熱力学平衡 \\
Lorentzian profile & Lorentz 型 \\
macrostate & マクロ状態 \\
macroturbulence & 巨視的乱流 \\
magnetic dipole transition (M1) & 磁気双極子遷移 \\
Maxwell's equations & Maxwell 方程式 \\
Maxwell--Boltzmann statistics & Maxwell--Boltzmann 統計 \\
mean free path & 平均自由行程 \\
microstate & ミクロ状態 \\
microturbulence & 微視的乱流 \\
mode density & モード密度 \\
modified blackbody & 修正黒体 \\
multi-temperature disk & 多温度円盤 \\
natural width & 自然幅 \\
neutral hydrogen & 中性水素 \\
neutrino & ニュートリノ \\
non-thermal radiation & 非熱的放射 \\
number operator & 数演算子 \\
number state & 数状態 \\
observable & 観測量 \\
occupation factor & 占有因子 \\
occupation number & 占有数 \\
opacity & 不透明度 \\
optical depth & 光学的厚さ \\
optically thick & 光学的に厚い \\
optically thin & 光学的に薄い \\
oscillator strength & 振動子強度 \\
partition function & 分配関数 \\
photoelectric effect & 光電効果 \\
photon & 光子 \\
photon gas & 光子気体 \\
Planck constant & Planck 定数 \\
Planck distribution & プランク分布 \\
Planck function & プランク関数 \\
plane wave & 平面波 \\
polarization & 偏光 \\
population inversion & 占有数反転 \\
power law & べき則 \\
Poynting vector & Poynting ベクトル \\
principal quantum number & 主量子数 \\
principle of maximum entropy & 最大エントロピー原理 \\
quantization & 量子化 \\
quantum number & 量子数 \\
radiation intensity & 放射強度 \\
radiative transfer equation & 放射輸送方程式 \\
Rayleigh--Jeans law & Rayleigh--Jeans 則 \\
recombination & 再結合 \\
redshift & 赤方偏移 \\
reionization & 再電離 \\
Rydberg constant & リュードベリ定数 \\
Saha equation & Saha の式 \\
scattering & 散乱 \\
scattering coefficient & 散乱係数 \\
selection rule & 選択則 \\
solid angle & 立体角 \\
source function & ソース関数 \\
specific intensity & 比強度 \\
spectral distortion & スペクトル歪み \\
spectral line & スペクトル線 \\
spectral type & スペクトル型 \\
spontaneous emission & 自発放出 \\
standard accretion disk & 標準降着円盤 \\
standing wave & 定在波 \\
Stefan--Boltzmann constant & Stefan--Boltzmann 定数 \\
stellar wind & 恒星風 \\
stimulated emission & 誘導放出 \\
synchrotron radiation & シンクロトロン放射 \\
term symbol & 項記号 \\
thermal bremsstrahlung & 熱的制動放射 \\
thermal equilibrium & 熱平衡 \\
thermal radiation & 熱的放射 \\
thermalization & 熱化 \\
thermalization depth & 熱化深さ \\
thermodynamic equilibrium & 熱力学的平衡 \\
transition & 遷移 \\
transverse wave & 横波 \\
ultraviolet catastrophe & 紫外線破綻 \\
vacuum fluctuations & 真空のゆらぎ \\
vacuum state & 真空状態 \\
Voigt profile & Voigt プロファイル \\
volume emissivity & 体積放射率 \\
wave equation & 波動方程式 \\
wave vector & 波数ベクトル \\
wavelength & 波長 \\
Wien's displacement law & Wien の変位則 \\
zero-point energy & 零点エネルギー \\
\end{longtable}
}




\end{document}
